% ============================================================
%  main.tex —— 《正交宇宙物理学：Greg Egan〈Orthogonal〉设定导论》
%  编译： xelatex main.tex  (连续两次，以生成目录与交叉引用)
% ============================================================
% ============================================================
%  preamble.tex —— 《正交宇宙物理学》共用导言区
%  编译： xelatex -> xelatex  (XeLaTeX + ctexbook)
% ============================================================
\documentclass[11pt, UTF8, scheme=chinese, fontset=windows]{ctexbook}

\usepackage{amsmath, amssymb, amsfonts, mathrsfs, bm}
\usepackage{graphicx}
\usepackage{geometry}
\geometry{a4paper, top=25mm, bottom=25mm, left=28mm, right=28mm}
\usepackage{booktabs}
\usepackage{array}
\usepackage{multirow}
\usepackage{longtable}
\usepackage{enumitem}
\usepackage{xcolor}
\usepackage{caption}
\usepackage[most]{tcolorbox}
\usepackage{fancyhdr}
\usepackage{hyperref}

\hypersetup{
  colorlinks=true,
  linkcolor=blue!50!black,
  urlcolor=teal!60!black,
  citecolor=green!45!black,
  bookmarksnumbered=true,
  pdftitle={正交宇宙物理学},
  pdfauthor={编译自 Greg Egan《Orthogonal》官方设定}
}

% ---- 图表 ----
\captionsetup{font=small, labelfont=bf, labelsep=quad}
\graphicspath{{images/}}

% 插图命令：\orthofig[宽度]{图文件名}{中文图注}
\newcommand{\orthofig}[3][0.62\textwidth]{%
  \begin{figure}[htbp]
    \centering
    \includegraphics[width=#1]{#2}
    \caption{#3}
  \end{figure}}

% 不浮动的小图（用于正文中就地的示意图）
\newcommand{\inlinefig}[3][0.5\textwidth]{%
  \begin{center}
    \includegraphics[width=#1]{#2}\\[2pt]
    {\small #3}
  \end{center}}

% ---- 自定义盒子 ----
\newtcolorbox{keybox}[1][]{
  colback=blue!4!white, colframe=blue!45!black,
  boxrule=0.6pt, arc=2pt, left=6pt, right=6pt, top=5pt, bottom=5pt,
  fonttitle=\bfseries, title=#1}

\newtcolorbox{notebox}[1][]{
  colback=orange!5!white, colframe=orange!65!black,
  boxrule=0.6pt, arc=2pt, left=6pt, right=6pt, top=5pt, bottom=5pt,
  fonttitle=\bfseries, title=#1}

% ---- 常用记号 ----
\newcommand{\Riem}{\mathrm{R}}
\newcommand{\Lor}{\mathrm{L}}
\newcommand{\dd}{\mathrm{d}}
\newcommand{\ii}{\mathrm{i}}
\newcommand{\ee}{\mathrm{e}}
\newcommand{\wms}{\omega_m}
\newcommand{\bA}{\mathbf{A}}
\newcommand{\bE}{\mathbf{E}}
\newcommand{\bB}{\mathbf{B}}
\newcommand{\bF}{\mathbf{F}}
\newcommand{\bj}{\mathbf{j}}
\newcommand{\bk}{\mathbf{k}}
\newcommand{\bp}{\mathbf{p}}
\newcommand{\bv}{\mathbf{v}}
\newcommand{\bu}{\mathbf{u}}
\newcommand{\bw}{\mathbf{w}}
\newcommand{\bx}{\mathbf{x}}
\newcommand{\be}{\mathbf{e}}
\newcommand{\bP}{\mathbf{P}}
\newcommand{\bq}{\mathbf{q}}
\newcommand{\bzero}{\mathbf{0}}
\newcommand{\half}{\tfrac{1}{2}}
\DeclareMathOperator{\Tr}{Tr}
\DeclareMathOperator{\diag}{diag}
\DeclareMathOperator{\sgn}{sgn}
\DeclareMathOperator{\Span}{span}
\DeclareMathOperator{\rank}{rank}

% ---- 页眉页脚 ----
\pagestyle{fancy}
\fancyhf{}
\fancyhead[LE,RO]{\small\thepage}
\fancyhead[LO]{\small\nouppercase{\rightmark}}
\fancyhead[RE]{\small\nouppercase{\leftmark}}
\renewcommand{\headrulewidth}{0.4pt}


\begin{document}

\frontmatter
\title{\Huge 正交宇宙物理学\\[6pt]
       \Large ——\ Greg Egan《Orthogonal》科幻设定导论\\(面向物理学研究生的硬核注记)}
\author{编译、翻译与整理自 Greg Egan 官方网站 \texttt{gregegan.net/ORTHOGONAL}}
\date{\today}
\maketitle

\chapter*{编译说明}
\addcontentsline{toc}{chapter}{编译说明}
本书是对 Greg Egan 为其长篇科幻三部曲 \emph{Orthogonal}（《正交》）所撰写的
官方物理设定注释（原文 8 万余词、逾百幅插图）的中文编译整理。
原著包含三部作品：\emph{The Clockwork Rocket}（《发条火箭》）、
\emph{The Eternal Flame}（《永恒之火》）与 \emph{The Arrows of Time}（《时间之箭》）。

按原作作者说明，正文各章所需数学不超过高中水平；标注为“\textbf{进阶}”的内容
（对应官网带 ``[Extra]'' 的页面）则假定读者具备本科高年级的物理与数学训练。
本书面向物理学研究生，故完整保留了进阶部分，并在必要处补足推导背景。

全书数学约定：采用自然单位制，取时空坐标 $(x,y,z,t)$，四维矢量的分量记作
$v^x,v^y,v^z,v^t$（$x,y,z,t$ 为标号而非幂次）。除特别说明外，黎曼（Riemannian）
度规取全正号 $(+,+,+,+)$，洛伦兹（Lorentzian，即我们宇宙的）度规取 $(-,+,+,+)$
或 $(+,-,-,-)$，文中给出公式时会明确标注。符号 $\omega_m$ 表示与粒子静止质量对应的
“质量频率”，即 $\omega_m = mc^2/\hbar$ 在自然单位下的形式。

本书所有插图均取自 Greg Egan 官方网站，版权归原作者所有，仅作学习与研究之用。

正文之外另设三份附录，供读者按需查阅：附录~A 汇列全书所用的数学工具（度规约定、正交群与
旋量、傅里叶方法、特殊函数与变分法）；附录~B 补齐正文中若干处跳步的关键推导；附录~C 以对照表
形式总览正交宇宙与我们所处宇宙在物理层面的系统性差异。

\tableofcontents
\listoffigures

\mainmatter

\chapter{导论：正交宇宙}

Greg Egan 的科幻三部曲《正交》（\emph{Orthogonal}）设定在一个时空几何与我们截然不同的宇宙中：那里的度规（metric）是黎曼型（Riemannian）的 $(+,+,+,+)$，四个维度地位完全平等，没有任何一个方向被单独挑出来充当"时间"。本章阐明这一改动的物理内涵，交代三部曲的剧情背景与本书的来龙去脉，并给出后续各章的路线图。本节及后续内容所依据的物理结论，全部来自 Egan 为小说撰写的官方物理注记（physics notes）；本书即是对这批注记的编译与整理。

\orthofig[0.55\textwidth]{images/Orthogonal200.png}{\emph{Orthogonal} 宇宙的设定：时空四维地位完全平等}

\section{一个符号之差}

我们所处的宇宙是洛伦兹型（Lorentzian）的：取光速 $c=1$，平直时空中的"勾股定理"（时空线元）写作

\begin{equation}
\dd s^2 = -\dd t^2 + \dd x^2 + \dd y^2 + \dd z^2 ,
\tag{0.1}
\end{equation}

其中时间项带着一个负号。这个负号绝非无关紧要的细节，它决定了全部因果结构：正是它把四维时空劈成类时（timelike）、类光（lightlike）与类空（spacelike）三类方向，划出光锥，赋予时间以特殊的地位。

Egan 的整个设定只改动了这一处：把这个减号换成加号，得到黎曼型（Riemannian）度规

\begin{equation}
\dd s^2 = \dd t^2 + \dd x^2 + \dd y^2 + \dd z^2 .
\tag{0.2}
\end{equation}

四维以此完全对称，不存在任何被视为"时间"的特殊方向，时空的度规处处正定，恰如一个普通的四维欧几里得空间——这一设定因此得名"正交"（orthogonal）。这一符号的翻转带来若干极具反直觉、却又自洽的后果：

\begin{itemize}
  \item \textbf{时空图可以如实画在纸上。}在洛伦兹时空中，光锥把时空图割裂开来，"同时"的截面随观察者而倾斜，凭平面几何的直觉去丈量常常出错；而在黎曼时空中，四维地位平等、度规处处正定，时空图就是一张货真价实的欧几里得图形，尺规直觉不再失效。
  \item \textbf{高速远行，对旅行者意味着更长而非更短的时间。}在洛伦兹时空中，运动时钟因时间膨胀而走得慢，远行者比留在家乡的人更年轻。黎曼时空中的情形恰好相反：一艘高速飞船的船上旅程会持续许多代人的时间，而母星上只过去了短短几年——旅行者经历的是更长的岁月。
  \item \textbf{不存在普适光速，也不存在速度上限。}由于没有特殊的时间方向，也就不存在一个所有观察者都认同的不变速度 $c$。光的传播速度不再唯一，它与光的波长（乃至产生方式）有关；相应地，也就不存在"任何速度都不能超过光速"这一上限。
  \item \textbf{光的产生会释放能量。}在我们的宇宙里，原子发光要付出能量、植物靠吸收光来获得养分；在正交宇宙中，发光过程本身是释放能量的——植物向漆黑的夜空发出自己的光，借此制造养分。
\end{itemize}

\orthofig{images/02__007.png}{恒定加速度下的世界线（World Lines of Constant Acceleration）}

这些差异会一路渗透到运动学、波、电磁学、热力学、引力乃至量子力学的每一个角落，正是本书后续各章要逐一展开的内容。

\section{三部曲与本书的由来}

《正交》由三卷构成，依次为《发条火箭》（\emph{The Clockwork Rocket}）、《永恒之火》（\emph{The Eternal Flame}）与《时间之箭》（\emph{The Arrows of Time}）。故事的主人公雅尔达（Yalda）生活在这样一个宇宙里。

开篇时雅尔达不到三岁。她把孱弱的外公达里奥（Dario）带进林中静养——达里奥连日萎靡不振，不肯离开一家人安睡的花圃。村中的莉维亚医生（Doctor Livia）诊断他患了严重的"光照不足"：农场的作物几乎是单色的，人体需要更宽的光谱；太阳光过于偏蓝、太快而让身体无法获取，田间的光又全是缓慢的红光，老人需要的恰是介于两极之间的赭黄、藤黄、藏红与金黄、翠玉与铜绿。

在这个宇宙中，\textbf{光没有普适的速度，光的产生会释放能量}；而\textbf{植物通过向黑暗的夜空发出自己的光来制造食物}。

雅尔达幼年时目睹了一连串奇异流星中的一颗，它们以空前巨大的速度闯入行星系统，被称为"伤人者"（Hurtlers）。事态很快明朗：她的世界危在旦夕，而应对"伤人者"所需的知识与技术，远远超出她所在文明已有的水平。看来唯一可行的出路是：派出一艘飞船以足够高的速度远航。这趟航行对船上的人而言要持续许多代人，可对故乡而言，仅仅过去几年——船员们将有机会发现母星急需的科学，并及时把它带回来，在灾难降临之前扭转乾坤。

《正交》讲的正是雅尔达及其后代在漫长使命的险途中求存、并努力为自己开辟有意义的生活，而毁灭的阴影始终笼罩着他们远在身后的世界的故事。全书分为三卷：

\begin{itemize}
  \item \textbf{第一卷：《发条火箭》}（\emph{The Clockwork Rocket}）
  \item \textbf{第二卷：《永恒之火》}（\emph{The Eternal Flame}）
  \item \textbf{第三卷：《时间之箭》}（\emph{The Arrows of Time}）
\end{itemize}

三部曲陆续由 Night Shade Books（旧金山/纽约）与 Orion/Gollancz（伦敦）出版：第一卷初版于 2011 年，第二卷于 2012 年，第三卷于 2013 年；此后有日文（早川书房）、意大利文（Mondadori/Urania）与俄文（Explorer Books）等译本。

围绕这一设定，作者撰写了一整套系统的物理注记：主页的"主论述"（main exposition）只需高中数学即可读懂，而标注"[进阶]"（[Extra]）的页面则深入得多，面向受过本科物理与数学训练、希望更详细钻研该主题的读者。全套注记逾八万词、配有一百余幅插图。本书即是在这些官方注记的基础上编译整理而成的中文专着，供物理学专业研究生阅读。

\section{本书结构}

本书依照循序渐进的原则编排，从最基础的运动学逐步推进到引力与量子领域。各章及其中一句话即可概括的核心结果如下：

\begin{itemize}
  \item \textbf{加与减}：从时空"勾股定理"中的加减号出发，建立洛伦兹与黎曼两种几何的基本对照。
  \item \textbf{双勾股定理}（The Dual Pythagorean Theorem）：给出一个简洁的几何证明，导出该宇宙中若干关键结果的根基。
  \item \textbf{几何与运动}（Geometry and Motion）：考察点积、无类时方向的时间、能量—动量矢量与守恒律，揭示"奇怪碰撞"等有别于相对论的碰撞规律。
  \item \textbf{几何与波}（Geometry and Waves）：讨论平面波、标量波与矢量波，说明波的叠加与偏振结构。
  \item \textbf{黎曼电磁学}（Riemannian Electromagnetism）：处理三种偏振、多普勒频移、电磁场张量，以及"波纹状库仑力"与非主流电子学。
  \item \textbf{黎曼热力学}（Riemannian Thermodynamics）：给出"比无穷热还热"的奇特结论，并推导黑体谱与理想气体性质。
  \item \textbf{黎曼广义相对论}（Riemannian General Relativity）：说明为何没有黑洞、没有引力波，并（很可能）没有大爆炸。
  \item \textbf{黎曼量子力学}（Riemannian Quantum Mechanics）：讨论固体结构、非理想谐振子、固体中的量子效应、黎曼原子与高分子。
  \item \textbf{世界观与文明设定}：梳理该宇宙中物理如何塑造其文明形态的诸般设定。
  \item \textbf{术语表}（Glossary）：汇总全书术语，便于随时查阅对照。
\end{itemize}

\orthofig{images/03__008.png}{略微不同的平面波之和（Sum of slightly different plane waves）}

按照官方注记的分类，本书同样保留了进阶部分的深入内容——这包括：双勾股定理的进阶讨论与对偶矢量；运动学中的对称性、自由粒子的哈密顿量与拉格朗日量；波论中关于对称群的表示、$SO(3)$ 与 $SO(4)$ 的双覆盖、自旋表示、四元数与李代数；电磁学的黎曼 Proca 方程、静电场与静磁场的各类算例、电磁能量流动与应力—能量张量；热力学的黑体谱与单粒子配分函数；广义相对论中的黎曼史瓦西解、引力红移与时间膨胀、宇宙学模型与引力波；量子力学中的黎曼狄拉克方程、黎曼伽马矩阵、平面波解与四元数方法；以及黎曼量子场论中的正则量子化、粒子与反粒子、传播子等内容。读者可根据需要选读。

\orthofig{images/03__009.png}{$SO(3)$ 中的不可收缩回路（Non-contractible loops in SO(3)）}

\orthofig{images/04__007.png}{点电荷周围的电通量（Electric Flux Around A Point Charge）}

\orthofig{images/05__005.png}{黎曼腔体辐射谱（Riemannian Cavity Radiation Spectrum）}

\orthofig{images/06__004.png}{黎曼引力散射（Riemannian gravitational scattering）}

\orthofig{images/07__011.png}{$SO(4)$ 不变的传播子（SO(4)-invariant propagator）}

\begin{keybox}[核心设定速记]
正交宇宙与我们宇宙的物理差异，可以浓缩为以下几点：

\begin{enumerate}
  \item \textbf{度规全正}：时空度规为黎曼型 $(+,+,+,+)$，而非洛伦兹型 $(-,+,+,+)$。
  \item \textbf{四维平等}：没有任何方向被单独指定为"时间"，不存在有特殊地位的类时方向。
  \item \textbf{时空图可如实作画}：正定的欧几里得式几何使尺规直觉在时空图中同样有效。
  \item \textbf{远行耗时而非省时}：高速星际旅行对旅行者意味着更长（许多代）的岁月，对故乡却只过数年。
  \item \textbf{无普适光速}：光速与波长（及产生方式）有关，不存在所有观察者共认的不变速度。
  \item \textbf{无速度上限}：不存在任何速度不得超越的极限。
  \item \textbf{发光释放能量}：光的产生过程释放能量，植物靠向夜空发光来制造养分。
  \item \textbf{引力与量子皆异}：没有黑洞、没有引力波，很可能也没有大爆炸；原子、固体与场论均须按黎曼几何重建。
\end{enumerate}
\end{keybox}

上面这一条条差异，都源自本章开篇处那一个符号的翻转。接下来的章节将逐一展开它们背后的数学与物理。

\chapter{加与减：正交宇宙物理速览}

本章从一条极简单的公式出发：只需把时空几何定律中的一个减号换成一个加号，
便能得到一个物理定律迥然不同的世界——\emph{黎曼宇宙}（Riemannian universe）。
我们将速览这一改动所引发的连锁后果：从固有的时间与光速，到能量与动量，
再到光的静止质量与物质的稳定性。读者会看到一些似曾相识的现象、一些奇异而优美
的结构，也会看到一些潜藏的危险。

\section{从减号到加号：时空中的勾股定理}

一百多年来我们已经知道，理解我们宇宙中时间的最佳方式，是把时间与熟悉的空间
三维结合起来，构成一个服从自身独特几何定律的四维\emph{时空}（space-time）。
牛顿（Newton）把时间视为绝对而普遍的量：它以完全相同的速率、整齐划一地为每个人
向前推进；而自爱因斯坦（Einstein）以来我们认识到，时间流逝所度量的，是我们在
四维时空中所走路径的一个属性，它与空间中一条路径的\emph{长度}非常相似。

没有人会对这一点感到惊讶：空间中连接同样两点的不同路径可以有不同的长度。
如果我从珀斯（Perth）走最直接的路线开车到悉尼（Sydney），而你中途绕道达尔文
（Darwin），那么我们的里程表读数不同，这毫不奇怪。我们现在明白，时空中的
“绕道”对时间流逝也有类似的作用。例如，若你留在地球上，“旅行”于 2050 年的
除夕与 2060 年的除夕之间；而我虽然同样在两次除夕与你相聚，却在中间这段时间
往返了一次半人马座$\alpha$（Alpha Centauri），那么我将老去约六年，而你老去十年。

空间中的绕道使路径更长，时空中的绕道却使它更短。其原因可以表述得很简单：
在度量时间间隔时，勾股定理的“时空版”用的是\emph{差}而不是\emph{和}——
是减号而不是加号。

勾股定理指出，直角三角形斜边的平方等于另外两条直角边平方之和。例如：

\orthofig{images/00__PM.001.png}{空间中的勾股定理：斜边平方等于两直角边平方之和}

而时空版则把平方之和换成平方之差：

\orthofig{images/00__PM.002.png}{勾股定理的时空版：用平方之差取代平方之和}

为使公式简洁，我们选取一种让时间与空间处于同等地位的单位：时间用“年”，
距离用“光年”。可以把时间间隔的平方写成

\begin{equation}
(\Delta \tau)^2 = (\Delta t)^2 - (\Delta x)^2 , \tag{1}
\end{equation}

其中 $\Delta\tau$ 为固有时（proper time）。利用两个这样的三角形，即可验证
前面关于半人马座$\alpha$之行的说法——它约有四光年之遥，而我此行的结果可能是
我只老去六年，而留在家中的你老去十年：

\orthofig{images/00__PM.003.png}{往返半人马座$\alpha$的时空三角形}

单程 $\Delta x = 4$ 光年、$\Delta t = 5$ 年，于是

\[
(\Delta \tau)^2 = 5^2 - 4^2 = 9 \;\Longrightarrow\; \Delta\tau = 3\ \text{年},
\]

往返两程合计 $6$ 年；而留在地球的人经历 $10$ 年。

当然，这些三角形在我们看来“比例不对”：那条本应长“$3$ 年”的斜边，
被画成明显长于代表“$5$ 年”的竖直边。但这正是要点所在：时空的“洛伦兹几何”
（Lorentzian geometry）并不等同于空间的几何，因此不可能在一张纸上把每个细节
都忠实地画出来。

现在可以很轻松地道出我们自己的宇宙与我一直在探索的替代宇宙之间的区别了。
在黎曼宇宙中，普通的勾股定理\emph{总是}成立，即使涉及时间间隔时也是如此。
这个宇宙的时空图可以完全忠实地画在纸上！那种朴素的直觉——认为一条“$5$ 年”
的边与一条“$4$ 光年”的边构成的三角形应当画成\emph{这样}：

\orthofig{images/00__PM.004.png}{黎曼宇宙中时空三角形的勾股关系}

如今是完全正确的。在黎曼宇宙中，绕道去往一颗遥远的恒星，对旅行者而言
总是意味着经历的时间更多，而不是更少。

星际航行对旅行者比对其他人耗时更长，听起来或许只是恼人；但由这同一处
简单改动所派生的差异远不止于此，它们波及从夜空的面貌到物质的稳定性等方方面面。

在我们宇宙中，(1) 式里的减号意味着：一旦距离（以光年计）超过时间（以年计），
平方之差就会变成一个负数。例如，设想用 $4$ 年穿行 $5$ 光年，公式给出的历时
平方便是 $-9$。既然没有实数的平方为负，这样一次假想的行程显然在性质上是不同的。
而这实际上是在告诉我们：没有任何东西能在如此短的时间内跨越如此大的距离。

在黎曼宇宙中，与之对应的行程毫无问题。对 $\Delta x = 5$ 光年、家乡系中
$\Delta t = 4$ 年的航行，船员测得的历时应为

\[
\Delta\tau = \sqrt{4^2 + 5^2} = \sqrt{41} \approx 6.4\ \text{年},
\]

而家乡的人只经历 $4$ 年，飞船被看作以光速的一又四分之一倍行进：

\[
v = \frac{5}{4}\,c = 1.25\,c .
\]

在黎曼宇宙中，不存在普适的速度上限。

事实上，黎曼宇宙中也不存在普适的\emph{光}速，因此“光年”和“光速”这些词
实际上只在我们自己的宇宙中才恰当。在黎曼宇宙中，代替“光速”的，是一个把时间
换算为距离的换算因子，它使勾股公式呈现出简洁漂亮的形式。正如你若用“公里”量
南北方向的距离、用“英里”量东西方向的距离，勾股定理就会失效；在黎曼宇宙中，
同样需要通过实验去测定：时间与距离各取什么单位，才能使该定理成立。

\section{没有普适光速的宇宙}

生活在一个光以不同速度传播的宇宙中会是什么样子？当然，光穿过玻璃之类的
透明介质时，在我们宇宙中也会这样；但在黎曼宇宙中，即使在真空里也会如此。
我们把不同波长的光感知为不同的颜色；在黎曼宇宙中，这些不同的波长将\emph{总是}
以不同的速度传播。

要找出波长与速度的关系，可以注意这样一点：无论光速如何，光波前在\emph{时空}
中的间距总是相同的。为什么？我们期望自己宇宙中的物理不会偏向三维空间里的任何
方向：若我们用一支激光笔朝东南方发射，我们会认为它产生的光，与朝正北或西北方
发射时产生的光性质完全相同。同样，黎曼宇宙中底层的物理也应当对刻画光速的
那个时空方向无动于衷。

假设一束速度适中的光脉冲的时空图如下：

\orthofig{images/00__PM.005.png}{时空中的光脉冲与波前}

这里我们为光脉冲画出的“路径”显示了它随时间的整体空间运动。正如水波的波峰
垂直于波传播的方向，标记这束光振荡峰值的波前，也垂直于它的路径。

光的\emph{波长}是这些波前在空间中的间距。当然，我们称为“空间”的那个方向
取决于我们自身的运动状态，因此我们测得的波长也会随之改变。相反，波前沿其
法向（即与波前垂直的方向）度量的\emph{间距}，是光本身的性质，与谁在测量无关。

要得到较慢光脉冲的波长，只需把前一幅图中的整族波前整体转动，
而不改变它们彼此之间的间距：

\orthofig{images/00__PM.006.png}{时空中较慢的光脉冲与波前}

由此可见，传播得越慢的光，波长越长。把它翻译成我们感知各种波长颜色的语言，
便可以说：红光将比紫光慢。

当你仰望黎曼宇宙的夜空时，从每颗恒星看到的红光，要比紫光多花数年乃至数百
年才能抵达你这里——所以除非这颗恒星完全没有横向运动，不同颜色将来自天空上
稍有差异的位置。于是每颗恒星都不再表现为单一的光点，其像会被展宽成一段
光谱，或一条“色谱拖尾”（colour trail），其中红、紫两端之间的张角直接度量了
该恒星的横向速度。

\begin{notebox}[补充：波长与速度的关系]
若记波前族的法向间距为 $d$（这是不变量），并设光在 $(x,t)$ 平面内的传播方向为
$(v,1)$，则在黎曼（欧氏）几何中，波前法向的单位分量为
$n_x = v/\sqrt{1+v^2}$。沿空间 $x$ 轴量得的波长为
\[
\lambda = \frac{d}{n_x} = \frac{d\sqrt{1+v^2}}{v},
\]
可见速度 $v$ 越小，波长 $\lambda$ 越长——与上面“转动波前族”的几何论证完全一致。
（此推导为补充说明，原文只给出转动波前族的定性论证。）
\end{notebox}

\section{能量-动量矢量}

除了没有普适的速度上限，关于黎曼宇宙中物质的行为还有什么可说的？要理解最
引人注目的差异，我们得先勾勒我们自己宇宙物理学中的几个简单概念。

对于运动的物体，有两个有用的量：它的动量（momentum）与动能（kinetic energy）。
刻画它们的最简单方式，也许是设想把一颗子弹射入某种屏障，该屏障以恒定力阻碍
其运动使之停下。子弹的\emph{动量}决定它停下来的快慢，而它的\emph{动能}决定它
停下来之前走过多远。这两个量都只取决于子弹的质量与速度。

自相对论问世以来，我们认识到：正如空间与时间最好被看成同一实体（时空）的
两个方面，运动物体的能量与动量也是同一个几何对象的两个方面，这个对象叫作
\emph{能量-动量矢量}（energy-momentum vector）。这听起来或许很专业，但实际上
它非常容易想象。对于静止的物体，能量-动量矢量就是一支长度等于物体质量、
指向时空\emph{时间}方向的箭头。

\orthofig{images/00__PM.007.png}{我们宇宙中的能量-动量矢量}

若使同一物体运动起来，原来的箭头只是被倾斜，以反映物体在给定时间内穿越的
空间大小：慢速物体略微倾斜，快速物体倾斜更大。

在上图中，箭头看上去似乎不只是被倾斜，还变长了。但这只是把时空画在纸面上
造成的错觉；若我们用时空版的勾股定理去计算倾斜箭头的长度，会发现它与那支竖直
箭头长度完全相同。

这些能量-动量矢量有什么用？若你测量物体能量-动量矢量在时间方向上的投影长度，
它告诉你物体的\emph{能量}。若你测量它在空间方向上的投影长度，它告诉你物体的
\emph{动量}。

\orthofig{images/00__PM.008.png}{我们宇宙中的总能量、动量与动能}

这里唯一稍显棘手之处在于：能量-动量矢量在时间方向上的投影给出的是物体的
\emph{总能量}（total energy），它不仅包含动能，还包含与物体质量相联系的能量。
因此，静止物体虽然有零动量（正如你所预期），却\emph{不是}零总能量；它的总能量
完全来自其质量（如前所述，质量决定了这些箭头的长度）。当同一物体被赋予运动时，
其总能量在“静质量能”之上的\emph{增量}，就是它的动能。

为什么我们没有引用那个著名的静质量能公式 $E=mc^2$？因为我们所用的这类单位里，
光速 $c$ 总等于 $1$（例如一年每光年），于是 $E=m$，无需 $c^2$。

那么在黎曼宇宙中又如何？我们仍然画一支长度等于物体质量的箭头，并根据物体的
速度给它一个相应的倾斜：

\orthofig{images/00__PM.009.png}{黎曼宇宙中的能量-动量矢量}

这一次，倾斜看起来就是普通的倾斜：我们肉眼就能看出两支箭头长度相同。时空的
特殊勾股定理已被普通勾股定理取代，因此我们可以在普通的纸上画出黎曼宇宙的
能量-动量矢量，并用普通的直尺量其长度。

这一改动如何影响运动物体的能量与动量？

\orthofig{images/00__PM.010.png}{黎曼宇宙中的总能量、动量与动能}

在黎曼宇宙中，倾斜箭头使它在时间方向上的投影\emph{小于}物体静止时的投影。
所以物体运动时的总能量比静止时\emph{更小}！我们仍然可以把物体的动能定义为
二者之差，但它现在代表的是一种能量的减少，而非增加。

声称运动物体比静止物体能量更低，听起来极其古怪；但事实上，动能与总能量之间
这种“上下颠倒”的特征，在最简单的情形里根本不会显现。我们熟悉的全部经典牛顿
力学，对慢速运动的物体仍然适用：你完全可以在黎曼宇宙里打一局台球，而察觉不到
任何不同。低速物体的动量与动能的近似公式，在两个宇宙中完全相同——只是动能
差一个负号，而台球桌上发生的任何事都不会揭示这个负号的意义。动能相对静质量能
是负的，而在我们宇宙中静质量能只在核反应中才显现。所以，黎曼宇宙这个古怪的
特性也许会在日常世界中一直隐藏着，只在玄奥的高能实验中才露出端倪？

\begin{equation}
E = \begin{cases}
\sqrt{m^2 + p^2} \approx m + \dfrac{p^2}{2m} & \text{(洛伦兹型)},\\[2ex]
\sqrt{m^2 - p^2} \approx m - \dfrac{p^2}{2m} & \text{(黎曼型)} .
\end{cases}
\end{equation}
式中 $p^2/(2m)=\half mv^2$ 即经典动能。两式在低速下完全对应，仅动能项的符号相反。

\section{光有静止质量}

但有一件远非玄奥的事情被我们忽略了：光！在我们自己的宇宙中，光没有静止质量，
而时空版的勾股定理允许光的能量-动量矢量具有相应的长度\emph{零}。为此，总能量
不必为零，只需等于动量（在光速为 $1$ 的单位下）即可。这样，矢量在空间方向与
时间方向上的投影，便在平方之差里相互抵消了。

\orthofig{images/00__PM.011.png}{我们宇宙中光的能量-动量矢量}

\begin{equation}
E^2 - p^2 = 0 \;\Longrightarrow\; E = p . \tag{2}
\end{equation}

在黎曼宇宙中，这是不可能的：一个矢量长度为零的唯一方式，是它在所有方向上的
投影都为零。静止质量为零就意味着完全没有能量与动量。所以在黎曼宇宙中，
光具有静止质量，和别的一切东西一样。

\begin{equation}
E^2 + p^2 = m^2 \;\Longrightarrow\; m>0 .
\end{equation}

\section{发光的能量学与物质的稳定性}

现在，任何产生光的过程都需要从别的形式（如动能或化学能）获取能量。但当动能
相对总能量为\emph{负}时，“获取”一词就不太恰当了。产生光时我们必须遵守能量
守恒定律，而既然光的总能量为正，那么产生光的那个体系的总能量就必须下降……
也就是说，它的动能必须\emph{增加}。

换言之，一个产生光的体系将经历动能的增加：它的各组成部分最终会比之前运动得
更快！

\orthofig{images/00__PM.012.png}{产生光的反应中的能量平衡}

举一个简单而具体的例子。假设我们做一个产生光的化学反应。当把化学物质混合在
一起时，它们大致是静止的（它们因热而带有一些动能，但我们一开始假定这可以
忽略）。为简单起见，我们还假定这两种化学物质质量相等。于是，在“反应前”的图中，
它们的能量-动量矢量只是两支相同的、指向时间方向的箭头。

“反应后”的图则需要包含所生成的光的能量-动量矢量，以及化学反应产物的能量-动量
矢量。我们假定产物与反应物质量相同；这未必成立，但正如我们将会看到的，
它当然是可能的。

要在“反应后”的图中把所有矢量安排得使总能量不变，产物的能量-动量矢量就必须
被倾斜。也就是说，产物必须运动起来：它们需要获得动能并降低总能量，以便为所
生成的光“垫付”总能量。

\begin{notebox}[反应前后能量配平的代数]
设两种反应物质量相等，均为静止（动量为零）：反应前总能量 $E_{\text{前}} = 2m$。
反应后设两产物质量仍为 $m$、携带反向动量 $\pm p$，并生成总能量为 $E_{\text{光}}>0$
的光。能量守恒要求
\[
E_{\text{前}} = E_{\text{后}} = E_{\text{光}} + 2\sqrt{m^2 - p^2} = 2m .
\]
由 $E_{\text{光}}>0$ 立得 $\sqrt{m^2 - p^2} < m$，即产物必须处于运动状态：其总能量
低于静质量能，从而为光提供了能量。（其中光与产物动量的分配由动量守恒
$\sum \mathbf{p} = \bzero$ 决定。）
\end{notebox}

化学反应中产生的动能表现为热，因此该反应将同时产生光与热。在我们宇宙中，
燃烧任何燃料也会产生光与热，但在那种情况下，它们是以化学能的改变为代价换来的，
而这又表现为产物与反应物之间（极其微小的）质量差异。而在这里，我们无需依赖任何
这样的变化，就制造出了光与热！我们特意令产物的质量与反应物完全相同，但即使把
产物质量取得略高或略低，反应也仍然能够发生。

在黎曼宇宙中，产生光的过程可以产生热，无论是否伴随化学能的改变。当然，化学能
有时是比热更有用的产物。黎曼宇宙中的植物不会吸收光来进行光合作用，反而会
\emph{发出}光，以储存由此生成的化学能。

产生光总是伴随动能的产生，这使黎曼宇宙成为一个危险的地方。生命需要掌握某些
精细的反应，才能安全地从这一过程中获取能量；就连最简单的物质基本单元也将面临
不稳定的威胁。在我们宇宙中，原子或分子可以吸收一个光子而进入一个结合得较松的
高能态；例如阳光中的紫外线可以把某些分子打散。在黎曼宇宙中，同类过程不再需要
外部辐射源：处于最紧束缚态的物质可以\emph{发出}自己的光，并借此获得动能来把
自身撕裂。由于这一效应，只有某些结构是稳定的。

\section{研究异宇宙的价值}

研究一个由不同于我们的定律所统治的宇宙，有什么价值？我们当然可以设想这样
一个宇宙或许真实存在——那里也有居民，也在思索：如果他们自己的某个加号变成了
减号，生活会变成什么样。但把我们的定律拆解开来，看看它们在新假设下如何运作，
也能为我们自己宇宙的物理学提供洞见。通过努力想象事物可能有何不同，我们可以
加深对许多事物的理解——无论是一个人类社会、一颗行星的地质、生命的化学，
还是最基本的物理定律。

黎曼宇宙充满了危险与惊奇。本文只是触及了它的皮毛，但希望读者已从中尝到了
一点它所能带来的乐趣。

\begin{keybox}[本章要点]
\begin{itemize}
  \item 把时空勾股定理的一个减号改为加号：洛伦兹型给出
        $(\Delta\tau)^2=(\Delta t)^2-(\Delta x)^2$，即“绕道使固有时更短”；
        黎曼型给出 $(\Delta\tau)^2=(\Delta t)^2+(\Delta x)^2$，即“绕道使时间更长”。
        时间与距离的换算单位必须由实验确定。
  \item 黎曼宇宙没有普适光速：波长越长（越红）的光传播越慢。波前在时空中的
        \emph{法向间距}是不变量，快慢之别仅表现为整族波前的转动；恒星因而呈现
        “色谱拖尾”，红紫张角直接度量其横向速度。
  \item 能量-动量矢量：长度等于质量，静止时指向时间方向；时间分量是总能量，
        空间分量是动量，动能 = 总能量 $-$ 静质量能。黎曼宇宙中运动物体的总能量
        比静止时\emph{更小}，故动能相对总能量为负；低速极限下与牛顿力学一致，
        仅动能差一个负号。
  \item 黎曼几何中零长矢量只能是零矢量，故光必有静止质量。于是“产生光必须付出
        能量”，发光体系的动能反而增加（表现为放热）；植物的光合作用逆转为
        “发光储能”；处于最紧束缚态的物质可自发光并借此获得动能、撕裂自身，
        使得只有特定结构才稳定。
\end{itemize}
\end{keybox}

\chapter{双勾股定理（The Dual Pythagorean Theorem）}

《正交》（\emph{Orthogonal}）宇宙中的许多物理, 都可以理解为由几何学中相当简单的想法所导出。
尽管由此得到的物理结论时常与我们自身宇宙中的结论大不相同, 几何本身却并无二致。
本章介绍其中一个极其简单的例子, 它却是小说中最重要的科学发现之一背后的基础;
同时, 它也是一个凭其在现实世界中的应用便值得掌握的结果。
我们将先给出这一定理及其几何证明, 随后在矢量空间的语言下, 说明它在对偶空间中的表述。

\section{勾股定理的两种形式}

我们先从一个几乎人人都会一眼解决的问题谈起。

\orthofig{images/01__001.png}{斜向道路的长度}

假设有人告诉我们: 一辆汽车沿一条笔直的道路行驶, 这条路既不完全南北向, 也不完全东西向。
汽车最终停在出发点以北 $3\,\mathrm{km}$、以东 $4\,\mathrm{km}$ 处。那么汽车究竟行驶了多远？

勾股定理（Pythagoras's Theorem）告诉我们, 直角三角形斜边长度的平方等于另外两条直角边长度的平方之和。
于是答案立刻可得: $3$ 与 $4$ 的平方分别为 $9$ 与 $16$, 它们的和为 $25$, 而 $25$ 正是 $5$ 的平方。
所以汽车行驶的总距离为 $5\,\mathrm{km}$。

\orthofig{images/01__002.png}{田地中平行等距的犁沟}

接下来, 考虑一块用平行等距的犁沟耕作过的田地, 这些犁沟既不完全南北向, 也不完全东西向。
从一根钉入地面的木桩出发, 如果我沿正北方向拉出一段绳子, 会发现它恰好穿过 $3$ 条完整的犁沟;
如果我把\textbf{同样长度的一段绳子}改沿正东方向拉出, 会发现它穿过 $4$ 条完整的犁沟。
问题在于: 如果我选择一个方向, 使绳子径直、垂直地横穿犁沟, 那么它会穿过多少条犁沟？

如果你依照上一题类推, 回答"五条", 那么你的答案完全正确——
但如果你以为这不过是勾股定理的又一次直接应用, 那就应当仔细看看插图。
这里\textbf{并不}存在边长为 $3,4,5$ 个单位的直角三角形!

然而这个答案并非误导性的巧合: 若我们把正北、正东方向穿过的犁沟数分别换成任意的 $n$ 与 $e$,
那么同样长度的绳子垂直穿过犁沟的数目 $f$ 将满足
\begin{equation}
f^{2} = e^{2} + n^{2}.
\end{equation}
但若它\textbf{不是}勾股定理的又一个实例, 这个关系究竟从何而来？

\orthofig{images/01__003.png}{田角处的直角三角形}

如果我们考察一个直角三角形, 它的直角顶点在木桩处, 斜边是经过木桩的第一条犁沟,
两条直角边分别沿正北与正东方向, 那么这两条边的长度将分别为 $r/3$ 与 $r/4$,
其中 $r$ 是绳子的长度。

连接该三角形的直角顶点与斜边、并与斜边成直角的那条线段, 其长度为 $r/5$。

\orthofig{images/01__004.png}{带斜边高的直角三角形}

现在不妨转向一般情形: 一个直角三角形, 两条直角边长分别为 $a$ 与 $b$, 斜边长为 $c$,
从直角顶点到斜边的垂直距离为 $d$。

勾股定理给出 $a,b,c$ 之间的关系:
\begin{equation}
c^{2} = a^{2} + b^{2},
\end{equation}
但 $a,b,d$ 之间又有什么关系呢？

虽然 $a,b,c$ 都是三角形的\textbf{边长}, 而 $a,b,d$ 乍看像是毫不相干的三元组, 实际上却并非如此。
它们的共同点在于: 三者都是三角形的\textbf{高}（altitude）——从三角形的某条边到其对顶点的垂直距离。
这为看出它们的关系提供了一条简单的门径: 三角形的面积等于底乘高（相对于该底的高）的一半,
换言之, 等于任意一条边的长度, 乘以从该边到对顶点的垂直距离, 再除以二。

这一事实给出了同一个三角形面积的两种不同表达式, 二者必须相等:
\begin{equation}
\tfrac{1}{2}\, a b = \tfrac{1}{2}\, c d.
\end{equation}
两边同时乘以 $2$, 再平方, 得到
\begin{align}
a^{2} b^{2} &= c^{2} d^{2}, \\
a^{2} b^{2} &= (a^{2} + b^{2})\, d^{2}, \\
\frac{1}{d^{2}} &= \frac{a^{2} + b^{2}}{a^{2} b^{2}},
\end{align}
其中从第一式到第二式用到了勾股定理, 从第二式到第三式则在等式两边同除以 $a^{2} b^{2} d^{2}$。
若将最后一式右端的分母通分消去, 即得

\begin{equation}
\frac{1}{d^{2}} = \frac{1}{a^{2}} + \frac{1}{b^{2}}. \tag{Dual Pythagorean Theorem}
\end{equation}

这一关系适用于任意直角三角形: 其中从两条直角边量得的高为 $a$ 与 $b$, 从斜边量得的高为 $d$。
这正是勾股定理的\textbf{对偶}（dual）形式——它与"平方和"版本 $c^{2}=a^{2}+b^{2}$ 互为对偶:
前者联系三条边长, 后者联系三条高。

现在我们可以\textbf{确切地}看清, 为什么同一段绳子在正北、正东方向穿过的犁沟数 $n,e$,
与垂直穿过的数目 $f$, 满足 $f^{2} = e^{2} + n^{2}$。
木桩处那个三角形对应的距离为
\begin{equation}
a = \frac{r}{e}, \qquad b = \frac{r}{n}, \qquad d = \frac{r}{f},
\end{equation}
代入双勾股定理, 得到
\begin{align}
\frac{1}{d^{2}} &= \frac{1}{a^{2}} + \frac{1}{b^{2}}, \\
\frac{f^{2}}{r^{2}} &= \frac{e^{2}}{r^{2}} + \frac{n^{2}}{r^{2}}, \\
f^{2} &= e^{2} + n^{2}.
\end{align}

\orthofig{images/01__005.png}{地形等高线}

这个结果为什么在物理学中如此重要——无论在《正交》宇宙中, 还是在我们自身的宇宙中？
它所告诉我们的是: 勾股定理的"平方和"公式, 不仅被长度这样的量所遵守,
也被那些"对偶"量所遵守——这些量计数的是\textbf{每}单位长度发生多少次事件, 或\textbf{每}单位长度变化多少。

我们可以把耕过的田地里的犁沟, 换成地形等高线图上的等海拔线。
于是双勾股定理告诉我们: 例如, 若我们所站的地形每向北行进 $1\,\mathrm{km}$ 下降 $30\,\mathrm{m}$,
每向东行进 $1\,\mathrm{km}$ 下降 $40\,\mathrm{m}$,
那么（至少在我们所观察的尺度上, 地形足够接近一个倾斜平面时）
它最陡的坡度, 亦即沿垂直横穿等海拔线方向的梯度, 将是每千米 $50\,\mathrm{m}$。

\orthofig{images/01__006.png}{波前}

然而, 这一结果最重要的应用大概要算波的几何。
若我们把犁沟换成\textbf{波前}（wavefront, 例如海浪的波峰, 或电磁波的波峰）,
那么在给定位移内出现的波前数目就称为\textbf{空间频率}（spatial frequency）,
以区别于计数给定时间内振荡次数的普通\textbf{时间频率}。
双勾股定理告诉我们: 一个波的总空间频率——沿波传播方向测得的每单位距离波前数——
其平方等于沿两个或更多相互垂直的方向做同样测量所得结果的平方之和。

在我们自身的宇宙中, 这已经是一个有用的知识。
而在《正交》的宇宙中, 它彻底革新了人们对物理学的理解。

对于不介意稍高一些数学水平的读者, 我们在下一节提供了关于本主题的进阶材料。

\section{对偶矢量}

\begin{notebox}[进阶内容]
本节面向愿意接受稍高数学水平的读者, 讨论\textbf{对偶矢量}（dual vectors）,
并从对偶空间的角度重新表述双勾股定理。
\end{notebox}

设 $V$ 是一个 $n$ 维矢量空间, $f$ 是从 $V$ 到实数集 $\mathbb{R}$ 的一个线性函数。

$f$ 的\textbf{核}（kernel）, 记作 $\ker f$, 是 $V$ 中使 $f$ 取零的那个子空间。
$\ker f$ 的维数必定至少为 $n-1$（除非 $f$ 是零函数, 此时它恰好等于 $n-1$）。
例如, 若 $n=2$ 且 $f \neq 0$, 则 $f$ 将在 $V$ 中一条过原点的直线上为零;
若 $n=3$ 且 $f \neq 0$, 则 $f$ 将在一个过原点的平面上为零, 依此类推。
这来自线性代数的一个基本定理, 即\textbf{秩-零化度定理}（Rank-Nullity Theorem）。

\orthofig{images/01__007.png}{二维矢量空间上的线性函数}

\orthofig{images/01__008.png}{二维矢量空间上的线性函数}

图中给出一个二维的例子; 这里函数 $f$ 在由矢量 $\bk$ 张成的过原点直线上为零。

若我们选取一个使 $f(\bv)$ 非零的矢量 $\bv$, 那么在 $V$ 中就有一条\textbf{不过原点}、
而是经过 $\bv$ 端点的直线, 其上 $f$ 恒等于 $f(\bv)$。
我们通过给 $\bv$ 加上 $\bk$ 的任意倍数来构造这条直线:
\begin{equation}
f(\bv + s \bk) = f(\bv) + s\, f(\bk) = f(\bv) = 1.
\end{equation}
最后一处等号之所以成立, 只是因为我们恰好选取了满足 $f(\bv)=1$ 的 $\bv$。
不过, $f$ 取常值的所有集合都具有这种形式: 平行于矢量 $\bk$ 的直线。
正如我们在图上把不同的矢量画成长短、方向各异的箭头,
我们也可以把不同的线性函数画成疏密、取向各异的\textbf{线束}（stacks of lines）。

越长的矢量画越大的箭头; 但在这一表示方案里, 一个更大的函数并不是画成间距更大的直线,
而是画成\textbf{排布得更紧密}的直线。
例如, 函数 $g = 2 f$ 的线束中, 每一对直线之间都会多出一条线,
因为当 $f$ 仅为 $1/2$ 时它就已经达到 $1$。

对任意矢量而言, $f$ 的值实际上是在计数该矢量的箭头穿过了线束中的多少个间隙。
例如, $\bw$ 的箭头穿过了 $f$ 的线束中的两个完整间隙, 故 $f(\bw)=2$。
我们还需顾及反向穿过间隙从而给出负值的情形（如 $\bp$）,
以及未穿过整数个间隙的情形（如 $\bq$）;
因此这更像是诉诸几何直观, 而非严格的数学方案。
不过你大概已经能看出, 它如何与主页上的例子相呼应——
在那里, 我们讨论的是在某个方向上移动给定距离时, 穿过了田间的多少条犁沟, 或地图上的多少条等高线。

若 $V$ 的维数大于 $2$, 那么在 $V$ 中我们得到的将不是一束直线, 而是维数为 $n-1$ 的一束平面或超平面。

现在, 考虑从 $V$ 到 $\mathbb{R}$ 的\textbf{所有}线性函数构成的集合。
我们可以把这个集合本身变成为一个 $n$ 维矢量空间, 它称为 $V$ 的\textbf{对偶空间}（dual space）, 记作 $V^{*}$。
$V^{*}$ 中的元素（例如 $f$）称为\textbf{对偶矢量}（dual vectors）。

要使一个集合成为实矢量空间, 我们需要能够把它的元素相加, 并乘以实数。
对 $V$ 上的线性函数集合, 这可以很自然地做到:
若 $f,g \in V^{*}$、$\bv \in V$、$s \in \mathbb{R}$,
我们通过规定它们作为函数取的值, 来定义和 $f+g$ 与标量倍 $s f$:
\begin{align}
(f+g)(\bv) &= f(\bv) + g(\bv), \\
(s f)(\bv) &= s\, f(\bv).
\end{align}

与任何矢量空间一样, 我们可以为 $V^{*}$ 选择一组基, 即 $n$ 个线性函数,
使得 $V^{*}$ 中的任何函数都可以用它们线性表出。
我们把这组基的元素记作 $\{\be^{1}, \be^{2}, \dots, \be^{n}\}$, 并写成
\begin{equation}
f = f_{1}\,\be^{1} + f_{2}\,\be^{2} + \dots + f_{n}\,\be^{n},
\end{equation}
其中 $f_{1}$ 等是 $f$ 关于这组基的\textbf{分量}。
在本注记中, 我们通常用上标标记矢量的分量, 却用\textbf{下标}标记对偶矢量的分量;
并且, 我们通常用下标标记基中的矢量, 却用\textbf{上标}标记基中的对偶矢量。

\orthofig{images/01__009.png}{对偶基}

若 $V$ 有一组特定的基 $\{\be_{1}, \be_{2}, \dots, \be_{n}\}$,
那么当 $V^{*}$ 的一组基 $\{\be^{1}, \be^{2}, \dots, \be^{n}\}$ 满足下述条件时,
我们说它\textbf{对偶于}（dual to）$V$ 的这组基:
\begin{equation}
\be^{i}(\be_{j}) = \delta^{i}_{j},
\end{equation}
其中 $\delta^{i}_{j}$ 称为\textbf{克罗内克 $\delta$ 符号}（Kronecker delta symbol）,
当 $i=j$ 时为 $1$, 当 $i \neq j$ 时为 $0$。
从几何上看, 这意味着: 对某个特定的 $i$, 函数 $\be^{i}$ 在 $V$ 中表现为一束直线或平面,
使得矢量 $\be_{i}$ 恰好穿过相邻两线之间的一个间隙,
而所有其他基矢量 $\be_{j}$ 都落在过原点的那条直线或平面之内, 完全不穿过该线束。
例如, 在左图中, 矢量 $\be_{2}$ 落在直线 $\be^{1}=0$ 之内, 而矢量 $\be_{1}$ 落在直线 $\be^{2}=0$ 之内。

给定 $V$ 的基 $\{\be_{1}, \be_{2}, \dots, \be_{n}\}$, 我们只需把原基中的每个矢量 $\be_{i}$ 代入 $f$,
即可求得 $V^{*}$ 中任意线性函数 $f$ 关于对偶基的分量。
因为两组基互为对偶, 只有特定的那个基函数 $\be^{i}$ 在 $\be_{i}$ 处非零,
于是求和中只剩下乘以它的那一项分量 $f_{i}$ 存活下来:
\begin{equation}
f(\be_{i}) = (f_{1}\be^{1} + f_{2}\be^{2} + \dots + f_{n}\be^{n})(\be_{i}) = f_{i}.
\end{equation}

若矢量 $\bv \in V$ 关于 $V$ 某组选定基的分量为 $v^{i}$,
而对偶矢量 $f \in V^{*}$ 关于对偶基的分量为 $f_{i}$, 则
\begin{equation}
f(\bv) = f(v^{i} \be_{i}) = f_{i}\, v^{i},
\end{equation}
其中我们用了\textbf{爱因斯坦求和约定}（Einstein summation convention）来简写对重复指标的求和
（即 $v^{i}\be_{i} = v^{1}\be_{1} + v^{2}\be_{2} + \dots + v^{n}\be_{n}$）。

假设 $V$ 上有一个点积。那么对 $V$ 中任意矢量 $\bw$, 我们可以定义一个 $V^{*}$ 中的线性函数 $f_{\bw}$:
\begin{equation}
f_{\bw}(\bv) = \bv \cdot \bw.
\end{equation}
反过来, 给定 $V^{*}$ 中的任意线性函数 $f$, 我们可以找到这样的矢量 $\bw_{f}$:
使得对每个 $\bv \in V$ 都有 $f(\bv) = \bv \cdot \bw_{f}$。
为此, 我们先选取 $V$ 的一组正交归一基 $\{\be_{1}, \be_{2}, \dots, \be_{n}\}$, 然后令
\begin{equation}
\bw_{f} = f_{i}\, \be_{i},
\end{equation}
这里仍使用爱因斯坦求和约定, 而分量 $f_{i}$ 是相对于与 $V$ 的正交归一基对偶的那组 $V^{*}$ 基取的。于是
\begin{align}
\bv \cdot \bw_{f}
&= \bv \cdot [f_{i} \be_{i}] \\
&= f_{i}\,[\bv \cdot \be_{i}] \\
&= f_{i}\, v^{i} = f(\bv).
\end{align}
因此, 我们可以把 $V$ 中的任一元素与 $V^{*}$ 中的唯一元素对应起来, 反之亦然。

我们用这种方式与 $f$ 相联系的矢量 $f_{i} \be_{i}$,
将正交于任何落在 $f$ 的核中的矢量 $\bk$,
因为 $\bk$ 与 $f_{i} \be_{i}$ 的点积恰好就是 $f(\bk)=0$。
由于我们用来可视化 $f$ 的线束中所有的直线（或平面等）都平行于过原点的那条核直线,
矢量 $f_{i} \be_{i}$ 将垂直于整个线束。

\begin{notebox}[补充: 与洛伦兹情形的对比]
上述 $V$ 与 $V^{*}$ 之间的对应, 本身是在一个\textbf{正定}（黎曼型, 度规全为正是 $(+,+,+,+)$）点积下得到的:
在正交归一基中, 这种对应使矢量与对应余矢量的分量\textbf{逐一相等}, 即 $v_{i} = v^{i}$,
因而两者可以被"等同"起来。
而在本宇宙常用的\textbf{洛伦兹}（Lorentzian）情形中, 度规取 $(-,+,+,+)$,
矢量与余矢量之间只相差一个\textbf{时间分量上的符号}: 若某矢量的分量为 $(v^{x}, v^{y}, v^{z}, v^{t})$,
则由度规降指标所得的对应余矢量分量为 $v_{a} = g_{ab}\, v^{b}$, 即
\[
v_{x} = v^{x}, \qquad v_{y} = v^{y}, \qquad v_{z} = v^{z}, \qquad v_{t} = -\,v^{t}.
\]
因此, 在洛伦兹情形中, 对偶矢量与原始矢量不能再像黎曼情形那样被简单地等同起来,
两者的时间分量相差一个符号。
\end{notebox}

我们可以给 $V^{*}$ 赋予一个点积, 办法是规定:
与 $V$ 的任意正交归一基对偶的任意一组基, 本身也是正交归一的。
这让我们得以经由其平方来定义对偶矢量 $f$ 的"长度"或大小:
\begin{equation}
|f|^{2} = f \cdot f
= (f_{1}\be^{1} + f_{2}\be^{2} + \dots + f_{n}\be^{n})
  \cdot (f_{1}\be^{1} + f_{2}\be^{2} + \dots + f_{n}\be^{n})
= (f_{1})^{2} + (f_{2})^{2} + \dots + (f_{n})^{2}.
\end{equation}
此处 $f_{i}$ 是相对于 $V^{*}$ 的基 $\{\be^{1}, \be^{2}, \dots, \be^{n}\}$ 取的分量,
而该基对偶于 $V$ 的一组正交归一基。

\orthofig{images/01__006.png}{波前}

按照我们对 $f$ 的几何解释, 这里每个分量 $f_{i} = f(\be_{i})$ 都是在计数基矢量 $\be_{i}$ 穿过了多少个"线束中的间隙"。
所以它们就像右图中所画的: 当我们沿 $x$ 轴与 $y$ 轴各移动一米时, 穿过了多少个波的周期。

但现在假设我们径直地、\textbf{垂直}地穿过线束,
用的单位矢量沿着 $f_{i} \be_{i}$ 的方向——确切地说, 是 $f_{i} \be_{i} / |f_{i} \be_{i}|$——
而我们知道它正交于整个线束。这样一次垂直穿越所得到的计数为
\begin{align}
f\!\left(\frac{f_{i} \be_{i}}{|f_{i} \be_{i}|}\right)
&= \frac{f_{i}\, f(\be_{i})}{|f_{i} \be_{i}|} \\
&= \frac{f_{i}\, f_{i}}{|f_{i} \be_{i}|} \\
&= \frac{|f|^{2}}{|f|} \\
&= |f|.
\end{align}

这正是双勾股定理的又一种表述!
用一个单位矢量径直穿过线束所得到的计数为 $|f|$,
而我们已经看到, 它的平方恰好等于分别沿 $n$ 个相互正交方向的单位矢量穿过该线束所得计数的平方之和。
换言之, 在对偶空间中, 那个"平方和"形式（勾股定理的原始版本）再次浮现——
这正是"双勾股"之名所传达的对偶关系: 矢量空间中的长度, 与对偶空间中的计数, 服从彼此对偶的两条勾股关系。

\chapter{几何与运动}

本章建立正交宇宙（\emph{Orthogonal}）中几何与运动学的基本框架。我们从洛伦兹几何（Lorentzian geometry）与黎曼几何（Riemannian geometry）的对比入手：前者以不变量 $x^{2}-t^{2}$ 刻画，后者以 $x^{2}+t^{2}$ 刻画，从而取消了类时、类空与类光方向之别。随后我们定义四维矢量的长度与点积，说明黎曼时空为何不存在任何优先的“时间方向”，进而引入能量-动量矢量与守恒定律，并以“奇异的碰撞”展示这一几何所导致的、在物理上不可回避的反直觉结论。本章进阶部分讨论不同观测者之间的能量比较、时空对称性，以及自由经典粒子的哈密顿量与拉格朗日量。

\section{洛伦兹几何与黎曼几何}

在我们居住的宇宙中，我们用来命名空间方向的词汇——上、下、左、右、前、后、南、北——全都由具体的物体来定义，诸如我们自己的身体，或我们脚下的大地。这样做只是出于方便，但归根结底我们别无选择：我们必须借助物体来辨认方向，因为空间本身的底层几何对所有方向一视同仁。

不那么明显的是，当我们谈论“未来”或“过去”时，我们实际上隐含着在指某个特定对象（或某群对象）的未来或过去。当我们说“向未来”移动一小时，我们通常指的是相对于周围的街道与建筑在空间中保持静止。但一辆疾驰的汽车在时空中划出的方向，与路边灯柱划出的方向一样，都是“未来”的一个同样合理的选择——同样地，时空的底层几何对所有这样的方向也一视同仁。

\orthofig{images/02__001.png}{洛伦兹时空中的方向}

然而，尽管你称之为“左”的方向可能正是我称之为“前”的方向，尽管当你认为我正以每秒一千公里的速度在你身旁飞掠而过时、你可能认为自己“静止不动、只在向未来前进”，但并\emph{非}时空中\emph{任意}两个方向都不可区分。归根结底，我们两人总会一致同意\emph{类时}（timelike）、\emph{类空}（spacelike）与\emph{类光}（null）方向之间的区分。

\emph{类时}方向是周围普通物体随时间划出的方向。我们或许会说一辆车正“向东行驶”，但我们当然知道它同时也在走向未来，因此在时空图上它会划出一条在走向未来时略微向东倾斜的线。

\emph{类空}方向“东”是我们可以画在图上的东西，但没有任何真实物体能沿一条\emph{只}向东、而不向时间推进的时空路径前进。

介于类时与类空方向之间的是光所划出的方向，这些被称为\emph{类光}（lightlike）或\emph{类零}（null）方向。

假设某人使用光速等于 1 的单位制（例如时间用年、距离用光年），测得时空中两个事件相隔距离 $x$、时间 $t$。那么对连接这两个事件的矢量，我们可以说：

\begin{itemize}
  \item 若 $x^{2}-t^{2}$ 为负，则该矢量是\emph{类时}的（例如图中从 A 到 B、A 到 C 的矢量）；
  \item 若 $x^{2}-t^{2}$ 为正，则该矢量是\emph{类空}的（例如图中从 A 到 E 的矢量）；
  \item 若 $x^{2}-t^{2}$ 等于零，则该矢量是\emph{类光}的（例如图中从 A 到 D 的矢量）。
\end{itemize}

狭义相对论的核心发现是：\emph{每一个}作匀速运动的观测者都会测得\emph{相同的} $x^{2}-t^{2}$ 值，因而他们会在把矢量归类为类时、类空或类光这一点上取得一致。使 $x^{2}-t^{2}$ 成为一个\emph{不变量}（invariant，即所有观测者都能认同的量）的这种时空几何，被称为\textbf{洛伦兹}几何。

\orthofig{images/02__002.png}{黎曼时空中的方向}

与之相对，在一个具有\textbf{黎曼}几何的宇宙中，所有人都认同的量是 $x^{2}+t^{2}$。这正是我们按勾股定理（Pythagoras's Theorem）在欧几里得空间中测得的两点间距离的平方。

（如果几何是弯曲的，这只有在 $x$ 与 $t$ 的测量尺度足够小时才成立，正如地球的弯曲表面只有在局限到足够小的区域时才服从欧几里得几何。）

这个平方距离 $x^{2}+t^{2}$ 恒为正（排除 $x$ 与 $t$ 同时为零的平凡情形）。于是旧的类时、类空、类光之别的消失了；在黎曼宇宙中，\emph{所有方向本质上都相同}。某个可被称为“东”的方向，与某个可被称为“未来”的方向之间，不存在任何根本差异。

在洛伦兹情形中，我们利用光速来确定 $x$ 与 $t$ 的测量单位。而在黎曼情形中，并不存在单一而普适的光速。那么我们如何选取测量距离与时间的恰当单位呢？我们只需调整时间单位与距离单位之间的比值，直到通过实验发现勾股定理始终成立——即便它同时涉及时间和空间。本注记中，我们为黎曼物理给出的所有方程，都假定已作出了这样的单位选择。

到目前为止，我们一直通过假定所关注的矢量只张成坐标系中两个方向来简化问题。而且，我们讨论过的矢量都还是\textbf{位移矢量}（displacement vector），它们像是我们可以设想的、从事件 A 指向事件 B 的箭头。本注记中我们还要讨论各种矢量，它们的分量与四个维度相关联，但度量的可能并非距离。因此一般而言，我们把矢量 $\bv$ 的分量记作 $v^{x}$、$v^{y}$、$v^{z}$ 与 $v^{t}$。（注意这些上标只是四个分量的\emph{标号}，与把某物取幂毫无关系。）四维矢量 $\bv$ 的长度平方由勾股定理从二维到四维的推广给出：

\begin{equation}
|\bv|^{2} = (v^{x})^{2} + (v^{y})^{2} + (v^{z})^{2} + (v^{t})^{2}
\end{equation}

即便观测者使用不同的坐标系、测得各个分量 $v^{x}$、$v^{y}$、$v^{z}$ 与 $v^{t}$ 的不同数值，所有人仍会认同矢量 $\bv$ 的长度 $|\bv|$。

\section{点积}

在黎曼几何中，使用不同坐标系的观测者除了认同矢量的长度，还会认同矢量之间的\emph{夹角}。这使我们能够定义一个同时与长度和夹角都相关的、不变量般的量。

设有两个矢量 $\bv$ 与 $\bw$，我们如下定义 $\bv$ 与 $\bw$ 的\textbf{点积}（dot product）：对四个维度 $x$、$y$、$z$、$t$ 中的每一个，把两个矢量的分量相乘，再全部相加。我们把它记作 $\bv\cdot\bw$，于是：

\begin{equation}
\bv\cdot\bw = v^{x}w^{x} + v^{y}w^{y} + v^{z}w^{z} + v^{t}w^{t}
\tag{黎曼点积}
\end{equation}

按此定义，矢量的长度平方就是它与自身的点积：

\begin{equation}
|\bv|^{2} = \bv\cdot\bv
\end{equation}

点积在旋转下保持不变。为看出这一点，考虑一个三角形，其中两条边是矢量 $\bv$ 与 $\bw$；第三条边将是 $\bv-\bw$。第三条边的长度平方为：

\begin{equation}
|\bv-\bw|^{2} = (\bv-\bw)\cdot(\bv-\bw) = |\bv|^{2} + |\bw|^{2} - 2\,\bv\cdot\bw
\end{equation}

因此 $\bv\cdot\bw = \bigl(|\bv|^{2} + |\bw|^{2} - |\bv-\bw|^{2}\bigr)/2$。由于旋转三角形时这些长度都不变，点积也不可能改变。

\orthofig{images/02__003.png}{点积}

若 $\bv$ 与 $\bw$ 之间的夹角为 $\theta$，则：

\begin{equation}
\bv\cdot\bw = |\bv|\,|\bw|\cos\theta
\end{equation}

这并非一个显而易见的论断，但很容易证明。设想最简单的情形：$\bw$ 位于 $x$ 轴正半轴上（于是 $w^{y}=w^{z}=w^{t}=0$，且 $w^{x}>0$），而 $\bv$ 位于 $xt$ 平面内（于是 $v^{y}=v^{z}=0$）。由于众多分量为零，我们有：

\begin{equation}
\bv\cdot\bw = v^{x}w^{x}, \qquad |\bw| = w^{x}
\end{equation}

由余弦的定义，我们又有：

\begin{equation}
v^{x} = |\bv|\cos\theta
\end{equation}

由此可知，在这一特定情形下：

\begin{equation}
\bv\cdot\bw = |\bv|\,|\bw|\cos\theta
\end{equation}

但对于任意给定的 $\bv$ 与 $\bw$，我们总可以旋转它们，使其满足以上所设的限制条件。旋转不会改变矢量的长度、也不改变它们之间的夹角 $\theta$，更不会改变它们的点积；因此若在旋转后结论成立，那么它在旋转之前也必然一直成立。

夹角的余弦在夹角小于 $90^\circ$ 时为正，恰为 $90^\circ$ 时为零，介于 $90^\circ$ 与 $180^\circ$ 之间时为负。于是：

\begin{itemize}
  \item 若 $\bv\cdot\bw > 0$，则 $\bv$ 与 $\bw$ 大致指向\emph{同一}方向（夹角小于 $90^\circ$）；
  \item 若 $\bv\cdot\bw = 0$，则 $\bv$ 与 $\bw$ 相互垂直；
  \item 若 $\bv\cdot\bw < 0$，则 $\bv$ 与 $\bw$ 大致指向\emph{相反}方向（夹角介于 $90^\circ$ 与 $180^\circ$ 之间）。
\end{itemize}

任一矢量 $\bv$ 都可以沿某个单位矢量 $\bu$ 分解为平行与垂直两个部分：

\begin{equation}
\bv = (\bu\cdot\bv)\,\bu + \bigl(\bv - (\bu\cdot\bv)\,\bu\bigr)
\end{equation}

第一部分显然平行于 $\bu$，因为它不过是数值 $\bu\cdot\bv$ 乘以 $\bu$。点积 $\bu\cdot\bv$ \emph{几乎}就是 $\bv$ 在 $\bu$ 上的投影长度；唯一的差别在于，若两矢量大致反向，则 $\bu\cdot\bv$ 为负，此时它是投影长度的相反数。

第二部分与 $\bu$ 垂直，这可以通过它与 $\bu$ 取点积看出：

\begin{equation}
\bu\cdot\bigl(\bv - (\bu\cdot\bv)\,\bu\bigr)
= (\bu\cdot\bv) - (\bu\cdot\bv)(\bu\cdot\bu)
= (\bu\cdot\bv) - (\bu\cdot\bv)
= 0
\end{equation}

我们将沿坐标系四个轴的单位矢量称为 $\be_{x}$、$\be_{y}$、$\be_{z}$、$\be_{t}$；这四个矢量被称为该坐标系的\textbf{基矢量}（basis vector）。于是对任意矢量 $\bv$，它在该坐标系中的分量由它与基矢量的点积给出：

\begin{equation}
v^{i} = \be_{i}\cdot\bv, \qquad i = x,\,y,\,z,\,t
\end{equation}

在黎曼宇宙中，所有这些矢量几何都非常简单：它不过是我们所熟悉的欧几里得空间几何再加上时间这一维。与之相反，在洛伦兹宇宙中，我们必须始终小心时间维度上的符号差异。例如，在洛伦兹宇宙中，两个四维矢量（或“四矢量”）的点积对时间坐标带有负号：

\begin{equation}
\bv\cdot\bw = v^{x}w^{x} + v^{y}w^{y} + v^{z}w^{z} - v^{t}w^{t}
\tag{洛伦兹点积}
\end{equation}

\section{没有类时方向的时间}

一个其几何中并未内蕴任何时间概念的宇宙，可能会让人联想到一幅类似我们宇宙中某次快照的景象：某一个单一的、冻结而静止的瞬间。但稍加思索便会发现，这完全说不通。当我们从全部四个维度来看时，我们自己的宇宙同样也是不变的；对我们而言构成\emph{变化}的，是穿过四维宇宙的、一个个相继的三维切片彼此不同这一事实。因此，如果支配黎曼宇宙内容的物理定律也允许相继的三维切片彼此不同，那么产生一种被感知为\emph{变化}的东西的可能性就同样存在。

\orthofig{images/02__005.png}{由世界线定义时间}

相对论告诉我们，并不存在唯一“正确的”方式把我们自己的宇宙切分为一个个相继的时间瞬间。处于相对运动中的不同人会以不同的方式做这种切分，每个人都把自己划出的\textbf{世界线}（world line，即他在时空中留下的路径或历史）当作时间方向的标记，而把与其垂直的所有方向当作空间方向。

现在，一个黎曼宇宙究竟包含何种物体或许并不显然，但结果它也会充满世界线。（其深层原因在于量子力学，但当我们接下来探讨波的时候，会一窥世界线是如何浮现的。）主要差别在于：洛伦兹宇宙中的世界线\emph{必须}沿类时方向延伸，而在黎曼宇宙中，它们可以沿的方向不受任何限制。

我们通常把时间看作普适的，而非我们自己世界线的一种属性，这是因为日常生活中我们周围一切人与物彼此之间的速度都非常小，于是我们所有的世界线都挤在一个紧密的线束中一起延伸，只是来回微小地穿插，但基本上互相平行。正是\emph{这一点}，在大多数实际情形下赋予了我们对时间的单一而共同的概念。对日常经验而言最重要的，是这种共享的运动，而非底层时空几何把某一类（大得多的！）方向归类为类时这一事实。黎曼宇宙中一个类似的、近乎平行的世界线线束，也会产生同样的实际感受，使线束的方向成为一种有用的共享时间概念。

因此，我们不妨假定自己能想象自身是黎曼宇宙的居民，拥有自己的世界线。正如我们在自己宇宙中所做的那样，我们通常会采用一个时间坐标，其轴指向我们世界线的方向，而空间坐标的轴与其正交。我们还将假定，我们有一套明确的方式，沿着我们的世界线选定某个方向作为“未来”。与我们自己的宇宙一样，这归根结底是熵的问题：未来将是熵增大的那个方向。

任何世界线与我们的世界线平行的物体，我们都会认为它是静止的，因为在我们坐标系中它的空间位置不随时间改变。任何世界线与我们的世界线成一定倾角的物体，我们都会认为它在运动。

\orthofig{images/02__006.png}{四维速度与四维加速度}

我们有自己的时间坐标 $t$，但相对论要求任何物体的时间流逝都沿着它的世界线来度量。我们把这样测得的时间称为\textbf{固有时}（proper time）$\tau$。在黎曼宇宙中，两个事件之间对某物体流逝的固有时，就是该物体世界线的长度——除了单位选择的问题之外，它与沿任何其他曲线测得的长度相同。

现在，考虑一个与物体世界线相切的单位矢量 $\bu$。我们称它为物体的\textbf{四维速度}（four-velocity），因为它是一个描述物体瞬时运动状态的四维矢量。我们可以通过用 $\bu$ 的空间分量除以其时间分量，来提取物体在某个特定坐标系中的普通（三维）速度 $\bv$；在图中，一切都被限制在 $xt$ 平面内，于是我们有：

\begin{equation}
v^{x} = u^{x}/u^{t}
\end{equation}

因为如果物体沿 $\bu$ 所标示的切线继续运动，它将在时间 $u^{t}$ 内行进距离 $u^{x}$。

当然，图中的物体并未保持恒定速度，它在加速。物体的\textbf{四维加速度}（four-acceleration），我们记作 $\mathbf{a}$，定义为物体的四维速度关于固有时 $\tau$ 的变化率。我们写作：

\begin{equation}
\mathbf{a} = \partial_{\tau}\bu
\end{equation}

不必被这里的记号吓到！符号“$\partial$”不过是“变化率”的一种简洁写法，而下标 $\tau$ 意味着我们谈论的是当 $\tau$ 变化时 $\bu$ 所发生的变化。

由于四维速度 $\bu$ 始终是单位矢量，它唯一可能发生的变化是方向上的，而非长度上的，这意味着四维加速度 $\mathbf{a}$ 必须与 $\bu$ 正交。

\orthofig{images/02__007.png}{恒定加速度}

假设我们取一个在我们的坐标系中初始静止的物体，并让它沿我们的 $x$ 方向受到恒定加速。确切地说，我们要求物体加速度的\emph{大小}为常量 $a$，且加速度四矢量始终留在 $xt$ 平面内。物体的世界线将起初平行于我们的 $t$ 轴，然后逐渐向 $x$ 轴倾斜。

在洛伦兹宇宙中，物体的运动会越来越快，但它永远达不到光速。

在牛顿宇宙中，物体的速度没有上限，但在任何有限时刻其速度仍是有限的。

在黎曼宇宙中，物体的世界线将是一段圆弧。这并不应令人惊讶，因为这里涉及的几何与数学非常熟悉——它与离心力的数学大致相同：对一个运动物体施以恒定的侧向拉力，会使它沿圆周路径运动。

我们可以把物体的 $(x,t)$ 坐标、四维速度与四维加速度显式地写成固有时 $\tau$ 的函数：

\begin{align}
x(\tau) &= \tfrac{1}{a}\,(1-\cos a\tau,\ \sin a\tau),\\
\bu(\tau) &= \partial_{\tau}x(\tau) = (\sin a\tau,\ \cos a\tau),\\
\mathbf{a}(\tau) &= \partial_{\tau}\bu(\tau) = a\,(\cos a\tau,\ -\sin a\tau).
\end{align}

于是，在有限时间之后——按我们的坐标是 $1/a$，按物体自身是 $\pi/(2a)$——物体的四维速度将平行于我们的 $x$ 轴，我们便会认为它具有无限大的速度。在那之后，它将在我们的坐标中反向于时间运动！

至于这种恒定加速在黎曼宇宙中\emph{实际上}究竟能在何种条件下维持，将取决于本应施加于物体的力的机制；与我们宇宙中一样，任何这样的过程都会受制于关于能量与熵的严格定律。但作为纯运动学的陈述，图示给出了\emph{如果}恒定加速能够实现时会发生什么。

\section{能量-动量矢量}

假设我们给周围每一个物体都关联一个箭头，其长度等于物体的质量 $m$，方向与物体的世界线相切。我们总选取这些箭头指向与我们个人的“未来”选择大致相符的方向；确切地说，我们总选取它们的指向，使其与我们的时间轴所成的角不超过 $90^\circ$。

\orthofig{images/02__008.png}{能量-动量矢量}

这个箭头就是物体的\textbf{能量-动量矢量}（energy-momentum vector），我们用符号 $\bP$ 表示它。我们已经定义了与物体世界线相切的\emph{单位}矢量，即四维速度 $\bu$，因此可以写出：

\begin{equation}
\bP = m\,\bu
\end{equation}

为简单起见，我们暂且假设物体世界线位于我们的 $xt$ 平面内，于是只需处理两个维度。我们把能量-动量矢量的 $t$ 分量称为物体的\textbf{总能量}（total energy）$E$，把它的 $x$ 分量称为物体的\textbf{动量}（momentum）$p$。于是：

\begin{equation}
P^{t} = E, \qquad P^{x} = p
\end{equation}

\begin{equation}
m^{2} = |\bP|^{2} = (P^{t})^{2} + (P^{x})^{2} = E^{2} + p^{2}
\end{equation}

物体的速度 $v$ 将是其四维速度的 $x$ 分量与 $t$ 分量之比 $u^{x}/u^{t}$。由于能量-动量矢量是四维速度的倍数，我们可以直接从能量-动量矢量的分量得到这个比值：

\begin{equation}
v = p/E
\end{equation}

利用上面最后两个方程，我们有：

\begin{align}
p &= Ev,\\
E^{2} &= m^{2} - p^{2} = m^{2} - E^{2}v^{2},\\
E^{2}(1 + v^{2}) &= m^{2},
\end{align}

从而给出用 $m$ 与 $v$ 表示的 $E$ 与 $p$：

\begin{equation}
E = \frac{m}{\sqrt{1+v^{2}}}, \qquad p = \frac{mv}{\sqrt{1+v^{2}}}
\end{equation}

你可能会认出，这些公式是对通常的相对论公式的一个简单改动，其中 $1+v^{2}$ 取代了 $1-v^{2}$。

利用 $1/\sqrt{1+v^{2}}$ 的泰勒级数，我们可以把这两个公式在 $v$ 较小时近似为：

\begin{equation}
\frac{1}{\sqrt{1+v^{2}}} \approx 1 - \tfrac{1}{2}v^{2},
\qquad
E \approx m - \tfrac{1}{2}mv^{2},
\qquad
p \approx mv
\end{equation}

\orthofig{images/02__009.png}{牛顿、洛伦兹与黎曼物理中的动能}

这里动量的近似值 $mv$，正是普通牛顿物理中物体的动量。而在牛顿物理中，物体的动能为：

\begin{equation}
K_{\text{Newtonian}} = \tfrac{1}{2}mv^{2}
\end{equation}

因此总能量的近似值为：

\begin{equation}
E \approx m - K_{\text{Newtonian}}
\end{equation}

在我们自己的宇宙中，我们认为物体具有一份静质量能量 $mc^{2}$（在 $c=1$ 的单位下等于 $m$），把物体的动能\emph{加}到它上面，即得总能量。而在黎曼宇宙中，我们必须从静质量能量中\emph{减去}动能。我们把精确的黎曼动能 $K$ 定义为：

\begin{equation}
K = m - E = m\left(1 - \frac{1}{\sqrt{1+v^{2}}}\right)
\end{equation}

在牛顿物理与洛伦兹物理中，随着物体速度增大，其动能可以无界增长。在牛顿物理中速度没有上限，因而 $K_{\text{Newtonian}}$ 可以任意大。在洛伦兹物理中

\begin{equation}
K_{\text{Lorentzian}} = m\left(\frac{1}{\sqrt{1-v^{2}}} - 1\right)
\end{equation}

也再次可以任意大。而在黎曼物理中，与牛顿物理一样，你能测得的物体速度没有上限；若物体的世界线垂直于你的世界线，其速度将是无穷大。但无论速度多大，其动能永远不会超过它的静质量 $m$。

那么势能呢？我们可以用通常的方式、借助作用在感兴趣粒子上的任何力来定义势能，于是动能与势能之和将被守恒。稍后我们将深入探讨黎曼版电磁学，看看一种静电势能将呈现何种形式。但我们可以立刻注意到一个奇特的事实。两个粒子彼此束缚所组成的系统，其势能通常比它们自由时更低，因为在束缚态中粒子“被困在一个能量势阱里”，四周是向下倾斜的阱壁。在洛伦兹物理中，这导致“质量亏损”：束缚系统的质量小于其各组成部分之和。在黎曼物理中，由于动能与势能相对于静质量能量具有相反的符号意义，质量亏损将朝相反的方向起作用：束缚系统的质量将\emph{大于}其各组成部分之和。

对于小速度，黎曼物理看上去与牛顿物理别无二致，动能“倒过来”这一事实未必显现。当固定数目的粒子只是彼此弹开时，动能不过是从一个粒子传递给另一个，或转化为势能。但在任何可以产生新粒子的情形中，动能与静质量能量的这种关系会产生惊人的后果：新粒子的产生不是由产生它的系统的动能或势能的\emph{减少}来平衡，而是由\emph{增加}来平衡。我们已经指出，在黎曼物理中，束缚系统的质量大于其各组成部分，因此任何处于束缚态、且能够（比如说）发光的物质，都将面临把它多余的质量转化为光、并在此过程中把自身瓦解的风险！这些后果——以及物质何以能在这种条件下保持稳定这一谜题——在小说中有详细得多的探讨。

\section{守恒定律}

在定义了能量-动量矢量之后，我们应当能够说明，在黎曼宇宙中，能量与动量被某个物体间的相互作用\emph{守恒}意味着什么。我们可以坚持选定一个特定的参考系，在其中给所有相关物体赋予能量-动量矢量——并确保它们都指向该参考系的未来——然后说，如果各物体在相互作用\emph{之前}所有能量-动量矢量之和等于相互作用\emph{之后}之和，那么就满足守恒定律。

\orthofig{images/02__010.png}{能量-动量守恒}

然而，还有一种优雅得多、完全无须援引任何特定参考系的做法。我们只需沿每个物体的世界线画一个箭头，令其长度等于该物体的质量，并选取所有箭头都\emph{背离}相互作用点。如果能量与动量守恒，那么所有这些箭头之和将为零。

确立了守恒律的这种对称表述之后，我们就能看出，它对任何一位真的前来、为相互作用赋予“之前”与“之后”能量-动量矢量的观测者意味着什么。

\orthofig{images/02__011.png}{能量-动量守恒}

原先的箭头全部背离相互作用点，因此观测者认为是与相互作用\emph{之前}的物体相关联的那些箭头，都需要被反向，以使它们指向该观测者的未来。但由于原先所有箭头之和为零，这些被反向的箭头之和将等于其余所有箭头之和。于是对这位观测者而言，无论他们把哪个方向称为“未来”，能量与动量在相互作用前后都是相同的。

观测者该如何处理一条\emph{恰好}垂直于其选定时间轴的世界线呢？在这种情况下，事实上唯一能做的就是任意地决定：要么把所有这类物体都视为从该相互作用中产生，此时它们的箭头保持不变，被当作相互作用\emph{之后}的能量-动量；要么把所有这类物体都视为抵达该相互作用，此时它们的箭头被反向，被当作相互作用\emph{之前}就已存在的能量-动量。无论哪种方式，能量与动量都将被看到是守恒的。

\section{奇异的碰撞}

\orthofig{images/02__012.png}{碰撞一}

\orthofig{images/02__013.png}{碰撞二}

假设黎曼宇宙中两个点状物体发生一次正碰。我们假定这是一次\textbf{完全弹性}（perfectly elastic）碰撞，没有能量转化为其他形式，两个物体都不碎裂，而且两者不会粘在一起。

在这些假定下，如果我们知道参与碰撞的每个粒子的一个能量-动量矢量，那么利用守恒定律就能求出其余矢量。例如在\textbf{右上方}的图中，如果我们知道第一个粒子的\emph{初始}能量-动量矢量 $\bP_{i,1}$ 与第二个粒子的 $\bP_{i,2}$，我们就能算出它们的\emph{末态}能量-动量矢量 $\bP_{f,1}$ 与 $\bP_{f,2}$ 必须是什么。在此图所示的碰撞中，能量与动量守恒，于是在黎曼宇宙中任何安排两个粒子以这种方式碰撞的人，都会预期结果大致如这里所描绘的那样。

在\textbf{右下方}的图中，我们重画了与之前完全相同的碰撞，只不过它被旋转了大约 $60^\circ$，并且我们不得不交换 $\bP_{f,2}$ 箭头的方向，以使它在新的参考系中具有正能量。这仍然是一次完全有效的碰撞，服从守恒定律……但对该图所示时间轴上的观测者而言，现在是三条世界线汇聚到一起，碰撞之后只留下一条。在这种观点下，第二个物体与自身的反物质版本相撞，两者湮灭，释放出的能量被第一个物体吸收，使其速度大大降低（总能量与动能的符号意义相反）。

两位观测者对同一组事件给出不同描述，这在相对论中是常有的事，但真正的转折出现在我们设想某个人\emph{安排}这样一次碰撞意味着什么的时候。在第一种观点中，我们可以说，任何按图示初始能量-动量矢量搭建实验的人，都能用守恒定律预测结果。但在第二种观点中，同一张图却告诉我们：任何成功安排两个粒子以能量-动量矢量 $\bP_{i,1}$ 与 $\bP_{i,2}$ 相撞的人，只有当这次成功伴随着第二个粒子的反物质版本在碰撞\emph{之前}抵达、一同参与碰撞时，才会成功。除非第三种成分也在场，否则没有人能让最初两个粒子以这些特定的速度相撞。

基于在我们自己宇宙中的经验，我们的直觉可能会反驳说，理应存在一种实验，其中只有三个粒子中的两个参与碰撞。但守恒定律说这是\emph{物理上不可能的}。在我们自己的宇宙中，我们习惯于物理定律约束我们能做什么，但从不是以这种方式！

\begin{keybox}[本节的结论]
这是本书最有冲击力的结论之一：在黎曼宇宙中，能量-动量守恒不仅告诉我们“会发生什么”，还会反过来限制“什么能够被安排出来”。同一张时空图存在两种自洽解读——两个粒子弹性散射，或者两个粒子中一个与自身的反物质湮灭——二者都满足守恒定律。因此，某个观测者认为只需两颗粒子的碰撞过程，在另一个观测者看来必然同时伴随着第三个粒子（反物质）的出现；两颗粒子以该方式、该速度碰撞这件事，在物理上不可能单独发生。
\end{keybox}

关于这一主题，尚有给不介意稍高数学水平的读者阅读的进阶材料。

\section{进阶：不同观测者之间的能量比较}

在几何与运动的主页中，我们把物体的能量-动量矢量定义为长度等于物体质量、指向沿物体世界线的矢量。当然，沿任一条线都有\emph{两个}方向，因此我们给出的规则是：选取与描述该物体的观测者的时间轴所成角不超过 $90^\circ$ 的那个方向。

当你致力于在单一参考系中描述一切时，这是一套有用的方案；但如果你有两位处于相对运动的观测者，就会出现一些情形，使得这条方向规则令他们沿物体的世界线选取\emph{相反}的方向，于是两种描述将不相容。

在洛伦兹宇宙中，这个问题不会出现，因为处于相对运动的普通观测者仍然认同任何矢量时间分量的\emph{符号}。我们当然可以想象时间反演的观测者，但他们构成一个独特的类别，对我们而言其\emph{所有}时间分量都取负。但在黎曼宇宙中，如果你在一组速度各异的物体的一批世界线上放置指向你正向时间的箭头，而我以相对于你相当高的速度运动，那么我可能在有些情形下认同你的箭头，在另一些情形下则不认同。如果我们各自把所有能量-动量矢量相加，所得的两个总矢量将毫无关联。我们当然在用不同的坐标系，但那并非问题所在；把一个矢量从一个坐标系变换到另一个坐标系是件直截了当的事。但若把这些撇开，仅把各个能量-动量矢量 $\bP_{i}$ 看作坐标无关的几何对象，你的总和可能是 $\bP_{1}+\bP_{2}+\bP_{3}+\bP_{4}$，而我的却是 $\bP_{1}+\bP_{2}-\bP_{3}-\bP_{4}$，于是它们是完全不同的矢量。

由于沿世界线选取哪个方向确实没有\emph{物理依据}\footnote{如果物体带有电荷，沿其世界线的两个方向就不再不可区分。这一点在黎曼电磁学一节中进一步讨论。}，四维空间中的能量-动量矢量在符号上是不确定的。好消息是，物体的\textbf{应力-能量张量}（stress-energy tensor）——它封装了关于能量与动量的密度和流动的信息——并不依赖于沿世界线选择方向。一个在其静止系中密度为 $\rho$ 的物体的应力-能量张量的坐标矩阵为：

\begin{equation}
T^{ab} = \rho\, u^{a} u^{b}
\end{equation}

其中 $\bu$ 是物体的四维速度。显然，在这里交换 $\bu$ 的符号毫无影响！因此，若一位观测者在其坐标系中计算应力-能量张量，另一位观测者在其自己的坐标系中计算它，两者将通过张量变换的正常几何规则相互关联。

于是，能量-动量守恒体现为：应力-能量张量（对一切事物，包括所有物质与辐射）的散度为零。

话虽如此，选定一个参考系并在该参考系中计算系统的总能量与总动量，仍然没有任何\emph{错}。我们可以在相对论性的哈密顿形式体系中做物理——下文将加以描述——其中一旦选定了时间坐标，我们就着手定义一个作为哈密顿量的正能量。在该形式体系中，总能量与总动量将是守恒量。唯一的错误在于，把这些守恒量合并成一个四维矢量，并指望该四维矢量能通过一次旋转与别人以同样方式得到的矢量相关联。

\section{进阶：对称性}

为简单起见，假设一个四维黎曼宇宙在一切方向上无限且完全平坦。这是一种理想化，但它非常有用；我们自己的宇宙有一种类似的理想化，称为闵可夫斯基时空（Minkowski spacetime），它是狭义相对论和大多数相对论量子物理的基础。我们把这个理想化的宇宙称为\textbf{欧几里得}四维空间，因为它服从欧几里得几何的公设。

[\textbf{注意}：在物理学文献中，“欧几里得”一词常被用于指经由 Wick 转动（Wick rotation）所产生的那些版本的物理定律。那些\emph{不是}适用于\emph{Orthogonal}宇宙的定律。]

我们选取一个事件 O 作为坐标原点，选取四个互相垂直的方向作为 $x$、$y$、$z$、$t$ 轴，并选定一套单位制。于是该宇宙中任何事件都可用一组四元实数 $(x,y,z,t)$ 来指定。所有这样的四元组的集合，数学家称之为 $\mathbb{R}^{4}$。

$\mathbb{R}^{4}$ 是一个矢量空间，对于 $\mathbb{R}^{4}$ 中任何矢量 $\bv=(v^{x},v^{y},v^{z},v^{t})$，我们用标准点积求其长度 $|\bv|$：

\begin{equation}
|\bv|^{2} = \bv\cdot\bv = (v^{x})^{2} + (v^{y})^{2} + (v^{z})^{2} + (v^{t})^{2}
\end{equation}

在一般的弯曲宇宙中，我们不能对两个事件的坐标做加减，但在这个平坦宇宙中，这完全说得通。从事件 $\mathbf{a}=(a^{x},a^{y},a^{z},a^{t})$ 到事件 $\mathbf{b}=(b^{x},b^{y},b^{z},b^{t})$ 的矢量就是 $\mathbf{b}-\mathbf{a}=(b^{x}-a^{x},\,b^{y}-a^{y},\,b^{z}-a^{z},\,b^{t}-a^{t})$，而两事件之间的距离为 $|\mathbf{b}-\mathbf{a}|$。

现在假设我们的欧几里得四维空间中充满了各种物体，它们满足本质上属于几何性质的定律（一个例子就是能量-动量守恒定律）。通过考察该空间的\textbf{对称性}（symmetry）集合——即那些改变物体位置或取向、却使其内禀几何性质（诸如两事件间的距离、两条世界线间的夹角、或两个波前之间的间距）保持不变的运算——我们可以更深入地理解这类几何定律可能是什么。

形式上，我们把函数 $f:\mathbb{R}^{4}\to\mathbb{R}^{4}$ 称为一个对称性，是指对任意两个事件 $\mathbf{a}$ 与 $\mathbf{b}$，$f(\mathbf{a})$ 与 $f(\mathbf{b})$ 之间的距离等于 $\mathbf{a}$ 与 $\mathbf{b}$ 之间的距离：

\begin{equation}
|f(\mathbf{b})-f(\mathbf{a})| = |\mathbf{b}-\mathbf{a}|
\end{equation}

直观上不难想到三类满足此条件的函数，用几何语言描述为：

\begin{itemize}
  \item \emph{平移}（translation）：$f$ 把一切沿固定方向移动固定距离；
  \item \emph{旋转}（rotation）：$f$ 至少固定一个事件，并刚性地把它周围的一切绕之转动；
  \item \emph{反射}（reflection）：反转图形的取向，同时保持距离与夹角不变。
\end{itemize}

平移很容易用数学描述。若 $\mathbf{s}$ 是任意矢量，我们可以定义：

\begin{equation}
T_{\mathbf{s}}(\mathbf{a}) = \mathbf{a} + \mathbf{s}
\end{equation}

$T_{\mathbf{s}}$ 显然是对称性，因为：

\begin{equation}
T_{\mathbf{s}}(\mathbf{b}) - T_{\mathbf{s}}(\mathbf{a}) = \mathbf{b}-\mathbf{a}
\end{equation}

由于平移可以把坐标原点移到任何其他事件，我们将专注于那些\emph{固定原点}的旋转与反射。它们总能与平移组合，给出最一般的对称性。

一个固定原点的对称性 $R$ 将是 $\mathbb{R}^{4}$ 上的线性函数，这意味着我们可以在施加它之前或之后对矢量做加减、数乘，结果相同：

\begin{equation}
R(\bv+\bw) = R(\bv) + R(\bw), \qquad R(s\bv) = s\,R(\bv)
\end{equation}

与 $\mathbb{R}^{4}$ 上任何线性函数相伴的是一个 $4\times4$ 矩阵 $R^{i}{}_{j}$，满足：

\begin{equation}
R(\bv) = R(v^{j}\be_{j}) = R^{i}{}_{j}\,v^{j}\,\be_{i}
\end{equation}

其中矢量分量 $v^{j}$ 与矩阵分量 $R^{i}{}_{j}$ 都是就标准基 $\{\be_{x},\be_{y},\be_{z},\be_{t}\}$ 而言的，并且我们使用了爱因斯坦求和约定。换言之，我们可以用矩阵去乘矢量的分量，从而再现 $\mathbb{R}^{4}$ 这一矢量空间上的任何线性函数。

若 $R$ 是一个固定原点的对称性，它将保持 $\mathbb{R}^{4}$ 上的标准点积不变，因此对任意两个矢量 $\bv$ 与 $\bw$：

\begin{equation}
R(\bv)\cdot R(\bw) = \bv\cdot\bw
\end{equation}

这不过是说 $R$ 将保持矢量的长度以及它们之间的夹角不变，因为这些量都依赖于点积。

令 $\bv=\be_{i}$、$\bw=\be_{j}$，其中 $\be_{i}$ 是 $\mathbb{R}^{4}$ 上标准基的第 $i$ 个基矢量，即一个第 $i$ 分量为 1、其余分量皆为零的四元组。于是我们有：

\begin{equation}
R(\be_{i})\cdot R(\be_{j}) = \be_{i}\cdot\be_{j}
\end{equation}

标准基是正交归一的，故 $\be_{i}\cdot\be_{j}=\delta_{ij}$（克罗内克 $\delta$ 符号，$i=j$ 时为 1、$i\neq j$ 时为 0）……这也正是 $4\times4$ 单位矩阵 $I_{4}$ 的 $i,j$ 分量：

\begin{equation}
R(\be_{i})\cdot R(\be_{j}) = \be_{i}\cdot\be_{j} = (I_{4})^{i}{}_{j}
\end{equation}

更有：

\begin{equation}
R(\be_{i})\cdot R(\be_{j}) = R^{k}{}_{i}R^{k}{}_{j} = (R^{\!\top})^{i}{}_{k}R^{k}{}_{j} = (R^{\!\top}R)^{i}{}_{j}
\end{equation}

其中 $R^{\!\top}$ 是 $R$ 的转置，即把 $R$ 的行变为列所得的矩阵。由于 $R^{\!\top}R$ 与 $I_{4}$ 的所有分量都一致，两者必相等：

\begin{equation}
R^{\!\top}R = I_{4}
\end{equation}

于是 $R^{\!\top}$ 就是 $R$ 的逆。

这给了我们一种刻画欧几里得四维空间所有线性对称性的方法：这类函数矩阵的\emph{转置}同时也是它的\emph{逆}。所有这类 $4\times4$ 矩阵的集合称为 $\mathrm{O}(4)$，即四维\textbf{正交群}（orthogonal group）。

举几个具体的例子，下面这个矩阵：

\begin{equation}
\begin{pmatrix}
\cos\theta & -\sin\theta & 0 & 0\\
\sin\theta & \cos\theta & 0 & 0\\
0 & 0 & \cos\phi & -\sin\phi\\
0 & 0 & \sin\phi & \cos\phi
\end{pmatrix}
\end{equation}

在 $xy$ 平面内旋转角度 $\theta$，在 $zt$ 平面内旋转角度 $\phi$。把它乘以它的转置，得到：

\begin{equation}
\begin{pmatrix}
\cos\theta & \sin\theta & 0 & 0\\
-\sin\theta & \cos\theta & 0 & 0\\
0 & 0 & \cos\phi & \sin\phi\\
0 & 0 & -\sin\phi & \cos\phi
\end{pmatrix}
\begin{pmatrix}
\cos\theta & -\sin\theta & 0 & 0\\
\sin\theta & \cos\theta & 0 & 0\\
0 & 0 & \cos\phi & -\sin\phi\\
0 & 0 & \sin\phi & \cos\phi
\end{pmatrix}
=
\begin{pmatrix}
1 & 0 & 0 & 0\\
0 & 1 & 0 & 0\\
0 & 0 & 1 & 0\\
0 & 0 & 0 & 1
\end{pmatrix}
\end{equation}

一个更简单的例子是下面这个矩阵：

\begin{equation}
\begin{pmatrix}
-1 & 0 & 0 & 0\\
0 & 1 & 0 & 0\\
0 & 0 & 1 & 0\\
0 & 0 & 0 & 1
\end{pmatrix}
\end{equation}

它是一个反转 $x$ 坐标的反射。把它乘以它的转置（也就是同一个矩阵），得到：

\begin{equation}
\begin{pmatrix}
-1 & 0 & 0 & 0\\
0 & 1 & 0 & 0\\
0 & 0 & 1 & 0\\
0 & 0 & 0 & 1
\end{pmatrix}
\begin{pmatrix}
-1 & 0 & 0 & 0\\
0 & 1 & 0 & 0\\
0 & 0 & 1 & 0\\
0 & 0 & 0 & 1
\end{pmatrix}
=
\begin{pmatrix}
1 & 0 & 0 & 0\\
0 & 1 & 0 & 0\\
0 & 0 & 1 & 0\\
0 & 0 & 0 & 1
\end{pmatrix}
\end{equation}

我们可以通过取矩阵的行列式来区分纯旋转与反射。旋转的行列式为 1，反射的行列式为 $-1$。所有转置为其逆\emph{且}行列式为 1 的 $4\times4$ 矩阵的集合称为 $\mathrm{SO}(4)$，即四维\textbf{特殊正交群}（special orthogonal group）。

把正交群 $\mathrm{O}(4)$ 的任意元素与平移组合所得的变换群，称为四维\textbf{欧几里得群}（Euclidean group）$\mathrm{E}(4)$。这个群给出四维欧几里得空间的所有对称性。若我们只限于旋转与平移——略去反射——我们就得到 $\mathrm{SE}(4)$，即\textbf{特殊欧几里得群}（special Euclidean group）。

把包含平移的欧几里得群的一般对称性用矩阵来表示有时会很方便。我们可以采用这样的约定：在 $\mathbb{R}^{4}$ 中一个事件的四个坐标末尾再补上一个坐标 1。按这种记法，我们可以把一个线性运算 $R$ 与一个平移矢量 $\mathbf{s}$ 合并成单个矩阵：

\begin{equation}
\begin{pmatrix}
R & \mathbf{s}\\
\bzero^{\!\top} & 1
\end{pmatrix}
\end{equation}

这是一个 $5\times5$ 矩阵；我们用符号 $R$ 缩写整个 $4\times4$ 矩阵，用符号 $\mathbf{s}$ 表示矩阵最后一列的前四个分量，用符号 $\bzero$ 表示矩阵最后一行中的前四个分量。用末尾带坐标 1 的矢量去乘它，即得：

\begin{equation}
\begin{pmatrix}
R & \mathbf{s}\\
\bzero^{\!\top} & 1
\end{pmatrix}
\begin{pmatrix}
\bv\\
1
\end{pmatrix}
=
\begin{pmatrix}
R\bv + \mathbf{s}\\
1
\end{pmatrix}
\end{equation}

我们可以通过把两个这样的对称性的矩阵相乘，来看出它们相继施加时如何相互作用，其中第一个对称性的矩阵在矩阵乘积中位于右侧：

\begin{equation}
\begin{pmatrix}
R_{2} & \mathbf{s}_{2}\\
\bzero^{\!\top} & 1
\end{pmatrix}
\begin{pmatrix}
R_{1} & \mathbf{s}_{1}\\
\bzero^{\!\top} & 1
\end{pmatrix}
=
\begin{pmatrix}
R_{2}R_{1} & R_{2}\mathbf{s}_{1} + \mathbf{s}_{2}\\
\bzero^{\!\top} & 1
\end{pmatrix}
\end{equation}

\section{进阶：自由经典粒子的哈密顿量与拉格朗日量}

在关于能量-动量矢量的入门一节中，我们考虑了一个速度为 $v$ 的自由经典粒子，并推导出它的总能量 $E$ 与动量 $p$（定义为能量-动量矢量的 $t$ 分量与 $x$ 分量）的公式：

\begin{equation}
E = \frac{m}{\sqrt{1+v^{2}}} \tag{1}
\end{equation}

\begin{equation}
p = \frac{mv}{\sqrt{1+v^{2}}} \tag{2}
\end{equation}

\begin{equation}
E^{2} = m^{2} - p^{2} \tag{3}
\end{equation}

我们现在的目标是用哈密顿力学与拉格朗日力学来描述同一个粒子。这不会告诉我们关于该粒子行为的任何新东西——我们已经知道它只是沿欧几里得四维空间中的一条直线世界线运动——但它将使我们能够选对符号与约定，让一切都自洽地运作，这在我们想要把哈密顿与拉格朗日形式体系用于更精密的目的时会大有帮助。

首先，我们假定已选定一套正交坐标 $x$、$y$、$z$、$t$。我们再假定粒子只沿 $x$ 方向运动来简化分析，于是只需考虑一个空间坐标；推广到三个空间坐标是直截了当的。

在哈密顿力学中，哈密顿量 $H$ 是把系统的总能量表示为 $n$ 个广义坐标 $q^{i}$（$i=1,\dots,n$，代表系统的自由度）、与这些坐标\textbf{共轭}（conjugate）的 $n$ 个动量 $p_{i}$、以及时间 $t$ 的函数。哈密顿方程给出动量与坐标的变化率\footnote{见参考文献 [1]。}：

\begin{equation}
\frac{dp_{i}}{dt} = -\frac{\partial H}{\partial q^{i}} \tag{4}
\end{equation}

\begin{equation}
\frac{dq^{i}}{dt} = \frac{\partial H}{\partial p_{i}} \tag{5}
\end{equation}

对于我们这个沿 $x$ 方向运动的自由粒子，$x$ 是唯一的坐标，将有一个共轭于它的动量 $p$。于是乍看之下，我们或许会以为可以把方程 \textbf{(3)} 中的 $E$ 直接代入得到 $H(x,p,t)=E$：

\begin{equation}
H(x,p,t) = \sqrt{m^{2}-p^{2}} \tag{6}
\end{equation}

而共轭于坐标 $x$ 的动量就简单地是能量-动量矢量的 $x$ 分量。

然而，这种做法行不通！用这套设定，哈密顿第二方程 \textbf{(5)} 会给出：

\begin{equation}
\begin{aligned}
\frac{dx}{dt} &= \frac{\partial H}{\partial p}
= \frac{\partial\bigl[\sqrt{m^{2}-p^{2}}\bigr]}{\partial p}
= -\frac{p}{\sqrt{m^{2}-p^{2}}}
= -v \qquad \textbf{[错！]}
\end{aligned}
\tag{7}
\end{equation}

我们的错误在于，想当然地认为按几何方式定义的动量——即能量-动量矢量的 $x$ 分量——必定与哈密顿形式体系中共轭于 $x$ 的动量相同。

但这很容易修正。如果我们定义一个“哈密顿动量”：

\begin{equation}
p_{H} = -p = -\frac{mv}{\sqrt{1+v^{2}}} \tag{8}
\end{equation}

作为共轭于坐标 $x$ 的动量，那么一切就都漂亮地自洽了：

\begin{equation}
H(x,p_{H},t) = \sqrt{m^{2}-p_{H}^{2}} \tag{9}
\end{equation}

\begin{equation}
\begin{aligned}
\frac{dx}{dt} &= \frac{\partial H}{\partial p_{H}}
= \frac{\partial\bigl[\sqrt{m^{2}-p_{H}^{2}}\bigr]}{\partial p_{H}}
= -\frac{p_{H}}{\sqrt{m^{2}-p_{H}^{2}}}
= v \qquad \textbf{[对！]}
\end{aligned}
\tag{10}
\end{equation}

这看似微不足道，但在此处把符号搞对，比等到我们遇到黎曼热力学中温度可正可负、正确选择未必那么显然的种种奇特之处时再来收拾要容易得多。

哈密顿第一方程 \textbf{(4)} 给出：

\begin{equation}
\frac{dp_{H}}{dt} = -\frac{\partial H}{\partial x} = 0 \tag{11}
\end{equation}

它只不过告诉我们粒子的动量是常量。

接下来我们来看自由粒子的拉格朗日量。拉格朗日量 $L$ 是这样一个量：它关于时间的积分——称为作用量（action）$S$——在粒子轨迹的变分下取极大或极小\footnote{见参考文献 [2]。}。这立刻给了我们一条线索：如果我们令 $L\,dt$ 正比于粒子世界线在时间间隔 $dt$ 内那一段的长度，那么作用量 $S=\int L\,dt$ 将正比于世界线的长度。由于两点之间最短的路径是直线，极小化作用量将产生直线的世界线，这正是自由粒子应有的样子。

拉格朗日量 $L$ 形式上定义为系统广义坐标 $q^{i}$、这些坐标的时间变化率 $\dot{q}^{i}=dq^{i}/dt$、以及时间 $t$ 的函数。对本系统，只有一个坐标 $x$，其时间变化率为 $v$，我们有：

\begin{equation}
\begin{aligned}
L(x,v,t)\,dt &= C\,ds,\\
L(x,v,t) &= C\,\frac{ds}{dt},
\end{aligned}
\tag{12}
\end{equation}

其中 $s$ 是沿粒子世界线的四维空间距离，$C$ 是尚待确定的常数。

对于在我们选定坐标系中以速度 $v$ 运动的粒子，沿其世界线的四维空间距离 $s$ 直接由勾股定理给出：它的平方等于流逝的时间 $t$ 与通过空间所走过的距离 $x=vt$ 的平方之和：

\begin{equation}
\begin{aligned}
s(t) &= \sqrt{t^{2} + v^{2}t^{2}}\\
&= t\sqrt{1+v^{2}}
\end{aligned}
\tag{13}
\end{equation}

这与 \textbf{(12)} 一起给出：

\begin{equation}
\begin{aligned}
\frac{ds}{dt} &= \sqrt{1+v^{2}},\\
L(x,v,t) &= C\sqrt{1+v^{2}}
\end{aligned}
\tag{14}
\end{equation}

一般的欧拉-拉格朗日方程为\footnote{见参考文献 [2]。}：

\begin{equation}
\frac{d}{dt}\left[\frac{\partial L}{\partial \dot{q}^{i}}\right] = \frac{\partial L}{\partial q^{i}} \tag{15}
\end{equation}

对本系统它变为：

\begin{equation}
\begin{aligned}
\frac{d}{dt}\left[\frac{\partial L}{\partial v}\right] &= \frac{\partial L}{\partial x},\\
\frac{d}{dt}\left[\frac{Cv}{\sqrt{1+v^{2}}}\right] &= 0,\\
\left[\frac{C}{m}\right]\frac{dp}{dt} &= 0
\end{aligned}
\tag{16}
\end{equation}

于是无论我们为常数 $C$ 取什么值，这都只会再次告诉我们粒子的动量是常量。不过，我们可以利用哈密顿力学与拉格朗日力学之间的关系来确定 $C$ 的正确取值——在后者中，哈密顿共轭动量是用拉格朗日量定义的\footnote{见参考文献 [3]。}：

\begin{equation}
p_{i} = \frac{\partial L}{\partial \dot{q}^{i}} \tag{17}
\end{equation}

对本系统，利用 \textbf{(8)} 与 \textbf{(14)}，它变为：

\begin{equation}
\begin{aligned}
p_{H} &= \frac{\partial L}{\partial v},\\
-\frac{mv}{\sqrt{1+v^{2}}} &= \frac{Cv}{\sqrt{1+v^{2}}},\\
C &= -m.
\end{aligned}
\end{equation}

于是我们得到拉格朗日量的最终表达式：

\begin{equation}
L(x,v,t) = -m\sqrt{1+v^{2}} \tag{18}
\end{equation}

作为交叉验证，我们可以确认哈密顿量本身与拉格朗日量之间的关系正确无误，一般而言\footnote{见参考文献 [1]。}：

\begin{equation}
H = p_{i}\dot{q}^{i} - L \tag{19}
\end{equation}

对本系统它变为：

\begin{equation}
\begin{aligned}
H &= p_{H}v - L\\
&= v\left(-\frac{mv}{\sqrt{1+v^{2}}}\right) - \left(-m\sqrt{1+v^{2}}\right)\\
&= m\left[\sqrt{1+v^{2}} - \frac{v^{2}}{\sqrt{1+v^{2}}}\right]\\
&= \frac{m}{\sqrt{1+v^{2}}}
\end{aligned}
\tag{20}
\end{equation}

与方程 \textbf{(1)} 一致。

因此，要用哈密顿与拉格朗日形式体系处理黎曼物理，须注意：

\begin{itemize}
  \item 自由粒子的哈密顿量 $H$ 是其总能量，但要用一个“哈密顿动量” $p_{H}$ 表示，而它是通常的、几何定义的动量的\emph{相反数}；
  \item 质量为 $m$ 的自由粒子的作用量 $S$ 等于 $-m$ 乘以它世界线的普通欧几里得长度，而拉格朗日量 $L$ 是这一作用量关于时间的导数。
\end{itemize}

\begin{notebox}[补充：与洛伦兹情形的符号差异]
为便于对照（以下属于标准物理知识的补充说明）：在洛伦兹情形中，自由粒子相应的量是哈密顿量 $H=\sqrt{m^{2}+p^{2}}$（即 $E=\sqrt{m^{2}+p^{2}}$）、拉格朗日量 $L=-m\sqrt{1-v^{2}}$，且共轭动量就是几何动量本身。而在黎曼情形中，根号内 $p^{2}$ 前是\emph{减号}（$H=\sqrt{m^{2}-p^{2}}$），共轭动量是几何动量的\emph{反号}（$p_{H}=-p$），拉格朗日量则是 $L=-m\sqrt{1+v^{2}}$——三个符号差异恰好彼此协调，使哈密顿方程与欧拉-拉格朗日方程给出正确的 $dx/dt=v$。
\end{notebox}

\section{参考文献}

\begin{itemize}
  \item [1] L. D. Landau and E. M. Lifshitz, \emph{Mechanics}, Butterworth-Heinemann, 1976, 第 40 节。
  \item [2] Landau and Lifshitz, \emph{op. cit.}, 第 2 节。
  \item [3] Landau and Lifshitz, \emph{op. cit.}, 第 7 节。
\end{itemize}

\chapter{几何与波}

\providecommand{\tr}{\operatorname{tr}}
\providecommand{\by}{\mathbf{y}}
\providecommand{\bb}{\mathbf{b}}

本章讨论正交宇宙（\emph{Orthogonal}）中的波与对称性。我们首先从理想化的平面波出发，说明在黎曼几何（Riemannian geometry）中"光"的传播矢量不再是零矢量，其空间频率 $\kappa$ 与时间频率 $\nu$ 满足 $\kappa^{2}+\nu^{2}=\nu_{\max}^{2}$；随后导出黎曼标量波方程（RSW）与黎曼矢量波方程（RVW）。进阶部分系统介绍对称群的表示论——从群表示的基本概念、可约与不可约表示、酉表示，到 SO(3)、SO(4) 的双覆盖以及 SU(2) 的自旋-$j$ 表示——为理解正交宇宙中粒子与波所服从的对称性奠定数学基础。

\section{平面波}

考虑一束穿过我们自己所处的洛伦兹宇宙的理想化光波：一根无限长、频率精确定义的正弦波。按惯例把这个频率记作 $\nu$。暂时忽略偏振问题，只用一个数 $A$ 来描述波在任意地点与时刻的强度。

设波沿坐标系的 $x$ 轴传播。在光速取为 1 的单位制下，这列波可表示为

\[
A(x,y,z,t)=A_{0}\sin\!\bigl(2\pi\nu\,(x-t)\bigr)
\]

\orthofig{images/03__001.png}{我们宇宙中的光波}

频率的倒数称为波的\textbf{周期}（period）$\tau$，即完成一个完整周期所需的时间。由于正弦函数在其宗量增加或减少 $2\pi$ 时精确重复，波 $A(x,y,z,t)$ 会在 $t$ 增加或减少 $\tau=1/\nu$ 时重复：

\[
A(x,y,z,t\pm\tau)=A(x,y,z,t)
\]

在光速取 1 的单位制下，光的\textbf{波长}（wavelength）与周期数值相等：$\lambda=1/\nu$。波 $A$ 会在 $x$ 增加或减少一个波长时重复：

\[
A(x\pm\lambda,y,z,t)=A(x,y,z,t)
\]

现在定义一个零矢量（null vector）$\bk$，它在我们坐标系中的分量为

\[
k^{x}=2\pi\nu,\qquad k^{y}=0,\qquad k^{z}=0,\qquad k^{t}=2\pi\nu
\]

我们称它为光波的\textbf{传播矢量}（propagation vector）。于是波的振幅可改写为

\[
A(x,y,z,t)=A_{0}\sin(k^{x}x+k^{y}y+k^{z}z-k^{t}t)
\]

表达式 $k^{x}x+k^{y}y+k^{z}z-k^{t}t$ 是把传播矢量 $\bk$ 的各分量与 $\bx=(x,y,z,t)$ 的相应分量相乘后再相加，这让人联想到黎曼几何中定义的点积，唯一的差别是 $t$ 分量乘积前带了一个负号。事实上它正是 $\bk\cdot\bx$ 的洛伦兹版本，并且是我们的洛伦兹宇宙中所有观测者——无论怎样运动——都会一致同意的量。

\orthofig{images/03__004.png}{红外极限}
\orthofig{images/03__002.png}{黎曼宇宙中的慢光波}
\orthofig{images/03__003.png}{黎曼宇宙中的快光波}
\orthofig{images/03__005.png}{紫外极限}

这提示我们，在\emph{黎曼}宇宙真空中传播的同类波应当取形式

\[
A(x,y,z,t)=A_{0}\sin(\bk\cdot\bx)
\]

其中点积是黎曼点积。在黎曼几何中\emph{不存在零矢量}：四维空间中的所有方向本质上都一样。因此表面上看，$\bk$ 可以是任意矢量。

不过我们总可以选取坐标系使 $\bk$ 落在 $xt$ 平面内，于是 $k^{y}=k^{z}=0$，从而

\[
A(x,y,z,t)=A_{0}\sin(\bk\cdot\bx)=A_{0}\sin(k^{x}x+k^{t}t)
\]

这列波的\textbf{波前}（wavefront）——即"波峰"——出现在 $\sin(\bk\cdot\bx)=1$ 处，也就是 $\bk\cdot\bx=\pi/2+2n\pi$（$n$ 为整数）处。若找到一点使 $\bk\cdot\bx$ 取这样的值，则在四维空间中存在三个与 $\bk$ 垂直的方向，沿它们移动不会改变 $\bk\cdot\bx$，所以波前是一系列等间距的三维区域，且这些区域与 $\bk$ 垂直。在我们的图中只画出四维空间的一张二维切片，于是这些区域表现为一列平行线。

这列波的波长 $\lambda$ 与周期 $\tau$ 是多少？不难看出

\[
A\!\left(x\pm\frac{2\pi}{k^{x}},y,z,t\right)=A(x,y,z,t),\qquad
A\!\left(x,y,z,t\pm\frac{2\pi}{k^{t}}\right)=A(x,y,z,t)
\]

所以波长与周期为

\[
\lambda=\frac{2\pi}{k^{x}},\qquad \tau=\frac{2\pi}{k^{t}}
\]

我们称波长的倒数为\emph{空间频率}（spatial frequency）$\kappa$；它完全类比于时间频率 $\nu$，告诉我们单位距离内可容纳多少个波周期。于是传播矢量 $\bk$ 的分量可以写成

\[
k^{x}=\frac{2\pi}{\lambda}=2\pi\kappa,\qquad k^{y}=0,\qquad k^{z}=0,\qquad k^{t}=\frac{2\pi}{\tau}=2\pi\nu
\]

并且空间频率 $\kappa$、时间频率 $\nu$ 与传播矢量大小 $|\bk|$ 之间满足

\[
|\bk|^{2}=(k^{x})^{2}+(k^{t})^{2}=(2\pi\kappa)^{2}+(2\pi\nu)^{2}
\]

于是

\[
\kappa^{2}+\nu^{2}=\frac{|\bk|^{2}}{4\pi^{2}}
\]

黎曼宇宙中不同速度的观测者都会同意 $|\bk|$（传播矢量大小）的取值，正如洛伦兹宇宙中的不同观测者都会同意与光波相伴的传播矢量是零矢量一样。因此我们假设，黎曼宇宙中相当于"光"的那种物理现象，由一批共享同一 $|\bk|$ 值的波组成。

固定 $|\bk|$ 的几何效果，是固定波前之间的间距——它不是在某一观测者所选定的坐标方向上量得的，而是垂直于波前、沿传播矢量方向量得的。黎曼空间中这个间距恒等于 $2\pi/|\bk|$。

〔我们稍后会看到，在现象的量子力学版本中，$|\bk|$ 的取值与\emph{这些波所对应粒子的静止质量}有关，所以"所有光具有相同 $|\bk|$"其实只是在坚持我们选定某个粒子来扮演我们宇宙中光子的角色。光子没有静止质量，但在黎曼宇宙中既然没有零矢量，也就不存在无静止质量的粒子，因此我们必须选取\emph{某个}非零值。〕

一旦 $|\bk|$ 固定，就可见随着空间频率 $\kappa$ 变大，时间频率 $\nu$ 必然变小，反之亦然，以使其平方和保持恒定。若两个频率之一为零，另一个就达到最大可能值，记作 $\nu_{\max}$，它与 $|\bk|$ 成正比：

\[
\nu_{\max}=\frac{|\bk|}{2\pi}
\]

于是 $\kappa$ 与 $\nu$ 的关系变为

\[
\kappa^{2}+\nu^{2}=\nu_{\max}^{2}
\]

正如我们在讨论双勾股定理时所见，这正是沿两根垂直轴测量的空间频率之间应有的关系。在黎曼宇宙中，时间频率与沿任何其他维度测得的频率没有区别，所以数学完全相同。频率 $\nu_{\max}$ 就是以最直接的方式、沿四维空间中垂直于波前的方向、单位距离内的波周期数。

\orthofig{images/03__006.png}{波长等量与速度的关系}

现在假设一个光脉冲的世界线沿传播矢量延伸，则光的速度 $v$ 等于 $\bk$ 的 $x$ 与 $t$ 分量之比：

\[
v=\frac{k^{x}}{k^{t}}=\frac{\kappa}{\nu}
\]

最慢的光 $v=0$，其波长无限长 $\lambda=\infty$，我们称之为"红外极限"，其时间频率为 $\nu_{\max}$。

最快的光 $v=\infty$，其波长最小 $\lambda=1/\nu_{\max}$，我们称之为"紫外极限"，其时间频率为零，因而周期无限长。

\orthofig{images/03__007.png}{平面波}

但我们凭什么假设光脉冲的世界线沿传播矢量延伸？

到目前为止所描述的波都是\textbf{平面波}（plane wave）：它们的波前是空间中的无限平面，也就是四维空间中无限超平面的切片。它们数学上非常简单，却显然高度理想化——而平面波的振幅当然\emph{不会}集中成任何类似平行于传播矢量的世界线的东西。右图中振幅最大的亮带\emph{垂直}于传播矢量，并且随着时间推移，这些波前沿与传播矢量\emph{相反}的方向移动！

但这些是占据整个四维空间的无限波。若要得到一个\emph{局域化}的波，就需要把不同频率的波叠加起来。这些波相互加强的位置不会以任何单一成分波的波前速度移动；它的速度是系统的群速度，而不是波前的相速度。

我们可以通过考察使两列频率略有差异的平面波保持同相所需的 $x$ 与 $t$ 之间的关系，来求波叠加后的群速度。设第一列波的空间、时间频率为 $\kappa,\nu$，第二列为 $\kappa+\Delta\kappa,\nu+\Delta\nu$，两列波相位保持相同要求

\[
\kappa x+\nu t=(\kappa+\Delta\kappa)x+(\nu+\Delta\nu)t
\]

即

\[
x=-\frac{\Delta\nu}{\Delta\kappa}\,t
\]

所以 $v_{\text{group}}=-\dd\nu/\dd\kappa$。把它用到 $\nu$ 与 $\kappa$ 的关系上：

\[
\nu=\sqrt{\nu_{\max}^{2}-\kappa^{2}},\qquad
v_{\text{group}}=-\frac{\dd\nu}{\dd\kappa}=\frac{\kappa}{\sqrt{\nu_{\max}^{2}-\kappa^{2}}}=\frac{\kappa}{\nu}
\]

这与我们直接从传播矢量得到的 $k^{x}/k^{t}$ 一致。

\orthofig{images/03__008.png}{略微不同的平面波之和}

左图由总共 61 列平面波叠加而成（并非所有标记传播矢量的箭头都足够明亮可见），它给出了黎曼宇宙中一个光脉冲历史的合理图像。更真实的情形是叠加一个连续的平面波谱，而非有限个波。

叠加结果大致平行于平均传播矢量，但它当然不是一条完美的锐利世界线，也不是等宽的"世界管"或"世界带"。不过脉冲随时间逐渐展宽这一事实完全在意料之中：黎曼宇宙中不同频率的光以不同速度传播，所以任何由许多不同频率构成的局域波都会如此逐渐色散。

\section{标量波}

上一节讨论了几种非常简单的波的数学表达式，但我们真正想要的是黎曼宇宙真空中\emph{一切}波——无论形状如何——都满足的方程。

我们自己宇宙中的波满足\textbf{偏微分方程}（partial differential equation），即关于波的强度沿各方向变化率之间的某种关系。若某量 $A$ 依赖于若干变量，例如我们在四维空间某个区域上放置的坐标 $x,y,z,t$，则 $A$ 对其中一个变量的\textbf{偏导数}（partial derivative），就是当只改变该变量、保持其余变量不变时 $A$ 的\emph{变化率}。

例如，考虑传播矢量为 $\bk$ 的黎曼平面波

\[
A=A_{0}\sin(\bk\cdot\bx)=A_{0}\sin(k^{x}x+k^{y}y+k^{z}z+k^{t}t)
\]

若固定 $y,z,t$ 而改变 $x$，$A$ 将按正弦波动。正弦函数对宗量的变化率就是同一宗量的余弦，但由于这里正弦函数的宗量中 $x$ 乘了因子 $k^{x}$，$A$ 的变化率也要乘上该因子；这只是微积分中称为链式法则的简单规则的一个例子。

$A$ 对 $x$ 的偏导数记作 $\partial_{x}A$。这个记号初次使用可能略显唬人，但其含义其实非常简单：设想空间中一条线，线上除 $x$ 外所有坐标都固定，然后考虑所要求导的量如何\emph{沿该线}变化，忽略它在别处的行为。用此记号可写

\[
\partial_{x}A=k^{x}A_{0}\cos(k^{x}x+k^{y}y+k^{z}z+k^{t}t)
\]

如果再对 $x$ 求\emph{这个}量的变化率呢？余弦函数对宗量的变化率是同一宗量的正弦的\emph{相反数}，同时链式法则又给出一个因子 $k^{x}$。二阶偏导数记作 $\partial_{x}^{2}A$，故

\[
\partial_{x}^{2}A=-(k^{x})^{2}A_{0}\sin(k^{x}x+k^{y}y+k^{z}z+k^{t}t)=-(k^{x})^{2}A
\]

对另外三个坐标的二阶偏导数为

\[
\partial_{y}^{2}A=-(k^{y})^{2}A,\qquad \partial_{z}^{2}A=-(k^{z})^{2}A,\qquad \partial_{t}^{2}A=-(k^{t})^{2}A
\]

四者相加得

\[
\partial_{x}^{2}A+\partial_{y}^{2}A+\partial_{z}^{2}A+\partial_{t}^{2}A
=-\bigl[(k^{x})^{2}+(k^{y})^{2}+(k^{z})^{2}+(k^{t})^{2}\bigr]A=-|\bk|^{2}A
\]

这里我们做的，是用在每个坐标方向对 $A$ 求二阶变化率的操作，把波的原始函数乘以一个正比于该方向频率平方的数；然后把\emph{所有}二阶变化率相加，就得到平方和，它等于 $|\bk|^{2}$——一个完全独立于传播矢量具体方向的因子。

现在 $|\bk|=2\pi\nu_{\max}$。为简洁起见（避免到处出现 $2\pi$ 因子），我们采用"角频率"记号，用符号 $\omega$ 表示，它等于普通频率 $\nu$ 的 $2\pi$ 倍。于是 $|\bk|=2\pi\nu_{\max}=\wms$，上式可改写为

\begin{equation}
\partial_{x}^{2}A+\partial_{y}^{2}A+\partial_{z}^{2}A+\partial_{t}^{2}A+\wms^{2}A=0 \tag{RSW}
\end{equation}

这个方程不再涉及原来那个平面波的任何具体细节，它只涉及一个在四维空间上变化的量 $A$，以及一个与波形状无关的数 $\wms$。

我们称此方程为\textbf{黎曼标量波方程}（Riemannian Scalar Wave equation, RSW）。"标量"（scalar）是指在四维空间每一点处的、所有观测者都能一致同意的一个数，这与矢量分量不同——后者依赖于坐标系的选择。此方程是我们自己宇宙物理学中出现的、称为亥姆霍兹（Helmholtz）方程的四维版本。

求导运算是线性的：若 $s$ 为常数，$A$、$B$ 为 $x$ 的函数，则

\[
\partial_{x}(sA)=s\,\partial_{x}A,\qquad \partial_{x}(A+B)=\partial_{x}A+\partial_{x}B
\]

二阶变化率同样如此，乘以常数如 $\wms^{2}$ 当然也如此。所以 RSW 方程是一个\textbf{线性方程}：取两个解相加，或把一个解乘以常数，结果仍是解。既然所有共享同一 $\wms$ 的黎曼平面波都满足 RSW 方程，那么由任意多个这样的平面波叠加而成的波也满足它。例如，上一节末尾由 61 列平面波叠加而成的波就满足 RSW 方程。

RSW 方程有一个——可能是灾难性的——问题。考虑波

\[
A=\sin\!\left[\tfrac{5}{4}\wms x\right]\exp\!\left[\tfrac{3}{4}\wms t\right]
\]

指数函数的变化率仍是指数函数，而二阶变化率给出原函数的\emph{正}倍数，不同于正弦、余弦给出的负倍数。于是

\[
\partial_{x}^{2}A=-\tfrac{25}{16}\wms^{2}A,\qquad \partial_{y}^{2}A=0,\qquad \partial_{z}^{2}A=0,\qquad \partial_{t}^{2}A=\tfrac{9}{16}\wms^{2}A
\]

把这些项相加可见，上面给出的这个波满足 RSW 方程。但它描述的却是一个随时间\emph{指数增长}的波！指数因子之所以可能出现，是因为它二阶变化率的正号被另一个坐标方向上大于 $\nu_{\max}$ 的空间频率所平衡——即 $\sin[(5/4)\wms x]$ 中的角频率 $(5/4)\wms$。仅仅声称"$\nu_{\max}$ 是最大可能频率"并不足以使它成立，而单靠 RSW 方程也无法强制这一限制。但如果允许这样的波存在，那么一个频率足够高的起初极微小的扰动，将迅速增长并吞没其他一切。

有一种办法可以避免这个问题，但它的发现是小说的重要情节，所以我在此不予剧透。

\section{矢量波}

我们宇宙中的光是电磁波。电磁场不是由逐点变化的单一数值刻画，而是涉及\emph{矢量}。因此我们想弄清，在黎曼宇宙中什么样的矢量波才说得通。

设固定一个四维矢量 $\bA_{0}$，它在我们坐标系中的分量为 $A_{0}^{x},A_{0}^{y},A_{0}^{z},A_{0}^{t}$。则可考虑矢量平面波

\[
\bA(x,y,z,t)=\bA_{0}\sin(\bk\cdot\bx)
\]

它与此前考虑的标量平面波非常相似。事实上，只要 $|\bk|=\wms$，$\bA$ 在我们坐标系中的\emph{每一个分量}都满足 RSW 方程。当然这些分量在不同坐标系中会不同……但若它们在一个坐标系中满足 RSW 方程，则在任何其他坐标系中也满足。（注意这里说的是不同的\emph{直角}坐标系，而非球坐标之类。）

这足够简单，但有一个小麻烦需要处理。虽然 $\bA$ 的分量 $A^{x},A^{y},A^{z},A^{t}$ 会随坐标系改变，但从上面写出的矢量波可以构造出一个标量波：

\[
D(x,y,z,t)=\bk\cdot\bA(x,y,z,t)=(\bk\cdot\bA_{0})\sin(\bk\cdot\bx)
\]

这个波是真正的\emph{标量}波——所有观测者无论怎样选取坐标系都会同意它的取值。这看似不重要，但我们在自己宇宙中的所有物理经验都表明，标量波与矢量波将对应不同的现象。在量子力学中，标量波对应自旋为零的粒子，矢量波对应携带某一单位自旋的粒子——光子便是一例。所以我们期望黎曼宇宙中相当于光的现象应是一个纯矢量波，不应附带如 $D$ 这样的标量波。

于是，我们对像 $\bA$ 这样的矢量平面波提出的要求是 $\bk\cdot\bA_{0}=0$。施加此条件意味着描述该波自身的矢量与传播矢量正交，从而把 $\bA_{0}$ 限制在一个与 $\bk$ 正交的三维子空间内。既然该子空间是三维的，黎曼光将有\emph{三}个独立的偏振态，而我们自己宇宙中的光只有两个。等我们详细考察黎曼电磁学时还会再谈这一点。

我们需要把"无标量"条件改写为可施加于任意矢量函数 $\bA(x,y,z,t)$ 的形式。对于平面波 $\bA(x,y,z,t)=\bA_{0}\sin(\bk\cdot\bx)$（其中要求 $\bk\cdot\bA_{0}=0$），可见

\[
\partial_{x}A^{x}+\partial_{y}A^{y}+\partial_{z}A^{z}+\partial_{t}A^{t}
=(A_{0}^{x}k^{x}+A_{0}^{y}k^{y}+A_{0}^{z}k^{z}+A_{0}^{t}k^{t})\cos(\bk\cdot\bx)
=(\bk\cdot\bA_{0})\cos(\bk\cdot\bx)=0
\]

所以我们宣布，下面一对方程构成\textbf{黎曼矢量波方程}（Riemannian Vector Wave equations, RVW）：

\begin{align}
\partial_{x}^{2}\bA+\partial_{y}^{2}\bA+\partial_{z}^{2}\bA+\partial_{t}^{2}\bA+\wms^{2}\bA &= \bzero \tag{RVW}\\[2pt]
\partial_{x}A^{x}+\partial_{y}A^{y}+\partial_{z}A^{z}+\partial_{t}A^{t} &= 0 \tag{Transverse}
\end{align}

其中第一式告诉我们 $\bA$ 的四个分量各自满足 RSW 方程，而\textbf{横波条件}（transverse condition）则保证这个波是纯矢量波。

对于不介意数学水平略高一些的读者，此主题还有进一步的补充材料。

\section{对称群的表示（Representations of the Symmetry Groups）}

在关于对称性的注记中，我们识别了欧几里得四维空间的若干对称群。其中最大的群是四维欧几里得群 $E(4)$，它包含一切保持点间距离的变换。该群之内有若干子群，包括：

\begin{itemize}
  \item 一切\emph{平移}组成的群，记作 $T(4)$；
  \item 正交群 $O(4)$，即一切保持原点不动的\emph{旋转与反射}；
  \item 特殊正交群 $SO(4)$，即一切保持原点不动的\emph{旋转}。
\end{itemize}

这些群中的变换可以从两个角度来看。若欧几里得四维空间中进行的各种过程服从物理定律——无论谈论的是经典粒子碰撞、波从一处传播到另一处，还是某个量子力学过程——那么我们要求的定律的对称性应当保证：把 $E(4)$ 中的任何变换施加于原来的过程，得到的新过程也服从定律。举一个简单的例子：若设想若干个粒子的世界线以服从能量-动量守恒的方式碰撞，那么我们应当能够把这些世界线一起旋转、反射和平移，而得到一个同样服从守恒律的新碰撞。

另一种看待这些对称性的方式，是把它们视为从一个坐标系到另一个坐标系的变换。若我使用某个坐标系、你使用另一个，只要我们都使用直角坐标——即 $x,y,z,t$ 四根轴彼此正交、四根轴采用同一尺度——那么我把 $\mathbb{R}^{4}$ 中的坐标四元组指定给四维空间中的实际物理点的方式会与你的不同，但我们两套坐标之间由一个 $E(4)$ 中的对称变换相联系。若我们描述同样的物理过程、只是给事物指定了不同的坐标，那么我们都应能对观测施加物理定律，并在任何客观问题上得出相同结论。某个碰撞中能量-动量是否守恒，这是个客观事实——我们两人用各自的坐标系描述同一组事件，应当对该问题得出相同结论。

无论从哪个角度，理解欧几里得空间的对称性都有助于我们理解能满足这些对称性的定律类型。而在理解我们自身宇宙的物理时，最有力的工具之一被证明是数学中称为\textbf{表示论}（representation theory）的分支。

\subsection{表示（Representations）}

设有群 $G$ 与矢量空间 $V$。考虑从 $V$ 到 $V$ 的线性函数（常称为"线性算符"），取其中一切有逆的函数，以函数复合为群乘法，可构成一个群，称为 $V$ 上的一般线性群 $GL(V)$。$G$ 在 $V$ 上的一个\textbf{表示}（representation）是一个函数 $\rho:G\to GL(V)$，满足对任意 $g,h\in G$：

\[
\rho(g)\rho(h)=\rho(gh)
\]

换言之，以下两个过程

\begin{itemize}
  \item 分别对 $g$ 和 $h$ 施加 $\rho$，再用 $GL(V)$ 中的群乘法合并结果；
  \item 先在 $G$ 中把 $g$ 与 $h$ 相乘，再用函数 $\rho$ 把 $gh$ 对应到 $GL(V)$ 的元素；
\end{itemize}

必须给出相同的最终结果，$\rho$ 才是 $G$ 的一个\emph{表示}。换一种说法，$\rho$ 是从 $G$ 到 $GL(V)$ 的一个群同态。

有些群本身就被定义（或等价于）$n\times n$ 矩阵群，例如 $SO(3)$ 与 $SO(4)$。此时它们显然在 $\mathbb{R}^{n}$ 与 $\mathbb{C}^{n}$ 上有表示（若矩阵为实值），或仅在 $\mathbb{C}^{n}$ 上有表示（若矩阵为复值），其中矢量空间上的线性算符就是矩阵与矢量的乘法。这称为该群的\textbf{基本表示}（fundamental representation）。

下面是另一种表示的实例。实数集合 $\mathbb{R}$ 以\emph{加法}为群运算构成一个群，记作 $(\mathbb{R},+)$。取二维矢量空间 $\mathbb{R}^{2}$，则 $GL(\mathbb{R}^{2})$ 由一切从 $\mathbb{R}^{2}$ 到 $\mathbb{R}^{2}$ 的可逆线性函数组成，也可以把它们看成有逆的 $2\times2$ 矩阵。定义表示 $\rho_{1}$ 如下：对 $(\mathbb{R},+)$ 中任意数 $r$，令 $\rho_{1}(r)$ 为把 $\mathbb{R}^{2}$ 中矢量旋转角度 $r$ 的矩阵：

\begin{equation}
\rho_{1}(r)=\begin{pmatrix}\cos r & -\sin r\\ \sin r & \cos r\end{pmatrix} \tag{1}
\end{equation}

不难看出 $\rho_{1}(r+s)=\rho_{1}(r)\rho_{1}(s)$，其中左边用的是 $(\mathbb{R},+)$ 中的加法，右边用的是 $GL(\mathbb{R}^{2})$ 中的矩阵乘法。其实无须真正做矩阵乘法就知道：先旋转角度 $s$ 再旋转角度 $r$ 等于旋转角度 $r+s$。

\subsection{可约与不可约表示}

在上面给出的 $(\mathbb{R},+)$ 在 $\mathbb{R}^{2}$ 上的表示 $\rho_{1}$（由式 (1) 定义）中，很明显我们"用满"了整个矢量空间：不存在 $\mathbb{R}^{2}$ 的某个子空间，使得当 $\rho_{1}$ 对群 $(\mathbb{R},+)$ 各元素给出的矩阵作用其上时，其中的矢量仍留在其中。

但现在考虑 $(\mathbb{R},+)$ 在 $\mathbb{R}^{2}$ 上的另一个表示：

\begin{equation}
\rho_{2}(r)=\begin{pmatrix} \exp(r) & 0\\ 0 & \exp(2r)\end{pmatrix} \tag{2}
\end{equation}

当乘以这样的矩阵时，$\mathbb{R}^{2}$ 中位于 $x$ 轴上的矢量仍留在 $x$ 轴上，位于 $y$ 轴上的矢量仍留在 $y$ 轴上。我们把表示 $\rho_{2}$ 称为\textbf{可约}（reducible），因为它实际上是由两个各自一维的表示"粘在一起"，而不是一个真正二维的表示。

与此相对，我们把先前的表示 $\rho_{1}$ 称为\textbf{不可约}（irreducible），因为不存在 $\mathbb{R}^{2}$ 中被所有 $\rho_{1}(r)$ 的作用保持的子空间。〔值得指出，同样的矩阵作用在\emph{复}数对 $\mathbb{C}^{2}$ 上时给出的是一个可约表示；用实数还是复数可以造成根本差别。〕

不可约表示是构造其他表示的基本积木，因此从纯数学角度值得关注。而在物理学中，事实证明每一个\emph{基本粒子}都对应着相关对称群的一个不同的\emph{不可约}表示！在关于矢量波的注记中，当我们添加一个条件以确保正确区分矢量波与标量波时，所做的正是确保我们挑出了一个不可约表示。

\subsection{酉表示（Unitary Representations）}

量子力学用复矢量空间中的矢量描述物理系统的态。量子态的矢量空间总是配备所谓的\textbf{内积}（inner product），它是点积概念向复数的推广。例如，在复数对的矢量空间 $\mathbb{C}^{2}$ 上，标准内积为

\[
\langle(a,b),(c,d)\rangle=a^{*}c+b^{*}d
\]

这与 $\mathbb{R}^{2}$ 上的标准点积相同，只是对第一个矢量的分量取复共轭。（有时内积按对第二个矢量取复共轭来定义，而非第一个。究竟共轭哪一个是约定问题，真正要紧的只是保持一致。）

注意这会给我们一个非负的\emph{实数}，尽管 $\bv$ 的分量可以是任意复数。我们这样使用内积，以给每个矢量指定长度或模 $|\bv|$。例如

\[
|(a,b)|^{2}=\langle(a,b),(a,b)\rangle=a^{*}a+b^{*}b=|a|^{2}+|b|^{2}
\]

其中 $|a|,|b|$ 指复数 $a,b$ 的绝对值。

若一个带内积的复矢量空间同时也是完备的，则称为希尔伯特空间（Hilbert space）。量子力学中使用的复矢量空间总是希尔伯特空间。

若两个单位矢量 $\bv$ 与 $\bw$ 表示某物理系统的两个不同量子态，则系统被制备在态 $\bv$ 时、通过"处于态 $\bw$"的检验的概率为

\[
P(\text{通过 }\bw\text{ 检验}\mid\text{制备为 }\bv)=|\langle\bv,\bw\rangle|^{2}
\]

现在，正如我们称保持点积的矩阵为\emph{正交}矩阵，我们称保持复内积的矩阵或线性算符为\textbf{酉}（unitary）的：

\[
\langle U(\bv),U(\bw)\rangle=\langle\bv,\bw\rangle
\]

如何判断一个矩阵描述的是不是酉算符？对于正交矩阵，条件是 $M^{T}=M^{-1}$，即矩阵的转置等于其逆。矩阵为酉的条件非常相似，只是用厄米伴随代替转置，即转置的复共轭，记作 $M^{*}$，矩阵为酉当且仅当 $M^{*}=M^{-1}$。

群 $G$ 在希尔伯特空间 $V$ 上的表示 $\rho$ 称为酉表示，若对每个 $g\in G$，$\rho(g)$ 都是 $V$ 上的酉线性算符。

对称群的酉表示在量子力学中扮演特殊角色，因为它们在保持内积的同时也保持概率。若你我在四维空间上使用不同的坐标，那么对于某个量子系统的同一组物理态，我们一般会用系统希尔伯特空间 $V$ 中\emph{不同}的复矢量来描述。但 $V$ 上应当存在相应对称群的一个表示，使我把你的矢量换算成我的——若该表示是\emph{酉}的，我们就会对矢量间所有内积、从而对所预言的概率达成一致。

不过实际上，由于概率都形如 $|\cdot|^{2}$，即便我的内积与你的内积相差一个绝对值为 1 的因子，二者仍然一致。绝对值为 1 的复数 $z$（$|z|=1$）称为一个\textbf{相位}（phase），所以我们也可以说我们的内积"相差一个相位"。任何纯虚数的指数其绝对值为 1，故相位常写成 $\exp(\ii\theta)$，$\theta$ 为实数。满足"通常表示定义的右端多乘一个相位"这一关系的函数 $\rho_{P}:G\to GL(V)$：

\[
\rho_{P}(g)\rho_{P}(h)=\exp\bigl(\ii\theta(g,h)\bigr)\,\rho_{P}(gh)
\]

称为 $G$ 在 $V$ 上的\textbf{投影表示}（projective representation）。

在我们自己的宇宙中确实存在一些物理系统服从相关对称群的投影表示而非"真正的"表示，我们也预期黎曼宇宙中会发生同样的事。不过这里不打算对处理投影表示的微妙之处给出严格论述，因为我们主要通过处理\textbf{双覆盖}（double cover）——即出现此问题的两个群的二重覆盖——来大体绕开它们。

\subsection{SO(3) 与 SO(4) 的双覆盖（Double Covers）}

在我们宇宙中，若你把坐标系绕某根轴连续旋转角度 $\theta$，则你用来描述自旋与该轴对齐的电子的态矢量将被乘以相位 $\exp(\ii\theta/2)$。当 $\theta$ 达到 $2\pi$、坐标系回到原处时，电子的态矢量已改变了一个因子 $\exp(\ii\pi)=-1$。这是因为电子的态矢量服从三维旋转群 $SO(3)$ 的\emph{投影表示}，而非真正的表示。

\subsubsection{SO(3) 与 SO(4) 的拓扑}

\orthofig{images/03__009.png}{SO(3) 中的不可收缩回路}
\orthofig{images/03__010.png}{SO(3) 中的可收缩回路}

描绘 $SO(3)$ 的一种方式，是把它想象成一个半径为 $\pi$ 的实心三维球，球心为 $O$。球内任一点 $P$ 可表示一个旋转，旋转角等于距离 $OP$，转轴为线段 $OP$；旋转方向取为：右手拇指从 $O$ 指向 $P$ 时四指卷曲的方向。在此方案下，球面上相对的两点表示相同的旋转，因为绕同一轴朝两个方向旋转 $\pi$ 效果相同。于是该球的任一直径实际上都是 $SO(3)$ 中的一条\emph{回路}：它的两个端点在球面上对应完全相同的旋转，故应当视为同一点，从而使回路闭合。但与普通空间中的回路不同，这条回路无法连续收缩为一点。直径的两个端点可以一起移动，只要始终保持彼此相对，但它们不能靠拢。

$SO(3)$ 中的回路分两类：可以连续形变收缩为一点的\textbf{可收缩}（contractible）回路，以及不能收缩的回路。把绕某轴的旋转角从 0 连续增大到 $2\pi$，得到一条不可收缩回路，它在我们的旋转球图像中就是一条直径——例如\textbf{右上图中的蓝色直线}。然而若继续把旋转角增大到 $4\pi$，就得到一条可收缩回路！\textbf{右下图中}近乎笔直的那对绿线说明了这一点（为清晰起见它们避免重复任何旋转）；沿绿—黄—红的顺序看，可见该路径收缩为一点。

我们说 $SO(3)$ 不是\textbf{单连通}（simply connected）的。在单连通拓扑空间中，从点 $A$ 到点 $B$ 的任意两条连续路径都可以连续地形变为彼此。在 $SO(3)$ 中，从 $A$ 到 $B$ 的路径总有两类，其中一类的路径无法形变为另一类的路径。

给定一个群的投影表示，我们总可以尝试通过把 $\rho_{P}(g)$ 再乘一个相位，把它变成真正的表示；这样做不会改变由矢量间内积得到的概率。但事实证明，只有当群是\emph{单连通}的，我们才能保证相位可以被完全消去。对于 $SO(3)$ 这样的群，我们不得不与某些投影表示打交道。

处理此问题的一个数学上方便的招数，是构造一个新群 $SO(3)^{\sim}$，它由从 $SO(3)$ 中的恒等元到群中各元素的\emph{路径等价类}组成，而不仅仅是群元素本身。若 $p$ 是某条具体路径，用 $[p]$ 表示一切与 $p$ 有相同端点、且可连续形变为 $p$ 的路径所构成的等价类。

若取从 $I_{3}$ 到 $SO(3)$ 中 $g$ 的一切连续路径，然后把可连续形变为彼此的路径视为等价，我们恰好得到从 $I_{3}$ 到 $g$ 的两个路径等价类。所以对 $SO(3)$ 中每个 $g$，新群 $SO(3)^{\sim}$ 都含有两个元素 $g_{1},g_{2}$，它们可视为从 $SO(3)$ 的恒等元连续移动到 $g$ 的两种不同方式。

如何在 $SO(3)^{\sim}$ 中相乘两个元素？形式上，一条穿过 $SO(3)$ 的路径是函数 $p:[0,1]\to SO(3)$，其中 $[0,1]$ 是 0 到 1 之间（含端点）的实数集合。取所要相乘的两个等价类的代表路径 $p_{1}$ 与 $p_{2}$，定义 $[p_{1}]$ 与 $[p_{2}]$ 之积为路径 $p_{3}(t)=p_{1}(t)\,p_{2}(t)$ 的等价类。这同时刻画了旋转的乘法与"路径类的乘法"的细节。

在 $SO(3)$ 中乘以路径类比听起来容易：对固定端点而言，它恰好像二元群 $\{1,-1\}$ 中的乘法，其中任何元素自乘都给出恒等元。

一个具体例子：考虑绕 $z$ 轴旋转角度从 0 到 $2\pi$ 所得到的路径

\[
p(t)=\begin{pmatrix}\cos(2\pi t) & -\sin(2\pi t) & 0\\ \sin(2\pi t) & \cos(2\pi t) & 0\\ 0 & 0 & 1\end{pmatrix}
\]

它是一条从恒等元回到恒等元的回路（即 $p(0)=p(1)=I_{3}$），但它\emph{不能}连续形变为恒等路径

\[
e(t)=\begin{pmatrix}1 & 0 & 0\\ 0 & 1 & 0\\ 0 & 0 & 1\end{pmatrix}
\]

然而 $[p][p]$ 是等价类 $[p^{2}]$，其中

\[
p^{2}(t)=\begin{pmatrix}\cos(4\pi t) & -\sin(4\pi t) & 0\\ \sin(4\pi t) & \cos(4\pi t) & 0\\ 0 & 0 & 1\end{pmatrix}
\]

该路径绕 $z$ 轴旋转角度从 0 到 $4\pi$，而它\emph{能}连续形变为恒等路径。所以 $[p][p]=[p^{2}]=[e]$。

不过，与其完全按上述定义处理 $SO(3)^{\sim}$，不如找一个与之同构的矩阵群来得更简单。暂时把这个尚未识别的矩阵群称为"$G_{3}$"——我们把 $G_{3}$ 与 $SO(3)^{\sim}$ 本身都称为 $SO(3)$ 的\textbf{双覆盖}（double cover）。双覆盖将是单连通群，所以我们只需处理它的真表示；但这样做时，我们实际上同时处理了 $SO(3)$ 的真表示\emph{与}投影表示。

为说明 $G_{3}$ 是 $SO(3)$ 的双覆盖，我们需要一个 2 对 1 的函数 $f_{3}:G_{3}\to SO(3)$，使得对一切 $g_{1},g_{2}\in G_{3}$：

\[
f_{3}(g_{1})f_{3}(g_{2})=f_{3}(g_{1}g_{2})
\]

这表明 $G_{3}$ 中的乘法刻画了 $SO(3)$ 中旋转乘法的含义，是群同态的一个例子。为表明 $G_{3}$ 中的乘法同时刻画了路径类乘法这一（非常简单的）结构，我们要求 $G_{3}$ 中被 $f_{3}$ 映到 $SO(3)$ 中恒等元 $I_{3}$ 的那两个元素：

\[
f_{3}(i_{1})=f_{3}(i_{2})=I_{3}
\]

的乘法表恰好与二元群 $\{1,-1\}$ 相同。换句话说，其中一个是 $G_{3}$ 中的恒等元，另一个自乘则给出恒等元。

注意，对 $SO(3)^{\sim}$ 本身我们很容易构造一个类似的函数 $f_{3^{\sim}}:SO(3)^{\sim}\to SO(3)$：对任意路径 $q:[0,1]\to SO(3)$，定义 $f_{3^{\sim}}([q])=q(1)$，而 $f_{3^{\sim}}$ 映到 $SO(3)$ 中恒等元的两个元素就是上面定义的 $[p]$ 与 $[e]$。

群 $SO(4)$ 也不是单连通的，且具有完全相同的诸多问题，因此我们也想为 $SO(4)$ 找一个双覆盖 $G_{4}$ 以及类似的同态 $f_{4}:G_{4}\to SO(4)$。$SO(4)$ 的拓扑是 $SO(3)\times S^{3}$：即它与 $SO(3)$ 与三维球面的直积同胚。为看出这点，对 $SO(4)$ 中给定元素 $R$，考虑它把 $t$ 方向单位矢量映到的矢量 $R\be_{t}$。一切这样的矢量将构成 $\mathbb{R}^{4}$ 中的单位三维球面，因为显然存在某个旋转把 $\be_{t}$ 映到 $\mathbb{R}^{4}$ 中任意单位矢量。若固定某个单位矢量 $\bu$，则 $SO(4)$ 中把 $\be_{t}$ 映到 $\bu$ 的那些元素 $R$ 构成 $SO(4)$ 的一个子群；不同 $\bu$ 对应的各子群彼此同构，而对 $\bu=\be_{t}$ 的选择，该子群按定义就是 $SO(3)$。这不是关于 $SO(4)$ 拓扑断言的严格证明，但应能使它直观可信。

\subsubsection{构造 2 对 1 同态}

为识别 $SO(3)$ 与 $SO(4)$ 的双覆盖，我们从四个 $2\times2$ 矩阵出发：

\[
H_{x}=\begin{pmatrix}0 & -\ii\\ -\ii & 0\end{pmatrix},\quad
H_{y}=\begin{pmatrix}0 & -1\\ 1 & 0\end{pmatrix},\quad
H_{z}=\begin{pmatrix}-\ii & 0\\ 0 & \ii\end{pmatrix},\quad
H_{t}=\begin{pmatrix}1 & 0\\ 0 & 1\end{pmatrix}
\]

这四个矩阵的行列式都为 1。把其中任意两个相乘，得到的要么还是这个集合中的某个矩阵，要么是其相反数。乘法表如下，其中取最左列的矩阵右乘顶行的矩阵。〔若这个乘法表看起来眼熟，也许是因为它给出了四元数（quaternions）的一个矩阵表示。〕

\begin{center}
\begin{tabular}{c|cccc}
 & $H_{x}$ & $H_{y}$ & $H_{z}$ & $H_{t}$\\\hline
$H_{x}$ & $-H_{t}$ & $H_{z}$ & $-H_{y}$ & $H_{x}$\\
$H_{y}$ & $-H_{z}$ & $-H_{t}$ & $H_{x}$ & $H_{y}$\\
$H_{z}$ & $H_{y}$ & $-H_{x}$ & $-H_{t}$ & $H_{z}$\\
$H_{t}$ & $H_{x}$ & $H_{y}$ & $H_{z}$ & $H_{t}$
\end{tabular}
\end{center}

对 $\mathbb{R}^{4}$ 中任一点 $\bx=(x,y,z,t)$，定义 $H(\bx)$ 为

\[
H(\bx)=x^{\mu}H_{\mu}=xH_{x}+yH_{y}+zH_{z}+tH_{t}
=\begin{pmatrix} t-z\ii & -y-x\ii\\ y-x\ii & t+z\ii\end{pmatrix}
\]

我们可以把 $2\times2$ 复矩阵的集合看成 8 维\emph{实}矢量空间：显然可以相加这样的矩阵、用实数相乘，其中一个基是八个在矩阵四个位置之一处为 1 或 $\ii$、其余为零的矩阵。于是 $H$ 是 $\mathbb{R}^{4}$ 与该 8 维矢量空间中一个四实维子空间之间的同构；我们把这个子空间也称为 $H$。由 $\{H_{x},H_{y},H_{z},H_{t}\}$ 的乘法表可见，$H$ 在矩阵乘法下封闭：$H$ 中任意两元素的矩阵乘积仍在 $H$ 中。

$H(\bx)$ 的行列式是 $\bx$ 的长度的平方：

\[
\det H(\bx)=(t-z\ii)(t+z\ii)-(-y-x\ii)(y-x\ii)=x^{2}+y^{2}+z^{2}+t^{2}=|\bx|^{2}
\]

而 $H(\bx)$ 的迹（对角线元素之和）是 $\bx$ 的 $t$ 坐标的两倍：

\[
\tr H(\bx)=(t-z\ii)+(t+z\ii)=2t
\]

$H(\bx)$ 的厄米伴随也在 $H$ 中。具体地，$H(\bx)^{*}=H((-x,-y,-z,t))$：

\[
H(\bx)^{*}=\begin{pmatrix} t+z\ii & y+x\ii\\ -y+x\ii & t-z\ii\end{pmatrix}
\]

而 $H(\bx)$ 与其厄米伴随之积是 $\bx$ 长度平方乘以 $2\times2$ 单位矩阵：

\[
H(\bx)H(\bx)^{*}=\begin{pmatrix}x^{2}+y^{2}+z^{2}+t^{2} & 0\\ 0 & x^{2}+y^{2}+z^{2}+t^{2}\end{pmatrix}=|\bx|^{2}I_{2}
\]

这意味着当 $|\bx|=1$ 时，$H(\bx)$ 是\emph{酉}的。事实上，一切行列式为 1 的 $2\times2$ 酉矩阵构成的群，即 \textbf{SU(2)}，正是 $H$ 中由 $\mathbb{R}^{4}$ 里单位矢量 $\bu$ 对应的矩阵 $H(\bu)$ 构成的子集。

除了 $\bx$ 的模，我们还可以从 $\bx,\by$ 在 $H$ 中的像提取它们的点积：

\[
\bx\cdot\by=\tfrac{1}{2}\tr\!\bigl(H(\bx)H(\by)^{*}\bigr)
\]

此结果可通过逐项验证 $H$ 矩阵（在上指标表示厄米伴随）得到：

\[
\tfrac{1}{2}\tr\!\bigl(H_{\mu}H^{\nu*}\bigr)=\delta_{\mu}^{\nu}
\]

这也使我们能把映射 $H$ 的逆写为

\[
H^{-1}(h)=\tfrac{1}{2}\tr\!\bigl(hH^{\nu*}\bigr)\,\be_{\nu}
\]

其中 $\{\be_{\nu}\}$ 是 $\mathbb{R}^{4}$ 的坐标基。

设 $g,h\in SU(2)$。我们如下定义函数 $f_{4}:SU(2)\times SU(2)\to SO(4)$——即从 $SU(2)$ 元素\emph{对}的集合到 $SO(4)$ 的函数：

\[
f_{4}(g,h)\,\bx=H^{-1}\!\bigl(g\,H(\bx)\,h^{-1}\bigr)
\]

由于 $g$ 与 $h^{-1}$ 都在 $SU(2)\subseteq H$ 中，乘积 $gH(\bx)h^{-1}$ 也在 $H$ 中，所以 $f_{4}(g,h)$ 是 $\mathbb{R}^{4}$ 上良定义的线性算符。利用 $\det(H(\bx))$ 与 $|\bx|^{2}$ 的关系，以及"矩阵乘积的行列式等于行列式之积"：

\[
|f_{4}(g,h)\bx|^{2}=\det\!\bigl(H(f_{4}(g,h)\bx)\bigr)
=\det\!\bigl(gH(\bx)h^{-1}\bigr)
=\det(g)\det\!\bigl(H(\bx)\bigr)\det(h^{-1})
=\det\!\bigl(H(\bx)\bigr)=|\bx|^{2}
\]

所以 $f_{4}(g,h)$ 把 $\mathbb{R}^{4}$ 中矢量映为长度相同的矢量，从而它要么是旋转、要么是反射。而结果总是旋转，即 $SO(4)$ 的元素。

$f_{4}$ 能给出 $SO(4)$ 的\emph{每一个}元素吗？稍作计算便不难表明，$f_{4}(g,h)$ 能给出六个坐标面内任意角度的旋转。通过复合这种形式的旋转，可以构造出任意旋转。

\begin{center}
\begin{tabular}{lll}
\toprule
$g$ & $h$ & 旋转所在的坐标面（角 $2\theta$）\\
\midrule
$\cos\theta\,H_{t}+\sin\theta\,H_{x}$ & $\cos\theta\,H_{t}+\sin\theta\,H_{x}$ & $yz$ 面\\
$\cos\theta\,H_{t}+\sin\theta\,H_{y}$ & $\cos\theta\,H_{t}+\sin\theta\,H_{y}$ & $xz$ 面\\
$\cos\theta\,H_{t}+\sin\theta\,H_{z}$ & $\cos\theta\,H_{t}+\sin\theta\,H_{z}$ & $xy$ 面\\
$\cos\theta\,H_{t}-\sin\theta\,H_{x}$ & $\cos\theta\,H_{t}+\sin\theta\,H_{x}$ & $xt$ 面\\
$\cos\theta\,H_{t}-\sin\theta\,H_{y}$ & $\cos\theta\,H_{t}+\sin\theta\,H_{y}$ & $yt$ 面\\
$\cos\theta\,H_{t}-\sin\theta\,H_{z}$ & $\cos\theta\,H_{t}+\sin\theta\,H_{z}$ & $zt$ 面\\
\bottomrule
\end{tabular}
\end{center}

我们按显然的方式把 $SU(2)\times SU(2)$ 变成一个群，即把对中各元素分别相乘：$(g_{1},h_{1})(g_{2},h_{2})=(g_{1}g_{2},h_{1}h_{2})$。用此乘法可得

\begin{align*}
f_{4}\bigl((g_{1},h_{1})(g_{2},h_{2})\bigr)\bx
&=f_{4}(g_{1}g_{2},h_{1}h_{2})\bx\\
&=H^{-1}\!\bigl(g_{1}g_{2}H(\bx)h_{2}^{-1}h_{1}^{-1}\bigr)\\
&=H^{-1}\!\bigl(g_{1}H\bigl(H^{-1}(g_{2}H(\bx)h_{2}^{-1})\bigr)h_{1}^{-1}\bigr)\\
&=f_{4}(g_{1},h_{1})f_{4}(g_{2},h_{2})\bx
\end{align*}

同时还有

\[
f_{4}(I_{2},I_{2})=I_{4},\qquad f_{4}(-I_{2},-I_{2})=I_{4}
\]

这表明 $SU(2)\times SU(2)$ 中有两个元素映到 $SO(4)$ 的恒等元，而这两个元素具有与 $\{1,-1\}$ 相同的群结构。这就满足了为表明 $SU(2)\times SU(2)$ 是 $SO(4)$ 的双覆盖所需要的条件。

我们是从讨论 $SO(3)$（而非 $SO(4)$）的双覆盖开始的，但上述一切都能轻易改造成此目的。若如下定义 $f_{3}:SU(2)\to SO(3)$：

\[
f_{3}(g)\,\bx=f_{4}(g,g)\,\bx=H^{-1}\!\bigl(g\,H(\bx)\,g^{-1}\bigr)
\]

则 $f_{3}(g)$ 仍在 $\mathbb{R}^{4}$ 上表现为旋转。但假设我们把 $\bx$ 限制到 $\mathbb{R}^{3}$，即令 $t$ 坐标为 0。由于 $t$ 与 $H(\bx)$ 的迹之间的关系，这意味着 $\tr H(\bx)=0$，再由迹的循环性质可得

\[
\tr H(f_{3}(g)\bx)=\tr\!\bigl(gH(\bx)g^{-1}\bigr)=\tr\!\bigl(g^{-1}gH(\bx)\bigr)=\tr H(\bx)=0
\]

所以 $f_{3}(g)\bx$ 的 $t$ 坐标也为 0。既然一切从 $\mathbb{R}^{4}$ 中 $t=0$ 三维子空间出发的矢量都留在该子空间内，$f_{3}(g)$ 给出三维中的旋转，即 $SO(3)$ 的元素。

总结起来，我们已确立：

\begin{itemize}
  \item $SU(2)$ 是 $SO(3)$ 的双覆盖；
  \item $SU(2)\times SU(2)$ 是 $SO(4)$ 的双覆盖。
\end{itemize}

\subsection{SU(2) 的表示}

既然已看到 $SU(2)$ 与 $SO(3)$、$SO(4)$ 的联系，下一步就是构造 $SU(2)$ 的不可约酉表示。由这些表示，我们可以进一步构造 $SO(3)$ 与 $SO(4)$ 的真表示和投影表示。

\subsubsection{自旋-0}

$SU(2)$ 最简单的不可约酉表示是每个群都有的一维\textbf{平凡表示}（trivial representation）：

\[
\rho_{0}(g)=I_{1}
\]

复数 $\mathbb{C}$ 是一个一复维矢量空间，我们与 $SU(2)$ 中每个元素 $g$ 关联的线性算符是 $1\times1$ 单位矩阵。这称为\textbf{自旋-0}（spin-0）表示，它描述标量在坐标变换下的变换方式：完全不变！

\subsubsection{自旋-$\tfrac12$}

$SU(2)$ 次简单的不可约酉表示是二维的：

\[
\rho_{\half}(g)=g
\]

换言之，我们直接用 $SU(2)$ 元素 $g$，通过通常的矩阵与矢量乘法，在 $\mathbb{C}^{2}$ 上定义线性算符。这称为\textbf{自旋-$\half$}（spin-$\half$）表示。

\subsubsection{自旋-1}

我们可以如下在 $\mathbb{C}^{3}$ 上构造 $SU(2)$ 的一个三维不可约酉表示：

\[
\rho_{1}(g)=f_{3}(g)
\]

我们用来表明 $SU(2)$ 是 $SO(3)$ 双覆盖的函数 $f_{3}$ 给出一个 $SO(3)$ 的元素，但没有理由不能把这个 $3\times3$ 实值矩阵当作复数三元组空间 $\mathbb{C}^{3}$ 上的线性算符。这称为\textbf{自旋-1}（spin-1）表示。

\subsubsection{自旋-$j$}

接下来，我们对任意半整数 $j$ 给出\textbf{自旋-$j$}（spin-$j$）表示的一般构造；这是 $SU(2)$ 的一个 $(2j+1)$ 维不可约酉表示。

从任意两个维数为 $n$ 与 $m$ 的矢量空间 $V$ 与 $W$，可以构造一个新矢量空间：$n\times m$ 矩阵的空间。矩阵可按显然方式相加、用实数或复数相乘。这个新矢量空间称为 $V$ 与 $W$ 的张量积（tensor product），记作 $V\otimes W$。

若 $V$ 有基 $\{\be_{1},\dots,\be_{n}\}$、$W$ 有基 $\{\be'_{1},\dots,\be'_{m}\}$，我们定义 $\be_{i}\otimes\be'_{j}$ 为一个 $n\times m$ 矩阵，除第 $i$ 行第 $j$ 列为 1 外其余皆为零。所有 $\be_{i}\otimes\be'_{j}$ 构成 $V\otimes W$ 的基。再定义 $V$ 中任意矢量 $\bv=v^{i}\be_{i}$ 与 $W$ 中任意矢量 $\bw=w^{j}\be'_{j}$ 之间的张量积为

\[
\bv\otimes\bw=v^{i}w^{j}\,\be_{i}\otimes\be'_{j}
\]

其中重复指标采用爱因斯坦求和约定。也就是说，我们把 $\bv$ 的分量竖排在一个 $n\times m$ 矩阵左侧、把 $\bw$ 的分量横排在该矩阵上方，再把矩阵中位于同一行同一列的矢量分量之积填入。当然 $V\otimes W$ 的一般元素并非都取这种形式，但总能把它写成张量积的线性组合。

可以如法构造任意多个矢量空间的张量积 $U\otimes V\otimes W\otimes\cdots$，得到可看作向三维以上延展的类矩阵数组的空间。若各矢量空间上都定义了内积，则可通过把各个内积相乘来定义张量积上的内积：

\[
\langle\bu\otimes\bv,\bw\otimes\bx\rangle=\langle\bu,\bw\rangle\langle\bv,\bx\rangle
\]

此内积通过在第一宗量上线性、在第二宗量上共轭线性，延拓到\emph{一切}张量对。

现在，我们可以在 $\mathbb{C}^{2}$ 的 $2j$ 重张量积上构造 $SU(2)$ 的一个表示：

\[
\rho_{\half\otimes 2j}(g)\,v_{1}\otimes v_{2}\otimes\cdots\otimes v_{2j}
=\rho_{\half}(g)v_{1}\otimes\rho_{\half}(g)v_{2}\otimes\cdots\otimes\rho_{\half}(g)v_{2j}
=gv_{1}\otimes gv_{2}\otimes\cdots\otimes gv_{2j}
\]

这就定义了 $\rho_{\half\otimes 2j}(g)$ 在 $\mathbb{C}^{2}$ 的 $2j$ 个元素的张量积上的作用，而 $\mathbb{C}^{2}\otimes\mathbb{C}^{2}\otimes\cdots$ 中其他元素的结果由线性性给出。注意

\[
\rho_{\half\otimes 2j}(-g)=(-1)^{2j}\rho_{\half\otimes 2j}(g)
\]

这是 $SU(2)$ 的一个酉表示，但当 $j\ge 1$ 时它\emph{不是}不可约的。例如 $j=1$ 时，$\mathbb{C}^{2}\otimes\mathbb{C}^{2}$ 中由反对称 $2\times2$ 矩阵构成的一维子空间在任何 $g$ 的 $\rho_{\half\otimes\half}(g)$ 作用下保持不变。设反对称矩阵为

\[
\begin{pmatrix}0 & a\\ -a & 0\end{pmatrix}=a\bigl(\be_{1}\otimes\be_{2}-\be_{2}\otimes\be_{1}\bigr)
\]

则

\[
\rho_{\half\otimes\half}(g)\bigl[a(\be_{1}\otimes\be_{2}-\be_{2}\otimes\be_{1})\bigr]
=a\,(g\be_{1})\otimes(g\be_{2})-a\,(g\be_{2})\otimes(g\be_{1})
\]

把 $g=\begin{pmatrix}p & q\\ r & s\end{pmatrix}$ 代入（$g\be_{1}=(p,r)$，$g\be_{2}=(q,s)$），并以 $\det g=ps-rq=1$ 化简，结果仍为 $\begin{pmatrix}0 & a\\ -a & 0\end{pmatrix}$ 的形式，故该子空间被保持。

所以 $\mathbb{C}^{2}$ 上自旋-$\half$ 表示的 $2j$ 重张量积并不能给出不可约表示。但它总是\emph{包含}一个我们感兴趣的特定不可约表示：作用在\textbf{对称}张量子空间上的那一个。正如我们说矩阵对称是指交换任意分量的指标得到同值的另一分量（$M^{ij}=M^{ji}$），我们说张量对称是指交换任意分量的\emph{任意}一对指标得到同值的另一分量，例如 $T^{ijklm}=T^{jiklm}$。但若能交换任意一对指标，也就能对指标作任意排列，故 $n$ 个指标的全部 $n!$ 种排列都必须给出同值的分量。

一个等价的条件是：对称张量可以写成若干项之和，其中每一项是同一组 $n$ 个矢量的张量积在所有排列上的求和。例如 $n=3$ 时，每一项形如

\[
\bu\otimes\bv\otimes\bw+\bv\otimes\bw\otimes\bu+\bw\otimes\bu\otimes\bv
+\bw\otimes\bv\otimes\bu+\bu\otimes\bw\otimes\bv+\bv\otimes\bu\otimes\bw
\]

对应某个矢量三元组 $\bu,\bv,\bw$。这 $n$ 个矢量不必互异；若它们有重复，某些张量积会在和中出现多次。我们可以简单地通过改变 $\be_{2}$ 在矢量组中出现的次数（从 0 到 $n$）来构造 $\mathbb{C}^{2}$ 的 $n$ 重张量积中对称张量的基，共得 $n+1$ 个矢量。例如 $n=3$ 时，$\mathbb{C}^{2}\otimes\mathbb{C}^{2}\otimes\mathbb{C}^{2}$ 中对称张量的一个正交归一基为

\begin{align*}
\Bigl\{&\be_{1}\otimes\be_{1}\otimes\be_{1},\;
\frac{\be_{2}\otimes\be_{1}\otimes\be_{1}+\be_{1}\otimes\be_{2}\otimes\be_{1}+\be_{1}\otimes\be_{1}\otimes\be_{2}}{\sqrt{3}},\\
&\frac{\be_{1}\otimes\be_{2}\otimes\be_{2}+\be_{2}\otimes\be_{1}\otimes\be_{2}+\be_{2}\otimes\be_{2}\otimes\be_{1}}{\sqrt{3}},\;
\be_{2}\otimes\be_{2}\otimes\be_{2}\Bigr\}
\end{align*}

由于自旋-$\half$ 表示的 $2j$ 重张量积 $\rho_{\half\otimes 2j}$ 只是把 $\rho_{\half}(g)$ 施加于张量积中的每个矢量，它作用在这样一项（排列后张量积之和）上的结果仍是同形式的一项。所以 $\rho_{\half\otimes 2j}$ 把对称张量映为对称张量。

我们定义 $SU(2)$ 的\textbf{自旋-$j$ 表示} $\rho_{j}$ 为 $\rho_{\half\otimes 2j}$ \emph{限制到对称张量子空间上}。该子空间维数为 $2j+1$，我们把它记作 $V_{j}$。

$SU(2)$ 的自旋-$j$ 表示是不可约酉表示，正如从其名称可以猜到，它们对应具有 $j$ 个角动量单位的基本粒子或其他量子力学系统。若 $j$ 为整数，该表示称为\textbf{玻色型}（bosonic）；若 $j$ 为半奇数，则称为\textbf{费米型}（fermionic）。

\subsubsection{等价表示}

本节开头给出了自旋-0、自旋-$\half$ 与自旋-1 表示的定义。显然前两个与我们对自旋-$j$ 的一般定义一致，但自旋-1 的定义看起来完全不同。我们说

\[
\rho_{1}(g)=f_{3}(g)
\]

其中 $f_{3}(g)$ 是 $\mathbb{R}^{3}$ 上的线性算符，定义为

\[
f_{3}(g)\bx=H^{-1}\!\bigl(g\,H(\bx)\,g^{-1}\bigr)
\]

我们随后把相关的实 $3\times3$ 矩阵当成 $GL(\mathbb{C}^{3})$ 的元素。但我们对自旋-$j$ 的一般定义在 $j=1$ 时化为

\[
\rho_{1}(g)\,v\otimes w=gv\otimes gw
\]

其中 $\bv,\bw\in\mathbb{C}^{2}$。我们把定义通过线性性延拓到 $\mathbb{C}^{2}\otimes\mathbb{C}^{2}$ 中任意张量，然后限制到对称张量上。这两个定义怎么可能描述同一件事？

它们显然\emph{不是}在描述\emph{完全}同一件事，但我们说两个表示\textbf{等价}（equivalent），若在两个表示所作用的矢量空间之间存在一个可逆线性函数，它与两个表示对易。也就是说，群 $G$ 在矢量空间 $V_{A}$ 上的表示 $\rho_{A}$ \emph{等价于}同一群在矢量空间 $V_{B}$ 上的表示 $\rho_{B}$，若存在可逆线性函数 $T:V_{A}\to V_{B}$，使得对每个 $g\in G$：

\[
T\rho_{A}(g)=\rho_{B}(g)T
\]

另一种写法是：对每个 $g\in G$ 与每个 $\bv\in V_{A}$，

\[
\rho_{A}(g)\bv=T^{-1}\rho_{B}(g)T\bv
\]

无论 $\rho_{A}(g)$ 对 $V_{A}$ 中 $\bv$ 做了什么，若用 $T$ 把 $\bv$ 送到 $V_{B}$、让 $\rho_{B}(g)$ 作用在其像上、再用 $T^{-1}$ 把结果带回 $V_{A}$，最终得到同一矢量。所以 $T$ 是一种翻译器，用来证明 $\rho_{A}(g)$ 与 $\rho_{B}(g)$ 在做同一件事——只是换了名字、换了地方。

按此定义，2 对 1 同态 $f_{3}:SU(2)\to SO(3)$ 与自旋-1 表示 $\rho_{1}$ 等价。

对量子力学角动量有一定经验的读者不妨尝试下面这道练习：

\begin{itemize}
  \item 把 $SU(2)$ 的一般元素 $g$ 写成 $H((x,y,z,t))$ 的形式，并施加条件 $x^{2}+y^{2}+z^{2}+t^{2}=1$。
  \item 计算 $f_{3}$ 为 $g$ 在标准基下给出的 $SO(3)$ 中的矩阵 $U_{\text{dc}}$。
  \item 计算 $\rho_{1}$ 为 $g$ 在正交归一基 $\{\be_{1}\otimes\be_{1},(\be_{1}\otimes\be_{2}+\be_{2}\otimes\be_{1})/\sqrt{2},\be_{2}\otimes\be_{2}\}$ 下给出的 $V_{1}$ 中的矩阵 $U_{1}$。
  \item 求一个矩阵 $T$ 使 $U_{\text{dc}}=T^{-1}U_{1}T$。\textbf{提示}：把 $V_{1}$ 中的基矢量看作对应 $z$ 轴自旋为 $-1,0,1$ 的量子态，把 $\mathbb{R}^{3}$ 中的基矢量看作对应 $x$ 轴、$y$ 轴、$z$ 轴自旋为 0 的量子态。若你会为由前者构造后者的态（对总自旋为 1 的粒子），那就给出了 $T$。
\end{itemize}

\subsubsection{对偶表示}

设群 $G$ 在某个复矢量空间 $V$ 上有表示 $\rho$。用 $V^{*}$ 表示一切从 $V$ 到复数 $\mathbb{C}$ 的线性函数构成的\emph{对偶矢量空间}。我们如下定义 $G$ 在 $V^{*}$ 上的\emph{对偶表示} $\rho^{*}$：

\[
(\rho^{*}(g)f)(v)=f\bigl(\rho(g^{-1})v\bigr),\qquad \forall f\in V^{*},\;v\in V,\;g\in G
\]

这里是通过指定新函数 $\rho^{*}(g)f$ 在任意矢量 $\bv$ 上的取值来定义它。在此定义中我们需要用 $g$ 的\emph{逆}来对 $\bv$ 施加 $\rho$，以便 $\rho^{*}$ 确实构成一个表示：

\begin{align*}
&(\rho^{*}(gh)f)(v)=f\bigl(\rho((gh)^{-1})v\bigr)\\
&\quad=f\bigl(\rho(h^{-1}g^{-1})v\bigr)
=f\bigl(\rho(h^{-1})\rho(g^{-1})v\bigr)\\
&\quad=(\rho^{*}(h)f)\bigl(\rho(g^{-1})v\bigr)
=\bigl(\rho^{*}(g)\rho^{*}(h)f\bigr)(v)
\end{align*}

结果表明，$SU(2)$ 的一切自旋-$j$ 表示都\emph{等价于}各自的对偶表示。

为说明这点，我们从自旋-$\half$ 开始。$n$ 维列维-奇维塔符号（Levi-Civita symbol）写作 $\varepsilon_{ij\cdots k}$（$n$ 个下标），定义为

\[
\varepsilon_{ij\cdots k}=\begin{cases}+1 & \{i,j,\dots,k\}\text{ 是 }\{1,2,\dots,n\}\text{ 的偶排列}\\ -1 & \{i,j,\dots,k\}\text{ 是 }\{1,2,\dots,n\}\text{ 的奇排列}\\ 0 & \text{其他}\end{cases}
\]

我们如下定义同构 $T_{\half}:\mathbb{C}^{2}\to\mathbb{C}^{2*}$：

\[
\bigl(T_{\half}(v)\bigr)(w)=\varepsilon_{ij}v^{i}w^{j}
\]

为确立 $\rho_{\half}$ 等价于 $\rho_{\half}^{*}$，需证明对一切 $g\in SU(2)$ 与 $\bv\in\mathbb{C}^{2}$，$\mathbb{C}^{2*}$ 中以下两个函数相同：

\[
T_{\half}\bigl(\rho_{\half}(g)v\bigr)=\rho_{\half}^{*}(g)T_{\half}(v)
\]

把它们作用于同一个 $\bw\in\mathbb{C}^{2}$，即成为

\[
\varepsilon_{ik}g^{i}_{j}v^{j}w^{k}=\varepsilon_{jm}(g^{-1})^{m}_{k}v^{j}w^{k}
\]

任一 $2\times2$ 矩阵 $A$ 的行列式可写为

\[
\det A=\varepsilon_{iq}A^{i}_{1}A^{q}_{2}
\]

由于 $g\in SU(2)$ 的行列式为 1，故

\[
\varepsilon_{iq}g^{i}_{1}g^{q}_{2}=1
\]

交换这里的下指标只改变符号：

\[
\varepsilon_{iq}g^{i}_{2}g^{q}_{1}=\varepsilon_{qi}g^{i}_{1}g^{q}_{2}=-\varepsilon_{iq}g^{i}_{1}g^{q}_{2}=-1
\]

重复下指标会给出零，因为 $\varepsilon_{iq}$ 反对称而 $g^{i}_{m}g^{q}_{m}$ 对称：

\[
\varepsilon_{iq}g^{i}_{1}g^{q}_{1}=\varepsilon_{iq}g^{i}_{2}g^{q}_{2}=0
\]

综上所述，这表明确维-奇维塔符号 $\varepsilon$ 在 $g$ 变换下不变：

\[
\varepsilon_{iq}g^{i}_{j}g^{q}_{m}=\varepsilon_{jm}
\]

在右端乘以 $(g^{-1})^{m}_{k}$ 得

\[
\varepsilon_{iq}g^{i}_{j}g^{q}_{m}(g^{-1})^{m}_{k}=\varepsilon_{jm}(g^{-1})^{m}_{k}
\;\Longrightarrow\;
\varepsilon_{iq}g^{i}_{j}\delta^{q}_{k}=\varepsilon_{jm}(g^{-1})^{m}_{k}
\;\Longrightarrow\;
\varepsilon_{ik}g^{i}_{j}=\varepsilon_{jm}(g^{-1})^{m}_{k}
\]

把它作用于 $\mathbb{C}^{2}$ 中任意一对矢量 $\bv,\bw$，即得所需结果：

\[
\varepsilon_{ik}g^{i}_{j}v^{j}w^{k}=\varepsilon_{jm}(g^{-1})^{m}_{k}v^{j}w^{k}
\]

所以我们证明了 $SU(2)$ 的自旋-$\half$ 表示通过 $T_{\half}$ 等价于它的对偶。

由于自旋-$j$ 表示是若干份自旋-$\half$ 的对称化张量积搭建而成，我们实际上已经完成了证明 $\rho_{j}$ 也等价于其对偶所需的大部分工作。例如对 $j=1$，我们定义 $T_{1}:\mathbb{C}^{2}\otimes\mathbb{C}^{2}\to(\mathbb{C}^{2}\otimes\mathbb{C}^{2})^{*}$ 为

\[
\bigl(T_{1}(v\otimes w)\bigr)(x\otimes y)=\varepsilon_{ij}v^{i}x^{j}\,\varepsilon_{kl}w^{k}y^{l}
\]

它通过线性性延拓到 $\mathbb{C}^{2}\otimes\mathbb{C}^{2}$ 的一切元素。实际表示空间 $V_{1}$ 由 $\mathbb{C}^{2}\otimes\mathbb{C}^{2}$ 中的对称张量构成，但我们只需把 $T_{1}$ 的定义域限制为 $T_{1}:V_{1}\to V_{1}^{*}$。每个 $\varepsilon$ 在任意行列式为 1 的矩阵 $g$ 变换下不变，这一事实足以表明 $T_{1}$ 与群作用对易：

\[
T_{1}\bigl(\rho_{1}(g)\,v\otimes w\bigr)=\rho_{1}^{*}(g)T_{1}(v\otimes w)
\]

可以选取 $V_{1}$ 的一个基 $\{\bb_{1},\bb_{2},\bb_{3}\}$：

\[
\bb_{1}=\frac{\ii\,(-\be_{1}\otimes\be_{1}+\be_{2}\otimes\be_{2})}{\sqrt{2}},\quad
\bb_{2}=\frac{-(\be_{1}\otimes\be_{1}+\be_{2}\otimes\be_{2})}{\sqrt{2}},\quad
\bb_{3}=\frac{\ii\,(\be_{1}\otimes\be_{2}+\be_{2}\otimes\be_{1})}{\sqrt{2}}
\]

它按内积正交归一：

\[
\langle\bb_{i},\bb_{j}\rangle=\delta_{ij}
\]

而且若我们用 $T_{1}$ 定义一个对称双线性配对，它在另一种意义下也是"正交归一"的：

\[
\bigl(T_{1}(\bb_{i})\bigr)(\bb_{j})=\delta_{ij}
\]

所以由 $T_{1}$ 得到的矢量配对正是基 $\{\bb_{i}\}$ 下的欧几里得点积或度规。我们知道 $\rho_{1}$ 的作用是酉的，保持 $V_{1}$ 上的复内积，而 $\rho_{1}$ 还保持这个欧几里得度规：

\[
\bigl(T_{1}(\rho_{1}(g)\bb_{i})\bigr)\bigl(\rho_{1}(g)\bb_{j}\bigr)=\delta_{ij}
\]

这些基元素可以通过一个映射 $T:V_{1}\to H$ 与矩阵 $\{H_{x},H_{y},H_{z}\}$ 相联系，使得 $H_{i}=T(\bb_{i})$：

\[
T(v\otimes w)=(\sqrt{2})\,v\otimes(w^{T}\varepsilon)
\]

这里我们定义了一个分量就是列维-奇维塔符号 $\varepsilon_{ij}$ 的矩阵 $\varepsilon$，并用行矢量 $\bw^{T}$ 与 $\varepsilon$ 的矩阵乘法产生一个新行矢量。列矢量与行矢量的张量积可看成矩阵，这与"矩阵乘法作为列矢量上的线性运算"的通常观念吻合。我们可以把先前的结果

\[
\varepsilon_{ik}g^{i}_{j}=\varepsilon_{jm}(g^{-1})^{m}_{k}
\]

用矩阵乘法（而非指标记号）表达为

\[
g^{T}\varepsilon=\varepsilon g^{-1}
\]

若我们不把双覆盖同态 $f_{3}$ 描述为作用在 $\mathbb{R}^{3}$ 上，而是描述为作用在 $\{H_{x},H_{y},H_{z}\}$ 张成的矩阵空间上，则有

\[
f_{3}(g)H=gHg^{-1}
\]

而 $\rho_{1}$ 的作用是

\[
\rho_{1}(g)\,v\otimes w=gv\otimes gw
\]

我们可以用 $T$ 表明这两个表示等价：

\begin{align*}
f_{3}(g)T(v\otimes w)
&=gT(v\otimes w)g^{-1}\\
&=(\sqrt{2})\,gv\otimes(w^{T}\varepsilon)g^{-1}
=(\sqrt{2})\,gv\otimes w^{T}g^{T}\varepsilon\\
&=(\sqrt{2})\,gv\otimes(gw)^{T}\varepsilon\\
&=T(gv\otimes gw)
=T\bigl(\rho_{1}(g)\,v\otimes w\bigr)
\end{align*}

\subsection{SO(3) 与 SO(4) 的表示}

\subsubsection{SO(3) 的自旋-$j$ 表示}

对每个\emph{整数} $j$，$SU(2)$ 的玻色型自旋-$j$ 表示给出 $SO(3)$ 的一个不可约、酉的\emph{真}表示，维数为 $2j+1$。

我们在 2 对 1 函数 $f_{3}$ 的帮助下得到这个表示。对 $SO(3)$ 中任一给定旋转 $R$，都有两个互为相反数的 $SU(2)$ 元素，比如 $g$ 与 $-g$，使得 $f_{3}(g)=f_{3}(-g)=R$。然而由于 $j$ 是整数，$2j$ 为偶数，偶数个负号在 $\rho_{j}$ 的定义中相互抵消，给出

\[
\rho_{j}(-g)=(-1)^{2j}\rho_{j}(g)=\rho_{j}(g)
\]

所以选择 $g$ 还是 $-g$ 没有差别，映射 $R\mapsto\rho_{j}(g)$ 给出 $SO(3)$ 的一个维数为 $2j+1$ 的真表示。

而当 $j$ 为\emph{半奇数}时——即费米型情形——选择 $g$ 还是 $-g$ 就有差别了，因为

\[
\rho_{j}(-g)=(-1)^{2j}\rho_{j}(g)=-\rho_{j}(g)
\]

此时从 $\rho_{j}$ 只能得到 $SO(3)$ 的一个\emph{投影}的 $(2j+1)$ 维表示。

\subsubsection{SO(4) 的自旋-$(j,k)$ 表示}

$SO(4)$ 的情况非常类似。首先，对两个半整数 $j$ 与 $k$ 的某种选择，我们用 $SU(2)$ 在维数为 $2j+1$ 与 $2k+1$ 的空间 $V_{j}$、$V_{k}$ 上的自旋-$j$、自旋-$k$ 表示 $\rho_{j}$、$\rho_{k}$，在张量积 $V_{j}\otimes V_{k}$ 上定义 $SU(2)\times SU(2)$ 的一个不可约酉表示 $\rho_{j,k}$：

\[
\rho_{j,k}(g,h)\,v\otimes w=\bigl(\rho_{j}(g)v\bigr)\otimes\bigl(\rho_{k}(h)w\bigr)
\]

我们随后通过线性性把它延拓到 $V_{j}\otimes V_{k}$ 的任意元素。若 $j+k$ 为整数，该表示称为玻色型；若 $j+k$ 为半奇数，则称为费米型。

对 $SO(4)$ 中任一给定旋转 $R$，都有两对 $SU(2)$ 元素，比如 $(g,h)$ 与 $(-g,-h)$，使得 $f_{4}(g,h)=f_{4}(-g,-h)=R$。若 $j+k$ 为\emph{整数}——玻色型情形——则在这两者之间的选择无关紧要，因为

\[
\rho_{j,k}(-g,-h)=(-1)^{2(j+k)}\rho_{j,k}(g,h)=\rho_{j,k}(g,h)
\]

所以此时映射 $R\mapsto\rho_{j,k}(g,h)$ 给出 $SO(4)$ 在 $V_{j}\otimes V_{k}$ 上的一个\emph{真}表示，维数为 $(2j+1)(2k+1)$。

但当 $j+k$ 为\emph{半奇数}时——费米型情形——我们有

\[
\rho_{j,k}(-g,-h)=(-1)^{2(j+k)}\rho_{j,k}(g,h)=-\rho_{j,k}(g,h)
\]

而 $\rho_{j,k}$ 只能给出 $SO(4)$ 的一个\emph{投影}表示。

我们先前指出，从 2 对 1 同态 $f_{3}:SU(2)\to SO(3)$ 在 $\mathbb{R}^{3}$ 上得到的 $SU(2)$ 表示，等价于 $SU(2)$ 在 $V_{1}$ 上的三维自旋-1 表示。那么 2 对 1 同态 $f_{4}:SU(2)\times SU(2)\to SO(4)$ 又如何呢？结果它等价于 $SU(2)\times SU(2)$ 在 $\mathbb{C}^{2}\otimes\mathbb{C}^{2}$ 上的四维自旋-$(\half,\half)$ 表示。

\subsection{四元数与 SO(4) 双覆盖的表示（Quaternions and Representations of the Double Cover of SO(4)）}

$\mathrm{SO}(4)$ 双覆盖最重要的三个表示，都可以用对四元数（quaternion）的运算来非常简洁而自然地描述\textsuperscript{[5][6]}（感谢 John Baez 向我解释如何把旋量（spinor）看作四元数）。

我们已经看到，二对一同态 $f_4:\mathrm{SU}(2)\times\mathrm{SU}(2)\to \mathrm{SO}(4)$ 可以通过构造一个由若干 $2\times2$ 复矩阵组成的四维实空间 $H$ 来定义。事实上，空间 $H$ 就是四元数的一个矩阵表示：一个定义了加法与乘法的四维向量空间，其加法与乘法像实数、复数那样定义，本质区别在于四元数的乘法不可交换。当然矩阵乘法也不可交换——而在 $H$ 中做矩阵乘法，正等同于做四元数乘法。

思考四元数的一种方式，是把它看作复数的推广：$-1$ 的平方根不再只有一个 $i$，而是有三个不同的四元数 $i,j,k$，满足
\[
i^2=j^2=k^2=-1,\qquad ijk=-1 .
\]
若把四元数 $i$ 认同为矩阵 $H_x$、$j$ 认同为 $H_y$、$k$ 认同为 $H_z$、$1$ 认同为 $H_t$，则一切都完美地对应起来：矩阵与四元数具有完全相同的乘法表。

$H$ 中的矩阵之乘法表为
\[
\begin{array}{c|cccc}
\cdot & H_x & H_y & H_z & H_t \\ \hline
H_x & -H_t & H_z & -H_y & H_x\\
H_y & -H_z & -H_t & H_x & H_y\\
H_z & H_y & -H_x & -H_t & H_z\\
H_t & H_x & H_y & H_z & H_t
\end{array}
\]
而四元数之乘法表为
\[
\begin{array}{c|cccc}
\cdot & i & j & k & 1 \\ \hline
i & -1 & k & -j & i\\
j & -k & -1 & i & j\\
k & j & -i & -1 & k\\
1 & i & j & k & 1
\end{array}
\]

我们将把实数四元组空间 $\mathbb{R}^4$、写成“实部加三个 $-1$ 的平方根的线性组合”的四元数、以及 $2\times2$ 矩阵空间 $H$，三者视为可互换：
\[
(a,b,c,d)\in\mathbb{R}^4 \;\longleftrightarrow\; a+bi+cj+dk
\;\longleftrightarrow\; H(a,b,c,d)=a H_t+b H_x+c H_y+d H_z .
\]

给定四元数 $q=a+bi+cj+dk$，其共轭 $q^*$ 为 $a-bi-cj-dk$，且对任意四元数 $p,q$ 有 $(pq)^*=q^*p^*$。取共轭对应于对相应矩阵取厄米（Hermitian）伴随。

四元数 $q$ 的实部 $\operatorname{Re}(q)$ 就是 $a$，它也等于 $\tfrac12(q+q^*)$。这对应于取相应矩阵的迹的一半。

我们定义四元数 $q$ 的长度或范数 $|q|$ 为相应向量的通常欧氏长度，故
\[
|q|^2=a^2+b^2+c^2+d^2=q q^*=q^* q .
\]

我们已指出，当把 $\mathbb{R}^4$ 中的向量 $\mathbf{x}$ 映射到矩阵空间 $H$ 时，$H(\mathbf{x})$ 的行列式等于 $|\mathbf{x}|$ 的平方。因此 $H$ 中的单位向量就是行列式为 $1$ 的 $2\times2$ 复矩阵，事实上它们恰好是矩阵群 $\mathrm{SU}(2)$ 的元素。于是 $\mathrm{SO}(4)$ 的双覆盖——即 $\mathrm{SU}(2)$ 中元素的有序对之集合——同样可以看作\textbf{单位四元数}的有序对之集合。

这使我们能够把二对一同态 $f_4$ 描述为 $\mathrm{SU}(2)\times\mathrm{SU}(2)$ 在四元数向量空间 $H$ 上的一个表示：
\begin{equation}
f_4(g,h)\,q=g\,q\,h^{-1}, \tag{4.2.1}
\end{equation}
其中 $\mathrm{SU}(2)$ 中的 $g,h$ 被理解为单位四元数。如上一节所述，这个二对一同态等价于 $\mathrm{SU}(2)\times\mathrm{SU}(2)$ 在 $\mathbb{C}^2\otimes\mathbb{C}^2$ 上的四维自旋 $(\tfrac12,\tfrac12)$ 表示，也称为\textbf{矢量表示}（vector representation）。除了它是对 $\mathrm{SO}(4)$ 的双覆盖而非 $\mathrm{SO}(4)$ 本身的表示这一点之外，这不过是用另一种话术在谈 $\mathbb{R}^4$ 中向量被 $\mathrm{SO}(4)$ 元素所做的普通旋转。$\mathrm{SO}(4)$ 中的每个旋转都能由某对单位四元数 $(g,h)$ 这样生成；唯一的复杂之处在于 $(-g,-h)$ 也产生完全相同的旋转，这正是 $\mathrm{SU}(2)\times\mathrm{SU}(2)$ 为 $\mathrm{SO}(4)$ 双覆盖所预期的。

到目前为止我们其实并未做任何新事情：不过是把矩阵空间 $H$ 认同为四元数，并在其上做与先前完全相同的运算。然而，不那么显然的是，我们还能证明 $\rho_{\frac12}$（$\mathrm{SU}(2)$ 的自旋 $\tfrac12$ 表示）——从而 $\mathrm{SU}(2)\times\mathrm{SU}(2)$ 的自旋 $(\tfrac12,0)$ 与自旋 $(0,\tfrac12)$ 表示——也可以用四元数来定义。

我们把 $\rho'_{\frac12}$ 定义为 $\mathrm{SU}(2)$ 在 $H$ 上的表示：
\[
\rho'_{\frac12}(g)\,q=g\,q .
\]
这里照例把 $\mathrm{SU}(2)$ 中的 $g$ 理解为单位四元数。我们可以通过任取 $\mathbb{C}^2$ 中一个单位向量 $\mathbf{v}_0$ 来构造 $H$ 与 $\mathbb{C}^2$ 之间的同构 $T$：
\[
T(q)=q\,v_0 ,
\]
其中通过把四元数 $q$ 认同为 $2\times2$ 复矩阵来作 $q$ 与 $\mathbf{v}_0$ 的乘法。

若要用 $T$ 把四元数认同为复向量空间 $\mathbb{C}^2$，我们需要给 $H$ 赋予一个复结构：即 $H$ 上对应 $\mathbb{C}^2$ 中乘以 $i$ 的线性算符 $J$。它的选取依赖于 $\mathbf{v}_0$ 的选择，一个简单的选择是 $\mathbf{v}_0=(0,1)$ 及复结构
\[
J(q)=q\,H_z ,
\]
即右乘 $H_z$。选择 $H_z$ 的原因在于 $\mathbf{v}_0$ 正是 $H_z$ 的本征值为 $i$ 的本征向量：
\[
H_z(0,1)=i\,(0,1).
\]
于是 $J^2(q)=q H_z^2=-q$，且当 $\mathbf{v}_0=(0,1)$ 时有
\[
T(J(q))=q H_z(0,1)=i\,q(0,1)=i\,T(q).
\]
利用 $T$ 容易证明两个表示 $\rho'_{\frac12}$ 与 $\rho_{\frac12}$ 等价：
\[
T(\rho'_{\frac12}(g)q)=g\,q\,v_0=\rho_{\frac12}(g)\,T(q).
\]
还需注意 $\rho'_{\frac12}$ 与复结构 $J$ 对易，因为 $J$ 所做的右乘与 $\rho'_{\frac12}$ 所做的左乘相互对易。

现在我们可以在 $H$ 上定义 $\mathrm{SU}(2)\times\mathrm{SU}(2)$ 的自旋 $(\tfrac12,0)$ 与自旋 $(0,\tfrac12)$ 表示：
\[
\rho'_{(\frac12,0)}(g,h)\,q=g\,q,\qquad
\rho'_{(0,\frac12)}(g,h)\,q=h\,q .
\]
它们分别称为\textbf{左手旋量表示}与\textbf{右手旋量表示}（left-handed／right-handed spinor representation）。我们采用统一记号，把 $H$ 上等价于 $f_4$ 的矢量表示写成
\[
\rho'_{(\frac12,\frac12)}(g,h)\,q=g\,q\,h^{-1}.
\]
若 $v$ 是一个向量（即按 $\rho'_{(\frac12,\frac12)}$ 变换的四元数），$r$ 是一个右手旋量（即按 $\rho'_{(0,\frac12)}$ 变换的四元数），则二者之积 $v\,r$ 将是一个左手旋量，按 $\rho'_{(\frac12,0)}$ 变换：
\begin{align*}
\big(\rho'_{(\frac12,\frac12)}(g,h)v\big)\big(\rho'_{(0,\frac12)}(g,h)r\big)
&=(g\,v\,h^{-1})(h\,r)\\
&=g\,(v\,r)\\
&=\rho'_{(\frac12,0)}(g,h)\,(v\,r).
\end{align*}
这意味着四元数乘法这一运算，充当了“矢量表示与右手旋量表示的张量积”到“左手旋量表示”之间的一个纠缠算符（intertwiner）——即一个从前者到后者、与所有表示都对易的线性映射。

同理，从矢量 $v$ 与左手旋量 $l$ 出发也能做类似的事，但必须把它们以 $v^* l$ 的形式组合：
\begin{align*}
\big(\rho'_{(\frac12,\frac12)}(g,h)v\big)^*\big(\rho'_{(\frac12,0)}(g,h)l\big)
&=(g\,v\,h^{-1})^*(g\,l)\\
&=(h\,v^*\,g^{-1})(g\,l)\\
&=h\,(v^* l)\\
&=\rho'_{(0,\frac12)}(g,h)\,(v^* l).
\end{align*}
可见 $v^* l$ 按右手旋量变换。

当把 $\mathbb{R}^4$ 映射到 $H$ 时，我们看到可以在 $H$ 上放置一个实内积：
\[
x\cdot y=\tfrac12\operatorname{tr}\!\big(H(x)H(y)^*\big).
\]
它可以改写成
\[
p\cdot q=\operatorname{Re}(p^* q)=\operatorname{Re}(p\,q^*)=\tfrac12(p\,q^*+q\,p^*).
\]
但现在把 $\mathbb{C}^2$ 映射到 $H$，我们真正需要的是 $H$ 上的一个\emph{复}内积，可定义为
\[
\langle p,q\rangle_S=\operatorname{Re}(p^* q)-\operatorname{Re}(p^* J(q))\,i
=\operatorname{Re}(p^* q)-\operatorname{Re}(p^* q H_z)\,i .
\]
下标 “$S$” 代表旋量（spinor）。容易看出，用任意单位四元数 $g\in\mathrm{SU}(2)$ 左乘都保持此内积，因为 $g^* g=|g|^2=1$：
\begin{align*}
\langle gp,gq\rangle_S
&=\operatorname{Re}\big((gp)^* gq\big)-\operatorname{Re}\big((gp)^* gq H_z\big)i\\
&=\operatorname{Re}\big(p^* g^* gq\big)-\operatorname{Re}\big(p^* g^* gq H_z\big)i\\
&=\langle p,q\rangle_S .
\end{align*}
因此我们关于 $\mathrm{SU}(2)$ 的表示 $\rho'_{\frac12}$ 在此复内积下是酉算符。我们还有
\begin{align*}
\langle gp,q\rangle_S
&=\operatorname{Re}\big((gp)^* q\big)-\operatorname{Re}\big((gp)^* q H_z\big)i\\
&=\operatorname{Re}\big(p^* g^* q\big)-\operatorname{Re}\big(p^* g^* q H_z\big)i\\
&=\langle p,g^* q\rangle_S .
\end{align*}
这告诉我们：若把左乘 $g$ 看作线性算符，则关于此内积的伴随算符就是左乘共轭 $g^*$。

\subsection{宇称（Parity）}

正交群 $\mathrm{O}(4)$ 由两部分组成：行列式为 $1$ 的正交矩阵，它们构成群 $\mathrm{SO}(4)$；以及行列式为 $-1$ 的正交矩阵，其中包含反射。从拓扑上看，这是分离的两块；必然如此，因为行列式是矩阵分量的多项式，因而是一个连续函数。它不可能在连通集合中间从 $1$ 跳到 $-1$，所以 $\mathrm{O}(4)$ 中行列式为 $1$ 与 $-1$ 的部分必然不连通。

[在我们自己的宇宙中，情况要复杂得多：洛伦兹变换群由\emph{四}块不连通的部分组成，因为除了空间反射以外还必须考虑时间反演，以及反射与时间反演的组合。但在黎曼宇宙中，时间反演不过是另一种反射。]

$\mathrm{O}(4)$ 的这两块可以通过与任意行列式为 $-1$ 的 $\mathrm{O}(4)$ 元素相乘而相互映射。我们取一个特定的线性算符，称为\textbf{宇称操作}（parity operation）$P$，其定义为
\begin{equation}
P(x,y,z,t)=(-x,-y,-z,t). \tag{4.2.2}
\end{equation}
若 $R\in\mathrm{SO}(4)$，则 $\det(R)=1$，而 $\det(PR)=\det(P)\det(R)=-1$。反之，任意 $\det(F)=-1$ 的 $F\in\mathrm{O}(4)$ 都可以与 $P$ 相乘得到 $\mathrm{SO}(4)$ 的元素。因此与 $P$ 相乘给出 $\mathrm{O}(4)$ 两个连通分量之间的一个可逆映射。显然 $P^2=I_4$，所以将此映射作用两次便回到出发点。

现在，我们想把 $\mathrm{SU}(2)\times\mathrm{SU}(2)$（即 $\mathrm{SO}(4)$ 的双覆盖）\emph{扩展}为一个构成 $\mathrm{O}(4)$ 双覆盖的新群。模仿 $\mathrm{SO}(4)$ 与 $\mathrm{O}(4)$ 的关系，我们通过加上一个 $\mathrm{SU}(2)\times\mathrm{SU}(2)$ 的“额外副本”来做这件事；这个副本是把 $\mathrm{SU}(2)\times\mathrm{SU}(2)$ 的元素与一个我们将称之为 $Q$ 的新群元素相乘而得到的。目标是让扩展群的第二部分充当 $\mathrm{O}(4)$ 中含反射部分的双覆盖。

我们把新群的元素要么写成 $(g,h)$（$\mathrm{SU}(2)\times\mathrm{SU}(2)$ 的元素），要么写成 $Q\cdot(g,h)$（群第二部分的元素）。后者只是一个无法进一步简化的形式表达式；这里 $Q$ 起到一种标记的作用。

除了扩展 $\mathrm{SU}(2)\times\mathrm{SU}(2)$，我们还想把二对一同态 $f_4:\mathrm{SU}(2)\times\mathrm{SU}(2)\to\mathrm{SO}(4)$ 扩展为从新群到 $\mathrm{O}(4)$ 的同态。我们把 $f_4$ 定义为
\[
f_4(g,h)\,x=H^{-1}\!\big(g\,H(x)\,h^{-1}\big),
\]
并作如下宣布来扩展它：
\[
f_4(Q)\,x=P\,x .
\]
先前我们注意到，对 $H(\mathbf{x})$ 取厄米伴随的效果是
\[
H(x)^*=H(Px),
\]
故有
\[
Px=H^{-1}\!\big(H(x)^*\big).
\]
这一点，加上要求扩展后的 $f_4$ 是群同态，告诉我们：
\begin{align*}
f_4\big(Q\cdot(g,h)\big)x
&=f_4(Q)\,f_4(g,h)\,x\\
&=P\,f_4(g,h)\,x\\
&=P\,H^{-1}\!\big(g\,H(x)\,h^{-1}\big)\\
&=H^{-1}\!\big((g\,H(x)\,h^{-1})^*\big)\\
&=H^{-1}\!\big(h^{*-1}H(x)^*\,g^*\big)\\
&\qquad\text{[矩阵乘积的伴随等于各矩阵伴随按相反顺序之积]}\\
&=H^{-1}\!\big(h\,H(Px)\,g^{-1}\big)\\
&\qquad\text{[因 $g,h\in\mathrm{SU}(2)$，其伴随即其逆]}\\
&=f_4(h,g)\,P\,x\\
&=f_4(h,g)\,f_4(Q)\,x\\
&=f_4\big((h,g)\cdot Q\big)x .
\end{align*}
由此我们可以断定
\begin{align*}
f_4\big(Q\cdot(g,h)\cdot Q^{-1}\big)
&=f_4\big(Q\cdot(g,h)\big)f_4(Q^{-1})\\
&=f_4\big((h,g)\cdot Q\big)f_4(Q^{-1})\\
&=f_4(h,g).
\end{align*}
由于 $f_4$ 是二对一的，且它映到同一点的两个元素总是相反数，于是推出 $Q\cdot(g,h)\cdot Q^{-1}=\pm(h,g)$。但对于 $(g,h)=(I_2,I_2)$（它是 $\mathrm{SU}(2)\times\mathrm{SU}(2)$ 的单位元，从而也是整个扩展群的单位元）的情形，唯一正确的选择是正号。又因 $\mathrm{SU}(2)\times\mathrm{SU}(2)$ 连通，且我们期望用 $Q$ 作共轭是一个连续映射，符号不能跳变。所以
\begin{equation}
Q\cdot(g,h)\cdot Q^{-1}=(h,g),\qquad
Q\cdot(g,h)=(h,g)\cdot Q . \tag{4.2.3}
\end{equation}
我们还有
\[
f_4(Q^2)=f_4(Q)^2=P^2=I_4,\qquad Q^2=\pm(I_2,I_2).
\]
两种选择在数学上都成立；哪一个为真只能由实验决定。在我们自己的宇宙中，结果发现宇称根本不是精确对称，因为 $\beta$ 衰变实验显示宇称被破坏。然而它仍是一个有用的概念，适用于广泛的现象，尽管它并非对全部物理都是完美对称。

若把 $\mathrm{O}(3)$ 视为 $\mathrm{O}(4)$ 中使时间轴固定的子群，并把 $\mathrm{SU}(2)$ 视为 $\mathrm{SU}(2)\times\mathrm{SU}(2)$ 中由形如 $(g,g)$ 的元素构成的子群，则以上所做的一切都可改造为 $\mathrm{O}(3)$ 的一个双覆盖。注意元素 $Q$ 与 $\mathrm{SU}(2)$ 的每个元素对易：
\[
Q\cdot(g,g)=(g,g)\cdot Q .
\]

\subsubsection{含宇称的表示（Representations Including Parity）}

设我们取 $\mathrm{SU}(2)\times\mathrm{SU}(2)$ 的某一个自旋 $(j,k)$ 表示，并想把它扩展到含 $Q$ 的群上。先前我们注意到，当 $Q$ 从 $\mathrm{SU}(2)\times\mathrm{SU}(2)$ 某元素的左侧移到右侧时，它交换这些元素的次序：
\[
Q\cdot(g,h)=(h,g)\cdot Q .
\]
此恒等式必须被表示保持：
\[
\rho_{j,k}(Q)\,\rho_{j,k}(g,h)=\rho_{j,k}(h,g)\,\rho_{j,k}(Q).
\]
把它作用到两个向量的张量积上：
\begin{align*}
\rho_{j,k}(Q)\,\rho_{j,k}(g,h)\,v\otimes w
&=\rho_{j,k}(h,g)\,\rho_{j,k}(Q)\,v\otimes w\\
\rho_{j,k}(Q)\,(\rho_j(g)v)\otimes(\rho_k(h)w)
&=\rho_{j,k}(h,g)\,\rho_{j,k}(Q)\,v\otimes w .
\end{align*}
若把自身限制在 $j=k$ 的特殊情形，并定义
\[
\rho_{j,j,p}(Q)\,v\otimes w=p\,w\otimes v
\]
（$p$ 是某个待定常数，暂时只把它写进表示的标签里），则我们恒等式的左右两边都变为
\[
p\,(\rho_j(h)w)\otimes(\rho_j(g)v).
\]
但当 $j\neq k$ 时该怎么办？我们需要把表示扩展为我们称之为 $(j,k)\oplus(k,j)$ 的东西，即原表示与另一个自旋互换的版本之直和。我们在 $(\mathbb{V}_j\otimes\mathbb{V}_k)\times(\mathbb{V}_k\otimes\mathbb{V}_j)$（其元素为分别取自两个张量积的有序对）上定义 $\rho_{(j,k)\oplus(k,j),p}$：
\begin{align}
\rho_{(j,k)\oplus(k,j),p}(g,h)\,(a\otimes b,\;c\otimes d)
&=\big((\rho_j(g)a)\otimes(\rho_k(h)b),\;\;(\rho_k(g)c)\otimes(\rho_j(h)d)\big), \tag{4.2.4}\\
\rho_{(j,k)\oplus(k,j),p}(Q)\,(a\otimes b,\;c\otimes d)
&=p\,(d\otimes c,\;b\otimes a). \notag
\end{align}
这同样能让我们满足那个交换 $\mathrm{SU}(2)\times\mathrm{SU}(2)$ 中两元素次序的恒等式。若不是因为有对称操作 $Q$，此表示显然可约，直和的每一部分都是不变子空间；但既然 $Q$ 交换两个子空间中的元素，整体表示就是不可约的。

那常数 $p$ 呢？$\rho_{j,j,p}(Q)^2$ 与 $\rho_{(j,k)\oplus(k,j),p}(Q)^2$ 的作用都只是乘以 $p^2$。若 $Q^2=(I_2,I_2)$，则对我们的任意一个表示有
\[
\rho(Q)^2=\rho(Q^2)=\rho(I_2,I_2)=I,
\]
故必有 $p^2=1$，即 $p=\pm1$。我们称 $p$ 为该表示的\textbf{内禀宇称}（intrinsic parity）。

另一方面，若 $Q^2=(-I_2,-I_2)$，则
\[
\rho_{(j,k)\oplus(k,j),p}(Q)^2=\rho_{(j,k)\oplus(k,j),p}(Q^2)
=\rho_{(j,k)\oplus(k,j),p}(-I_2,-I_2)=(-1)^{2(j+k)}I .
\]
在玻色情形（$j+k$ 为整数），内禀宇称的选择与之前相同。在费米情形（$j+k$ 等于半奇数），内禀宇称是虚数 $\pm i$。

把 $\mathrm{SU}(2)$ 的表示扩展到含 $Q$ 的情形要简单得多。既然 $Q$ 与每个元素对易，它必被表示为恒等算符的倍数，故有
\[
\rho_{j,p}(Q)=p\,I,
\]
其中 $p$ 的选择取决于相同的判据。

最后，我们可以像构造 $\mathrm{SO}(3)$、$\mathrm{SO}(4)$ 的表示那样，用扩展后的二对一同态 $f_3$ 与 $f_4$ 来构造 $\mathrm{O}(3)$ 与 $\mathrm{O}(4)$ 的表示。在玻色情形这些是真正的表示，在费米情形则是射影表示（projective representation）。

\subsection{欧氏群的表示（Representations of the Euclidean Group）}

到目前为止，我们讨论的所有表示，都是欧氏群中使坐标原点固定的子群的表示。现在我们想前进到完整的群，它包含\emph{平移}（translation）——把四维空间中每一点都位移某个常矢量的对称操作。

\subsubsection{函数空间上的表示（Representations on Function Spaces）}

为找出整个欧氏群的表示，我们首先要引入一种新的向量空间：$\mathbb{R}^4$ 上的复值函数空间。取从 $\mathbb{R}^4$ 到 $\mathbb{C}$ 的函数之集，以显然的方式定义加法与标量乘法，便可把它当作向量空间。对任意这样的函数 $A,B$ 及任意复数 $c$：
\[
(A+B)(x)=A(x)+B(x),\qquad (cA)(x)=c\,A(x).
\]
通过定义函数间的内积
\[
\langle A,B\rangle=\int_{\mathbb{R}^4}A(x)^* B(x)\,d^4x,
\]
我们可以把 $\mathbb{R}^4$ 上复值函数的一个子集做成希尔伯特空间，记作 $L^2(\mathbb{R}^4)$。$L^2(\mathbb{R}^4)$ 由那些其模方
\[
|A|^2=\langle A,A\rangle=\int_{\mathbb{R}^4}A(x)^* A(x)\,d^4x=\int_{\mathbb{R}^4}|A(x)|^2 d^4x
\]
有限的函数组成。这样的函数称为\textbf{平方可积}（square-integrable）函数。

我们可以在 $L^2(\mathbb{R}^4)$ 上定义 $\mathrm{E}(4)$ 的一个表示 $\rho_F$：
\[
\big(\rho_F(g)A\big)(x)=A(g^{-1}x).
\]
可以把函数 $A$ 想象为 $\mathbb{R}^4$ 中一个“涂着”复数的刚性物体。当我们用对称操作 $g$ 滑动这个物体时，$\mathbb{R}^4$ 中每一点现在都关联到最终落在该点的那个新数。这里需要用 $g^{-1}$ 而非 $g$，才能得到一个表示：
\begin{align*}
\big(\rho_F(gh)A\big)(x)
&=A\big((gh)^{-1}x\big)\\
&=A\big(h^{-1}g^{-1}x\big)\\
&=\big(\rho_F(h)A\big)(g^{-1}x)\\
&=\big(\rho_F(g)\rho_F(h)A\big)(x).
\end{align*}
由于欧氏群在 $\mathbb{R}^4$ 上的作用保持四维体积，给出内积的积分不变。所以这是一个酉表示。它在整个 $L^2(\mathbb{R}^4)$ 上不是不可约的，但我们可以找到一个不可约子空间：对任意固定值 $\omega_m$ 的黎曼标量波方程（Riemannian Scalar Wave equation, RSW）之解：
\begin{equation}
\partial_x^2 A+\partial_y^2 A+\partial_z^2 A+\partial_t^2 A+\omega_m^2 A=0. \tag{RSW}
\end{equation}
要看出满足 RSW 方程的函数 $A$ 意味着 $\rho_F(g)A$ 也满足它，设 $g^{-1}$ 作用在 $\mathbf{x}=x^i\mathbf{e}_i$ 上为
\[
g^{-1}x=(R^i_{\ j}x^j+a^i)\,\mathbf{e}_i,
\]
其中 $R^i_{\ j}$ 是某正交矩阵的分量，$a^i$ 是 $\mathbb{R}^4$ 中某向量的分量，并对重复指标采用爱因斯坦求和约定。则
\begin{align*}
&\partial_x^2(\rho_F(g)A)(x)+\partial_y^2(\rho_F(g)A)(x)
+\partial_z^2(\rho_F(g)A)(x)+\partial_t^2(\rho_F(g)A)(x)\\
&=\partial_x^2\big(A(g^{-1}x)\big)+\partial_y^2\big(A(g^{-1}x)\big)
+\partial_z^2\big(A(g^{-1}x)\big)+\partial_t^2\big(A(g^{-1}x)\big)\\
&=\sum_{p}\partial_{x^p}^2\big(A(g^{-1}x)\big)\\
&=\sum_{p,q,r}R^q_{\ p}R^r_{\ p}\,\partial_{y^q}\partial_{y^r}A(y)\Big|_{y=g^{-1}x}\\
&=\sum_{q,r}(RR^{T})^q_{\ r}\,\partial_{y^q}\partial_{y^r}A(y)\Big|_{y=g^{-1}x}\\
&=\sum_{r}\partial_{y^r}^2 A(y)\Big|_{y=g^{-1}x}\\
&=-\omega_m^2 A(y)\Big|_{y=g^{-1}x}
=-\omega_m^2\,(\rho_F(g)A)(x).
\end{align*}
这里我们做的是：从 $A$ 对坐标 $x^p$ 的导数（$A$ 只通过 $g^{-1}\mathbf{x}$ 见到这些坐标）转换到对直接传给 $A$ 的坐标 $y^r$ 的导数；这让我们能够引用原函数所满足的方程。实际上，我们是在说明：在一个直角坐标系中满足 RSW 方程的函数，在任何其他直角坐标系中进行计算时也同样满足它。

那向量值函数呢？我们可以在 $\mathbb{C}^n$ 值函数空间上定义内积
\[
\langle A,B\rangle=\int_{\mathbb{R}^4}\langle A(x),B(x)\rangle\,d^4x,
\]
其中积分号内的内积就是 $\mathbb{C}^n$ 本身的标准内积。于是 $\mathbb{C}^n$ 值函数的模方为
\[
|A|^2=\langle A,A\rangle=\int_{\mathbb{R}^4}\langle A(x),A(x)\rangle\,d^4x=\int_{\mathbb{R}^4}|A(x)|^2 d^4x .
\]
把 $g$ 的“非平移部分”作用到函数在 $\mathbb{R}^4$ 中每一点所赋予的四元复数组上，我们可以在平方可积 $\mathbb{C}^4$ 值函数空间 $\mathbb{C}^4\otimes L^2(\mathbb{R}^4)$ 上得到一个酉表示。若 $\mathbf{A}\in\mathbb{C}^4\otimes L^2(\mathbb{R}^4)$ 且 $g^{-1}\mathbf{x}=(R^i_{\ j}x^j+a^i)\mathbf{e}_i$，我们定义
\[
\big(\rho_V(g)A\big)(x)
=(R^{-1})^i_{\ j}A^j(g^{-1}x)\,\mathbf{e}_i
=R^j_{\ i}A^j(g^{-1}x)\,\mathbf{e}_i .
\]
若要求 $\mathbf{A}$ 的每个分量都满足 RSW 方程，则 $\rho_V(g)\mathbf{A}$ 的分量也满足。但正如在向量波注记中所提示的，这\emph{并不}足以给出不可约表示，因为我们还可以通过施加\textbf{横波条件}（transverse condition）
\begin{equation}
\partial_x A^x+\partial_y A^y+\partial_z A^z+\partial_t A^t=0 \tag{Transverse}
\end{equation}
把函数限制到一个更小的不变子空间上。设 $\mathbf{A}$ 满足横波条件，则对 $\rho_V(g)\mathbf{A}$ 有
\begin{align*}
&\partial_x(\rho_V(g)A)^x(x)+\partial_y(\rho_V(g)A)^y(x)
+\partial_z(\rho_V(g)A)^z(x)+\partial_t(\rho_V(g)A)^t(x)\\
&=\sum_{p,j}\partial_{x^p}R^j_{\ p}A^j(g^{-1}x)\\
&=\sum_{p,j,q}R^q_{\ p}\partial_{y^q}R^j_{\ p}A^j(y)\Big|_{y=g^{-1}x}\\
&=\sum_{j,q}(RR^T)^q_{\ j}\partial_{y^q}A^j(y)\Big|_{y=g^{-1}x}\\
&=\sum_{j}\partial_{y^j}A^j(y)\Big|_{y=g^{-1}x}=0 .
\end{align*}
所以 $\rho_V(g)\mathbf{A}$ 也满足横波条件。

在这个情形中我们设法在大表示里找到了一个不可约表示，但过程仍有些神秘。要想系统地确定欧氏群\emph{所有}的不可约酉表示，需要援引一种更精细的方法。

\textbf{一个警告}：本节余下部分包含大量抽象的群论构造。若你享受这种细致程度，那很好——但若你觉得它难啃，跳过即可！它并非后续材料的必要前提。

\subsubsection{群作用、轨道与迷向子群（Group Actions, Orbits and Isotropy Subgroups）}

首先给出几个我们需要的简单定义\textsuperscript{[2]}。

群 $G$ 在集合 $M$ 上的一个\textbf{作用}（action）$\alpha$，为群的每个元素 $g$ 指派一个可逆函数 $\alpha(g):M\to M$，并要求 $\alpha(e)x=x$，其中 $e$ 是群的单位元。若 $\alpha_L$ 是\textbf{左作用}，则对所有 $x\in M$ 及 $g,h\in G$，它必须满足
\[
\alpha_L(g)\alpha_L(h)x=\alpha_L(gh)x .
\]
若 $\alpha_R$ 是\textbf{右作用}，则条件为
\[
\alpha_R(h)\alpha_R(g)x=\alpha_R(gh)x .
\]
最显然的例子是群作用在\emph{自身}上，通过左乘或右乘：
\[
\alpha_L(g)a=ga,\qquad \alpha_R(g)a=ag .
\]
群也可以通过\textbf{共轭}（conjugation）作用在自身上，既可作为左作用也可作为右作用：
\[
\alpha_L(g)a=gag^{-1},\qquad \alpha_R(g)a=g^{-1}ag .
\]
群在向量空间上的表示是左作用，且附加要求：每个群元素所对应的函数必须是线性算符。

若有一个 $G$ 在 $M$ 上的作用 $\alpha$，则 $M$ 中任意元素 $x$ 的\textbf{轨道}（orbit）是集合 $\{\alpha(g)x\mid g\in G\}$。它由 $M$ 中所有能被某个群元素作用把 $x$ 移到的点组成。例如，若取 $\mathrm{SO}(3)$ 在 $\mathbb{R}^3$ 上以矩阵乘法作用（即 $\mathrm{SO}(3)$ 元素只是旋转 $\mathbb{R}^3$ 中的向量），则除原点外 $\mathbb{R}^3$ 中任一点的轨道是一个球面，而原点的轨道只是原点本身。

$G$ 在 $M$ 上的作用 $\alpha$ 称为\textbf{可迁}（transitive）的，如果对任意两个元素 $m_1,m_2$，都存在某个 $g\in G$ 使
\[
\alpha(g)m_1=m_2 .
\]
换句话说，若作用在 $M$ 上可迁，则任意点的轨道就是整个集合 $M$。

我们把 $G$ 在 $M$ 上的\emph{所有轨道之集}记作 $M/G$。在 $\mathrm{SO}(3)$ 与 $\mathbb{R}^3$ 的例子中，轨道集 $\mathbb{R}^3/\mathrm{SO}(3)$ 同构于非负实数集，等于各球面轨道的半径。

\textbf{迷向子群}（isotropy subgroup）$G_m$ 是 $G$ 中使 $M$ 中某点 $m$ 在此作用下固定的元素之集：
\[
G_m=\{h\in G\mid \alpha(h)m=m\}.
\]
例如，若 $G=\mathrm{SO}(3)$、$M=\mathbb{R}^3$，则使 $z$ 轴固定的迷向子群 $G_{(0,0,1)}$ 是 $\mathrm{SO}(3)$ 的一个同构于二维转动群 $\mathrm{SO}(2)$ 的子群。

若 $M$ 中两点 $x,y$ 位于某作用 $\alpha$ 的同一轨道上，则它们在该作用下的迷向子群彼此同构；即存在可逆群同态 $\phi:G_x\to G_y$。对左作用，若 $y=\alpha_L(g)x$，则同构为 $\phi(h)=ghg^{-1}$。对右作用，若 $y=\alpha_R(g)x$，则同构为 $\phi(h)=g^{-1}hg$。

例如，$\mathrm{SO}(3)$ 中分别使 $x$ 轴与 $y$ 轴固定的两个子群彼此同构，取 $g$ 为把 $x$ 轴映到 $y$ 轴的任意旋转并用 $\phi(h)=ghg^{-1}$ 即可证明。

\subsubsection{半直积（Semidirect Products）}

设有一个群 $G$ 及 $G$ 在向量空间 $\mathbb{R}^4$ 上的一个表示 $\sigma$。则通过定义群运算
\[
(g,x)(h,y)=\big(gh,\;\sigma(g)y+x\big),
\]
我们可以把 $G\times\mathbb{R}^4$（有序对 $(g,\mathbf{x})$ 之集）变成一个群 $S$。群单位元为 $(I,\mathbf{0})$，逆元为
\[
(g,x)^{-1}=\big(g^{-1},\;-\sigma(g)^{-1}x\big).
\]
当 $G=\mathrm{O}(4)$ 且 $\sigma$ 为矩阵乘法时，$S$ 正是欧氏四维空间一切对称操作构成的欧氏群 $\mathrm{E}(4)$。$S$ 的群乘法规则，就是我们在用扩展矩阵表示平移矢量时见到的欧氏对称操作的复合规则：
\[
\begin{pmatrix}R_2 & \mathbf{s}_2\\ \mathbf{0} & 1\end{pmatrix}
\begin{pmatrix}R_1 & \mathbf{s}_1\\ \mathbf{0} & 1\end{pmatrix}
=
\begin{pmatrix}R_2 R_1 & R_2\mathbf{s}_1+\mathbf{s}_2\\ \mathbf{0} & 1\end{pmatrix}.
\]
形如 $(I,\mathbf{x})$ 的元素构成平移子群 $\mathrm{T}(4)$，而形如 $(g,\mathbf{0})$ 的元素构成作为 $\mathrm{E}(4)$ 子群的 $\mathrm{O}(4)$。

我们说 $\mathrm{E}(4)$ 是 $\mathrm{O}(4)$ 与 $\mathrm{T}(4)$ 的\textbf{半直积}。这一术语在群论中有更一般的定义，但对我们而言最关心的事实是：$\mathrm{E}(4)$ 的每个元素都能唯一地写成 $\mathrm{O}(4)$ 与 $\mathrm{T}(4)$ 各一个元素的有序对。

我们还关心把 $\mathrm{O}(4)$ 换成另一个群时会发生什么——例如 $\mathrm{SO}(4)$、或其双覆盖 $\mathrm{SU}(2)\times\mathrm{SU}(2)$、或 $\mathrm{SU}(2)\times\mathrm{SU}(2)$ 扩展到 $\mathrm{O}(4)$ 双覆盖后所得的群。若使用这些双覆盖之一，它在 $\mathbb{R}^4$ 上的表示 $\sigma$ 将是二对一同态 $f_4$，而所得半直积 $S$ 本身将是一个双覆盖，或是 $\mathrm{E}(4)$ 的双覆盖，或是其排除反射的子群 $\mathrm{SE}(4)$ 的双覆盖。

\subsubsection{向量丛（Vector Bundles）}

我们现在定义一个称为\textbf{向量丛}（vector bundle）的数学结构\textsuperscript{[3]}。它由一个集合 $E$（丛）、一个称为丛的\textbf{底空间}（base space）的集合 $M$、以及一个\textbf{投影}（projection）$\pi:E\to M$ 组成，后者把 $E$ 中每个元素映射到 $M$ 中一个元素。

\orthofig{images/03__011.png}{圆上的向量丛}
\orthofig{images/03__012.png}{圆上的莫比乌斯丛}

对 $M$ 中每一点 $m$，$E$ 中投影到 $m$ 的子集是一个向量空间；我们把 $\pi^{-1}(m)$ 作为一个向量空间记作 $E_m$。我们称 $E$ 是\textbf{在}集合 $M$ \textbf{之上}的向量丛，并可把它想象为一族向量空间，每一个“坐落在”$M$ 中不同点之上。我们把每个向量空间 $E_m$ 称为 $m$ 上的\textbf{纤维}（fibre）。

一个简单的向量丛例子是 $\mathbb{R}^3$ 中中心在 $z$ 轴上的无限长圆柱面，连同到该圆柱面与 $xy$ 平面相交所得圆的投影。“在”圆上每点“之上”都躺着一个一维实向量空间，由投影到该点的那条直线组成。

\textbf{平凡}丛（trivial bundle）是指仅由底空间 $M$ 与某向量空间 $V$ 的点对组成、即丛为 $E=M\times V$ 的丛。我们的圆柱面就是一例。但丛不必是平凡的。例如我们可以构造一个具有莫比乌斯带拓扑的丛。它与圆柱面有相同的底空间，且每个 $E_m$ 也是一维实向量空间，但它显然\emph{不}是同一个集合。（我们无法把无限宽的莫比乌斯带嵌入三维空间，因此该插图只是想给拓扑一个直观印象。）

若存在 $G$ 的分别作用在 $E$ 与 $M$ 上的左作用 $\alpha_E$ 与 $\alpha_M$，且它们与投影 $\pi:E\to M$ 相容——即先作用 $E$ 再投影到 $M$，与先投影到 $M$ 再作用 $M$ 效果相同——则向量丛 $E$ 称为\textbf{齐性 $G$-丛}（homogeneous $G$-bundle）。换言之，对所有 $g\in G$ 与 $\mathbf{w}\in E$：
\[
\alpha_M(g)\pi(w)=\pi\big(\alpha_E(g)w\big).
\]
这意味着 $\alpha_E(g)$ 把纤维 $E_m$ 映入纤维 $E_{\alpha_M(g)m}$。我们要求这个映射是\emph{线性}的。

在任意向量丛 $E$ 上，我们可以定义丛的\textbf{截面}（section）向量空间 $\Gamma(E)$。$\Gamma(E)$ 中每个截面 $f$ 是一个从底空间 $M$ 到 $E$ 的函数，使得 $f$ 把每一点 $m$ 指派给“坐落在”该点之上的向量空间 $E_m$ 的一个元素。另一种说法是：
\[
\pi(f(m))=m
\]
对 $M$ 中所有 $m$ 成立。

通过逐点把两个截面相加、并把截面乘以标量（在每一点乘以其值），$\Gamma(E)$ 成为向量空间。截面有点像 $M$ 上的向量值函数——只不过 $M$ 中每点的值都落在不同的向量空间 $E_m$ 中。若我们在圆柱面上画一条曲线，使其与圆柱面上每条竖直线恰好相交一次，那就给出一个截面。

若 $E$ 是齐性 $G$-丛，我们可以用作用 $\alpha_E$ 与 $\alpha_M$ 得到 $G$ 在截面向量空间 $\Gamma(E)$ 上的一个表示。若把这个表示记为 $r$，则对每个 $g\in G$、每个 $f\in\Gamma(E)$ 及每个 $m\in M$：
\begin{equation}
\big(r(g)f\big)(m)=\alpha_E(g)\,f\big(\alpha_M(g^{-1})m\big). \tag{4.2.5}
\end{equation}
我们可以对所得的新函数施加投影 $\pi$，来检验它确实是\emph{截面}：
\begin{align*}
\pi\big((r(g)f)(m)\big)
&=\pi\big(\alpha_E(g)f(\alpha_M(g^{-1})m)\big)\\
&=\alpha_M(g)\pi\big(f(\alpha_M(g^{-1})m)\big)\\
&=\alpha_M(g)\alpha_M(g^{-1})m\\
&=m .
\end{align*}

\subsubsection{向量丛上的表示（Representations on Vector Bundles）}

以下内容主要基于 Sternberg\textsuperscript{[4]}，并按需要从洛伦兹物理改造为黎曼物理。

我们先把向量空间 $\mathbb{R}^4$ 看作一个\textbf{阿贝尔}（Abelian）群：即运算（此处为向量加法）可交换的群。我们可以把这个群叫作 $\mathrm{T}(4)$，但改记为 $(\mathbb{R}^4,+)$，以提醒我们可以把 $\mathrm{T}(4)$ 中每个平移简单地看成 $\mathbb{R}^4$ 中的一个向量。

虽然我们不予证明，但阿贝尔群在\emph{复}向量空间上的每个不可约表示都是一维的。一维复向量空间上的线性算符就是乘以一个复数，而若该表示是酉的，这个数的逆必等于其复共轭，故其绝对值必须为 $1$。$(\mathbb{R}^4,+)$ 的一个不可约酉表示 $\chi$ 本质上是一个从 $(\mathbb{R}^4,+)$ 到绝对值为 $1$ 的复数集的函数。这些复数也称为\emph{相位}（phase），每个相位都可写成 $\exp(i\theta)$ 的形式，其中 $\theta$ 为某实数。

$\mathbb{R}^4$ 中任意向量 $\mathbf{k}$ 给出群 $(\mathbb{R}^4,+)$ 的一个一维酉表示 $\chi_{\mathbf{k}}$：
\[
\chi_{\mathbf{k}}(x)=\exp(i\,\mathbf{k}\cdot x).
\]
我们可以通过验证
\[
\chi_{\mathbf{k}}(x)\chi_{\mathbf{k}}(y)=\exp(i\,\mathbf{k}\cdot x)\exp(i\,\mathbf{k}\cdot y)=\exp(i\,\mathbf{k}\cdot(x+y))=\chi_{\mathbf{k}}(x+y)
\]
看出这是一个表示。事实上，$(\mathbb{R}^4,+)$ 的\emph{全部}不可约酉表示都是这种形式。于是我们可以用第二种方式看待 $\mathbb{R}^4$ 本身：把它看作加法群 $(\mathbb{R}^4,+)$ 的全部不可约酉表示之集。

设 $S$ 是某群 $G$ 与 $(\mathbb{R}^4,+)$ 经由 $G$ 在 $\mathbb{R}^4$ 上一个保持点积的表示 $\sigma$ 所得之半直积。若 $G$ 是 $\mathrm{O}(4)$ 或 $\mathrm{SO}(4)$，则 $\sigma$ 就是基本矩阵表示；若 $G$ 是双覆盖，则 $\sigma$ 就是二对一同态 $f_4$。存在一个 $S$ 在其子群 $(\mathbb{R}^4,+)$ 上的右作用 $\alpha$，定义如下：其中 $s=(g,\mathbf{y})$ 是 $S$ 的任意元素，$\mathbf{x}$ 是 $(\mathbb{R}^4,+)$ 的任意元素，当要把 $\mathbf{x}$ 显式作为半直积 $S$ 的元素时写成 $(I,\mathbf{x})$。
\begin{align*}
\alpha(s)x
&=s^{-1}xs\\
&=(g,y)^{-1}(I,x)(g,y)\\
&=\big(g^{-1},-\sigma(g)^{-1}y\big)(I,x)(g,y)\\
&=\big(g^{-1},-\sigma(g)^{-1}y\big)(g,y+x)\\
&=\big(I,\sigma(g)^{-1}x\big)\\
&=\sigma(g)^{-1}x .
\end{align*}
由这个右作用 $\alpha$，可得到 $S$ 在 $\mathbb{R}^4$（作为加法群 $(\mathbb{R}^4,+)$ 的不可约酉表示之集）上的一个左作用 $\beta$：
\begin{align*}
(\beta(s)\chi_{\mathbf{k}})(x)
&=\chi_{\mathbf{k}}(\alpha(s)x)\\
&=\chi_{\mathbf{k}}\big(\sigma(g)^{-1}x\big)\\
&=\exp\big(i\,\mathbf{k}\cdot(\sigma(g)^{-1}x)\big)\\
&=\exp\big(i\,(\sigma(g)\mathbf{k})\cdot x\big)\\
&\qquad\text{[$\sigma(g)$ 保持点积，且把它同时作用于 $\mathbf{k}$ 与 $\sigma(g)^{-1}x$]}\\
&=\chi_{\sigma(g)\mathbf{k}}(x).
\end{align*}
我们通过 $\beta$ 在表示集 $\chi_{\mathbf{k}}$ 上的作用来定义它，但我们将模糊 $\chi_{\mathbf{k}}$ 与 $\mathbf{k}$ 的区别，并把 $\beta$ 也用作 $\mathbb{R}^4$（仅仅看作向量空间）上的作用：
\[
\beta(s)\mathbf{k}=\beta(g,y)\mathbf{k}=\sigma(g)\mathbf{k}.
\]
任意非零向量 $\mathbf{k}$ 在此作用 $\beta$ 下的轨道都是 $\mathbb{R}^4$ 中的一个三维球面，因为 $\sigma(g)$ 保持 $\mathbf{k}$ 的长度——而且（就我们的目的而言）群 $G$ 总是至少为 $\mathrm{SO}(4)$。零向量的轨道只是它自己，一个单点。于是我们可以用半径来对所有轨道分类。[在洛伦兹物理中情况更复杂：那时 $\mathbf{k}$ 依其类时、类空或类光而具有不同类型的轨道。详情见 Sternberg 中所引章节。]

现在，设 $\rho$ 是半直积群 $S$ 在某向量空间 $V$ 上的一个不可约酉表示。$S$ 或是欧氏群 $\mathrm{E}(4)$、或是特殊欧氏群 $\mathrm{SE}(4)$、或是其中之一的双覆盖。

若把 $\rho$ 限制到 $S$ 的平移子群 $(\mathbb{R}^4,+)$，则 $\rho$ 给出 $(\mathbb{R}^4,+)$ 的一个酉表示，但除非 $V$ 是一维的，这个被限制的表示\emph{不会}不可约——因为 $(\mathbb{R}^4,+)$ 是阿贝尔群，其不可约表示都是一维的。于是一般而言，$V$ 将由\emph{许多}在限制后的 $\rho$ 作用下不变的一维子空间张成。

在每个这样的一维子空间中，对所有 $\mathbf{x}\in(\mathbb{R}^4,+)$ 及子空间中每个 $\mathbf{w}$，将有
\[
\rho(x)\mathbf{w}=\chi_{\mathbf{k}}(x)\mathbf{w},
\]
其中向量 $\mathbf{k}$ 标识出 $\rho$ 的限制等价于 $(\mathbb{R}^4,+)$ 的哪一个不可约酉表示。由于无法保证给定的向量 $\mathbf{k}$ 只出现在单个一维子空间里，我们把 $V$ 中使所有 $\mathbf{w}$ 都满足
\[
\rho(x)\mathbf{w}=\chi_{\mathbf{k}}(x)\mathbf{w}
\]
的整个子空间写作 $E(\mathbf{k})$。这样我们就分解了 $V$，得到子空间 $E(\mathbf{k})$，其中表示 $\rho$ 的\emph{平移}部分与某个特定向量 $\mathbf{k}$ 的 $\chi_{\mathbf{k}}$ 相同。

\orthofig{images/03__013.png}{由表示构造的向量丛}

若用 $\rho(s)$（其中 $s=(g,\mathbf{y})$ 是 $S$ 的任意元素）作用在 $E(\mathbf{k})$ 中某个 $\mathbf{w}$ 上会怎样？对任意 $\mathbf{x}\in(\mathbb{R}^4,+)$ 有
\begin{align*}
\rho(x)\rho(s)\mathbf{w}
&=\rho(s)\rho(s^{-1})\rho(x)\rho(s)\mathbf{w}\\
&=\rho(s)\rho(s^{-1}xs)\mathbf{w}\\
&=\rho(s)\chi_{\mathbf{k}}(\alpha(s)x)\mathbf{w}\\
&=\rho(s)(\beta(s)\chi_{\mathbf{k}})(x)\mathbf{w}\\
&=\rho(s)\chi_{\sigma(g)\mathbf{k}}(x)\mathbf{w}\\
&=\chi_{\sigma(g)\mathbf{k}}(x)\,\rho(s)\mathbf{w}.
\end{align*}
这告诉我们 $\rho(s)\mathbf{w}$ 落在 $E(\sigma(g)\mathbf{k})$ 中，其中 $g$ 来自 $s=(g,\mathbf{y})$。换句话说，$\rho(g,\mathbf{y})$ 把 $E(\mathbf{k})$ 映入 $E(\sigma(g)\mathbf{k})$。

于是对应于各个 $E(\mathbf{k})$ 的所有向量 $\mathbf{k}$ 必定来自 $S$ 在 $\mathbb{R}^4$ 上（等价地，$G$ 的作用 $\sigma$）作用 $\beta$ 下的同一条轨道，因而它们都有相同的长度。若该长度为 $k=|\mathbf{k}|$，我们定义 $\mathbf{k}_0=(0,0,0,k)\in\mathbb{R}^4$。则每个 $\mathbf{k}$ 都来自 $\mathbf{k}_0$ 的轨道，而该轨道是 $\mathbb{R}^4$ 中半径为 $k$ 的一个三维球面。我们把这个三维球面称为 $M_k$。

把在 $V$ 中找到的所有 $E(\mathbf{k})$ 合起来，我们写
\[
E=\bigcup_{k}E(k),
\]
其中 “$\bigcup$” 指不交并（disjoint union），即把所有 $E(\mathbf{k})$ 看作完全不同的集合、把它们的所有元素收集到一个大集合 $E$ 中。（具体地说，这意味着尽管所有子空间 $E(\mathbf{k})$ 共享原点——即 $V$ 的原点——但在 $E$ 中我们把每个 $E(\mathbf{k})$ 都看作有各自不同的原点。）

$V$ 中一般向量不会落在这些子空间的\emph{任何}一个里，但 $\rho(s)$ 总把每个 $E(\mathbf{k})$ 映到另一个这样的子空间 $E(\sigma(g)\mathbf{k})$——所以 $\rho$ 给出 $S$ 在 $E$ 上的一个作用，恰如向量丛上的作用，纤维被映到纤维！我们通过定义投影 $\pi$ 把每个 $E(\mathbf{k})$ 的元素都映到 $\mathbf{k}$ 本身，来把这一族向量空间 $E(\mathbf{k})$ 组成的集合 $E$ 做成 $M_k$ 之上的向量丛。

表示 $\rho$ 给出 $S$ 在 $E$ 上的一个左作用，我们还有 $S$ 在 $M_k$ 上的一个左作用，即 $\beta$。这两个作用与我们的投影 $\pi$ 相容；对所有 $s\in S$ 及所有 $\mathbf{w}\in E$：
\[
\beta(s)\pi(w)=\pi\big(\rho(s)w\big).
\]
我们可以用这些作用得到 $S$ 在截面向量空间 $\Gamma(E)$ 上的一个表示。若把这个表示记为 $r$，则对每个 $s\in S$、每个 $f\in\Gamma(E)$ 及每个 $\mathbf{k}\in M_k$：
\[
\big(r(s)f\big)(k)=\rho(s)\,f\big(\beta(s^{-1})k\big).
\]
事实上，$S$ 在 $\Gamma(E)$ 上的表示 $r$ \emph{等价于} $S$ 在 $V$ 上的原表示 $\rho$！给定 $V$ 中任意向量 $\mathbf{v}$，可把它唯一地写成分别落在各 $E(\mathbf{k})$ 中向量之和。这就给每个 $\mathbf{k}$ 指派了 $E(\mathbf{k})$ 中一个向量，于是可把它看作 $\Gamma(E)$ 中一个截面。同样，给定 $\Gamma(E)$ 中任意截面，可以把各 $E(\mathbf{k})$ 中的向量作为 $V$ 中向量加起来。在这个 $V$ 与 $\Gamma(E)$ 的同一化下，两表示等价。

现在，一般而言 $\rho$ 会把向量从一个 $E(\mathbf{k})$ 移到另一个，但若把它限制到某个合适的子群，我们就能考察它仅作用在\emph{一个}这样的空间上的情况。我们定义 $S_{\mathbf{k}_0}$ 为 $S$ 中所有满足
\[
\beta(h,x)\mathbf{k}_0=\sigma(h)\mathbf{k}_0=\mathbf{k}_0
\]
的 $(h,\mathbf{x})$ 组成的迷向子群。我们还定义 $G_{\mathbf{k}_0}$ 为 $G$ 中所有满足 $\sigma(h)\mathbf{k}_0=\mathbf{k}_0$ 的 $h$ 组成的迷向子群。显然 $G_{\mathbf{k}_0}$ 恰由 $S_{\mathbf{k}_0}$ 各元素 $(h,\mathbf{x})$ 的非平移部分 $h$ 组成。事实上，$S_{\mathbf{k}_0}$ 是 $G_{\mathbf{k}_0}$ 与 $(\mathbb{R}^4,+)$ 的半直积。同一条轨道上所有点的迷向子群彼此同构，所以选择 $\mathbf{k}_0$ 而非轨道上其他 $\mathbf{k}$ 对这些子群的结构没有影响。

限制到 $S_{\mathbf{k}_0}$ 时，$\rho$ 把 $E(\mathbf{k}_0)$ 映到自身，给出 $S_{\mathbf{k}_0}$ 的一个表示。但表示的平移部分我们已经了解：$E(\mathbf{k})$ 的定义说明，在 $E(\mathbf{k}_0)$ 上 $\rho$ 限制到平移 $(I,\mathbf{y})$ 由 $\chi_{\mathbf{k}_0}(\mathbf{y})$ 描述，即一致地乘以相位 $\exp(i\,\mathbf{k}_0\cdot\mathbf{y})$。所以让我们集中于非平移部分，抽取出 $G_{\mathbf{k}_0}$ 在 $E(\mathbf{k}_0)$ 上的一个表示 $\eta$。对任意 $h\in G_{\mathbf{k}_0}$ 及 $\mathbf{w}\in E(\mathbf{k}_0)$，定义
\[
\eta(h)\mathbf{w}=\rho(h,0)\mathbf{w}.
\]
$\rho$ 是表示这一事实足以使 $\eta$ 是表示，而 $\rho$ 不可约意味着 $\eta$ 也不可约。后者或许不那么显然：若存在 $E(\mathbf{k}_0)$ 的某个真子空间 $U$ 被 $G_{\mathbf{k}_0}$ 保持，我们可以定义
\[
U^{+}=\text{由}\ \{\rho(s)\mathbf{w}\mid s\in S,\ \mathbf{w}\in U\}\ \text{张成的}\ V\ \text{的子空间}.
\]
$U^{+}$ 不是 $V$ 的全部，因为按假设 $G_{\mathbf{k}_0}$ 的任何元素都不能把 $\mathbf{w}$ 移出 $U$，$S_{\mathbf{k}_0}$ 的任何元素也不行（它至多再乘一个相位），而 $S_{\mathbf{k}_0}$ 之外的任何 $S$ 元素都把每个 $\mathbf{w}$ 映入某个 $\mathbf{k}\neq\mathbf{k}_0$ 的 $E(\mathbf{k})$。但 $U^{+}$ 是不变子空间，因为用 $\rho$ 作用在形如 $\rho(s)\mathbf{w}$ 的向量的任意线性组合上，仍给出同形向量的线性组合。因此我们关于 $U$ 的假设将与 $\rho$ 的不可约性矛盾。

关于使 $\eta$ 成为 $E(\mathbf{k}_0)$ 上不可约表示的迷向子群 $G_{\mathbf{k}_0}$，我们能说些什么？我们将假设 $k\neq0$，于是 $\mathbf{k}_0$ 是欧氏四维空间中的非零向量。保持 $\mathbf{k}_0=(0,0,0,k)$ 固定的 $\mathrm{SO}(4)$ 子群就是 $\mathrm{SO}(3)$。我们可以通过从全群与子群两边都去掉“特殊”（special）的 “S”，和／或把两者都做成双覆盖，来得到各种变体。为了得到真正的表示，最简单的做法是用双覆盖而非 $\mathrm{SO}(4)$ 或 $\mathrm{O}(4)$ 本身，于是两个最有用的情形是：

\begin{itemize}
\item 若 $G=\mathrm{SU}(2)\times\mathrm{SU}(2)$（$\mathrm{SO}(4)$ 的双覆盖），则 $G_{\mathbf{k}_0}$ 是 $\mathrm{SU}(2)$（$\mathrm{SO}(3)$ 的双覆盖）。更确切地说，它是 $\mathrm{SU}(2)\times\mathrm{SU}(2)$ 中由形如 $(h,h)$ 的元素对组成的子群，同构于 $\mathrm{SU}(2)$。
\item 若 $G$ 是 $\mathrm{SU}(2)\times\mathrm{SU}(2)$ 扩展到 $\mathrm{O}(4)$ 双覆盖所得的群，则 $G_{\mathbf{k}_0}$ 是 $\mathrm{SU}(2)$ 扩展到 $\mathrm{O}(3)$ 双覆盖所得的群。
\end{itemize}

这些群的不可约酉表示我们已经考察过了。对 $\mathrm{SU}(2)$，它们是自旋 $j$ 表示。对 $\mathrm{O}(3)$ 的双覆盖，还有一个内禀宇称的选择。

在对 $\rho$（一个假定的 $S$ 的不可约表示）作了“反向工程”之后，我们现在能看出如何自己\emph{构造}同类的表示。步骤如下：

\begin{itemize}
\item 要找到 $\mathrm{SE}(4)$（排除反射）双覆盖的表示，令群 $G$ 等于 $\mathrm{SU}(2)\times\mathrm{SU}(2)$，即 $\mathrm{SO}(4)$ 的双覆盖。
\item 要找到 $\mathrm{E}(4)$（含反射）双覆盖的表示，令群 $G$ 等于 $\mathrm{SU}(2)\times\mathrm{SU}(2)$ 扩展到 $\mathrm{O}(4)$ 双覆盖所得的群。
\item 定义 $S$ 为 $G$ 与 $(\mathbb{R}^4,+)$ 的半直积。为此所需的 $G$ 在 $\mathbb{R}^4$ 上的表示——在一般情形中我们称之为 $\sigma$——将是 $G$ 到 $\mathrm{SO}(4)$ 或 $\mathrm{O}(4)$ 的二对一同态 $f_4$。
\item 选取一个正数 $k$，定义 $\mathbf{k}_0=(0,0,0,k)$，并定义 $\chi_{\mathbf{k}_0}(\mathbf{y})=\exp(i\,\mathbf{k}_0\cdot\mathbf{y})$。
\item 选取迷向子群 $G_{\mathbf{k}_0}$ 的一个不可约酉表示 $\eta$。这将是 $\mathrm{SU}(2)$ 的一个自旋 $j$ 表示，可能还带上一个内禀宇称 $p$。把 $\eta$ 作用其上的向量空间 $\mathbb{C}^{2j+1}$ 记为 $W$。
\item 通过 $\xi(h,\mathbf{y})=\chi_{\mathbf{k}_0}(\mathbf{y})\eta(h)$ 定义 $S_{\mathbf{k}_0}$（$G_{\mathbf{k}_0}$ 与 $(\mathbb{R}^4,+)$ 的半直积）在 $W$ 上的一个表示 $\xi$。
\item 我们要一个齐性 $S$-丛 $E$，其底空间为 $M_k$，每个向量空间 $E(\mathbf{k})$ 同构于 $W$。但我们不能把 $E$ 定义为 $M_k\times W$（平凡丛），因为没有自然的作用让整个群 $S$ 作用其上。不过很容易找到 $S$ 在自身上的作用——而且，存在 $S_{\mathbf{k}_0}$ 在 $S$ 上的一个作用，使轨道 $S/S_{\mathbf{k}_0}$ 能与 $M_k$ 建立一一对应。定义 $S_{\mathbf{k}_0}$ 在 $S$ 上的一个左作用 $\gamma_0$ 为 $\gamma_0(s_0)s=s\,s_0^{-1}$。把 $S_{\mathbf{k}_0}$ 的元素取逆后放到 $S$ 元素的\emph{右}边看似古怪，但这是本质的。然后定义映射 $\pi_0:S\to M_k$：
\[
\pi_0(g,\mathbf{x})=f_4(g)\mathbf{k}_0 .
\]
有了 $\pi_0$，$S$ 中在作用 $\gamma_0$ 下的同一条轨道 $S/S_{\mathbf{k}_0}$ 上的所有元素都被认同为 $M_k$ 的同一元素。既已知道如何把 $S$“缩小”为 $M_k$，我们在集合 $S\times W$ 上照搬这一技巧。定义 $S_{\mathbf{k}_0}$ 在 $S\times W$ 上的一个左作用 $\gamma$：
\[
\gamma(h,\mathbf{y})\big((g,\mathbf{x}),\mathbf{w}\big)=\big((g,\mathbf{x})(h,\mathbf{y})^{-1},\ \xi(h,\mathbf{y})\mathbf{w}\big).
\]
把 $((g,\mathbf{x}),\mathbf{w})$ 在此作用下的轨道写作 $[((g,\mathbf{x}),\mathbf{w})]$。现在我们定义丛 $E$ 为 $S\times W$ 在 $S_{\mathbf{k}_0}$ 的作用 $\gamma$ 下的\emph{轨道之集}，记作 $E=(S\times W)/S_{\mathbf{k}_0}$。再定义投影 $\pi:E\to M_k$ 为
\[
\pi\big([((g,\mathbf{x}),\mathbf{w})]\big)=f_4(g)\mathbf{k}_0,
\]
这使 $E$ 成为 $M_k$ 之上的一个丛。
\item 我们的投影 $\pi$ 定义是否无歧义，即结果与轨道中所取的点无关？当遍历轨道时，$S$ 的元素 $(g,\mathbf{x})$ 变为 $(g,\mathbf{x})(h,\mathbf{y})^{-1}$，其非平移部分（投影所依赖的）为 $gh^{-1}$。当 $h$ 遍历迷向子群 $G_{\mathbf{k}_0}$ 时，$gh^{-1}$ 取哪些值？定义 $G(\mathbf{k}_0,\mathbf{k})$ 为 $G$ 中在作用 $f_4$ 下把 $\mathbf{k}_0$ 映到 $\mathbf{k}$ 的子集：
\[
G(\mathbf{k}_0,\mathbf{k})=\{a\in G\mid f_4(a)\mathbf{k}_0=\mathbf{k}\}.
\]
若 $g\in G(\mathbf{k}_0,\mathbf{k})$，则 $gh^{-1}$ 也在其中，因为 $h^{-1}$ 使 $\mathbf{k}_0$ 固定。设 $a$ 也在 $G(\mathbf{k}_0,\mathbf{k})$ 中，则 $f_4(g)\mathbf{k}_0=\mathbf{k}$、$f_4(a)\mathbf{k}_0=\mathbf{k}$。于是
\[
\mathbf{k}_0=f_4(a^{-1})\mathbf{k}=f_4(a^{-1})f_4(g)\mathbf{k}_0=f_4(a^{-1}g)\mathbf{k}_0,
\]
这告诉我们 $a^{-1}g\in G_{\mathbf{k}_0}$。令 $h=a^{-1}g$，则 $h^{-1}=g^{-1}a$，且 $gh^{-1}=a$。所以在每条轨道中 $gh^{-1}$ 所取的值恰是集合 $G(\mathbf{k}_0,\mathbf{k})$，因此无论选轨道中哪一点，它都投影到同一个 $\mathbf{k}$。
\item 定义 $E(\mathbf{k})$ 为 $\pi^{-1}(\mathbf{k})=\{[((a,\mathbf{x}),\mathbf{w})]\mid a\in G(\mathbf{k}_0,\mathbf{k}),\ \mathbf{x}\in\mathbb{R}^4,\ \mathbf{w}\in W\}$。要把纤维 $E(\mathbf{k})$ 变成向量空间，需要能进行通常的向量相加与标量乘法。为此，选取某函数 $a:M_k\to G$，使 $a(\mathbf{k})$ 总在 $G(\mathbf{k}_0,\mathbf{k})$ 中，然后定义映射 $\phi_{a(\mathbf{k})}:W\to E(\mathbf{k})$ 与 $\psi_{a(\mathbf{k})}:E(\mathbf{k})\to W$ 为
\[
\phi_{a(\mathbf{k})}(\mathbf{w})=[((a(\mathbf{k}),\mathbf{0}),\mathbf{w})],\qquad
\psi_{a(\mathbf{k})}\big([((g,\mathbf{x}),\mathbf{w})]\big)=\xi\big((a(\mathbf{k}),\mathbf{0})^{-1}(g,\mathbf{x})\big)\mathbf{w}.
\]
$\psi_{a(\mathbf{k})}$ 是良定义的，因为若在同一轨道中另取一个元素，我们得到相同结果。这正是把 $(h,\mathbf{y})$ 取逆后放到右侧所带来的好处：
\begin{align*}
&\psi_{a(\mathbf{k})}\big([((g,\mathbf{x})(h,\mathbf{y})^{-1},\ \xi(h,\mathbf{y})\mathbf{w})]\big)\\
&=\xi\big((a(\mathbf{k}),\mathbf{0})^{-1}(g,\mathbf{x})(h,\mathbf{y})^{-1}\big)\,\xi(h,\mathbf{y})\mathbf{w}\\
&=\xi\big((a(\mathbf{k}),\mathbf{0})^{-1}(g,\mathbf{x})\big)\mathbf{w}.
\end{align*}
这两个映射 $\phi_{a(\mathbf{k})}$ 与 $\psi_{a(\mathbf{k})}$ 实际上互为逆映射：
\begin{align*}
&\phi_{a(\mathbf{k})}\Big(\psi_{a(\mathbf{k})}\big([((g,\mathbf{x}),\mathbf{w})]\big)\Big)\\
&=\phi_{a(\mathbf{k})}\Big(\xi\big((a(\mathbf{k}),\mathbf{0})^{-1}(g,\mathbf{x})\big)\mathbf{w}\Big)\\
&=\big[\big((a(\mathbf{k}),\mathbf{0}),\ \xi\big((a(\mathbf{k}),\mathbf{0})^{-1}(g,\mathbf{x})\big)\mathbf{w}\big)\big]\\
&=\big[\gamma\big((g,\mathbf{x})^{-1}(a(\mathbf{k}),\mathbf{0})\big)\big((a(\mathbf{k}),\mathbf{0}),\ \xi((a(\mathbf{k}),\mathbf{0})^{-1}(g,\mathbf{x}))\mathbf{w}\big)\big]\\
&=\big[\big((a(\mathbf{k}),\mathbf{0})(a(\mathbf{k}),\mathbf{0})^{-1}(g,\mathbf{x}),\ \xi((g,\mathbf{x})^{-1}(a(\mathbf{k}),\mathbf{0}))\xi((a(\mathbf{k}),\mathbf{0})^{-1}(g,\mathbf{x}))\mathbf{w}\big)\big]\\
&=[((g,\mathbf{x}),\mathbf{w})],
\end{align*}
以及
\begin{align*}
\psi_{a(\mathbf{k})}\big(\phi_{a(\mathbf{k})}(\mathbf{w})\big)
&=\psi_{a(\mathbf{k})}\big([((a(\mathbf{k}),\mathbf{0}),\mathbf{w})]\big)\\
&=\xi\big((a(\mathbf{k}),\mathbf{0})^{-1}(a(\mathbf{k}),\mathbf{0})\big)\mathbf{w}=\mathbf{w}.
\end{align*}
我们用 $\phi_{a(\mathbf{k})}$ 与 $\psi_{a(\mathbf{k})}$ 把 $E(\mathbf{k})$ 与 $W$ 认同。于是若 $\mathbf{c},\mathbf{d}\in E(\mathbf{k})$、$z$ 是复标量，通过定义
\[
\mathbf{c}+\mathbf{d}=\phi_{a(\mathbf{k})}\big(\psi_{a(\mathbf{k})}(\mathbf{c})+\psi_{a(\mathbf{k})}(\mathbf{d})\big),\qquad
z\,\mathbf{c}=\phi_{a(\mathbf{k})}\big(z\,\psi_{a(\mathbf{k})}(\mathbf{c})\big)
\]
把 $E(\mathbf{k})$ 变成向量空间。
\item $\phi_{a(\mathbf{k})}$、$\psi_{a(\mathbf{k})}$ 与 $\xi(h,\mathbf{0})=\eta(h)$（$G_{\mathbf{k}_0}$ 在 $W$ 上的表示）之间有一些简单关系。若 $h\in G_{\mathbf{k}_0}$、$a(\mathbf{k})\in G(\mathbf{k}_0,\mathbf{k})$，则
\[
\phi_{a(\mathbf{k})h}=\phi_{a(\mathbf{k})}\circ\xi(h,\mathbf{0}),\qquad
\psi_{a(\mathbf{k})h}=\xi(h,\mathbf{0})^{-1}\circ\psi_{a(\mathbf{k})}.
\]
由此可知，我们在 $G(\mathbf{k}_0,\mathbf{k})$ 中选取哪个 $a(\mathbf{k})$ 对 $E(\mathbf{k})$ 上向量空间运算的定义没有影响。若把 $E(\mathbf{k})$ 中两元素 $\mathbf{c}$ 与 $\mathbf{d}$ 之和重新定义为 $\mathbf{c}+\mathbf{d}=\phi_{a(\mathbf{k})h}(\psi_{a(\mathbf{k})h}(\mathbf{c})+\psi_{a(\mathbf{k})h}(\mathbf{d}))$，它就变成
\begin{align*}
&\phi_{a(\mathbf{k})}\Big(\xi(h,\mathbf{0})\big(\psi_{a(\mathbf{k})h}(\mathbf{c})+\psi_{a(\mathbf{k})h}(\mathbf{d})\big)\Big)\\
&\quad=\phi_{a(\mathbf{k})}\Big(\xi(h,\mathbf{0})\big(\xi(h,\mathbf{0})^{-1}(\psi_{a(\mathbf{k})}(\mathbf{c})+\psi_{a(\mathbf{k})}(\mathbf{d}))\big)\Big)\\
&\quad=\phi_{a(\mathbf{k})}\big(\psi_{a(\mathbf{k})}(\mathbf{c})+\psi_{a(\mathbf{k})}(\mathbf{d})\big),
\end{align*}
与原定义一致。
\item 定义 $S$ 在 $E$ 上的作用 $\rho$ 为
\[
\rho(b,\mathbf{z})\,[((g,\mathbf{x}),\mathbf{w})]=[((b,\mathbf{z})(g,\mathbf{x}),\mathbf{w})],
\]
并定义 $S$ 在 $M_k$ 上的作用 $\beta$ 为 $\beta(b,\mathbf{z})\mathbf{k}=f_4(b)\mathbf{k}$。这些作用与我们的投影 $\pi$ 相容：
\begin{align*}
\pi\big(\rho(b,\mathbf{z})[((g,\mathbf{x}),\mathbf{w})]\big)
&=\pi\big([((b,\mathbf{z})(g,\mathbf{x}),\mathbf{w})]\big)\\
&=f_4(bg)\mathbf{k}_0=f_4(b)f_4(g)\mathbf{k}_0\\
&=\beta(b,\mathbf{z})\,\pi\big([((g,\mathbf{x}),\mathbf{w})]\big).
\end{align*}
\item 设我们以 $\rho\big((a(\mathbf{k}),\mathbf{0})(h,\mathbf{y})(a(\mathbf{k}),\mathbf{0})^{-1}\big)$ 作用在纤维 $E(\mathbf{k})$ 上，其中 $(h,\mathbf{y})\in S_{\mathbf{k}_0}$。元素 $(a(\mathbf{k}),\mathbf{0})(h,\mathbf{y})(a(\mathbf{k}),\mathbf{0})^{-1}$ 将在使 $\mathbf{k}$ 固定的迷向子群 $S_{\mathbf{k}}$ 中——该子群同构于 $S_{\mathbf{k}_0}$。不难看出
\[
\rho\big((a(\mathbf{k}),\mathbf{0})(h,\mathbf{y})(a(\mathbf{k}),\mathbf{0})^{-1}\big)\,[((g,\mathbf{x}),\mathbf{w})]
=\phi_{a(\mathbf{k})}\Big(\xi(h,\mathbf{y})\,\psi_{a(\mathbf{k})}\big([((g,\mathbf{x}),\mathbf{w})]\big)\Big).
\]
这表明 $\rho$ 限制到纤维 $E(\mathbf{k})$ 与子群 $S_{\mathbf{k}}$ 上的作用“等价于”我们起初所用的 $S_{\mathbf{k}_0}$ 的不可约表示 $\xi$。严格说来我们只应谈完全同一个群的等价表示，但既然已用显式同构把 $S_{\mathbf{k}}$ 与 $S_{\mathbf{k}_0}$ 认同，我们就是在说明：$S_{\mathbf{k}}$ 在每个纤维上的表示本质上只是 $S_{\mathbf{k}_0}$ 在 $W$ 上表示 $\xi$ 的一个版本。
\item 通过以下公式定义 $S$ 在丛截面向量空间 $\Gamma(E)$ 上的表示 $r$：对每个 $s\in S$、$f\in\Gamma(E)$ 及 $\mathbf{k}\in M_k$，
\[
\big(r(s)f\big)(\mathbf{k})=\rho(s)\,f\big(\beta(s^{-1})\mathbf{k}\big).
\]
这将是我们所寻求的\textbf{$S$ 的不可约酉表示}。呼！$\mathrm{E}(4)$ 或 $\mathrm{SE}(4)$ 双覆盖的这个表示完全由我们对以下量的选择确定：
\begin{itemize}
\item 三维球面的半径 $k$；
\item $\mathrm{SU}(2)$ 的不可约表示 $\eta$ 的自旋 $j$；
\item 以及对 $\mathrm{E}(4)$ 的情形，$\eta$ 的内禀宇称 $p$。
\end{itemize}
\end{itemize}

\subsubsection{自旋 0 的丛表示（The Spin-0 Bundle Representation）}

在 $\mathrm{SU}(2)$ 的自旋 0 表示的情形，向量空间 $W$ 就是复数集，$\eta(h)$ 是恒等算符，且我们还有
\[
\xi(h,y)=\chi_{\mathbf{k}_0}(y)\eta(h)=\exp(i\,\mathbf{k}_0\cdot y).
\]
我们用来把每个纤维与 $W$ 认同的映射 $\psi_{a(\mathbf{k})}:E(\mathbf{k})\to W$ 变为
\begin{align*}
\psi_{a(\mathbf{k})}\big([((g,x),w)]\big)
&=\xi\big((a(\mathbf{k}),0)^{-1}(g,x)\big)w\\
&=\chi_{\mathbf{k}_0}\big(f_4(a(\mathbf{k})^{-1})x\big)w\\
&=\exp\big(i\,\mathbf{k}_0\cdot(f_4(a(\mathbf{k})^{-1})x)\big)w\\
&=\exp\big(i\,(f_4(a(\mathbf{k}))\mathbf{k}_0)\cdot x\big)w\\
&=\exp(i\,\mathbf{k}\cdot x)\,w .
\end{align*}
这完全不依赖于我们对 $a(\mathbf{k})$ 的选择！于是我们可以定义一个从整个丛 $E$ 到平凡丛的映射 $\psi:E\to M_k\times\mathbb{C}$：
\[
\psi\big([((g,x),w)]\big)=\big(k,\ \exp(i\,\mathbf{k}\cdot x)\,w\big),
\]
这里为清晰起见只写了 $\mathbf{k}$，但它可由 $E$ 中的点显式给出：$\mathbf{k}=\pi([((g,\mathbf{x}),\mathbf{w})])=f_4(g)\mathbf{k}_0$。我们可以用 $\psi$ 把丛的一个截面变成从 $M_k$ 到复数的函数。同样，我们也能经由映射 $\phi:M_k\times\mathbb{C}\to E$ 把这样的函数变成丛的截面，其中
\[
\phi(k,w)=[((a(k),0),w)].
\]
看起来这似乎依赖于 $a(\mathbf{k})$，但实际上无论怎样都映到完全相同的轨道。这是因为 $\xi(h,\mathbf{y})$ 现在只依赖于平移部分 $\mathbf{y}$，所以只要 $a_1$ 与 $a_2$ 都在 $G(\mathbf{k}_0,\mathbf{k})$ 中，作用 $\gamma(a_2^{-1}a_1,\mathbf{0})$ 就把 $((a_1,\mathbf{0}),w)$ 变成 $((a_2,\mathbf{0}),w)$，表明它们在同一条轨道上。

若我们用 $\phi$ 把函数 $w:M_k\to\mathbb{C}$ 变成一个截面，用 $S$ 的某元素 $(b,\mathbf{z})$ 经由表示 $r$ 作用到该截面上，再用 $\psi$ 把截面变回 $M_k$ 上的函数，我们就得到 $S$ 在 $M_k$ 上函数空间中的一个表示，我们也把它称为 $r$。对任意这样的函数 $w$：
\begin{equation}
\big(r(b,z)w\big)(k)=\exp(i\,\mathbf{k}\cdot z)\,w\big(f_4(b)^{-1}k\big). \tag{4.2.6}
\end{equation}
我们预期 $M_k$ 上自旋 0 丛的截面表示与 $\omega_m=k$ 的 RSW 方程之解有关。在关于标量波的注记中，我们考察过该方程的实值平面波解，但为当前目的我们需要复值平面波：
\[
A_k(x)=A_0\exp(-i\,\mathbf{k}\cdot x).
\]
我们可以通过构造 $\mathbb{R}^4$ 上的一个函数 $A_w$（它等于对平面波解集、以 $M_k$ 上复值函数 $w(\mathbf{k})$ 为权重的积分）来在两个表示间架起桥梁。换句话说，我们把 $M_k$ 中每个向量 $\mathbf{k}$ 视为\emph{一个平面波的传播矢量}，并取 $w(\mathbf{k})$ 为该平面波对整个 $\mathbb{R}^4$ 上函数之傅里叶合成（Fourier synthesis）的贡献：
\[
A_w(x)=\int_{M_k}w(k)\exp(-i\,\mathbf{k}\cdot x)\,d^3k .
\]
$A_w$ 在我们先前讨论过的 $\mathbb{R}^4$ 上函数空间的表示 $\rho_F$ 下如何变换？
\begin{align*}
\big(\rho_F(b,z)A_w\big)(x)
&=A_w\big((b,z)^{-1}x\big)=A_w\big(f_4(b^{-1})(x-z)\big)\\
&=\int_{M_k}w(k)\exp\big(-i\,\mathbf{k}\cdot(f_4(b^{-1})(x-z))\big)\,d^3k\\
&=\int_{M_k}w(f_4(b)^{-1}k')\exp\big(-i\,(f_4(b)^{-1}k')\cdot(f_4(b^{-1})(x-z))\big)\,d^3k'\\
&\qquad\text{[换积分变量为 $k'=f_4(b)k$]}\\
&=\int_{M_k}w(f_4(b)^{-1}k')\exp\big(-i\,k'\cdot(x-z)\big)\,d^3k'\\
&=\int_{M_k}\exp(i\,k'\cdot z)\,w(f_4(b)^{-1}k')\exp(-i\,k'\cdot x)\,d^3k'\\
&=A_{r(b,z)w}(x).
\end{align*}
对每个正值 $\omega_m$，RSW 方程的平方可积解构成 $\mathbb{R}^4$ 上函数的一个不可约子空间，而 $\rho_F$ 限制到该子空间上等价于 $k=\omega_m$ 的自旋 0 丛表示。自旋 0 丛表示就是“空间上函数表示”的傅里叶变换，说明 RSW 方程的解在写成平面波解之积分时如何变换。

若想包含宇称操作 $P$，我们可以自由选择内禀宇称为 $1$ 或 $-1$ 的表示。

\subsubsection{自旋 1 的丛表示（The Spin-1 Bundle Representation）}

$\mathrm{SU}(2)$ 的自旋 1 表示是三维的，它给出 $\mathrm{SO}(3)$ 的一个真表示，等价于\textbf{基本表示}（fundamental representation）——其中每个 $3\times3$ 实矩阵按矩阵乘法作用在 $\mathbb{C}^3$ 中的向量上。

这里的谜题是：这种三维性如何嵌入\emph{四维}欧氏群的表示？四维空间中的向量在 $\mathrm{SO}(4)$ 下以自然方式变换，那 $\mathrm{SO}(3)$ 又是如何登场的呢？

我们在向量波的注记中见到了一些线索：在那里我们发现可以把黎曼矢量波方程（Riemannian Vector Wave, RVW）的解限制到另外服从一个“横波条件”的解，该条件要求平面波的向量\emph{值}与波的\emph{传播矢量}正交。而正如在讨论函数空间上的表示时所指出的，横波条件从更大的解空间中挑出一个不变子空间。所以，既然已把自旋 0 丛与标量平面波联系起来，我们可以试着把自旋 1 丛与\emph{矢量}平面波联系起来，丛纤维的三维性来自描述波的值的向量应与波的传播矢量正交这一事实。

\orthofig{images/03__014.png}{二维球面的切丛}

事实上——撇开我们为表示所作的丛构造不谈——对于任何像 $M_k$ 这样的三维球面，已经有一个非常自然的三维向量丛的选取，即三维球面的\textbf{切丛}（tangent bundle）。若退一步想想 $\mathbb{R}^3$ 中的二维球面，类似的丛可以很容易地可视化。在球面上每一点，都有一个由所有在该点与曲面相切的向量组成的二维向量空间。这个\textbf{切空间}（tangent space）在不同点是\emph{不同的}向量空间，所以这样向量的整体自然被设想为一个向量丛。有些读者会知道，微分几何给出了切丛的一个精确、内蕴于二维球面或三维球面本身的定义，不依赖于把这些球面看作嵌在 $\mathbb{R}^3$ 或 $\mathbb{R}^4$ 中。但就我们的目的而言，把三维球面 $M_k$ 看作嵌在 $\mathbb{R}^4$ 中、其切向量为字面意义上与三维球面相切、因而与从球心指向表面的向量 $\mathbf{k}$ 正交的\emph{四维}向量，其实非常有用。在这幅图像中，我们可以把 $\mathbf{k}$ 当作平面波的传播矢量，而切向量 $\mathbf{A}_0$ 给出波的四维矢量值：
\[
A_k(x)=A_0\exp(-i\,\mathbf{k}\cdot x),\qquad \mathbf{k}\cdot\mathbf{A}_0=0 .
\]
只差一个小步骤：我们需要允许“复切向量”。这意味着不再把 $\mathbf{A}_0$ 限制在 $\mathbb{R}^4$ 中，而是允许它取 $\mathbb{C}^4$ 中的值。只要我们仍有 $\mathbf{k}\cdot\mathbf{A}_0=0$，数学就同样漂亮地成立。我们称允许这些向量取复值的向量丛为 $M_k$ 的\textbf{复化切丛}（complexified tangent bundle）。

令人惊讶的是，三维球面的切丛实际上是平凡的！二维球面则不是，这与我们可以把任何三维球面认同为群 $\mathrm{SU}(2)$ 有关。我们最初遇到 $\mathrm{SU}(2)$ 是在构造从 $\mathrm{SU}(2)\times\mathrm{SU}(2)$ 到 $\mathrm{SO}(4)$ 的二对一同态时，我们看到 $\mathrm{SU}(2)$ 是我们称之为 $H$ 的矩阵四维向量空间中单位向量的三维球面。现在，用我们那个从 $\mathbb{R}^4$ 到该矩阵空间的映射 $H$，我们有
\[
H(\mathbf{k}_0)=H(0,0,0,k)=k\,I_2 .
\]
若 $G=\mathrm{SU}(2)\times\mathrm{SU}(2)$，我们可以选取 $a(\mathbf{k})=(H(\mathbf{k}/k),I_2)$：
\begin{align*}
f_4(a(k))\mathbf{k}_0
&=H^{-1}\big(H(k/k)\,H(\mathbf{k}_0)\,I_2\big)\\
&=H^{-1}\big(H(k)/k\,(k\,I_2)\,I_2\big)\\
&=\mathbf{k}.
\end{align*}
除了把 $\mathbf{k}_0$ 映到 $\mathbf{k}$ 之外，算符 $f_4(a(\mathbf{k}))$ 还可用来把 $\mathbf{k}_0$ 处的切空间搬到 $\mathbf{k}$ 处的切空间——而且由于这里一切都涉及 $\mathbb{R}^4$ 中坐标的光滑函数，对所有 $\mathbf{k}\in M_k$ 这样做，就能取 $\mathbf{k}_0$ 处切空间的一个基，把它变成覆盖整个三维球面的光滑向量场的一个基。若切丛\emph{不}平凡，任何这样做的尝试都必然在某处涉及不连续（例如，想象在莫比乌斯带上做同样的事）。

我们需要证明自旋 1 丛实际上与 $M_k$ 的复化切丛相同。利用 $a(\mathbf{k})=(H(\mathbf{k}/k),I_2)$ 及 $\eta(h)=f_3(h)=f_4(h,h)$，我们用来把每个纤维与此情形中的表示空间（即 $\mathbb{C}^3$）认同的映射 $\psi_{a(\mathbf{k})}$ 变为
\begin{align*}
\psi_{a(k)}\big([((g,x),w)]\big)
&=\xi\big((a(k),0)^{-1}(g,x)\big)w\\
&=\chi_{\mathbf{k}_0}\big(f_4(a(k)^{-1})x\big)\,f_4(a(k)^{-1}g)w\\
&=\exp(i\,\mathbf{k}\cdot x)\,f_4(a(k)^{-1}g)w .
\end{align*}
由于 $\mathbf{w}\in\mathbb{C}^3$，且 $f_4(a(\mathbf{k})^{-1}g)$ 是保持 $t$ 轴的旋转，这个向量也将落在 $\mathbb{C}^3$ 中。若再用 $f_4(a(\mathbf{k}))$ 作用于这个向量，它将给出 $\mathbb{C}^4$ 中一个与 $\mathbf{k}$ 正交的子空间中的向量。我们定义映射 $\psi':E\to M_k\times\mathbb{C}^4$ 为
\begin{align*}
\psi'\big([((g,x),w)]\big)
&=\big(k,\ f_4(a(k))\,\psi_{a(k)}([((g,x),w)])\big)\\
&=\big(k,\ f_4(a(k))\exp(i\,\mathbf{k}\cdot x)f_4(a(k)^{-1}g)w\big)\\
&=\big(k,\ \exp(i\,\mathbf{k}\cdot x)\,f_4(g)w\big),
\end{align*}
这里我们知道 $f_4(g)\mathbf{w}$ 与 $\mathbf{k}$ 正交，因为 $\mathbf{w}\in\mathbb{C}^3$（即与 $\mathbf{k}_0=(0,0,0,k)$ 正交），而 $f_4(g)$ 保持点积，同时把 $\mathbf{k}_0$ 映到 $\mathbf{k}$。逆映射 $\phi':M_k\times\mathbb{C}^4\to E$ 由
\[
\phi'(k,u)=\big[((a(k),0),\ f_4(a(k)^{-1})u)\big]
\]
给出。这里 $f_4(a(\mathbf{k})^{-1})\mathbf{u}$ 落在 $\mathbb{C}^3$ 中当且仅当 $\mathbf{k}\cdot\mathbf{u}=0$，所以 $\phi'$ 只在 $M_k\times\mathbb{C}^4$ 中我们与 $M_k$ 复化切丛认同的那部分上有定义，而非在整个 $M_k\times\mathbb{C}^4$ 上。

若我们用 $\phi'$ 把复化切丛的一个截面 $f(\mathbf{k})=(\mathbf{k},\mathbf{u}(\mathbf{k}))$ 变成自旋 1 丛的一个截面，用 $S$ 的某元素 $(b,\mathbf{z})$ 经由表示 $r$ 作用到该截面，再用 $\psi'$ 把结果变回切丛的一个截面，则切丛截面上的表示为
\[
\big(r(b,z)f\big)(k)=\big(k,\ \exp(i\,\mathbf{k}\cdot z)\,f_4(b)\,u\big(f_4(b)^{-1}k\big)\big).
\]
我们也可把它重新解释为在满足 $\mathbf{u}(\mathbf{k})\cdot\mathbf{k}=0$ 的\emph{函数} $\mathbf{u}:M_k\to\mathbb{C}^4$ 空间上的表示，并写成
\[
\big(r(b,z)u\big)(k)=\exp(i\,\mathbf{k}\cdot z)\,f_4(b)\,u\big(f_4(b)^{-1}k\big).
\]
截面上的表示与我们从 $S$ 在 $M_k\times\mathbb{C}^4$ 上自然选取的作用
\[
\rho'(b,z)(k,u)=\big(f_4(b)k,\ \exp\big(i\,(f_4(b)k)\cdot z\big)f_4(b)u\big)
\]
所得的相同。

最后，我们可以验证它与 $\rho_V$（我们在 $\mathbb{R}^4$ 上平方可积向量值函数空间上的表示，限制到满足黎曼矢量波方程与横波条件的函数子空间）等价。对每个满足 $\mathbf{u}(\mathbf{k})\cdot\mathbf{k}=0$ 的函数 $\mathbf{u}:M_k\to\mathbb{C}^4$，我们通过傅里叶合成定义 $\mathbb{R}^4$ 上的一个函数：
\[
A_u(x)=\int_{M_k}u(k)\exp(-i\,\mathbf{k}\cdot x)\,d^3k .
\]
$\mathbf{A}_{\mathbf{u}}$ 在 $\rho_V$ 下如何变换？
\begin{align*}
\big(\rho_V(b,z)A_u\big)(x)
&=f_4(b)A_u\big((b,z)^{-1}x\big)=f_4(b)A_u\big(f_4(b^{-1})(x-z)\big)\\
&=f_4(b)\int_{M_k}u(k)\exp\big(-i\,\mathbf{k}\cdot(f_4(b^{-1})(x-z))\big)\,d^3k\\
&=f_4(b)\int_{M_k}u(f_4(b)^{-1}k')\exp\big(-i\,(f_4(b)^{-1}k')\cdot(f_4(b^{-1})(x-z))\big)\,d^3k'\\
&\qquad\text{[换积分变量为 $k'=f_4(b)k$]}\\
&=f_4(b)\int_{M_k}u(f_4(b)^{-1}k')\exp\big(-i\,k'\cdot(x-z)\big)\,d^3k'\\
&=f_4(b)\int_{M_k}\exp(i\,k'\cdot z)\,u(f_4(b)^{-1}k')\exp(-i\,k'\cdot x)\,d^3k'\\
&=A_{r(b,z)u}(x).
\end{align*}
所以自旋 1 丛表示可以看作“$\mathbb{R}^4$ 上向量值函数某子空间上的表示”的傅里叶变换。横波条件
\[
\partial_x A^x+\partial_y A^y+\partial_z A^z+\partial_t A^t=0
\]
的傅里叶变换对应于条件 $\mathbf{u}(\mathbf{k})\cdot\mathbf{k}=0$，因为在任意方向上偏导数的傅里叶变换对应于把变换结果乘以该方向上的频率分量，即 $\partial_x A^x\to k_x u^x$，等等。

与自旋 0 情形一样，若我们要的不仅是 $\mathrm{SE}(4)$ 而是 $\mathrm{E}(4)$ 的表示，我们可以自由选择内禀宇称为 $1$ 或 $-1$ 的表示。若选后者，则 $P$ 的作用会把一切再乘一个因子 $-1$。

\subsubsection{自旋 $\tfrac12$ 的丛表示（The Spin-$\tfrac12$ Bundle Representation）}

$\mathrm{SU}(2)$ 的自旋 $\tfrac12$ 表示是二维的，所以在其截面上我们得到欧氏群双覆盖表示的丛 $E$ 将具有同构于 $\mathbb{C}^2$ 的纤维 $E(\mathbf{k})$。照例底空间 $M_k$ 是一个三维球面。

把这个丛可视化很棘手，但正如对 $M_k$ 的复化切丛一样，如果我们把它看作形如 $M_k\times\mathbb{C}^n$（某 $n$）的平凡丛的一个“子丛”，会有帮助。沿循自旋 1 丛中所见的模式，我们预期把自旋 $\tfrac12$ 丛表示与 $\mathbb{R}^4$ 上 $\mathbb{C}^n$ 值函数的某子空间上的表示联系起来。在自旋 1 情形，若把表示限制到非平移群 $G$、并限制到只在原点非零的函数，我们最终得到 $\mathrm{SO}(4)$ 或 $\mathrm{O}(4)$ 双覆盖的四维矢量表示。若在自旋 $\tfrac12$ 情形照做，类似的结果会是什么？$\mathrm{O}(4)$ 双覆盖最小的费米表示是自旋 $(\tfrac12,0)\oplus$ 自旋 $(0,\tfrac12)$，它也是四维的。所以我们将试着把自旋 $\tfrac12$ 丛认同为平凡丛 $M_k\times\mathbb{C}^4$ 的一个子丛，使得结果得到在原点处（在旋转与反射下）按自旋 $(\tfrac12,0)\oplus(0,\tfrac12)$ 变换的 $\mathbb{C}^4$ 值函数。

$\mathrm{SU}(2)\times\mathrm{SU}(2)$ 在 $\mathbb{C}^4=\mathbb{C}^2\oplus\mathbb{C}^2$ 上的自旋 $(\tfrac12,0)\oplus(0,\tfrac12)$ 表示为
\[
\tau(g_1,g_2)(v_1,v_2)=(g_1 v_1,\ g_2 v_2),
\]
其中 $g_1,g_2\in\mathrm{SU}(2)$、$\mathbf{v}_1,\mathbf{v}_2\in\mathbb{C}^2$。

我们定义映射 $\psi'':E\to M_k\times\mathbb{C}^4$ 为
\[
\psi''\big([((g,x),w)]\big)=\big(k,\ \exp(i\,\mathbf{k}\cdot x)\,\tau(g)(w,w)\big),
\]
其中照例 $\mathbf{k}=\pi([((g,\mathbf{x}),\mathbf{w})])=f_4(g)\mathbf{k}_0$。这里为 $\mathrm{SU}(2)\times\mathrm{SU}(2)$ 的元素我们只写 $g$，而不显式写成 $\mathrm{SU}(2)$ 中两个独立元素的对。此定义无歧义，因为若取同一等价类的另一代表元，我们得到相同结果：
\begin{align*}
&\psi''\big([((g,x)(h,y)^{-1},\ \xi(h,y)w)]\big)\\
&=\psi''\big([((gh^{-1},\ -f_4(gh^{-1})y+x),\ \exp(i\,\mathbf{k}_0\cdot y)\eta(h)w)]\big)\\
&=\Big(f_4(gh^{-1})\mathbf{k}_0,\ \exp\big(i\,(f_4(gh^{-1})\mathbf{k}_0)\cdot(-f_4(gh^{-1})y+x)\big)\exp(i\,\mathbf{k}_0\cdot y)\,\tau(gh^{-1})(\eta(h)w,\eta(h)w)\Big)\\
&=\Big(f_4(g)\mathbf{k}_0,\ \exp(-i\,\mathbf{k}_0\cdot y)\exp\big(i\,(f_4(g)\mathbf{k}_0)\cdot x\big)\exp(i\,\mathbf{k}_0\cdot y)\,\tau(g)\tau(h^{-1})(hw,hw)\Big)\\
&=\Big(f_4(g)\mathbf{k}_0,\ \exp\big(i\,(f_4(g)\mathbf{k}_0)\cdot x\big)\,\tau(g)(w,w)\Big)\\
&=\big(k,\ \exp(i\,\mathbf{k}\cdot x)\,\tau(g)(w,w)\big).
\end{align*}

接下来定义 $S$ 在 $M_k\times\mathbb{C}^4$ 上的左作用 $\rho''$：
\[
\rho''(b,z)\big(k,(v,w)\big)=\big(f_4(b)k,\ \exp\big(i\,(f_4(b)k)\cdot z\big)\,\tau(b)(v,w)\big).
\]
自旋 $\tfrac12$ 丛 $E$ 上通常的作用 $\rho$ 为
\[
\rho(b,z)\,[((g,x),w)]=[((b,z)(g,x),w)].
\]
$\psi''$ 与这些丛作用对易：
\[
\rho''(b,z)\,\psi''\big([((g,x),w)]\big)=\psi''\big(\rho(b,z)[((g,x),w)]\big).
\]
由作用 $\rho''$ 得到 $S$ 在 $M_k\times\mathbb{C}^4$ 的截面 $f(\mathbf{k})=(\mathbf{k},\mathbf{u}(\mathbf{k}))$ 上的表示：
\[
\big(r(b,z)f\big)(k)=\big(k,\ \exp(i\,\mathbf{k}\cdot z)\,\tau(b)\,u(f_4(b)^{-1}k)\big).
\]
更简洁地，在函数 $\mathbf{u}:M_k\to\mathbb{C}^4$ 上的表示为
\begin{equation}
\big(r(b,z)u\big)(k)=\exp(i\,\mathbf{k}\cdot z)\,\tau(b)\,u\big(f_4(b)^{-1}k\big). \tag{4.2.7}
\end{equation}
当然 $\psi''$ 的像不是整个 $M_k\times\mathbb{C}^4$，所以我们需要限制函数 $\mathbf{u}$。在自旋 1 情形，子丛由条件 $\mathbf{u}(\mathbf{k})\cdot\mathbf{k}=0$ 指定，我们需要找出自旋 $\tfrac12$ 情形中 $\mathbf{u}(\mathbf{k})$ 与 $\mathbf{k}$ 满足的类似条件。首先注意纤维 $E(\mathbf{k}_0)$ 的像是
\[
\psi''\big(E(\mathbf{k}_0)\big)=\{\psi''([((I,0),w)])\mid w\in\mathbb{C}^2\}=\{(\mathbf{k}_0,(w,w))\mid w\in\mathbb{C}^2\}.
\]
用 $\rho''(b,\mathbf{z})$ 作用在形如上述的元素上，得到
\[
\psi''(E)=\Big\{\Big(f_4(b)\mathbf{k}_0,\ \exp\big(i\,(f_4(b)\mathbf{k}_0)\cdot z\big)\,\tau(b)(w,w)\Big)\ \Big|\ w\in\mathbb{C}^2,\ b\in\mathrm{SU}(2)\times\mathrm{SU}(2),\ z\in\mathbb{R}^4\Big\}.
\]
若 $b=(b_1,b_2)$，其中 $b_1,b_2\in\mathrm{SU}(2)$，则
\[
f_4(b_1,b_2)\mathbf{k}_0=H^{-1}\big(b_1 H(\mathbf{k}_0)b_2^{-1}\big)=k\,H^{-1}(b_1 b_2^{-1}),
\]
故 $H(f_4(b_1,b_2)\mathbf{k}_0)=k\,b_1 b_2^{-1}$，且
\[
\tau(b_1,b_2)(w,w)=(b_1 w,\ b_2 w).
\]
把最后两个结果合起来，对 $\psi''(E)$ 中每个 $(\mathbf{k},(\mathbf{u}_1,\mathbf{u}_2))$ 有
\[
H(\mathbf{k})\,\mathbf{u}_2=k\,\mathbf{u}_1 .
\]
左乘伴随 $H(\mathbf{k})^*$ 得
\begin{align*}
H(\mathbf{k})^* H(\mathbf{k})\,\mathbf{u}_2&=k\,H(\mathbf{k})^*\mathbf{u}_1\\
k^2\,H(\mathbf{k}/k)^* H(\mathbf{k}/k)\,\mathbf{u}_2&=k\,H(\mathbf{k})^*\mathbf{u}_1\\
k^2\,\mathbf{u}_2&=k\,H(\mathbf{k})^*\mathbf{u}_1\\
H(\mathbf{k})^*\mathbf{u}_1&=k\,\mathbf{u}_2 .
\end{align*}
这些方程可以打包成一个矩阵方程：
\[
\begin{pmatrix}-k\,I_2 & H(\mathbf{k})\\ H(\mathbf{k})^* & -k\,I_2\end{pmatrix}
\begin{pmatrix}\mathbf{u}_1\\ \mathbf{u}_2\end{pmatrix}
=\begin{pmatrix}\mathbf{0}\\ \mathbf{0}\end{pmatrix}.
\]
我们也可把它写成
\[
\big(k_\mu\gamma^\mu-k\big)u(k)=0,
\]
其中对 $\mu$ 求和，并定义矩阵 $\gamma^\mu$ 为
\[
\gamma^\mu=\begin{pmatrix}0 & H_\mu\\ H_\mu^* & 0\end{pmatrix}.
\]
我们不会写出由此定义得到的那些具体的矩阵，因为它们依赖于我们在把自旋 $\tfrac12$ 丛与 $M_k\times\mathbb{C}^4$ 的某子丛认同时所选取的基；而在后面讨论狄拉克方程时，我们将集中于对特定应用有用的两个特殊基的选取。$\gamma$ 矩阵真正重要的性质是一个不依赖于基选取的性质：
\[
\gamma^\mu\gamma^\nu+\gamma^\nu\gamma^\mu=2\,\delta^{\mu\nu}I_4 .
\]
从 $H$ 矩阵的定义与性质不难验证这一点。

在傅里叶合成
\[
A_u(x)=\int_{M_k}u(k)\exp(-i\,\mathbf{k}\cdot x)\,d^3k
\]
中对每个平面波施加条件
\[
\big(k_\mu\gamma^\mu-k\big)u(k)=0,
\]
等同于要求 $\mathbf{A}_{\mathbf{u}}(\mathbf{x})$ 满足微分方程
\[
i\,\gamma^\mu\partial_\mu A_u(x)-k\,A_u(x)=0 .
\]
若把 $k$（传播矢量的模）认同为粒子的质量 $m$，这恰恰就是狄拉克方程（Dirac equation）。于是狄拉克方程的解构成 $\mathbb{R}^4$ 上 $\mathbb{C}^4$ 值函数的一个子空间，它按自旋 $\tfrac12$ 丛表示变换。

\subsection{射影正交群（The Projective Orthogonal Groups）}

关于欧氏四维空间中的对称性，我们还需要考察一个进一步的变体。设有一个\emph{中性}经典粒子在四维空间中自由穿行，其世界线是一条无限长的直线。我们可以给这个粒子赋予一个沿其世界线指向某特定方向的能动量矢量，但只有在我们选定一个参考系、并出于方便或约定决定能动量矢量在该参考系中应有正的时间坐标时才能如此。实际上并\emph{不存在}这样一个与粒子本身相关联的矢量。世界线确实在四维空间中有一个方向——在相差一个符号的意义上——粒子的质量也决定了其能动量的大小，但我们所选的任何能动量矢量都可以被其相反矢量同样好地取代。

这表明对一个自由中性粒子而言，$\mathrm{SO}(4)$ 其实并不是正确的对称群。我们真正想做的是，把“反转所有矢量”这一操作视为对该系统\emph{完全没有影响}。当然，对带电粒子就不合适了，因为它的电荷在此操作下会改变符号。

本质上，我们想做的，是把 $\mathbb{R}^4$ 上通常的恒等算符 $I_4$ \emph{以及} $-I_4$（反转所有矢量的方向）\emph{二者}都对应于某取代 $\mathrm{SO}(4)$ 的新对称群的单位元。我们能做到这一点，不是通过限制到 $\mathrm{SO}(4)$ 的某个子群，而是通过构造一个商群（quotient group）。取 $\mathrm{SO}(4)$ 的子群 $\mathrm{SZO}(4)=\{I_4,-I_4\}$，然后构造新群 $\mathrm{PSO}(4)=\mathrm{SO}(4)/\mathrm{SZO}(4)$，它由位于 $\mathrm{SZO}(4)$ 同一陪集内的等价类组成。子群的陪集，是用群元素从左边或右边乘该子群所得的集合；此处 $\mathrm{SZO}(4)$ 的元素与一切对易，故左、右陪集没有差别。于是每个陪集形如 $g\,\mathrm{SZO}(4)=\{g,-g\}$（$g\in\mathrm{SO}(4)$）。我们将它写作 $\pm g$；对任意 $g\in\mathrm{SO}(4)$，$\pm g$ 是 $\mathrm{PSO}(4)$ 的一个元素。我们用 $\mathrm{SO}(4)$ 中的乘法来乘这些元素，故
\[
(\pm g)(\pm h)=\pm(gh).
\]
显然 $\pm I_4$ 充当单位元，逆元为
\[
(\pm g)^{-1}=\pm(g^{-1}).
\]
我们构造出的这个新群 $\mathrm{PSO}(4)$ 称为四维射影特殊正交群（projective special orthogonal group in four dimensions）。我们可以对 $\mathrm{O}(4)$ 作类似构造，得到射影一般正交群 $\mathrm{PGO}(4)=\mathrm{O}(4)/\{I_4,-I_4\}$，它包含反射。谈 $\mathrm{PGO}(3)$（商群 $\mathrm{O}(3)/\{I_3,-I_3\}$）也说得通。由于 $-I_3$ 的行列式为 $-1$，它不是 $\mathrm{SO}(3)$ 的元素，所以不存在 $\mathrm{PSO}(3)$ 这样的东西——尽管按某些定义，$\mathrm{PSO}(3)$ 就直接定义为 $\mathrm{SO}(3)/\{I_3\}=\mathrm{SO}(3)$。

由
\[
\pi(g)=\pm g
\]
定义的映射 $\pi:\mathrm{SO}(4)\to\mathrm{PSO}(4)$ 是一个从 $\mathrm{SO}(4)$ 到 $\mathrm{PSO}(4)$ 的二对一同态，所以 $\mathrm{SO}(4)$ 是 $\mathrm{PSO}(4)$ 的双覆盖，就像 $\mathrm{SU}(2)\times\mathrm{SU}(2)$ 是 $\mathrm{SO}(4)$ 的双覆盖一样。更有甚者，我们可以把二对一同态 $\pi$ 与二对一同态 $f_4$ 复合，得到一个从 $\mathrm{SU}(2)\times\mathrm{SU}(2)$ 到 $\mathrm{PSO}(4)$ 的四对一同态 $\alpha$：
\[
\alpha(g,h)=\pi\big(f_4(g,h)\big)=\pm f_4(g,h).
\]
$\mathrm{SU}(2)\times\mathrm{SU}(2)$ 的四个元素 $(g,h)$、$(-g,h)$、$(g,-h)$、$(-g,-h)$ 在 $\alpha$ 下的像都落在 $\mathrm{PSO}(4)$ 中同一点。经由 $\alpha$，我们可以像用 $\mathrm{SU}(2)\times\mathrm{SU}(2)$ 的自旋 $(j,k)$ 酉表示得到 $\mathrm{SO}(4)$ 的表示那样，用它们得到 $\mathrm{PSO}(4)$ 的（一般而言是\emph{射影}的）表示。回忆自旋 $(j,k)$ 表示定义为
\[
\rho_{j,k}(g,h)\,x\otimes y=(\rho_j(g)x)\otimes(\rho_k(h)y),
\]
且
\[
\rho_j(-g)=(-1)^{2j}\rho_j(g).
\]
这表明当 $j$ 与 $k$ 都为整数时，四个算符 $\rho_{j,k}(\pm g,\pm h)$ 全都相同，所以 $\rho_{j,k}$ 将给出 $\mathrm{PSO}(4)$ 的一个真表示，而不只是射影表示。

\section{李代数（Lie Algebras）}

$n$ 维旋转群 SO($n$) 是所谓\textbf{李群}（Lie group）的一个例子：群元素可以被赋予坐标，使得群运算（两个元素相乘、元素取逆）都是这些坐标的光滑函数。为避免出现类似地球极点与 $180^\circ$ 经线处的坐标退化问题，有时需要使用不止一张坐标“图卡”（chart）。以三维旋转群 SO(3) 为例，它可以设想为一个半径 $\pi$ 的实心三维球，球面上相对的两点彼此认同：从球心指向球内某点的矢量给出旋转轴，该矢量的长度给出旋转角。这个集合可以用一张以恒等元素为中心的坐标图卡覆盖，再用另一张以某个 $\pi$ 旋转为中心的图卡覆盖。

现在设 $n$ 维空间中的一个刚体随时间转动，其质心固定于坐标原点。在任意时刻 $t$，都存在某个矩阵 $R(t)$，把物体上每一点在 $t=0$ 时的位置映到它在 $t$ 时的新位置。若物体光滑地转动，则 $R(t)$ 的分量是 $t$ 的光滑函数，我们可以定义 $R'(t)$——其各分量是 $R$ 的各分量对 $t$ 的导数所构成的矩阵。

$R'(t)$ 显然与转动体的\textbf{角速度}（angular velocity）有关。在三维中我们惯于把转动与某一根转轴对应，并把角速度描述为沿该轴的矢量，其长度与转动速率成正比。但在四维中，转动不会留下唯一的固定轴——它要么不固定任何矢量，要么使一整个平面保持不变。因此一般而言，角速度最好用一个矩阵 $\Omega$ 来描述：当它作用于从转心指向物体上某点的位移矢量时，给出该点的瞬时线速度：

\begin{equation}
\Omega\,\mathbf{x} = \mathbf{v}
\end{equation}

物体上在时刻 $t_1$ 位于 $\mathbf{x}_1$ 的点，其原始位置（$t=0$ 时）为 $\mathbf{x}_0 = R(t_1)^{-1}\mathbf{x}_1$，于是同一点的一般位置可写为

\begin{equation}
\mathbf{x}(t) = R(t)\,R(t_1)^{-1}\,\mathbf{x}_1
\end{equation}

其一般线速度为

\begin{equation}
\mathbf{v}(t) = R'(t)\,R(t_1)^{-1}\,\mathbf{x}_1
\end{equation}

特别地，在 $t=t_1$ 时：

\begin{equation}
\mathbf{v}(t_1) = R'(t_1)\,R(t_1)^{-1}\,\mathbf{x}_1
\end{equation}

于是我们看到：

\begin{equation}
\Omega(t_1) = R'(t_1)\,R(t_1)^{-1}
\end{equation}

$R'(t)$ 与角速度有关，但为了得到 $\Omega$，我们必须“撤销”到目前为止所发生的净转动。或者，如果我们把时间坐标的原点选得使 $t_1=0$，那么 $R(t_1)^{-1}$ 就是单位矩阵，于是 $\Omega(0)=R'(0)$。

无论如何，既然我们知道所有矩阵 $R(t)$ 都是正交矩阵——即其转置等于其逆——我们有：

\begin{align}
R(t)\,R(t)^{T} &= I_n \\
\partial_t\!\left(R(t)\,R(t)^{T}\right) &= 0 \\
R'(t)\,R(t)^{T} + R(t)\,R'(t)^{T} &= 0 \\
R'(t)\,R(t)^{T} &= -R(t)\,R'(t)^{T} \\
\Omega(t) &= -\Omega(t)^{T}
\end{align}

这表明角速度 $\Omega$ 是一个\textbf{反对称矩阵}（antisymmetric matrix）：它的转置等于原矩阵的相反数。

例如，设

\begin{equation}
R(t) = \begin{pmatrix}
\cos(\omega t) & -\sin(\omega t) & 0 \\
\sin(\omega t) & \cos(\omega t) & 0 \\
0 & 0 & 1
\end{pmatrix}
\end{equation}

这些是绕 $z$ 轴以角速度 $\omega$ 转动的物体随时间的旋转。我们对时间求导再令时间为零，得到角速度矩阵：

\begin{equation}
R'(t) = \begin{pmatrix}
-\omega\sin(\omega t) & -\omega\cos(\omega t) & 0 \\
\omega\cos(\omega t) & -\omega\sin(\omega t) & 0 \\
0 & 0 & 0
\end{pmatrix}
\end{equation}

\begin{equation}
\Omega(0) = R'(0) = \begin{pmatrix}
0 & -\omega & 0 \\
\omega & 0 & 0 \\
0 & 0 & 0
\end{pmatrix}
\end{equation}

我们把 $n\times n$ 反对称矩阵的集合称为与 SO($n$) 相对应的李代数，并按惯例把群名改为小写来命名它：so($n$)。李代数是一个矢量空间：我们可以把元素相加，或乘以标量（此处为实数），结果仍是同一李代数中的元素。反对称矩阵相加或乘以实数，所得矩阵仍反对称。

给定李代数中的两个元素 $\mathbf{v}$ 与 $\mathbf{w}$，我们还可以施行称为\textbf{李括号}（Lie bracket）或\textbf{交换子}（commutator）的运算：

\begin{equation}
[\mathbf{v},\mathbf{w}] = \mathbf{v}\mathbf{w} - \mathbf{w}\mathbf{v}
\end{equation}

其中乘积 $\mathbf{v}\mathbf{w}$ 与 $\mathbf{w}\mathbf{v}$ 按通常的矩阵乘法构成。容易看出两个反对称矩阵的李括号仍是反对称的：

\begin{align}
[\mathbf{v},\mathbf{w}]^{T} &= (\mathbf{v}\mathbf{w} - \mathbf{w}\mathbf{v})^{T} \\
&= (\mathbf{v}\mathbf{w})^{T} - (\mathbf{w}\mathbf{v})^{T} \\
&= \mathbf{w}^{T}\mathbf{v}^{T} - \mathbf{v}^{T}\mathbf{w}^{T} \\
&= \mathbf{w}\mathbf{v} - \mathbf{v}\mathbf{w} \\
&= -[\mathbf{v},\mathbf{w}]
\end{align}

一般而言，我们可以把与任一李群相联系的李代数，设想为所有穿过该群恒等元素的光滑路径的可能切矢所组成的矢量空间。对旋转而言，可以把群中的路径设想为某个刚体随时间的转动、把恒等元素处的切矢设想为描述该刚体在某一时刻的角速度，来使这一概念更具体；但同样的想法也适用于像 SU(2) 这样更抽象的群。

给定李代数 $g$ 的一个元素 $\mathbf{v}$，存在一个从 $g$ 到李群 $G$ 的函数，称为\textbf{指数映射}（exponential map），它把 $\mathbf{v}$ 的倍数映为一条经过恒等元素、切矢为 $\mathbf{v}$ 的路径。这可以非常一般地做到，但我们将集中讨论 $G$ 与 $g$ 的元素都是矩阵的情形。此时指数映射就是通常的指数函数，只是被推广到矩阵。这既可以用幂级数来做：

\begin{equation}
\exp(A) = I_n + A + \frac{A^{2}}{2!} + \frac{A^{3}}{3!} + \frac{A^{4}}{4!} + \cdots
\end{equation}

也可以在许多情形下，选择一组使矩阵 $A$ 呈对角形式 $\diag(\lambda_1,\dots,\lambda_n)$ 的基，然后只需对各对角元取指数得到 $\diag(\exp(\lambda_1),\dots,\exp(\lambda_n))$，再变换回原来的基。

例如，若我们对反对称角速度矩阵 $\Omega$ 的某一倍数取指数，就得到一个正交旋转矩阵：

\begin{equation}
R(t) = \exp(t\Omega)
\end{equation}

它给出经过旋转群的一条路径。我们在恒等元素处求这条路径的切矢，只需对 $t$ 求导再令 $t=0$：

\begin{align}
R'(t) &= \Omega\exp(t\Omega) \\
R'(0) &= \Omega
\end{align}

我们已经指出，李代数 so($n$) 由 $n\times n$ 反对称矩阵组成。那么与群 SU($n$) 相联系的李代数 su($n$) 又是什么呢？SU($n$) 中的矩阵是酉矩阵，即其转置的复共轭（也称为厄米共轭）等于其逆。于是，若有一条经过 SU($n$) 的路径 $M(t)$ 且 $M(0)=I_n$，并把厄米共轭记作 $M(t)^{*}$，我们有：

\begin{align}
M(t)\,M(t)^{*} &= I_n \\
\partial_t\!\left(M(t)\,M(t)^{*}\right) &= 0 \\
M'(t)\,M(t)^{*} + M(t)\,M'(t)^{*} &= 0 \\
M'(t)\,M(t)^{*} &= -M(t)\,M'(t)^{*} \\
M'(0) &= -M'(0)^{*}
\end{align}

因此李代数元素 $m=M'(0)$ 的厄米共轭等于原元素的相反数。具有这一性质的矩阵称为\textbf{反厄米矩阵}（skew-Hermitian）。

但 su($n$) 并不包含\emph{全部} $n\times n$ 反厄米矩阵。原来，矩阵的迹与其指数的行列式之间存在一个关系：

\begin{equation}
\det(\exp(A)) = \exp(\Tr A)
\end{equation}

对对角矩阵不难看出为什么成立。若 $A=\diag(\lambda_1,\lambda_2,\dots,\lambda_n)$，则

\begin{align}
\Tr A &= \lambda_1 + \lambda_2 + \dots + \lambda_n \\
\exp(A) &= \diag(\exp(\lambda_1),\exp(\lambda_2),\dots,\exp(\lambda_n)) \\
\det(\exp(A)) &= \exp(\lambda_1)\times\exp(\lambda_2)\times\dots\times\exp(\lambda_n) \\
&= \exp(\Tr A)
\end{align}

由于 SU($n$) 的元素行列式为 1，su($n$) 的元素的迹的指数必须为 1，即迹为零。因此 su($n$) 由\textbf{无迹、反厄米的 $n\times n$ 矩阵}组成。

\subsection{so(3) 与 su(2)}

我们知道群 SU(2) 构成群 SO(3) 的双覆盖。当一个李群是另一个的李群的双覆盖时，它们的李代数总是“相同的”，即我们可以找到二者之间的一个李代数同构：一个保持李括号的可逆线性映射。

我们之前如下定义了 $f_3:\mathrm{SU}(2)\to\mathrm{SO}(3)$：

\begin{equation}
f_3(g)\,\mathbf{x} = H^{-1}\!\left(g\,H(\mathbf{x})\,g^{-1}\right)
\end{equation}

其中对 $\mathbb{R}^{3}$ 中的任一点 $\mathbf{x}=(x,y,z)$，我们定义：

\begin{equation}
H(\mathbf{x}) = \begin{pmatrix}
-z\,\ii & -y - x\,\ii \\
y - x\,\ii & z\,\ii
\end{pmatrix}
\end{equation}

在此三维情形里，$H(\mathbf{x})$ 无迹且反厄米，因而是 su(2) 的一个元素。所以我们有 $H:\mathbb{R}^{3}\to\mathrm{su}(2)$。

设 SU(2) 中的群元素 $g$ 由 su(2) 的元素 $\mathbf{v}$ 的倍数取指数得到：

\begin{equation}
g(t) = \exp(t\mathbf{v})
\end{equation}

我们便可以这样求出与 su(2) 中 $\mathbf{v}$ 相对应的 so(3) 元素：对 $f_3(g)\mathbf{x}$ 关于 $t$ 求导，再令 $t=0$。

\begin{align}
\partial_t\!\left(f_3(g(t))\,\mathbf{x}\right)
&= \partial_t\!\left(H^{-1}\!\left(g(t)\,H(\mathbf{x})\,g(t)^{-1}\right)\right) \\
&= H^{-1}\!\left(\partial_t\!\left(g(t)\,H(\mathbf{x})\,g(t)^{-1}\right)\right) \\
&= H^{-1}\!\left(\partial_t\!\left(\exp(t\mathbf{v})\,H(\mathbf{x})\,\exp(-t\mathbf{v})\right)\right) \\
&= H^{-1}\!\left(\mathbf{v}\exp(t\mathbf{v})\,H(\mathbf{x})\,\exp(-t\mathbf{v}) - \exp(t\mathbf{v})\,H(\mathbf{x})\,\mathbf{v}\exp(-t\mathbf{v})\right)
\end{align}

令 $t=0$，得：

\begin{align}
\left.\partial_t\!\left(f_3(g(t))\,\mathbf{x}\right)\right|_{t=0}
&= H^{-1}\!\left(\mathbf{v}\,H(\mathbf{x}) - H(\mathbf{x})\,\mathbf{v}\right) \\
&= H^{-1}\!\left([\mathbf{v},H(\mathbf{x})]\right)
\end{align}

现在我们可以定义映射 $df_3:\mathrm{su}(2)\to\mathrm{so}(3)$，即双覆盖的\emph{导数}：

\begin{equation}
df_3(\mathbf{v})\,\mathbf{x} = H^{-1}\!\left([\mathbf{v},H(\mathbf{x})]\right)
\end{equation}

让我们列出把 $H$ 作用于 $\mathbb{R}^{3}$ 的基矢量所得到的矩阵，并把它们用作 su(2) 的一组基：

\begin{align}
H_x = H(\mathbf{e}_x) &= \begin{pmatrix} 0 & -\ii \\ -\ii & 0 \end{pmatrix} \\
H_y = H(\mathbf{e}_y) &= \begin{pmatrix} 0 & -1 \\ 1 & 0 \end{pmatrix} \\
H_z = H(\mathbf{e}_z) &= \begin{pmatrix} -\ii & 0 \\ 0 & \ii \end{pmatrix}
\end{align}

这些矩阵之间的交换子（李括号）为：

\begin{align}
[H_x,H_y] &= 2H_z \\
[H_y,H_z] &= 2H_x \\
[H_z,H_x] &= 2H_y
\end{align}

若交换李括号内元素的顺序则结果反号，而任一元素与自身的李括号自然为零。由此推出：

\begin{align}
df_3(H_x/2)\,\mathbf{e}_x &= H^{-1}\!\left([H_x/2,H_x]\right) = 0 \\
df_3(H_x/2)\,\mathbf{e}_y &= H^{-1}\!\left([H_x/2,H_y]\right) = \mathbf{e}_z \\
df_3(H_x/2)\,\mathbf{e}_z &= H^{-1}\!\left([H_x/2,H_z]\right) = -\mathbf{e}_y
\end{align}

于是 $df_3$ 把 su(2) 中的 $H_x/2$ 映为 so(3) 中生成绕 $x$ 轴旋转的元素 $J_x$（即若对 $J_x$ 的任一倍数取指数，就得到绕 $x$ 轴的旋转）。对所有基元素显式写出，我们有：

\begin{align}
df_3(H_x/2) = J_x &= \begin{pmatrix} 0 & 0 & 0 \\ 0 & 0 & -1 \\ 0 & 1 & 0 \end{pmatrix} \\
df_3(H_y/2) = J_y &= \begin{pmatrix} 0 & 0 & 1 \\ 0 & 0 & 0 \\ -1 & 0 & 0 \end{pmatrix} \\
df_3(H_z/2) = J_z &= \begin{pmatrix} 0 & -1 & 0 \\ 1 & 0 & 0 \\ 0 & 0 & 0 \end{pmatrix}
\end{align}

各坐标轴旋转生成元 $J$ 矩阵之间的交换子为：

\begin{align}
[J_x,J_y] &= J_z \\
[J_y,J_z] &= J_x \\
[J_z,J_x] &= J_y
\end{align}

那些烦人的因子 2 与 $1/2$ 最终都相互抵消，映射 $df_3:\mathrm{su}(2)\to\mathrm{so}(3)$ 保持李括号：

\begin{equation}
[df_3(\mathbf{v}), df_3(\mathbf{w})] = df_3([\mathbf{v},\mathbf{w}])
\end{equation}

\subsection{so(4)}

与四维欧氏空间中旋转所构成的李群 SO(4) 相联系的李代数 so(4)，是一个六维矢量空间。我们可以通过数 $4\times 4$ 反对称矩阵的自由度看出这一点（例如矩阵对角线上方有 6 个元素）。这些元素的一组基由六个坐标平面（每对坐标轴确定一个平面）中各自的旋转生成元组成。注意我们把坐标排序为 $t,x,y,z$，因此这些矩阵的第一行与第一列涉及 $t$。

\begin{align}
J_{xy} &= \begin{pmatrix} 0 & 0 & 0 & 0 \\ 0 & 0 & -1 & 0 \\ 0 & 1 & 0 & 0 \\ 0 & 0 & 0 & 0 \end{pmatrix} \\
J_{yz} &= \begin{pmatrix} 0 & 0 & 0 & 0 \\ 0 & 0 & 0 & 0 \\ 0 & 0 & 0 & -1 \\ 0 & 0 & 1 & 0 \end{pmatrix} \\
J_{zx} &= \begin{pmatrix} 0 & 0 & 0 & 0 \\ 0 & 0 & 0 & 1 \\ 0 & 0 & 0 & 0 \\ 0 & -1 & 0 & 0 \end{pmatrix} \\
J_{xt} &= \begin{pmatrix} 0 & 1 & 0 & 0 \\ -1 & 0 & 0 & 0 \\ 0 & 0 & 0 & 0 \\ 0 & 0 & 0 & 0 \end{pmatrix} \\
J_{yt} &= \begin{pmatrix} 0 & 0 & 1 & 0 \\ 0 & 0 & 0 & 0 \\ -1 & 0 & 0 & 0 \\ 0 & 0 & 0 & 0 \end{pmatrix} \\
J_{zt} &= \begin{pmatrix} 0 & 0 & 0 & 1 \\ 0 & 0 & 0 & 0 \\ 0 & 0 & 0 & 0 \\ -1 & 0 & 0 & 0 \end{pmatrix}
\end{align}

这里 $J_{ij}$ 把 $\mathbf{e}_i$ 映到 $\mathbf{e}_j$，把 $\mathbf{e}_j$ 映到 $-\mathbf{e}_i$。与上面所列坐标指标顺序相反的配对给出相反的矩阵，例如 $J_{tz}=-J_{zt}$。这个李代数可以这样概括：当 $i,j,k$ 互不相同时，

\begin{equation}
[J_{ij}, J_{kj}] = J_{ik}
\end{equation}

除这个基本恒等式及其带适当符号变化的重新排列外，所有其他的交换子都为零。

这已经足够简单，但原来还存在另一组基的选择，它揭示了 so(4) 的一个关键性质。

\begin{align}
A_x = \tfrac12(J_{yz}+J_{xt}) &= \begin{pmatrix} 0 & \half & 0 & 0 \\ -\half & 0 & 0 & 0 \\ 0 & 0 & 0 & -\half \\ 0 & 0 & \half & 0 \end{pmatrix} \\
A_y = \tfrac12(J_{zx}+J_{yt}) &= \begin{pmatrix} 0 & 0 & \half & 0 \\ 0 & 0 & 0 & \half \\ -\half & 0 & 0 & 0 \\ 0 & -\half & 0 & 0 \end{pmatrix} \\
A_z = \tfrac12(J_{xy}+J_{zt}) &= \begin{pmatrix} 0 & 0 & 0 & \half \\ 0 & 0 & -\half & 0 \\ 0 & \half & 0 & 0 \\ -\half & 0 & 0 & 0 \end{pmatrix} \\
B_x = \tfrac12(J_{yz}-J_{xt}) &= \begin{pmatrix} 0 & -\half & 0 & 0 \\ \half & 0 & 0 & 0 \\ 0 & 0 & 0 & -\half \\ 0 & 0 & \half & 0 \end{pmatrix} \\
B_y = \tfrac12(J_{zx}-J_{yt}) &= \begin{pmatrix} 0 & 0 & -\half & 0 \\ 0 & 0 & 0 & \half \\ \half & 0 & 0 & 0 \\ 0 & -\half & 0 & 0 \end{pmatrix} \\
B_z = \tfrac12(J_{xy}-J_{zt}) &= \begin{pmatrix} 0 & 0 & 0 & -\half \\ 0 & 0 & -\half & 0 \\ 0 & \half & 0 & 0 \\ \half & 0 & 0 & 0 \end{pmatrix}
\end{align}

三个 $A$ 矩阵由某个 $J$ 矩阵与它的\textbf{对偶}（dual，即第一个矩阵未曾触及的那个平面中旋转的生成元）之和的一半构成。（注意这里的“对偶”与前面用过的含义不同，前面指的是矢量与矢量的线性函数之间的关系。）\textbf{B} 矩阵遵循同样的配方，但用的是差而非和。结果是：若把“取对偶”看作 so(4) 上的一个线性算符（存在一整套优美且自洽的语境使这说得通，但在此展开会离题太远），则 $A$ 矩阵是“自对偶”的，即等于自身的对偶；而 $B$ 矩阵是“反自对偶”的，即等于自身对偶的相反数。

在这组基下，交换子为：

\begin{align}
[A_x,A_y] &= A_z, \quad \text{及其循环置换} \\
[B_x,B_y] &= B_z, \quad \text{及其循环置换} \\
[A_i,B_j] &= 0, \quad \text{对所有 } i,j
\end{align}

因此 $A$ 矩阵与 $B$ 矩阵各自构成的三维李子代数的结构与 so(3) 或 su(2) 完全相同，而两个子代数之间的所有交换子都为零。我们把这一结果写成：

\begin{equation}
\mathrm{so}(4) = \mathrm{so}(3)\oplus\mathrm{so}(3)
\end{equation}

so(4) 如此分裂为两个 so(3)（或 su(2)）副本，与 $\mathrm{SU}(2)\times\mathrm{SU}(2)$ 是 SO(4) 的双覆盖这一事实直接相关。我们之前如下定义了函数 $f_4:\mathrm{SU}(2)\times\mathrm{SU}(2)\to\mathrm{SO}(4)$：

\begin{equation}
f_4(g,h)\,\mathbf{x} = H^{-1}\!\left(g\,H(\mathbf{x})\,h^{-1}\right)
\end{equation}

其中对 $\mathbb{R}^{4}$ 中的 $\mathbf{x}=(x,y,z,t)$，我们定义 $H(\mathbf{x})$ 为：

\begin{equation}
H(\mathbf{x}) = \begin{pmatrix}
t - z\,\ii & -y - x\,\ii \\
y - x\,\ii & t + z\,\ii
\end{pmatrix}
\end{equation}

若令 $g(t)=\exp(t\mathbf{v})$、$h(t)=\exp(t\mathbf{w})$，其中 $\mathbf{v},\mathbf{w}$ 是 su(2) 的元素，则我们可以对 $t$ 求导再令 $t=0$，得到映射 $df_4:\mathrm{su}(2)\oplus\mathrm{su}(2)\to\mathrm{so}(4)$，定义为：

\begin{equation}
df_4(\mathbf{v},\mathbf{w})\,\mathbf{x} = H^{-1}\!\left(\mathbf{v}\,H(\mathbf{x}) - H(\mathbf{x})\,\mathbf{w}\right)
\end{equation}

利用下面的乘法表（表中第 $i$ 行第 $j$ 列为 $H_i H_j$）：

\begin{equation}
\begin{array}{c|cccc}
 & H_x & H_y & H_z & H_t \\ \hline
H_x & -H_t & H_z & -H_y & H_x \\
H_y & -H_z & -H_t & H_x & H_y \\
H_z & H_y & -H_x & -H_t & H_z \\
H_t & H_x & H_y & H_z & H_t
\end{array}
\end{equation}

我们发现：

\begin{align}
df_4(0, H_x/2)\,\mathbf{e}_x &= H^{-1}\!\left(-H_x^{2}/2\right) = \half\,\mathbf{e}_t \\
df_4(0, H_x/2)\,\mathbf{e}_y &= H^{-1}\!\left(-H_y H_x/2\right) = \half\,\mathbf{e}_z \\
df_4(0, H_x/2)\,\mathbf{e}_z &= H^{-1}\!\left(-H_z H_x/2\right) = -\half\,\mathbf{e}_y \\
df_4(0, H_x/2)\,\mathbf{e}_t &= H^{-1}\!\left(-H_t H_x/2\right) = -\half\,\mathbf{e}_x
\end{align}

换言之，$df_4(0,H_x/2)=A_x$。依此类推，我们得到：

\begin{align}
df_4(0, H_x/2) &= A_x, \quad df_4(0, H_y/2) = A_y, \quad df_4(0, H_z/2) = A_z \\
df_4(H_x/2, 0) &= B_x, \quad df_4(H_y/2, 0) = B_y, \quad df_4(H_z/2, 0) = B_z
\end{align}

于是 $df_4$ 把 su(2) 的一个副本映入 so(4) 的自对偶部分，把另一个副本映入反自对偶部分。

\subsection{李代数表示（Lie Algebra Representations）}

每当我们有李群在矢量空间 $V$ 上的一个表示 $\rho$，就可以对其求导，得到相应李代数在 $V$ 上的表示 $d\rho$。正如李群表示保持群元素之间的乘法：

\begin{equation}
\rho(gh) = \rho(g)\rho(h)
\end{equation}

李代数表示保持李括号：

\begin{equation}
d\rho([\mathbf{v},\mathbf{w}]) = [d\rho(\mathbf{v}), d\rho(\mathbf{w})]
\end{equation}

我们按通常方式求导：令 $g=\exp(t\mathbf{v})$（$\mathbf{v}$ 为李代数中某元素），对 $t$ 求导，再令 $t=0$。通过这种方法，我们可以方便地从最重要的李群表示得到李代数表示。

我们先从 SU(2) 的前三个自旋-$j$ 表示出发，把它们变成 su(2) 的自旋-$j$ 表示。

\subsubsection{自旋-0}

su(2) 的自旋-0 表示如你所料是平凡的：由于群表示始终给出恒等算符，其导数为零。

\begin{equation}
d\rho_0(\mathbf{v}) = 0
\end{equation}

这里的右边 0 是作用在复数上、把它们乘以零的线性算符。

\subsubsection{自旋-$\tfrac12$}

su(2) 的自旋-$\tfrac12$ 表示作用在 $\mathbb{C}^{2}$，即二维复矢量空间上。

\begin{equation}
d\rho_{\frac12}(\mathbf{v}) = \mathbf{v}
\end{equation}

该表示就是把 su(2) 的原始元素 $\mathbf{v}$ 按通常的矩阵乘法作用在 $\mathbb{C}^{2}$ 上。

\subsubsection{自旋-1}

su(2) 在 $\mathbb{C}^{3}$ 上的自旋-1 表示终于需要一点工作了，但这部分工作我们大多已经做过。

\begin{equation}
d\rho_1(\mathbf{v})\,\mathbf{x} = df_3(\mathbf{v})\,\mathbf{x} = H^{-1}\!\left([\mathbf{v},H(\mathbf{x})]\right)
\end{equation}

严格说来，按 $H(\mathbf{x})$ 的定义，这给出的是 $\mathbb{R}^{3}$ 而非 $\mathbb{C}^{3}$ 上的线性算符；但与自旋-1 群表示一样，我们可以直接取该算符的实 $3\times 3$ 矩阵并作用于 $\mathbb{C}^{3}$。这一表示可以概括为：

\begin{equation}
d\rho_1(H_k/2) = J_k, \quad k=x,y,z
\end{equation}

这里的 $H_k$ 是我们先前列出的 su(2) 基元素，而 $J_k$ 是 $\mathbb{R}^{3}$ 中绕各坐标轴旋转的生成元。

在以上所有情形中，我们都可以通过用同构 $df_3:\mathrm{su}(2)\to\mathrm{so}(3)$ 的逆把 so(3) 的任一元素映回 su(2)（规则为 $J_k\to H_k/2$），再对结果施以 su(2) 的表示，从而得到 so(3) 而非 su(2) 的表示。so(3) 的自旋-0 表示当然就是零，而 so(3) 的自旋-1 表示则直接把 so(3) 的矩阵元素作用于 $\mathbb{C}^{3}$。

\subsubsection{so(4) 的自旋-$(j,k)$ 表示}

我们之前定义了群 $\mathrm{SU}(2)\times\mathrm{SU}(2)$（即 SO(4) 的双覆盖）的自旋-$(j,k)$ 表示。对某两个半整数 $j$ 与 $k$ 的选取，我们使用 $\rho_j$ 与 $\rho_k$（SU(2) 在维数分别为 $2j+1$ 与 $2k+1$ 的空间 $V_j$ 与 $V_k$ 上的自旋-$j$ 与自旋-$k$ 表示），来定义 $\mathrm{SU}(2)\times\mathrm{SU}(2)$ 在张量积 $V_j\otimes V_k$ 上的一个不可约酉表示 $\rho_{j,k}$：

\begin{equation}
\rho_{j,k}(g,h)\;x\otimes y = (\rho_j(g)\,x)\otimes(\rho_k(h)\,y)
\end{equation}

随后我们把它按线性扩展到 $V_j\otimes V_k$ 的任意元素。

这一李群表示的导数是 $\mathrm{su}(2)\oplus\mathrm{su}(2)$ 在 $V_j\otimes V_k$ 上的李代数表示：

\begin{equation}
d\rho_{j,k}(\mathbf{v},\mathbf{w})\;x\otimes y = (d\rho_j(\mathbf{v})\,x)\otimes y + x\otimes(d\rho_k(\mathbf{w})\,y)
\end{equation}

我们给出几个具体例子，使用目前已描述的 su(2) 的自旋-$j$ 表示。对 $j=\tfrac12$ 与 $k=0$，$V_j=\mathbb{C}^{2}$、$V_k=\mathbb{C}$，故 $V_j\otimes V_k$ 就是 $\mathbb{C}^{2}$，我们可以在所有张量积中略去第二个因子：

\begin{align}
d\rho_{\frac12,0}(\mathbf{v},\mathbf{w})\;x\otimes y &= (d\rho_{\frac12}(\mathbf{v})\,x)\otimes y + x\otimes(d\rho_0(\mathbf{w})\,y) \\
d\rho_{\frac12,0}(\mathbf{v},\mathbf{w})\,x &= \mathbf{v}\,x
\end{align}

其中最后一式右边是矩阵乘法。类似地，对 $j=0$ 与 $k=\tfrac12$：

\begin{equation}
d\rho_{0,\frac12}(\mathbf{v},\mathbf{w})\,x = \mathbf{w}\,x
\end{equation}

我们目前所描述的，是 $\mathrm{su}(2)\oplus\mathrm{su}(2)$ 的表示。李代数 so(4) 与 $\mathrm{su}(2)\oplus\mathrm{su}(2)$ 同构，我们可以用函数 $df_4:\mathrm{su}(2)\oplus\mathrm{su}(2)\to\mathrm{so}(4)$

\begin{equation}
df_4(\mathbf{v},\mathbf{w})\,\mathbf{x} = H^{-1}\!\left(\mathbf{v}\,H(\mathbf{x}) - H(\mathbf{x})\,\mathbf{w}\right)
\end{equation}

来认同这两个代数。这相当于认同基元素：

\begin{align}
(0, H_k/2) &\leftrightarrow A_k, \quad k=x,y,z \\
(H_k/2, 0) &\leftrightarrow B_k, \quad k=x,y,z
\end{align}

例如，如果我们想把自旋-$(0,\tfrac12)$ 表示用于生成元 $J_{xy}$，就要经过以下步骤：

\begin{align}
J_{xy} &= A_z + B_z \\
df_4^{-1}(J_{xy}) &= (H_z/2, H_z/2) \\
d\rho_{0,\frac12}(H_z/2, H_z/2)\,x &= \half H_z\,x
\end{align}

so(4) 的自旋-$(0,\tfrac12)$ 表示从 so(4) 投影到由自对偶的 $A$ 矩阵张成的 so(3) 副本，再按通常方式把它映到 su(2)。类似地，so(4) 的自旋-$(\tfrac12,0)$ 表示做同样的事情，但投影到由反自对偶的 $B$ 矩阵张成的子空间。我们把所有可能性列成表，因为在考察黎曼狄拉克方程（Riemannian Dirac equation）时它们会派上用场。

\begin{equation}
\begin{array}{c|c|c}
 & \text{spin-}(\tfrac12,0) & \text{spin-}(0,\tfrac12) \\ \hline
J_{xy} = A_z + B_z \;\to\; & \half H_z & \half H_z \\
J_{yz} = A_x + B_x \;\to\; & \half H_x & \half H_x \\
J_{zx} = A_y + B_y \;\to\; & \half H_y & \half H_y \\
J_{xt} = A_x - B_x \;\to\; & -\half H_x & \half H_x \\
J_{yt} = A_y - B_y \;\to\; & -\half H_y & \half H_y \\
J_{zt} = A_z - B_z \;\to\; & -\half H_z & \half H_z
\end{array}
\end{equation}

\section*{参考文献}

\begin{itemize}
  \item[{[1]}] \emph{Gauge Fields, Knots and Gravity}, John Baez \& Javier P. Muniain, World Scientific, 1994, pp.~179--182.
  \item[{[2]}] \emph{Group Theory and Physics}, S. Sternberg, Cambridge University Press, 1994, Section 1.3.
  \item[{[3]}] Sternberg, \emph{op.~cit.}, Section 3.2.
  \item[{[4]}] Sternberg, \emph{op.~cit.}, Sections 3.3, 3.8, 3.9 and 4.9.
  \item[{[5]}] \emph{Spinors and Trialities}, John Baez.
  \item[{[6]}] \emph{Symplectic, Quaternionic, Fermionic}, John Baez.
\end{itemize}

\chapter{黎曼电磁学}

本章把黎曼宇宙中"光"的概念发展为一套完整的电磁学理论。我们将看到，四维矢量势 $\bA$ 给出电磁场张量 $F$，电荷受力的规律与洛伦兹宇宙形式上完全平行，但偏振模式、多普勒频移、静电力的空间振荡乃至一切静电与磁现象都表现出本质性的差异。理解这些差异，是理解黎曼宇宙中物质如何相互作用的关键。

\section{三种偏振}

在我们所在的宇宙中，真空中的光是完全\textbf{横波}（transverse wave）：构成光波的振荡电场与磁场始终与该波在空间中传播的方向正交。

\orthofig{images/04__001.png}{横波与纵波}

在黎曼宇宙中，如矢量波一节的讨论，等价的光波在某种意义上也总是"横"的：描述振荡场的四维矢量 $\bA$ 与四维传播矢量 $\bk$ 正交。若波写作
\[
A(x) = A_0 \sin(\bk \cdot x),
\]
则必须有 $\bk \cdot \bA_0 = 0$。施加该条件后，波矢矢量仍允许三个独立方向。例如，设传播矢量 $\bk$ 位于 $xt$ 平面内（如右图）。那么当 $\bA_0$ 沿 $y$ 轴、沿 $z$ 轴、或位于 $xt$ 平面内且与 $\bk$ 成直角时，这三个方向互相垂直，且都与 $\bk$ 正交。（第二种可能性在图中未画出，因为为了显示 $t$ 方向我们省略了 $z$ 方向。）

然而把这些矢量投影到三维空间后，前两个方向\textbf{仍}保持与波传播方向正交，第三个方向却变得与它平行（或反平行）。在这个意义上，黎曼宇宙中"光"的等价物可以包含一个\textbf{纵模}（longitudinal mode）：场沿波传播的方向振荡。（图中我们展示的是向 $xy$ 平面的投影。由于所用矢量本就没有 $z$ 分量，其结果与投影到三维空间相同。）

必须强调，这三种\textbf{偏振}（polarisation）模式——两种看起来横、一种看起来纵——当我们把光本身作为四维空间中的平面波来看待时，是完全相同的。在每种情形下，$\bk$ 与 $\bA_0$ 都互相正交。只有当我们选定某个特定方向作为时间轴，并\emph{相对于该选择}去区分四维空间中的方向时，谈论"纵偏振"才有意义。尽管如此，这里与我们的宇宙仍存在非常真实的差别——我们的宇宙只有两种偏振。在黎曼宇宙中，如果偏振片的等价物移除了光的一个横向分量，那么"正交偏振片"（在我们的宇宙中会完全阻断光）仍会透射随机偏振光源中的纵分量。

\section{多普勒频移（Doppler Shift）}

\orthofig{images/04__024.png}{洛伦兹宇宙中的多普勒频移}

多普勒频移（Doppler shift）是一个众所周知的效应：观察者的运动改变了他们测得的光波波长与频率。

右图显示同一列波前扫过洛伦兹宇宙中三位不同的观察者，光从右侧射来。"参考观察者"假定相对光源静止，另外两位观察者则分别朝向、背离光源运动。

参考观察者测得相邻波前之间的周期为 OR。朝光运动的观察者在事件 O 和 F 处遇到波前，但图中所示的 OF 的欧几里得长度并非观察者经历的时间，后者需用洛伦兹版勾股定理计算。为便于比较，我们用了双曲线——洛伦兹宇宙中一条等固有时（proper time）曲线，正如圆弧是欧几里得空间中的等距曲线——把该间隔移至参考观察者的世界线上。在那里可以明确看到 $\mathrm{OF}'$ 小于 OR，故向前运动的观察者测得光的周期变短、频率升高。这称为\textbf{蓝移}（blue shift），因为蓝光处于可见光谱的高频端。

类似地，向后运动的观察者在事件 O 与 B 处遇到波前；为比较我们把该间隔移至 $\mathrm{OB}'$，此处它显然长于 OR。因此向后运动的观察者测得光的周期变长、频率降低。这称为\textbf{红移}（red shift）。

设 F 的坐标为 $(x_F, t_F)$。若观察者速度为 $v$，则 $x_F = v\,t_F$。记 $\tau_R$ 为参考观察者测得的周期，即 OR。由于波前从 F 到 R 的速度为 1（采用我们的单位），有
\[
\tau_R = t_F + x_F = t_F\,(1+v),
\]
\[
t_F = \frac{\tau_R}{1+v}.
\]
于是朝前运动的观察者测得的光周期由洛伦兹版勾股定理给出：
\[
\tau_F = \sqrt{t_F^{2} - x_F^{2}}
       = t_F\sqrt{1 - v^{2}}
       = \tau_R\frac{\sqrt{1 - v^{2}}}{1+v}
       = \tau_R\sqrt{\frac{1 - v}{1+v}}.
\]
洛伦兹蓝移的\emph{频率}之比恰为周期之比的倒数：
\[
\frac{\nu_F}{\nu_R} = \frac{\tau_R}{\tau_F} = \sqrt{\frac{1+v}{1-v}}.
\]
洛伦兹蓝移因子不依赖于光的具体频率。设速度 $v = 0.6$，即光速的百分之六十，则蓝移因子为 2，来自正前方的一切光相对静止时的频率都翻倍；同理，来自正后方的光频率减半。

黎曼情形稍复杂，但让我们仍从同样的情景开始：一位观察者接近光源，另一位背离光源——两者速度都低于光本身的速度。

\orthofig{images/04__025.png}{黎曼宇宙中的多普勒频移}

此情形下从图中可清晰看到，朝前运动的观察者测得的光周期 $\mathrm{OF}'$ \emph{大于} OR，而向后运动的观察者测得的周期 $\mathrm{OB}'$ \emph{小于} OR。

这是否意味着接近光的观察者看到"红移"？这里我们必须停下来追问：在黎曼宇宙中用何种语言描述可见光谱？这并非关于异宇宙生物能否体验相同的"红""蓝"感受的哲学疑虑——只要我们保持一致，尽可把他们的色觉随意翻译成我们的词汇。但无论黎曼宇宙居民对光的感觉如何，我们都不可能做出尊重光的一切客观物理属性的完美翻译。在我们的宇宙中，红光比蓝光既有更长的波长、也有更长的周期。在黎曼宇宙中，\emph{没有任何一种光}同时具有比另一种光更长的波长\emph{和}更长的周期。关系总是相反的：波长越长的光，周期反而越短。

因此我们必须作出选择。小说《Orthogonal》中采取的任意方案是，把可见光中最短的\emph{波长}译为"紫"，于是我们把"蓝移"一词定义为"朝向更短的波长"。

确定了术语之后，我们继续来量化最容易的东西，而这仍然是周期或频率之比而非波长之比。若该光相对参考观察者的速度为 $c$，则波前沿与传播矢量相反的方向以速度 $1/c$ 运动，因为与斜率为 $c$ 的直线垂直的直线的斜率是 $-1/c$。黎曼版的计算成为
\[
\tau_R = t_F - \frac{x_F}{1/c} = t_F\,(1 - vc),
\]
\[
t_F = \frac{\tau_R}{1 - vc},
\]
\[
\tau_F = \sqrt{t_F^{2} + x_F^{2}}
       = t_F\sqrt{1 + v^{2}}
       = \tau_R\frac{\sqrt{1 + v^{2}}}{1 - vc}.
\]

\orthofig{images/04__026.png}{多普勒频移后的频率}
\orthofig{images/04__028.png}{多普勒频移后的空间频率}

于是黎曼"蓝移"对接近光源的观察者而言，反而涉及观测频率的\emph{降低}：
\[
\frac{\nu_F}{\nu_R} = \frac{\tau_R}{\tau_F} = \frac{1 - vc}{\sqrt{1 + v^{2}}}.
\]
若观察者速度 $v$ 达到 $1/c$，则其测得的光频率降为零。图中世界线为 OG 的观察者就是例子：该观察者看到光速为无穷大，波长为最小可能值。

若观察者朝前运动的速度大于 $1/c$（例如世界线为 OH 的观察者），来自光源的传播矢量在他看来指向过去，因此在与光的任何相互作用中，他都会认为\emph{自己}是光源，传播矢量方向反转。

把同一公式中的 $v$ 取负值即可得到背离光源的观察者的频率比。从图中可见，当 $v = -c$ 时会发生特别的事：世界线为 OC 的观察者与光同步，使光的相对速度为零。对这样的观察者，频率达到最大值 $\nu_{\max}$——即可能的最大频率，因为他正径直穿过波前。

当观察者开始超过光速、比光更快地向后运动时又如何？我们的频率比公式给出一个正值，此时它小于最大值，故光应当仍可见。但当你以比光本身更快的速度逃离光源时，如何还能接收到来自光源的光？我们的假定是：光在一段延展的时间内被持续产生（比如星光就是如此），因此这不是观察者与某一单个光脉冲之间的赛跑。若观察者运动快于所讨论的光（如世界线为 OD 的观察者），他会开始追上"旧"光——那些已有足够领先量的光正等着他。并且他会看到这些光来自他行进的方向，而非光源本身的方向。

除了多普勒频移改变光的频率，几何关系还意味着运动的观察者测得的光线之间的\emph{夹角}与参考观察者不同，这一效应称为\textbf{光行差}（aberration）。在黎曼情形下，我们可用普通欧几里得几何理解多普勒频移与光行差。

设参考观察者看到天空中均匀分布的大量光源发出某一特定颜色的光。所有入射光线在四维空间中构成一个锥面，其轴为参考观察者的世界线。若把所有传播矢量投影到参考观察者所认为的空间中，它们会对称地展开成一个圆（下图中灰色所示）。

位于同一位置、以某速度 $v$ 向右运动的观察者，看到的光来向对应于同一锥面上这些光线在他们自己空间中的投影——该空间相对参考观察者的空间是倾斜的。这会产生两个效应：改变传播矢量投影的长度，并改变这些投影之间的角度。传播矢量在观察者空间上投影的\emph{长度}正比于光的空间频率（如最初考察波时所讨论的），因此更长的投影意味着更短的波长。

\orthofig[0.72\textwidth]{images/04__027.png}{黎曼宇宙中的多普勒频移与光行差}

该图展示了观察者取各种速度时的结果，假定所有入射光在参考观察者参考系中的速度为 0.75。图中颜色与小说用于可见色调的翻译方案一致，虚线表示紫外与红外。

第一个例子取 $v = 0.25$：来自观察者行进方向的光发生蓝移，光源在天空中显得被压缩挤在一起；来自相反方向的光发生红移，在天空中散开。用这种说法——把红移、蓝移理解为波长而非频率——这段描述听起来与洛伦兹版本一模一样。

一旦 $v = c$，一切光看起来都来自前方。当 $v = 1/c = 4/3$（未画出）时，正前方的光线将逆着时间到达，因而不可见；当 $v = 2$ 时，锥面的可觉察部分因此被略去。

对 $v = 2$，我们看到一件在洛伦兹物理中绝不可能发生的奇事：来自同一表观方向、具有两种不同颜色、实际上来自两个不同光源的光。在红、黄、绿光线与紫外光混杂之处，较红的光实际上来自后方（按参考观察者而言），但运动观察者的倾斜使其看起来来自前方，并与真正来自前方的光叠加在一起。

当 $v = \infty$ 时，参考观察者认为来自前方的一切光如今都来自运动观察者的未来，因而不可见；此时可见的光全部来自后方，并被挤入一片天空区域，其角宽度恰好等于入射光锥的角宽度。

请记住，我们这里描述的一切都只涉及在参考观察者看来\emph{只有一种颜色}的光。即便对一位相对恒星平均运动静止的观察者而言，在黎曼宇宙中大多数恒星也会呈现为\emph{色彩拖尾}（colour trails），因为恒星在发出目前到达观察者的红光之时与稍后发出同样正在到达的、更快的紫光之时之间，已经移动了一点。多普勒频移与光行差如何改变这些恒星拖尾，在小说第一卷中有所讨论；书中预告片可见这些效应的图示，另有一个相关主题的视频教程。

\begin{notebox}[术语约定]
黎曼宇宙中波长与周期的关系与我们的宇宙相反，无法同时把"更长的波长"和"更长的周期"都译为"更红"。小说采取任意方案：把最短波长译为"紫"，并约定\textbf{蓝移＝朝向更短的波长}。
\end{notebox}

\section{黎曼电磁场}

至此我们一直在谈论一个"振荡场" $\bA$，却未确切说明它是什么。事实证明，可以完全类比我们自己宇宙的洛伦兹电磁学，发展出一整套黎曼电磁学理论，其中包括我们可以合理地称为电场与磁场的矢量场，以及一种类似电荷的物质属性——它使相关物质受到电磁力，并具备产生电磁波的能力。

在这一框架中，我们一直称为 $\bA$ 的场就是所谓的\textbf{四维矢量势}（four-vector potential）。正如电势是这样一种量：它在空间中的变化率给出电场；四维矢量势则是一个矢量，其导数给出完整的电磁场。具体地说，电磁场 $F$ 的分量为
\[
F_{ij} = \partial_i A_j - \partial_j A_i,
\]
其中 $i$、$j$ 取遍 $x$、$y$、$z$ 与 $t$。换言之，$F$ 由一个 $4\times4$ 矩阵描述，它由四维矢量势 $\bA$ 导出。

这个矩阵的意义是什么？设有带电量为 $q$ 的粒子，沿一条世界线运动，其四维切线是单位矢量 $\bu$，即粒子的\emph{四维速度}（four-velocity）。若把 $\bu$ 乘以电荷 $q$ 与矩阵 $F$，就得到带电粒子所受的力 $\mathbf{f}$：
\[
\mathbf{f} = q F \bu.
\]
当然，在牛顿力学中力是三维矢量，但在相对论物理中（无论黎曼还是洛伦兹），把力视为四维矢量更为合理。正如牛顿的著名方程
\[
\mathbf{f} = m\,\mathbf{a}
\]
（力等于质量乘加速度），力的四维矢量 $\mathbf{f}$ 与加速度的四维矢量 $\mathbf{a}$ 相联系……后者又与四维速度通过
\[
\mathbf{a} = \partial_{\tau}\bu
\]
相联系。四维加速度是四维速度相对\emph{固有时}（proper time）$\tau$ 的变化率，$\tau$ 是沿粒子世界线测得的时间。把所有这些片段放在一起，我们得到关于电磁场 $F$ 对带电量 $q$、质量 $m$ 的粒子运动之影响的如下描述：
\[
\partial_{\tau}\bu = \frac{q}{m} F \bu.
\]
四维速度 $\bu$ 的模长恒为 1。任何与 $\bu$ 自身平行的改变都会改变其模长，因此描述其变化率的矢量必须始终与 $\bu$ 正交。这意味着 $F$ 必须属于一类特殊的矩阵：用 $F$ 乘任何矢量，所得结果都与原矢量正交。不难看出，这对任何\emph{反对称矩阵}（antisymmetric matrix）都成立——即等于自身转置之负的矩阵，或 $F_{ij} = -F_{ji}$。于是，使用爱因斯坦求和约定，有
\[
\bu \cdot (F\bu) = u^{i} F_{ij} u^{j} = u^{i} (-F_{ji}) u^{j} = -\bu \cdot (F\bu),
\]
这表明点积必须为零。而用 $\bA$ 定义 $F$ 的方式，保证了它\emph{确实}是反对称的。

那么，带电粒子在\emph{恒定}电磁场中运动时将走怎样的世界线？最容易考虑的情形是矩阵 $F$ 取特别简单的形式：
\[
F_{ij} = b_i c_j - b_j c_i,
\]
其中 $\mathbf{b}$ 与 $\mathbf{c}$ 是任意两个\emph{不平行}的矢量。（若 $\mathbf{b}$ 是 $\mathbf{c}$ 的倍数，则右端变为零。）我们把此式缩写为
\[
F = \mathbf{b} \wedge \mathbf{c}.
\]
符号"$\wedge$"称为"楔积"（wedge product）。此时可把 $F$ 与矢量 $\bu$ 的矩阵乘积用 $\mathbf{b}$ 与 $\mathbf{c}$ 重写：
\[
F\bu = (\mathbf{c}\cdot\bu)\,\mathbf{b} - (\mathbf{b}\cdot\bu)\,\mathbf{c}.
\]
显然这是始终位于 $\mathbf{b}$、$\mathbf{c}$ 张成平面内的矢量；它也只依赖于 $\bu$ 落在该平面内的分量：若 $\bu$ 完全与该平面正交（即 $\mathbf{b}\cdot\bu = \mathbf{c}\cdot\bu = 0$），则 $F\bu$ 为零。

\orthofig{images/04__002.png}{恒定电磁场中的粒子}

结果一般形如右图。当 $\bu$ 既不完全落在 $\mathbf{b}$、$\mathbf{c}$ 张成的平面内，又不完全与之正交时，粒子会以恒定速率保持在与该平面正交的方向上运动，而它在平行于该平面的方向上的运动则被弯成一个圆。于是粒子世界线整体是一条螺旋线。

若该平面由你视为空间的两个方向张成（换言之，平面与你的世界线正交），那么运动的圆周部分在你看来就是空间中的运动：粒子要么在一个固定平面内打转，要么在空间中沿螺旋前进。这正是我们的宇宙中与\textbf{磁场}相联系的行为。事实上，在一种略为老式的粒子加速器——回旋加速器（cyclotron）中，正是用恒定磁场使粒子如此打转。磁场的描述约定是：赋予它一个与粒子打转平面正交的空间方向。

若该平面包含你视为时间的方向，则世界线的弯曲将表现为粒子在平面内另一（空间）方向上加速。在我们宇宙中，带电粒子被如此加速时，我们认为这是\textbf{电场}的作用。在洛伦兹物理中，加速粒子的世界线是双曲线而非圆，但除此之外效应相同。我们赋予电场一个落在该平面内的空间方向。

于是形如 $F = \mathbf{b}\wedge\mathbf{c}$ 的电磁场被视为：
\begin{itemize}
\item 若 $\mathbf{b}$ 与 $\mathbf{c}$ 都是空间方向，即一个同时垂直于 $\mathbf{b}$、$\mathbf{c}$ 的磁场；
\item 若该平面包含你的时间方向，即一个空间方向落在 $\mathbf{b}$、$\mathbf{c}$ 张成平面内的电场。
\end{itemize}
若该平面不干净地落入这两类中的任何一类，$F$ 将被视为同时包含电场与磁场。

让我们用这些描述来看黎曼光的横模与纵模中发生了什么。一般地，若
\[
A(x) = A_0 \sin(\bk\cdot x),
\]
则电磁场由矩阵 $F$ 描述：
\[
F(x) = (\bk \wedge A_0)\cos(\bk\cdot x).
\]
设光传播矢量位于 $xt$ 平面内，在 $x$ 方向速度为 1，即 $\bk = \mathbf{e}_x + \mathbf{e}_t$。（其实它必须是满足条件 $|\bk| = \wms$ 的该矢量的某一倍数，但 $\wms$ 的取值在下面并不重要。）

首先考察一个我们视为横模的波，取 $\bA_0 = \mathbf{e}_y$。$\bk$ 与 $A_0$ 张成的平面既不完全空间化，也不包含我们的时间轴 $\mathbf{e}_t$，故该波将\emph{同时}包含电场与磁场。但我们可以把 $F$ 拆成在我们的坐标系中分别对应磁场与电场的两部分：
\[
\begin{aligned}
F(x) &= \big((\mathbf{e}_x + \mathbf{e}_t)\wedge \mathbf{e}_y\big)\cos(\bk\cdot x) \\
     &= (\mathbf{e}_x \wedge \mathbf{e}_y)\cos(\bk\cdot x) + (\mathbf{e}_t \wedge \mathbf{e}_y)\cos(\bk\cdot x).
\end{aligned}
\]
第一部分，平面由 $\mathbf{e}_x$、$\mathbf{e}_y$ 张成，是一个指向 $z$ 方向的振荡磁场。第二部分，平面由 $\mathbf{e}_t$、$\mathbf{e}_y$ 张成，是一个指向 $y$ 方向的振荡电场。这与我们自己宇宙中的光非常相似：电场、磁场与光的传播方向三者两两垂直。

其次考察一个我们视为纵模的波，取 $\bA_0 = \mathbf{e}_x - \mathbf{e}_t$。
\[
\begin{aligned}
F(x) &= \big((\mathbf{e}_x + \mathbf{e}_t)\wedge(\mathbf{e}_x - \mathbf{e}_t)\big)\cos(\bk\cdot x) \\
     &= 2\,\mathbf{e}_t \wedge \mathbf{e}_x \cos(\bk\cdot x).
\end{aligned}
\]
在化简楔积时，我们利用了任何矢量与自身的楔积为零、以及 $\mathbf{b}\wedge\mathbf{c} = -\mathbf{c}\wedge\mathbf{b}$ 这两点。结果是沿 $x$ 方向——即平行于光本身运动的方向——的一个振荡的\emph{纯电场}。

\begin{keybox}[黎曼电磁场的要点]
四维矢量势 $\bA$ 通过 $F_{ij} = \partial_i A_j - \partial_j A_i$ 给出反对称的电磁场张量 $F$；其反对称性保证四维速度模长恒定。带电荷 $q$、质量 $m$ 的粒子的运动满足 $\partial_\tau \bu = (q/m)F\bu$。当 $F$ 的平面由两个空间方向张成时是磁场，当包含时间方向时是电场，否则两者兼有。
\end{keybox}

\section{波纹状的库仑力（The Corrugated Coulomb Force）}

黎曼矢量波方程（RVW）支配真空中电磁场的四维矢量势 $\bA$——即不存在任何可能产生该场的带电物质时的情形。但我们可以轻易地把 RVW 方程推广，加入一个描述电荷存在的"源"项，如下所示：
\begin{equation}
\partial_x^{2}\bA + \partial_y^{2}\bA + \partial_z^{2}\bA + \partial_t^{2}\bA + \wms^{2} \bA + \bj = \mathbf{0}. \tag{RVWS}
\end{equation}

我们在此方程中加入的四维矢量 $\bj$ 有非常简单的含义。设有某种带电物质存在于所关注的区域，一位相对该物质静止的观察者测得其电荷密度为 $\rho$。若该物质的四维速度为 $\bu$，我们定义
\[
\bj = \rho\,\bu.
\]
现在，$\bu$ 是切于物体世界线的单位矢量。但在世界线上任一点都有\emph{两个}这样的矢量，指向相反的方向。我们如何选取正确的那一个？若某观察者以其中一种选择作为其时间轴，带电物质会显得带正电；若取另一方向，同一物质会显得带负电。改变选择会把 $\rho$ 与 $\bu$ 同乘以 $-1$，从而抵消这一变化。因此选哪一个都无所谓。

我们把矢量 $\bj$ 称为\textbf{四维电流}（four-current）。若一位并不相对带电物质静止的观察者测量 $\bj$ 的分量，他将同时看到一个电荷密度 $j^t$，以及一个描述电荷运动的\textbf{电流密度}（current density）$j^x$、$j^y$、$j^z$。

该方程最简单也最重要的解之一称为\textbf{库仑势}（Coulomb potential）。这是带电量为 $q$ 的粒子占据单个点、并在某坐标系中静止时的势。当然这是一种理想化，它有沿粒子世界线出现无穷大电荷密度这样一个略显麻烦的数学特征。尽管如此，求解并不太难。如何做到这一点的细节见该主题的补充注记；这里我们把它作为已知结果接受。若 $r$ 为到粒子的距离，则四维矢量势为
\[
\bA(r) = \frac{q}{4\pi r}\cos(\wms r)\,\mathbf{e}_t.
\]
当 $\bA$ 只有这样的一个 $t$ 分量、且 $\bA$ 自身不随时间变化时，更简单的做法是忘掉四维矢量势，只用电势 $\phi$ 来工作，它是 $A_t$ 的相反数。

\emph{库仑势：}
\[
\phi(r) = \begin{cases}
-\dfrac{q}{4\pi r}\cos(\wms r), & \text{（Riemannian）}\\[8pt]
\dfrac{q}{4\pi r}, & \text{（Lorentzian）}
\end{cases}
\]
电势是带一单位电荷的正"试探粒子"的势能——正如任何随处变化的势能一样，它产生一个力。力在任一方向的分量等于 $\phi$ 在该方向变化率的相反数（粒子将朝势能\emph{增加}的相反方向加速，正如石头会沿山\emph{滑下}）：
\[
\mathbf{f} = -\big(\partial_x\phi\,\mathbf{e}_x + \partial_y\phi\,\mathbf{e}_y + \partial_z\phi\,\mathbf{e}_z\big),
\]
用梯度符号"$\nabla$"缩写为
\[
\mathbf{f} = -\nabla\phi.
\]
我们声称的 $\phi$ 与 $A_t$ 之间的简单关系，可通过检验二者对试探粒子给出完全相同的力来验证。

\orthofig{images/04__003.png}{正电荷周围黎曼与洛伦兹电势的比较}
\orthofig{images/04__004.png}{负电荷周围黎曼与洛伦兹电势的比较}

$\phi$ 是什么样子？右上图显示在黎曼与洛伦兹电磁学中正电荷周围电势的情形，下图显示负电荷周围的电势。

黎曼情形与洛伦兹情形之间有两个显著差异。在我们自己的宇宙中，众所周知\emph{异种电荷相吸}，\emph{同种电荷相斥}。但在黎曼宇宙中，正电荷附近的电势是一个井而非一座山，因此若两个正电荷足够靠近，它们会相吸；异种电荷若足够靠近则会相斥。

第二个差异是，我们必须指明两电荷之间的距离，才能说它们相吸还是相斥。在洛伦兹宇宙中，静电力只是随距离逐渐减弱；在黎曼宇宙中，力的方向还随距离\emph{交替}变化。从 $\phi$ 解的形式可以看出，这种交替的空间频率等于光的最大可能频率 $\nu_{\max}$（因为 $\wms = 2\pi\nu_{\max}$）。等价地，其起伏的波长等于光的最小可能波长 $\lambda_{\min} = 1/\nu_{\max}$。

不难验证我们对 $\bA$ 的解满足 RVWS 方程；这种交替电势其实不大令人惊讶，因为我们已经看到，紫外极限下的平面波也涉及一个不随时间变化、而以频率 $\nu_{\max}$ 在空间中振荡的电磁场。该方程要求所有方向上频率的平方和为常数，因此任何不随时间变化的静场\emph{必然}在空间中振荡。

但我们怎能确信方程本身正确？若把方程中 $\bj$ 的符号反号，我们就会得到原解的相反情形，从而得到近距相斥的同种电荷。是什么阻止我们提出黎曼宇宙中电磁学定律取\emph{那种}形式呢？

最简单的回答是回想我们在能量与动量注记中讨论过的一件事：在黎曼宇宙中相对论总能量的意义与势能和动能相反。在洛伦兹宇宙中，电磁场拥有能量，正比于场强的平方。给定一个正电荷与一个负电荷，当二者相互靠近时该能量会减少，因为它们会越来越多地\emph{抵消}彼此的场。但在黎曼宇宙中，使势能最小意味着使场的能量\emph{最大}。对于足够靠近的同种电荷，把它们拉得更近会使它们形状相同的场叠加得越来越匹配，\emph{减少}两粒子场之间的抵消，从而\emph{增加}整体场能。

\section{不受欢迎的电子学（Unpopular Electronics）}

在我们的宇宙中，环绕带电粒子（如电子）的微小电场能够轻易累积到产生可见效应的程度。把塑料笔在合适的衣物上摩擦十几次，它就会带上多余电荷，产生强到足以克服重力吸起一小片纸屑的电场。（纸张本身并无净电荷，但笔的电场使纸表面的电荷分离。）多余电荷在笔表面如何精确分布无关紧要；只要\emph{存在}多余电荷，其每一部分都会协同增强笔周围的电场强度——而这个电场在时间上持久、在空间上一致，在很长的区域内大致指向同一方向，并保持到电荷缓慢泄漏为止。

\orthofig{images/04__005.png}{随机排列电荷产生的电势}
\orthofig{images/04__006.png}{晶格排列电荷产生的电势}

在黎曼宇宙中，由于库仑势在空间中的振荡，多余电荷的效应非常不同。大量相同粒子杂乱排列（如右图）时，对附近某点电势的贡献既有正也有负，彼此大量抵消。当然抵消不会完美，但残余的场将是一系列密布的峰谷，而非洛伦兹情形下那种单一、持久的斜坡。

有序排列的电荷阵列（如右下图）在某些点可能比随机排列抵消得更少，但同样会产生起伏的电势地形。正试探粒子若从正确的初始位置静止出发，确实可能被这堆电荷排斥，但该效应不像洛伦兹物理中的等价物那样容易复现。更有甚者，任何置于该场中、由许多粒子（而非单一粒子）组成的物体，将经历在诸多略有不同位置上取决于场的各个力的总和，从而导致进一步的抵消。

其结果是：静电现象将对电荷的细节、微观排列十分敏感，且总体而言远不如在我们宇宙中那样容易探测和利用。

类似的效应也适用于磁性。在我们的世界里，仅需把大量小磁偶极子的方向对齐，它们的场便足以互相增强；永磁体正是这样工作的。但在黎曼宇宙中，即便偶极子本身全部对齐，当你穿越空间时，每个偶极子产生的场也会在两个相反方向之间来回切换，且在任意一点，来自距离略有不同的各个偶极子的贡献之间都会发生抵消。

不过，尽管这些问题排除了那种让我们最初窥见电磁学本质的、容易的低技术实验，它们并非完全不可逾越。我们描述过的静场（不随时间变化）其波长极小，正是这一点使一切如此精细、如此容易抵消；而一旦引入\emph{时间}上的振荡，这种小波长会逐渐让位于更大的波长。在直流电路中需要极高精度才能工作的电子元件，若改用频率足够高的\textbf{交流电}（alternating current），对不完美的容忍度将大得多。

什么叫"频率足够高"？只有当频率 $\omega$ 达到 $0.995\,\wms$ 时，场的波长才会增长到最小波长的十倍，因此要获得任何可用的空间不精确容忍度，代价就是必须处理黎曼宇宙中能传播的、本质上频率最高的电磁辐射。这将带来它自己的技术挑战——但它们与在我们宇宙中处理高频的困难并不完全相同，后者有些实际上与波长有关。

对不介意稍高数学水平的读者，这个主题另有补充材料。

\section{黎曼 Proca 方程（The Riemannian Proca Equation）}

如果光子有质量，我们自身宇宙中的电磁学会有何不同？1930 年代，罗马尼亚物理学家 Alexandru Proca 将麦克斯韦方程组（Maxwell's equations）推广，建立了一套大质量粒子产生类电磁力的理论，用以解释弱核力，这是开创性的工作。Proca 的名声似乎不及其应得，但他的结果曾被 Wolfgang Pauli 在 1946 年的诺贝尔获奖演讲中提到。由与弱力的联系不难猜到：给传递力的粒子赋予\emph{静止质量}会缩短其作用程。若在我们自身宇宙中光子有质量，库仑势将随距离指数衰减。

但正如我们已经看到的，黎曼库仑势并不会指数衰减；相反，它会随空间振荡。从洛伦兹几何变为黎曼几何造成了一切差别。

为得到 Proca 方程的黎曼版本，我们从带源项 $\bj$（我们称之为\emph{四维电流}，four-current）的黎曼矢量波方程出发，再加上我们对任何矢量波 $\bA$ 施加的横波条件——它用来排除那些伪装成矢量波的标量波解。

\begin{equation}
\partial_{x}^{2}\bA + \partial_{y}^{2}\bA + \partial_{z}^{2}\bA + \partial_{t}^{2}\bA + \wms^{2}\bA + \bj = \mathbf{0} \tag{RVWS}
\end{equation}
\begin{equation}
\partial_{x} A^{x} + \partial_{y} A^{y} + \partial_{z} A^{z} + \partial_{t} A^{t} = 0 \tag{Transverse}
\end{equation}

从这对方程立刻可以得到一个好结果：
\begin{equation}
\partial_{x} j^{x} + \partial_{y} j^{y} + \partial_{z} j^{z} + \partial_{t} j^{t} = 0 \tag{1}
\end{equation}
它由横波条件、以及四维电流 $\bj$ 等于 $\bA$ 与其导数的线性组合这一事实推出。这相当于一条\emph{电荷守恒}（conservation of charge）的陈述：某点电荷密度随时间增大的速率 $\partial_t j^t$，是电流密度散度 $\partial_x j^x + \partial_y j^y + \partial_z j^z$ 的相反数，后者描述单位时间内流出该点周围小区域的电荷净量。

我们此前指出，能量-动量四维矢量在不同观测者计算下会有严重问题，因为不存在客观的手段来决定它应沿物体世界线指向何方。四维电流没有这个毛病，因为它被定义为 $\bj = \rho\bu$，其中 $\rho$ 是带电物质静止系中的电荷密度；若我们把 $\bu$ 反号，$\rho$ 也随之反号，因为一个时间反演的正电荷看起来是负的，反之亦然。（当然，把"正"与"负"这两个标签分配给电荷只是约定，但这是一个可以全局地、一劳永逸地作出的选择。）

正如寻常电磁学，我们用 $\bA$ 定义电磁场 $F$：
\begin{equation}
F_{ab} = \partial_{a} A_{b} - \partial_{b} A_{a} \tag{2}
\end{equation}
这里的 $A_a$ 是矢量 $\bA$ 对应的对偶矢量（dual vector）的分量。留意这一区分是个好习惯，尽管在黎曼空间的幺正直角坐标中，矢量分量（如 $A^{a}$）与相应对偶矢量分量（如 $A_{a}$）完全相同。在洛伦兹时空中，这几乎成立，但并非完全成立：我们有 $A_{x} = A^{x}$、$A_{y} = A^{y}$、$A_{z} = A^{z}$，但 $A_{t} = -A^{t}$。

假定我们选三个坐标，记作 $a$、$b$、$c$。那么仅凭 $F$ 的定义、以及导数可交换（即 $\partial_{a}\partial_{b} = \partial_{b}\partial_{a}$）这一事实，就有：
\begin{align}
\partial_{a} F_{bc} + \partial_{b} F_{ca} + \partial_{c} F_{ab}
&= \partial_{a}(\partial_{b} A_{c} - \partial_{c} A_{b})
 + \partial_{b}(\partial_{c} A_{a} - \partial_{a} A_{c})
 + \partial_{c}(\partial_{a} A_{b} - \partial_{b} A_{a}) \nonumber\\
&= 0. \tag{3}
\end{align}

由 $F$ 的定义还可推出：
\begin{align}
\partial_{b} F^{ab}
&= \partial_{b}(\partial^{a} A^{b} - \partial^{b} A^{a}) \nonumber\\
&= \partial^{a}(\partial_{b} A^{b}) - \partial_{b}\partial^{b} A^{a} \nonumber\\
&= -\partial_{b}\partial^{b} A^{a}. \tag{4}
\end{align}
这里使用了爱因斯坦求和约定，且 $\partial_{b} A^{b}$ 由横波条件为零。

把结果 (4) 代入带源项的黎曼矢量波方程 (RVWS)，即得\emph{黎曼 Proca 方程}（Riemannian Proca equation）。我们还有方程 (3)，它仅由 $F$ 的定义推出，因而为黎曼版与洛伦兹版的电磁学\emph{共有}。下面给出四维形式的麦克斯韦方程组以作比较。在此处及此后的一切中，我们为洛伦兹方程选取光速为 1、真空介电常数 $\varepsilon_{0}$ 为 1 的单位制，并采用 $(-+++)$ 符号的洛伦兹度规，而某些文献用 $(+---)$ 符号。

\begin{keybox}[黎曼 Proca 方程 与 麦克斯韦方程组]
\textbf{[黎曼 Proca 方程]（Riemannian Proca Equation）}
\begin{align}
\partial_{b} F^{ab} - \wms^{2} A^{a} - j^{a} &= 0 \tag{Riemannian}\\
\partial_{a} F_{bc} + \partial_{b} F_{ca} + \partial_{c} F_{ab} &= 0 \tag{Common}
\end{align}

\textbf{[麦克斯韦方程组]（Maxwell's Equations）}
\begin{align}
\partial_{b} F^{ab} - j^{a} &= 0 \tag{Lorentzian}\\
\partial_{a} F_{bc} + \partial_{b} F_{ca} + \partial_{c} F_{ab} &= 0 \tag{Common}
\end{align}
\end{keybox}

\section{高斯定律与安培定律的变迁（What Becomes of Gauss and Ampère）}

黎曼电磁学的四维方程很简洁，但要在各种情形中看清正在发生什么、并与洛伦兹对应物作比较，给出三维版本会有所帮助；在三维版本中，我们不谈电磁场 $F$，而是用两个三维矢量场——电场 $\bE$ 与磁场 $\bB$——来描述一切。

\subsection{共同的基础（Common Ground）}

我们从一些定义开始。电场 $\bE$ 的分量取为电磁场 $F$ 中第一个指标为空间指标、第二个指标为 $t$ 的那些分量；而磁场 $\bB$ 的每个分量是 $F$ 中指标为\emph{另外两个}空间方向、并保持循环顺序 $xyz$ 的分量。

\begin{keybox}[电场、磁场与电磁场的定义]
\textbf{[电场 $\bE$]}
\begin{equation}
(E_{x}, E_{y}, E_{z}) = (F_{xt}, F_{yt}, F_{zt}) \tag{Common}
\end{equation}

\textbf{[磁场 $\bB$]}
\begin{equation}
(B_{x}, B_{y}, B_{z}) = (F_{yz}, F_{zx}, F_{xy}) \tag{Common}
\end{equation}

\textbf{[电磁场 $F$]}（此处采用 $(-+++)$ 符号的洛伦兹度规）
\begin{equation}
F_{ab} = \begin{pmatrix}
0 & -E_{x} & -E_{y} & -E_{z}\\
E_{x} & 0 & B_{z} & -B_{y}\\
E_{y} & -B_{z} & 0 & B_{x}\\
E_{z} & B_{y} & -B_{x} & 0
\end{pmatrix} \tag{Common}
\end{equation}
\end{keybox}

注意在上面的矩阵中，$F$ 的第一个指标指行、第二个指标指列，而 $t$ 分量显示在\emph{第一}行与\emph{第一}列。于是例如 $F_{xt}$ 是第二行第一列的元素。

我们把\emph{电标势}（electric scalar potential）$\phi$ 定义为（对偶矢量）四维矢量势 $\bA$ 的时间分量的\emph{相反数}，并把三维\emph{磁矢势} $\bA_{(3)}$ 取为 $\bA$ 的剩余部分。

\begin{keybox}[电势与磁矢势]
\textbf{[电势 $\phi$]}
\begin{equation}
\phi = -A_{t} \tag{Common}
\end{equation}

\textbf{[磁矢势 $\bA_{(3)}$]}
\begin{equation}
(A_{(3)x}, A_{(3)y}, A_{(3)z}) = (A_{x}, A_{y}, A_{z}) \tag{Common}
\end{equation}
\end{keybox}

最后我们用四维电流矢量 $\bj$ 定义\emph{电荷密度}（charge density）$\rho$ 与三维\emph{电流密度}（current density）$\bj_{(3)}$。

\begin{keybox}[电荷密度与电流密度]
\textbf{[电荷密度 $\rho$]}
\begin{equation}
\rho = j^{t} \tag{Common}
\end{equation}

\textbf{[电流密度 $\bj_{(3)}$]}
\begin{equation}
(j_{(3)}^{x}, j_{(3)}^{y}, j_{(3)}^{z}) = (j^{x}, j^{y}, j^{z}) \tag{Common}
\end{equation}
\end{keybox}

这些定义——正如我们给出的、上下指标恰如此处所列——无论我们做的是黎曼物理还是洛伦兹物理都相同。但若你想与关于洛伦兹版本的文献作比较，请记住：对 $t$ 指标进行升或降会得到原量的相反数。另外请注意，某些文献对洛伦兹度规使用 $(+---)$ 符号，而此处的洛伦兹公式使用 $(-+++)$。

连同由方程 (2) 用四维矢量势 $\bA$ 定义 $F$ 这一做法，这些定义让我们把电场描述为：其势 $\phi$ 的梯度的相反数、再减去磁矢势的时间变化率；而磁场描述为：其势 $\bA_{(3)}$ 的旋度。同样，这完全如同常规电磁学。

\begin{keybox}[由势给出的场]
\begin{equation}
\bE = -\nabla\phi - \partial_{t}\bA_{(3)} \tag{Common}
\end{equation}
\begin{equation}
\bB = \nabla\times\bA_{(3)} \tag{Common}
\end{equation}
\end{keybox}

接下来，考虑一个带电量为 $q$、四维速度为 $\bu$ 的粒子所受的四维力 $\mathbf{f}$：
\begin{equation}
\mathbf{f} = q F \mathbf{u} \tag{5}
\end{equation}
四维力是粒子能量-动量矢量 $\mathbf{P}$ 对固有时 $\tau$ 的变化率，因此这等价于：
\begin{equation}
\partial_{\tau}\mathbf{P} = q F \mathbf{u} \tag{6}
\end{equation}
现在，$\mathbf{P}$ 的空间部分就是三维动量 $\bp$，而四维速度 $\bu$ 的空间部分并不完全是普通速度 $\bv$。普通速度描述粒子的空间坐标对\emph{坐标时间} $t$ 的变化率，而 $\bu$ 的空间部分给出对\emph{固有时} $\tau$ 的变化率——因此 $\bu$ 的空间部分是 $(dt/d\tau)\bv$。不过，我们可以改用 $\bp$ 对坐标时间的变化率，从而吸收掉这个 $(dt/d\tau)$ 因子。

我们最终得到的称为洛伦兹力定律（Lorentz force law）。这同样是黎曼版与洛伦兹版共有的（尽管相对论性运动对粒子动量 $\bp$ 的影响当然\emph{不}相同）。

\begin{keybox}[洛伦兹力定律]
\begin{equation}
\partial_{t}\mathbf{p} = q(\bE + \bv\times\bB) \tag{Common}
\end{equation}
\end{keybox}

接下来，我们把由定义 [方程 (3)] 强加于 $F$ 的条件，翻译为对 $\bE$ 与 $\bB$ 的推论。当我们在方程 (3) 中取 $a$、$b$、$c$ 为 $x$、$y$、$z$ 时，它告诉我们 $\bB$ 的散度必须为零。这称为磁高斯定律（Gauss's Law For Magnetism），它陈述了不存在磁单极子（这在常规电磁学中成立，当然也存在磁单极子确实存在的思辨性理论）。

当我们在方程 (3) 中取 $a$、$b$、$c$ 为 $t$ 与两个空间坐标——对三对空间坐标中的每一对——时，它告诉我们 $\bE$ 的旋度与 $\bB$ 的时间变化率之和为零。这一结果称为麦克斯韦-法拉第方程（Maxwell-Faraday equation），它描述了磁场随时间变化时电场如何被产生出来。

\begin{keybox}[磁高斯定律与麦克斯韦-法拉第方程]
\begin{equation}
\nabla\cdot\bB = 0 \tag{Common}
\end{equation}
\begin{equation}
\nabla\times\bE + \partial_{t}\bB = 0 \tag{Common}
\end{equation}
\end{keybox}

\subsection{改变（Changes）}

最后我们来到黎曼电磁学与洛伦兹电磁学的差别，它们来自把麦克斯韦方程组中涉及场源的那个方程替换为黎曼 Proca 方程。

当我们在黎曼 Proca 方程或麦克斯韦方程中把指标 $a$ 取为 $t$ 时，我们得到两个版本的高斯定律：在麦克斯韦情形中，它告诉我们电通量线只在电荷处起始与终止。在黎曼 Proca 情形中，这不再成立：通量线会从真空中冒出来，在这一点上电势的表现恰如一个电荷密度。

\begin{keybox}[高斯定律]
\begin{equation}
\nabla\cdot\bE = \wms^{2}\phi - \rho \tag{Riemannian}
\end{equation}
\begin{equation}
\nabla\cdot\bE = \rho \tag{Lorentzian}
\end{equation}
\end{keybox}

当我们在黎曼 Proca 方程或麦克斯韦方程中把指标 $a$ 取为每个空间坐标时，我们得到两个版本的安培-麦克斯韦定律（Ampère-Maxwell Law），它描述了磁场由电流、变化的电场——以及在黎曼情形中，直接由磁矢势——所产生。

\begin{keybox}[安培-麦克斯韦定律]
\begin{equation}
\nabla\times\bB + \partial_{t}\bE = \wms^{2}\bA_{(3)} + \bj_{(3)} \tag{Riemannian}
\end{equation}
\begin{equation}
\nabla\times\bB - \partial_{t}\bE = \bj_{(3)} \tag{Lorentzian}
\end{equation}
\end{keybox}

\section{静电学的例子（Examples from Electrostatics）}

\subsection{库仑势（The Coulomb Potential）}

我们在电磁学注记的主页讨论过黎曼库仑势。现在我们有了推导该势的工具。

\orthofig{images/04__007.png}{点电荷周围的电通量线}

我们感兴趣的是静止于我们坐标系中的一个点电荷 $q$ 周围的场。该情形不随时间变化、在空间上完全径向对称，因此它实际上是一个一维问题：一切都是到电荷距离 $r$ 的函数。

黎曼电磁学给这一情形带来的奇特转折是：电通量线——在洛伦兹电磁学中总在电荷处起始与终止——现在可以在真空中间终止，而这一效应依赖于电势。右图中箭头表示电场的方向——但此处画出的是通量线而\emph{非}矢量：场的强度由通量线的疏密表示，而非其长度。

与洛伦兹情形共有的主要数学困难在于：粒子所在位置处的电荷密度是\emph{无穷}的。绕过它的办法是对电荷密度作空间积分；对一个包含该粒子的区域积分将给出有限值 $q$。但我们的方程都是微分方程，因此首先要把其中一个化为合适的形式。为此我们利用散度定理（divergence theorem）——纯数学中的一条结论：对任意矢量场 $\bE$ 与任意空间区域，$\bE$ 与区域表面外法向的点积的积分，等于 $\bE$ 的散度在该区域上的体积分：
\begin{equation}
\int_{\text{Surface}} \bE\cdot\mathbf{n} = \int_{\text{Volume}} \nabla\cdot\bE \tag{7}
\end{equation}

我们取矢量场 $\bE$ 为电场，取积分区域为围绕点电荷、半径为 $r$ 的球面。我们预期电势 $\phi$ 只是 $r$ 的函数，且由问题的径向对称性，预期电场沿径向指向或背离电荷。鉴于问题是静态的，我们可以仅用电势表示 $\bE$：
\begin{align}
\bE(r) &= -\nabla\phi(r) \nonumber\\
&= -\phi'(r)\,\be_{r} \tag{8}
\end{align}
其中 $\be_{r}$ 是指离电荷的单位矢量。黎曼版的高斯定律给出 $\nabla\cdot\bE$ 作为电荷密度 $\rho$ 与电势 $\phi$ 的函数。我们把这一点、连同 (8)，代入 (7) 中作用于电荷周围球形区域的积分。我们还利用这样的事实：$\rho$ 在包含电荷的区域上的任何体积分都简单地给出 $q$。两边除以 $4\pi$ 后，得到：
\begin{equation}
-r^{2}\phi'(r) = \wms^{2}\int_{0}^{r}\phi(s)\,s^{2}\,ds - \frac{q}{4\pi} \tag{9}
\end{equation}

现在，根据处理常规静电学的经验，我们猜测：若把 $\phi$ 写成一个新函数 $f$ 除以 $r$，事情会更简单：
\begin{equation}
\phi(r) = \frac{f(r)}{r} \tag{10}
\end{equation}
\begin{equation}
\phi'(r) = \frac{f'(r)}{r} - \frac{f(r)}{r^{2}} \tag{11}
\end{equation}

用新函数 $f$ 表示，方程 (9) 变为 (12)；我们取 (12) 在 $r=0$ 处求值，得到 (12a)：
\begin{equation}
f(r) - r f'(r) = \wms^{2}\int_{0}^{r} f(s)\,s\,ds - \frac{q}{4\pi} \tag{12}
\end{equation}
\begin{equation}
f(0) = -\frac{q}{4\pi} \tag{12a}
\end{equation}

(12) 对 $r$ 求导给出 (13)，稍加改写给出 (14)：
\begin{equation}
-r f''(r) = \wms^{2} f(r)\,r \tag{13}
\end{equation}
\begin{equation}
f''(r) + \wms^{2} f(r) = 0 \tag{14}
\end{equation}

方程 (14) 是一个非常著名的微分方程，其实值通解为：
\begin{equation}
f(r) = C_{1}\cos(\wms r) + C_{2}\sin(\wms r) \tag{15}
\end{equation}
方程 (12a) 告诉我们 $C_{1} = -q/(4\pi)$。

$C_{2}$ 呢？我们为 $C_{2}$ 选取任何值都会给出问题的一个有效解，但既然这一项与点电荷 $q$ 无关，我们令 $C_{2}$ 为零。正如常规电磁学，一个问题的通解往往包含某种仅仅路过关注区域的辐射——在此即碰巧在电荷静止系中静止的径向对称辐射。

于是我们推导出了黎曼库仑势。下面也给出相应的电场 $\bE = -\nabla\phi$。

\begin{keybox}[库仑势与库仑场]
\textbf{[库仑势]}
\[
\phi(r) = \begin{cases}
-\dfrac{q}{4\pi r}\cos(\wms r), & \text{(Riemannian)}\\[8pt]
\dfrac{q}{4\pi r}, & \text{(Lorentzian)}
\end{cases}
\]

\textbf{[库仑场]}
\[
\bE(r) = \begin{cases}
-\dfrac{q}{4\pi r^{2}}
\bigl[\cos(\wms r) + \wms r\sin(\wms r)\bigr]\be_{r}, & \text{(Riemannian)}\\[9pt]
\dfrac{q}{4\pi r^{2}}\,\be_{r}, & \text{(Lorentzian)}
\end{cases}
\]
\end{keybox}

\subsection{黎曼电磁学的格林函数（A Green's Function for Riemannian Electromagnetism）}

单个静止点电荷的库仑势原则上让我们能通过对场源的积分求出任何静态电荷分布的电场。然而，拥有一个更基本的、黎曼电动力学方程的解会很有用：它联系于一个瞬时"脉冲"电荷——它在四维空间某个事件处产生，随即消失。显然这种行为违反电荷守恒，但把该解沿任意多个具有完整历史的电荷的世界线积分，就能找到一个尊重电荷守恒的解。

这样一个基本解称为\emph{格林函数}（Green's function）。

我们先寻找黎曼标量波方程（Riemannian Scalar Wave equation）的四维旋转对称解（无源项）。这就是四维的亥姆霍兹方程（Helmholtz equation），施加四维旋转对称后，我们得到关于单变量 $s$（四维空间中到原点的距离）的函数 $G$ 的常微分方程：
\begin{equation}
G''(s) + \frac{3}{s}G'(s) + \wms^{2}G(s) = 0 \tag{16}
\end{equation}
该方程的通解为：
\begin{equation}
G(s) = C_{1}\frac{J_{1}(\wms s)}{s} + C_{2}\frac{Y_{1}(\wms s)}{s} \tag{17}
\end{equation}
其中 $J_{1}$ 与 $Y_{1}$ 分别是第一类与第二类贝塞尔函数（Bessel function）。尽管这是无源方程在 $s>0$ 时的解，但当 $s$ 趋于零时贝塞尔函数 $Y_{1}$ 趋于负无穷，这正暗示了与点电荷相联系的格林函数所应有的那种奇异行为。

借助变量替换，我们可以显式地把 $G$ 沿一个静止点电荷的整条世界线积分——把沿世界线的时间坐标 $t$ 换成 $s = \sqrt{r^{2}+t^{2}}$，即从世界线上某个事件到离点电荷空间距离为 $r$ 的某个事件之间的四维距离。利用 $t = \sqrt{s^{2}-r^{2}}$ 与 $dt = (s/t)\,ds$，有：
\begin{align}
\int_{-\infty}^{\infty}G\!\left(\sqrt{r^{2}+t^{2}}\right)dt
&= 2\int_{r}^{\infty}G(s)\,\frac{s}{\sqrt{s^{2}-r^{2}}}\,ds \nonumber\\
&= 2\int_{r}^{\infty}\frac{C_{1}J_{1}(\wms s) + C_{2}Y_{1}(\wms s)}{\sqrt{s^{2}-r^{2}}}\,ds \nonumber\\
&= 2\,\frac{C_{1}\sin(\wms r) - C_{2}\cos(\wms r)}{\wms r} \tag{18}
\end{align}

把它与一个带电量为 $q$ 的点粒子的黎曼库仑势相匹配，可令 $C_{1}=0$、$C_{2}=q\wms/(8\pi)$。

我们是对标量势 $\phi$ 做的这一计算，但若用四维矢量来表达，结果会最有用。用四维矢量的语言说，粒子世界线的每一无穷小段对四维矢量势 $\bA$ 贡献一个平行于粒子四维速度 $\bu$ 的部分。我们加上一个负号，因为 $\phi = -A_{t}$。

\begin{keybox}[格林函数]
粒子带电量为 $q$。其世界线 $\mathbf{y}(\tau)$ 以固有时 $\tau$ 参数化。其四维速度 $\bu(\tau) = \partial_{\tau}\mathbf{y}(\tau)$。四维矢量势 $\bA$ 在事件 $\mathbf{x}$ 处取值。
\begin{equation}
d\bA(\mathbf{x}) = -\bu(\tau)\,\frac{q\wms}{8\pi}\,
\frac{Y_{1}\bigl(\wms|\mathbf{x}-\mathbf{y}(\tau)|\bigr)}{|\mathbf{x}-\mathbf{y}(\tau)|}\,d\tau \tag{Riemannian}
\end{equation}
\end{keybox}

我们不作洛伦兹对应物的推导，因为那需要大段绕路来解释所有细节与差别。我们只指出：所谓李纳-维谢尔势（Liénard-Wiechert potential）在给定事件处，只依赖电荷世界线与该事件的过去光锥之交点处电荷的位置与四维速度。换言之，正如你所料，在洛伦兹物理中 $\bA$ 只受粒子从过去以光速传播过来的信息的影响。

我们此处给出的黎曼格林函数不区分过去与未来。对静电与静磁问题这没问题，但我们须记住：若把它用于产生电磁波的情形，它将给出既含入射波又含出射波的解。

\subsection{电偶极子（Electric Dipoles）}

电偶极子（electric dipole）由两个强度相等、一正一负、并保持固定间距的点电荷组成。若两电荷彼此靠近，它们会倾向于互相抵消库仑势，但会剩下一个特征性的\emph{偶极场}（dipole field）。

\orthofig{images/04__008.png}{电偶极子}

我们可以通过研究一个极限情形来简化对偶极场形状的理解：让两电荷彼此越来越靠近，同时每个电荷的强度越来越大。若我们定义矢量 $\bp$（\emph{偶极矩}，dipole moment）为从负电荷指向正电荷的位移矢量\emph{乘以}电荷的（正）强度，那么我们取这样的极限：$\bp$ 保持恒定有限，而间距趋于零、每个电荷的强度趋于无穷。

得到该极限的最简单方式是：沿与所选偶极矩相反的方向对库仑势求\emph{导数}。所得电势见右图，电势与电场的公式见下表。这里 $\mathbf{r}$ 是从偶极子位置到我们求场之点的三维矢量，$r$ 是其大小。

\begin{keybox}[电偶极子]
\textbf{[电偶极子电势]}
\[
\phi(\mathbf{r}) = \begin{cases}
-\dfrac{\bp\cdot\mathbf{r}}{4\pi r^{3}}
\bigl[\cos(\wms r) + \wms r\sin(\wms r)\bigr], & \text{(Riemannian)}\\[9pt]
\dfrac{\bp\cdot\mathbf{r}}{4\pi r^{3}}, & \text{(Lorentzian)}
\end{cases}
\]

\textbf{[电偶极子电场]}
\[
\bE(\mathbf{r}) = \begin{cases}
-\dfrac{3(\bp\cdot\mathbf{r})\mathbf{r} - r^{2}\bp}{4\pi r^{5}}
\bigl[\cos(\wms r) + \wms r\sin(\wms r)\bigr]
+ \dfrac{\wms^{2}(\bp\cdot\mathbf{r})\mathbf{r}}{4\pi r^{3}}\cos(\wms r), & \text{(Riemannian)}\\[12pt]
\dfrac{3(\bp\cdot\mathbf{r})\mathbf{r} - r^{2}\bp}{4\pi r^{5}}, & \text{(Lorentzian)}
\end{cases}
\]
\end{keybox}

若你在看到黎曼偶极势时有一种\emph{似曾相识}（déjà vu）之感，那你多半在\emph{常规}电磁学中见过\emph{振荡}偶极子势的非常相似的图形。事实上，静态黎曼偶极子的场，恰与常规电磁学中对一个振荡偶极子相关的入射与出射辐射求和所构造的驻波的空间部分相同。

\subsection{带电球壳（Charged Spherical Shells）}

假设总电荷 $Q$ 均匀分布在半径为 $R$ 的球壳上。洛伦兹电磁学中有一个众所周知的结果：球外电势与位于球心处点电荷的电势完全相同，而球内电势为常数。然而在黎曼电磁学中结果大不相同！无论是对表面各处的贡献作显式积分，还是用适当形式的高斯定律，我们都得到如下结果。

\begin{keybox}[均匀带电球壳]
半径为 $R$、总电荷为 $Q$ 的球壳。

\textbf{[电势（Riemannian）]}
\[
\phi(r) = \begin{cases}
-\dfrac{Q\cos(\wms R)\sin(\wms r)}{4\pi\wms R r}, & r<R\\[9pt]
-\dfrac{Q\sin(\wms R)\cos(\wms r)}{4\pi\wms R r}, & r>R
\end{cases}
\]

\textbf{[电势（Lorentzian）]}
\[
\phi(r) = \begin{cases}
\dfrac{Q}{4\pi R}, & r<R\\[9pt]
\dfrac{Q}{4\pi r}, & r>R
\end{cases}
\]

\textbf{[电场（Riemannian）]}
\[
\bE(r) = \begin{cases}
-\dfrac{Q\cos(\wms R)}{4\pi\wms R r^{2}}
\bigl[\sin(\wms r) - \wms r\cos(\wms r)\bigr]\be_{r}, & r<R\\[10pt]
-\dfrac{Q\sin(\wms R)}{4\pi\wms R r^{2}}
\bigl[\cos(\wms r) + \wms r\sin(\wms r)\bigr]\be_{r}, & r>R
\end{cases}
\]

\textbf{[电场（Lorentzian）]}
\[
\bE(r) = \begin{cases}
\mathbf{0}, & r<R\\[9pt]
\dfrac{Q}{4\pi r^{2}}\,\be_{r}, & r>R
\end{cases}
\]
\end{keybox}

\orthofig{images/04__009.png}{带电球壳}

在黎曼情形中，球壳的外部电势是点电荷电势乘以一个因子 $\sin(\wms R)/(\wms R)$，而内部电势则交换了 $r$ 与 $R$ 的角色。

尽管内部电势一般不是常数，但对某些 $R$ 值，内部\emph{或}外部电势会为零。当 $\wms R$ 为 $\pi/2$ 的奇数倍、等价地当 $R$ 为光的最小波长 $\lambda_{\min}$ 的四分之一波长的奇数倍时，内部电势为零。当 $\wms R$ 为 $\pi$ 的整数倍、等价地当 $R$ 为 $\lambda_{\min}$ 一半的整数倍时，外部电势为零。

当然，这些精确的相消对电荷分布的确切几何非常敏感。不过一般而言，外部电势会比点电荷的电势大幅减弱。

\subsection{带电实心球（Charged Solid Spheres）}

我们可以对球壳的结果积分，得到电荷 $Q$ 均匀分布于整个实心球体时的电势与电场。

\begin{keybox}[均匀带电实心球]
半径为 $R$、总电荷为 $Q$ 的实心球。

\textbf{[电势（Riemannian，$r<R$）]}
\begin{equation}
\phi(r) = \frac{3Q}{4\pi\wms^{2}R^{3}}
\left\{1 - \bigl[\cos(\wms R) + \wms R\sin(\wms R)\bigr]\frac{\sin(\wms r)}{\wms r}\right\}
\end{equation}

\textbf{[电势（Lorentzian，$r<R$）]}
\begin{equation}
\phi(r) = \frac{Q}{8\pi R}\left[3 - \left(\frac{r}{R}\right)^{2}\right]
\end{equation}

\textbf{[电场（Riemannian，$r<R$）]}
\begin{equation}
\bE(r) = -\frac{3Q\bigl[\cos(\wms R) + \wms R\sin(\wms R)\bigr]}{4\pi\wms^{3}R^{3}r^{2}}
\bigl[\sin(\wms r) - \wms r\cos(\wms r)\bigr]\be_{r}
\end{equation}

\textbf{[电场（Lorentzian，$r<R$）]}
\begin{equation}
\bE(r) = \frac{Q r}{4\pi R^{3}}\,\be_{r}
\end{equation}
\end{keybox}

\orthofig{images/04__029.png}{带电实心球}

在洛伦兹情形中，与球壳一样，实心球外的电势与电场就是集中于球心处点电荷的电势与电场。在实心球内部，电场来自比你更靠近球心的那部分球体，因此它随到球心的距离线性增大，而电势是到球心距离的二次函数。

在黎曼情形中，外部电势与电场是点电荷的值乘以一个依赖球体大小的因子：
\[
\frac{3\bigl[\sin(\wms R) - \wms R\cos(\wms R)\bigr]}{\wms^{3}R^{3}}
\]
该因子随 $R$ 振荡，其第一个零点在 $R\approx 0.715\,\lambda_{\min}$ 处。

内部电势由一个依赖 $R$ 但不振荡的平坦项，加上一个同时随 $R$ 与 $r$ 振荡的项构成。适当地选取 $R$ 可使振荡部分为零，从而使电势在整个球体内保持平坦，其第一个零点在 $R\approx 0.445\,\lambda_{\min}$ 处。

\subsection{电容器（Capacitors）}

假设我们有两个同心带电球壳，带有等量异号电荷。这一构置构成一个带电\emph{电容器}（capacitor）。电子电路中现实世界的电容器通常比这复杂得多，但这个简单几何允许我们做一些精确计算，以演示电容在黎曼宇宙中如何运作。

\begin{keybox}[球形电容器]
内壳半径 $R_{1}$，总电荷 $-Q$；外壳半径 $R_{2}$，总电荷 $+Q$。

\textbf{[电势（Riemannian，$r<R_{2}$）]}
\begin{equation}
\phi(r) = \frac{Q\sin(\wms r)\bigl[R_{2}\cos(\wms R_{1}) - R_{1}\cos(\wms R_{2})\bigr]}{4\pi\wms r R_{1}R_{2}}
\end{equation}

\textbf{[电势（Lorentzian，$r<R_{2}$）]}
\begin{equation}
\phi(r) = \frac{Q(R_{1} - R_{2})}{4\pi R_{1}R_{2}}
\end{equation}

\textbf{[电场（Riemannian，$r<R_{2}$）]}
\begin{equation}
\bE(r) = \frac{Q\bigl[\sin(\wms r) - \wms r\cos(\wms r)\bigr]
\bigl[R_{2}\cos(\wms R_{1}) - R_{1}\cos(\wms R_{2})\bigr]}{4\pi\wms r^{2}R_{1}R_{2}}\,\be_{r}
\end{equation}

\textbf{[电场（Lorentzian，$r<R_{2}$）]}
\begin{equation}
\bE(r) = \mathbf{0}
\end{equation}

\textbf{[电容（Riemannian）]}
\begin{equation}
C = \frac{8\pi\wms R_{1}^{2}R_{2}^{2}}
{4R_{1}R_{2}\sin(\wms R_{1})\cos(\wms R_{2})
 - R_{1}^{2}\sin(2\wms R_{2}) - R_{2}^{2}\sin(2\wms R_{1})}
\end{equation}

\textbf{[电容（Lorentzian）]}
\begin{equation}
C = \frac{4\pi R_{1}R_{2}}{R_{2} - R_{1}}
\end{equation}
\end{keybox}

\orthofig{images/04__010.png}{球形电容器}

在洛伦兹情形中，电势总是从内壳上的负值升到外壳上的零，器件两端的电压被定义为正值：
\[
V_{\text{Lorentzian}} = \phi(R_{2}) - \phi(R_{1}) = \frac{Q(R_{2} - R_{1})}{4\pi R_{1}R_{2}}
\]
总正电荷与电压差之间的比例常数称为器件的\emph{电容}（capacitance）$C$：
\[
C_{\text{Lorentzian}} = \frac{Q}{V_{\text{Lorentzian}}} = \frac{4\pi R_{1}R_{2}}{R_{2} - R_{1}}
\]

在黎曼情形中，两壳之间的电压差仍正比于总电荷，我们可以用同样的方式定义电容，但公式（见上表）复杂得多，且对由光的最小波长所设定的长度尺度敏感。黎曼电容原则上可以是正的、负的、甚至无穷大。无穷大电容意味着你可以向该器件注入任意多的电荷，而两壳自身之间不会建立起电压——尽管电场仍会增大。负电容意味着所带正电荷过剩的球壳处于比所带负电荷过剩的球壳更低的电势，因此若把两者连通，正电荷的球壳会吸入更多正电荷。当你把普通电容器短路时，它会放电；当你把具有负电容的电容器短路时，它的电荷会\emph{增加}。

显然这可能引发失控过程，而我们（高度简化的）分析中没有任何东西能表明它何时会终止。但在一个更详细的、包含所有涉及材料性质的负电容器电路模型中，最终会出现一些复杂性来中止电荷的积聚。同样，黎曼库仑势允许\emph{同种电荷相吸}的情形这一单纯事实，似乎威胁到"宇宙中所有正电荷最终可能聚集到一处"的可能性——但那种场景忽略了量子力学效应对同种带电粒子聚集所施加的限制。

同样重要的是要注意：我们研究的是一种理想化情形，即球壳完全光滑、电荷均匀分布，其尺度远小于光的最小波长。任何大于该尺度的起伏都会产生一个混合了正电容与负电容的器件，从而导致那种在黎曼宇宙中缓和一切静电现象的相消。

此外，整个分析都假定电荷与电压的任何变化都发生得非常缓慢。我们将在后一节中处理交流电中的电容器。

\section{静磁学算例（Examples from Magnetostatics）}

上一节我们建立了黎曼电磁学的场方程与四维矢量势语言。本节把它用于三类最基本的静磁构型——无限长直线电流、磁偶极子与螺线管——并顺带导出毕奥-萨伐尔定律的黎曼版本。我们将反复看到同一个模式：黎曼静态解与洛伦兹振荡解的\emph{空间部分}具有相同的形状，差别只在于黎曼解在空间中以最大频率 $\wms$ 振荡，从而在宏观尺度上不断自我抵消。

\subsection{线电流（Linear Current）}

设有一恒定电流 $I$ 沿一根无限长、无限细的直导线流动。黎曼版安培-麦克斯韦定律（Ampère-Maxwell Law）给出磁场旋度 $\nabla\times\bB$ 作为电流密度与三维磁势 $\bA_{(3)}$ 的函数。但由于我们把电流视为集中于一根无限细的导线上，最方便的做法是借助开尔文-斯托克斯定理（Kelvin-Stokes theorem）把该定律改写为积分形式——它把一个矢量场旋度在一曲面上的积分，与该场沿曲面边界曲线的线积分联系起来：
\begin{equation}
\int_{\text{Surface}} (\nabla \times \bB)\cdot\mathbf{n} = \int_{\text{Boundary}} \bB\cdot\mathbf{t} \tag{19}
\end{equation}
其中 $\mathbf{n}$ 是曲面的单位法矢，$\mathbf{t}$ 是构成曲面边界之曲线的单位切矢；若以我们所选的 $\mathbf{n}$ 定义“上方”，则 $\mathbf{t}$ 沿曲面逆时针方向环绕。

若取一个以导线为中心、垂直于导线的半径为 $r$ 的圆盘作为曲面，由对称性可预期磁势 $\bA_{(3)}$ 平行于导线，且只是到导线距离 $r$ 的函数。若令导线沿 $z$ 轴放置，我们就有：
\begin{equation}
\bA_{(3)}(r) = A(r)\,\be_z \tag{20}
\end{equation}
\begin{equation}
\begin{aligned}
\bB(r) &= \nabla \times \bA_{(3)}(r) \\
       &= \partial_y A(r)\,\be_x - \partial_x A(r)\,\be_y \\
       &= A'(r)\Big[\frac{y}{r}\,\be_x - \frac{x}{r}\,\be_y\Big] \\
       &= -A'(r)\,\be_\phi
\end{aligned}
\tag{21}
\end{equation}
其中 $\be_\phi$ 是绕导线逆时针指向的单位矢量场。再应用开尔文-斯托克斯定理，即式 (19)，以及安培-麦克斯韦定律，得到：
\begin{equation}
I + 2\pi \wms^2 \int_0^r A(s)\,s\,ds = -2\pi r\,A'(r) \tag{22}
\end{equation}
两端同除以 $2\pi$、对 $r$ 求导并稍作整理，得：
\begin{equation}
A''(r) + \frac{A'(r)}{r} + \wms^2 A(r) = 0 \tag{23}
\end{equation}
该微分方程的通解为：
\begin{equation}
A(r) = C_1 J_0(\wms r) + C_2 Y_0(\wms r) \tag{24}
\end{equation}
其中 $J_0$、$Y_0$ 分别是第一类与第二类贝塞尔函数（Bessel function）。对这些贝塞尔函数求导给出：
\begin{equation}
A'(r) = -\wms\big[C_1 J_1(\wms r) + C_2 Y_1(\wms r)\big] \tag{25}
\end{equation}
在此结果下，式 (22) 右端在 $r\to0$ 时的极限为：
\begin{equation}
\lim_{r\to0}\big(-2\pi r\,A'(r)\big) = -4 C_2 \tag{26}
\end{equation}
而式 (22) 左端取同一极限则恰好就是电流 $I$。于是 $C_2 = -I/4$。$C_1$ 仍待定，但与推导库仑势时一样，我们把 $C_1$ 项取为来自过去、与电流 $I$ 无关的静止辐射场。

\begin{align}
\text{线电流磁势}\quad
\bA_{(3)}(r) &= -\frac{I}{4}\,Y_0(\wms r)\,\be_z && \text{（Riemannian）}\\
\bA_{(3)}(r) &= -\frac{I}{2\pi}\,\log(r)\,\be_z && \text{（Lorentzian）}\\[4pt]
\text{线电流磁场}\quad
\bB(r) &= -\frac{I\,\wms}{4}\,Y_1(\wms r)\,\be_\phi && \text{（Riemannian）}\\
\bB(r) &= \frac{I}{2\pi r}\,\be_\phi && \text{（Lorentzian）}
\end{align}

由于贝塞尔函数是振荡的，导线周围磁场反向变化的长度尺度，与点电荷周围电场的交替尺度相仿。

\orthofig{images/04__011.png}{直线导体周围的磁通量}

由于在紧邻电流处洛伦兹情形与黎曼情形中磁场的方向相同，而洛伦兹力定律在两种情形下也相同，理论上两根足够靠近（且足够细）的平行载流导线之间会感受到吸引力。然而正如静电力一样，场的空间振荡会对任何宽度超过振荡波长的物体造成显著抵消。

这个例子再一次体现了黎曼静态解与洛伦兹振荡解空间部分之间的联系：导线周围的黎曼场，等同于常规电磁学中一根\emph{振荡}导线周围驻波的空间部分。当然现实中的振荡电流通常伴随纯外向行波，但在存在一列等强度入射波的情况下，就会产生形状恰好如此的驻波。

\subsection{毕奥-萨伐尔定律（The Biot-Savart Law）}

在常规静磁学中，毕奥-萨伐尔定律（Biot-Savart Law）给出沿细导线流动的恒定电流 $I$ 所产生的磁场：
\begin{equation}
\bB = \frac{I}{4\pi}\int \frac{\mathbf{t} \times \mathbf{r}}{r^3}\,dl \tag{27}
\end{equation}
这里积分变量 $l$ 是沿导线的长度，$\mathbf{r}$ 是从导线上的一个线元指向磁场 $\bB$ 待求点的三维位移矢量，$\mathbf{t}$ 是导线的单位切矢。

我们将利用此前导出的黎曼格林函数来得到其黎曼对应物。

导线长度为 $dl$ 的每个线元都视为同时包含运动电荷与静止电荷，二者的电量均为 $dq = \rho\,dl$，其中 $\rho$ 是导线内的线电荷密度。运动电荷对格林函数积分的贡献为 $dq\,\bu\,d\tau$——其中 $\bu$ 是该线元中每个运动电荷的四维速度——但我们知道该矢量的时间分量会被导线中存在的等量静止电荷精确抵消，因为导线整体视为电中性。$\bu\,d\tau$ 的空间部分正是 $\bv\,dt$，其中 $\bv$ 是运动电荷的普通三维速度，而 $t$ 是导线静止参考系中的坐标时。又因为流过导线的电流 $I = \rho v$——用矢量表示为 $I\mathbf{t} = \rho\bv$，其中 $\mathbf{t}$ 是导线单位切矢——我们有：
\begin{equation}
\begin{aligned}
(dq\,\bu\,d\tau)_{\text{net}} &= \rho\,dl\,\bv\,dt \\
 &= I\,\mathbf{t}\,dt\,dl
\end{aligned}
\tag{28}
\end{equation}
于是我们可把格林函数对 $t$ 积分，而把 $I\,\mathbf{t}\,dl$ 视为常量；这个积分与我们此前用来从格林函数得到库仑势的积分相同。因此毫不意外，由该积分得到的磁势看起来就像一个库仑势，而对其取旋度得到的磁场，其大小（方向则不然）与库仑电场相同。

\begin{align}
\text{毕奥-萨伐尔定律（磁势）}\quad
\bA_{(3)}(\mathbf{r}) &= \frac{I}{4\pi}\int \frac{\mathbf{t} \cos(\wms r)}{r}\,dl && \text{（Riemannian）}\\
\bA_{(3)}(\mathbf{r}) &= \frac{I}{4\pi}\int \frac{\mathbf{t}}{r}\,dl && \text{（Lorentzian）}\\[4pt]
\text{毕奥-萨伐尔定律（磁场）}\quad
\bB(\mathbf{r}) &= \frac{I}{4\pi}\int \big[\cos(\wms r) + \wms r\sin(\wms r)\big]\frac{\mathbf{t}\times\mathbf{r}}{r^3}\,dl && \text{（Riemannian）}\\
\bB(\mathbf{r}) &= \frac{I}{4\pi}\int \frac{\mathbf{t}\times\mathbf{r}}{r^3}\,dl && \text{（Lorentzian）}
\end{align}

用毕奥-萨伐尔定律对无限长直导线周围的磁势作显式积分，所得结果与我们此前得到的公式一致。

\subsection{磁偶极子（Magnetic Dipoles）}

\textbf{磁偶极子}（magnetic dipole）是产生某种简单而高度对称磁场的系统。一小圈环流电流，或具有量子力学自旋的带电粒子，都属于此类；许多具有更复杂磁场的系统，从远处看也表现为磁偶极子。

\orthofig{images/04__012.png}{磁偶极子周围的磁通量}

对于一圈电流环，\textbf{磁矩}（magnetic moment）——记作 $\boldsymbol\mu$——定义为垂直于电流环的矢量，其大小等于环面积与环流电流强度之积。约定是：当拇指与磁矩矢量对齐时，电流沿右手手指方向环绕。纯偶极场既可取有限电流环场中的主导项（即随距离衰减最慢的项），也可取环面积趋于零而电流趋于无穷、二者乘积保持有限这一极限情形下的场。

在洛伦兹电磁学中，磁偶极子的磁场恰好取与\emph{电}偶极子的\emph{电}场相同的数学形式。然而这在黎曼情形下不可能成立，因为磁场 $\bB$ 必须在处处满足 $\nabla\cdot\bB = 0$——也就是说磁通线构成不中断的闭合回路——但这一条对黎曼电场即使在真空中也不成立，而电偶极子场的电场线会从远离偶极子本身的地方起始并终止。

我们可以用毕奥-萨伐尔定律求小电流环极限下的磁偶极势。与电偶极子一样，我们对某个适当量求导以取极限。此处我们对\emph{半}个电流环积分——即对某个线元与正对其的线元（该处电流方向相反）的贡献之和积分。在小环极限下，该和式恰好是 $1/r$ 或 $\cos(\wms r)/r$ 跨过电流环的方向导数在环心处的取值，再乘以环的直径与环的切矢。

\begin{align}
\text{磁偶极势}\quad
\bA_{(3)}(\mathbf{r}) &= \frac{\boldsymbol\mu\times\mathbf{r}}{4\pi r^3}\big[\cos(\wms r) + \wms r\sin(\wms r)\big] && \text{（Riemannian）}\\
\bA_{(3)}(\mathbf{r}) &= \frac{\boldsymbol\mu\times\mathbf{r}}{4\pi r^3} && \text{（Lorentzian）}\\[6pt]
\text{磁偶极场}\quad
\bB(\mathbf{r}) &= \frac{3(\boldsymbol\mu\cdot\mathbf{r})\mathbf{r} - r^2\boldsymbol\mu}{4\pi r^5}\big[\cos(\wms r) + \wms r\sin(\wms r)\big] \nonumber\\
 &\qquad - \frac{\wms^2(\boldsymbol\mu\times\mathbf{r})\times\mathbf{r}}{4\pi r^3}\cos(\wms r) && \text{（Riemannian）}\\
\bB(\mathbf{r}) &= \frac{3(\boldsymbol\mu\cdot\mathbf{r})\mathbf{r} - r^2\boldsymbol\mu}{4\pi r^5} && \text{（Lorentzian）}
\end{align}

在洛伦兹电磁学中，尽管并非所有材料都能被磁化，但使大量磁偶极子（一般是电子自旋）得以联合产生强得多磁场的条件并不十分苛刻：只要一群偶极子的磁矩矢量相互平行，它们对外磁场的贡献就会彼此增强。但由于黎曼磁偶极场在极小的长度尺度上就反向，在任意一群偶极子中，它们的场之间会发生大量抵消——合成场同样呈现那种空间振荡。在黎曼宇宙中，不可能存在我们那种能在长距离上沿单一方向持续施力的永磁体。

\subsection{螺线管与电感（Solenoids and Inductance）}

\textbf{螺线管}（solenoid）是螺旋状盘绕的导线。当有恒定电流通过时，我们将近似求出线圈内外的磁场，假定螺线管足够长，以致可以忽略两端处究竟发生什么。实际上我们分析的是一个无限长螺线管，它比有限长更易处理，因为可以近似认为它同时具有沿轴的平移对称性与绕轴的旋转对称性。

具有这种柱对称、且磁场沿 $z$ 轴指向的磁势与磁场的最一般解为：
\begin{equation}
\begin{aligned}
\bA_{(3)}(r) &= \big[a J_1(\wms r) + b Y_1(\wms r)\big]\be_\phi\\
\bB(r) &= \big[a \wms J_0(\wms r) + b \wms Y_0(\wms r)\big]\be_z
\end{aligned}
\tag{29a,b}
\end{equation}

\orthofig{images/04__013.png}{螺线管的磁通量}

但我们需要允许线圈内外的解不同，于是有四个待定系数：$a_{\text{int}}$、$b_{\text{int}}$、$a_{\text{ext}}$、$b_{\text{ext}}$。解在 $r = 0$ 处必须有限，这要求 $b_{\text{int}} = 0$；我们又要求 $\bA_{(3)}$ 在线圈半径 $r = R$ 处连续。第三个关系来自对一条细长竖直矩形回路应用安培定律，该回路包围螺线管单位高度上 $n$ 匝绕组中的电流；它告诉我们，紧贴线圈内侧与外侧的 $\bB$ 场之差等于该电流。

要完全定出解还需第四个方程，这要多花点功夫。沿线圈的一条竖直窄条对毕奥-萨伐尔定律给出的 $\bA_{(3)}$ 贡献积分并不难，但要得到\emph{绕}线圈一周积分的精确表达式却难以处理。不过我们可以先对线圈中心附近一小距离处 $\bA_{(3)}$ 的贡献求出关于 $r$ 的一阶泰勒展开，再对\emph{该展开}绕整个线圈积分；把它与由一般解得到的等价泰勒展开相匹配，便给出 $a_{\text{int}}$ 的值，然后解其余方程即可确定全部系数。结果 $a_{\text{ext}} = 0$，于是内部解与外部解都各只剩一项。

在常规电磁学中，无限长螺线管外的磁场为零，但在黎曼情形下一般并非如此。

\begin{align}
\text{长螺线管（磁势）}\quad
\bA_{(3)}(r) &= \begin{cases}
-\tfrac12 n I \pi R\,Y_1(\wms R)\,J_1(\wms r)\,\be_\phi, & r < R\\[4pt]
-\tfrac12 n I \pi R\,J_1(\wms R)\,Y_1(\wms r)\,\be_\phi, & r > R
\end{cases} && \text{（Riemannian）}\\
\bA_{(3)}(r) &= \begin{cases}
\dfrac{n I r}{2}\,\be_\phi, & r < R\\[6pt]
\dfrac{n I R^2}{2r}\,\be_\phi, & r > R
\end{cases} && \text{（Lorentzian）}\\[6pt]
\text{长螺线管（磁场）}\quad
\bB(r) &= \begin{cases}
-\tfrac12 n I \pi \wms R\,Y_1(\wms R)\,J_0(\wms r)\,\be_z, & r < R\\[4pt]
-\tfrac12 n I \pi \wms R\,J_1(\wms R)\,Y_0(\wms r)\,\be_z, & r > R
\end{cases} && \text{（Riemannian）}\\
\bB(r) &= \begin{cases}
n I\,\be_z, & r < R\\[4pt]
0, & r > R
\end{cases} && \text{（Lorentzian）}
\end{align}

\begin{align}
\text{长螺线管（线圈内总磁通）}\quad
\Phi &= -n I \pi^2 R^2 J_1(\wms R) Y_1(\wms R) && \text{（Riemannian）}\\
\Phi &= n I \pi R^2 && \text{（Lorentzian）}\\[6pt]
\text{长螺线管（电感，管长 $l$）}\quad
L &= -n^2 \pi^2 R^2 l\,J_1(\wms R) Y_1(\wms R) && \text{（Riemannian）}\\
L &= n^2 \pi R^2 l && \text{（Lorentzian）}
\end{align}

上表中我们给出了穿过螺线管的总磁通；它是磁场 $\bB$ 在垂直于轴的横截面上的面积分。

若通过螺线管的电流开始变化，磁场也会变化，于是据麦克斯韦-法拉第定律（Maxwell-Faraday Law）将产生一个电场，其旋度正比于磁场的时间变化率。再由开尔文-斯托克斯定理，电场沿任何包围该变化磁场的回路的积分，将正比于磁场变化率在该回路所围面积上的积分。而该面积分正是穿过回路的磁通的时间变化率。因此每个包围着变化磁通的回路周围都会有一个正比于磁通变化率的电动势。事实上，比例常数恰好是 $-1$：
\[
\text{EMF} = -\frac{d\Phi}{dt}.
\]
把这一论证用于构成螺线管的各线圈，若通过螺线管的电流变化，则螺线管两端将产生正比于电流变化率的电压。其比例常数的相反数称为螺线管的\textbf{电感}（inductance）$L$：
\[
\text{EMF} = -L\,\frac{dI}{dt}.
\]
长度为 $l$（因而共有 $nl$ 匝线圈）的螺线管的黎曼电感为
\begin{equation}
L_{\Riem} = n l\,\frac{\Phi}{I} = -n^2 \pi^2 R^2 l\,J_1(\wms R) Y_1(\wms R),
\end{equation}
而洛伦兹值为
\begin{equation}
L_{\Lor} = n l\,\frac{\Phi}{I} = n^2 \pi R^2 l.
\end{equation}
黎曼电感中贝塞尔函数之积可正可负，因而电感可取两种符号。负电感如同负电容，会导致失控效应：负电感中电流的增大会产生驱动电流继续增大的电压，直到材料损坏或其他效应给电流的增长踩下刹车。

但与电容器一样，我们此处的模型高度理想化。正、负电感在线圈几何上的差别约为光的一个最小波长，因此若线圈导线比这更粗、或偏离完美圆形超过这一距离，螺线管实际就将同时包含正、负电感——通常会在两者之间造成显著抵消。

此外，我们此处的全部公式都假定了一种可近似为恒定电流的情形。携带交流电流的螺线管我们将在后面一节处理。

\section{电磁能流（Electromagnetic Energy Flow）}

我们在负电容或负电感系统中见到的那类失控效应，在我们自己的宇宙中显然会违反能量守恒；但在黎曼宇宙中，由于与物质（包括电磁场）相联系的能量与动能、势能意义相反，要精确弄清究竟发生了什么就更为棘手。我们需要能定量描述电磁场中储存与输运的能量。为此，我们得先绕道讨论黎曼电磁学的拉格朗日处理。

\subsection{黎曼电磁学的拉格朗日量（Lagrangian for Riemannian Electromagnetism）}

像电磁学这样的场论的拉格朗日量（Lagrangian）是一个依赖于场及其导数的量 $L$，当场满足适当方程时，它在四维空间某一区域上的积分在变分下取驻值。若对 $L$ 积分得到所谓的作用量（action）$S$：
\[
S(A_k) = \int L(A_k),
\]
则当 $\bA$ 满足场方程时，$S$ 应在 \emph{一阶}上不因 $\bA$ 的任何微小变化而改变，就像一元函数在局部极大或极小处那样。

若把拉格朗日量表示为场分量 $A_k$ 及其导数 $\partial_j A_k$ 的函数，则——只要场在积分区域边界上为零、或满足循环边界条件——作用量取驻值的要求等价于欧拉-拉格朗日方程（Euler-Lagrange equation）：
\[
\partial_j\!\left(\frac{\partial L}{\partial(\partial_j A_k)}\right) = \frac{\partial L}{\partial A_k}.
\]
我们把黎曼普罗卡拉格朗日量（Riemannian Proca Lagrangian）$L_{\text{RP}}$ 分为两部分：场拉格朗日量 $L_{\text{field}}$ 与相互作用项 $L_{\text{inter}}$。下面同时给出洛伦兹对应物\footnote{脚注 [1] 见原文参考文献。}。

\begin{align}
\text{黎曼普罗卡拉格朗日量}\quad
L_{\text{field}} &= \tfrac14 F_{ij} F^{ij} - \tfrac12 \wms^2 A_a A^a \nonumber\\
 &= \tfrac12\big(|\bB|^2 + |\bE|^2\big) - \tfrac12 \wms^2\big(|\bA_{(3)}|^2 + \phi^2\big)\\
L_{\text{inter}} &= -A_k j^k = -\bA_{(3)}\cdot\bj_{(3)} + \phi\rho\\
L_{\text{RP}} &= L_{\text{field}} + L_{\text{inter}} && \text{（Riemannian）}
\end{align}

\begin{align}
\text{麦克斯韦拉格朗日量}\quad
L_{\text{field}} &= -\tfrac14 F_{ij} F^{ij} = -\tfrac12\big(|\bB|^2 - |\bE|^2\big)\\
L_{\text{inter}} &= A_k j^k = \bA_{(3)}\cdot\bj_{(3)} - \phi\rho\\
L_{\text{Maxwell}} &= L_{\text{field}} + L_{\text{inter}} && \text{（Lorentzian）}
\end{align}

完整拉格朗日量的欧拉-拉格朗日方程，分别给出黎曼普罗卡方程或麦克斯韦方程。

\subsection{电磁学的应力-能量张量（The Stress-Energy Tensor for Electromagnetism）}

我们可以借助下式求出黎曼电磁场的应力-能量张量（stress-energy tensor），记作 $T$\footnote{脚注 [2] 见原文参考文献。}：

\begin{align}
\text{由场拉格朗日量给出的应力-能量张量}\quad
T_{ab} &= -L_{\text{field}}\,g_{ab} + 2\,\partial_{g^{ab}} L_{\text{field}} && \text{（Riemannian）}\\
T_{ab} &= L_{\text{field}}\,g_{ab} - 2\,\partial_{g^{ab}} L_{\text{field}} && \text{（Lorentzian）}
\end{align}

这里 $g_{ab}$ 与 $g^{ab}$ 是四维空间\textbf{度规张量}（metric tensor）的分量，分别带两个下标或两个上标。在正交归一坐标中，这些分量的矩阵就是 $4\times4$ 单位矩阵——即 $a=b$ 时为 1，否则为 0。但若把四维势场对偶矢量形式的分量 $A_k$ 视为拉格朗日量的基本变量，则每当我们升指标以得到像 $A_a A^a$ 这样的项时，都用到了 $g^{ab}$（采用爱因斯坦求和约定）：
\[
A_a A^a = A_a\big(g^{ab} A_b\big).
\]
因此，若把拉格朗日量视为四维势分量 $A_k$ \emph{与}度规张量分量 $g^{ab}$ 的函数，则对度规求导（在\emph{实际}度规处取值）就给出应力-能量张量中的第二项。

详细解释这一构造为何成立会偏离主题太远，但它最终与爱因斯坦引力方程——把一个由度规导出的张量与任何存在物质的应力-能量张量联系起来——本身可由适当拉格朗日量导出的方式相合。关键点是：如此构造出的完整应力-能量张量（包含全部物质）的散度为零，这意味着能量与动量守恒。

我们将用电磁场 $F$ 与四维势 $\bA$ 来表达这一计算结果，同时也用三维场 $\bB$、$\bE$、$\phi$ 与 $\bA_{(3)}$ 来表达。

\begin{equation}
\begin{gathered}
\textbf{[黎曼电磁应力-能量张量]}\quad
T_{ab} = -L_{\text{field}}\,g_{ab} + F_{ac} F_b^{\ c} - \wms^2 A_a A_b\\
\resizebox{0.98\textwidth}{!}{$\displaystyle T_{ab} = \begin{pmatrix}
\mathcal{A} & B_y E_z - B_z E_y + \wms^2 A_x\phi & B_z E_x - B_x E_z + \wms^2 A_y\phi & B_x E_y - B_y E_x + \wms^2 A_z\phi\\[4pt]
B_y E_z - B_z E_y + \wms^2 A_x\phi & \mathcal{B} + E_x^2 - B_x^2 - \wms^2 A_x^2 & E_x E_y - B_x B_y - \wms^2 A_x A_y & E_x E_z - B_x B_z - \wms^2 A_x A_z\\[4pt]
B_z E_x - B_x E_z + \wms^2 A_y\phi & E_x E_y - B_x B_y - \wms^2 A_x A_y & \mathcal{B} + E_y^2 - B_y^2 - \wms^2 A_y^2 & E_y E_z - B_y B_z - \wms^2 A_y A_z\\[4pt]
B_x E_y - B_y E_x + \wms^2 A_z\phi & E_x E_z - B_x B_z - \wms^2 A_x A_z & E_y E_z - B_y B_z - \wms^2 A_y A_z & \mathcal{B} + E_z^2 - B_z^2 - \wms^2 A_z^2
\end{pmatrix}$}
\end{gathered}
\end{equation}
其中
\[
\mathcal{A} = \frac{|\bE|^2 - |\bB|^2 + \wms^2\big(|\bA_{(3)}|^2 - \phi^2\big)}{2},
\qquad
\mathcal{B} = \frac{|\bB|^2 - |\bE|^2 + \wms^2\big(|\bA_{(3)}|^2 + \phi^2\big)}{2}.
\]

\begin{equation}
\begin{gathered}
\textbf{[洛伦兹电磁应力-能量张量]}\quad
T_{ab} = L_{\text{field}}\,g_{ab} + F_{ac} F_b^{\ c}\\
\resizebox{0.98\textwidth}{!}{$\displaystyle T_{ab} = \begin{pmatrix}
\dfrac{|\bE|^2 + |\bB|^2}{2} & B_y E_z - B_z E_y & B_z E_x - B_x E_z & B_x E_y - B_y E_x\\[6pt]
B_y E_z - B_z E_y & \mathcal{C} - B_x^2 - E_x^2 & -E_x E_y - B_x B_y & -E_x E_z - B_x B_z\\[4pt]
B_z E_x - B_x E_z & -E_x E_y - B_x B_y & \mathcal{C} - B_y^2 - E_y^2 & -E_y E_z - B_y B_z\\[4pt]
B_x E_y - B_y E_x & -E_x E_z - B_x B_z & -E_y E_z - B_y B_z & \mathcal{C} - B_z^2 - E_z^2
\end{pmatrix}$}
\end{gathered}
\end{equation}
其中 $\mathcal{C} = \dfrac{|\bE|^2 + |\bB|^2}{2}$。

电磁场单独的 $T$ 在 $\bj$ 不为零时其散度并不为零，而是有：
\[
\partial_b T^a{}^b + F^a{}_c\,j^c = 0.
\]
第二项对应于作用在电流上的四维力密度，而后者又将是带电物质自身应力-能量张量的散度。因此，电磁场的应力-能量张量与它作用的物质的应力-能量张量之和将为零。

\subsection{能量密度与坡印廷矢量（Energy Density and the Poynting Vector）}

应力-能量张量看起来可能有点吓人，但眼下我们先忽略第一行与第一列之外的各项，它们描述压强与切应力。我们关心的项是 $T_{tt}$，它给出电磁场中的能量密度 $u$；以及矢量 $\mathbf{S} = (T^{tx}, T^{ty}, T^{tz})$，称为\textbf{坡印廷矢量}（Poynting vector），它描述单位面积上的能流率。（注意我们必须升一个 $t$ 指标才能得到坡印廷矢量，这在洛伦兹情形下会改变符号。）

\begin{align}
\text{电磁能量密度}\quad
u &= \frac{|\bE|^2 - |\bB|^2 + \wms^2\big(|\bA_{(3)}|^2 - \phi^2\big)}{2} && \text{（Riemannian）}\\
u &= \frac{|\bE|^2 + |\bB|^2}{2} && \text{（Lorentzian）}\\[6pt]
\text{坡印廷矢量}\quad
\mathbf{S} &= \bB \times \bE + \wms^2 \phi\,\bA_{(3)} && \text{（Riemannian）}\\
\mathbf{S} &= \bE \times \bB && \text{（Lorentzian）}
\end{align}

下面我们看几个简单例子中的能量密度与能流。

\subsubsection{平面波中的能量（Energy in Plane Waves）}

对于平面波，我们在四维空间中有如下描述：
\[
\bA(\bx) = \bA_0 \sin(\bk\cdot\bx), \qquad
F(\bx) = (\bk \wedge \bA_0)\cos(\bk\cdot\bx),
\]
其中 $|\bk| = \wms$ 且 $\bA_0\cdot\bk = 0$。由此可计算应力-能量张量：
\begin{equation}
\begin{aligned}
T_{ab} &= L_{\text{field}}\,g_{ab} + F_{ac} F_b^{\ c} - \wms^2 A_a A_b\\
T &= A_0^2\,\bk\otimes\bk\,\cos(\bk\cdot\bx)^2 + \wms^2\Big[\bA_0\otimes\bA_0 - \frac{A_0^2}{2} I_4\Big]\cos(2\bk\cdot\bx)
\end{aligned}
\end{equation}
若把 $T$ 在一个周期上取平均，则 $\cos(2\bk\cdot\bx)$ 变为零，而 $\cos(\bk\cdot\bx)^2$ 变为 $1/2$，于是：
\[
\langle T\rangle = \tfrac12 A_0^2\,\bk\otimes\bk.
\]
这恰好是我们预期一团四维速度为 $\bu = \bk/\wms$、且在其静止系中质量-能量密度为 $\tfrac12 A_0^2\wms^2$ 的均匀物质云应有的应力-能量张量。若如此定义 $\bu$，并定义单位矢量 $\mathbf{a}_0 = \bA_0/A_0$，则应力-能量张量可写作：
\begin{equation}
T = A_0^2 \wms^2\Big[\bu\otimes\bu\,\cos(\bk\cdot\bx)^2 + \big(\mathbf{a}_0\otimes\mathbf{a}_0 - \tfrac12 I_4\big)\cos(2\bk\cdot\bx)\Big].
\end{equation}
设光具有时间角频率 $\omega = k_t = \wms u_t$。则能量密度 $u$（勿与四维速度 $\bu$ 或其任何分量混淆）为：
\begin{equation}
\begin{aligned}
u = T_{tt} &= A_0^2\Big[\omega^2\cos(\bk\cdot\bx)^2 + \wms^2\big(a_{0,t}^2 - \tfrac12\big)\cos(2\bk\cdot\bx)\Big]\\
&= \tfrac12 A_0^2\Big[\omega^2 + \Big(\omega^2 + \big(2a_{0,t}^2 - 1\big)\wms^2\Big)\cos(2\bk\cdot\bx)\Big]
\end{aligned}
\end{equation}
显然存在使能量密度在某些时刻为负的 $\omega$ 与 $a_{0,t}$ 取值：例如，若 $a_{0,t} = 0$ 且 $\omega < \wms/\sqrt{2}$。但任意周期上的平均能量密度仍为正：
\[
\langle u\rangle = \tfrac12 A_0^2 \omega^2.
\]
从 $\langle T\rangle$ 可以看出，坡印廷矢量 $\mathbf{S}$ 的同类平均将平行于传播矢量 $\bk$ 的空间投影，而后者又平行于与四维速度 $\bu = \bk/\wms$ 对应的普通速度 $\bv$。具体地：
\[
\langle \mathbf{S}\rangle = \tfrac12 A_0^2 \omega^2\,\bv.
\]

\subsubsection{电容器中的能量（Energy in Capacitors）}

我们可以把电场能量密度公式用于此前分析过的球形电容器。在洛伦兹情形下，电容器外的电场为零，能量密度只依赖于场，因此可直接积分得到有限结果。

在黎曼情形下，情况稍棘手。势与电场都延伸到电容器之外，由它们算出的能量密度直到无穷远都非零。给定半径 $S \gg R_2$ 的球内所含能量关于 $S$ 是周期性的，且这些周期的峰-峰间距不随距离减小，因此到无穷远的积分未定义。但我们可以把周期部分置零、取其剩余部分的渐近值，从而得到一个合理的有限结果。

\[
\text{球形电容器：内壳半径 } R_1\text{，总电量 } -Q;\quad \text{外壳半径 } R_2\text{，总电量 } +Q.
\]

球形电容器（电容）：
\begin{align}
C &= \frac{8 \pi \wms R_1^2 R_2^2}{4 R_1 R_2 \sin(\wms R_1)\cos(\wms R_2) - R_1^2 \sin(2\wms R_2) - R_2^2 \sin(2\wms R_1)} && \text{（Riemannian）}\\
C &= \frac{4 \pi R_1 R_2}{R_2 - R_1} && \text{（Lorentzian）}
\end{align}

\begin{align}
\text{球形电容器（电场能量密度）}\quad
u(r) &= \frac{|\bE(r)|^2 - \wms^2 \phi(r)^2}{2} \quad \text{（见电容器的场计算）} && \text{（Riemannian）}\\
u(r) &= \begin{cases}
0, & r < R_1\\[4pt]
\dfrac{Q^2}{32 \pi^2 r^4}, & R_1 < r < R_2\\[6pt]
0, & r > R_2
\end{cases} && \text{（Lorentzian）}
\end{align}

球形电容器（电场总能量）：
\begin{align}
\langle U\rangle &= \int_0^{R_2} 4 \pi r^2 u(r)\,dr + \int_{R_2}^{S} 4 \pi r^2 u(r)\,dr \quad (\text{对 } S \gg R_2 \text{ 取平均}) \nonumber\\
&= \frac{Q^2}{16 \pi \wms R_1^2 R_2^2}\Big[R_1^2 \sin(2\wms R_2) + R_2^2 \sin(2\wms R_1) - 4 R_1 R_2 \sin(\wms R_1)\cos(\wms R_2)\Big] \nonumber\\
&= -\frac{Q^2}{2C} && \text{（Riemannian）}\\
U &= \int_{R_1}^{R_2} 4 \pi r^2 u(r)\,dr = \frac{Q^2}{8\pi}\Big[\frac{1}{R_1} - \frac{1}{R_2}\Big] = \frac{Q^2}{2C} && \text{（Lorentzian）}
\end{align}

我们在黎曼与洛伦兹两种情形下得到的结果，都与预期的电容器势能相符——若把电容器从零电荷充到总电量 $Q$ 所需的能量积分起来：
\[
\text{势能} = \int_0^Q V(q)\,dq = \int_0^Q \frac{q}{C}\,dq = \frac{Q^2}{2C}.
\]
在洛伦兹情形下，这恰好是储存在电场中的能量。而在黎曼情形下，它恰恰相反！原因当然在于黎曼宇宙中势能与电磁场能意义相反。

\subsubsection{电感器中的能量（Energy in Inductors）}

螺线管中储存能量的计算与电容器遵循同一总体模式。在洛伦兹情形下，有限体积内存在恒定磁场，场的总能量极易计算。

在黎曼情形下，我们不能忽略螺线管外的场，且无穷区域上的积分不收敛；但若积分到半径 $S$，所包围的总能量在极大值与极小值之间周期性变化，且在 $S$ 很大时趋近于固定值。下表中我们用余弦函数给出 $S$ 的贝塞尔函数之积的渐近表达式。对该余弦项取一个周期上的平均值（只需把该项置零即可轻易得到），所得结果便与由电感得到的能量相符。

\begin{align}
\text{长螺线管（电感）}\quad
L &= -n^2 \pi^2 R^2 l\,J_1(\wms R) Y_1(\wms R) && \text{（Riemannian）}\\
L &= n^2 \pi R^2 l && \text{（Lorentzian）}
\end{align}

长螺线管（磁场能量密度）：
\begin{align}
u(r) &= \begin{cases}
\tfrac18 n^2 I^2 \pi^2 R^2 \wms^2 Y_1(\wms R)^2\big(J_1(\wms r)^2 - J_0(\wms r)^2\big), & r < R\\[4pt]
\tfrac18 n^2 I^2 \pi^2 R^2 \wms^2 J_1(\wms R)^2\big(Y_1(\wms r)^2 - Y_0(\wms r)^2\big), & r > R
\end{cases} && \text{（Riemannian）}\\
u(r) &= \begin{cases}
\tfrac12 n^2 I^2, & r < R\\[4pt]
0, & r > R
\end{cases} && \text{（Lorentzian）}
\end{align}

长螺线管（磁场总能量）：
\begin{align}
\langle U\rangle &= 2 \pi l\Big[\int_0^{R} u(r)\,r\,dr + \int_{R}^{S} u(r)\,r\,dr\Big] \quad (\text{对 } S \gg R \text{ 取平均}) \nonumber\\
&= -\tfrac14 n^2 I^2 \pi^3 R^3 l \wms\Big[Y_1(\wms R)^2 J_0(\wms R) J_1(\wms R) \nonumber\\
 &\qquad - J_1(\wms R)^2\big[Y_0(\wms R) Y_1(\wms R) - (S/R)\,Y_0(\wms S) Y_1(\wms S)\big]\Big] \nonumber\\
&\approx -\tfrac14 n^2 I^2 \pi^3 R^3 l \wms\Big[Y_1(\wms R)^2 J_0(\wms R) J_1(\wms R) \nonumber\\
 &\qquad - J_1(\wms R)^2\big[Y_0(\wms R) Y_1(\wms R) - \frac{\cos(2\wms S)}{\pi \wms R}\big]\Big] \nonumber\\
&= \tfrac12 n^2 I^2 \pi^2 R^2 l\,J_1(\wms R) Y_1(\wms R) = -\tfrac12 L I^2 && \text{（Riemannian）}\\
U &= \pi R^2 l\,u(0) = \tfrac12 n^2 I^2 \pi R^2 l = \tfrac12 L I^2 && \text{（Lorentzian）}
\end{align}

对电感器而言，势能可通过把电流从零升至某最终恒定值 $I$ 所需做的功求得。当我们在时间 $dt$ 内把电流从 $i$ 变为 $i + di$ 时，我们让电量 $i\,dt$ 逆着电压 $V = L\,di/dt$ 移动。于是：
\[
\text{势能} = \int_0^I V(t)\,i\,dt = \int_0^I L\,\frac{di}{dt}\,i\,dt = \tfrac12 L I^2.
\]
如我们所料，这样算出的势能在洛伦兹情形下与磁场中的总能量相符，但在黎曼情形下与磁场中的能量相反。

\section{由静磁解导出的振荡解（Oscillating Solutions Derived From Magnetostatic Ones）}

设我们有一个黎曼普罗卡方程的静磁解，四维势为 $\bA_{\text{MS}}$、源四维电流为 $\bj_{\text{MS}}$。我们说“静磁”是指 $\bA_{\text{MS}}$ 与 $\bj_{\text{MS}}$ 都不随时间变化，且场纯粹是磁性的，即 $A_{\text{MS}}^t = 0$。我们已经考察过三个这样的解：恒定线电流、磁偶极子、以及带恒定电流的螺线管。

现在假设我们取该解，并把 $\wms$——黎曼光的最大角频率——替换为一个更小的值 $k$，从而得到 $\bA_{\text{MS},k}$ 与 $\bj_{\text{MS},k}$，它们满足方程
\[
\partial_x^2 A_{\text{MS},k} + \partial_y^2 A_{\text{MS},k} + \partial_z^2 A_{\text{MS},k} + k^2 A_{\text{MS},k} + j_{\text{MS},k} = 0.
\]
然后我们构造一个振荡解：
\[
\bA = \bA_{\text{MS},k}\,\cos(\omega t), \qquad \bj = \bj_{\text{MS},k}\,\cos(\omega t),
\]
其时间角频率为 $\omega$，满足
\begin{equation}
k^2 + \omega^2 = \wms^2.
\end{equation}
新的 $\bA$ 与 $\bj$ 将满足 RVWS 方程：
\begin{equation}
\begin{aligned}
&\partial_x^2 \bA + \partial_y^2 \bA + \partial_z^2 \bA + \partial_t^2 \bA + \wms^2 \bA + \bj\\
&\quad = \cos(\omega t)\Big[\partial_x^2 A_{\text{MS},k} + \partial_y^2 A_{\text{MS},k} + \partial_z^2 A_{\text{MS},k} - \omega^2 A_{\text{MS},k} + \wms^2 A_{\text{MS},k} + j_{\text{MS},k}\Big]\\
&\quad = \cos(\omega t)\Big[\partial_x^2 A_{\text{MS},k} + \partial_y^2 A_{\text{MS},k} + \partial_z^2 A_{\text{MS},k} + k^2 A_{\text{MS},k} + j_{\text{MS},k}\Big] = 0.
\end{aligned}
\end{equation}
横波条件又如何呢？我们的静磁解满足该条件，且没有时间分量：
\[
\partial_x A_{\text{MS},k}^x + \partial_y A_{\text{MS},k}^y + \partial_z A_{\text{MS},k}^z = 0.
\]
把 $\bA_{\text{MS},k}$ 乘以 $\cos(\omega t)$ 后这依然成立，且当然仍有 $A^t = 0$。因此我们的新振荡解确为黎曼普罗卡方程的真正解。

在上述一切中，我们同样可以用 $\sin(\omega t)$ 而非 $\cos(\omega t)$。我们在此构造中用的是 $t$ 方向，还是四维空间中任何其他既使解不变、又使其四维势分量为零的方向，也都没有区别。

我们可以用非常类似的过程，从原来的黎曼静磁解得到同类的振荡\emph{洛伦兹}解。在洛伦兹电磁学中，四维势不出现在麦克斯韦方程中，其唯一的物理意义通过电磁场 $F$ 体现。但不同的四维势 $\bA$ 可以给出完全相同的 $F$，因此我们可以在不改变物理的前提下对 $\bA$ 作某些种类的改变；这称为\textbf{规范自由度}（gauge freedom）。处理规范自由度的一种方便做法，是给 $\bA$ 额外施加一个它必须满足的条件，而在不同情境下有多种选择可使计算更简便。其中一种选择称为洛伦兹规范条件（Lorenz gauge condition）——是“Lorenz”而非“Lorentz”，这是两个完全不同的人！——它要求：
\[
\partial_x A^x + \partial_y A^y + \partial_z A^z + \partial_t A^t = 0.
\]
这是我们施加于每个黎曼矢量波的横波条件的洛伦兹版本。因此，若我们在洛伦兹规范下做洛伦兹电磁学，两类电磁学之间的联系就变得清晰得多；在那里麦克斯韦方程等价于下列关于四维势的方程：
\begin{align}
\partial_x^2 \bA + \partial_y^2 \bA + \partial_z^2 \bA - \partial_t^2 \bA + \bj &= \bzero \tag{LVWS}\\
\partial_x A^x + \partial_y A^y + \partial_z A^z + \partial_t A^t &= 0 \tag{Lorenz}
\end{align}
若取我们原来的黎曼静磁解 $\bA_{\text{MS}}$（对应四维电流 $\bj_{\text{MS}}$），可如下得到振荡洛伦兹解。我们把 $\wms$ 替换为任意频率 $\omega$，得到 $\bA_{\text{MS},\omega}$ 与 $\bj_{\text{MS},\omega}$，再把它们乘以 $\cos(\omega t)$：
\[
\bA_L = \bA_{\text{MS},\omega}\,\cos(\omega t), \qquad \bj_L = \bj_{\text{MS},\omega}\,\cos(\omega t).
\]
这些函数将满足带源的洛伦兹矢量波方程 (LVWS)：
\[
\begin{aligned}
&\partial_x^2 \bA_L + \partial_y^2 \bA_L + \partial_z^2 \bA_L - \partial_t^2 \bA_L + \bj_L\\
&\quad = \cos(\omega t)\Big[\partial_x^2 A_{\text{MS},\omega} + \partial_y^2 A_{\text{MS},\omega} + \partial_z^2 A_{\text{MS},\omega} + \omega^2 A_{\text{MS},\omega} + j_{\text{MS},\omega}\Big] = 0.
\end{aligned}
\]
由于两个四维势都没有时间分量，$\bA_{\text{MS},\omega}$ 满足横波条件这一点，足以让 $\bA_L$ 满足洛伦兹条件。

\subsection{线形交流电流（Linear Alternating Current）}

把刚才描述的方法用于线性导体中恒定电流的四维势，就得到一根载有交流电流的线性导体周围振荡驻波场的解。

线形交流驻波解（电流 $I_0\cos(\omega t)$ 沿 $z$ 轴流动；黎曼解中 $k^2 + \omega^2 = \wms^2$）：
\begin{align}
\bA_{(3)}(r) &= -\frac{I_0}{4}\,Y_0(kr)\cos(\omega t)\,\be_z && \text{（Riemannian）}\\
\bA_{(3)}(r) &= -\frac{I_0}{4}\,Y_0(\omega r)\cos(\omega t)\,\be_z && \text{（Lorentzian）}
\end{align}

\begin{align}
\text{线形交流·磁场与电场}\quad
\bB(r) &= -\frac{I_0 k}{4}\,Y_1(kr)\cos(\omega t)\,\be_\phi\\
\bE(r) &= -\frac{I_0 \omega}{4}\,Y_0(kr)\sin(\omega t)\,\be_z && \text{（Riemannian）}\\
\bB(r) &= -\frac{I_0 \omega}{4}\,Y_1(\omega r)\cos(\omega t)\,\be_\phi\\
\bE(r) &= -\frac{I_0 \omega}{4}\,Y_0(\omega r)\sin(\omega t)\,\be_z && \text{（Lorentzian）}
\end{align}

驻波解在空间中形式固定，只在时间上振荡。这正是导线处于圆柱形腔中时我们预期的波。但若我们想要行波解呢？驻波可以写成入射行波与出射行波之和或差，反之入射波与出射波也可由这些驻波之和或差还原；因此只要我们能找到第二个驻波解，就应能构造出行波。

对于第二个驻波解，我们回到最初对线电流的计算，使用那个完全与电流强度无关的无源解。这相当于把上述势中的贝塞尔函数 $Y_0$ 换成 $J_0$。若再把新解相对原解移相 90 度，即把 $\cos(\omega t)$ 因子换成 $\sin(\omega t)$，然后把两个解相加，就得到一个出射行波。由于我们加上的第二个解是无源的，无需改动电流；这只是同一根导线、同一电流在不同边界条件下的波。

在洛伦兹情形下，我们需要\emph{相减}而非相加这第二个解，才能得到出射波。

线形交流出射行波解（黎曼解中 $k^2 + \omega^2 = \wms^2$）：
\begin{align}
\bA_{(3)}(r) &= -\frac{I_0}{4}\big[Y_0(kr)\cos(\omega t) + J_0(kr)\sin(\omega t)\big]\be_z && \text{（Riemannian）}\\
\bA_{(3)}(r) &= -\frac{I_0}{4}\big[Y_0(\omega r)\cos(\omega t) - J_0(\omega r)\sin(\omega t)\big]\be_z && \text{（Lorentzian）}
\end{align}

\begin{align}
\text{线形交流·磁场与电场}\quad
\bB(r) &= -\frac{I_0 k}{4}\big[Y_1(kr)\cos(\omega t) + J_1(kr)\sin(\omega t)\big]\be_\phi\\
\bE(r) &= -\frac{I_0 \omega}{4}\big[Y_0(kr)\sin(\omega t) - J_0(kr)\cos(\omega t)\big]\be_z && \text{（Riemannian）}\\
\bB(r) &= -\frac{I_0 \omega}{4}\big[Y_1(\omega r)\cos(\omega t) - J_1(\omega r)\sin(\omega t)\big]\be_\phi\\
\bE(r) &= -\frac{I_0 \omega}{4}\big[Y_0(\omega r)\sin(\omega t) + J_0(\omega r)\cos(\omega t)\big]\be_z && \text{（Lorentzian）}
\end{align}

线形交流的坡印廷矢量：
\begin{align}
\mathbf{S}(r) &= \frac{I_0^2 k \omega}{16}\big[Y_1(kr)\cos(\omega t) + J_1(kr)\sin(\omega t)\big]\nonumber\\
&\qquad\times\big[Y_0(kr)\sin(\omega t) - J_0(kr)\cos(\omega t)\big]\be_r && \text{（Riemannian）}\\
\mathbf{S}(r) &= -\frac{I_0^2 \omega^2}{16}\big[Y_1(\omega r)\cos(\omega t) - J_1(\omega r)\sin(\omega t)\big]\nonumber\\
&\qquad\times\big[Y_0(\omega r)\sin(\omega t) + J_0(\omega r)\cos(\omega t)\big]\be_r && \text{（Lorentzian）}\\
\langle \mathbf{S}(r)\rangle &= \frac{I_0^2 \omega}{16 \pi r}\,\be_r && \text{（Common）}
\end{align}

\begin{equation}
\text{线形交流·平均辐射功率（每单位导线长度）}\quad \langle P\rangle = \frac{I_0^2 \omega}{8} \tag{Common}
\end{equation}
从 $\langle \mathbf{S}(r)\rangle$——即在一个时间周期上取平均的坡印廷矢量——我们能最清楚地看出这些是出射行波：它是一个显然为正的数值乘以从导线径向向外的单位矢量。

在黎曼情形下，空间角频率 $k$ 从最终结果中消失，这乍看或许有点令人困惑；毕竟我们预期这些波的速度为 $k/\omega$，而能量流密度应为该速度乘以能量密度。但事实证明，能量密度反比于 $k$，这从大 $r$ 处四维势的形式不难看出——由于贝塞尔函数的渐近展开，它反比于 $k$ 的平方根：
\[
\bA(r) \approx \frac{I_0 \cos(kr + \omega t + \pi/4)}{2\sqrt{2 \pi k r}}\,\be_z.
\]
于是我们此前对平面波能流的分析，会由这列平面波给出与精确解导出的相同的平均坡印廷矢量。

在洛伦兹情形下，辐射功率意味着必须做功以维持电流幅度不变。在黎曼情形下，常规意义上的功必须\emph{由}电流来做，以阻止它变得更大！这虽然奇怪，但考虑到电磁场中的能量与动能、势能意义相反，这正是我们应当预期的。

\subsection{振荡磁偶极子（Oscillating Magnetic Dipoles）}

我们将用同样的方法，基于此前对静态偶极子的结果，构造振荡磁偶极子的场。我们不再展示两个驻波解，而是直接跳到出射行波。

振荡磁偶极子的出射行波解（偶极矩 $\boldsymbol\mu\cos(\omega t)$；黎曼解中 $k^2 + \omega^2 = \wms^2$）：
\begin{align}
\bA_{(3)}(\mathbf{r}) &= \frac{\boldsymbol\mu\times\mathbf{r}}{4\pi r^3}\big[\cos(kr + \omega t) + kr\sin(kr + \omega t)\big] && \text{（Riemannian）}\\
\bA_{(3)}(\mathbf{r}) &= \frac{\boldsymbol\mu\times\mathbf{r}}{4\pi r^3}\big[\cos(\omega(r-t)) + \omega r\sin(\omega(r-t))\big] && \text{（Lorentzian）}
\end{align}

\begin{align}
\text{振荡磁偶极子·磁场与电场}\quad
\bB(\mathbf{r}) &= \frac{3(\boldsymbol\mu\cdot\mathbf{r})\mathbf{r} - r^2\boldsymbol\mu}{4\pi r^5}\big[\cos(kr + \omega t) + kr\sin(kr + \omega t)\big] \nonumber\\
 &\qquad - \frac{k^2(\boldsymbol\mu\times\mathbf{r})\times\mathbf{r}}{4\pi r^3}\cos(kr + \omega t)\\
\bE(\mathbf{r}) &= \frac{\omega\,\boldsymbol\mu\times\mathbf{r}}{4\pi r^3}\big[\sin(kr + \omega t) - kr\cos(kr + \omega t)\big] && \text{（Riemannian）}\\
\bB(\mathbf{r}) &= \frac{3(\boldsymbol\mu\cdot\mathbf{r})\mathbf{r} - r^2\boldsymbol\mu}{4\pi r^5}\big[\cos(\omega(r-t)) + \omega r\sin(\omega(r-t))\big] \nonumber\\
 &\qquad - \frac{\omega^2(\boldsymbol\mu\times\mathbf{r})\times\mathbf{r}}{4\pi r^3}\cos(\omega(r-t))\\
\bE(\mathbf{r}) &= \frac{\omega\,\boldsymbol\mu\times\mathbf{r}}{4\pi r^3}\big[\omega r\cos(\omega(r-t)) - \sin(\omega(r-t))\big] && \text{（Lorentzian）}
\end{align}

\begin{align}
\text{振荡磁偶极子·一个周期上平均的坡印廷矢量}\quad
\langle \mathbf{S}\rangle &= \frac{k^3 \omega\,\big((\boldsymbol\mu\cdot\boldsymbol\mu) - (\boldsymbol\mu\cdot\be_r)^2\big)}{32 \pi^2 r^2}\,\be_r && \text{（Riemannian）}\\
\langle \mathbf{S}\rangle &= \frac{\omega^4\big((\boldsymbol\mu\cdot\boldsymbol\mu) - (\boldsymbol\mu\cdot\be_r)^2\big)}{32 \pi^2 r^2}\,\be_r && \text{（Lorentzian）}\\[6pt]
\text{振荡磁偶极子·一个周期上平均的总功率}\quad
\langle P\rangle &= \frac{k^3 \omega\,(\boldsymbol\mu\cdot\boldsymbol\mu)}{12 \pi} && \text{（Riemannian）}\\
\langle P\rangle &= \frac{\omega^4\,(\boldsymbol\mu\cdot\boldsymbol\mu)}{12 \pi} && \text{（Lorentzian）}
\end{align}

\orthofig{images/04__014.png}{振荡磁偶极子的功率}

若看黎曼四维势在大 $r$ 处的渐近形式，我们有：
\[
\bA(r) \approx \frac{k\sin(kr + \omega t)}{4\pi r}\,\boldsymbol\mu\times\be_r.
\]
偏振始终是横向的，四维势环绕偶极子轴指向。其大小在垂直于偶极子方向最大，在轴上本身降为零。辐射功率的角分布与洛伦兹情形完全一致。

在黎曼情形下，一个周期上平均的能量密度正比于 $k^2 \omega^2$。乘以波速 $k/\omega$，就给出表中与右图所示功率的 $k^3 \omega$ 频率依赖关系。

\subsection{螺线管中的交流电流（Alternating Current in a Solenoid）}

我们可以用同样的方法，把带恒定电流螺线管的静磁描述改造成带交流电流者。为得到一个无源静磁解，我们把恒定电流螺线管解外部部分的因子 $Y_1$ 换成 $J_1$，并让同一解一直延续到 $z$ 轴。把两个驻波解组合起来，就得到出射行波解。

长螺线管（交流）的出射行波解（半径 $R$，电流 $I_0\cos(\omega t)$，单位长度 $n$ 匝；黎曼解中 $k^2 + \omega^2 = \wms^2$）：
\begin{align}
\bA_{(3)}(r) &= \begin{cases}
-\tfrac12 \pi I_0 n R\,J_1(kr)\big[J_1(kR)\sin(\omega t) + Y_1(kR)\cos(\omega t)\big]\be_\phi, & r < R\\[4pt]
-\tfrac12 \pi I_0 n R\,J_1(kR)\big[J_1(kr)\sin(\omega t) + Y_1(kr)\cos(\omega t)\big]\be_\phi, & r > R
\end{cases} && \text{（Riemannian）}\\
\bA_{(3)}(r) &= \begin{cases}
\tfrac12 \pi I_0 n R\,J_1(\omega r)\big[J_1(\omega R)\sin(\omega t) - Y_1(\omega R)\cos(\omega t)\big]\be_\phi, & r < R\\[4pt]
\tfrac12 \pi I_0 n R\,J_1(\omega R)\big[J_1(\omega r)\sin(\omega t) - Y_1(\omega r)\cos(\omega t)\big]\be_\phi, & r > R
\end{cases} && \text{（Lorentzian）}
\end{align}

长螺线管（交流）的磁场与电场：
\begin{align}
\bB(r) &= \begin{cases}
-\tfrac12 \pi k I_0 n R\,J_0(kr)\big[J_1(kR)\sin(\omega t) + Y_1(kR)\cos(\omega t)\big]\be_z, & r < R\\[4pt]
-\tfrac12 \pi k I_0 n R\,J_1(kR)\big[J_0(kr)\sin(\omega t) + Y_0(kr)\cos(\omega t)\big]\be_z, & r > R
\end{cases}\\
\bE(r) &= \begin{cases}
\tfrac12 \pi I_0 n \omega R\,J_1(kr)\big[J_1(kR)\cos(\omega t) - Y_1(kR)\sin(\omega t)\big]\be_\phi, & r < R\\[4pt]
\tfrac12 \pi I_0 n \omega R\,J_1(kR)\big[J_1(kr)\cos(\omega t) - Y_1(kr)\sin(\omega t)\big]\be_\phi, & r > R
\end{cases} && \text{（Riemannian）}\\
\bB(r) &= \begin{cases}
\tfrac12 \pi I_0 n \omega R\,J_0(\omega r)\big[J_1(\omega R)\sin(\omega t) - Y_1(\omega R)\cos(\omega t)\big]\be_z, & r < R\\[4pt]
\tfrac12 \pi I_0 n \omega R\,J_1(\omega R)\big[J_0(\omega r)\sin(\omega t) - Y_0(\omega r)\cos(\omega t)\big]\be_z, & r > R
\end{cases}\\
\bE(r) &= \begin{cases}
-\tfrac12 \pi I_0 n \omega R\,J_1(\omega r)\big[J_1(\omega R)\cos(\omega t) + Y_1(\omega R)\sin(\omega t)\big]\be_\phi, & r < R\\[4pt]
-\tfrac12 \pi I_0 n \omega R\,J_1(\omega R)\big[J_1(\omega r)\cos(\omega t) + Y_1(\omega r)\sin(\omega t)\big]\be_\phi, & r > R
\end{cases} && \text{（Lorentzian）}
\end{align}

长螺线管（交流）在一个周期上平均的坡印廷矢量：
\begin{align}
\langle \mathbf{S}(r)\rangle &= \begin{cases}
0, & r < R\\[4pt]
\dfrac{\pi I_0^2 n^2 R^2 \omega\,J_1(kR)^2}{4r}\,\be_r, & r > R
\end{cases} && \text{（Riemannian）}\\
\langle \mathbf{S}(r)\rangle &= \begin{cases}
0, & r < R\\[4pt]
\dfrac{\pi I_0^2 n^2 R^2 \omega\,J_1(\omega R)^2}{4r}\,\be_r, & r > R
\end{cases} && \text{（Lorentzian）}
\end{align}

长螺线管（交流）在一个周期上平均的总辐射功率（管长 $l$）：
\begin{align}
\langle P_{\text{Radiated}}\rangle &= \tfrac12 \pi^2 I_0^2 l n^2 R^2 \omega\,J_1(kR)^2 \nonumber\\
&\approx \tfrac18 \pi^2 I_0^2 l n^2 R^4 \wms k^2 \quad (\textbf{低 } k \textbf{ 极限}) && \text{（Riemannian）}\\
\langle P_{\text{Radiated}}\rangle &= \tfrac12 \pi^2 I_0^2 l n^2 R^2 \omega\,J_1(\omega R)^2 \nonumber\\
 &\approx \tfrac18 \pi^2 I_0^2 l n^2 R^4 \omega^3 \quad (\textbf{低 } \omega \textbf{ 极限}) && \text{（Lorentzian）}
\end{align}

长螺线管（交流）线圈内的总磁通：
\begin{align}
\Phi &= -\pi^2 I_0 n R^2 J_1(kR)\big[J_1(kR)\sin(\omega t) + Y_1(kR)\cos(\omega t)\big] && \text{（Riemannian）}\\
\Phi &= \pi^2 I_0 n R^2 J_1(\omega R)\big[J_1(\omega R)\sin(\omega t) - Y_1(\omega R)\cos(\omega t)\big] && \text{（Lorentzian）}
\end{align}

长螺线管（交流）长度 $l$ 的线圈上的电压：
\begin{align}
V &= \pi^2 I_0 l n^2 R^2 \omega\,J_1(kR)\big[Y_1(kR)\sin(\omega t) - J_1(kR)\cos(\omega t)\big] \nonumber\\
&\approx -\pi I_0 l n^2 R^2 \wms \sin(\omega t) \quad (\textbf{低 } k \textbf{ 极限}) && \text{（Riemannian）}\\
V &= \pi^2 I_0 l n^2 R^2 \omega\,J_1(\omega R)\big[Y_1(\omega R)\sin(\omega t) + J_1(\omega R)\cos(\omega t)\big] \nonumber\\
 &\approx -\pi I_0 l n^2 R^2 \omega \sin(\omega t) \quad (\textbf{低 } \omega \textbf{ 极限}) && \text{（Lorentzian）}
\end{align}

长螺线管（交流）长度 $l$ 的线圈上平均消耗的电功率：
\begin{align}
\langle P_{\text{Radiated}}\rangle &= -\tfrac12 \pi^2 I_0^2 l n^2 R^2 \omega\,J_1(kR)^2 \nonumber\\
&\approx -\tfrac18 \pi^2 I_0^2 l n^2 R^4 \wms k^2 \quad (\textbf{低 } k \textbf{ 极限}) && \text{（Riemannian）}\\
\langle P_{\text{Radiated}}\rangle &= \tfrac12 \pi^2 I_0^2 l n^2 R^2 \omega\,J_1(\omega R)^2 \nonumber\\
 &\approx \tfrac18 \pi^2 I_0^2 l n^2 R^4 \omega^3 \quad (\textbf{低 } \omega \textbf{ 极限}) && \text{（Lorentzian）}
\end{align}

黎曼解的第一个有趣特征在于：如今是空间角频率 $k$ 而非恒定电流时的 $\wms$ 决定了螺线管几何的尺度。螺线管的直流行为会对任何可与光的最小波长相比拟的瑕疵极端敏感——而一个现实的器件，其导线宽度可能跨越若干波长，以致整个结构包含一系列几乎彼此抵消的负电感与正电感——而如今我们有可能获得大得多的波长，从而得到一个既无抵消、又对线圈精确形状不那么敏感的系统。

\orthofig{images/04__015.png}{交流螺线管的电流、电压与功率}

我们在处理直流螺线管时指出，它可以具有正电感或负电感，因而可以\emph{要么}阻碍、\emph{要么}帮助电流变化。然而在交流语境下这一区分不那么重要；重要的是一个完整周期上消耗的功率，而我们选择出射波这一点保证了黎曼螺线管将充当电功率的\emph{源}，而洛伦兹对应物则需要消耗功率。

右图给出三种尺寸线圈下黎曼与洛伦兹情形的电流、电压与功率。这里 $J_1$、$Y_1$ 是缩写，在黎曼情形下代表 $J_1(kR)$、$Y_1(kR)$，在洛伦兹情形下代表 $J_1(\omega R)$、$Y_1(\omega R)$。我们此处对电压所用的符号约定是：普通电阻器上的电压恰与电流同相，因而按 $VI$ 之积算出的功率是正在耗散的电能。

在洛伦兹情形下，电压与电流的相位差从不超过 90 度。在低频直流极限下，$J_1(\omega R) Y_1(\omega R) \approx -1/\pi$（到一阶），电压恰好超前电流 90 度。若考虑一个至少几毫米大小、携带千赫兹或更低频交流的电感器，其波长远大于螺线管尺寸，因此那个“极限情形”实际上对许多常见交流电路都是良好近似。不过随着频率升高，$Y_1(\omega R)$ 最终变为零，使电压与电流同相，继而变为正，于是电压滞后于电流。但无论 $J_1(\omega R)$ 与 $Y_1(\omega R)$ 取何值，每个周期平均耗散的功率总是非负。

在黎曼情形下，电压与电流的相位差从不\emph{小于} 90 度。直流极限涉及最大可能的空间频率，以及高度依赖线圈几何的行为。只有在高（时间）频率极限下波长才变大，此时 $J_1(kR) Y_1(kR) \approx -1/\pi$，电压超前电流 90 度。但在所有频率与线圈尺寸下，平均耗散的功率都是非正——因为在黎曼情形下，任何辐射出去的能量都必须伴随常规能量的增加。

\section{由静电解导出的振荡解（Oscillating Solutions Derived From Electrostatic Ones）}

本节沿用与上一节平行的思路，把\emph{静电}解改造成振荡解。与静磁情形不同的是：静电势来自四维势的时间分量，当我们把一个纯静电势乘以 $\cos(\omega t)$ 时，它自身会产生非零的时间导数，而四维势若只有时间分量，就没有空间分量来抵消这一导数，于是不能直接满足横波条件（transverse condition）。因此需要额外引入一个空间分量。我们先处理振荡电偶极子，再处理带交流电流的球形电容器。

\subsection{振荡电偶极子（Oscillating Electric Dipoles）}

我们把纯静电势 $\phi_{\text{ES}}$ 改用一个新常数 $k$（而非 $\wms$）重新标定，再乘以 $\cos(\omega t)$；这样得到的解会以原电荷密度乘以 $\cos(\omega t)$ 为源满足黎曼矢量波方程（Riemannian Vector Wave equation, RVWS），其中照例有 $k^{2}+\omega^{2}=\wms^{2}$。但它\emph{不}满足横波条件：此时四维势的时间分量为 $-\phi_{\text{ES},k}\cos(\omega t)$，具有非零的时间导数，而四维势没有空间分量，其自身的导数不足以使散度之和为零。

对静态电偶极子，有一个简单的办法绕过这一困难。静态偶极势恰是库仑势沿偶极轴（取为 $z$ 轴）方向空间导数的相反数。于是，若令四维势的 $z$ 分量等于库仑势（同样用常数 $k$ 而非 $\wms$ 标定）乘以 $\omega\sin(\omega t)$，它在 $z$ 方向的空间导数便恰好抵消时间分量带来的时间导数，从而满足横波条件。新添的项也以同样方式修改后的源满足 RVWS，且电荷自动守恒。具体而言，这给源增添了一个\emph{点状振荡电流}（pointlike oscillating current），其相位比偶极强度的振荡落后九十度。

与之前一样，可用此法构造两个驻波，再组合为出射行波；也同样可以改造成洛伦兹解。

\medskip
\noindent\textbf{振荡电偶极子·出射行波解}：电偶极矩为 $\bp\cos(\omega t)$；对黎曼解，$k^{2}+\omega^{2}=\wms^{2}$。

对应的四维势为
\begin{align}
\text{（Riemannian）}\quad
\phi(\mathbf{r}) &= -\frac{\bp\cdot\mathbf{r}}{4\pi r^{3}}\bigl[\cos(kr+\omega t)+kr\sin(kr+\omega t)\bigr] \nonumber\\
A_{(3)}(\mathbf{r}) &= -\frac{\bp\,\omega}{4\pi r}\sin(kr+\omega t)\\
\text{（Lorentzian）}\quad
\phi(\mathbf{r}) &= \frac{\bp\cdot\mathbf{r}}{4\pi r^{3}}\bigl[\cos\bigl(\omega(r-t)\bigr)+\omega r\sin\bigl(\omega(r-t)\bigr)\bigr] \nonumber\\
A_{(3)}(\mathbf{r}) &= \frac{\bp\,\omega}{4\pi r}\sin\bigl(\omega(r-t)\bigr)
\end{align}

相应的电场与磁场为
\begin{align}
\text{（Riemannian）}\quad
\bE(\mathbf{r}) &= -\frac{3(\bp\cdot\mathbf{r})\mathbf{r}-r^{2}\bp}{4\pi r^{5}}\bigl[\cos(kr+\omega t)+kr\sin(kr+\omega t)\bigr] \nonumber\\
&\quad + \frac{k^{2}(\bp\cdot\mathbf{r})\mathbf{r}+\omega^{2}r^{2}\bp}{4\pi r^{3}}\cos(kr+\omega t) \nonumber\\
\bB(\mathbf{r}) &= \frac{\omega\,\bp\times\mathbf{r}}{4\pi r^{3}}\bigl[kr\cos(kr+\omega t)-\sin(kr+\omega t)\bigr]\\
\text{（Lorentzian）}\quad
\bE(\mathbf{r}) &= \frac{3(\bp\cdot\mathbf{r})\mathbf{r}-r^{2}\bp}{4\pi r^{5}}\bigl[\cos\bigl(\omega(r-t)\bigr)+\omega r\sin\bigl(\omega(r-t)\bigr)\bigr] \nonumber\\
&\quad - \frac{\omega^{2}(\bp\times\mathbf{r})\times\mathbf{r}}{4\pi r^{3}}\cos\bigl(\omega(r-t)\bigr) \nonumber\\
\bB(\mathbf{r}) &= \frac{\omega\,\bp\times\mathbf{r}}{4\pi r^{3}}\bigl[\sin\bigl(\omega(r-t)\bigr)-\omega r\cos\bigl(\omega(r-t)\bigr)\bigr]
\end{align}

一个周期上平均的坡印廷矢量（Poynting vector）与总功率为
\begin{align}
\text{（Riemannian）}\quad
\langle\mathbf{S}(r)\rangle &= \frac{k\omega^{3}(\bp\cdot\bp)+k^{3}\omega(\bp\cdot\be_{r})^{2}}{32\pi^{2}r^{2}}\,\be_{r}\\
\text{（Lorentzian）}\quad
\langle\mathbf{S}(r)\rangle &= \frac{\omega^{4}\bigl[(\bp\cdot\bp)-(\bp\cdot\be_{r})^{2}\bigr]}{32\pi^{2}r^{2}}\,\be_{r}
\end{align}
\begin{align}
\text{（Riemannian）}\quad
\langle P\rangle &= \frac{(k^{3}\omega+3k\omega^{3})(\bp\cdot\bp)}{24\pi}\\
\text{（Lorentzian）}\quad
\langle P\rangle &= \frac{\omega^{4}(\bp\cdot\bp)}{12\pi}
\end{align}

\orthofig{images/04__016.png}{振荡电偶极子的辐射功率随频率的变化}

黎曼解在这里的功率-频率关系与振荡磁偶极子略有不同。它也是我们迄今见到的第一个能产生\emph{纵向偏振波}（longitudinally polarised waves）的显式源。

对于很大的 $r$，黎曼四维势可写成
\begin{equation}
\bA(\mathbf{r}) \approx \frac{\sin(kr+\omega t)}{4\pi r}\bigl[k(\bp\cdot\be_{r})\be_{t}-\omega\,\bp\bigr].
\end{equation}
可把它分解为横波与纵波两部分：
\begin{align}
\bA_{T}(\mathbf{r}) &\approx \frac{\sin(kr+\omega t)}{4\pi r}\,\omega\bigl[(\bp\cdot\be_{r})\be_{r}-\bp\bigr],\\
\bA_{L}(\mathbf{r}) &\approx \frac{\sin(kr+\omega t)}{4\pi r}(\bp\cdot\be_{r})\bigl[k\,\be_{t}-\omega\,\be_{r}\bigr].
\end{align}
横波部分没有时间分量，且与波在空间中的传播方向 $\be_{r}$ 正交。利用我们对平面波能量的分析，记 $\theta$ 为偶极矢量与波在空间中传播方向之间的夹角，则横模与纵模中一个时间周期上平均的局部能量密度为
\begin{align}
\langle u_{T}(r)\rangle &\approx \frac{1}{32\pi^{2}r^{2}}(\bp\cdot\bp)\,\omega^{4}\sin^{2}\theta,\\
\langle u_{L}(r)\rangle &\approx \frac{1}{32\pi^{2}r^{2}}(\bp\cdot\bp)(\omega^{2}+k^{2})\,\omega^{2}\cos^{2}\theta.
\end{align}
由此可见，横波在垂直于偶极轴的方向最强，在轴上降为零；纵波则相反：在轴上最强，垂直方向降为零。横波的能量角分布与洛伦兹情形一致。

把这两个能量密度乘以波速 $k/\omega$ 并对整个球面积分，得到两种形式各自的总功率：
\begin{align}
\langle P_{T}\rangle &= \frac{k\omega^{3}(\bp\cdot\bp)}{12\pi},\\
\langle P_{L}\rangle &= \frac{(k\omega^{3}+k^{3}\omega)(\bp\cdot\bp)}{24\pi}.
\end{align}
两者相加正是表中给出的总辐射功率。

\subsection{电容器中的交流电流（Alternating Current in a Capacitor）}

为了考察球形电容器由交流电流对两壳层充放电时的行为，我们须把电流放进图像中。在不破坏球对称的前提下，唯一可行的做法是让电流对称地从一壳层直接流到另一壳层、穿过二者之间的空隙。字面意义上这样做显然会破坏电容器的功能，但可以把这一模型看作如下装置的近似：大量平板电容器通过位于器件\emph{旁侧}（而非真正穿过器件内部）的导线充放电；把大量此类电路排列成围绕某一中心点扇形展开的样式，将产生与本模型非常相似的场。

我们先由球壳上的振荡电荷求黎曼四维势：改造静电解，把 $\wms$ 换成 $k$ 并乘以 $\cos(\omega t)$。然后再加入一项四维势，描述在两壳层之间往返流动的电流；它可先用毕奥-萨伐尔定律作为静磁解求出，再用标准方法转成振荡解。由于所加电流恰好对应球壳上的振荡电荷，电荷守恒成立，故四维势满足横波条件，给出合法解。此后照例需要把它与一个无源驻波组合，得到出射行波解。

由于四维势是径向对称的，不存在磁场。在洛伦兹情形下，这意味着根本不可能有辐射；而在黎曼情形下，存在纯横波的辐射。

\medskip
\noindent\textbf{球形电容器（交流）·出射行波解}：内壳半径 $R_{1}$、总电荷 $-Q_{0}\cos(\omega t)$；外壳半径 $R_{2}$、总电荷 $+Q_{0}\cos(\omega t)$；对黎曼解，$k^{2}+\omega^{2}=\wms^{2}$。

\paragraph{势。}
黎曼解的四维势为
\begin{align}
\text{（Riemannian）}\quad
&\varphi(r) = \frac{Q_{0}\sin(kr)}{4\pi kr R_{1}R_{2}}\bigl[R_{2}\cos(kR_{1}+\omega t) \nonumber\\
&\qquad - R_{1}\cos(kR_{2}+\omega t)\bigr], \qquad r<R_{1}; \nonumber\\
&\varphi(r) = \frac{Q_{0}}{4\pi kr R_{1}R_{2}}\bigl[R_{2}\sin(kR_{1})\cos(kr+\omega t) \nonumber\\
&\qquad - R_{1}\sin(kr)\cos(kR_{2}+\omega t)\bigr], \qquad R_{1}<r<R_{2}; \nonumber\\
&\varphi(r) = \frac{Q_{0}\cos(kr+\omega t)}{4\pi kr R_{1}R_{2}}\bigl[R_{2}\sin(kR_{1})-R_{1}\sin(kR_{2})\bigr], \qquad r>R_{2}.
\end{align}
\begin{align}
\text{（Riemannian）}\quad
&\bA_{(3)}(r) = \frac{Q_{0}\omega\bigl(kr\cos(kr)-\sin(kr)\bigr)}{4\pi k^{3}r^{2}R_{1}R_{2}} \nonumber\\
&\qquad \bigl[R_{2}\sin(kR_{1}+\omega t)-R_{1}\sin(kR_{2}+\omega t)\bigr]\be_{r}, \qquad r<R_{1}; \nonumber\\
&\bA_{(3)}(r) = \frac{Q_{0}\omega}{4\pi k^{3}r^{2}R_{1}R_{2}}\Bigl[\cos(\omega t)\bigl(kr\cos(kr)-\sin(kr)\bigr) \nonumber\\
&\qquad\qquad \bigl(R_{2}\sin(kR_{1})-R_{1}\sin(kR_{2})\bigr) \nonumber\\
&\qquad\qquad + \sin(\omega t)\Bigl(R_{1}\cos(kR_{2})\bigl(\sin(kr)-kr\cos(kr)\bigr) \nonumber\\
&\qquad\qquad - R_{2}\sin(kR_{1})\bigl(kr\sin(kr)+\cos(kr)\bigr) \nonumber\\
&\qquad\qquad + kR_{1}R_{2}\Bigr)\Bigr]\be_{r}, \qquad R_{1}<r<R_{2}; \nonumber\\
&\bA_{(3)}(r) = \frac{Q_{0}\omega\bigl(kr\cos(kr+\omega t)-\sin(kr+\omega t)\bigr)}{4\pi k^{3}r^{2}R_{1}R_{2}} \nonumber\\
&\qquad \bigl[R_{2}\sin(kR_{1})-R_{1}\sin(kR_{2})\bigr]\be_{r}, \qquad r>R_{2}.
\end{align}
洛伦兹解的四维势为
\begin{align}
\text{（Lorentzian）}\quad
\varphi(r) &= \frac{Q_{0}(R_{1}-R_{2})\cos(\omega t)}{4\pi R_{1}R_{2}}, \qquad r<R_{1}; \nonumber\\
\varphi(r) &= \frac{Q_{0}(r-R_{2})\cos(\omega t)}{4\pi r R_{2}}, \qquad R_{1}<r<R_{2}; \nonumber\\
\varphi(r) &= 0, \qquad r>R_{2}; \nonumber\\
\bA_{(3)}(r) &= 0.
\end{align}

\paragraph{电场。} 黎曼解
\begin{align}
\text{（Riemannian）}\quad
&\bE(r) = \frac{Q_{0}\wms^{2}\bigl(kr\cos(kr)-\sin(kr)\bigr)}{4\pi k^{3}r^{2}R_{1}R_{2}} \nonumber\\
&\qquad \bigl[R_{1}\cos(kR_{2}+\omega t)-R_{2}\cos(kR_{1}+\omega t)\bigr]\be_{r}, \qquad r<R_{1}; \nonumber\\
&\bE(r) = \frac{Q_{0}}{4\pi k^{3}r^{2}R_{1}R_{2}}\Bigl[\cos(\omega t)\Bigl(\wms^{2}\bigl(R_{1}\cos(kR_{2})(kr\cos(kr)-\sin(kr)) \nonumber\\
&\qquad\qquad + R_{2}\sin(kR_{1})(kr\sin(kr)+\cos(kr))\bigr)-\omega^{2}kR_{1}R_{2}\Bigr) \nonumber\\
&\qquad\qquad + \wms^{2}\sin(\omega t)\bigl(kr\cos(kr)-\sin(kr)\bigr) \nonumber\\
&\qquad\qquad \bigl(R_{2}\sin(kR_{1})-R_{1}\sin(kR_{2})\bigr)\Bigr]\be_{r}, \qquad R_{1}<r<R_{2}; \nonumber\\
&\bE(r) = \frac{Q_{0}\wms^{2}\bigl(kr\sin(kr+\omega t)+\cos(kr+\omega t)\bigr)}{4\pi k^{3}r^{2}R_{1}R_{2}} \nonumber\\
&\qquad \bigl[R_{2}\sin(kR_{1})-R_{1}\sin(kR_{2})\bigr]\be_{r}, \qquad r>R_{2}.
\end{align}
洛伦兹解
\begin{align}
\text{（Lorentzian）}\quad
\bE(r) &= 0, \qquad r<R_{1}; \nonumber\\
\bE(r) &= -\frac{Q_{0}\cos(\omega t)}{4\pi r^{2}}\,\be_{r}, \qquad R_{1}<r<R_{2}; \nonumber\\
\bE(r) &= 0, \qquad r>R_{2}.
\end{align}

\paragraph{坡印廷矢量（一个周期上平均）。}
\begin{align}
\text{（Riemannian）}\quad
\langle\mathbf{S}(r)\rangle &= 0, \qquad r<R_{1}; \nonumber\\
\langle\mathbf{S}(r)\rangle &= \frac{Q_{0}^{2}\wms^{2}\omega}{32\pi^{2}k^{3}r^{3}R_{1}^{2}R_{2}}\bigl[\bigl(r\sin(kR_{1})-R_{1}\sin(kr)\bigr) \nonumber\\
&\qquad \bigl(R_{2}\sin(kR_{1})-R_{1}\sin(kR_{2})\bigr)\bigr]\be_{r}, \qquad R_{1}<r<R_{2}; \nonumber\\
\langle\mathbf{S}(r)\rangle &= \frac{Q_{0}^{2}\wms^{2}\omega}{32\pi^{2}k^{3}r^{2}R_{1}^{2}R_{2}^{2}}\bigl[R_{2}\sin(kR_{1})-R_{1}\sin(kR_{2})\bigr]^{2}\be_{r}, \qquad r>R_{2}.
\end{align}
\begin{equation}
\text{（Lorentzian）}\quad \langle\mathbf{S}(r)\rangle = 0.
\end{equation}

\paragraph{总辐射功率（一个周期上平均）。}
\begin{align}
\text{（Riemannian）}\quad
\langle P_{\text{Radiated}}\rangle &= \frac{Q_{0}^{2}\wms^{2}\omega}{8\pi k^{3}R_{1}^{2}R_{2}^{2}}\bigl[R_{2}\sin(kR_{1})-R_{1}\sin(kR_{2})\bigr]^{2} \nonumber\\
&\approx \frac{Q_{0}^{2}(R_{2}^{2}-R_{1}^{2})^{2}\,\wms^{3}k^{3}}{288\pi} \qquad (\textbf{低 } k \textbf{ 极限});\\
\text{（Lorentzian）}\quad
\langle P_{\text{Radiated}}\rangle &= 0.
\end{align}

\paragraph{两壳层之间的电压。}
\begin{align}
\text{（Riemannian）}\quad
V &= \frac{Q_{0}}{4\pi k^{3}R_{1}^{2}R_{2}^{2}}\Bigl[\wms^{2}\sin(\omega t)\bigl(R_{2}\sin(kR_{1})-R_{1}\sin(kR_{2})\bigr)^{2} \nonumber\\
&\qquad - \cos(\omega t)\Bigl(\wms^{2}\bigl(R_{2}^{2}\sin(kR_{1})\cos(kR_{1}) \nonumber\\
&\qquad\qquad + R_{1}\cos(kR_{2})(R_{1}\sin(kR_{2})-2R_{2}\sin(kR_{1}))\bigr) \nonumber\\
&\qquad\qquad + \omega^{2}kR_{1}R_{2}(R_{1}-R_{2})\Bigr)\Bigr] \nonumber\\
&\approx -\frac{Q_{0}(R_{2}-R_{1})\bigl(3+\wms^{2}R_{1}(R_{2}-R_{1})\bigr)\cos(\omega t)}{12\pi R_{1}R_{2}} \nonumber\\
&\qquad (\textbf{低 } k \textbf{ 极限});\\
\text{（Lorentzian）}\quad
V &= \frac{Q_{0}(R_{2}-R_{1})\cos(\omega t)}{4\pi R_{1}R_{2}}.
\end{align}

\paragraph{平均消耗的电功率。}
\begin{align}
\text{（Riemannian）}\quad
\langle P_{\text{Electrical}}\rangle &= -\frac{Q_{0}^{2}\wms^{2}\omega}{8\pi k^{3}R_{1}^{2}R_{2}^{2}}\bigl[R_{2}\sin(kR_{1})-R_{1}\sin(kR_{2})\bigr]^{2} \nonumber\\
&\approx -\frac{Q_{0}^{2}(R_{2}^{2}-R_{1}^{2})^{2}\,\wms^{3}k^{3}}{288\pi} \qquad (\textbf{低 } k \textbf{ 极限});\\
\text{（Lorentzian）}\quad
\langle P_{\text{Electrical}}\rangle &= 0.
\end{align}

如常，黎曼情形下电路消耗的电功率与总辐射功率相反，因此须从电路中\emph{提取}功，才能让振荡电流的峰值振幅保持不变。电压/电流相位差与振荡频率之间的精确关系会很复杂，但电路消耗的功率始终为负（最坏也是零）这一事实意味着：电压与电流的相位差永远至少为 90 度。

在洛伦兹情形下，几何结构杜绝了一切辐射损失，故电压与电流的相位差始终精确为 90 度。

\section{共振电路（Resonant Circuits）}

\subsection{洛伦兹电路中的共振（Resonance in Lorentzian Circuits）}

在我们自身宇宙所适用的基础电路理论中，通常假定电容器与电感器具有不依赖电流频率的固定电容值与电感值。这是一个合理假设，因为在洛伦兹宇宙中，中等的\emph{时间}频率对应的波长远大于典型电子元件的尺寸。

但这并不意味着含有这些元件的电路行为与频率无关。对电容器而言，电容 $C$ 固定的是器件中存储的电荷与极板间电压之比；但当我们考察电压与\emph{电流}（而非电荷）的关系时，电流频率会通过振荡电荷的导数进入关系式。下文中，我们用 $Q_{0}$、$I_{0}$、$V_{0}$ 表示瞬时值按谐波变化的一个振荡电荷、电流或电压的振幅。

\begin{align}
Q &= Q_{0}\sin(\omega t),\\
V &= \frac{Q}{C} = \frac{Q_{0}}{C}\sin(\omega t),\\
I &= \frac{dQ}{dt} = \omega Q_{0}\cos(\omega t),\\
V_{0} &= \frac{I_{0}}{\omega C}.
\end{align}

对电感器，$L$ 固定电压与电流变化率之比，于是
\begin{align}
I &= I_{0}\cos(\omega t),\\
\frac{dI}{dt} &= -\omega I_{0}\sin(\omega t),\\
V &= L\frac{dI}{dt} = -L\omega I_{0}\sin(\omega t),\\
V_{0} &= (L\omega)I_{0}.
\end{align}

若如下定义\textbf{容抗}（capacitative reactance）$X_{C}$ 与\textbf{感抗}（inductive reactance）$X_{L}$：
\begin{align}
X_{C} &= \frac{1}{\omega C},\\
X_{L} &= \omega L,
\end{align}
则 $X_{C}$、$X_{L}$ 扮演着类似电阻的角色，有
\begin{align}
V_{0} &= I_{0}R, && \text{（电阻器）}\\
V_{0} &= I_{0}X_{C}, && \text{（电容器）}\\
V_{0} &= I_{0}X_{L}, && \text{（电感器）}.
\end{align}

然而，这三种情形中电压的瞬时值各不相同：电阻器上电压与电流同相；电容器上电压滞后电流 90 度（若电流取余弦，则电压是正弦的正倍数）；电感器上电压超前电流 90 度（若电流取余弦，则电压是正弦的负倍数）。这意味着，若三者串联、因而流过同一电流，则电容器上的电压将与电感器上的电压相差 180 度，也就是恰好与之相反。于是净电抗
\begin{equation}
X = X_{L}-X_{C}
\end{equation}
决定了这两个元件上的合成电压，它与电流相差 90 度。

接着定义包含电阻效应的\textbf{阻抗}（impedance）$Z$ 以及总相位差 $\phi$：
\begin{align}
Z &= \sqrt{R^{2}+X^{2}},\\
\phi &= \arctan\!\left(\frac{X}{R}\right),\\
\cos\phi &= \frac{R}{Z},\\
\sin\phi &= \frac{X}{Z}.
\end{align}
这两个量让我们能描述串联的三种元件（电容器、电感器与电阻器）上的合成电压：
\begin{align}
V &= RI_{0}\cos(\omega t)-XI_{0}\sin(\omega t) \nonumber\\
&= I_{0}Z\bigl[\cos\phi\cos(\omega t)-\sin\phi\sin(\omega t)\bigr] \nonumber\\
&= I_{0}Z\cos(\omega t+\phi).
\end{align}

\orthofig{images/04__017.png}{洛伦兹共振电路中电流随频率的变化}

显然，当 $X$ 为零时阻抗取极小。若把此处的角频率记为 $\omega_{\text{res}}$，则
\begin{align}
X &= X_{L}-X_{C} = 0,\\
\omega_{\text{res}}L-\frac{1}{\omega_{\text{res}}C} &= 0,\\
\omega_{\text{res}} &= \frac{1}{\sqrt{LC}}.
\end{align}
在共振频率 $\omega_{\text{res}}$ 处，电感与电容恰好彼此抵消，电流振幅 $I_{0}$ 达到一个完全由电阻决定的峰值。

举一例：设有一螺线管，长 10 cm、半径 5 cm、每毫米一匝。按国际单位制，其直流电感为 0.987 毫亨。与之串联一个球形电容器，内半径 5 cm、外半径 5.01 cm，其直流电容为 2.787 纳法。把螺线管、电容器与一个 1000 欧姆电阻串联。公式给出共振角频率，对应普通频率 $\nu=95.96$ 千赫兹。

图中给出给定电压下、随频率变化流过这三者的电流；频率轴为对数刻度，纵轴已归一化，使单独流过 1000 欧姆电阻的电流取值为 1。正如预期，在 $10^{5}$ 赫兹附近有一个峰。在把电感与电容视为与频率无关的近似之外，还会有其它共振，但它们要到 GHz 频段才会出现，那里波长开始接近螺线管的尺寸。

现在设想：我们不给这三个元件加振荡电压，而是把电容器充到电荷 $Q_{i}$，然后合上电路，让电流流过。会发生什么？

回路中所有电压之和必须为零：
\begin{align}
V_{C}+V_{L}+V_{R} &= 0,\\
\frac{Q}{C}+L\frac{dI}{dt}+RI &= 0,\\
\frac{Q}{C}+L\frac{d^{2}Q}{dt^{2}}+R\frac{dQ}{dt} &= 0,\\
\frac{d^{2}Q}{dt^{2}}+2\beta\frac{dQ}{dt}+\omega_{\text{res}}^{2}Q &= 0,
\end{align}
其中我们定义了 $\beta=R/(2L)$。给定 $Q=Q_{i}$、且 $t=0$ 时 $dQ/dt=0$，并假设 $\beta<\omega_{\text{res}}$，该微分方程的解为
\begin{equation}
Q(t) = Q_{i}\,e^{-\beta t}\left[\cos\!\bigl(\sqrt{\omega_{\text{res}}^{2}-\beta^{2}}\,t\bigr)+\frac{\beta}{\sqrt{\omega_{\text{res}}^{2}-\beta^{2}}}\sin\!\bigl(\sqrt{\omega_{\text{res}}^{2}-\beta^{2}}\,t\bigr)\right].
\end{equation}
它描述一个作指数衰减的振荡函数。振荡频率小于电路对驱动电压以最小阻抗响应的共振频率 $\omega_{\text{res}}$；不过随着电阻减小，振荡会趋近该频率。

\subsection{黎曼电路中的共振（Resonance in Riemannian Circuits）}

上一节的概念可以搬用到黎曼宇宙中的类似情形，但有若干重大变化。首先，很难再合理地假定 $L$ 与 $C$ 本身不依赖 $\omega$。在黎曼宇宙中，\emph{静态}场的波长取\emph{最小值}，且只有随时间频率增大而变长。波长的增长来得很晚，而且非常突然：直到 $\omega=0.866\,\wms$ 时波长才达到最小值的两倍，在 $\omega=0.995\,\wms$ 时达到最小值十倍，在 $\omega=0.99995\,\wms$ 时达到一百倍。因此，除极高的频率外，任何频率下工作的电感器或电容器，其电流-电压关系都由场波长与元件几何的相互作用决定，从而远比上文为洛伦兹情形给出的电抗-频率公式复杂地依赖于频率。

此外，黎曼电感器与电容器发出的电磁辐射会给它们带来显著的、依赖频率的\emph{负电阻}。这就在阻抗的电阻部分与相位差中引入了依赖频率的项：
\begin{align}
X &= X_{L}-X_{C},\\
R_{\text{tot}} &= R-R_{\text{rad}},\\
Z &= \sqrt{R_{\text{tot}}^{2}+X^{2}},\\
\phi &= \arctan\!\left(\frac{X}{R_{\text{tot}}}\right),
\end{align}
其中除普通电阻 $R$ 外，这里的一切都依赖频率（而就连这一点也是简化假设）。因此，我们再也无法保证 $X=0$ 处的频率会给出最小阻抗 $Z$。

\orthofig{images/04__018.png}{黎曼电感器与电容器辐射的功率}

尽管如此，这些额外的复杂性意味着即便一个很简单的电路也能有有趣的行为。设有一个与上一节所述完全相同的螺线管：长 10 cm、半径 5 cm、每毫米一匝。为把它代入黎曼物理，我们取 $\wms=2\pi\times10^{15}$ Hz。

图中给出了该螺线管的电抗与电阻。注意这里所有的波长都对应\emph{极其}接近 $\wms$ 的时间频率。由于螺线管单独的电抗就跨越零点，电路无需再加电容器；若把这个螺线管与一个普通电阻器接成回路，使该电阻器在螺线管电抗为零的最长波长处恰好平衡它的负电阻，则仅含这两个元件的闭合回路就会在该波长处共振，原则上可无限期地维持电流。螺线管会发出电磁波，把\emph{常规}能量带入电路，而电阻器会把该能量变成热。这不会违反任何热力学定律：能量守恒，因为电磁场能量与热/动能具有相反的符号；熵增加，因为有辐射产生。

令人惊讶的是，这个极简电路的行为中还内建了某种稳定性。若电流开始指数增长，那将意味着其频率展宽；而尽管共振点并不\emph{完全}位于最小电阻对应的波长处，这里时间频率的差异却如此之小，以至于只需极小的指数增长率就足以把频率展宽到使螺线管向电路馈能的速率下降。当然，若电流开始下降，同样的效应会加剧其阻尼；因此需要额外的调节机制（例如一个非线性电阻器，其电阻随电流增大而增大），才能让电路以恒定速率采集能量。

\section{弯曲黎曼空间中的电磁学（Electromagnetism in Curved Riemannian Space）}

以上关于电磁学的一切，都是在平坦空间（洛伦兹情形下为平坦时空）中、用笛卡尔坐标表述的。但我们并不真的指望黎曼宇宙处处完美平坦，正如不指望我们自身的宇宙如此；因此，理解如何把方程改写以适应弯曲空间会有所帮助。这还有一个附带好处：让我们能轻松处理平坦空间中的非笛卡尔坐标。

若你此前没做过弯曲时空中的计算，下面这份简明小结也许会让你有点晕。这里仅给出所需的要点。

在广义相对论的洛伦兹物理中，把方程从平坦时空搬到弯曲时空的经验法则是：把偏导数换成协变导数（covariant derivative）。在平坦空间中对矢量场求导时，我们默认把不同点处的矢量看作属于同一矢量空间；若说一个矢量场导数为零、因而是常量，这个断言只有在能把 A 点的矢量与 B 点的矢量作比较时才有意义。但在诸如地球曲面这样的弯曲面上，我们如何把伦敦地面上各种可能速度构成的矢量空间与内罗毕的同类矢量空间作比较？即便离开地球表面、把这些矢量视为三维的，也无法把一个位置处的所有速度与另一个位置处的速度一一配对——而在弯曲宇宙中，我们根本无从“离开表面”。

解决办法是给“导数”这一概念补充一个称为列维-奇维塔联络（Levi-Civita connection）的几何结构，它给出沿曲线对矢量作\textbf{平行移动}（parallel transport）的概念：即沿一条曲线行进时，可以按该联络把起点处的一个矢量“携带”着一起走，并使之始终保持与原方向“平行”。列维-奇维塔联络的优点在于与度规相容：度规定义了弯曲空间上的点积，而列维-奇维塔联络让你在平行移动两个矢量时保持其点积不变。协变导数计算的是矢量场\emph{相对于}列维-奇维塔联络的导数：若你沿曲线用列维-奇维塔联络平行移动一个矢量，这就是判定“此矢量不变”的标准，协变导数据此识别任何变化。

具体地说，设弯曲空间上有一个矢量场 $\mathbf{v}$，其在某坐标基下的分量为 $v^{b}$。则此矢量场沿某一坐标方向 $a$ 的协变导数为
\begin{equation}
\nabla_{a} v^{b} = \partial_{a} v^{b} + \Gamma^{b}_{\;ca} v^{c},
\end{equation}
其中 $\Gamma$ 就是列维-奇维塔联络，它告诉我们如何修正偏导数，以得到尊重平行移动的导数。若 $g_{ab}$ 与 $g^{ab}$ 是度规在我们的坐标系中的分量，则列维-奇维塔联络 $\Gamma$ 的分量（常称为克里斯托费尔符号，Christoffel symbols）为
\begin{equation}
\Gamma^{b}_{\;ca} = \tfrac12\,g^{bk}\bigl[\partial_{a}g_{kc}+\partial_{c}g_{ka}-\partial_{k}g_{ca}\bigr].
\end{equation}
注意 $\Gamma$ 对其后两个指标对称：$\Gamma^{b}_{\;ca}=\Gamma^{b}_{\;ac}$。

我们可以把平行移动的概念从矢量推广到任何张量。例如，若用列维-奇维塔联络把矢量 $\mathbf{v}$ 与 $\mathbf{w}$ 从 A 点平行移动到 B 点，得到 B 点的 $\mathbf{v}'$ 与 $\mathbf{w}'$，则秩 (2,0) 张量的平行移动按如下方式定义：A 点的 $\mathbf{v}\otimes\mathbf{w}$ 变为 B 点的 $\mathbf{v}'\otimes\mathbf{w}'$。对偶矢量则要求：若 A 点的对偶矢量 $\bm{\alpha}$ 满足 $\bm{\alpha}(\mathbf{v})=c$，则 $\bm{\alpha}$ 从 A 到 B 的平行移动得到 $\bm{\alpha}'$，使 $\bm{\alpha}'(\mathbf{v}')=c$。

这些要求给出我们所需张量的协变导数公式：
\begin{align}
\nabla_{a} A_{b} &= \partial_{a} A_{b} - \Gamma^{h}_{\;ba} A_{h},\\
\nabla_{a} F_{bc} &= \partial_{a} F_{bc} - \Gamma^{h}_{\;ba} F_{hc} - \Gamma^{h}_{\;ca} F_{bh},\\
\nabla_{a} F^{bc} &= \partial_{a} F^{bc} + \Gamma^{b}_{\;ha} F^{hc} + \Gamma^{c}_{\;ha} F^{bh}.
\end{align}
把第二式用于度规、并利用 $\Gamma$ 的定义，可得 $\nabla_{a} g_{bc}=0$。本质上，我们对 $\Gamma$ 的定义正是为了得到这一结果：$\Gamma$ 是使度规本身被判为常量的联络。

若把电磁学方程中的偏导数换成协变导数，就得到：
\begin{align}
\nabla_{b} F^{ab} - \wms^{2} A^{a} - j^{a} &= 0 && \text{（Riemannian）} \tag{Proca}\\
\partial_{a} F_{bc} + \partial_{b} F_{ca} + \partial_{c} F_{ab} &= 0 && \text{（Common）}\\
\nabla_{b} F^{ab} - j^{a} &= 0 && \text{（Lorentzian）}\\
\partial_{a} F_{bc} + \partial_{b} F_{ca} + \partial_{c} F_{ab} &= 0 && \text{（Common）}
\end{align}

为什么黎曼与洛伦兹电磁学共有的那条方程里仍然是偏导数，而非协变导数？若用协变导数写出该方程，并利用 $F_{bc}$ 反对称、而 $\Gamma$ 对后两个指标对称这一点，所有修正项会两两抵消，只剩下偏导数。

在电磁场 $F$ 与四维势 $\bA$ 的关系中，协变导数的修正项也同样抵消：
\begin{align}
F_{ab} &= \nabla_{a} A_{b} - \nabla_{b} A_{a} \nonumber\\
&= \partial_{a} A_{b} - \partial_{b} A_{a} && \text{（Common）}.
\end{align}
由此可见，只要这样用 $\bA$ 定义 $F$，黎曼 Proca 方程与麦克斯韦方程组中共有的那条方程就仍会自动满足，因为这里什么都没变，出现的偏导数与平坦时空情形完全相同。

接下来这一步稍显微妙。在平坦空间或时空中，偏导数\emph{对易}：求两次导数时先后次序无关紧要。在弯曲空间中，协变导数则不然；事实上，“曲率”这一概念本身正与协变导数不对易这一点紧密相连。

设沿两个不同坐标方向（指标为 $a$、$b$）以两种次序对矢量场 $\mathbf{v}$ 求协变导数，二者之差为
\begin{equation}
\nabla_{a}\nabla_{b}\mathbf{v} - \nabla_{b}\nabla_{a}\mathbf{v} = R^{h}_{\;cab}\,v^{c}\,\be_{h},
\end{equation}
其中 $\be_{h}$ 是指标为 $h$ 的坐标方向上的基矢量，而四指标张量 $R$ 就是所谓的黎曼曲率张量（Riemann curvature tensor）——它得名于我们一路提到的同一位 Georg Friedrich Bernhard Riemann（尽管这个张量在洛伦兹弯曲时空中与在黎曼弯曲空间中同样有用）。把协变导数用列维-奇维塔联络显式展开，即可把黎曼曲率张量的分量写成
\begin{equation}
R^{h}_{\;cab} = \partial_{a}\Gamma^{h}_{\;cb} - \partial_{b}\Gamma^{h}_{\;ca} + \Gamma^{h}_{\;ka}\Gamma^{k}_{\;bc} - \Gamma^{h}_{\;kb}\Gamma^{k}_{\;ca}.
\end{equation}

现在，设四维势 $\bA$ 满足横波条件或洛伦兹规范（Lorenz gauge condition）的协变导数版本：
\begin{equation}
\nabla_{b} A^{b} = 0,
\end{equation}
其中照例对重复指标使用爱因斯坦求和约定（Einstein Summation Convention）。若协变导数对易，则表达式 $\nabla_{b}\nabla_{a} A^{b}$ 将为零……但它们并不对易，于是我们得到
\begin{equation}
\nabla_{b}\nabla_{a} A^{b} = \nabla_{b}\nabla_{a} A^{b} - \nabla_{a}\nabla_{b} A^{b} = R^{b}_{\;cba} A^{c} = R_{ca} A^{c},
\end{equation}
其中二指标张量 $R$ 称为里奇曲率张量（Ricci curvature tensor），由对黎曼曲率张量“缩并”（即对其中两个指标求和）而得。

若利用这一结果，用四维势 $\bA$ 来计算 $\nabla_{b} F^{ab}$——它同时出现在黎曼 Proca 方程与麦克斯韦方程组中——则得到
\begin{align}
\nabla_{b} F^{ab} &= g^{\alpha a} g^{\beta b}\, \nabla_{b} F_{\alpha\beta} \nonumber\\
&= g^{\alpha a} g^{\beta b}\, \nabla_{b}(\nabla_{\alpha} A_{\beta} - \nabla_{\beta} A_{\alpha}) \nonumber\\
&= g^{\alpha a} g^{\beta b}\, (\nabla_{b}\nabla_{\alpha} A_{\beta} - \nabla_{b}\nabla_{\beta} A_{\alpha}) \nonumber\\
&= g^{\alpha a}\, (\nabla_{b}\nabla_{\alpha} A^{b} - \nabla_{b}\nabla^{b} A_{\alpha}) \nonumber\\
&= g^{\alpha a}\, (R_{c\alpha} A^{c} - \nabla_{b}\nabla^{b} A_{\alpha}) \nonumber\\
&= R_{c}^{\;a} A^{c} - \nabla_{b}\nabla^{b} A^{a}.
\end{align}
现在可以把一切用四维势 $\bA$ 表示：
\begin{align}
\nabla_{b}\nabla^{b} A^{a} - R_{c}^{\;a} A^{c} + \wms^{2} A^{a} + j^{a} &= \mathbf{0} \tag{RVWS}\\
\nabla_{c} A^{c} &= 0 \tag{Transverse}\\
\nabla_{b}\nabla^{b} A^{a} - R_{c}^{\;a} A^{c} + j^{a} &= \mathbf{0} \tag{LVWS}\\
\nabla_{c} A^{c} &= 0 \tag{Lorenz}
\end{align}
这里前两式是弯曲空间中的黎曼矢量波方程（RVWS），后两式是弯曲时空中洛伦兹规范下的麦克斯韦方程组（LVWS）。

既然谈到弯曲空间中的波方程，我们还可以给出一个修改过的标量波方程。标量的协变导数就是同方向的偏导数，标量的梯度无需借助度规或列维-奇维塔联络即可定义。然而，黎曼标量波方程中出现的、沿所有坐标方向的\emph{二阶导数}之和，只有在一般情况下把它按“梯度的散度”用协变导数计算时，才与坐标系无关：
\begin{align}
(\operatorname{grad} A)_{j} &= \partial_{j} A,\\
\operatorname{div}\operatorname{grad} A &= g^{ij}\bigl(\partial_{i}\partial_{j} A - \Gamma^{k}_{\;ji}\partial_{k} A\bigr).
\end{align}
这个运算（拉普拉斯算符的推广）称为拉普拉斯-贝尔特拉米算符（Laplace-Beltrami operator）。当度规只有对角分量时（许多坐标系都是如此），可以很方便地把度规行列式算作各对角元之积。记 $|g|$ 为度规行列式的绝对值，可以证明拉普拉斯-贝尔特拉米算符还有一个等价表达式：
\begin{equation}
\operatorname{div}\operatorname{grad} A = \frac{1}{\sqrt{|g|}}\,\partial_{i}\bigl[\sqrt{|g|}\,g^{ij}\,\partial_{j} A\bigr],
\end{equation}
它比费力去算克里斯托费尔符号更好用。即便你此前没见过这个方程，盯着它看久了，大概也能辨认出它藏在球坐标或柱坐标下拉普拉斯算符的公式背后。

于是有
\begin{align}
\operatorname{div}\operatorname{grad} A + \wms^{2} A + j &= \mathbf{0} \tag{RSWS}\\
\operatorname{div}\operatorname{grad} A + j &= \mathbf{0} \tag{LSWS}
\end{align}
前者是弯曲空间中带源的黎曼标量波方程，后者是弯曲时空中带源的洛伦兹标量波方程。

\medskip
\noindent\textbf{延伸阅读：} \emph{Gravitation} 一书第 16.3 与 22.4 节，作者 Charles Misner、Kip Thorne 与 John Wheeler，W.H. Freeman，San Francisco，1973 年。

\section{边界条件（Boundary Conditions）}

正如我们初次讨论黎曼波方程时提到的，这些方程有个严重问题：它们允许在某一方向上具有高于 $\wms$ 的角频率、同时在另一方向上作指数变化的解。我们据以导出黎曼标量波方程的那些良好、行为端正的平面波，其四维频率的平方和等于一个固定总量 $\nu_{\max}^{2}$，故任何单个方向的频率都不能超过 $\nu_{\max}$；但方程本身无法排除带有指数因子的解，例如 $\cos(kx)\,e^{\alpha t}$，只要 $k^{2}-\alpha^{2}=\wms^{2}$，它就满足黎曼标量波方程。

若你读过《正交》第一卷，就会知道如何避免这些指数解。若你虽未读此书、却已冒着剧透警告读到这里，那么这是你决定是否继续读下去的最后机会。

若黎曼宇宙是有限而无边界的，则要求波方程的解连续、且具有连续导数，就会排除带指数因子的解。周期函数天生就能在沿闭合曲线绕行一周后与自身光滑衔接，而指数函数做不到这一点。（从真空中的自由波过渡到带源的场时，情形会略微微妙一些，我们将在以下各节看到一些例子。）

至今我们大多把黎曼宇宙当作一个无限的、完全平坦的四维空间，同时指出这只是一个近似，类似于把洛伦兹宇宙近似为平坦闵可夫斯基时空这一有用的近似。本着同样的精神，我们可以考察黎曼宇宙的两个\emph{理想化}模型：它们有限，但仍对曲率作简化假设。在其一，即 4-环面（4-torus）模型中，黎曼宇宙保持完全平坦；在另一，即 4-球面（4-sphere）模型中，宇宙具有恒定正曲率。

\subsection{4-环面（The 4-torus）}

取平坦四维空间中一块呈矩形超棱柱（hyperprism）的区域，在其上放置坐标 $(x,y,z,t)$，取值范围分别为 $-L^{x}/2$ 到 $L^{x}/2$、$-L^{y}/2$ 到 $L^{y}/2$、$-L^{z}/2$ 到 $L^{z}/2$、$-L^{t}/2$ 到 $L^{t}/2$。然后我们宣布这个超棱柱的八个三维超面都与其对面“粘合”。例如，所有点 $(x,y,z,-L^{t}/2)$ 都与对应点 $(x,y,z,L^{t}/2)$ 等同。这是如下二维操作的四维版本：在平面上取一个矩形，把它的对边等同起来做成一个环面。

不过须强调：整个四维空间仍然\emph{完全平坦}；我们并没有把超棱柱“卷进”任何更高维空间，而只是规定这个黎曼宇宙模型在所有方向上有限、且其拓扑形如所述，这种拓扑称为\textbf{4-环面}（4-torus）。我们对拓扑的选择并不要求四维空间处处曲率为零，但当然允许它。

下文我们把这个模型宇宙称为 $\text{T}^{4}$。我们假定整个四维空间是平坦的，并已选用了如上所述那样的坐标。原点位置、以及坐标从 $L^{i}/2$ 跳到 $-L^{i}/2$ 的位置，当然都没有物理上的特殊性；用原坐标得到的任何黎曼物理解，若把一切整体平移一个任意位移矢量，仍然成立。然而，$\text{T}^{4}$ 宇宙的形状所施加的边界条件并不具旋转对称性，因此若取一个解并施加任意旋转，它便不再满足这些边界条件。

$\text{T}^{4}$ 上任何足够好的标量函数 $A(x,y,z,t)$ 都可写成傅里叶级数
\begin{equation}
A(x,y,z,t) = \sum_{i,j,k,l} a_{i,j,k,l}\, f_{i,j,k,l}(x,y,z,t),
\end{equation}
其中求和遍历 $i,j,k,l$ 的一切整数值（正、负与零），而
\begin{align}
f_{i,j,k,l}(x,y,z,t) &= f_{i}\!\left(\frac{x}{L^{x}}\right) f_{j}\!\left(\frac{y}{L^{y}}\right) f_{k}\!\left(\frac{z}{L^{z}}\right) f_{l}\!\left(\frac{t}{L^{t}}\right),\\
f_{n}(u) &= \sin(2\pi n u), && n>0,\\
f_{n}(u) &= \cos(2\pi n u), && n<0,\\
f_{0}(u) &= \frac{1}{\sqrt2}.
\end{align}
我们把函数 $f_{i,j,k,l}$ 称为 $\text{T}^{4}$ 的\textbf{傅里叶基函数}（Fourier basis functions）。以 $\text{T}^{4}$ 上的积分为函数间的内积
\begin{equation}
\langle f,g\rangle = \int_{\text{T}^{4}} fg,
\end{equation}
则不同基函数彼此正交，且都有相同的模方：$V/16$，其中 $V=L^{x}L^{y}L^{z}L^{t}$ 是 $\text{T}^{4}$ 的总体积（四维体积）。

每个基函数都是一个驻波，在 $x$、$y$、$z$、$t$ 方向上绕整个宇宙宽度分别经历 $|i|$、$|j|$、$|k|$、$|l|$ 个周期。给定函数 $A(x,y,z,t)$，可以显式算出傅里叶系数 $a_{i,j,k,l}$：
\begin{equation}
a_{i,j,k,l} = \frac{16}{V}\int_{\text{T}^{4}} f_{i,j,k,l}(x,y,z,t)\,A(x,y,z,t).
\end{equation}

现在想知道：哪些（如果有的话）$f_{i,j,k,l}$ 满足无源黎曼标量波方程？把该微分方程作用于一个傅里叶基函数，得到一个代数方程：
\begin{equation}
\left(\frac{i}{L^{x}}\right)^{2}+\left(\frac{j}{L^{y}}\right)^{2}+\left(\frac{k}{L^{z}}\right)^{2}+\left(\frac{l}{L^{t}}\right)^{2} = \nu_{\max}^{2},
\end{equation}
其中 $\nu_{\max}=\wms/(2\pi)$。若 $L^{i}$ 与 $\nu_{\max}$ 只是随机选的数，则\emph{没有}整数 $i,j,k,l$ 能满足此式。于是有两种情形要考虑：无源 RSW 方程无解的一般情形，以及 $L^{i}$ 与 $\nu_{\max}$ 取值恰好容许某些解存在的特殊情形。

\subsubsection{允许无源解的特殊情形（Special Case Allowing Sourceless Solutions）}

为举出特殊情形的例子，设所有 $L^{i}=1$、$\nu_{\max}=\sqrt{90}$。于是任何平方和为 90 的整数 $i,j,k,l$ 都会给出满足 RSW 方程的傅里叶基函数 $f_{i,j,k,l}$。若把一切排列与正负号选择都算上，这样的整数四元组共有 1872 个，但它们都可从下列九个方程导出：
\begin{align*}
0^{2}+0^{2}+3^{2}+9^{2} &= 90, & 0^{2}+1^{2}+5^{2}+8^{2} &= 90, & 0^{2}+4^{2}+5^{2}+7^{2} &= 90,\\
1^{2}+2^{2}+2^{2}+9^{2} &= 90, & 1^{2}+2^{2}+6^{2}+7^{2} &= 90, & 1^{2}+3^{2}+4^{2}+8^{2} &= 90,\\
2^{2}+5^{2}+5^{2}+6^{2} &= 90, & 3^{2}+3^{2}+6^{2}+6^{2} &= 90, & 3^{2}+4^{2}+4^{2}+7^{2} &= 90.
\end{align*}
在一个宇宙尺度与光的最小波长之比可与我们的可观测宇宙尺度（以远紫外光波长计约为 $10^{34}$）相比拟的黎曼宇宙中，对合适的 $L^{i}$ 与 $\nu_{\max}$ 取值，解的数目会极其巨大。我们不去深究计数解所需的数论（可参阅 Mathworld 的 Sum of Squares function 页面略窥一斑），但从直观上看，在宇宙尺度上离散解的数目很容易大到看起来连续。换言之，尽管这种黎曼宇宙中的无源平面波只能取有限多个特定的传播矢量，实际可选者却多到仿佛是一个包含所有方向的连续统。

由于无源解都由\emph{有限}多个傅里叶基函数构成，它们处处光滑有限。其各方向频率都不会超过 $\nu_{\max}$，而且它们同样可以写成有限多个平面波的叠加——这正是我们最初设想的构造波方程一般解的方式。

对带源的标量波方程（其源记作 $H$），我们能说什么？
\begin{equation}
\partial_{x}^{2}A(\mathbf{x}) + \partial_{y}^{2}A(\mathbf{x}) + \partial_{z}^{2}A(\mathbf{x}) + \partial_{t}^{2}A(\mathbf{x}) + \wms^{2} A(\mathbf{x}) + H(\mathbf{x}) = 0 \tag{RSWS}
\end{equation}
若把函数 $A$（系数为 $a$）与标量源 $H$（系数为 $h$）都作傅里叶展开，则有
\begin{equation}
a_{i,j,k,l}\left[\left(\frac{i}{L^{x}}\right)^{2}+\left(\frac{j}{L^{y}}\right)^{2}+\left(\frac{k}{L^{z}}\right)^{2}+\left(\frac{l}{L^{t}}\right)^{2}-\nu_{\max}^{2}\right] = \frac{h_{i,j,k,l}}{4\pi^{2}}.
\end{equation}
对那些满足无源方程的 $i,j,k,l$（即使方括号中表达式为零者），源的傅里叶系数 $h_{i,j,k,l}$ 必须为零，解才可能整体存在，而 $a_{i,j,k,l}$ 可自由选取。对其它取值，解上方程得到
\begin{equation}
a_{i,j,k,l} = \frac{h_{i,j,k,l}}{4\pi^{2}\left[\left(\dfrac{i}{L^{x}}\right)^{2}+\left(\dfrac{j}{L^{y}}\right)^{2}+\left(\dfrac{k}{L^{z}}\right)^{2}+\left(\dfrac{l}{L^{t}}\right)^{2}-\nu_{\max}^{2}\right]}.
\end{equation}
因此，源将决定所有\emph{不}对应无源解的系数，之后我们可以自由地再加上任意想要的无源解。

\subsubsection{无无源解的一般情形（Generic Case With No Sourceless Solutions）}

对一般的 $L^{i}$ 与 $\nu_{\max}$ 取值，没有任何傅里叶基函数会解无源黎曼标量波方程。此时源的傅里叶分量没有被强制为零者，我们总可以用
\begin{equation}
a_{i,j,k,l} = \frac{h_{i,j,k,l}}{4\pi^{2}\left[\left(\dfrac{i}{L^{x}}\right)^{2}+\left(\dfrac{j}{L^{y}}\right)^{2}+\left(\dfrac{k}{L^{z}}\right)^{2}+\left(\dfrac{l}{L^{t}}\right)^{2}-\nu_{\max}^{2}\right]}
\end{equation}
得到一个解（假定傅里叶级数收敛）。

\subsubsection{平面电荷（Planar Charge）}

举一个极简例子：设有一片静止的、单位电荷密度的平面薄片，平分 $\text{T}^{4}$ 黎曼宇宙，位于 $yz$ 平面内。此时四维势时间分量的源就是沿 $x$ 坐标的一维狄拉克 $\delta$ 函数。由于一切都只依赖 $x$，我们把其它三个维度从傅里叶系数下标与积分中略去，并简单地把 $L^{x}$ 记作 $L$。

源的（非零）傅里叶系数为
\begin{align}
h_{0} &= \frac{\sqrt2}{L},\\
h_{i} &= \frac{2}{L}, && i<0.
\end{align}
这样的源只有在 $\nu_{\max}L$ 不是整数时才可能实现，故我们假定如此。四维势时间分量的解的非零傅里叶系数为
\begin{align}
a_{0} &= -\frac{1}{2\sqrt2\,\pi^{2} L\,\nu_{\max}^{2}},\\
a_{i} &= \frac{L}{2\pi^{2}(i^{2}-L^{2}\nu_{\max}^{2})}, && i<0.
\end{align}

与其去显式求和这个傅里叶级数，我们改用另一种方法求解。利用问题的对称性以及黎曼版的推广高斯定律，很容易确定：在未加任何边界条件时，与单位平面电荷关联的四维势为
\begin{equation}
A_{t,\text{src}} = -\frac{\sin(\wms|x|)}{2\wms}.
\end{equation}
但还存在一个具有相同对称性的无源解，我们可以任意倍地加上它：
\begin{equation}
A_{t,\text{nsrc}} = \frac{\cos(\wms x)}{2\wms}.
\end{equation}
这两个函数都是 $x$ 的偶函数（即对所有 $x$，在 $\pm x$ 处取值相同），所以任何解在 $x=\pm L/2$ 处都连续。但偶函数在 $\pm x$ 处的导数\emph{相反}，故只有当该处导数为零时，解才能在 $x=\pm L/2$ 处光滑地与自己衔接。通过调节通解 $A_{t,\text{src}}+C\,A_{t,\text{nsrc}}$ 中的常数 $C$，我们可以保证在 $x=\pm L/2$ 处导数为零。结果化简为
\begin{equation}
A_{t,\text{bc}} = -\frac{\cos\!\bigl(\pi\nu_{\max}(L-2|x|)\bigr)}{4\pi\nu_{\max}\sin(\pi\nu_{\max}L)}.
\end{equation}
$A_{t,\text{bc}}$ 的傅里叶系数恰是我们上面写出的那些，故两种方法一致。这个解所描述的是势中的一个相位移动，它让势能光滑地绕过宇宙一圈，同时在平面电荷处保留恰到好处的间断，以满足该处的推广高斯定律。

\subsubsection{线电荷（Linear Charge）}

设有位于 $\text{T}^{4}$ 黎曼宇宙 $z$ 轴上的一条静止的、单位电荷密度的线。此时四维势时间分量的源是 $x$ 与 $y$ 坐标上的狄拉克 $\delta$ 函数。我们把 $z$ 与 $t$ 坐标从傅里叶系数中略去，为简单起见设 $L^{x}=L^{y}=L$。源的非零傅里叶系数为
\begin{align}
h_{0,0} &= \frac{2}{L^{2}},\\
h_{i,0} = h_{0,i} &= \frac{2\sqrt2}{L^{2}}, && i<0,\\
h_{i,j} &= \frac{4}{L^{2}}, && i,j<0.
\end{align}
这样的源只有在 $L^{2}\nu_{\max}^{2}$ 不是两个整数平方和时才可能实现，故我们假定如此。四维势时间分量的解的非零傅里叶系数为
\begin{align}
a_{0,0} &= -\frac{1}{2\pi^{2} L^{2}\nu_{\max}^{2}},\\
a_{i,0} = a_{0,i} &= \frac{1}{\sqrt2\,\pi^{2}(i^{2}-L^{2}\nu_{\max}^{2})}, && i<0,\\
a_{i,j} &= \frac{1}{\pi^{2}(i^{2}+j^{2}-L^{2}\nu_{\max}^{2})}, && i,j<0.
\end{align}

\orthofig{images/04__019.png}{$\text{T}^{4}$ 宇宙中线电荷的势：整个宇宙（$L$ 仅为几个波长），边界效应显著}
\orthofig{images/04__020.png}{$\text{T}^{4}$ 宇宙中线电荷的势：大宇宙（宽约千倍）的局部，场已趋于径向对称}

可以显式算出一个指标上的求和，把傅里叶级数化简为对另一指标的求和。整个表达式没有闭式，但指标数减半会让结果在数值上容易处理得多：
\begin{align}
A_{t} &= \sum_{j=0}^{\infty} f_{-j}(0)\,\beta_{j}\!\left(\frac{x}{L}\right) f_{-j}\!\left(\frac{y}{L}\right),\\
\beta_{j}(u) &= \frac{\cosh\!\bigl(\alpha_{j}(1-2|u|)\bigr)}{2\alpha_{j}\sinh(\alpha_{j})},\\
\alpha_{j} &= \pi\sqrt{j^{2}-L^{2}\nu_{\max}^{2}}.
\end{align}
注意 $\alpha_{j}$ 起初是虚数——直到 $j/L$ 超过 $\nu_{\max}$——而在它为虚数期间，函数 $\beta_{j}(u)$ 是振荡的，因为虚数 $ix$ 的 $\cosh$ 就是 $x$ 的 $\cos$。

一旦 $\alpha_{j}$ 变为实数，$\beta_{j}(u)$ 便从 $u=0$ 处的正极大值单调下降到 $u=1/2$ 处（也是正的）的极小值，后者对应距源半个宇宙之遥的点。这种下降并非字面上的指数衰减——因为指数衰减绝不会平缓到一个极小值——但非常相似。因此这些非振荡项随离源距离迅速衰减。

右图给出垂直于电荷线的平面内零势的等值线，展示边界条件如何使场的形状发生畸变。[由于未画出非零等值线，此处没有关于场强的信息——这里等值线的间距基本上就是波长。] \textbf{上图}展示整个宇宙，其参数选择使 $L$ 仅为几个波长，效应非常显著。\textbf{下图}展示同一尺寸（以波长计）的区域，但此时它只是一个宽一千倍的宇宙中的一小部分，靠近电荷处场已开始变得更接近径向对称。所以，尽管看场为了满足边界条件而丧失径向对称很有趣，但在一个现实尺度的宇宙中——宽度至少约 $10^{30}$ 个波长——这些效应不太可能被经验探测到。

\subsubsection{线形交流电流（Linear Alternating Current）}

引入这些边界条件的最初动机，是为了避免高频波中的指数暴涨。我们已经看到：若 $\text{T}^{4}$ 宇宙中能够存在无源波，则它们保证不会超过波方程中出现的名义最大频率。于是显而易见的问题是：若在 $\text{T}^{4}$ 宇宙中有某种以高于最大频率的频率振荡的源，会发生什么？

\orthofig{images/04__021.png}{$\text{T}^{4}$ 宇宙中交流电流的磁势}

最容易分析的源是线形交流电流。若电流沿 $\text{T}^{4}$ 黎曼宇宙的 $z$ 轴流动，并以频率 $l_{\text{AC}}/L^{t}$（对某整数 $l_{\text{AC}}$）振荡，则源与解在时间方向上共享单一的傅里叶分量，我们可以把它分离出去，并以与上一问题几乎相同的方式处理解的空间依赖。差别在于：此前只含 $j^{2}$ 项的指标平方和中，现在会额外加上一个常数项 $l_{\text{AC}}^{2}$。同前，为简单起见设宇宙在各方向（包括所选的时间方向）宽度相同，均为 $L$。于是有
\begin{align}
A_{z} &= \cos\!\left(\frac{2\pi l_{\text{AC}} t}{L}\right)\sum_{j=0}^{\infty} f_{-j}(0)\,\beta_{j}\!\left(\frac{x}{L}\right) f_{-j}\!\left(\frac{y}{L}\right),\\
\beta_{j}(u) &= \frac{\cosh\!\bigl(\alpha_{j}(1-2|u|)\bigr)}{2\alpha_{j}\sinh(\alpha_{j})},\\
\alpha_{j} &= \pi\sqrt{j^{2}+l_{\text{AC}}^{2}-L^{2}\nu_{\max}^{2}}.
\end{align}
若电流振荡频率 $l_{\text{AC}}/L$ 超过 $\nu_{\max}$，则 $\alpha_{j}$ 定义中根号内的表达式始终为正，故对一切 $j$，$\alpha_{j}$ 都是实数。如上一节所讨论，当 $\alpha_{j}$ 为实数时，函数 $\beta_{j}(u)$ 以一种非常类似指数衰减的方式迅速下降，同时平缓下来、在绕过宇宙一半处达到零导数。

我们无法给出对 $j$ 的无限和的闭式，但图中展示了大数量项之和。可以看出，高频源伴随的场只在源附近很有限处显著，随距离下降的速度远快于频率低于 $\nu_{\max}$ 的线形交流电流周围的辐射场。

\subsubsection{柯西数据与预言（Cauchy Data and Predictions）}

在我们自身宇宙中，真空中的光由洛伦兹波方程支配；若我们知道某一时刻 $t_{0}$ 空间某区域 $\mathcal{R}$ 内电磁场的值及其时间导数，就能预言此后某段时间内的场值。当然，电磁波总能从侧面进入该区域，故随着时间从 $t_{0}$ 推进，我们能够预言的区域会以光速缩小；但原则上总会有一部分确定的时空，我们的初始数据能让我们预言其中场将是什么。

这类数据——某时刻空间某区域上函数的值及其时间导数——称为\textbf{柯西数据}（Cauchy data）。洛伦兹波方程（二阶双曲型微分方程）的本性就在于，我们可以用柯西数据得到其解在未来某段时间的取值。

另一个可用柯西数据的双曲型方程例子，是弹性弦小位移的波方程。设弦有限、两端固定，则若我们知道某时刻整条弦的位移及其时间导数，原则上就能预言弦此后运动的全部未来。更进一步，即便我们的知识仅限于弦的一部分，由于波以某一最大速度 $c_{\max}$ 传播，我们仍能自信地对一段区域作出预言，只不过该区域会从我们有数据的部分逐渐缩小，两端以速率 $c_{\max}$ 被“啃掉”。

与此相反，黎曼波方程是\emph{椭圆型}微分方程。要在某区域内求解椭圆型微分方程，通常需要知道区域\emph{整个边界}上解的值。我们自身宇宙中椭圆型微分方程的例子不涉及时空区域，而涉及空间区域。例如，固体材料中达到的平衡温度服从一个椭圆型微分方程——拉普拉斯方程——要确定材料某区域内的温度，一般需要知道该区域整个边界上的温度。若只告诉你一块铁立方体某个面上的温度——再加上从该面指向立方体内部的温度导数，即柯西数据——这\emph{不}是计算立方体内各处温度的可靠办法。

举例来说，设与我们有数据的那个面相对的面上布满了高低温交替、间距很小的条纹。我们的数据也许只描述这些条纹的一个极其微弱、被抹平的版本。温度从我们的面推进到对面时会伴随温差的指数式增长，从而极大地放大数据中任何不精确之处，以致仅有这些被抹平的条纹及其导数，对面温度的确切取值就变得极难把握。但若反过来\emph{给定}立方体每个面上的温度，则在该区域内插值出满足拉普拉斯方程的温度分布，将可靠得多。

在一个无限的黎曼宇宙中，由柯西数据对黎曼波方程作预言，其难度就相当于只凭一个面的数据去计算立方体中的温度。既然方程是椭圆型的，我们或许会得出结论：只能对它的解作\emph{事后推断}（postdiction）——即收集某空间区域内场初值、以及过一段时间后的终值、连同该时间段内区域边界上发生的情况等数据，再用相关四维空间边界上的所有这些信息，事后来计算区域内部场的时间演化。这样的局面能让我们检验物理定律，却会让我们极难预见和筹谋未来。

但在诸如 $\text{T}^{4}$ 这样的\emph{有限}黎曼宇宙中，情况没那么糟。$\text{T}^{4}$ 中的无源波只有有限多个傅里叶基函数能对总波作出贡献，因此只要我们能够确定它们全部的系数，就能知道该波的全部历史。若再纳入一个源——它本身也应服从一条同样一般形式的方程——问题会更复杂，但原理相同。

为简单起见，我们处理一个无源标量波。设我们知道某一时刻整个空间上该波的值及其时间导数。我们选择坐标，使有数据的时刻为 $t=0$；当然，“时间”可以是环面可以绕行的四个方向中的任意一个。

设某傅里叶基函数 $f_{i,j,k,l}$ 满足无源波方程。若 $l\le 0$，该函数依赖时间的因子 $f_{l}(t/L^{t})$ 是余弦或常数函数，因此在 $t=0$ 处非零，于是只要对 $t=0$ 的数据作三维傅里叶分析，就能确定 $f_{i,j,k,l}$ 的系数。若 $l>0$，依赖时间的因子是正弦，故它在 $t=0$ 处为零；但它的时间导数在 $t=0$ 处非零，于是可以由我们已有的 $t=0$ 时间导数数据作傅里叶分析来确定其系数。因此在数据与其时间导数之间，我们能算出对无源波有贡献的每个基函数的系数，从而能算出该波在任何时刻（将来或过去）的值。

当然，指望黎曼宇宙中有人掌握整个宇宙的电磁场信息是荒唐的。但反过来说，当我们在自身宇宙中预言未来五分钟会发生什么时，我们从不曾拥有周围五光分（约九千万公里）范围内的完备信息。尽管如此，我们至今仍设法检验了科学理论，并能足够好地预言未来以存活下来。事实是，我们身处一个足够有序、平静的时间和地点，以致\emph{通常}可以假定周围最重要的电磁辐射源是像太阳这样近旁、行为已知且相当可预测的物体。物理定律\emph{允许}来自未知源的、会让我们完全措手不及的突然巨量辐射涌入，这并未摧毁我们做科学或规划未来的能力。

在 Max Tegmark 的经典论文《“万有理论”是否仅仅是终极集合论？》中，Tegmark 提出：支配一个没有类时维度之宇宙的椭圆型偏微分方程将使其不可能作出预言，从而令他所称的“自省子系统”（self-aware substructures）很难有效运作。但在一个\emph{有限}黎曼宇宙中，若存在局部环境相对平静有序的区域——正是我们自身演化与繁荣所依赖的那类条件——那么，为了作出预言而严格需要横跨整个宇宙的柯西数据这一点，就不会比在我们自身宇宙中、为了知道未来五分钟将发生什么而严格需要一个半径九千万公里区域的柯西数据更能成为决定生命能达到何种成就的因素。

\subsection{4-球面（The 4-sphere）}

五维空间中一个实心超球面的边界，是一个有限、无边界的四维空间，称为 4-球面，或 $\text{S}^{4}$。4-球面无须嵌入任何更高维空间，也无须具有均匀曲率；但为简单起见，我们考察一个具有此拓扑的黎曼宇宙，它\emph{确实}具备一个嵌入平坦五维空间的 4-球面会从该空间继承的全部几何性质。若取超球面的半径为 $R$，这便固定了总体积（四维体积）为
\begin{equation}
V = \frac{8}{3}\pi^{2}R^{4},
\end{equation}
并把 4-球面内任何测地线的最大长度固定为 $2\pi R$。里奇标量曲率（它衡量空间内一个实心球的体积随半径增大而不如欧氏空间中增长得快这一程度）在每一点处都是 $12/R^{2}$。

$\text{S}^{4}$ 作为模型宇宙的好处在于它比 $\text{T}^{4}$ 更对称。考察 $\text{S}^{4}$ 中保持某点不动的对称性，它们恰好是与欧氏四维空间中相同的群 $O(4)$。而代替欧氏空间平移的，我们只需把群扩展到 $O(5)$。

代价是必须处理一个弯曲的四维空间。与 $\text{T}^{4}$ 不同，具有 $\text{S}^{4}$ 拓扑的空间不可能处处完全平坦。为什么？偶数 $n$ 的 $\text{S}^{n}$ 的欧拉数（Euler characteristic）总是 2（这可通过数一数超立方体的各部分简单地证明）。广义高斯-博内定理（Generalised Gauss-Bonnet Theorem）把空间曲率相关函数的一个积分与欧拉数等同起来，而若曲率为零，该积分就为零——与欧拉数的已知值矛盾。

我们可以在 $\text{S}^{4}$ 上放置一种极坐标，含四个角坐标：
\begin{equation}
0\le\xi\le\pi,\qquad 0\le\psi\le\pi,\qquad 0\le\theta\le\pi,\qquad 0\le\phi\le2\pi,
\end{equation}
它们把平坦五维空间中半径 $R$ 的 4-球面上一点参数化为
\begin{equation}
\bigl(R\cos\xi,\ R\sin\xi\cos\psi,\ R\sin\xi\sin\psi\cos\theta,\ R\sin\xi\sin\psi\sin\theta\cos\phi,\ R\sin\xi\sin\psi\sin\theta\sin\phi\bigr).
\end{equation}
用这些坐标，度规是对角的，非零分量为
\begin{align}
g_{\xi\xi} &= R^{2},\\
g_{\psi\psi} &= R^{2}\sin^{2}\xi,\\
g_{\theta\theta} &= R^{2}\sin^{2}\xi\sin^{2}\psi,\\
g_{\phi\phi} &= R^{2}\sin^{2}\xi\sin^{2}\psi\sin^{2}\theta,
\end{align}
由此得到度规行列式的平方根
\begin{equation}
\sqrt{|g|} = R^{4}\sin^{3}\xi\,\sin^{2}\psi\,\sin\theta.
\end{equation}

正如 $\text{T}^{4}$ 上标量函数可写成傅里叶级数，$\text{S}^{4}$ 上行为良好的标量函数可以展开为四维球谐函数（spherical harmonics）之和：
\begin{align}
Y^{m}_{j,k,l}(\xi,\psi,\theta,\phi) &= \Phi_{m}(\phi)\,\Theta^{m}_{j}(\theta)\,\Psi^{j}_{k}(\psi)\,\Xi^{k}_{l}(\xi)\big/R^{4},\\
\Phi_{m}(\phi) &= \frac{\sin(m\phi)}{\sqrt\pi}, && m>0,\\
\Phi_{m}(\phi) &= \frac{\cos(m\phi)}{\sqrt\pi}, && m<0,\\
\Phi_{0}(\phi) &= \frac{1}{\sqrt{2\pi}},\\
\Theta^{m}_{j}(\theta) &= \sqrt{\frac{(j+\tfrac12)(j-|m|)!}{(j+|m|)!}}\,P^{|m|}_{j}(\cos\theta),\\
\Psi^{j}_{k}(\psi) &= \sqrt{\frac{(k+1)(k+j+1)!}{(k-j)!}}\,\frac{P^{-j-\frac12}_{k+\frac12}(\cos\psi)}{\sqrt{\sin\psi}},\\
\Xi^{k}_{l}(\xi) &= \sqrt{\frac{(l+3/2)(l-k)!}{(l+k+2)!}}\,\frac{P^{k+1}_{l+1}(\cos\xi)}{\sin\xi}.
\end{align}
这里 $P$ 是第一类连带勒让德函数（associated Legendre function of the first kind）。指标 $m,j,k,l$ 为整数，受如下约束：
\begin{equation}
0\le|m|\le j\le k\le l.
\end{equation}

\orthofig{images/04__022.png}{球谐函数中各因子的图像}

函数 $\Phi_{m}(\phi)$ 具有简单的三角形式，但另外三个坐标的函数长什么样？它们都遵循大致相同的模式：当其上方指标取其最大可能值时，它们从坐标为零时的零值，上升到 $\pi/2$ 处的单个极大或极小，再在坐标达到 $\pi$ 时回落为零。随着该指标的值下降，它们会从零回到零的过程中多出一个极值。当指标降到零时，极值的计数照旧递增，只不过此时函数在其取值范围的端点处不再为零。

对给定的 $l$，这些四维球谐函数的总数可由对其它指标的约束求出：
\begin{equation}
N(l) = \frac{(l+1)(l+2)(2l+3)}{6}.
\end{equation}
指标不同的球谐函数彼此正交，即它们的乘积在 $\text{S}^{4}$ 上按体积 $\sqrt{|g|}$ 加权积分后为零。我们这里纳入的因子还确保每个球谐函数平方的积分为 1。

哪些球谐函数满足我们在弯曲空间一节导出的球面上的无源黎曼标量波方程？不难证明球谐函数是拉普拉斯-贝尔特拉米算符的本征函数，即
\begin{equation}
\operatorname{div}\operatorname{grad} Y^{m}_{j,k,l} = \left[-\frac{l(l+3)}{R^{2}}\right] Y^{m}_{j,k,l}.
\end{equation}
因此，若 $\wms^{2}R^{2}$ 是形如 $l(l+3)$ 的整数，则该 $l$ 对应的全部 $N(l)$ 个球谐函数都满足无源方程。若 $\wms^{2}R^{2}$ 不是这种形式的整数，则不存在无源解。于是我们的处境与 $\text{T}^{4}$ 上非常相似：一般情况下，最大频率与宇宙尺度取一般值时不允许存在无源解；但若几何容许无源解存在，它们也由有限多个模式构成。这里我们能非常容易地数出模式数，无须操心 $\text{T}^{4}$ 所需的那类数论细节。由于当 $\wms^{2}R^{2}=l(l+3)$ 时有 $N(l)$ 个模式，对大 $l$ 有
\begin{align}
l &\approx \wms R,\\
N(l) &\approx N(\wms R) \approx \frac{(\wms R)^{3}}{3}.
\end{align}
当然，在任何尺度哪怕大致可与我们的可观测宇宙相比的宇宙中，这都是一个极大的数目。但事实上 $\text{S}^{4}$ 的对称性意味着：若存在无源解，就存在局部看起来像传播矢量为\emph{任意}值的平面波的解，而不只是一个虽大却离散的选择集。对位于 $\xi=\psi=\theta=\pi/2$、$\phi$ 坐标为任意值的观测者，考察球谐函数 $Y^{l}_{l,l,l}$，其中 $l$ 是使 $l(l+3)=\wms^{2}R^{2}$ 的那个特定整数。观测者将处于 $\xi$、$\psi$、$\theta$ 各函数的极宽而平坦的极值处，而函数在 $\phi$ 方向按 $\cos(l\phi)\approx\cos(\wms R\phi)$ 变化，局部看起来恰是我们为欧氏四维空间描述过的那种平面波。但对观测者位置的任意选择以及任意传播矢量选择，我们都能简单地选取满足所述条件的坐标，并在那些坐标下构造同样的解。

若考虑带源 $H$ 的 RSW 方程，并把源的球谐系数写作 $h^{m}_{j,k,l}$、解 $A$ 的写作 $a^{m}_{j,k,l}$，则有
\begin{equation}
a^{m}_{j,k,l}\bigl[l(l+3)-\wms^{2}R^{2}\bigr] = h^{m}_{j,k,l}\,R^{2}.
\end{equation}
若存在某个 $l$ 使 $l(l+3)=\wms^{2}R^{2}$，则源不能包含该 $l$ 值的任何球谐函数，而解可以任意多地包含这些球谐函数。对其它 $l$ 值，源的系数确定解的系数：
\begin{equation}
a^{m}_{j,k,l} = \frac{h^{m}_{j,k,l}\,R^{2}}{l(l+3)-\wms^{2}R^{2}}.
\end{equation}

\subsubsection{4-球面上的格林函数（A Green's Function for the 4-Sphere）}

设有 $\text{S}^{4}$ 上一个点状、$\delta$ 函数式的源突发。与它关联的解是什么？我们预期它看起来有点像我们此前在欧氏四维空间中为电偶极-电荷突然脉动找到的格林函数，后者正比于 $Y_{1}(\wms s)/s$，其中 $Y_{1}$ 是第二类贝塞尔函数。

为简单起见，我们只限于标量波方程。若把源放在坐标系的极点，即 $\xi=\psi=\theta=0$ 而 $\phi$ 无定义处（正如地球南北极点处经度无定义），则唯一非零的球谐系数将具有 $m=j=k=0$，因为其它取值都会使球谐函数 $Y^{m}_{j,k,l}$ 在极点处为零。于是非零系数为
\begin{align}
h^{0}_{0,0,l} &= -\frac{\sqrt{(l+1)(l+2)(2l+3)}}{4\pi},\\
a^{0}_{0,0,l} &= -\frac{R^{2}\sqrt{(l+1)(l+2)(2l+3)}}{4\pi\bigl[l(l+3)-\wms^{2}R^{2}\bigr]}.
\end{align}
我们假定不存在使 $l(l+3)=\wms^{2}R^{2}$ 的整数 $l$。解为
\begin{equation}
A(\xi) = -R^{2}\sum_{l=0}^{\infty}\frac{(2l+3)\,P^{1}_{l+1}(\cos\xi)}{8\pi^{2}\bigl[l(l+3)-\wms^{2}R^{2}\bigr]\sin\xi}.
\end{equation}

\orthofig{images/04__023.png}{4-球面上的格林函数}

右图给出这个求和的一个数值近似。该函数在其定义域的前半部分行为正如我们所料：随离源距离振荡并衰减；但随后在接近对极点时发生一次令人不安的复现（resurgence）。这是 $\text{S}^{4}$ 对称性的产物；在具有相同拓扑但较不均匀的空间中，这一效应会弱得多。

值得注意的是，即便有这种完美对称性，对极点处的场也是\emph{有限}的，不同于源点处的场。$\xi=\pi$ 处的和可以显式算出：
\begin{equation}
A(\pi) = -\frac{R^{2}\bigl(\wms^{2}R^{2}+2\bigr)}{16\pi\cos\!\left(\dfrac{\pi}{2}\sqrt{4\wms^{2}R^{2}+9}\right)}.
\end{equation}
这里分母中的余弦只有在 $\wms^{2}R^{2}$ 违反整数假定时才可能为零，故此处的场是有限的。

\subsubsection{柯西数据与预言（Cauchy Data and Predictions）}

对 $\text{T}^{4}$ 宇宙，我们发现：若有单一时刻横跨宇宙宽度的柯西数据（其中时间轴可以是绕 4-环面一周的四个坐标中的任意一个），我们就能确定一个自由波的有限多个系数的值，从而重构其全部历史。

对 $\text{S}^{4}$，我们可以在任何“大 3-球面”（great 3-sphere），即任何半径 $R$ 的 3-球面上，用柯西数据做同样的事。假定几何容许无源波，则满足无源波方程的全部球谐函数 $Y^{m}_{j,k,l}(\xi,\psi,\theta,\phi)$ 都共享同一个 $l$ 值。我们选取一个坐标，使有数据的那个 3-球面上 $\xi=\pi/2$。对于那些在 $\xi=\pi/2$ 处取极大或极小的球谐函数，可以由该 3-球面上的场值求出其系数；而那些在该处为零的球谐函数，其 $\xi$ 方向导数将在该处取极大或极小，于是可以由场的导数求出其系数。因此，由 3-球面上的柯西数据，我们可以重构解的全部历史。

若我们的数据在一个更小的 3-球面上，即描述为超曲面 $\xi=\xi_{0}$（对某个 $\xi_{0}<\pi/2$），情况又如何？只要我们确实知道 $\xi_{0}$ 的值，因子 $\Xi^{k}_{l}(\xi)$ 与 $\partial_{\xi}\Xi^{k}_{l}(\xi)$ 在 3-球面上就是已知量（而且它们绝不会同时为零），因此原则上我们总应能算出解的全部系数。

这引出如下奇特的观察：\emph{原则上}我们甚至能从最小 3-球面上的柯西数据重构整个解。毕竟，这样一个 3-球面是\emph{两个}有限区域的边界：按通常理解其内部，以及 $\text{S}^{4}$ 宇宙的其余部分——正如北极圈既是北极周围区域的边界，也是地球表面其余部分的边界。但在实践中，对远小于 $\pi/2$ 的 $\xi$，$\Xi^{k}_{l}(\xi)$ 的值相比其在 $\pi/2$ 处的取值会变得\emph{极其}小，而且球谐函数中其它各因子的峰也会越来越密，以致从小 3-球面向外外推到整个宇宙，将对数据的精度提出高得难以承受的要求。

\section{参考文献（References）}

\begin{itemize}
  \item[{[1]}] \emph{Classical Electrodynamics}，John David Jackson，John Wiley \& Sons，1999 年。第 12.7 节给出普通电磁学的拉格朗日量，第 12.8 节给出\emph{洛伦兹} Proca 电动力学的拉格朗日量。（注意 Jackson 使用了与我们不同的单位制，以及洛伦兹度规的 $(+---)$ 号差，因此比较这些公式时需要小心。）
  \item[{[2]}] \emph{Gravitation}，Charles Misner、Kip Thorne 与 John Wheeler，W.H. Freeman，San Francisco，1973 年。第 21.3 节。
\end{itemize}

\chapter{黎曼热力学}

本章把统计力学的基本框架移植到勾股定理符号反转的黎曼宇宙（Riemannian universe）。由于动量空间中粒子可用的态数被质量频率 $\wms=mc^{2}/\hbar$ 所截断，黎曼粒子的相空间体积随能量先增后减，由此自然引入"比无限热更热"的负温度概念。以此为起点，我们计算黎曼黑体谱（同时给出纯经典与量子两版）、单粒子配分函数，以及理想气体——固定粒子数的稀薄气体与能产生粒子对的费米子气体——的能量、压强、熵与粒子分布。

\section{比无限热更热（Hotter Than Infinitely Hot）}

一个粒子有多少种方式拥有给定的能量 $E$？能量固定了其动量矢量 $\bp$ 的大小 $p$，但该矢量仍可指向三维空间中的任意方向，因而自由地落在某个球面上。若我们并不精确知道粒子能量，只知道它介于 $E$ 与 $E+\delta E$ 之间（$\delta E$ 很小），那么 $p$ 将介于由这些能量所确定的两值 $p$ 与 $p+\delta p$ 之间，动量矢量便会落在一个球壳内。当然该球壳内有无限多个不同矢量，所以问这个能量区间内粒子"有多少种方式"在经典力学中并不恰当；量子力学把情形变成真正离散的有限可能性。但在经典物理中，考虑粒子所在区域（动量空间、普通空间或两者）的\emph{体积}是有用的。

\orthofig{images/05__001.png}{动量空间中的球壳}

若把粒子的状态视为某个抽象空间（相空间，phase space）中的一点，其坐标同时包含粒子位置与三个方向上的动量，那么"能量介于 $E$ 与 $E+\delta E$ 之间"便意味着状态落在相空间的某个区域内。若假设粒子无论能量如何都被限制在体积为 $V$ 的固定普通空间区域，则它在相空间中所占的体积取决于与其能量相容的可能动量矢量——即刚刚描述的球壳内的那些矢量。

为什么我们关心相空间中的体积，而不是别的状态描述？原因是一条应用数学结论——刘维尔定理（Liouville's Theorem），它描述系统状态在相空间中某一点的概率随时间的变化。借助刘维尔定理可以证明：若我们有大量系统副本，其状态均匀散布于相空间，则随着时间推移、描述各系统的点四处运动，它们仍将以同样的均匀方式分布。若另选某种随意的状态描述，这一性质并不成立；只有状态在\emph{相空间}中运动的特定数学结构才使其为真。

当我们要以统计方式处理系统时，这一点很重要。对一个含数十亿粒子的复杂系统（例如把热气体与冷固体接触），我们不会去追踪每个粒子的位置与动量，而是等待一切停止变化、沉入平衡。一旦系统达到平衡，刘维尔定理告诉我们，可把它视为在与其已知信息相容的相空间任一位置\emph{等可能}。若这听起来理所当然，请记住：许多同样"理所当然"的说法其实是错的——例如，一个系统通常\emph{并非}在我们已知的能量范围内的各个取值上等可能。

\orthofig{images/05__002.png}{洛伦兹与黎曼相空间体积}

回到单粒子例子，我们来计算：若已知粒子能量介于 $E$ 与 $E+\delta E$ 之间、位置限于体积为 $V$ 的普通空间区域，则它必落在的相空间区域体积为多少。

在\emph{洛伦兹}物理中，自由粒子的总能量大于等于其静止质量 $m$，能量与动量的关系为
\begin{equation}
p^{2}=E^{2}-m^{2},
\end{equation}
由此可得
\begin{equation}
\frac{dp}{dE}=\frac{E}{p}=\frac{E}{\sqrt{E^{2}-m^{2}}}.
\end{equation}
若粒子能量落在 $E$ 到 $E+\delta E$ 的小范围内，则粒子状态所在相空间区域的体积为
\begin{align}
\Omega_{\Lor}
&= 4\pi p^{2}\,\delta p\,V \nonumber\\
&= 4\pi p^{2}\frac{dp}{dE}\,\delta E\,V \nonumber\\
&= 4\pi\sqrt{E^{2}-m^{2}}\,E\,\delta E\,V .
\end{align}

在\emph{黎曼}物理中，自由粒子的总能量介于零与该粒子的静止质量 $m$ 之间，我们有
\begin{equation}
p^{2}=m^{2}-E^{2},\qquad
\frac{dp}{dE}=-\frac{E}{p}=-\frac{E}{\sqrt{m^{2}-E^{2}}}.
\end{equation}
动量随能量增大而减小这一事实本身即很重要；为求相空间体积，我们取 $\delta p=(dp/dE)\,\delta E$ 的绝对值，得
\begin{align}
\Omega_{\Riem}
&= 4\pi p^{2}\left|\frac{dp}{dE}\right|\delta E\,V \nonumber\\
&= 4\pi\sqrt{m^{2}-E^{2}}\,E\,\delta E\,V .
\end{align}

如 $\Omega$ 的曲线所示，这两个函数非常不同。但值得注意的是，当能量接近静止质量 $m$ 时，它们几乎是彼此的镜像。这是因为在低速下，两个宇宙中的粒子实际上都服从牛顿力学；黎曼宇宙的唯一区别是总能量随动能增加而减小。

最显著的差别在于 $\Omega_{\Lor}$ 总是随能量增大。无论粒子能量为何，只要再加一点能量，总会使可供它使用的相空间体积更大。与此相反，$\Omega_{\Riem}$ 只在某个特定能量处达到最大可用相空间体积。计算两种情形下相空间体积对能量的变化率，得到
\begin{align}
\frac{d\Omega_{\Lor}}{dE}&=4\pi\,\delta E\,V\,\frac{2E^{2}-m^{2}}{\sqrt{E^{2}-m^{2}}},\\
\frac{d\Omega_{\Riem}}{dE}&=4\pi\,\delta E\,V\,\frac{m^{2}-2E^{2}}{\sqrt{m^{2}-E^{2}}},
\end{align}
故黎曼相空间体积在 $E=m/\sqrt2$ 处取极大。也许有人以为最大的相空间体积应联系于最大的可能动量 $p=m$，它发生在 $E=0$ 处；但尽管那使动量空间中的球壳半径最大，我们也不能忘记球壳厚度在变化。在 $E=0$ 处 $dp/dE$ 也为零，故 $\delta p$（球壳厚度）为零。实际的极大体积是球壳半径与其厚度之间折中的结果。

若改用归一化速度 $v_0$，黎曼的能量-速度关系可写为 $E=m/\sqrt{1+v_0^{2}}$（与 $E^{2}+p^{2}=m^{2}$ 相容），于是相空间体积取极大的能量 $E=m/\sqrt2$ 恰好对应 $v_0=1$。

相空间体积随能量这种变化有何意义？为理解它，设有两个系统，已知其能量 $E_{1}$、$E_{2}$，以及各自相空间体积随能量的变化 $\Omega_{1}(E_{1})$、$\Omega_{2}(E_{2})$。若把两系统看作一个系统，可定义一个新的相空间，其坐标包含两系统所有粒子的位置与动量。该相空间中被复合系统占据的体积就是各自体积之积：
\begin{equation}
\Omega_{total}(E_{1},E_{2})=\Omega_{1}(E_{1})\,\Omega_{2}(E_{2}).
\end{equation}
现在假定允许能量在两系统间流动，总能量 $E_{1}+E_{2}$ 保持不变，但各系统的能量可以变化。若有一定量的能量 $Q$ 从系统 1 流向系统 2，则
\begin{equation}
\Omega_{total}(E_{1}-Q,\,E_{2}+Q)=\Omega_{1}(E_{1}-Q)\,\Omega_{2}(E_{2}+Q),
\end{equation}
它对总相空间体积的影响由下式描述：
\begin{equation}
\frac{d\Omega_{total}(E_{1}-Q,E_{2}+Q)}{dQ}
=-\Bigl(\frac{d\Omega_{1}}{dE_{1}}\Bigr)\Omega_{2}+\Omega_{1}\Bigl(\frac{d\Omega_{2}}{dE_{2}}\Bigr),
\end{equation}
这里用到了"乘积的导数等于依次对每个因子求导再相加"的法则。

若能量从系统 1 流向系统 2，总相空间体积 $\Omega_{total}$ 会上升、下降还是保持不变？若为每个子系统定义一个新量——它的\emph{温度}（temperature）$T$，就更容易判断：
\begin{equation}
T_{1}=\frac{\Omega_{1}}{d\Omega_{1}/dE_{1}},\qquad
T_{2}=\frac{\Omega_{2}}{d\Omega_{2}/dE_{2}}.
\end{equation}
只要知道 $T_{1}$ 与 $T_{2}$ 的关系，我们便可把它们分母中的导数乘过去，判断 $d\Omega_{total}/dQ$ 的符号。先假定 $d\Omega_{1}/dE_{1}$ 与 $d\Omega_{2}/dE_{2}$ \emph{都为正}；由于 $\Omega$ 本身恒正，这等价于 $T_{1}$ 与 $T_{2}$ 都为正。于是有：

\begin{itemize}
  \item 若 $T_{1}>T_{2}>0$，则 $\Omega_{1}(d\Omega_{2}/dE_{2})>\Omega_{2}(d\Omega_{1}/dE_{1})$，故 $d\Omega_{total}/dQ$ 为正（即能量从系统 1 流向系统 2 使总相空间体积增大）；
  \item 若 $T_{2}>T_{1}>0$，则 $\Omega_{1}(d\Omega_{2}/dE_{2})<\Omega_{2}(d\Omega_{1}/dE_{1})$，故 $d\Omega_{total}/dQ$ 为负。
\end{itemize}

也就是说，能量从高温系统流向低温系统，与热力学第二定律一致。

你也许听说过温度是一种只对含大量粒子的系统才有意义的"涌现现象"，但严格说来这是不对的：我们按上述方式定义的温度对任何系统都成立，哪怕单个粒子。真正依赖于粒子数众多的是与温差相伴的能量流动方向的近乎确定性。若两个孤立的粒子相互作用，偶然出现较热的粒子实际上从较冷的粒子吸取能量并不奇怪；但在一个含 $10^{23}$ 个粒子的系统中，能量从冷流向热将是不可能的（概率小到天文数字）。

拥有\emph{负}温度会有什么后果？我们为两个正温度系统所作、关于能量流动方向的论证，依赖于把一个不等式乘以 $(d\Omega_{1}/dE_{1})(d\Omega_{2}/dE_{2})$。若两个系统\emph{都}具有负温度，这两个导数的乘积仍然为正，故完全同样的论证依然成立：能量仍然从温度较高者流向温度较低者。

这意味着，若两个都是负温度的系统相互作用……把它们合并后总相空间体积仍在"高 $T$ 到低 $T$"的流动方向上增大。对于黎曼粒子，$\Omega_{\Riem}$ 在 $E=m/\sqrt2$ 处取极大，故 $E<m/\sqrt2$ 时 $d\Omega/dE>0$（$T>0$），而 $E>m/\sqrt2$ 时 $d\Omega/dE<0$（$T<0$）。于是温度随能量的变化为：能量从 $0$ 增大到 $m/\sqrt2$ 时，$T$ 由 $0^{+}$ 升至 $+\infty$；越过 $E=m/\sqrt2$ 后温度变为负值，从 $-\infty$ 一路升到 $E\to m$ 时的 $0^{-}$。因此负温度并不代表"更冷"，而是比无限热\emph{更热}——这正是本节标题"比无限热更热"的含义。

\begin{notebox}[比无限热更热]
在黎曼宇宙中，相空间体积 $\Omega_{\Riem}(E)=4\pi\sqrt{m^{2}-E^{2}}\,E\,\delta E\,V$ 在 $E=m/\sqrt2$ 处达到极大。其后果是温度由 $T=\Omega/(d\Omega/dE)$ 给出时，能量高于 $m/\sqrt2$ 的部分对应\emph{负}温度：系统越接近最大可能能量 $E=m$，温度越接近 $0^{-}$，而其能量高于"无限热"所对应的状态。负温度系统会向任何正温度系统放能，故比无限热更热。
\end{notebox}

关于本主题的额外材料，提供给了不介意稍高数学门槛的读者，见以下各节。

\section{黎曼黑体谱（The Riemannian Black Body Spectrum）}

黑体谱（black body spectrum）描述热辐射中能量在不同频率的光之间的分布——热辐射是物体热能自由地与光交换时产生的辐射。我们不打算计算黎曼宇宙中各类物质与电磁场交换能量的性质；下文只假设一个能量可以在不同频率的光之间自由流动的情境，如同分析我们自身宇宙中同类情形时通常假定的那种完美发射体与吸收体所提供的情境。

我们从某个封闭腔（cavity）内的电磁场入手，腔壁完全不透明，场与周围材料处于热平衡。沿用 Reif\textsuperscript{[1]} 对洛伦兹情形的处理，我们假定施加在场上的精确边界条件无关紧要。为简单起见，考察一个边长为 $L$、体积为 $V=L^{3}$ 的立方腔，并施加周期性边界条件：立方体任一面上场的取值，与其相对面上的对应点相同。

满足这些边界条件的黎曼标量波方程（Riemannian Scalar Wave equation）的任一解，都可写成如下形式平面波的线性组合的实部：
\begin{align}
f_{i,j,k}(x,y,z,t)&=\exp\!\bigl(-\ii\,\bk_{i,j,k}\cdot(x,y,z,t)\bigr),\\
\bk_{i,j,k}&=\frac{2\pi}{L}\bigl(i\,\be_{x}+j\,\be_{y}+k\,\be_{z}\bigr)+\omega_{i,j,k}\,\be_{t},\\
\omega_{i,j,k}&=\sqrt{\wms^{2}-\frac{4\pi^{2}}{L^{2}}\bigl(i^{2}+j^{2}+k^{2}\bigr)},
\end{align}
其中 $i$、$j$、$k$ 为整数。它们可取正、负或零，但 $\omega_{i,j,k}$ 定义中根号内的表达式必须为正，这把可能取值限制为有限多种，对应于空间频率不得超过 $\nu_{max}$ 的限制。

当然我们要的是矢量解而非标量解。对每一组整数，都能找到三个互相正交、且都垂直于传播矢量 $\bk_{i,j,k}$ 的单位矢量，从而给出每个模式的三种偏振。显式地：
\begin{align}
\mathbf{f}^{(1)}_{i,j,k}&=\frac{1}{\wms}\Bigl[\frac{2\pi}{L}(i\,\be_{t}-j\,\be_{z}+k\,\be_{y})-\omega_{i,j,k}\,\be_{x}\Bigr]f_{i,j,k},\\
\mathbf{f}^{(2)}_{i,j,k}&=\frac{1}{\wms}\Bigl[\frac{2\pi}{L}(-i\,\be_{z}-j\,\be_{t}+k\,\be_{x})+\omega_{i,j,k}\,\be_{y}\Bigr]f_{i,j,k},\\
\mathbf{f}^{(3)}_{i,j,k}&=\frac{1}{\wms}\Bigl[\frac{2\pi}{L}(-i\,\be_{y}+j\,\be_{x}+k\,\be_{t})-\omega_{i,j,k}\,\be_{z}\Bigr]f_{i,j,k},
\end{align}
其中每个式子里都省略了 $(x,y,z,t)$ 变量。

现在，令电磁场的四维势 $\bA$ 等于这些矢量模式的线性组合：
\begin{equation}
\bA=A_{(1)}\mathbf{f}^{(1)}_{i,j,k}+A_{(2)}\mathbf{f}^{(2)}_{i,j,k}+A_{(3)}\mathbf{f}^{(3)}_{i,j,k},
\end{equation}
则利用我们此前推出的平面波能量密度公式可以证明，立方腔内的总能量为
\begin{equation}
U=\half V\bigl(A_{(1)}^{2}+A_{(2)}^{2}+A_{(3)}^{2}\bigr)\omega_{i,j,k}^{2}.
\end{equation}
若向 $\bA$ 添加其它 $i,j,k$ 取值的模式，能量简单地相加；这是不同波数三角函数正交性的结果。因此完全一般的能量表达式就是这类平方项之和，包含每个模式的系数。

[若愿意，也可以改用\emph{驻波}（standing wave）模式而非行波模式。对每一对反向行进的行波模式 $\mathbf{f}^{(p)}_{i,j,k}$ 与 $\mathbf{f}^{(p)}_{-i,-j,-k}$，可构造两个驻波模式 $(\mathbf{f}^{(p)}_{i,j,k}\pm\mathbf{f}^{(p)}_{-i,-j,-k})/\sqrt2$。例外是 $i=j=k=0$，它本身已是驻波。这些模式的总能量与行波模式的相同。]

角频率介于 $\omega$ 与 $\omega+d\omega$ 之间的模式有多少个？若假定腔足够大，$i$、$j$、$k$ 取值很大，可以把它们当作连续量而得到很好的近似。令 $r^{2}=i^{2}+j^{2}+k^{2}$，模式数等于三倍（三种偏振）乘半径 $r$、厚度 $|dr|$ 的球壳体积，其中
\begin{equation}
\omega^{2}=\wms^{2}-\frac{4\pi^{2}}{L^{2}}r^{2}.
\end{equation}
于是得到模式密度
\begin{equation}
\rho(\omega)\,d\omega=3\times 4\pi r^{2}\,|dr|
=3V\,\frac{\omega\sqrt{\wms^{2}-\omega^{2}}}{2\pi^{2}}\,d\omega .
\end{equation}

\orthofig{images/05__004.png}{洛伦兹黑体发射谱}

黑体谱的正确处理当然是量子的，稍后我们将推导其量子版本；但值得先做一个纯\emph{经典}的计算。在洛伦兹宇宙中，由于光的频率没有上限，任何腔中的模式总数都无限。在量子力学发展之前这非常令人困惑：经典统计力学中称为能均分定理（equipartition theorem）的结果提示，这类模式每一个都应分得总能量的等份，若有无限多个则毫无意义。因此高频处没有令人满意的经典黑体谱计算。而在黎曼宇宙中，有限的模式数在整个谱上都给出合理的经典结果。

能均分定理指出：在温度为 $T$ 的平衡系统中，凡是其能量贡献正比于某个相空间坐标或动量平方的自由度，平均能量都是 $\half kT$，其中 $k$ 为玻尔兹曼常数（Boltzmann's constant），它只设定温度标度。然而，尽管我们腔中的每个模式看似只是单个坐标，它们实际上同时涉及一个坐标与一个动量：模式的瞬时振幅（以频率 $\omega_{i,j,k}$ 振荡）与其变化率（相应的动量）。正如谐振子（如弹簧上的质量块）中势能正比于振子位移的平方、动能正比于位移时间变化率的平方，两部分能量之和恒定；但对处于热平衡的一堆这样的振子，每个振子的平均能量是 $kT$，而不是 $\half kT$。

因此，频率 $\omega$ 处腔场的能量密度为
\begin{equation}
u_{classical}(\omega)\,d\omega=3kT\,\frac{\omega\sqrt{\wms^{2}-\omega^{2}}}{2\pi^{2}}\,d\omega .
\end{equation}
经典能量密度谱总在
\begin{equation}
\omega_{peak}=\frac{\wms}{\sqrt2}\approx 0.707\,\wms
\end{equation}
处取峰，与温度无关。整个谱上的总能量密度为 $u_{classical}(\omega)$ 对 $\omega$ 从 $0$ 到 $\wms$ 的积分：
\begin{equation}
u_{classical,\,total}=\frac{\wms^{3}\,kT}{2\pi^{2}}.
\end{equation}

在洛伦兹情形中，从腔中发射出的辐射谱形状与腔\emph{内}的辐射谱形状完全相同，因为任何从小孔逃逸的辐射都以同一速度——光速——行进。但在黎曼宇宙中，每个频率的光以不同的速度行进，使黑体发射谱与腔内的能量分布大不相同。

为计算黑体单位面积的发射功率，除了对每个角频率乘以它的速度
\begin{equation}
v(\omega)=\frac{\sqrt{\wms^{2}-\omega^{2}}}{\omega}
\end{equation}
外，还需对所有空间方向积分；对给定取向的孔，这给出额外的因子 $1/4$。于是
\begin{align}
I_{classical}(\omega)\,d\omega
&=\frac14 v(\omega)\,u_{classical}(\omega)\,d\omega \nonumber\\
&=\frac{3kT\,(\wms^{2}-\omega^{2})}{8\pi^{2}}\,d\omega .
\end{align}
所以经典发射谱就是一条倒置的抛物线，峰在 $\omega=0$ 处。温度的唯一作用是提高总发射功率，完全不改变谱形。这当然与我们自己宇宙的黑体谱形成鲜明对比：那里更热光源的光，其频率峰位于比更冷光源更短的波长，且总发射功率按温度的四次方标度，而非我们这里发现的线性标度。

当然，我们目前为止的黎曼计算都是经典的；要作真正比较需要量子版本。我们预期黎曼光服从与光子相同的一般量子统计规律：玻色子（具体说，自旋为 1 的粒子），不带电，因而自身即其反粒子。每个频率为 $\nu$ 的黎曼光子能量为 $h\nu$（$h$ 为普朗克常数，Planck's constant），或写作 $\hbar\omega$（$\hbar=h/2\pi$，$\omega$ 为角频率）。

量子统计力学中的相关计算\textsuperscript{[2]}告诉我们，腔场中每个模式平均被若干个黎曼光子占据，其数目是模式能量与温度的函数：
\begin{equation}
n(\omega)=\frac{1}{\exp\!\bigl(E(\omega)/kT\bigr)-1}=\frac{1}{\exp\!\bigl(\hbar\omega/kT\bigr)-1}.
\end{equation}
对高温（或 $\hbar\to0$ 的极限），此式变为 $n(\omega)\approx kT/(\hbar\omega)$，这恰好断言了能均分定理的同一件事：每个模式平均能量为 $kT$。注意这个表达式对负温度根本没有意义：这类系统不可能存在负温度下的热平衡。

把 $n(\omega)$ 乘以能量 $\hbar\omega$ 与态密度 $\rho(\omega)$，即得腔能量密度谱的量子版本：
\begin{equation}
u_{quantum}(\omega)\,d\omega
=\frac{3\hbar\omega^{2}\sqrt{\wms^{2}-\omega^{2}}}{2\pi^{2}\bigl(\exp(\hbar\omega/kT)-1\bigr)}\,d\omega .
\end{equation}
与经典版本一样，需再乘以同频光速的四分之一，以得到发射功率谱：
\begin{equation}
I_{quantum}(\omega)\,d\omega
=\frac{3\hbar\omega\,(\wms^{2}-\omega^{2})}{8\pi^{2}\bigl(\exp(\hbar\omega/kT)-1\bigr)}\,d\omega .
\end{equation}

\orthofig{images/05__005.png}{黎曼黑体腔谱}
\orthofig{images/05__006.png}{黎曼黑体发射谱}

在这些黎曼图中，我们取 $\wms=2\pi\times10^{15}$ Hz，然后使用真实世界的普朗克常数与玻尔兹曼常数值。可以看到，随着温度升高，黎曼经典谱与量子谱\emph{按比例地}彼此靠近：在最宽间隔处绝对差大致不变，但整个谱变大，故这只占能量的更小份额。

解读这些图时，务必牢记：在黎曼宇宙中，大多数物质的正温度都对应极端条件——例如气体中的相对论速度——而且正温度\emph{越低}，那些条件就\emph{越极端}。所以这里较高的温度反而代表最不奇异的态。

这提示我们：量子与经典黑体谱的差异在黎曼宇宙中会是相对次要的现象，不太可能在量子力学的发展中扮演它在我们自身历史中所扮演的角色。把这里的经典/量子差异与先前洛伦兹情形的图相比，洛伦兹紫外灾难（ultraviolet catastrophe）一目了然，而黎曼的"蓝色下垂（blue sag）"并不起眼，要在整个谱上作精确测量才能探测到。

然而，腔辐射还有另一种"灾难"：任何温度下辐射的平均压强都是无限的！为什么？对时间角频率 $\omega$ 接近零的腔模，为使其能量 $\half V A^{2}\omega^{2}$ 等于能均分定理确定的统一值 $kT$，其振幅 $A$ 必须变得越来越大的；而当能量被低的 $\omega$ 限制住时，它们的空间频率却接近 $\wms$，故其动量与压强（同样正比于 $A^{2}$）无界。

当然，完美黑体与完美热平衡都只是理想化，所以这并不表明任何真实物理量会变成无穷。但它确实清楚表明：在光能自由地在所有可能模式之间交换能量的任何系统中，辐射压都会极高。

\section{单粒子配分函数}

上一节为腔辐射枚举的那些态，本质上就是任何黎曼粒子可用的态。三种偏振的因子 $3$ 可能换成另一个因子（取决于粒子的自旋），$\wms$ 的取值则取决于粒子质量；除此之外，计数单粒子可能态的方法完全相同。当然，对玻色子，每个态可有任意多粒子；对费米子，每个态至多一个；但识别态本身的第一步对玻色子与费米子并无差别。

处理具有指定温度的平衡系统时，把关于系统态的信息打包到称为配分函数（partition function）的东西中很有用（另见文献\textsuperscript{[3]}）。它的一般定义为
\begin{equation}
Z=\sum_{r}\exp(-\beta E_{r}),\qquad \beta=\frac{1}{kT},
\end{equation}
指标 $r$ 只是系统态的标号，$E_{r}$ 表示各态的能量。我们定义了一个具有逆能量量纲的温度参数 $\beta$，它比温度本身更有用。$\beta$ 在无限温度处为零，而当温度从上方、下方趋近零时分别趋于 $+\infty$ 与 $-\infty$。

配分函数中的每一项正比于系统处于某个态的概率。概率本身可由该项除以所有项之和（即配分函数本身）得到，故
\begin{equation}
P(r)=\text{系统处于态 }r\text{ 的概率}=\frac{\exp(-\beta E_{r})}{Z}.
\end{equation}
配分函数最有用之处在于，系统的各种平均性质都能轻易从中提取。例如，平均能量为
\begin{align}
E_{mean}
&=\sum_{r}P(r)E_{r}
=\frac{\sum_{r}\exp(-\beta E_{r})E_{r}}{Z} \nonumber\\
&=-\frac{\partial_{\beta}Z}{Z}
=-\partial_{\beta}\ln Z .
\end{align}

质量为 $m$（故最大角频率为 $m/\hbar$）、被限制在体积 $V$ 内的单个无自旋黎曼粒子的配分函数为
\begin{align}
Z_{m}
&=\sum_{i,j,k}\exp(-\beta E_{i,j,k})
=\sum_{i,j,k}\exp(-\beta\hbar\,\omega_{i,j,k}) \nonumber\\
&=\sum_{i,j,k}\exp\!\left(-\beta\hbar\sqrt{(m/\hbar)^{2}-\frac{4\pi^{2}}{V^{2/3}}\bigl(i^{2}+j^{2}+k^{2}\bigr)}\right),
\end{align}
其中求和遍历使根号内表达式非负的所有整数 $i,j,k$（正负皆可）。我们把上一节的 $\omega_{i,j,k}$ 改用粒子质量 $m$（而非该粒子的最大角频率）表达，并把腔边长 $L$ 换成 $V^{1/3}$。在态数很大的极限下，可把它近似为 $(i,j,k)$ 空间中一个球形区域上的积分，用球极坐标、以 $r^{2}=i^{2}+j^{2}+k^{2}$ 描述最为方便：
\begin{equation}
Z_{m}=4\pi\int_{0}^{R}
\exp\!\left(-\beta\hbar\sqrt{(m/\hbar)^{2}-\frac{4\pi^{2}}{V^{2/3}}r^{2}}\right)r^{2}\,dr,
\end{equation}
其中 $R=V^{1/3}m/(2\pi\hbar)$ 是能量为零时的 $r$ 值。

作变量替换 $r=Rs$，上式变为
\begin{equation}
Z_{m}=3Z_{0}\int_{0}^{1}\exp\!\bigl(-\beta m\sqrt{1-s^{2}}\bigr)s^{2}\,ds,
\end{equation}
其中 $Z_{0}=(4/3)\pi R^{3}=Vm^{3}/(6\pi^{2}\hbar^{3})$ 为总态数。

再变换到角变量 $\gamma$，令 $s=\sin\gamma$，得到
\begin{equation}
Z_{m}=3Z_{0}\int_{0}^{\pi/2}\exp(-\beta m\cos\gamma)\sin^{2}\gamma\,\cos\gamma\,d\gamma .
\end{equation}
这最后一种形式可以解释为：对粒子能量-动量矢量所在的半 3-球面（half-3-sphere）的积分。

令人惊讶的是，这个积分可以写成显式公式：

\begin{keybox}[单粒子配分函数（Single-Particle Partition Function）]
\begin{equation}
Z_{m}=Z_{0}\left[1+\frac{3\pi\bigl[L_{2}(\beta^{*})-I_{2}(\beta^{*})\bigr]}{2\beta^{*}}\right],
\qquad
Z_{0}=\frac{Vm^{3}}{6\pi^{2}\hbar^{3}},
\qquad
\beta^{*}=\beta m \quad(\text{黎曼}).
\end{equation}
\end{keybox}

这里 $L_{2}$ 为修正斯特鲁维函数（modified Struve function），$I_{2}$ 为第一类修正贝塞尔函数（modified Bessel function of the first kind）。

\orthofig{images/05__007.png}{黎曼配分函数}

右图给出了配分函数，以及一个适用于 $\beta$ 的\emph{小值}（对应很高温度）的二次近似：
\begin{equation}
\lim_{\beta\to0}Z_{m}\approx Z_{0}\left[1-\frac{3\pi\beta^{*}}{16}+\frac{\beta^{*2}}{5}\right].
\end{equation}
当然，如需更高精度，完全可以把幂级数取到更高阶。此外，还有一个适用于 $\beta$ 的\emph{大负值}（对应小的负温度）的渐近形式，一旦 $\beta^{*}$ 足够大就开始有用：
\begin{equation}
\lim_{\beta\to-\infty}Z_{m}\approx
\frac{3Z_{0}}{\beta^{*2}}\left[1-\frac{\exp(-\beta^{*})\,(8\beta^{*}+15)}{8\sqrt{-2\beta^{*}/\pi}}\right].
\end{equation}

\section{理想气体}

\subsection{固定粒子数的稀薄气体}

设有由固定数目 $N$ 个粒子组成的气体，足够稀薄，使得任一量子态的平均粒子数远小于 $1$。此时粒子是玻色子还是费米子无关紧要，我们处于一种经典极限。不过，粒子仍应被正确地视为\emph{不可分辨}（indistinguishable），这避免了真正经典处理中出现的一些问题。

在此情形下，我们可以用上一节导出的单粒子配分函数来表达整团气体的配分函数 $Z_{G}$。借助 Reif\textsuperscript{[4]} 的一些一般计算，有
\begin{equation}
\ln Z_{G}=N-N\ln N+N\ln Z_{m}.
\end{equation}
我们以自然对数的形式给出配分函数，因为这在计算气体的各种宏观性质时最直接有用。

\subsubsection{能量}

我们首先要算的是气体的平均能量：
\begin{align}
E_{mean}
&=-\partial_{\beta}\ln Z_{G}
=-N\,\partial_{\beta}\ln Z_{m}
=-N\,\frac{\partial_{\beta}Z_{m}}{Z_{m}} \nonumber\\
&=3Nm\left(\frac{1}{\beta^{*}}+
\frac{\pi L_{1}(\beta^{*})-\pi I_{1}(\beta^{*})+2}
{3\pi I_{2}(\beta^{*})-3\pi L_{2}(\beta^{*})-2\beta^{*}}\right),
\end{align}
其中 $\beta^{*}=\beta m$。

\orthofig{images/05__008.png}{稀薄黎曼气体的平均能量}

无限温度极限可轻易从配分函数的二次近似得到：
\begin{equation}
\lim_{\beta\to0}E_{mean}=\frac{3\pi Nm}{16}\approx 0.589\,Nm .
\end{equation}
如右图所示，对大的负 $\beta$（对应小的负温度），粒子运动缓慢，平均能量接近它们静止质量能量之和；对小的 $\beta$（对应大的温度），粒子运动更快，总能量更少——正如我们从总能量与动能的关系所预期的；而对大的正 $\beta$（对应小的正温度），气体处于高度相对论态，其动能接近静止质量，使总能量几乎降到零。

\subsubsection{压强}

接下来计算气体的平均压强。用配分函数表示压强的通常公式\textsuperscript{[5]}利用了如下事实：当系统能量依赖其体积时，能量对体积的变化率的\emph{相反数}就是系统施加的压强。（设想系统缓慢膨胀，对外做正功。所作之功是力乘以膨胀的线性距离，也就是压强乘以体积变化。）于是
\begin{align}
p_{mean}
&=\sum_{r}P(r)p_{r}
=\frac{\sum_{r}\exp(-\beta E_{r})\,p_{r}}{Z} \nonumber\\
&=\frac{\sum_{r}\exp(-\beta E_{r})\,(-\partial_{V}E_{r})}{Z}
=\frac{\partial_{V}\sum_{r}\exp(-\beta E_{r})}{\beta Z}
=\frac{\partial_{V}\ln Z}{\beta}.
\end{align}

然而，这对黎曼情形\emph{并不}正确！一个次要差别是：在黎曼物理中，能量对体积的变化率就是压强本身，而非其相反数；这本质上与如下原因是同一个——在自由粒子的哈密顿处理中，共轭于每个空间坐标的动量是通常动量的\emph{相反数}，如我们此前所讨论的。

但即便在这里改一个符号，我们仍会留下一个方程：它看似熟悉，实际上却毫无意义。

不妨还是把这个计算做完，看看问题出在哪里。由于 $Z_{m}$ 只是 $V$ 的线性函数（经由 $Z_{0}$），有
\begin{equation}
\partial_{V}\ln Z_{m}=\frac{\partial_{V}Z_{m}}{Z_{m}}=\frac{Z_{m}/V}{Z_{m}}=\frac{1}{V}.
\end{equation}
转到整团气体的配分函数，并用标准压强公式、同时对黎曼的能量/体积/压强关系改一个符号，得到
\begin{equation}
p_{mean}=-\frac{\partial_{V}\ln Z_{G}}{\beta}
=-\frac{N\,\partial_{V}\ln Z_{m}}{\beta}
=-\frac{N}{\beta V}\qquad[\text{错误！}]
\end{equation}
这与下式等价：
\begin{equation}
p_{mean}V=-NkT\qquad[\text{错误！}]
\end{equation}
这只是理想气体的标准状态方程，差一个符号。虽然对低负温度而言这差不多是我们所\emph{预期}的，而且"温度趋于无限时压强应趋于无限"表面上似乎也说得通……但这个方程意味着一旦跨入正温度，压强就会变成负的！生活在一个用相对论粒子轰击容器壁、反而让容器处于\emph{张力}而非压强之下的世界里固然有趣，但那完全不合物理。

这个谜题的解答是：我们的配分函数 $Z_{m}$ 带有一种对容器体积的额外依赖，而在导出"用配分函数表示压强"的标准公式时所依据的洛伦兹物理中，这种依赖根本不存在。当我们改变容器的体积时，改变的不仅是各粒子态的能量，还有态的\emph{数目}。在黎曼情形下，更大的容器允许在达到最小波长之前有更多（有限个）粒子态；而在洛伦兹情形下，由于没有最小波长，态的数目永远是无限的。

解决此问题最简单的办法，是回到 $Z_{m}$ 的某个早期形式：
\begin{equation}
Z_{m}=4\pi\int_{0}^{R}
\exp\!\left(-\beta\hbar\sqrt{(m/\hbar)^{2}-\frac{4\pi^{2}}{V^{2/3}}r^{2}}\right)r^{2}\,dr,
\qquad R=\frac{V^{1/3}m}{2\pi\hbar}.
\end{equation}
由此我们可以找出并减去 $V$ 经由积分上限 $R$ 所产生的影响，只留下我们想要的效应：能级的变化。$\partial_{V}Z_{m}$ 中来自积分上限的项，就是被积函数（含整体常数 $4\pi$）在 $r=R$ 处的值乘以 $\partial_{V}R$。它等于
\begin{equation}
4\pi\exp(0)\,R^{2}\times\frac13 V^{-2/3}\frac{m}{2\pi\hbar}
=\frac{(4/3)\pi R^{3}}{V}
=\frac{Z_{0}}{V},
\qquad Z_{0}=\frac{(4/3)\pi R^{3}}{1}=\frac{Vm^{3}}{6\pi^{2}\hbar^{3}}.
\end{equation}

\orthofig{images/05__009.png}{稀薄黎曼气体的压强}

若重复我们原来的计算，但从 $\partial_{V}Z_{m}$ 中减去这个虚假贡献、只留下平均压强，则得到
\begin{align}
p_{mean}
&=-N\,\frac{(Z_{m}/V)-(Z_{0}/V)}{Z_{m}\,\beta} \nonumber\\
&=-\frac{N}{\beta V}\left[1-\frac{Z_{0}}{Z_{m}}\right],
\end{align}
这等价于
\begin{equation}
p_{mean}V=-NkT\left[1-\frac{Z_{0}}{Z_{m}}\right]\qquad[\text{正确！}]
\end{equation}
当温度为负时 $Z_{m}>Z_{0}$，方括号中的修正因子为正，故压强为正。事实上，对低负温度，$Z_{m}$ 极大，修正因子非常接近 $1$，于是我们得到传统的理想气体定律，只是按 $T$ 的符号作了调整。

但当温度为正时 $Z_{m}<Z_{0}$，方括号中的修正因子为负；不过此时 $-1/\beta$ 也（因 $\beta>0$）为负，两者相乘使压强重新保持为正，从而恢复了物理上合理的结论。取 $\beta\to0$ 的极限，即无限温度，得
\begin{equation}
\lim_{\beta\to0}p_{mean}=\frac{3\pi Nm}{16V}.
\end{equation}
当 $\beta$ 趋于正无穷时，温度从上方趋近零，进入越来越极端的相对论粒子速度领域，压强随之趋于无穷，正如你所预期的。

\subsubsection{熵}

系统的熵也可以从配分函数得到\textsuperscript{[6]}：
\begin{equation}
S=k(\ln Z+\beta E_{mean}).
\end{equation}
对我们的稀薄气体，
\begin{equation}
S_{G}=k(\ln Z_{G}+\beta E_{mean}),\qquad
\ln Z_{G}=N-N\ln N+N\ln Z_{m},\qquad
E_{mean}=-N\frac{\partial_{\beta}Z_{m}}{Z_{m}}.
\end{equation}

\orthofig{images/05__010.png}{稀薄黎曼气体的熵}

我们不再把熵完整地写成含对数、贝塞尔与斯特鲁维函数的表达式；如图所示，它在 $\beta=0$（即无限温度）处有一个有限的峰。

在不同温度下，$N$ 个粒子如何按能量分布？在我们稀薄气体的假设下，标号为 $r$ 的任一态的平均粒子数可证为\textsuperscript{[4]}
\begin{equation}
n_{r}=\frac{N\exp(-\beta E_{r})}{\sum_{s}\exp(-\beta E_{s})}=\frac{N\exp(-\beta E_{r})}{Z_{m}}.
\end{equation}
我们关心按能量 $E=\hbar\omega$ 的分布，因此可以利用我们为腔辐射计算过的、对 $\omega$ 的态密度，只须去掉偏振因子 $3$（因为我们为理想气体设想的是一个无自旋粒子）：
\begin{equation}
\rho(\omega)\,d\omega=\frac{V\,\omega\sqrt{\wms^{2}-\omega^{2}}}{2\pi^{2}}\,d\omega,
\qquad
\rho(E)\,dE=\frac{V\,E\sqrt{m^{2}-E^{2}}}{2\pi^{2}\hbar^{3}}\,dE .
\end{equation}

\orthofig{images/05__011.png}{稀薄黎曼气体按能量的粒子数分布}

把它乘以该能量下每个态的平均占据数，得到
\begin{align}
n(E)\,dE
&=N\,\frac{V}{2\pi^{2}\hbar^{3}}\,E\sqrt{m^{2}-E^{2}}\,\frac{\exp(-\beta E)}{Z_{m}}\,dE \nonumber\\
&=N\,\frac{3Z_{0}}{m^{3}Z_{m}}\,E\sqrt{m^{2}-E^{2}}\,\exp(-\beta E)\,dE .
\end{align}
用粒子质量 $m$ 的倍数 $\varepsilon$ 来表示能量更为方便，即 $E=\varepsilon m$：
\begin{equation}
n(\varepsilon)\,d\varepsilon
=N\,\frac{3Z_{0}}{Z_{m}}\,\varepsilon\sqrt{1-\varepsilon^{2}}\,\exp(-\beta^{*}\varepsilon)\,d\varepsilon .
\end{equation}
右图给出了三个 $\beta^{*}$ 取值下的分布。在较低温度——无论正负——粒子的能量相对受限，而在无限温度下它们的分布更宽。

无限温度的稀薄气体有一个有趣的性质：它的应力-能量张量（stress-energy tensor）是 SO(4) 不变的，因此在这种气体内部\emph{以任意速度}运动的任何观测者，测得周围环境完全相同的压强与能量密度。我们可以这样看出这一点：无限温度下，气体的平均能量密度与平均压强都等于 $3\pi Nm/(16V)$。于是应力-能量张量只是单位矩阵的一个倍数，坐标变换不会改变它。

不过，细看按能量的粒子分布会发现，这种 SO(4) 不变性是以相当微妙的方式达成的。$\beta=0$（即无限温度）时的分布为
\begin{equation}
n_{\infty}(\varepsilon)\,d\varepsilon=3N\,\varepsilon\sqrt{1-\varepsilon^{2}}\,d\varepsilon .
\end{equation}
若改用角变量 $\gamma$（粒子 4-速度与时间轴的夹角），则 $\varepsilon=\cos\gamma$，于是
\begin{equation}
n_{\infty}(\gamma)\,|d\gamma|=3N\,\cos\gamma\,\sin^{2}\gamma\,|d\gamma| .
\end{equation}
3-球面上旋转不变的测度（在对其余与三维空间中不同方向有关的坐标积分掉之后）正是 $\sin^{2}\gamma\,d\gamma$，所以这说明粒子的 4-速度\emph{并非}如人们对 SO(4) 不变态所预期的那样均匀分布在 3-球面上。与均匀分布相比，粒子更倾向于分布在较低速度、即更靠近时间轴、$\cos\gamma$ 较大之处。

那么无限温度气体又是如何变得 SO(4) 不变的呢？考察一个能量为 $E=m\cos\gamma$ 的粒子，其 4-速度 $\bu$ 的时间分量 $u^{t}$ 将是 $\cos\gamma$。若粒子静止系中的能量密度为 $\rho$，则其应力-能量张量为 $\rho\,\bu\otimes\bu$。这意味着该粒子对静止容器系中总能量的贡献为 $V\rho\,(u^{t})^{2}=V\rho\cos^{2}\gamma$。令其等于 $E=m\cos\gamma$，就得到一个随粒子速度升高、$\cos\gamma$ 下降而增大的静止系能量密度：
\begin{equation}
\rho=\frac{m}{V\cos\gamma}.
\end{equation}
这个因子看似费解，因为在那运动粒子的静止系中，测得的容器宽度是\emph{增大}的；这是黎曼版的、更熟悉的洛伦兹相对论长度收缩（length contraction）的对应物。但同理，若有若干粒子以相同速度穿越容器，在它们的共同静止系中，测得它们比容器静止的观测者所测的更靠拢。

所以，对一个快速穿越气体的观测者而言，低 $\cos\gamma$ 粒子相对稀少这一情况，被它们看起来更密集这一点所补偿。尽管粒子速度的分布看起来不同，总的能量密度（与压强）却保持不变。

\subsection{含粒子对产生的费米子气体}

对缓慢运动的粒子气体，自然假设其粒子数保持固定。但当粒子以相对论速度运动时，这是一个相当人为的约束，因为随机的高速碰撞有时可能把粒子散射到过去（我们此前考察过的一类碰撞）。这是否需要反粒子存在，取决于粒子的本性（有些粒子自身即其反粒子，如同我们宇宙中的光子），但它必定涉及粒子总数的\emph{变化}。

设有由自旋为 $\tfrac12$ 的粒子组成的理想气体。它们是费米子，服从费米-狄拉克统计（Fermi-Dirac statistics）：每个量子态至多被一个粒子占据。假设存在一种机制可以产生粒子-反粒子\emph{对}，所需能量由气体内能的变化来平衡。这里我们不关心粒子对产生与湮灭过程的具体细节；我们的目标只是弄清这样的系统能否在正、负温度两个区域都处于热力学平衡。在腔辐射的处理中，我们发现由数目不定、可自由产生的玻色子构成的系统总有\emph{正}温度。我们的费米子气体能产生新粒子，是否也会如此？

由于粒子与反粒子成对地产生和湮灭，它们数目之差将保持恒定。我们把波动的粒子数与反粒子数分别记为 $N_{P}$ 与 $N_{A}$，固定的粒子超出量记为 $N$：
\begin{equation}
N=N_{P}-N_{A}.
\end{equation}
我们容器中可供粒子使用的每个单粒子量子态 $r$，都或被 $0$、或被 $1$ 个粒子占据，以及或被 $0$、或被 $1$ 个反粒子占据；泡利不相容原理（Pauli Exclusion principle）禁止每种粒子在同一态中多于一个，但并不阻止一个粒子与一个反粒子处于同一态。我们把态 $r$ 中粒子与反粒子的数目分别记为 $n_{P,r}$ 与 $n_{A,r}$。

整团气体的一个态，就是对每个 $n_{P,r}$ 与 $n_{A,r}$ 赋以 $0$ 或 $1$，并使得粒子对反粒子的总超出量为 $N$。我们把给定 $N$ 值的气体配分函数记为 $Z_{G}(N)$：
\begin{equation}
Z_{G}(N)=\sum_{\{n_{P,r},\,n_{A,r}\}:\,N_{P}-N_{A}=N}
\exp\!\left(-\beta\sum_{r}E_{r}\,(n_{P,r}+n_{A,r})\right).
\end{equation}
这里外层求和遍历粒子超出量为 $N$ 的所有气体态，而每个气体态的总能量取为其中所有粒子能量之和。

体积为 $V$ 的容器中单粒子态数是有限的；我们在计算单粒子配分函数时已求出这个数，即 $Z_{0}$，但在此精确地说应取其两倍，以计入我们费米子的两个自旋态：
\begin{equation}
N_{max}=2Z_{0}=\frac{Vm^{3}}{3\pi^{2}\hbar^{3}}.
\end{equation}
若我们为粒子备有 $N_{max}$ 个可用态、为反粒子同样的数目，则粒子超出量为 $N$ 的态数可通过对反粒子数求和得到：
\begin{align}
\bigl|\{n_{P,r},n_{A,r}:N_{P}-N_{A}=N\}\bigr|
&=\sum_{i=0}^{N_{max}-N}\binom{N_{max}}{N+i}\binom{N_{max}}{i} \nonumber\\
&=\binom{2N_{max}}{N_{max}+N} \nonumber\\
&\approx\frac{2^{2N_{max}}}{\sqrt{\pi N_{max}}}\exp\!\left(-\frac{N^{2}}{N_{max}}\right).
\end{align}

假定我们固定 $\beta$，只考虑气体配分函数 $Z_{G}(N)$ 随 $N$ 变化的方式。外层求和的项数由上面导出的高斯函数给出。而求和中各项的平均值不会随 $N$ 迅速变化，因为无论 $N$ 取何值，它们都以大致相同的方式取样各态的能量。于是我们可以\emph{非常粗略}地说，$Z_{G}(N)$ 本身将表现得像计数项数的高斯函数的一个倍数，即
\begin{equation}
Z_{G}(N)\approx C(\beta)\exp\!\left(-\frac{N^{2}}{N_{max}}\right),
\end{equation}
其中常数 $C(\beta)$ 依赖 $\beta$，我们已把其它依赖 $N_{max}$ 的因子吸收进了它。

现在，若把 $Z_{G}(N)$ 乘以一个关于 $N$ 的指数函数 $\exp(\alpha N)$，那会使高斯峰发生位移：
\begin{align}
Z_{G}(N)\exp(\alpha N)
&\approx C(\beta)\exp\!\left(-\frac{N^{2}}{N_{max}}+\alpha N\right) \nonumber\\
&=C(\beta)\exp\!\left(\frac{N_{max}\alpha^{2}}{4}\right)
\exp\!\left(-\frac{(N-N_{max}\alpha/2)^{2}}{N_{max}}\right).
\end{align}
若我们想让它例如在 $N_{0}$ 处取峰，则令
\begin{equation}
\alpha=\frac{2N_{0}}{N_{max}}.
\end{equation}
我们可以用一个对 $N$ 不加约束的求和、并借助高斯峰的位置挑出 $N=N_{0}$ 这一点，来近似配分函数 $Z_{G}(N_{0})$：
\begin{equation}
\sum_{N=0}^{N_{max}}Z_{G}(N)\exp(\alpha N)\approx Z_{G}(N_{0})\exp(\alpha N_{0})\,\Delta N,
\end{equation}
其中 $\Delta N$ 度量峰的宽度。我们把左端的求和称为"巨配分函数（grand partition function）"\textsuperscript{[7]}；它对气体所有可能态求和，不论粒子数如何，只在每个能量上附加一项。由于没有了约束，我们可以把它的计算进行得相当远：
\begin{align}
Z_{GRAND}
&=\sum_{N=0}^{N_{max}}Z_{G}(N)\exp(\alpha N) \nonumber\\
&=\sum_{\{n_{P,r},n_{A,r}\}}
\exp\!\left(-\beta\sum_{r}E_{r}(n_{P,r}+n_{A,r})\right)\exp(\alpha N) \nonumber\\
&=\sum_{\{n_{P,r},n_{A,r}\}}
\exp\!\left(-\sum_{r}\bigl[n_{P,r}(\beta E_{r}-\alpha)+n_{A,r}(\beta E_{r}+\alpha)\bigr]\right) \nonumber\\
&=\prod_{r}\left[
\left(\sum_{n_{P,r}=0}^{1}\exp(-n_{P,r}(\beta E_{r}-\alpha))\right)
\left(\sum_{n_{A,r}=0}^{1}\exp(-n_{A,r}(\beta E_{r}+\alpha))\right)\right] \nonumber\\
&=\prod_{r}\left[
\bigl(1+\exp(-(\beta E_{r}-\alpha))\bigr)
\bigl(1+\exp(-(\beta E_{r}+\alpha))\bigr)\right].
\end{align}

我们可以这样近似气体配分函数的对数：
\begin{align}
\ln Z_{G}(N_{0})
&\approx-\alpha N_{0}-\ln\Delta N+\ln Z_{GRAND} \nonumber\\
&\approx-\alpha N_{0}+\sum_{r}\Bigl[
\ln\bigl(1+\exp(-(\beta E_{r}-\alpha))\bigr)
+\ln\bigl(1+\exp(-(\beta E_{r}+\alpha))\bigr)\Bigr].
\end{align}
再作一些类似文献\textsuperscript{[7]}中的计算，可以证明各态的平均占据数，对粒子与反粒子分别为
\begin{equation}
n_{P,r,mean}=\frac{1}{\exp(\beta E_{r}-\alpha)+1},
\qquad
n_{A,r,mean}=\frac{1}{\exp(\beta E_{r}+\alpha)+1}.
\end{equation}
与黎曼光子态占据数的情形不同，这些表达式不会对负温度变得无意义；事实上，无论 $\alpha$ 与 $\beta$ 取何值，它们都介于 $0$ 与 $1$ 之间，正是费米子所要求的。要为给定的 $\beta$ 与 $N_{0}$ 求出 $\alpha$，我们需要解
\begin{equation}
N_{0}=\sum_{r}\bigl(n_{P,r,mean}-n_{A,r,mean}\bigr)
=\sum_{r}\frac{\sinh\alpha}{\cosh\alpha+\cosh(\beta E_{r})},
\end{equation}
以定出 $\alpha$。由于这个求和的积分形式似乎没有闭式解析解，一般需要数值求解。我们的初始近似 $\alpha=2N_{0}/N_{max}$ 可以作为迭代求解的起点。

不过，若取 $N_{0}=0$，考虑粒子与反粒子数\emph{精确相等}的气体，我们就能展示这类系统的主要定性特征。由对称性，此时对所有 $\beta$ 值都有 $\alpha=0$，于是
\begin{equation}
\ln Z_{G}\approx 2\sum_{r}\ln\bigl(1+\exp(-\beta E_{r})\bigr),
\qquad
n_{P,r,mean}=n_{A,r,mean}=\frac{1}{\exp(\beta E_{r})+1}.
\end{equation}

\orthofig{images/05__012.png}{费米子对气体的总能量}

右图给出系统的平均能量，其形状大致与稀薄气体的平均能量曲线相同，但涉及的能量与粒子质量乘以单粒子态总数可比。

平均能量由数值积分算出，采用与单粒子配分函数相同的那组变量 $r$ 与 $s$：
\begin{align}
E_{mean}
&=\int_{0}^{R}\bigl(-\partial_{\beta}\ln Z_{G}\bigr)\,4\pi r^{2}\,dr \nonumber\\
&=2\int_{0}^{R}\bigl(-\partial_{\beta}\ln(1+\exp(-\beta E_{r}))\bigr)\,4\pi r^{2}\,dr \nonumber\\
&=6N_{max}m\int_{0}^{1}
\frac{s^{2}\sqrt{1-s^{2}}}{\exp(\beta^{*}\sqrt{1-s^{2}})+1}\,ds .
\end{align}
对 $\beta^{*}=0$ 与 $\beta^{*}\to-\infty$，指数分别趋于 $1$ 与 $0$，该积分可精确求出，分别给出 $3\pi N_{max}m/16$ 与 $3\pi N_{max}m/8$。

\orthofig{images/05__013.png}{费米子对气体的总粒子数}

下一张图给出平均总粒子数（含反粒子）。在低负温度下系统接近饱和，几乎每一个量子态都被一个粒子和一个反粒子填满。在无限温度下这些态平均半满，而随着我们进入越来越低的正温度，数目逐渐下降。这里所绘的积分为
\begin{align}
N_{total,mean}
&=\int_{0}^{R}\bigl(n_{P,r,mean}+n_{A,r,mean}\bigr)\,4\pi r^{2}\,dr \nonumber\\
&=2\int_{0}^{R}\frac{1}{1+\exp(\beta E_{r})}\,4\pi r^{2}\,dr \nonumber\\
&=6N_{max}\int_{0}^{1}
\frac{s^{2}}{\exp(\beta^{*}\sqrt{1-s^{2}})+1}\,ds .
\end{align}
同样，对 $\beta^{*}=0$ 与 $\beta^{*}\to-\infty$，指数趋于 $1$ 与 $0$，积分可精确求出，分别给出 $N_{max}$ 与 $2N_{max}$。

\orthofig{images/05__014.png}{费米子对气体的平均单粒子能量}

第三张图给出平均单粒子能量，它就是第一条曲线除以第二条曲线。对从负温度一直到无限温度，它几乎保持不变，约为 $3\pi m/16$，即对所有态的平均能量。在小的负温度下，这是因为几乎所有态都被永久占据；而在较大的负温度下，平均占据数虽下降，但粒子仍均匀分布在各态之间。

只有当我们进入低正温度时，才不仅粒子数下降，而且粒子开始占据较低能量的态。若处理的系统比反粒子多出 $N_{0}$ 个粒子（或反之），则能量将趋于一个渐近值，即 $N_{0}$ 个粒子填满了最低的 $N_{0}$ 个能态。但由于我们取 $N_{0}=0$，低正温度下的渐近能量为零。

务必记住：这里没有任何东西告诉我们任何粒子对产生过程的机制或速率；这些结果只是把该过程的存在并已达平衡视作理所当然。粒子对产生发生的速率很可能强烈依赖温度；若在某个温度 $T$ 下速率很低，那么该温度下的平衡态可能完全不能说明我们会实际观测到什么。所以我们不应把这些负温度费米子对气体——其中几乎每一个量子态都被占据——的存在，当作任意负温度费米子气体的近期命运！

\orthofig{images/05__015.png}{费米子对气体的平均压强}

不过，对一个具有正温度的高度相对论气体，即便粒子之间的经典碰撞也能导致粒子数变化，所以这里的描述更可能是适用的。

系统的平均压强为
\begin{equation}
p_{mean}=-\frac{1}{\beta}\,\partial_{V}\ln Z_{G},
\end{equation}
但有一个提醒：我们必须避免在求稀薄气体压强时已识别的那个陷阱，要在对单粒子态积分\emph{之前}就对 $V$ 求导。具体计算从略。

在无限温度下，我们求得的压强换算到每个粒子，与稀薄气体的结果相同。所以至少在高温度下，当原气体含 $N$ 个粒子时，含粒子对产生的费米子气体给出的每粒子压强与稀薄气体一致。

\section{参考文献}

\begin{itemize}
  \item[{[1]}] \emph{Fundamentals of Statistical and Thermal Physics}, Frederick Reif, McGraw-Hill, 1965, 第 9.13、9.14、9.15 节。
  \item[{[2]}] Reif, \emph{op. cit.}, 第 9.3、9.5 节。
  \item[{[3]}] Reif, \emph{op. cit.}, 第 6.5、6.6 节。
  \item[{[4]}] Reif, \emph{op. cit.}, 第 9.8 节。
  \item[{[5]}] Reif, \emph{op. cit.}, 第 6.5 节，方程 (6.5.12)。
  \item[{[6]}] Reif, \emph{op. cit.}, 第 6.6 节，方程 (6.6.5)。
  \item[{[7]}] Reif, \emph{op. cit.}, 第 9.6 节。
\end{itemize}

\chapter{黎曼量子力学}

在勾股定理符号反转的黎曼宇宙（Riemannian universe）中，物质能否稳定存在？经典图像给出的答案是否定的：处于势能谷（energy valley）中的带电粒子一旦振动就会辐射，而黎曼物理中动能与辐射能的符号相反，辐射反而使它加速，最终逃逸出谷。本章用量子力学来回答这一问题。我们先分析晶体点阵（crystal lattice）中单个粒子所感受到的、由无穷多近邻叠加而成的势井（energy well），再讨论非完美谐振子（imperfect harmonic oscillator）的经典困境、固体中的量子效应与稳定性判据，进而考察两粒子"黎曼原子"（Riemannian atom）与可能的大分子结构；章末以两节进阶材料给出本书采用的单位约定，以及黎曼相对论量子力学中能量与动量算符的构造。

\section{固体的结构（The Structure of Solids）}

在我们的宇宙中，静电库仑力（Coulomb force）服从简单的平方反比律（inverse-square law）：同种电荷相斥、异种电荷相吸，故正电荷附近的正电荷会感受到一个形如陡峭光滑山坡的“能量地形（energy landscape）”，其坡度随距离逐渐变平；负电荷附近则是一个能量坑（energy pit），即把前者上下颠倒。

在黎曼宇宙中，相应的力在近处使\emph{同种电荷相吸、异种电荷相斥}。但这只在某一距离以内成立，因为荷电粒子周围的能量地形之上，叠加了在平方反比衰减上规则排列的一串波峰与波谷。

\orthofig{images/04__003.png}{正电荷周围黎曼与洛伦兹静电势}
\orthofig{images/04__004.png}{负电荷周围黎曼与洛伦兹静电势}

右图给出能量随距离的变化。在中心粒子周围的三维空间中，这些波峰与波谷将形成同心球壳（concentric shells），从而粒子有可能落入彼此的能量波谷而被束缚在一起。

但其中包含一个麻烦。若用经典物理分析（即不计量子力学效应），能量谷中的粒子不会静止，因为它通常拥有足够能量来回运动；而这样运动就会发射电磁辐射。在我们宇宙中，这会使它损失动能、运动衰减；但在黎曼宇宙中，动能与电磁辐射能的符号相反，因此在能量谷中振荡的粒子在辐射时反而会\emph{加速}。若它这样不断获得动能，最终会彻底逃出谷。

粒子得以保持束缚的途径之一是：它振荡的固有频率大于电磁辐射的最大频率 $\nu_{max}$。带电粒子以高于 $\nu_{max}$ 的频率振荡时不会产生辐射。

然而，要产生单一纯频率的振荡，能量谷必须具有完美的形状，即势能作为到中心距离的函数是一条无瑕的抛物线（parabola）。这是非常苛刻的要求，库仑力当然无法满足。对“不完美”的能量谷，即使振荡的主频率高于 $\nu_{max}$，粒子运动也仍含有较低频率的分量。

那么，黎曼宇宙中任何种类的物质如何保持稳定？答案在于量子力学。稍后我们会讨论细节，但为此需要构造一个简单而具体的例子来演示这一效应。

与其考察只有两个粒子由静电力束缚的黎曼版氢原子，不如从另一极端入手，考虑一个粒子被大量规则排列的近邻所包围时所感受到的能量地形。

我们要特别考察的这种排列有多种名称：晶体学家称之为“面心立方系统（face-centred cubic system）”，数学家则称之为“$A_{3}$ 点阵（$A_{3}$ lattice）”。若取一组立方体网格，在每个立方体的每个顶点处以及每个面的中心放一个粒子，就得到 $A_{3}$ 点阵。数学上，取所有满足 $x+y+z$ 为偶整数的点 $(x,y,z)$，便得到一种每边长为 $2$、最近邻距离为 $\sqrt{2}$ 的点阵。当然，我们可以把整个点阵乘上任意因子 $s$ 随意缩放。

\orthofig{images/07__001.png}{$A_{3}$ 或面心立方点阵}
\orthofig{images/07__002.png}{堆叠的橘子与 $A_{3}$ 结构}

左上图中蓝色平面展示了“切开”同一点阵的另一种方式。这些平面被分成等边三角形，顶点处是粒子。这种排列正是球形物体（如橘子或台球）在一个平面内自然堆叠的方式，而相邻平面彼此契合，就像一层层橘子嵌入下一层的凹陷中一样。

$A_{3}$ 点阵中每个粒子有十二个最近邻：同一“橘子层”内六个，上一层三个，下一层三个。由于最近邻距离为 $\sqrt{2}$，其距离平方为 $2$。在我们最初的方案中，每个点阵点具有坐标 $(x,y,z)$ 且 $x+y+z$ 为偶整数，而由
\begin{equation}
(x+y+z)^{2}-(x^{2}+y^{2}+z^{2})=2(xy+yz+xz)
\end{equation}
可知 $(x,y,z)$ 到 $(0,0,0)$ 的距离平方也必为偶整数。由此，$(0,0,0)$ 的近邻距离依次为
\[
\sqrt{2},\ \sqrt{4},\ \sqrt{6},\ \sqrt{8},\ \sqrt{10}\ \dots
\]
点阵中每一点的近邻都呈相同构型。处于逐渐增大距离上的粒子数目，就是把 $2,4,6,8,10\dots$ 写成三个整数平方和的方式数（允许零与负数，且不同的排列顺序分别计数）。前几项计数为
\[
12,\ 6,\ 24,\ 12,\ 24,\ 8,\ 48,\ 6,\ 36,\ 24\ \dots
\]
但这一数列没有简单的公式。

现在假设我们有一团按 $A_{3}$ 点阵排列的正电荷粒子。若这团只有 $4$ 个粒子，它们可以彼此等距，我们还能选择合适的整体尺度使它们都恰好落在彼此最近邻能量波谷的正中心。但对一大团粒子，粒子间距将包含上面列出的各种近邻距离，即所有偶整数的平方根。显然它们\emph{并不}都是某一取值的整数倍，因此不可能让它们各自都恰好位于能量波谷底部。于是整体点阵间距将是一个折中值：它在“让最近邻距离接近库仑势的谷-谷间距 $\lambda_{min}$”与“这一选择对其余更远近邻的影响”之间权衡能量、使总能量最小。

结果是对非常大的团簇，当
\begin{equation}
s=\frac{\sqrt{3}}{2}\lambda_{min}\approx 0.866\,\lambda_{min}
\end{equation}
时团簇势能最小，而最近邻距离为
\begin{equation}
d=\sqrt{2}\,s=\sqrt{\tfrac{3}{2}}\,\lambda_{min}\approx 1.2247\,\lambda_{min}.
\end{equation}
由于这几乎恰好位于 $\lambda_{min}$ 上方四分之一波长处，最近邻对能量的贡献会非常小！尽管如此，包含更多更远近邻的求和结果显示，在这个 $s$ 取值下能量井达到最深。

此时引入精细结构常数（fine structure constant）$\alpha$ 会很有用，它是一个无量纲量，定义为（在时间与空间采用相同单位、即本书通篇所用单位下）
\begin{equation}
\alpha=\frac{e^{2}}{4\pi\hbar\varepsilon_{0}} .
\end{equation}
这里 $e$ 是构成点阵的粒子的电荷，$\varepsilon_{0}$ 是决定静电力强度的介电常数（permittivity constant），$\hbar$ 是普朗克常数除以 $2\pi$。在我们自己的宇宙中 $\alpha\approx 1/137$，但黎曼宇宙中当然没有理由与之相同。

我们再定义
\begin{equation}
m_{\mathrm{photon}}=\hbar\,\wms .
\end{equation}
给定角频率的光，其每个\emph{量子}（quantum）的能量等于该角频率乘以 $\hbar$。因此 $\hbar\wms$ 就是最大可能角频率的光所对应的能量。而正如我们考察黎曼宇宙中的能量与动量时所见，粒子的最大总能量就是其静止质量，所以 $\hbar\wms$ 正是一个黎曼光子的静止质量。

把这两个常数相乘，乘积 $\alpha m_{\mathrm{photon}}$ 等于相距 $\lambda_{min}$ 的两个粒子势能的 $2\pi$ 倍。这使得它成为任何“粒子间距与 $\lambda_{min}$ 挂钩”情形下一个有用的能量标度。

\orthofig{images/07__003.png}{$A_{3}$ 点阵中井深随粒子数的变化}

右图的双对数图给出了一粒子在一系列近似球形的 $A_{3}$ 团簇中心所感受到的能量井的深度。纵轴以 $\alpha m_{\mathrm{photon}}$ 为单位，横轴以 $\lambda_{min}$ 为单位。

作为对比的红线给出同样电荷\emph{均匀}分布于同一区域时井可能达到的最大深度。把黎曼库仑势对实心球积分所得出的精确深度是振荡的，在这一红线与相反的能量（此时井变成峰）之间来回循环。

由于并未采样每个团簇尺寸，能量随团簇增大的振荡被追踪得有些不均匀——但清晰的趋势是井深与团簇尺寸成正比，平均而言约为均匀分布电荷情形的两倍：
\begin{equation}
D_{centre}\approx \frac{8}{3\sqrt{3}\,\pi}\left(\frac{R}{\lambda_{min}}\right)\alpha\,m_{\mathrm{photon}} .
\end{equation}
这不是我们会在自己宇宙的晶体固体中期待的现象：在那里超过某一点后，晶体尺寸对深处电子所感受的条件毫无影响！人们可能以为这种长程影响必定源于“点阵完全由正电荷构成、没有负电荷抵消其累积效应”。但事实上，在点阵中加入负电荷在这一方面几乎不会有差别：它们将位于自己的势能极小点阵上（这些点在正电荷看来像是能量峰），而它们所产生的势能图样实际上会\emph{加深}正电荷的能量井。

\orthofig{images/07__004.png}{能量井形状与完美抛物线的对比}

能量井的深度很重要，但形状呢？任何有限团簇的井都不会完美地旋转对称，且在点阵内各处也不相同；但考虑一个理想化的旋转对称井——势能只依赖于到中心的距离——仍然有用。为近似井的形状，我们可以利用均匀带电球壳内部势能曲线的结果：总势是许多这样一项之和，而每一项都是下式的常数倍
\begin{equation}
U_{s}(r)=-\frac{\sin(\wms r)}{\wms r}.
\end{equation}
既然我们能把点阵近似成一系列带电球壳，中心井的形状必是这一 $U_{s}(r)$ 的某个固定倍数。

$U_{s}(r)$ 本身当然不是抛物线。右图中的蓝线给出 $U_{s}(r)$，而红线是在 $r=0$ 处与它具有相同二阶导数的抛物线。即使把它理想化为旋转对称，井形状的“不完美”也意味着：在这口井中来回运动的经典粒子不会只有单一纯振动频率。因此即使我们保证其振荡主频率大于 $\nu_{max}$，在经典物理中也还不足以阻止它通过运动的较低频率分量向外辐射。

\section{不完美谐振子（Imperfect Harmonic Oscillators）}

当粒子在一口能量井中运动、且（相对井底的）势能正比于到井中心距离的\emph{平方}时，把它拉向中心的力就简单地正比于距离。这样的系统称为谐振子（harmonic oscillator）。弹簧上振动的重物、重力下摆动的单摆都是例子——不过与多数真实物理系统一样，其能量井只是近似抛物，振荡较小时描述最佳。

假设有一口完美抛物、三维径向对称的能量井，粒子在其中无摩擦、无辐射、无任何复杂因素地运动。若力与距离之间的比例常数为 $k$，则势能 $U$ 与力 $\bF$ 为
\begin{align}
U&=\half k\,r^{2},\\
\bF&=-\partial_{r}U\,\be_{r}=-k\,r\,\be_{r}.
\end{align}
把这一径向力分解为沿 $x$、$y$、$z$ 轴的分量。若粒子坐标为 $(x,y,z)$，则径向单位矢量为
\begin{equation}
\be_{r}=\left(\frac{x}{r},\frac{y}{r},\frac{z}{r}\right),
\end{equation}
于是力的各分量为
\begin{equation}
(F_{x},F_{y},F_{z})=(-kx,-ky,-kz).
\end{equation}
利用“力等于质量乘加速度”求每个坐标对时间的二阶变化率，得
\begin{equation}
\bigl(\partial_{t}^{2}x(t),\partial_{t}^{2}y(t),\partial_{t}^{2}z(t)\bigr)
=\bigl(-\tfrac{k}{m}x(t),-\tfrac{k}{m}y(t),-\tfrac{k}{m}z(t)\bigr).
\end{equation}
这三个方程的解为
\begin{equation}
(x(t),y(t),z(t))=
\bigl(A_{x}\cos(\omega t+\phi_{x}),\,
A_{y}\cos(\omega t+\phi_{y}),\,
A_{z}\cos(\omega t+\phi_{z})\bigr),
\end{equation}
其中 $\omega=\sqrt{k/m}$。

\orthofig{images/07__005.png}{抛物与非抛物能量井中的轨道}
\orthofig{images/07__006.png}{做轨道运动的粒子的 $y$ 坐标}

粒子将沿每个坐标方向以相同的频率 $\omega$ 来回振荡。振荡振幅 $A_{x}$、$A_{y}$、$A_{z}$ 一般不同，相位 $\phi_{x}$、$\phi_{y}$、$\phi_{z}$ 也各异。但尽管这些参数有差异，由于\emph{频率}相同，经过时间 $2\pi/\omega$ 后粒子总会精确回到同一点。它走一条\emph{闭合轨道}（closed orbit），一般情形下是一个以井中心为中心的椭圆，如右图中的红色曲线。

若能量井不严格是抛物型呢？力将不再严格正比于 $r$，且力的各分量也不再只依赖于对应坐标。一切都变得更复杂，唯一闭合的轨道是完全没有径向运动的圆轨道。右上图中的蓝色曲线给出能量井为 $U_{s}(r)$（我们用于点阵中粒子势能的函数）时的一段轨道。

右下图为粒子运动的 $y$ 分量随时间的变化，以及描述该运动频率成分的傅里叶变换（Fourier transform）。轨道会进动（precession）——即离井中心最远的点随时间旋转，其角频率记为 $\omega_{p}$——这一事实把频谱中最强的峰从抛物井下的单一值 $\omega_{o}$（轨道频率）分裂为两个峰，位于 $\omega_{o}\pm\omega_{p}$。

现在，这些峰又高又窄（注意这些分量的强度是按对数标度绘制的！），因此若 $\omega_{o}-\omega_{p}>\wms$，粒子的运动就只有极少量处于能够辐射的频率上。但“极少量”并非零，随着时间推移粒子仍会获得能量，最终被驱出井。

\section{固体中的量子效应（Quantum Effects in Solids）}

到目前为止我们的处理完全是经典的，但要真正理解这一情形，需要用量子力学。黎曼宇宙中运动不太快的粒子的经典物理由牛顿力学描述，与在我们的宇宙中一样；而由于\emph{非相对论}量子力学是牛顿力学的量子版本，我们可以直接使用它的许多结果。特别地，对抛物能量井，量子谐振子的行为与我们宇宙中完全相同。

\orthofig{images/07__007.png}{能量本征函数}

量子力学中，粒子的态不是用确定的位置与速度描述，而是用波函数（wave function）。它给粒子可能所在的每个空间点联系一个复数，称为\emph{振幅}（amplitude）；在那里找到粒子的概率正比于振幅的模方。

具有确定能量的波函数，其形状随时间保持不变。右图给出一些例子。完整的波函数在三维空间中变化；这里只画出它们对到井中心距离 $r$ 的依赖关系。

数 $n$ 与 $l$ 是\emph{量子数}（quantum number），用来区分三维量子谐振子的各种可能波函数。量子数 $n$ 与粒子总能量相关：
\begin{equation}
E_{n}=\left(n+\tfrac{3}{2}\right)\hbar\omega .
\end{equation}
这里 $n$ 是非负整数，$\hbar$ 是普朗克常数除以 $2\pi$，$\omega$ 是等效经典谐振子的振荡角频率。于是粒子具有一串等间距的\emph{能级}（energy level）。最低可能的能量态 $n=0$ 并不位于井底，而是位于井底上方 $\tfrac{3}{2}$ 个 $\hbar\omega$ 处，其后所有能级间距均为一个 $\hbar\omega$。

量子数 $l$ 描述粒子具有的总角动量。从波函数图可见，随 $l$ 增大，粒子更可能出现在离井中心更远之处。然而谐振子的特别之处在于：$l$ 对粒子的总能量没有任何影响。量子数 $l$ 被限制在 $n$ 与 $0$ 之间，且当 $n$ 为偶数（奇数）时 $l$ 也为偶数（奇数）。第三个量子数通常记为 $m$，可取从 $-l$ 到 $l$ 的任意整数值，它规定沿某个选定坐标轴的角动量大小；$m$ 对能量与径向波函数都没有影响，但它与 $l$ 一起决定了波函数在偏离该轴不同角度上的变化方式。

总而言之，三维谐振子中粒子可占据 $(n+1)(n+2)/2$ 个不同的量子态，它们全部具有相同的能量 $E_{n}$。

\orthofig{images/07__008.png}{固体点阵井中的能级}

现在，当能量井不严格是抛物型时，所有波函数的形状都会有微小改变。其最重要的后果是：角动量量子数 $l$ 现在\emph{确实}影响粒子的总能量。于是抛物情形下由单一值构成的能级，现在会分裂成若干个彼此靠近的能级。

右图给出若干粒子质量与井深下的能级。每种情形下的红色能级对应完美抛物能量井；蓝色能级对应同样深度、形状为 $U_{s}(r)$（我们预期点阵中粒子势能的近似形状）的井。

量子力学中，带电粒子不再像经典物理描述的那样连续辐射。相反，若它处于能量为 $E_{A}$ 的某个态，而另有一个能量为 $E_{B}$ 的态可供其跃迁，它可以发射一个能量等于两者之差的光子，并从态 $A$ 跳到态 $B$。在我们宇宙中，原子中的电子就是这样发光的。但在黎曼物理中有两点不同。首先，若 $E_{A}$ 与 $E_{B}$ 采用“动能为正”的约定，则光子能量符号相反。因此只有当 $E_{B}>E_{A}$ 时才会发生跃迁，带电粒子在发射光子的同时跳到\emph{更高}的能级。

另一点不同是，黎曼物理中存在最大可能光子能量 $m_{\mathrm{photon}}$ 需要考虑。若 $E_{B}-E_{A}>m_{\mathrm{photon}}$，则单个光子无法携带足够能量促成跃迁。原则上粒子可以发射两个、三个或四个光子，但所需光子数越多，跃迁概率就越小。

结果是：若能量井中粒子所处的量子态，其上下一档态比它高出若干个最大光子能量，它的处境就是\emph{稳定的}。量子力学中很少有情形能永远稳定，但能隙越大，现状就越稳定。例如，图中第一口井里处于最低能量态的粒子，其上方能隙为 $5.2\,m_{\mathrm{photon}}$，因此需要同时发射\emph{六个}光子才能跃迁到更高态。

关键在于：从抛物能量井变成我们在点阵中预期的那种不完美形状，未必会破坏这种稳定性。能级会有一些移动与分裂，但对许多井而言，某些能级上方仍存在若干个最大光子能量的能隙。某个这样的能级\emph{下方}即使存在小得多的能隙也无妨——粒子可以从环境中吸收一个光子下到更低能级，再重新发射光子跳回原处；这类跃迁只是让固体与光发生相互作用，不会对其稳定性产生不利影响。（事实上，即便能隙小于 $m_{\mathrm{photon}}$，也存在一些微妙问题决定哪些跃迁能用单个光子完成。）

究竟哪些固体会在量子力学意义下稳定，取决于诸多参数，包括构成点阵的粒子质量以及点阵中能量井的深度。井深取决于点阵尺寸、黎曼静电力的强度以及每个粒子的电荷。我们此前的计算给出
\begin{equation}
D_{centre}\approx \frac{8}{3\sqrt{3}\,\pi}\left(\frac{R}{\lambda_{min}}\right)\alpha\,m_{\mathrm{photon}},
\end{equation}
其中 $R$ 是团簇半径，$D_{centre}$ 是团簇正中心处井的深度。我们将修改这一公式以计入两点小麻烦。第一，我们希望确保连团簇边缘的井也稳定，为此先把这一估计折半。第二，我们希望允许每口井中有多个粒子存在。因此把估计乘以 $W$，即每口井中的粒子数，因为每个粒子对贯穿整个点阵的场都贡献大致相同的份额
\begin{equation}
D\approx \frac{4W}{3\sqrt{3}\,\pi}\left(\frac{R}{\lambda_{min}}\right)\alpha\,m_{\mathrm{photon}} .
\end{equation}
当然，与我们宇宙中一样，一个现实的固体可能由许多小晶体组成，而非从一端到另一端保持完美点阵结构；此时 $R$ 应是其中一个晶体的半径，而非整块固体的尺度。

设点阵中典型能量井的深度为 $D$。势的最佳抛物近似为
\begin{equation}
U_{p}(r)=D\left[\frac{1}{6}\wms^{2}r^{2}-1\right].
\end{equation}
除了共享深度 $D$ 外，此函数在 $r=0$ 处还与点阵势 $U_{s}(r)$ 具有相同的二阶导数。此井中经典粒子所受的力为
\begin{equation}
\bF=-\partial_{r}U_{p}(r)\,\be_{r}=-\frac{D}{3}\wms^{2}r\,\be_{r},
\end{equation}
于是力常数 $k$ 等于 $\tfrac{D}{3}\wms^{2}$，质量为 $m$ 的粒子在井中的固有振荡频率为
\begin{equation}
\omega=\sqrt{\frac{k}{m}}=\wms\sqrt{\frac{D}{3m}} .
\end{equation}
这意味着井中能量台阶大小 $\hbar\omega$ 与最大光子能量 $m_{\mathrm{photon}}=\hbar\wms$ 之比为
\begin{equation}
\frac{\hbar\omega}{m_{\mathrm{photon}}}=\frac{\omega}{\wms}=\sqrt{\frac{D}{3m}} .
\end{equation}
因此，忽略能级的移动与分裂，点阵的稳定性将取决于 $\sqrt{D/(3m)}$ 的取值。若 $\sqrt{D/(3m)}\le 1$，点阵将\emph{不}稳定；而这个量越大，点阵越稳定。右图说明了这一效应：当粒子质量 $m$ 增大或 $D$ 减小时，不同 $n$ 对应的能级靠得更近。

若我们希望保证
\begin{equation}
\frac{\hbar\omega}{m_{\mathrm{photon}}}=\sqrt{\frac{D}{3m}}>S
\end{equation}
对某个“安全因子（safety factor）”$S$ 成立，则粒子质量与光子质量之比必须小于某个临界值，即 \textbf{(1a)}。若该比值固定，则只要点阵尺寸大于 \textbf{(1b)} 给出的阈值，就能达到同样水平的稳定性。

\begin{keybox}[点阵稳定性条件（Lattice stability conditions）]
\begin{equation}
\frac{m}{m_{\mathrm{photon}}}<\frac{4W\alpha}{9\sqrt{3}\,\pi S^{2}}\left(\frac{R}{\lambda_{min}}\right),\tag{1a}
\end{equation}
\begin{equation}
\frac{R}{\lambda_{min}}>\frac{9\sqrt{3}\,\pi S^{2}}{4W\alpha}\left(\frac{m}{m_{\mathrm{photon}}}\right).\tag{1b}
\end{equation}
\end{keybox}

粒子质量较小有利于稳定性，因为这增大了能级间距，需要更多光子才能跨越该间距。但质量也可能\emph{过}小。完美谐振子的能量井会一直延伸下去，井壁越升越高；而点阵中的井宽度与深度都有限。若能隙相对于井深过大，井中占据最高能级的粒子就有跳出井（而非跃迁到更高的束缚态）的风险。事实上，若能级间距过大，井中甚至可能装不下我们所需的 $W$ 个态。

假设构成点阵的粒子是费米子（fermion），因此与我们宇宙中的电子一样不能共占同一量子态，那么这 $W$ 个粒子将填满若干能级。量子谐振子中，对每个量子数 $n$ 有 $(n+1)(n+2)/2$ 个不同态；若再计入不同自旋的态则数量翻倍，因此到并包括任一给定 $n$ 为止的态总数为
\begin{equation}
\sum_{i=0}^{n}(i+1)(i+2)=\frac{(n+1)(n+2)(n+3)}{3}.
\end{equation}
对总共 $W$ 个粒子所占据的最高能级精确求解非常繁琐，但我们可以作粗略近似，取
\begin{equation}
n\approx \sqrt[3]{3W}.
\end{equation}
“第 $n$ 个能级能装进井内”这一要求给出了质量比的下限，即下面的 \textbf{(1c)}，以及对应的 $R$ 下限 \textbf{(1d)}。

现在，随着点阵增大，质量下限 \textbf{(1c)} 越来越低，而质量上限 \textbf{(1a)} 越来越高；因此无论 $\alpha$ 与 $S$ 取何值，足够大的点阵都会稳定。反过来，对足够小的点阵，我们会到达上限与下限相遇之处，此时两个条件无法同时满足。条件 \textbf{(1e)} 就是要求点阵足够大、以至于不会出现这种情况。

\begin{keybox}[点阵稳定性条件（续）（Lattice stability conditions{,} continued）]
\begin{equation}
\frac{m}{m_{\mathrm{photon}}}>\frac{3^{7/6}\pi}{4\alpha W^{1/3}}\left(\frac{\lambda_{min}}{R}\right),\tag{1c}
\end{equation}
\begin{equation}
\frac{R}{\lambda_{min}}>\frac{3^{7/6}\pi}{4\alpha W^{1/3}}\left(\frac{m_{\mathrm{photon}}}{m}\right),\tag{1d}
\end{equation}
\begin{equation}
\frac{R}{\lambda_{min}}>\frac{3^{11/6}\pi S}{4W^{2/3}\alpha}.\tag{1e}
\end{equation}
\end{keybox}

\orthofig{images/07__009.png}{点阵的稳定性条件}

举个例子，若 $\alpha=1/137$（我们自己宇宙中精细结构常数的值），且 $S=4$（即确保所有粒子都需同时发射至少四个光子才能跳出其基态），则条件 \textbf{(1e)} 要求
\begin{equation}
\frac{R}{\lambda_{min}}>\frac{3225}{W^{2/3}}.
\end{equation}
右图给出稳定点阵可能范围的粗略图景。我们忽略了能级分裂的效应（对最高能级它可能起失稳作用），也忽略了同一井中粒子之间的相互作用（它可能通过降低某些能级而改善稳定性）。尽管如此，此处显现的总体趋势应当成立：对给定粒子质量，\emph{向每口井中加入更多粒子}能使更小的点阵保持稳定。

\section{黎曼原子？（Riemannian Atoms?）}

讨论了含数百万粒子的点阵之后，现在来分析只用两个全同粒子构成的系统。原则上，它们可能以两种不同的方式束缚在一起。它们可能相隔大致 $\lambda_{min}$ 的距离，使各自处于对方的第一个“能量波谷”中；或者它们可能比 $\lambda_{min}/4$ 更近，从而处于彼此的“能量坑”中——即在静电力开始在吸引与排斥之间往复切换之前出现的、无底的静电吸引井。

但从我们对点阵的计算可以看出，若一对粒子处于彼此的能量波谷中，它们不太可能稳定。尽管具体情形与点阵不同，影响能量井深度与能级间距的一般因素却是相同的。对 $W=1$、$R=\lambda_{min}$ 的条件 \textbf{(1e)} 变为
\begin{equation}
\alpha>\frac{3^{11/6}\pi}{4}S\approx 5.89\,S .
\end{equation}
即便假设我们可以自由挑选黎曼宇宙中所有基本常数的值，若为了让此式成立而选取很大的 $\alpha$，那只会增加粒子在不同能级间跃迁的\emph{速率}，从而又要求我们选取更大的 $S$ 以保证稳定。

不过，若两个粒子靠得近得多，我们就能完全忽略库仑力中的波纹，把静电吸引当作平方反比力处理。这类量子力学系统的能级非常著名，因为除若干常数改变外，它们就是氢原子的能级。当然氢原子由一个带负电的电子被束缚在重得多、带正电的质子周围，但由平方反比吸引所决定的波函数数学在两种情形下本质相同。

若我们套用氢原子的量子力学，会发现最简单的两粒子黎曼原子的能级为
\begin{equation}
E_{1}=-\frac{m\alpha^{2}}{4},\qquad
E_{n}=\frac{E_{1}}{n^{2}} .
\end{equation}
这里 $m$ 是每个粒子的质量，$\alpha$ 是精细结构常数，量子数 $n$ 是大于等于 $1$ 的整数。对氢原子原则上 $n$ 没有上限，但对我们的黎曼原子，若 $n$ 过大，波函数将延伸到能量坑之外，我们所用的平方反比力近似就失效。事实上，我们至少需要到 $n=2$ 这一近似依然合理。波函数的距离尺度可通过套用氢原子的玻尔半径（Bohr radius）得到；我们场景下的合适取值为
\begin{equation}
a_{0}=\frac{2\hbar}{m\alpha}.
\end{equation}
要使 $n=2$ 的波函数足够局域化，需要
\begin{equation}
a_{0}<\frac{\lambda_{min}}{64}.
\end{equation}
这里的因子 $1/64$ 有些随意，但它使粒子间距大于 $\lambda_{min}/4$（势首次过零之处）的概率低于 $0.1\%$。这对应一个关于粒子质量与光子质量之比的条件
\begin{equation}
\frac{m}{m_{\mathrm{photon}}}>\frac{64}{\pi\alpha}.\tag{2a}
\end{equation}

那么，这个原子处于 $n=1$ 态时保持稳定的条件是什么？第一与第二能级之差为
\begin{equation}
\Delta E=E_{2}-E_{1}=-\frac{3}{4}E_{1}=\frac{3m\alpha^{2}}{16}.
\end{equation}
关键在于这一能量差包含多少个光子质量：
\begin{equation}
\frac{\Delta E}{m_{\mathrm{photon}}}=\frac{3\alpha^{2}}{16}\left(\frac{m}{m_{\mathrm{photon}}}\right).
\end{equation}
为使原子免于自发发射光子，我们要求此比值超过某个安全因子 $S$，于是
\begin{equation}
\frac{m}{m_{\mathrm{photon}}}>\frac{16S}{3\alpha^{2}}.\tag{2b}
\end{equation}
对 $S=4$、$\alpha=1/137$，这意味着 $m/m_{\mathrm{photon}}>400{,}000$！所以两粒子原子所需的粒子质量，远大于使尺寸合理的点阵稳定的那类粒子。

若原子中不是两个粒子而是 $Z$ 个粒子，能级将乘以 $Z^{2}$。然而，假设这些粒子是费米子，则这 $Z$ 个粒子将填满若干能级。量子数为 $n$ 的不同量子态有 $2n^{2}$ 个，因此最高的 $n$ 值由下式给出
\begin{equation}
Z=2(1+4+9+\dots+n^{2})=\frac{n(n+1)(2n+1)}{3}\approx\frac{2n^{3}}{3},
\end{equation}
\begin{equation}
n\approx\sqrt[3]{\frac{3Z}{2}}.
\end{equation}
（注意这些计算\emph{仅对大 $Z$ 是粗略近似}，不能用于我们的两粒子原子。）我们需要的不是在 $n=1$ 与 $n=2$ 之间留出安全能隙，而是在 $n$ 与 $n+1$ 之间：
\begin{align}
\Delta E&=Z^{2}(E_{n+1}-E_{n})\nonumber\\
&=Z^{2}(-E_{1})\left[\frac{1}{n^{2}}-\frac{1}{(n+1)^{2}}\right]\nonumber\\
&=Z^{2}\frac{m\alpha^{2}}{4}\frac{2n+1}{(n+1)^{2}n^{2}}\nonumber\\
&\approx 2Z^{2}\frac{m\alpha^{2}}{4n^{3}}\nonumber\\
&\approx \frac{Zm\alpha^{2}}{3}.
\end{align}
要求对此能隙有安全因子 $S$，意味着
\begin{equation}
\frac{m}{m_{\mathrm{photon}}}>\frac{3S}{Z\alpha^{2}}.\tag{3}
\end{equation}
对 $S=4$、$\alpha=1/137$，即便 $Z=1000$ 的原子也要求 $m/m_{\mathrm{photon}}>225$。

\section{大分子（Macromolecules）}

在固体晶体点阵与黎曼原子这两种情形之间，应当存在一系列可由任意给定质量粒子构成的结构。稳定性的最明显判据是：单个粒子的能级中包含上方有足够能隙将其隔离的态。我们已经看到，通过种种手段加深能量井可以做到这一点，但这样做可能需要把成百上千个粒子纳入结构之中。

为使粒子系统稳定，系统整体的能级也有必要具有同样的能隙，从而使“分子”的集体振动、弯折与旋转同样免于失控的光发射。转动态的能级正比于相应量子数的平方；由于整数平方之间的间距越来越大，转动能谱最终会在每个态上方形成极大能隙。振动与弯折运动则更为棘手：除了要求分子每个自由度各自的能谱有合适能隙外，还需要结构具有足够对称性，使整份能级清单不至于变成一片看似随机的、彼此紧挨的密集谱线——一组一千个全同谐振子、能隙都恰好为 $\Delta E$，其能谱的最小能隙就是 $\Delta E$；但一组一千个彼此能隙各异的谐振子，则可能通过让一个振子跃迁到更高能级、同时让另一个落到更低能级，从而获得比任何单个振子的量子更小的能量。

要精确刻画哪些分子能对所有这些威胁都保持稳定，将是计算量子化学中的一大壮举，因此我们只能勾画一个大致图景。大型的、类聚合物的、甚至可能呈片状的分子会有最好的机会；其结果可能是：除稀薄气体外，最常见的物质状态都具有相对高度长程有序。于是，黎曼宇宙中相当多的化学过程可能并非如我们熟悉的、在随机相互作用的小分子沸腾液体中进行，而更像在液晶（liquid crystal）中发生。点阵内部可见的那种规则能量井排列，实际上会延伸到点阵之外一段距离，许多稳定大分子也可能如此；它们的相互作用将表现为并排沿着能量地形中的锯齿彼此棘轮式地滑动、相互出入。

有关本主题的额外材料，提供给不介意稍高数学门槛的读者。

\section{进阶：单位（Units）}

本书通篇采用时间与距离用相同单位度量的体系。这等价于在我们宇宙中令光速 $c$ 等于 $1$（例如，以米度量距离，并用光走一米所需的时间作为相应的时间单位）。在黎曼宇宙中，这相当于选取时间单位使得勾股定理（Pythagoras's Theorem）成立，即使三角形的一边牵涉的是时间间隔而非空间。在小说《正交》（\emph{Orthogonal}）中，人们经验性地发现，这等同于令蓝光的速度等于 $1$。

在本节中我们将更进一步，选取质量与能量的单位使“约化普朗克常数” $\hbar=h/(2\pi)$ 等于 $1$。此时质量与能量用具有长度倒数（或空间频率）——也就是时间倒数（或时间频率）——量纲的单位来度量。

我们这一特殊选择意味着，频率 $\nu$ 与能量之间的普朗克关系 $E=h\nu$ 变为 $E=2\pi\nu=\omega$，其中 $\omega$ 是波的角频率；空间频率 $\kappa$ 与动量之间的关系则变为 $p=2\pi\kappa=k$，其中 $k$ 是角空间频率。黎曼标量波方程（Riemannian Scalar Wave equation）中出现的最大角频率 $\wms$ 于是就直接等于相应粒子的静止质量。

\section{进阶：相对论能量与动量算符（Relativistic Energy and Momentum Operators）}

在迄今讨论的非相对论量子力学中，我们只是把通常的薛定谔方程（Schrödinger equation）应用于带电粒子间作用力所对应的势能，其依据是：只要把动能视为正、并选取与之相容的势能符号，黎曼宇宙中的非相对论经典动力学就与牛顿力学完全相同。

对于黎曼\emph{相对论}量子力学，我们需要略作不同的处理。量子力学的通常表述与相应经典力学的哈密顿形式（Hamiltonian form）紧密相关；而在黎曼情形中，哈密顿意义下与各坐标共轭的动量，是同一方向上相对论动量的\emph{相反数}。

在非相对论量子力学中，我们把能量等同于算符 $\ii\partial_{t}$，但把动量的空间分量等同于算符 $-\ii\partial_{x}$、$-\ii\partial_{y}$、$-\ii\partial_{z}$。洛伦兹相对论量子力学中也是如此；但在黎曼情形中这说不通：我们需要把空间与时间坐标等同对待。为了保持与哈密顿经典力学的通常对应，我们将保留对应关系
\begin{equation}
E\to \ii\partial_{t},
\end{equation}
并把相对论动量等同于其它坐标上同类算符，而\emph{不带}通常的负号：
\begin{equation}
p_{x}\to \ii\partial_{x},\qquad p_{y}\to \ii\partial_{y},\qquad p_{z}\to \ii\partial_{z}.
\end{equation}
这也会影响我们用于角动量的算符。若用我们的黎曼相对论动量 $\bp$ 定义三维矢量算符 $\mathbf{L}=\mathbf{r}\times\bp$，它变为
\begin{align}
\mathbf{L}&=(x,y,z)\times(\ii\partial_{x},\ii\partial_{y},\ii\partial_{z})\nonumber\\
&=\ii\bigl(y\partial_{z}-z\partial_{y},\ z\partial_{x}-x\partial_{z},\ x\partial_{y}-y\partial_{x}\bigr).
\end{align}
与线动量一样，它也与通常算符相差一个负号。因此，其各分量之间的对易子（commutator）与非相对论量子力学中的不同：
\begin{equation}
[L_{i},L_{j}]=-\ii\varepsilon_{ijk}L_{k},
\end{equation}
其中列维-奇维塔符号（Levi-Civita symbol）$\varepsilon_{ijk}$ 当 $(i,j,k)$ 是 $(x,y,z)$ 的偶排列时等于 $1$，奇排列时等于 $-1$，含重复坐标时等于 $0$。同样，此式右端与通常情形相差一个符号。

现在，为使粒子的\emph{内禀}（intrinsic）角动量与这整套方案相容，相应的算符必须与轨道角动量 $\mathbf{L}$ 具有相同的对易代数。这意味着，对自旋 $\tfrac{1}{2}$ 粒子角动量各分量而言，合适的算符不再是泡利自旋矩阵（Pauli spin matrices），而是其\emph{相反数}。细节将随出现而展开；此处只是提醒：要对通常约定中的这一细小扭转有心理准备。

\section{进阶：黎曼 Dirac 方程（The Riemannian Dirac Equation）}

在我们自身的宇宙中，Dirac 方程（Dirac equation）是描述电子、正电子这类粒子的相对论性波动方程：这些粒子具有质量、电荷以及 $\half$ 自旋。这类粒子的波函数有四个分量，但它们并不是时空矢量的分量，而且在洛伦兹群的对称变换下，其变换方式与时空矢量截然不同。黎曼宇宙中对应的方程极为相似，事实上用\textbf{伽马矩阵}（gamma matrices）表述时可以写成完全一样的形式：伽马矩阵是四个 $4\times4$ 矩阵，作用在波函数分量上。但两类时空的号差（signature）不同，所要求的伽马矩阵的精确性质也不同。

Dirac 方程本身为
\begin{equation}
i\gamma^{\mu}\partial_{\mu}\psi - m\psi = 0 \tag{Dirac Equation}
\end{equation}
其中 $m$ 为粒子质量，$\psi$ 为四分量波函数，$I_{4}$ 为 $4\times4$ 单位矩阵，且我们采用 $\hbar=1$ 的单位制。我们对伽马矩阵施加的条件是
\begin{equation}
\{\gamma^{\mu},\gamma^{\nu}\}\equiv\gamma^{\mu}\gamma^{\nu}+\gamma^{\nu}\gamma^{\mu}=2g^{\mu\nu}I_{4} \tag{Common}
\end{equation}
其中在黎曼情形 $g^{\mu\nu}=\delta^{\mu\nu}$ 为
\begin{equation}
g^{\mu\nu}=\delta^{\mu\nu}=
\begin{pmatrix}
1&0&0&0\\
0&1&0&0\\
0&0&1&0\\
0&0&0&1
\end{pmatrix}
\qquad(\text{黎曼})
\end{equation}
而在洛伦兹情形 $g^{\mu\nu}=\eta^{\mu\nu}$，采用 $(+{}-{}-{}-)$ 号差：
\begin{equation}
\eta^{\mu\nu}=
\begin{pmatrix}
1&0&0&0\\
0&-1&0&0\\
0&0&-1&0\\
0&0&0&-1
\end{pmatrix}
\qquad(\text{洛伦兹})
\end{equation}

\begin{notebox}[关于度规号差的约定]
在本编注的绝大部分内容中，洛伦兹情形的类空维度取正号，从而 $\eta$ 具有 $(-+++)$ 号差。但由于几乎所有的量子场论文献都使用 $(+---)$ 号差，为便于比较，此处我们改用该约定。
\end{notebox}

在 Dirac 方程中，我们对指标 $\mu$ 施行爱因斯坦求和约定（Einstein summation convention）；$\mu$ 遍历四个时空维度，把四个伽马矩阵分别与对应方向上的导数组合起来。时空任意一点处的波函数 $\psi$ 是一个四分量复矢量，每个伽马矩阵按普通的矩阵乘法作用于 $\psi$ 的一个导数。但 $\psi$ 的四个分量\textbf{并不}由四个时空维度来编号。我们可以用若干附加指标（\textbf{旋量}（spinor）指标，区别于 $\mu$ 这样的时空指标）显式地写出这个矩阵乘法：
\begin{equation}
i(\gamma^{\mu})^{s}_{\,r}\,\partial_{\mu}\psi^{r}-m\psi^{s}=0
\end{equation}
这里求和约定同时应用于 $r$ 与 $\mu$；对旋量指标 $s$ 的四个取值，我们显式地得到四个方程。

在陈述对伽马矩阵施加的条件时，记号 $\{\gamma^{\mu},\gamma^{\nu}\}$ 称为一个\textbf{反交换子}（anticommutator）。它只是一个方便的简写，其定义为
\begin{equation}
\{A,B\}\equiv AB+BA
\end{equation}
在定义反交换子时 $A$、$B$ 可以是任意对象，但对我们而言它们通常是矩阵或其它线性算符（例如微分算符）。我们很快还会用到更常见的\textbf{交换子}（commutator）记号：
\begin{equation}
[A,B]\equiv AB-BA
\end{equation}

若 $\psi$ 是黎曼 Dirac 方程的一个解，我们便有
\begin{align*}
i\gamma^{\mu}\partial_{\mu}\psi - m\psi &= 0,\\
(i\gamma^{\mu}\partial_{\mu}-m)\psi &= 0,\\
(-i\gamma^{\nu}\partial_{\nu}-m)(i\gamma^{\mu}\partial_{\mu}-m)\psi &= 0,\\
(\gamma^{\nu}\gamma^{\mu}\partial_{\nu}\partial_{\mu}+m^{2})\psi &= 0,\\
\bigl(\tfrac12(\gamma^{\mu}\gamma^{\nu}+\gamma^{\nu}\gamma^{\mu})\partial_{\nu}\partial_{\mu}+m^{2}\bigr)\psi &= 0,\\
(\delta^{\nu\mu}\partial_{\nu}\partial_{\mu}+m^{2})\psi &= 0 .
\end{align*}
倒数第三行到倒数第二行之所以成立，是因为尽管\textbf{单项} $\gamma^{\mu}\gamma^{\nu}\neq\gamma^{\nu}\gamma^{\mu}$，但由求和约定得到的、遍历所有此类项的总和在交换 $\mu$ 与 $\nu$ 时是对称的。最后一行告诉我们：四分量场 $\psi$ 的每个分量都将服从黎曼标量波方程（Riemannian Scalar Wave equation）。

\subsection{黎曼伽马矩阵（Riemannian Gamma Matrices）}

满足条件 $\{\gamma^{\mu},\gamma^{\nu}\}=2\delta^{\mu\nu}I_{4}$ 的伽马矩阵有无穷多种选法，因为给定任何一组满足该性质的四个矩阵，我们总可以换到 $\mathbb{C}^{4}$（这些矩阵所作用的四维复空间）的另一组基，从而得到四个新矩阵。

但某些选法会使事情更简单。我们将给出两个例子，即所谓的 Weyl 基（Weyl basis）与 Dirac 基（Dirac basis）。首先，回顾一下在角动量与三维转动的量子力学中出现的泡利自旋矩阵（Pauli spin matrices）是方便的：
\begin{equation}
\sigma_{x}=
\begin{pmatrix}0&1\\1&0\end{pmatrix},\qquad
\sigma_{y}=
\begin{pmatrix}0&-i\\ i&0\end{pmatrix},\qquad
\sigma_{z}=
\begin{pmatrix}1&0\\0&-1\end{pmatrix}
\qquad(\text{Common})
\end{equation}
半个泡利矩阵之间的交换子为
\begin{equation}
[\tfrac12\sigma_{i},\tfrac12\sigma_{j}]=i\varepsilon_{ijk}\,\tfrac12\sigma_{k}.
\end{equation}
这里的 $\varepsilon_{ijk}$ 所做的是：挑出与 $i$、$j$ 都不同的那个 $k$，并选取适当的符号，从而给出
\begin{equation}
[\tfrac12\sigma_{x},\tfrac12\sigma_{y}]=i\,\tfrac12\sigma_{z},
\end{equation}
以及它的循环置换。正如我们在关于李代数 $\mathrm{so}(3)$ 的一节中所讨论的，绕坐标轴转动的生成元 $J_{i}$ 具有交换子代数
\begin{equation}
[J_{i},J_{j}]=\varepsilon_{ijk}J_{k}.
\end{equation}
在描述李代数 $\mathrm{su}(2)$ 与 $\mathrm{so}(3)$ 之间的同构时，我们用了 $\mathrm{su}(2)$ 的一组基 $\{H_{x},H_{y},H_{z}\}$，由三个无迹反厄米矩阵（其伴随等于原矩阵的相反数）组成。在做量子力学时，更常见的是使用厄米矩阵，因为量子力学中任何实值可观测量都对应一个厄米算符。泡利矩阵是厄米矩阵，与我们那组 $\mathrm{su}(2)$ 基的关系为
\begin{equation}
\sigma_{k}=iH_{k},\qquad H_{k}=-i\sigma_{k}.
\end{equation}
由于 $H_{k}$ 是 $\mathrm{su}(2)$ 的元素，它们的任意实系数线性组合都可以取指数，得到 $\mathrm{SU}(2)$ 的一个元素。厄米矩阵 $\tfrac12\sigma_{k}$ 就是那些本征矢为确定自旋态的量子力学算符（对非相对论量子力学中的 $\half$ 自旋粒子而言），相应的本征值为 $\pm\tfrac12$。若想把一个厄米矩阵取指数后得到一个酉算符（即保持量子波函数范数的算符），我们需要先把它乘以 $\pm i$，例如
\begin{equation}
U_{x}(\theta)=\exp(-i\theta\sigma_{x}/2),
\end{equation}
这是非相对论量子力学中 $\half$ 自旋粒子绕 $x$ 轴转动角 $\theta$ 的酉算符。

当我们从三维转到四维时，李代数 $\mathrm{so}(4)$ 有六个生成元 $J_{\mu\nu}$，对应由坐标方向 $\mu$ 与 $\nu$ 张成的六个平面内的转动。为固定矩阵 $J_{\mu\nu}$ 的符号，我们采用如下约定：把 $J_{\mu\nu}$ 的正数倍取指数，得到 $\mathrm{SO}(4)$ 中把 $\mu$ 轴转向 $\nu$ 轴的元素。这些生成元的交换子为
\begin{equation}
[J_{\mu\nu},J_{\rho\nu}]=J_{\mu\rho}
\end{equation}
（适用于两平面恰好共用一根坐标轴 $\nu$ 的情形）。除去此式及其带适当符号变化的置换之外，所有其它交换子都为零。这是合理的：若两平面\textbf{没有}公共轴，则在两者中的转动完全独立；若它们共享两根轴，则二者是同一平面，而同一平面内的转动可交换。

从非相对论量子力学转向\textbf{黎曼}相对论量子力学时，还有一个关键改动。正如我们需要对线动量使用不同于通常的算符，我们也需要对角动量这样做。我们此前讨论过，对于轨道角动量，只要用线动量来定义它，这一改动会自动出现；但对于内禀角动量，我们需要选取通常泡利矩阵的\textbf{相反数}，才能得到粒子自旋的正确符号。我们将在关于自旋与角动量守恒的一节中看到这一点的确认。

我们的第一组黎曼伽马矩阵可以借助泡利矩阵写成如下形式（\textbf{Weyl 基}）：
\begin{align}
\gamma^{t} &=
\begin{pmatrix}0&I_{2}\\ I_{2}&0\end{pmatrix}
=\sigma_{x}\otimes I_{2}, \nonumber\\
\gamma^{x} &=
\begin{pmatrix}0&-i\sigma_{x}\\ i\sigma_{x}&0\end{pmatrix}
=\sigma_{y}\otimes\sigma_{x}, \nonumber\\
\gamma^{y} &=
\begin{pmatrix}0&-i\sigma_{y}\\ i\sigma_{y}&0\end{pmatrix}
=\sigma_{y}\otimes\sigma_{y}, \nonumber\\
\gamma^{z} &=
\begin{pmatrix}0&-i\sigma_{z}\\ i\sigma_{z}&0\end{pmatrix}
=\sigma_{y}\otimes\sigma_{z}, \nonumber\\
\gamma^{\mu} &=
\begin{pmatrix}0&H_{\mu}\\ H_{\mu}^{*}&0\end{pmatrix}
\tag{Riemannian}
\end{align}
这里记号 $\sigma_{x}\otimes I_{2}$ 表示把第一个矩阵 $\sigma_{x}$ 中的每个元素替换为该元素与第二个矩阵 $I_{2}$ 的乘积，从而得到一个 $4\times4$ 矩阵。最后一行给出了适用于任意指标 $\mu$ 的一般形式（此时我们按之前的定义取 $H_{t}=I_{2}$；但注意，与另外三个 $H$ 矩阵不同，$H_{t}$ 不属于 $\mathrm{su}(2)$）。

这些伽马矩阵全都是厄米的、无迹的，且行列式为 $1$。每个矩阵的平方都是单位矩阵，而所有不相同的矩阵对都反交换：即它们的反交换子为零。若把它们与 Peskin 和 Schroeder 那本广为使用的量子场论教科书中 Weyl 基的伽马矩阵选择\textsuperscript{[1]}相比较，我们的 $\gamma^{t}$ 相同，而这里其它矩阵是洛伦兹版本的 $-i$ 倍。这反映了这样一个事实：Peskin 与 Schroeder 对度规采用 $(+---)$ 号差，故按我们的约定，他们的 $t$ 坐标已具有类空标志差。

我们可以用这些伽马矩阵导出群 $\mathrm{SO}(4)$ 的一个表示，所用公式取自文献\textsuperscript{[1]}：
\begin{equation}
S^{\mu\nu}=(i/4)\,[\gamma^{\mu},\gamma^{\nu}].
\end{equation}
这里每个 $S^{\mu\nu}$ 都是作用在波函数 $\psi$ 在各点处所属四维复矢量空间上的厄米矩阵。它是 $\mathrm{so}(4)$ 元素 $J_{\mu\nu}$ 在某个尚待确定的四维李代数表示中的像，再乘以一个使其成为厄米矩阵所需的因子 $\pm i$。当然，在这里看到泡利矩阵的相反数并不令我们意外，因为我们知道那个负号对于使内禀自旋与其它一切相容是必要的。

在 Weyl 基中，相应的（黎曼）自旋矩阵为
\begin{align}
S^{xy} &=
\begin{pmatrix}-\sigma_{z}/2&0\\ 0&-\sigma_{z}/2\end{pmatrix}
=-\tfrac12 I_{2}\otimes\sigma_{z}, \nonumber\\
S^{yz} &=
\begin{pmatrix}-\sigma_{x}/2&0\\ 0&-\sigma_{x}/2\end{pmatrix}
=-\tfrac12 I_{2}\otimes\sigma_{x}, \nonumber\\
S^{zx} &=
\begin{pmatrix}-\sigma_{y}/2&0\\ 0&-\sigma_{y}/2\end{pmatrix}
=-\tfrac12 I_{2}\otimes\sigma_{y}, \nonumber\\
S^{xt} &=
\begin{pmatrix}\sigma_{x}/2&0\\ 0&-\sigma_{x}/2\end{pmatrix}
=\tfrac12\sigma_{z}\otimes\sigma_{x}, \nonumber\\
S^{yt} &=
\begin{pmatrix}\sigma_{y}/2&0\\ 0&-\sigma_{y}/2\end{pmatrix}
=\tfrac12\sigma_{z}\otimes\sigma_{y}, \nonumber\\
S^{zt} &=
\begin{pmatrix}\sigma_{z}/2&0\\ 0&-\sigma_{z}/2\end{pmatrix}
=\tfrac12\sigma_{z}\otimes\sigma_{z}
\tag{Riemannian}
\end{align}
从这些矩阵是分块对角的这一事实可以看出，它们所定义的四维表示是\textbf{可约的}（reducible）：它可以分裂为两个二维子空间，该表示的作用永远不会把它们混合起来。

那么，这里被粘在一起的两个 $\mathrm{so}(4)$ 表示是什么呢？考察我们此前描述过的自旋-$(\tfrac12,0)$ 与自旋-$(0,\tfrac12)$ 表示（其详细表格列于那一节末尾），我们看到：作用于 $\psi$ 前两个分量的表示对应于自旋-$(\tfrac12,0)$，作用于后两个分量的表示对应于自旋-$(0,\tfrac12)$；唯一的差别是那一讨论中 $\mathrm{su}(2)$ 的 $H$ 矩阵被换成了泡利 $\sigma$ 矩阵的相反数。

李代数表示只是群表示的导数，因此 $\mathrm{SO}(4)$ 的双覆盖应通过群表示自旋-$(\tfrac12,0)\oplus$自旋-$(0,\tfrac12)$ 作用在 $\psi$ 上。而正如我们在讨论包含宇称的表示时所见，这可以推广为 $\mathrm{O}(4)$ 双覆盖的一个表示，其中宇称运算（把 $x$、$y$、$z$ 变号）会交换该表示的两个二维子空间。

若用 $\rho$ 作为自旋-$(\tfrac12,0)\oplus$自旋-$(0,\tfrac12)$ 表示的缩写，则（在来源于 $\mathrm{SO}(4)$ 双覆盖中两个元素之选择的符号不确定度之内）有 $\rho:\mathrm{SO}(4)\to\mathrm{SU}(4)$。也就是说，在不计符号的意义下，$\rho$ 为 $\mathrm{SO}(4)$ 中的每个转动指定一个行列式为 $1$ 的 $4\times4$ 酉矩阵。熟悉洛伦兹情形 Dirac 方程的读者会记得，$\mathrm{SO}(3,1)$ 的对应表示中，“推动”（boost，即涉及相对运动的坐标变换）的表示\textbf{不是}酉的；但在黎曼情形中我们不面临这一麻烦。

显式地，$\rho$ 把 $\mathrm{SO}(4)$ 中的转动 $R$ 映为 $\mathrm{SU}(4)$ 中的酉算符 $\rho(R)$，使得
\begin{equation}
\rho(\exp(J^{\mu\nu}))=\exp(iS^{\mu\nu}).
\end{equation}
该映射的导数是 $\mathrm{so}(4)$ 到 $\mathrm{su}(4)$ 的李代数同构：
\begin{equation}
d\rho(J^{\mu\nu})=iS^{\mu\nu}.
\end{equation}
有了这个定义，就不难验证：对 $\mathrm{SO}(4)$ 中的任意转动 $R$，
\begin{equation}
\rho(R^{-1})\,\gamma^{\nu}\,\rho(R)=R^{\nu}_{\ \zeta}\,\gamma^{\zeta} \tag{1}
\end{equation}
或者，用指标记号把这个矩阵乘法完全展开：
\begin{equation}
\rho(R^{-1})^{q}_{\,p}\,(\gamma^{\nu})^{p}_{\,r}\,\rho(R)^{r}_{\,s}
=(R^{\nu}_{\ \zeta}\gamma^{\zeta})^{q}_{\,s} \tag{2}
\end{equation}
验证此式最容易的办法是求导，并利用公式 $(J^{\alpha\beta})^{\nu}_{\ \zeta}=\delta^{\alpha\zeta}\delta^{\beta\nu}-\delta^{\alpha\nu}\delta^{\beta\zeta}$：
\begin{align*}
[\gamma^{\nu},S^{\alpha\beta}] &= -i\,(J^{\alpha\beta})^{\nu}_{\ \zeta}\,\gamma^{\zeta},\\
[\gamma^{\nu},(i/4)[\gamma^{\alpha},\gamma^{\beta}]] &= i\,(\delta^{\alpha\nu}\gamma^{\beta}-\delta^{\beta\nu}\gamma^{\alpha}),\\
\tfrac14[\gamma^{\nu},[\gamma^{\alpha},\gamma^{\beta}]] &= \delta^{\alpha\nu}\gamma^{\beta}-\delta^{\beta\nu}\gamma^{\alpha},\\
\tfrac14[\gamma^{\nu},\,2\delta^{\alpha\beta}I_{4}-2\gamma^{\beta}\gamma^{\alpha}]
&= \tfrac14\bigl((\gamma^{\alpha}\gamma^{\nu}+\gamma^{\nu}\gamma^{\alpha})\gamma^{\beta}
+\gamma^{\beta}(\gamma^{\alpha}\gamma^{\nu}+\gamma^{\nu}\gamma^{\alpha})\\
&\qquad-(\gamma^{\beta}\gamma^{\nu}+\gamma^{\nu}\gamma^{\beta})\gamma^{\alpha}
-\gamma^{\alpha}(\gamma^{\beta}\gamma^{\nu}+\gamma^{\nu}\gamma^{\beta})\bigr),\\
-\tfrac12(\gamma^{\nu}\gamma^{\beta}\gamma^{\alpha}-\gamma^{\beta}\gamma^{\alpha}\gamma^{\nu})
&= \tfrac14\bigl(\gamma^{\nu}\gamma^{\alpha}\gamma^{\beta}
+\gamma^{\beta}\gamma^{\alpha}\gamma^{\nu}\\
&\qquad-\gamma^{\nu}\gamma^{\beta}\gamma^{\alpha}
-\gamma^{\alpha}\gamma^{\beta}\gamma^{\nu}\bigr),\\
0 &= \gamma^{\nu}\gamma^{\alpha}\gamma^{\beta}
+\gamma^{\nu}\gamma^{\beta}\gamma^{\alpha}
-\gamma^{\alpha}\gamma^{\beta}\gamma^{\nu}
-\gamma^{\beta}\gamma^{\alpha}\gamma^{\nu},\\
0 &= \gamma^{\nu}(2\delta^{\alpha\beta}I_{4})-(2\delta^{\alpha\beta}I_{4})\gamma^{\nu},\\
0 &= 0 .
\end{align*}
倒着走这些步骤，就得到原命题的推导。

方程 (1) 与 (2) 告诉我们：若用 $\rho(R)$ 变换 $\gamma$ 矩阵的旋量坐标，其结果等同于把这四个 $\gamma$ 矩阵彼此之间作旋转——就好像它们是受 $\mathrm{SO}(4)$ 元素 $R$ 作用的时空矢量的分量一样。

这样做的目的在于理解当我们旋转坐标时 Dirac 方程会发生什么。若我们做坐标变换，从 $x^{\mu}$ 变到 $x'^{\nu}=R^{\nu}_{\ \mu}x^{\mu}$，则坐标导数 $\partial_{\mu}$ 被替换为 $\partial'_{\nu}=(R^{-1})^{\mu}_{\ \nu}\partial_{\mu}$，而旋量波函数 $\psi^{s}$ 被变换为 $\psi'^{p}=\rho(R)^{p}_{\,s}\psi^{s}$。把这些都用新坐标写出 Dirac 方程，有
\begin{equation}
i(\gamma^{\nu})^{p}_{\,r}\,\partial'_{\nu}\psi'^{r}-m\psi'^{p}=0. \tag{3}
\end{equation}
把变换后的导数与波函数用原来的量表示代入：
\begin{equation}
i(\gamma^{\nu})^{p}_{\,r}(R^{-1})^{\mu}_{\ \nu}\,\partial_{\mu}\,\rho(R)^{r}_{\,s}\psi^{s}
-m\rho(R)^{p}_{\,s}\psi^{s}=0. \tag{4}
\end{equation}
在左边乘以 $\rho(R^{-1})$，得
\begin{align*}
&\rho(R^{-1})^{q}_{\,p}\Bigl[i(\gamma^{\nu})^{p}_{\,r}(R^{-1})^{\mu}_{\ \nu}\partial_{\mu}\rho(R)^{r}_{\,s}\psi^{s}
-m\rho(R)^{p}_{\,s}\psi^{s}\Bigr]=0,\\
&i\rho(R^{-1})^{q}_{\,p}(\gamma^{\nu})^{p}_{\,r}(R^{-1})^{\mu}_{\ \nu}\partial_{\mu}\rho(R)^{r}_{\,s}\psi^{s}
-m\rho(R^{-1})^{q}_{\,p}\rho(R)^{p}_{\,s}\psi^{s}=0,\\
&i\rho(R^{-1})^{q}_{\,p}(\gamma^{\nu})^{p}_{\,r}(R^{-1})^{\mu}_{\ \nu}\partial_{\mu}\rho(R)^{r}_{\,s}\psi^{s}
-m\psi^{q}=0 .
\end{align*}
利用 \textbf{(2)}，上式变为
\begin{align}
&i(R^{\nu}_{\ \zeta}\gamma^{\zeta})^{q}_{\,s}(R^{-1})^{\mu}_{\ \nu}\partial_{\mu}\psi^{s}-m\psi^{q}=0, \nonumber\\
&i(\gamma^{\mu})^{q}_{\,s}\partial_{\mu}\psi^{s}-m\psi^{q}=0. \tag{5}
\end{align}
方程 \textbf{(5)} 就是原来 $\psi$ 的 Dirac 方程。于是我们证明了 Dirac 方程是 \textbf{$\mathrm{SO}(4)$ 不变的}：用 $R$ 旋转坐标后得到的解，正是原方程的解经 $\rho(R)$ 变换后的结果。

值得强调的一点是：当我们变换 $\psi$ 时，并不把 $\gamma$ 矩阵当作某种本身也该被变换的场。相反，伽马矩阵把导数的时空坐标变换与 $\psi$ 相应的旋量变换精巧地结合起来，使得所有变化彼此抵消。理解伽马矩阵的最佳方式是把它看成一个\textbf{缠绕算符}（intertwiner）：它是从一个矢量空间到另一个矢量空间的线性映射，且与同一群在这两个空间上的表示相交换。项 $\partial_{\mu}\psi^{s}$ 在两个表示的张量积下变换：一个是来自时空坐标导数的、基本矢量表示的对偶表示，另一个是自旋-$(\tfrac12,0)\oplus$自旋-$(0,\tfrac12)$ 的 Dirac 旋量表示。把方程 \textbf{(2)} 的两边都乘以 $(R^{-1})^{\mu}_{\ \nu}\rho(R)^{w}_{\,q}$，重新排列某些项的顺序（此处可以随意进行，因为带全部指标的项只是标量），再把所得结果的两边都作用于 $\partial_{\mu}\psi^{s}$，我们得到
\begin{equation}
(\gamma^{\nu})^{w}_{\,r}\,\rho(R)^{r}_{\,s}\,(R^{-1})^{\mu}_{\ \nu}\,\partial_{\mu}\psi^{s}
=\rho(R)^{w}_{\,q}\,(\gamma^{\mu})^{q}_{\,s}\,\partial_{\mu}\psi^{s}. \tag{6}
\end{equation}
在右边，$\gamma$ 先把张量积中的一个项 $\partial_{\mu}\psi^{s}$ 映为带指标 $q$ 的 Dirac 旋量，然后 $\rho(R)$ 作用于它，给出另一个带指标 $w$ 的 Dirac 旋量。在左边，$R^{-1}$ 作用于对偶坐标指标 $\mu$，而 $\rho(R)$ 作用于旋量坐标 $s$，\textbf{然后} $\gamma$ 接过结果（此时它仍位于张量积空间中），把它映为带指标 $w$ 的 Dirac 旋量。我们既可以先作用表示、再作用映射 $\gamma$，也可以先作用 $\gamma$、再作用表示，而得到完全相同的结果，这正是我们说 $\gamma$ 是一个缠绕算符的含义。

我们选择 Weyl 基作为伽马矩阵，是为了尽可能容易地看出 $\psi$ 所受的是 $\mathrm{SO}(4)$ 的哪个表示。另一个将用于下一节的基是 Dirac 基。在这个基中，只有 $\gamma^{t}$ 改变，其它矩阵与 Weyl 基中完全相同（\textbf{Dirac 基}）：
\begin{align}
\gamma^{t} &=
\begin{pmatrix}I_{2}&0\\ 0&-I_{2}\end{pmatrix}
=\sigma_{z}\otimes I_{2}, \nonumber\\
\gamma^{x} &=
\begin{pmatrix}0&-i\sigma_{x}\\ i\sigma_{x}&0\end{pmatrix}
=\sigma_{y}\otimes\sigma_{x}, \nonumber\\
\gamma^{y} &=
\begin{pmatrix}0&-i\sigma_{y}\\ i\sigma_{y}&0\end{pmatrix}
=\sigma_{y}\otimes\sigma_{y}, \nonumber\\
\gamma^{z} &=
\begin{pmatrix}0&-i\sigma_{z}\\ i\sigma_{z}&0\end{pmatrix}
=\sigma_{y}\otimes\sigma_{z}
\tag{Riemannian}
\end{align}
由此得到的自旋矩阵在只涉及空间坐标时不变，因此它们同样是分块对角的；但生成涉及 $t$ 坐标之转动的那些矩阵会混合旋量的两半。这并不意味着 $\mathrm{SO}(4)$ 的表示变了：Dirac 基中的表示与 Weyl 基中的表示等价。Dirac 基中的（黎曼）自旋矩阵为
\begin{align}
S^{xy} &=
\begin{pmatrix}-\sigma_{z}/2&0\\ 0&-\sigma_{z}/2\end{pmatrix}
=-\tfrac12 I_{2}\otimes\sigma_{z}, \nonumber\\
S^{yz} &=
\begin{pmatrix}-\sigma_{x}/2&0\\ 0&-\sigma_{x}/2\end{pmatrix}
=-\tfrac12 I_{2}\otimes\sigma_{x}, \nonumber\\
S^{zx} &=
\begin{pmatrix}-\sigma_{y}/2&0\\ 0&-\sigma_{y}/2\end{pmatrix}
=-\tfrac12 I_{2}\otimes\sigma_{y}, \nonumber\\
S^{xt} &=
\begin{pmatrix}0&-\sigma_{x}/2\\ -\sigma_{x}/2&0\end{pmatrix}
=-\tfrac12\sigma_{x}\otimes\sigma_{x}, \nonumber\\
S^{yt} &=
\begin{pmatrix}0&-\sigma_{y}/2\\ -\sigma_{y}/2&0\end{pmatrix}
=-\tfrac12\sigma_{x}\otimes\sigma_{y}, \nonumber\\
S^{zt} &=
\begin{pmatrix}0&-\sigma_{z}/2\\ -\sigma_{z}/2&0\end{pmatrix}
=-\tfrac12\sigma_{x}\otimes\sigma_{z}
\tag{Riemannian}
\end{align}

\subsection{平面波解（Plane Wave Solutions）}

\subsubsection{Dirac 基下的平面波解}

假设我们寻找如下形式的 Dirac 方程平面波解：
\begin{equation}
\psi(x)=u\,\exp(-ik\cdot x),
\end{equation}
其中 $\bx$ 是坐标四矢量。在我们的单位制下，传播矢量 $\bk$ 恰与平面波对应的能量-动量矢量相同，当然 $|\bk|=m$。旋量 $u$ 是 $\mathbb{C}^{4}$ 的常元素。

Dirac 方程随即变为关于 $u$ 的线性方程：
\begin{equation}
(k_{\mu}\gamma^{\mu}-mI_{4})u=0.
\end{equation}
只要 $|\bk|=m$，这里的矩阵行列式总为零，因此总存在非平凡解。找解最容易的办法是取 $\bk=m\be_{t}$，即静止波的能量-动量矢量。在 Dirac 基中，这给出
\begin{equation}
\begin{pmatrix}0&0\\ 0&-2mI_{2}\end{pmatrix}u=0,
\end{equation}
它被任何后两个分量为零的旋量所满足，例如
\begin{equation}
u_{0}=(\xi,0),
\end{equation}
其中 $\xi$ 是 $\mathbb{C}^{2}$ 的任意元素，而此处的 $\mathbf{0}$ 也应视为位于 $\mathbb{C}^{2}$ 中。

若我们用 $\mathrm{SO}(4)$ 矩阵 $\exp(-\theta J_{zt})$ 变换 $zt$ 平面内的坐标，则在新坐标中，原矢量 $m\be_{t}$ 将变成沿正 $z$ 方向、以速度 $\tan\theta$ 运动的波的能量-动量矢量。为求出与之相伴的旋量，我们把 $u_{0}$ 乘以 $\exp(-i\theta S^{zt})$：
\begin{align*}
k &= \exp(-\theta J_{zt})\,m\be_{t}=m(\cos\theta\,\be_{t}+\sin\theta\,\be_{z}),\\
u(k) &= \exp(-i\theta S^{zt})\,u_{0}
=\bigl(\xi\cos(\theta/2),\; i\sigma_{z}\xi\sin(\theta/2)\bigr).
\end{align*}
平面波运动得越快，Dirac 旋量的前一半越小、后一半越大。在无限速度 $\theta=\pi/2$ 处，两半大小相等。

现在，人们容易设想可以干脆把转动推到 $\theta=\pi/2$ 之外，以获得用坐标变换把一个黎曼电子变成一个黎曼正电子的描述，但 Dirac 方程还应付不了这一点。问题在于：传播矢量指向时间反方向的平面波会给出负能量，而对正电子我们理应与对电子一样测到正能量。要正确处理粒子与反粒子的关系，我们确实需要更精细的量子场论形式。

\subsubsection{Weyl 基下的平面波解}

换到 Weyl 基又会怎样？我们寻找如下形式的解：
\begin{equation}
\psi(x)=(\xi_{L},\xi_{R})\exp(-ik\cdot x).
\end{equation}
把它代入 Dirac 方程，得到关于 $\xi_{L}$、$\xi_{R}$ 的线性方程：
\begin{equation}
(k_{\mu}\gamma^{\mu}-mI_{4})(\xi_{L},\xi_{R})=0,
\end{equation}
利用用 $H$ 矩阵表示的 Weyl 基伽马矩阵形式，可把它写成
\begin{align}
k^{\mu}H_{\mu}\,\xi_{R}-m\xi_{L}&=0,\\
k^{\mu}H_{\mu}^{*}\,\xi_{L}-m\xi_{R}&=0.
\end{align}
正如我们从 Dirac 基中的情形所预期的那样，这两个方程并不独立。由于 $k^{\mu}H_{\mu}$ 是 $|\bk|$ 乘以一个酉矩阵（即逆等于厄米伴随的矩阵），我们有
\begin{equation}
(k^{\mu}H_{\mu})(k^{\mu}H_{\mu}^{*})=|\bk|^{2}I_{2}=m^{2}I_{2}.
\end{equation}
因此，若把第二个矩阵 Dirac 方程左边乘以 $k^{\mu}H_{\mu}$，再整体除以 $m$，我们便重新得到第一个方程。

对静止粒子，$\bk=m\be_{t}$，解具有形式
\begin{equation}
\psi(x)=(\xi,\xi)\exp(-imt).
\end{equation}
若在 $zt$ 平面内推动它（当然要用 Weyl 基的自旋矩阵），我们有
\begin{align*}
k &= \exp(-\theta J_{zt})\,m\be_{t}=m(\cos\theta\,\be_{t}+\sin\theta\,\be_{z}),\\
(\xi_{L},\xi_{R}) &= \exp(-i\theta S^{zt})(\xi,\xi)
=\bigl(\exp(-\tfrac12 i\theta\sigma_{z})\xi,\ \exp(\tfrac12 i\theta\sigma_{z})\xi\bigr).
\end{align*}
矩阵 $\exp(\pm\tfrac12 i\theta\sigma_{z})$ 是对角的，其元素等于相位 $\exp(\pm\tfrac12 i\theta)$，因此与 Dirac 基的情形不同，推动不改变旋量两半的范数。

在 Weyl 基中，我们可以轻松地为任意 $\bk$ 构造平面波解：
\begin{equation}
(\xi_{L}(k),\xi_{R}(k))
=\Bigl(\sqrt{\frac{k^{\mu}H_{\mu}}{m}}\,\xi,\
\sqrt{\frac{k^{\mu}H_{\mu}^{*}}{m}}\,\xi\Bigr).
\end{equation}
把它代入第一个矩阵 Dirac 方程 $k^{\mu}H_{\mu}\xi_{R}-m\xi_{L}=0$，并把两边乘以 $\sqrt m$，得
\begin{align*}
k^{\mu}H_{\mu}\sqrt{k^{\mu}H_{\mu}^{*}}\,\xi-m\sqrt{k^{\mu}H_{\mu}}\,\xi &= 0,\\
\sqrt{k^{\mu}H_{\mu}}\sqrt{(k^{\mu}H_{\mu})(k^{\mu}H_{\mu}^{*})}\,\xi-m\sqrt{k^{\mu}H_{\mu}}\,\xi &= 0,\\
m\sqrt{k^{\mu}H_{\mu}}\,\xi-m\sqrt{k^{\mu}H_{\mu}}\,\xi &= 0,
\end{align*}
从而证实它确实是解。如果你对不得不对矩阵开平方感到头疼，其实这类平方根可以很容易地写出：
\begin{equation}
\sqrt{k^{\mu}H_{\mu}}=\frac{k^{\mu}H_{\mu}+mH_{t}}{\sqrt{2(m+k^{t})}}.
\end{equation}

假设我们选取 $\mathbb{C}^{2}$ 的一组正交归一基 $\{\xi_{s}\}$。一个有用的例子是某个泡利自旋矩阵的归一化本征矢——按惯例通常选 $\sigma_{z}$，此时我们用标记 $s=\pm\tfrac12$。于是对任意四动量，我们可以写出黎曼 Dirac 方程的一对归一化平面波解：

\begin{keybox}[Weyl 基下黎曼 Dirac 方程的平面波解]
$\bk$ 为四动量，$|\bk|=m$；$\{\xi_{s}\}$ 是 $\mathbb{C}^{2}$ 的一组正交归一基；$H_{\mu}$ 是作为矩阵的单位四元数（定义见此）；求和 $\sum_j k^{j}H_{j}$ 遍历空间指标 $x$、$y$、$z$。
\begin{equation}
\psi_{\bk,s}(\bx)=u(\bk,s)\exp(-i\bk\cdot\bx) \tag{Plane wave}
\end{equation}
\begin{equation}
u(\bk,s)=\frac{\bigl([(m+k^{t})H_{t}+k^{j}H_{j}]\xi_{s},\ [(m+k^{t})H_{t}-k^{j}H_{j}]\xi_{s}\bigr)}{2\sqrt{m(m+k^{t})}} \tag{4-spinor part}
\end{equation}
\begin{equation}
u(\bk,s_{1})^{\dagger}u(\bk,s_{2})=\xi_{s_{1}}^{\dagger}\xi_{s_{2}}=\delta_{s_{1}s_{2}} \tag{Orthonormality}
\end{equation}
\end{keybox}

\subsection{四元数形式的 Dirac 方程（The Dirac Equation For Quaternions）}

如果我们沿用上一节关于 Weyl 基平面波的结果，并把由 $H_{\mu}$ 张成的 $2\times2$ 复矩阵空间 $\mathbb{H}$ 与二复维空间 $\mathbb{C}^{2}$ 都等同于我们此前描述过的四元数，那么就可以把联系旋量 $\xi_{L}$、$\xi_{R}$ 与能量-动量矢量 $\bk$ 的第一个方程，当作三个四元数之间的方程：
\begin{equation}
k^{\mu}H_{\mu}\,\xi_{R}-m\xi_{L}=0
\quad\longrightarrow\quad
Q_{V}\,q_{R}-m\,q_{L}=0,
\end{equation}
其中 $Q_{V}=k^{\mu}H_{\mu}$ 是像矢量那样变换的四元数，$q_{L}$ 是像左手旋量那样变换的四元数，$q_{R}$ 是像右手旋量那样变换的四元数。换言之，$\mathrm{SU}(2)\times\mathrm{SU}(2)$ 的元素 $(g,h)$ 对这三个四元数的作用为
\begin{align}
\rho_{(\tfrac12,\tfrac12)}(g,h)\,Q_{V} &= g\,Q_{V}\,h^{-1},\\
\rho_{(\tfrac12,0)}(g,h)\,q_{L} &= g\,q_{L},\\
\rho_{(0,\tfrac12)}(g,h)\,q_{R} &= h\,q_{R}.
\end{align}
乘积 $Q_{V}q_{R}$ 的变换方式与 $q_{L}$ 相同，从而使该方程不变。沿用此前的讨论，若在 $\mathbb{C}^{2}$ 中选取一个单位矢量 $\bv_{0}$，我们可以显式地把 $\xi_{L}=q_{L}\bv_{0}$、$\xi_{R}=q_{R}\bv_{0}$ 等同起来。此时，把我们这个四元数方程右乘以 $\bv_{0}$，就得到相应的旋量方程。

于是我们看到：平面波的 Dirac 方程可以解释为三个四元数之间极其简单的线性关系，这三个四元数分别在不同的 $\mathrm{SU}(2)\times\mathrm{SU}(2)$ 表示下变换，一个矢量代表粒子的能量-动量，一对左手、右手旋量代表波本身。

\begin{keybox}[四元数形式的黎曼 Dirac 方程]
$(q_{L},q_{R})$ 是作为一对四元数的 Dirac 旋量；$Q_{V}$ 是平面波解的能量-动量矢量：$\psi(\bx)=(q_{L},q_{R})\exp(-iQ_{V}\cdot\bx)$。
\begin{equation}
Q_{V}\,q_{R}-m\,q_{L}=0 \tag{Riemannian}
\end{equation}
\end{keybox}

为更好地理解这一四元数描述，我们来看看通常二分旋量形式下 $\half$ 自旋粒子的角动量如何用四元数语言翻译。首先要选定一个\textbf{复结构}（complex structure）：即 $\mathbb{H}$ 上的一个线性算符 $J$，它对应于 $\mathbb{C}^{2}$ 上乘以 $i$。$J$ 的平方必须等于乘以 $-1$，且 $J$ 应与 $\mathbb{H}$ 上任何对应于 $\mathbb{C}^{2}$ 复线性算符的线性算符相交换。要满足这些要求，我们可以令 $J$ 等于\textbf{右}乘某个平方为 $-1$ 的四元数，而把其它线性算符实现为\textbf{左}乘。具体地，我们取
\begin{equation}
J(q)=q\,H_{z}.
\end{equation}
给定这一选择，我们取 $\bv_{0}=(0,1)$，它是 $H_{z}$ 的、本征值为 $i$ 的本征矢。随后我们通过 $T:\mathbb{H}\to\mathbb{C}^{2}$ 把 $\mathbb{H}$ 与 $\mathbb{C}^{2}$ 等同起来，其中
\begin{align*}
T(q) &= q\,\bv_{0},\\
T(J(q)) &= qH_{z}(0,1)=i\,q(0,1)=i\,T(q).
\end{align*}
$J$ 的选择在 $\mathbb{H}$ 中挑出两个平面，我们可以把它们看成通常旋量的两个分量所在的复平面：$zt$ 平面与 $xy$ 平面。用一个四元数 $q$ 右乘任意复数 $a+bi$（$a$、$b$ 为实数），即实现为右乘 $a+bH_{z}$，而这一运算把这两个平面各自映入自身。

我们给 $\mathbb{H}$ 配备旋量内积 $\langle\cdot,\cdot\rangle_{S}$：
\begin{equation}
\langle p,q\rangle_{S}=\operatorname{Re}_{H}(p^{*}q)-\operatorname{Re}_{H}(p^{*}J(q))\,i
=\operatorname{Re}_{H}(p^{*}q)-\operatorname{Re}_{H}(p^{*}qH_{z})\,i,
\end{equation}
其中 $\operatorname{Re}_{H}$ 表示四元数的实部，以区别于复数的实部（后者记作 $\operatorname{Re}_{\mathbb{C}}$）。我们用 “$*$” 表示四元数的共轭，用 “conj” 表示普通复共轭。若把四元数表示为 $2\times2$ 复矩阵，则四元数共轭就是取这些矩阵的厄米伴随或复共轭转置，故 $q^{*}=\operatorname{conj}(q^{T})$。

我们的旋量内积实际上就是把 $\mathbb{C}^{2}$ 上的标准内积搬到 $\mathbb{H}$ 上，使得
\begin{equation}
\langle p,q\rangle_{S}=\langle T(p),T(q)\rangle.
\end{equation}
为看出这一点，先看实部：
\begin{align*}
\operatorname{Re}_{\mathbb{C}}\langle T(p),T(q)\rangle
&= \tfrac12\bigl(\langle T(p),T(q)\rangle+\langle T(q),T(p)\rangle\bigr)\\
&= \tfrac12\bigl(\langle p\bv_{0},q\bv_{0}\rangle+\langle q\bv_{0},p\bv_{0}\rangle\bigr)\\
&= \tfrac12\bigl(\langle \bv_{0},p^{*}q\bv_{0}\rangle+\langle \bv_{0},q^{*}p\bv_{0}\rangle\bigr)\\
&= \bigl\langle \bv_{0},\tfrac12(p^{*}q+q^{*}p)\bv_{0}\bigr\rangle\\
&= \operatorname{Re}_{H}(p^{*}q)\,\langle \bv_{0},\bv_{0}\rangle\\
&= \operatorname{Re}_{H}(p^{*}q).
\end{align*}
从倒数第三行到倒数第二行，是因为 $\tfrac12(p^{*}q+q^{*}p)=\operatorname{Re}_{H}(p^{*}q)$ 只是一个纯实数——用矩阵来说，就是同一个实数乘以 $2\times2$ 单位矩阵。至于虚部：
\begin{align*}
\operatorname{Im}_{\mathbb{C}}\langle T(p),T(q)\rangle
&= \operatorname{Re}_{\mathbb{C}}\bigl(-i\langle T(p),T(q)\rangle\bigr)\,i\\
&= -\operatorname{Re}_{\mathbb{C}}\langle T(p),iT(q)\rangle\,i\\
&= -\operatorname{Re}_{\mathbb{C}}\langle T(p),T(qH_{z})\rangle\,i\\
&= -\operatorname{Re}_{H}(p^{*}qH_{z})\,i.
\end{align*}

现在，二分旋量形式中通常的自旋本征态 $\xi_{k\pm}$ 满足方程（对 $k=x,y,z$）：
\begin{equation}
-\sigma_{k}\,\xi_{k\pm}=\pm\xi_{k\pm}.
\end{equation}
注意此处我们用了泡利矩阵的\textbf{相反数}，原因如前所述。在四元数形式中，我们需要把泡利矩阵 $\sigma_{k}$（它们不是 $\mathbb{H}$ 的元素）替换为基四元数 $H_{k}=-i\sigma_{k}$，并把 $\mathbb{C}^{2}$ 旋量 $\xi_{k\pm}$ 替换为满足 $\xi_{k\pm}=q_{k\pm}\bv_{0}$ 的四元数旋量 $q_{k\pm}$。这把我们引向四元数方程
\begin{equation}
H_{k}\,q_{k\pm}=\pm\,q_{k\pm}H_{z}.
\end{equation}
若把此式乘以 $-i$，把 $\bv_{0}$ 分别右乘两边，并利用 $\bv_{0}$ 是 $H_{z}$ 的、本征值为 $i$ 的本征矢这一事实，我们就重新得到此前关于 $\mathbb{C}^{2}$ 中旋量的方程。

这些四元数方程对每一种符号选择都有解，其解是 $\mathbb{H}$ 中的整张平面；当然，这些平面在与任意复数相乘时映入自身，对应 $\mathbb{C}^{2}$ 中的一复维本征空间。

对 $H_{x}$，正本征空间由 $H_{t}+H_{y}$ 与 $H_{x}+H_{z}$ 张成；负本征空间由 $H_{t}-H_{y}$ 与 $H_{x}-H_{z}$ 张成。

对 $H_{y}$，正本征空间由 $H_{t}-H_{x}$ 与 $H_{y}+H_{z}$ 张成；负本征空间由 $H_{t}+H_{x}$ 与 $H_{y}-H_{z}$ 张成。

对 $H_{z}$，正本征空间由 $H_{t}$ 与 $H_{z}$ 张成；负本征空间由 $H_{x}$ 与 $H_{y}$ 张成。

（所有这一切都可推广到 $\mathrm{SO}(4)$ 双覆盖的左右手旋量表示；在那里我们在 $\mathbb{R}^{4}$ 的各个坐标平面内测量自旋。各种自旋算符都是（半）$H$ 矩阵或其相反数，如本节末尾所列举，因此我们所需的一切都能从上文描述的本征空间中取得。例如，对自旋-$(\tfrac12,0)$ 左手旋量表示，$zt$ 坐标平面内的自旋由 $-\tfrac12 H_{z}$ 测量，故其正、负本征空间恰是上文 $H_{z}$ 的负、正本征空间。）

$\mathbb{C}^{2}$ 中的任意旋量都会成为沿空间中\textbf{某个}方向所测自旋的正本征态。在四元数形式中，对四元数 $q$ 而言，这个方向极易描述。按我们所选的复结构，它就是
\begin{equation}
s(q)=q\,H_{z}\,q^{-1}.
\end{equation}
显然它满足
\begin{equation}
s(q)\,q=q\,H_{z}.
\end{equation}
自旋轴 $s(q)$ 永远是纯空间矢量，因为 $H_{z}$ 一开始就是空间性的，而由 $\rho_{(\tfrac12,\tfrac12)}(g,g)Q_{V}=gQ_{V}g^{-1}$ 产生的所有 $\mathrm{SO}(4)$ 转动都保持时间轴不变，因而实际上都是 $\mathrm{SO}(3)$ 中的转动。我们还能极容易地找到自旋恰与 $q$ 相反的自旋旋量四元数，记作 $\operatorname{rev}_{S}(q)$：
\begin{align*}
\operatorname{rev}_{S}(q) &= q\,H_{y},\\
s\bigl(\operatorname{rev}_{S}(q)\bigr)
&= qH_{y}H_{z}H_{y}^{-1}q^{-1}
= qH_{y}H_{z}(-H_{y})q^{-1}\\
&= -qH_{x}H_{y}q^{-1}
= -s(q).
\end{align*}
按我们“用右乘 $H_{z}$ 实现乘以 $i$”的复结构，自旋反转函数 $\operatorname{rev}_{S}$ 是共轭线性的：
\begin{align*}
\operatorname{rev}_{S}(iq) &= \operatorname{rev}_{S}(qH_{z})=qH_{z}H_{y}\\
&= -qH_{y}H_{z}=-i\,\operatorname{rev}_{S}(q).
\end{align*}
容易看出 $\operatorname{rev}_{S}$ 与 $\mathrm{SU}(2)\times\mathrm{SU}(2)$ 的作用相交换，因为该作用从左边乘四元数，而 $\operatorname{rev}_{S}$ 从右边乘。

$\mathrm{SO}(3)$ 中任何保持矢量 $s(q)$ 不变的转动，都可以由 $\rho_{(\tfrac12,\tfrac12)}(g,g)$ 产生，其中 $g=\cos(\theta)H_{t}+\sin(\theta)s(q)$。这样的转动会以左乘 $g$ 的方式作用于旋量 $q$（无论它是左手还是右手），给出
\begin{align*}
g\,q &= \bigl(\cos(\theta)H_{t}+\sin(\theta)s(q)\bigr)q\\
&= \bigl(\cos(\theta)H_{t}+\sin(\theta)\,qH_{z}q^{-1}\bigr)q\\
&= \cos(\theta)\,q+\sin(\theta)\,qH_{z}\\
&= \bigl(\cos(\theta)+i\sin(\theta)\bigr)q
= \exp(i\theta)\,q .
\end{align*}
换言之，这个转动仅仅改变旋量 $q$ 的整体相位。

在粒子可以以某个任意速度运动的相对论情境中，我们不能再把转动看成具有一根轴。取而代之，我们用一个单位四元数对 $(s_{L},s_{R})$ 来标识一个转动，它通过 $v\to s_{L}\,v\,s_{R}^{-1}$ 作用。若我们做一次对应转动 $(g,h)$ 的任意坐标变换，则原来的转动被变换为 $(g\,s_{L}\,g^{-1},\;h\,s_{R}\,h^{-1})$。这来自三个操作的复合：(1) 反演坐标变换，回到原转动所在的坐标；(2) 施行原转动；(3) 施加坐标变换，回到新坐标。

若某粒子的状态是 Dirac 旋量 $(q_{L},q_{R})$ 且 $q_{L}\neq0$、$q_{R}\neq0$，我们可以定义一对用于转动的单位四元数：
\begin{align}
s_{L}(q_{L}) &= q_{L}H_{z}q_{L}^{-1},\\
s_{R}(q_{R}) &= q_{R}H_{z}q_{R}^{-1}.
\end{align}
注意，若施行一次对应 $(g,h)$ 的坐标变换，并使用旋量变换律 $q_{L}\to gq_{L}$、$q_{R}\to hq_{R}$，则转动 $\bigl(s_{L}(q_{L}),s_{R}(q_{R})\bigr)$ 的变换方式恰如一个转动在坐标变换下应有的方式。由这些定义还可推出 $s_{L}(q_{L})^{2}=s_{R}(q_{R})^{2}=-H_{t}$，故把该转动施行两次，$v\to s_{L}(q_{L})^{2}\,v\,s_{R}(q_{R})^{-2}$，恰好给出恒等变换。

显然，$s_{L}(q_{L})$ 与 $s_{R}(q_{R})$ 将满足
\begin{align*}
s_{L}(q_{L})\,q_{L} &= q_{L}H_{z},\\
s_{R}(q_{R})\,q_{R} &= q_{R}H_{z}.
\end{align*}
其物理含义是：由 $(s_{L}(q_{L}),s_{R}(q_{R}))$ 这一对产生的矢量转动，将伴随 Dirac 旋量的一次变换 $(q_{L},q_{R})\to(s_{L}(q_{L})q_{L},\,s_{R}(q_{R})q_{R})=(q_{L}H_{z},q_{R}H_{z})$，而这仅仅相当于按我们设在 $\mathbb{H}$ 上的复结构把整个 Dirac 旋量乘以 $i$。

在粒子静止的参照系中，能量-动量矢量 $Q_{V}=mH_{t}$，于是方程
\begin{equation}
Q_{V}\,q_{R}-m\,q_{L}=0
\end{equation}
给出 $q_{R}=q_{L}$，进而有 $s_{L}(q_{L})=s_{R}(q_{R})$，相应的转动就是三维的。因为我们知道该转动的平方是恒等变换，所以它必定是某个特定平面内的半周转动。这一事实在我们换到粒子运动的参照系时不会改变，因此一般而言，我们可以把 $(s_{L}(q_{L}),s_{R}(q_{R}))$ 视为标识出粒子自旋所关联的\textbf{转动平面}，用以取代三维情境中所用的自旋轴 $s(q)$。

更进一步，我们可以辨认出两个与该平面正交（且彼此正交）、并为该转动所保持不变的矢量：能量-动量矢量 $Q_{V}=m\,q_{L}q_{R}^{-1}$，以及矢量
\begin{equation}
s_{LR}(q_{L},q_{R})=q_{L}H_{z}q_{R}^{-1}.
\end{equation}
当粒子静止且 $q_{R}=q_{L}$ 时，后者与自旋轴 $s(q)$ 重合。

若把自旋反转函数 $\operatorname{rev}_{S}$ 应用于 Dirac 旋量中的两个四元数，则与自旋几何相关的全部四个四元数都被反转：$s_{L}(q_{L})$、$s_{R}(q_{R})$、$Q_{V}$ 以及 $s_{LR}(q_{L},q_{R})$。

\subsection{拉格朗日量（The Lagrangian）}

构造一个能给出黎曼 Dirac 方程的拉格朗日量是容易的。

\begin{keybox}[Dirac 方程的拉格朗日量]
\begin{align}
L_{RD} &= \psi^{\dagger}(i\gamma^{\mu}\partial_{\mu}\psi-m\psi)
=i\psi^{\dagger}\gamma^{\mu}\partial_{\mu}\psi-m\psi^{\dagger}\psi \tag{Riemannian}\\
L_{LD} &= \psi^{\dagger}\gamma^{t}(i\gamma^{\mu}\partial_{\mu}\psi-m\psi)
=i\psi^{\dagger}\gamma^{t}\gamma^{\mu}\partial_{\mu}\psi-m\psi^{\dagger}\gamma^{t}\psi \tag{Lorentzian}
\end{align}
\end{keybox}

这里 $\psi^{\dagger}$ 是 $4\times1$ 列矢量 $\psi$ 的厄米伴随（即共轭转置），因此它是一个 $1\times4$ 行矢量，我们按矩阵乘法把它乘进其右侧的对象中。

在洛伦兹情形中，代表 $\psi$ 上洛伦兹变换的矩阵一般不是酉的，其后果是 $\psi^{\dagger}\psi$ 不是一个标量：它在洛伦兹变换下不不变。因此有必要把 $\psi^{\dagger}$ 换成 $\psi^{\dagger}\gamma^{t}$；这给出一个洛伦兹不变的内积 $\psi^{\dagger}\gamma^{t}\psi$，并使我们能构造出一个洛伦兹不变的拉格朗日量。（在量子场论文献中，$\psi^{\dagger}\gamma^{t}$ 被简写为顶上加一横的 $\psi$。）但这些小麻烦只属于洛伦兹情形。在黎曼情形中，所有自旋矩阵都是厄米的，$\mathrm{SO}(4)$ 的表示是酉的，且 $\psi^{\dagger}\psi$ 是 $\mathrm{SO}(4)$ 不变的：
\begin{align*}
\psi &\to U\psi,\\
\psi^{\dagger} &\to (U\psi)^{\dagger}=\psi^{\dagger}U^{\dagger}=\psi^{\dagger}U^{-1},\\
\psi^{\dagger}\psi &\to \psi^{\dagger}U^{-1}U\psi=\psi^{\dagger}\psi.
\end{align*}
$L_{RD}$ 的第一项也是 $\mathrm{SO}(4)$ 不变的，这借助伽马矩阵那节的方程 \textbf{(2)} 可以很容易地证明。

我们借助欧拉-拉格朗日方程（Euler-Lagrange equations）从拉格朗日量得到 Dirac 方程。我们把 $\psi^{\dagger}$ 视为独立场，写出
\begin{align*}
\partial_{\mu}\bigl[\partial_{\partial_{\mu}\psi^{\dagger}}L_{RD}\bigr] &= \partial_{\psi^{\dagger}}L_{RD},\\
0 &= i\gamma^{\mu}\partial_{\mu}\psi-m\psi,
\end{align*}
其中方括号里的表达式为零，因为拉格朗日量不依赖 $\psi^{\dagger}$ 的任何导数。改用 $\psi$ 本身作为场做对应的计算会得到相同结果，只是更麻烦：
\begin{align*}
\partial_{\mu}\bigl[\partial_{\partial_{\mu}\psi}L_{RD}\bigr] &= \partial_{\psi}L_{RD},\\
\partial_{\mu}\bigl[i\psi^{\dagger}\gamma^{\mu}\bigr] &= -m\psi^{\dagger},\\
i\,\partial_{\mu}\psi^{\dagger}\,\gamma^{\mu}+m\psi^{\dagger} &= 0,\\
-\bigl(i\,\partial_{\mu}\psi^{\dagger}\,\gamma^{\mu}+m\psi^{\dagger}\bigr)^{\dagger} &= 0,\\
i\gamma^{\mu}\partial_{\mu}\psi-m\psi &= 0 .
\end{align*}

\subsection{守恒流（Conserved Current）}

现在我们有了 Dirac 方程的拉格朗日量，便可以用诺特定理（Noether's Theorem）来找出一个守恒流。假设我们把 Dirac 旋量 $\psi$ 乘以任意一个模为 $1$ 的复数，即一个“相位” $\exp(i\alpha)$（$\alpha$ 为实值）。由于 $\psi^{\dagger}$ 最终带上相反的相位，拉格朗日量不变。

诺特定理为拉格朗日量的每一种这样的对称性配一个守恒流：
\begin{align*}
j^{\mu} &= \bigl(\partial_{\partial_{\mu}\psi}L_{RD}\bigr)\,\bigl(\partial_{\alpha}\exp(i\alpha)\psi\bigr)\big|_{\alpha=0}\\
&= -\psi^{\dagger}\gamma^{\mu}\psi .
\end{align*}
我们可以把它乘以任意常数，该流仍是守恒的。具体地，我们定义

\begin{keybox}[Dirac 方程的流]
$e$ 为粒子所带电荷。
\begin{align}
j^{\mu} &= e\,\psi^{\dagger}\gamma^{\mu}\psi \tag{Riemannian}\\
j^{\mu} &= e\,\psi^{\dagger}\gamma^{t}\gamma^{\mu}\psi \tag{Lorentzian}
\end{align}
\end{keybox}

称它为“守恒流”，意思是其散度为零，因而它不会凭空产生或消失：
\begin{align*}
\partial_{\mu}j^{\mu}
&= e\,\partial_{\mu}(\psi^{\dagger}\gamma^{\mu}\psi)\\
&= e\bigl((\partial_{\mu}\psi^{\dagger})\gamma^{\mu}\psi+\psi^{\dagger}\gamma^{\mu}(\partial_{\mu}\psi)\bigr)\\
&= e\bigl((im\psi^{\dagger})\psi+\psi^{\dagger}(-im\psi)\bigr)\\
&= 0 .
\end{align*}

考察 Dirac 基下的静止平面波解：
\begin{equation}
\psi(x)=(\xi,0)\exp(-im\be_{t}\cdot x),
\end{equation}
其中 $\xi$ 是任意二分量旋量。对这个解，我们简单地有
\begin{equation}
\bj=e(\xi^{\dagger}\xi)\,\be_{t},
\end{equation}
因为除 $\gamma^{t}$ 之外的伽马矩阵的结构使得所有其它分量都为零。当然，对无限平面波我们面临通常的归一化问题，但眼下我们忽略它，只关注 $\bj$ 的方向：在这个特定情形中它与粒子的传播矢量同向——并且若换到一个波在其中运动的参照系，它仍将同向。因此 $\bj$ 的行为恰如我们期待的电流密度。

\subsection{与电磁学的耦合（Coupling With Electromagnetism）}

我们可以把 Dirac 方程的拉格朗日量与先前为黎曼 Proca 方程（Riemannian Proca equation，即支配电磁学的方程）导出的拉格朗日量结合起来。然后，利用用 Dirac 场表示的电流密度表达式，我们便得到一个完全以基本场表示的总拉格朗日量。[注意：对洛伦兹版本，为便于与量子场论教科书比较，本页采用 $(+---)$ 号差。这意味着其中一些洛伦兹方程与电磁学注记里的版本并不完全一致。]

\begin{keybox}[量子电动力学的拉格朗日量]
$m_{ph}$ 为光子质量，$m$ 为 Dirac 粒子的质量，$e$ 为其电荷；电磁场 $F_{\mu\nu}=\partial_{\mu}A_{\nu}-\partial_{\nu}A_{\mu}$。
\begin{align}
L_{RQED} &= \psi^{\dagger}(i\gamma^{\mu}\partial_{\mu}\psi-m\psi)+\tfrac14 F_{\mu\nu}F^{\mu\nu}-\tfrac12 m_{ph}^{2}A_{\mu}A^{\mu}-A_{\mu}j^{\mu} \nonumber\\
&= \psi^{\dagger}(i\gamma^{\mu}\partial_{\mu}\psi-m\psi)+\tfrac14 F_{\mu\nu}F^{\mu\nu}-\tfrac12 m_{ph}^{2}A_{\mu}A^{\mu}-eA_{\mu}\psi^{\dagger}\gamma^{\mu}\psi \tag{Riemannian}\\[4pt]
L_{LQED} &= \psi^{\dagger}\gamma^{t}(i\gamma^{\mu}\partial_{\mu}\psi-m\psi)-\tfrac14 F_{\mu\nu}F^{\mu\nu}-A_{\mu}j^{\mu} \nonumber\\
&= \psi^{\dagger}\gamma^{t}(i\gamma^{\mu}\partial_{\mu}\psi-m\psi)-\tfrac14 F_{\mu\nu}F^{\mu\nu}-eA_{\mu}\psi^{\dagger}\gamma^{t}\gamma^{\mu}\psi \tag{Lorentzian}
\end{align}
\end{keybox}

这些拉格朗日量的欧拉-拉格朗日方程同时给出耦合到 Dirac 源的电磁场方程，以及存在电磁场时的 Dirac 方程。

\begin{keybox}[量子电动力学的场方程]
$m_{ph}$ 为光子质量，$m$ 为 Dirac 粒子的质量，$e$ 为其电荷；电磁场 $F_{\mu\nu}=\partial_{\mu}A_{\nu}-\partial_{\nu}A_{\mu}$。
\begin{align}
&\partial_{\mu}\partial^{\mu}A^{\nu}+m_{ph}^{2}A^{\nu}+e\,\psi^{\dagger}\gamma^{\nu}\psi=0, \nonumber\\
&i\gamma^{\mu}\partial_{\mu}\psi-m\psi-eA_{\mu}\gamma^{\mu}\psi=0 \tag{Riemannian}\\[4pt]
&\partial_{\mu}F^{\mu\nu}-e\,\psi^{\dagger}\gamma^{t}\gamma^{\nu}\psi=0, \nonumber\\
&i\gamma^{\mu}\partial_{\mu}\psi-m\psi-eA_{\mu}\gamma^{\mu}\psi=0 \tag{Lorentzian}
\end{align}
\end{keybox}

\subsection{自旋与角动量守恒（Spin and Conservation of Angular Momentum）}

电磁场中的黎曼 Dirac 方程为
\begin{equation}
i\gamma^{\mu}\partial_{\mu}\psi-m\psi-eA_{\mu}\gamma^{\mu}\psi=0.
\end{equation}
让我们把它改写为显式的哈密顿形式：在左边乘以 $\gamma^{t}$（其平方，如同所有黎曼伽马矩阵的平方一样，是单位矩阵），并把时间导数分离出来：
\begin{align}
i\partial_{t}\psi &= H\psi,\\
H &= -i\gamma^{t}\gamma^{j}\partial_{j}+m\gamma^{t}+eA_{\mu}\gamma^{t}\gamma^{\mu}.
\end{align}
这里对重复指标 $j$ 的求和只遍历空间坐标，而对 $\mu$ 的求和遍历全部四维。

我们把在左边乘以 $\gamma^{t}$ 后得到的伽马矩阵称为 \textbf{alpha 矩阵}（alpha matrices）；与伽马矩阵一样，它们在 Weyl 基下形式非常简单。

\begin{keybox}[Weyl 基下的黎曼 alpha 矩阵]
\begin{align}
\alpha^{\mu} &= \gamma^{t}\gamma^{\mu}
=\begin{pmatrix}H_{\mu}^{*}&0\\ 0&H_{\mu}\end{pmatrix} \tag{Riemannian}
\end{align}
\end{keybox}

现在，若一个可观测量的算符与哈密顿量相交换，它就描述一个守恒量。若我们有一个处于径向对称静电场中的黎曼电子，我们预期其角动量守恒，因此让我们计算该情形下的哈密顿量与轨道角动量算符 $\mathbf{L}=\mathbf{r}\times\mathbf{p}$ 之间的交换子。在势为 $V(r)$ 的静电场中，$A_{t}=-V(r)$，其余所有 $A_{j}=0$，从而给出哈密顿量
\begin{align*}
H &= -i\alpha^{j}\partial_{j}+m\gamma^{t}-eV(r)\\
&= -\alpha^{j}p_{j}+m\gamma^{t}-eV(r),
\end{align*}
其中 $p_{j}=i\partial_{j}$ 是线动量的第 $j$ 个分量；别忘了，对黎曼相对论动量，我们必须使用不同的算符约定。轨道角动量算符的分量为
\begin{equation}
L_{i}=\varepsilon_{ijk}x^{j}p_{k}.
\end{equation}
为求它与哈密顿量的交换子，我们先计算：
\begin{align*}
[x^{j}p_{k},p_{s}]
&= -[x^{j}\partial_{k},\partial_{s}]\\
&= -\bigl(x^{j}\partial_{k}\partial_{s}-\partial_{s}(x^{j}\partial_{k})\bigr)\\
&= -\bigl(x^{j}\partial_{k}\partial_{s}-x^{j}\partial_{s}\partial_{k}-\delta_{s}^{j}\partial_{k}\bigr)\\
&= -i\,\delta_{s}^{j}\,p_{k}.
\end{align*}
这给出
\begin{equation}
[L_{i},p_{s}]=-i\,\varepsilon_{isk}\,p_{k}.
\end{equation}
$L_{i}$ 显然与哈密顿量的第二项相交换，因为那一项只涉及乘以一个固定矩阵。第三项呢？对任何只依赖径向距离平方 $r^{2}=x^{2}+y^{2}+z^{2}=x^{l}x^{l}$ 的函数 $f$，以及任意函数 $g$，我们有
\begin{align*}
[x^{j}p_{k},f(x^{l}x^{l})]\,g
&= i\,[x^{j}\partial_{k},f(x^{l}x^{l})]\,g\\
&= i\bigl(x^{j}\partial_{k}(f(x^{l}x^{l})g)-f(x^{l}x^{l})x^{j}\partial_{k}g\bigr)\\
&= i\bigl(2g\,x^{j}x^{k}\,f'(x^{l}x^{l})+f(x^{l}x^{l})x^{j}\partial_{k}g-f(x^{l}x^{l})x^{j}\partial_{k}g\bigr)\\
&= 2i\,x^{j}x^{k}\,f'(x^{l}x^{l})\,g .
\end{align*}
把结果写成不含函数 $g$ 的线性算符：
\begin{equation}
[x^{j}p_{k},f(x^{l}x^{l})]=2i\,x^{j}x^{k}\,f'(x^{l}x^{l}).
\end{equation}
由于该表达式在指标 $j$ 与 $k$ 上对称，而 $L_{i}=\varepsilon_{ijk}x^{j}p_{k}$ 中的 $\varepsilon$ 是反对称的，故 $[L_{i},f(x^{l}x^{l})]$ 的总和为零；又因为哈密顿量的第三项可以写成形如 $f$ 的函数，所以该项与 $L_{i}$ 相交换。

这样一来，只剩下哈密顿量的第一项。于是我们有
\begin{align*}
[L_{i},H] &= [L_{i},-\alpha^{j}p_{j}]\\
&= -\alpha^{j}[L_{i},p_{j}]\\
&= i\,\varepsilon_{ijk}\,\alpha^{j}p_{k}.
\end{align*}
由于此式非零，轨道角动量将不守恒。然而，如果我们以正确的方式把自旋角动量归给电子本身，\textbf{总}角动量仍然可以守恒。

我们此前描述过代表转动生成元的厄米自旋矩阵 $S^{\mu\nu}$。若定义空间自旋矩阵 $\mathbf{S}=(S^{yz},S^{zx},S^{xy})$，则在 Weyl 基中有 $S_{i}=-\tfrac12 I_{2}\otimes\sigma_{i}=-\tfrac12 i\,I_{2}\otimes H_{i}$。

\begin{keybox}[Weyl 基下的空间自旋矩阵]
\begin{align}
S_{i} &= \begin{pmatrix}-\tfrac12 iH_{i}&0\\ 0&-\tfrac12 iH_{i}\end{pmatrix} \tag{Riemannian}
\end{align}
\end{keybox}

带空间指标的 $H$ 矩阵之间的交换子为
\begin{equation}
[H_{j},H_{k}]=2\,\varepsilon_{jkl}\,H_{l}.
\end{equation}
现在，在带\textbf{空间指标}的 alpha 矩阵中，伴随矩阵 $H_{j}^{*}=-H_{j}$。由此我们可以轻松算出空间自旋矩阵与空间 alpha 矩阵的交换子：
\begin{equation}
[S_{i},\alpha^{j}]=-i\,\varepsilon_{ijl}\,\alpha^{l}.
\end{equation}
矩阵 $S_{i}$ 将与哈密顿量除第一项之外的所有项相交换，因此
\begin{align*}
[S_{i},H] &= -[S_{i},\alpha^{j}]\,p_{j}\\
&= i\,\varepsilon_{ijl}\,\alpha^{l}p_{j}\\
&= i\,\varepsilon_{ikj}\,\alpha^{j}p_{k}\quad[\text{重标指标}]\\
&= -i\,\varepsilon_{ijk}\,\alpha^{j}p_{k}.
\end{align*}
这正是轨道角动量分量 $L_{i}$ 的交换子的相反数，于是对总角动量我们有
\begin{equation}
[L+S,H]=0.
\end{equation}
我们最终为总角动量得到了一个守恒量，这一事实确认了我们对自旋矩阵所作的选法是正确的！

\section{进阶：电子-光子自旋耦合（Electron-Photon Spin Coupling with Quaternions）}

我们此前讨论过，\emph{四元数对}（pairs of quaternions）$(q_L,q_R)$ 如何可被解释为 Dirac 旋量，以及 $\tfrac12$ 自旋粒子的角动量如何在该形式中描述。我们把每一份四元数视为一个二复维的希尔伯特空间，用以取代通常的旋量空间 $\mathbb{C}^{2}$，办法是给它配备一个\emph{复结构}（complex structure）——而我们具体的选择是用四元数 $H_z$ 的\emph{右乘}来实现乘以 $\ii$（即 $-1$ 的平方根）。在这些约定下，Weyl 基中自由平面波的 Dirac 方程变得异常简单：
\begin{equation*}
q_V\,q_R - m\,q_L = 0 .
\end{equation*}
这里 $(q_L,q_R)$ 是 Dirac 旋量，其第一个分量按左手旋量变换、第二个按右手旋量变换，而 $q_V$ 是给出平面波 4-动量的矢量四元数。四元数值的平面波为
\begin{equation*}
\psi(x)=(q_L,q_R)\exp(-\ii\,q_V\cdot x).
\end{equation*}

我们也能发展出一套用四元数描述自旋为 1 的粒子（例如黎曼光子）的角动量的理论。$\mathrm{SU}(2)\times\mathrm{SU}(2)$（即 $\mathrm{SO}(4)$ 的双覆盖）的矢量表示其实就是一个作用在四元数上的表示，其中群 $\mathrm{SU}(2)$ 对应单位四元数，而 $\mathrm{SU}(2)\times\mathrm{SU}(2)$ 中的一对单位四元数 $(g,h)$ 通过
\begin{equation*}
(g,h):\ q\mapsto g\,q\,h^{-1}=g\,q\,h^*
\end{equation*}
作用在四元数自身之上。若把四元数 $\mathbb{H}$ 与四维欧氏空间 $\mathbb{R}^{4}$ 等同，则该表示对应于用 $\mathrm{SO}(4)$ 的元素（行列式为 $1$ 的实正交 $4\times4$ 矩阵）去乘 $\mathbb{R}^{4}$ 中的矢量。这在做经典物理、只想谈论旋转实四矢量时是没问题的，但自旋为 1 的粒子的角动量是复希尔伯特空间中的矢量。我们如何把两者联系起来？

答案是给四元数\emph{复化}（complexify）。正如把一般复数写成 $a+b\,\ii$（$a$、$b$ 为实数），我们把一般的\emph{复化四元数}（complexified quaternion）写成 $q+s\,\ii$，其中 $q$、$s$ 为四元数。这里的 $\ii$ \emph{不}应被当作四元数内部任何一个 $-1$ 的平方根，也不等同于我们在\emph{旋量}四元数上所用的「右乘 $H_z$」。（只有当我们把一个复化矢量四元数乘以一个旋量以得到另一个旋量时，才把因子 $\ii$ 换算成右乘 $H_z$。）在复化四元数本身内部，$\ii$ 只是一个与一切相交换、并满足 $\ii^{2}=-1$ 的符号。于是若把两个复化四元数相乘，便有
\begin{equation*}
(q+s\,\ii)(p+r\,\ii)=(qp-sr)+(sp+qr)\,\ii .
\end{equation*}

处理复化四元数时，我们采用如下约定：「$\operatorname{Re}_H$」表示四元数的实部，以区别于复数的实部（后者记作「$\operatorname{Re}_{\mathbb{C}}$」）；用「$*$」表示四元数的共轭，用「$\operatorname{conj}$」表示普通复共轭。例如在下面这个式子中，我们把所有与 $H_t$ 正交的四元数分量取相反数，却\emph{不}改变 $\ii$ 的系数的符号：
\begin{equation*}
\bigl((H_t+H_x)+(H_t+H_y)\,\ii\bigr)^*=(H_t-H_x)+(H_t-H_y)\,\ii .
\end{equation*}

四元数上定义有一个实内积（即点积）：
\begin{equation*}
q\cdot s=\operatorname{Re}_H(q^*s)=\tfrac12(q^*s+s^*q).
\end{equation*}
它在 $\mathrm{SU}(2)\times\mathrm{SU}(2)$ 的作用下不变；若用单位四元数对 $(g,h)$ 作用，得
\begin{align*}
q\cdot s &\to (g\,q\,h^{-1})\cdot(g\,s\,h^{-1})\\
&= \tfrac12\bigl((g\,q\,h^{-1})^*(g\,s\,h^{-1})+(g\,s\,h^{-1})^*(g\,q\,h^{-1})\bigr)\\
&= \tfrac12\bigl(h\,q^*\,g^{-1}g\,s\,h^{-1}+h\,s^*\,g^{-1}g\,q\,h^{-1}\bigr)\\
&= \tfrac12\bigl(h\,q^*\,s\,h^{-1}+h\,s^*\,q\,h^{-1}\bigr)\\
&= \tfrac12\,h\,(q^*s+s^*q)\,h^{-1}\\
&= h\,\operatorname{Re}_H(q^*s)\,h^{-1}\\
&= \operatorname{Re}_H(q^*s)\\
&= q\cdot s .
\end{align*}

我们把这个点积延拓到复化四元数上，只须要求它对每个向量复线性：
\begin{equation*}
(q+s\,\ii)\cdot(p+r\,\ii)=(q\cdot p-s\cdot r)+(s\cdot p+q\cdot r)\,\ii .
\end{equation*}
我们还需要复化四元数的希尔伯特空间上的一个内积，它对第二个变量线性、而对第一个变量\emph{共轭线性}（在复结构的意义上！）：
\begin{align*}
\langle q+s\,\ii,\,p+r\,\ii\rangle_V
&= \operatorname{conj}(q+s\,\ii)\cdot(p+r\,\ii)\\
&= (q-s\,\ii)\cdot(p+r\,\ii)\\
&= (q\cdot p+s\cdot r)+(q\cdot r-s\cdot p)\,\ii .
\end{align*}
我们给这个内积加上下标「$V$」，意为矢量（vector），以区别于旋量四元数上的内积：
\begin{equation*}
\langle p,q\rangle_S=\operatorname{Re}_H(p^*q)-\operatorname{Re}_H(p^*qH_z)\,\ii .
\end{equation*}
既然已证明点积在 $\mathrm{SU}(2)\times\mathrm{SU}(2)$ 作用下不变，我们的复内积 $\langle\cdot,\cdot\rangle_V$ 也将不变——换言之，$\mathrm{SU}(2)\times\mathrm{SU}(2)$ 的这个表示是酉的，正如它必须如此。

到目前为止，我们大体上采取了一种纯粹主义的做法，把四元数当作基本的数学对象处理，这给出许多计算中有用的公式。但有时值得回想：我们可以把普通四元数视为 $2\times2$ 复矩阵的一个子空间，它由四个矩阵 $H_\mu$ 的实线性组合张成——在这个语境中，我们可以把\emph{复化}四元数视为任意的 $2\times2$ 复矩阵，它们可由 $H_\mu$ 的复线性组合得到。

把复化矢量四元数看作 $2\times2$ 复值矩阵、把旋量四元数看作位于由 $H_\mu$ 的实线性组合所张成的较小子空间中，则矢量和旋量内积可写成
\begin{align*}
\langle p,q\rangle_V &= \tfrac12\operatorname{tr}\bigl(\operatorname{conj}(p^{T})\,q\bigr),\\
\langle p,q\rangle_S &= \bigl\langle p,\,q-\ii\,qH_z\bigr\rangle_V
= \tfrac12\operatorname{tr}\bigl(\operatorname{conj}(p^{T})(q-\ii\,qH_z)\bigr).
\end{align*}

当我们把复化四元数用作黎曼光子自旋的希尔伯特空间时，我们把不含 $\ii$ 因子的「纯」四元数取为线偏振（linear polarisation）态，对应经典黎曼矢量平面波中四维势矢量 $\bA$ 的方向。在光子静止的参照系中，四元数 $H_x$、$H_y$、$H_z$ 是平面波三种线偏振的态矢量，它们全都与其传播矢量正交（在此情形传播矢量为 $H_t$，因为我们处于静止系）。

三个矢量 $H_x$、$H_y$、$H_z$ 也可视为一个自旋为 1 的粒子的态，其自旋在 $x$、$y$、$z$ 方向上的分量分别为\emph{零}。这与自旋为 1 的粒子（有质量，不同于我们宇宙中的光子）的非相对论量子力学精确吻合：那里选取 $\mathbb{C}^{3}$ 的一组基，由在三个正交方向上自旋分量均为零的态组成，就可使 $\mathrm{SU}(2)$ 在 $\mathbb{C}^{3}$ 上的自旋-1 表示与 $\mathrm{SO}(3)$ 在 $\mathbb{R}^{3}$ 上的矢量表示等同。在这组基下，作用在 $\mathbb{C}^{3}$ 上、代表转动的自旋-1 酉矩阵变为实值——事实上变得与转动的 $\mathrm{SO}(3)$ 矩阵完全相同。

粒子沿某根轴自旋分量为 $+1$ 或 $-1$ 的态又如何？我们可以模仿经典波中圆偏振（circular polarisation）与线偏振之间的关系来找出这些态。在经典平面波中，比方说绕 $z$ 轴的圆偏振，出现在沿 $x$ 轴线偏振的波与另一个沿 $y$ 轴线偏振、且与前者相差 $90^\circ$ 相位的波相叠加之时。对复指数而言，$90^\circ$ 的相移可由乘以 $\ii$ 产生。于是我们预期沿 $z$ 轴自旋为一个单位的归一化光子态为 $(H_x\pm\ii H_y)/\sqrt2$。不过，由于那个促使我们为费米子采用泡利自旋矩阵\emph{相反数}的扭转，我们在这里必须仔细选择符号：\emph{和}是负自旋本征态，而\emph{差}是正自旋本征态。

\begin{center}
\begin{tabular}{ll}
\toprule
复化矢量四元数 & 角动量分量 \\
\midrule
$H_x$ & $m_x=0$ \\
$H_y$ & $m_y=0$ \\
$H_z$ & $m_z=0$ \\
$(H_y\pm\ii H_z)/\sqrt2$ & $m_x=-1,\,1$ \\
$(H_z\pm\ii H_x)/\sqrt2$ & $m_y=-1,\,1$ \\
$(H_x\pm\ii H_y)/\sqrt2$ & $m_z=-1,\,1$ \\
$H_t$ & $m_x=m_y=m_z=0$ \\
\bottomrule
\end{tabular}
\end{center}

从这些态出发，我们可以通过投影到各自旋本征态并乘以自旋，来构造沿每根轴的自旋算符。例如：
\begin{align*}
J_z(v)
&= -\frac{H_x+\ii H_y}{\sqrt2}\Bigl\langle\frac{H_x+\ii H_y}{\sqrt2},\,v\Bigr\rangle_V
+\frac{H_x-\ii H_y}{\sqrt2}\Bigl\langle\frac{H_x-\ii H_y}{\sqrt2},\,v\Bigr\rangle_V\\
&= \ii\bigl(H_x\langle H_y,v\rangle_V-H_y\langle H_x,v\rangle_V\bigr)\\
&= \ii\bigl(H_x(H_y\cdot v)-H_y(H_x\cdot v)\bigr).
\end{align*}
从倒数第二行到最后一行的过渡中，我们利用了这样一个事实：当第一个自变量不含虚部时，复内积与点积变得完全相同。

我们将上式写作
\begin{equation*}
J_z=\ii\,(H_x\otimes H_y-H_y\otimes H_x),
\end{equation*}
并约定它通过每个张量积中右端项的点积作用于任何复化四元数。类似地有
\begin{align*}
J_x &= \ii\,(H_y\otimes H_z-H_z\otimes H_y),\\
J_y &= \ii\,(H_z\otimes H_x-H_x\otimes H_z).
\end{align*}
除去一个因子 $-\ii$ 外，若把这些算符视为作用在 $\mathbb{R}^{3}$ 上的矩阵，它们就正是生成绕各自轴转动的 $\mathrm{so}(3)$ 元素。

当然，我们所在的完整希尔伯特空间有四复维，而非三复维。无论我们为静止光子可用的三维自旋子空间选取什么基，都需要第四个矢量（例如 $H_t$），以便处理光子运动的普遍情形。

为了描述黎曼光子与电子之间的相互作用，我们需要能够构造出适用于这些粒子的 $\mathrm{SU}(2)\times\mathrm{SU}(2)$ 表示之间的缠绕算符（intertwiner），无论是单个粒子，还是比方说由一个光子和一个电子组成的系统。当我们处理这样的复合系统时，适用于粒子联合量子态的表示是分别适用于各粒子的表示的张量积：
\begin{equation*}
\rho_{\text{photon and electron}}(g,h)\,\phi_{\text{photon}}\otimes\psi_{\text{electron}}
=(\rho_{\text{photon}}(g,h)\phi_{\text{photon}})\otimes(\rho_{\text{electron}}(g,h)\psi_{\text{electron}}).
\end{equation*}
适用于单个粒子的表示是不可约的，但这些张量积表示不是。确切地说，它们等价于两个或更多不可约表示的直和（direct sum）。这意味着张量积矢量空间内将存在在表示下不变的子空间——表示不会把其中任何矢量移出该子空间——而每个子空间都将等价于某个不可约表示。最简单的例子是 $\mathrm{SU}(2)$ 的两个自旋-$\tfrac12$ 表示的张量积，它包含一个等价于 $\mathrm{SU}(2)$ 的自旋-0 表示的子空间，以及另一个等价于 $\mathrm{SU}(2)$ 的自旋-1 表示的子空间。我们可以用一种简写把它写成 $\tfrac12\otimes\tfrac12=0\oplus1$。我们实际上是把自旋-$j$ 表示定义为位于 $2j$ 份自旋-$\tfrac12$ 表示张量积之内，因此那段讨论是更详细地看清这一想法的好地方。

更一般地，若把 $\mathrm{SU}(2)$ 的自旋-$j$ 表示与自旋-$k$ 表示作张量积，所得空间会分裂为自旋从 $|j-k|$ 到 $j+k$、以 $1$ 为步长的一系列子空间：
\begin{equation*}
j\otimes k=|j-k|\oplus|j-k|+1\oplus\cdots\oplus j+k .
\end{equation*}
对于 $\mathrm{SU}(2)\times\mathrm{SU}(2)$ 的表示（由一对自旋描述），我们有
\begin{equation*}
(j_1,j_2)\otimes(k_1,k_2)
=(|j_1-k_1|,|j_2-k_2|)\oplus(|j_1-k_1|+1,|j_2-k_2|)\oplus\cdots\oplus(j_1+k_1,j_2+k_2),
\end{equation*}
其中左手自旋以整数步长取遍 $|j_1-k_1|$ 与 $j_1+k_1$ 之间的每一个值，右手自旋独立地以整数步长取遍 $|j_2-k_2|$ 与 $j_2+k_2$ 之间的每一个值。

现在，我们的目标是描述诸如电子吸收光子这类过程，即从由电子和光子组成的复合系统出发，最终只剩下电子。我们可以从更基本的构件出发把它搭建起来：不从光子和完整的 Dirac 旋量开始，而从光子和一个左手或右手旋量开始。在这些情形中，四元数乘法给出一个极其简单的缠绕算符，连接初始复合系统的表示与适用于相反手性旋量的表示。换言之，我们可以描述一个左手旋量吸收光子变成右手旋量，或一个右手旋量吸收光子变成左手旋量。

在继续之前，我们把基四元数作为自旋量的解释列成表格。复结构为右乘 $H_z$。

\begin{center}
\begin{tabular}{lll}
\toprule
基四元数 & 作为自旋量态的解释 & 角动量分量 \\
\midrule
$H_t$ & UP & $m_z=\tfrac12$ \\
$H_z$ & $\ii$ UP & $m_z=\tfrac12$ \\
$H_y$ & $-\text{DOWN}$ & $m_z=-\tfrac12$ \\
$H_x$ & $-\ii\,\text{DOWN}$ & $m_z=-\tfrac12$ \\
\bottomrule
\end{tabular}
\end{center}

整体相位的选取是任意的，但这个指定给出了一个一致的方案，可把所有这些四元数解释为 $\tfrac12$ 自旋粒子的自旋态。回忆一下，由 $H_t$ 和 $H_z$ 张成的整个二维实平面应被看成一个复平面，由 $H_x$ 和 $H_y$ 张成的平面也是如此。

\subsection{电子吸收光子（Electron Absorbs Photon）}

\orthofig{images/07__013.png}{自旋量吸收光子}

我们先假定由一个右手旋量 $q_R$ 描述的粒子吸收一个光子 $q_V$，变成左手旋量 $q_L$。眼下我们探索的只是自旋，所以完全不关心这些粒子的能量与动量。在由单位四元数对 $(g,h)$ 描述的转动下，这些四元数各自以不同方式变换：
\begin{equation*}
q_R\to h\,q_R,\qquad q_V\to g\,q_V\,h^{-1},\qquad q_L\to g\,q_L .
\end{equation*}
然而可以看到，\emph{乘积} $q_V q_R$ 恰好像左手旋量那样变换：
\begin{equation*}
(q_V q_R)\to(g\,q_V\,h^{-1})(h\,q_R)=g\,(q_V q_R).
\end{equation*}
如果「右手旋量变成左手旋量」这件事用四元数做起来显得有点神秘，我们可以通过考察 $\mathrm{SU}(2)\times\mathrm{SU}(2)$ 的表示来看清\emph{为什么}这是必需的：
\begin{equation*}
(\tfrac12,\tfrac12)\otimes(0,\tfrac12)=(\tfrac12,0)\oplus(\tfrac12,1).
\end{equation*}
$(\tfrac12,\tfrac12)$ 表示是光子，$(0,\tfrac12)$ 表示是右手旋量，而 $(\tfrac12,0)$ 表示是张量积中恰好像左手旋量那样变换的一个子空间。张量积中剩下的子空间按 $(\tfrac12,1)$ 表示变换，涉及总自旋 $3/2$；我们假定不存在具有该自旋的基本粒子，因此这一表示只能由右手旋量与光子保持为分离实体来实现。

作为从描述「光子与右手旋量系统」的空间到描述左手旋量的空间的缠绕算符的候选者，我们提出
\begin{equation*}
T_1(q_V\otimes q_R)=q_V q_R/2 .
\end{equation*}
那个特定的归一化因子 $\tfrac12$ 的理由稍后说明。我们来看看这个缠绕算符对各种初始复合系统给出什么。
\begin{align*}
T_1(\text{Photon with }m_z=0,\ \text{UP spinor})
&= T_1(H_z\otimes H_t)\\
&= H_z H_t/2=H_z/2\\
&= [\ii/2]\ \text{UP}.
\end{align*}
从而，一个自旋 UP 的右手旋量可以吸收一个沿 $z$ 轴自旋为零的光子，产生一个自旋 UP 的左手旋量，该过程的振幅为 $\ii/2$。
\begin{align*}
T_1(\text{Photon with }m_z=0,\ \text{DOWN spinor})
&= T_1(H_z\otimes(-H_y))\\
&= -H_z H_y/2=H_x/2\\
&= [-\ii/2]\ \text{DOWN}.
\end{align*}
从而，一个自旋 DOWN 的右手旋量可以吸收一个沿 $z$ 轴自旋为零的光子，产生一个自旋 DOWN 的左手旋量，振幅为 $-\ii/2$。
\begin{align*}
T_1(\text{Photon with }m_z=-1,\ \text{UP spinor})
&= T_1\bigl((H_x+\ii H_y)/\sqrt2\otimes H_t\bigr)\\
&= (H_x+\ii H_y)H_t/(2\sqrt2)\\
&= (H_x+\ii H_y)/(2\sqrt2)\\
&= (H_x+H_y H_z)/(2\sqrt2)\\
&= (2H_x)/(2\sqrt2)\\
&= [-\ii/\sqrt2]\ \text{DOWN}.
\end{align*}
从第四行到第五行，我们利用了旋量四元数的复结构，把最初作为复化\emph{矢量}四元数中一个独立数字出现的因子 $\ii$，换算成了右乘 $H_z$。结论是：一个自旋 UP 的右手旋量可以吸收一个沿 $z$ 轴自旋为 $-1$ 的光子，产生一个自旋 DOWN 的左手旋量，振幅为 $-\ii/\sqrt2$。
\begin{align*}
T_1(\text{Photon with }m_z=-1,\ \text{DOWN spinor})
&= T_1\bigl((H_x+\ii H_y)/\sqrt2\otimes(-H_y)\bigr)\\
&= -(H_x+\ii H_y)H_y/(2\sqrt2)\\
&= -(H_z-\ii H_t)/(2\sqrt2)\\
&= -(H_z-H_z)/(2\sqrt2)\\
&= 0 .
\end{align*}
一个自旋 DOWN 的右手旋量与一个沿 $z$ 轴自旋为 $-1$ 的光子，产生任何左手旋量的振幅\emph{为零}。
\begin{align*}
T_1(\text{Photon with }m_z=1,\ \text{UP spinor})
&= T_1\bigl((H_x-\ii H_y)/\sqrt2\otimes H_t\bigr)\\
&= (H_x-\ii H_y)/(2\sqrt2)\\
&= (H_x-H_y H_z)/(2\sqrt2)\\
&= (H_x-H_x)/(2\sqrt2)\\
&= 0 .
\end{align*}
一个自旋 UP 的右手旋量与一个沿 $z$ 轴自旋为 $1$ 的光子，产生任何左手旋量的振幅\emph{为零}。
\begin{align*}
T_1(\text{Photon with }m_z=1,\ \text{DOWN spinor})
&= T_1\bigl((H_x-\ii H_y)/\sqrt2\otimes(-H_y)\bigr)\\
&= -(H_x-\ii H_y)H_y/(2\sqrt2)\\
&= -(H_z+\ii H_t)/(2\sqrt2)\\
&= -(H_z+H_t H_z)/(2\sqrt2)\\
&= -(2H_z)/(2\sqrt2)\\
&= [-\ii/\sqrt2]\ \text{UP}.
\end{align*}
一个自旋 DOWN 的右手旋量可以吸收一个沿 $z$ 轴自旋为 $1$ 的光子，产生一个自旋 UP 的左手旋量，振幅为 $-\ii/\sqrt2$。

最后我们注意到：
\begin{equation*}
T_1(H_t\otimes\text{UP})=\tfrac12\,\text{UP},\qquad
T_1(H_t\otimes\text{DOWN})=\tfrac12\,\text{DOWN}.
\end{equation*}
$H_t$ 对于静止光子的自旋而言并非一个可能的状态，但它确实具有 $z$ 分量为零的自旋，即 $J_z(H_t)=0$，因此它能让我们自旋量的自旋方向保持不变是合理的。

正如我们希望的那样，凡是无法在守恒 $z$ 分量自旋的同时产生某个左手旋量的情形，我们得到的振幅都为零。不过对于其余情形，概率——振幅的模平方——并不为 $1$。例如，一个自旋 UP 的自旋量与一个沿 $z$ 轴自旋为零的光子，并不一定以概率 $1$ 产生另一个自旋 UP 的自旋量。为什么不是？因为这个态\emph{并非}总自旋的本征态。它\emph{是}自旋 $z$ 分量 $m_z=\tfrac12$ 的本征态，但复合系统总自旋 $j$ 取 $3/2$ 而非 $\tfrac12$ 的振幅并不为零，这从张量积空间的分解即可看出：
\begin{equation*}
(\tfrac12,\tfrac12)\otimes(0,\tfrac12)=(\tfrac12,0)\oplus(\tfrac12,1).
\end{equation*}
我们为缠绕算符选取的归一化，使得产生一个自旋 UP 自旋量的所有不同方式的振幅模平方之和为 $1$：
\begin{align*}
&|T_1(\text{Photon with }m_z=0,\ \text{UP spinor})|^{2}
+|T_1(\text{Photon with }m_z=1,\ \text{DOWN spinor})|^{2}
+|T_1(H_t\otimes\text{UP})|^{2}\\
&\quad= |\ii/2|^{2}+|-\ii/\sqrt2|^{2}+|\tfrac12|^{2}
=\tfrac14+\tfrac12+\tfrac14=1 .
\end{align*}
对产生自旋 DOWN 的不同方式，同样成立。这让我们可以把从缠绕算符得到的每一个振幅，看作我们的某个初始态与复合系统希尔伯特空间（在按 $(\tfrac12,0)$ 表示变换的子空间内）中两个归一化矢量之一的\emph{内积}。对于这两个矢量，总自旋 $j$ 都是 $\tfrac12$，而 $z$ 分量 $m_z$ 分别为 $\tfrac12$ 或 $-\tfrac12$。所以按这种方式归一化的缠绕算符，给出的是测量 $j$ 与 $m_z$ 并发现它们取这些特定值的振幅。

对一个左手旋量吸收光子、变成右手旋量的情形，我们可以做类似处理。定义
\begin{equation*}
T_2(q_V\otimes q_L)=-q_V^*\,q_L/2 .
\end{equation*}
负号在稍后会证明是一个方便的选取。记住这里的 $*$ 运算共轭四元数，但对我们复化四元数中的数字 $\ii$ 没有影响。在转动 $(g,h)$ 下，结果像右手旋量那样变换：
\begin{align*}
(q_V^*\,q_L) &\to (g\,q_V\,h^{-1})^*(g\,q_L)\\
&= (h\,q_V^*\,g^{-1})(g\,q_L)\\
&= h\,(q_V^*\,q_L),
\end{align*}
其中用了 $g^{-1}=g^*$、$h^{-1}=h^*$，因为它们是单位四元数。

我们可以把左手和右手旋量的两个缠绕算符组合起来，给出一个黎曼 Dirac 粒子吸收光子的缠绕算符：
\begin{equation*}
T_3\bigl(q_V\otimes(q_L,q_R)\bigr)=\bigl(q_V q_R,\ -q_V^*\,q_L\bigr)/2 .
\end{equation*}
由于 Dirac 粒子同时具有左手和右手分量，它在吸收光子之前与之后适用相同的表示。不过请记住，这仍然与能量和动量的守恒无关；我们只看\emph{自旋}。

\subsection{与自旋零的耦合（Coupling With Spin Zero）}

在继续考察其它涉及光子和电子的相互作用之前，我们描述一些把光子与电子的表示连到 $\mathrm{SO}(4)$ 平凡表示的缠绕算符。平凡表示（trivial representation）描述自旋为 0 的粒子如何变换：它在坐标变换下完全不变。与其自旋相联系的希尔伯特空间是一复维的，即它就是复数 $\mathbb{C}$。因此，一个描述两个光子合并产生一个自旋-0 粒子的缠绕算符，将是从两份复化四元数的张量积到复数的函数，且当把同一个转动施加于这两个光子时函数值不变。

这样的函数极容易找到：它就是点积！加上一个合适的归一化因子，我们定义
\begin{equation*}
T_4(q_{V1}\otimes q_{V2})=q_{V1}\cdot q_{V2}/2 .
\end{equation*}
「当把同一个转动施加于 $q_{V1}$ 与 $q_{V2}$ 时点积的值不变」这一点几乎就是转动是什么的\emph{定义}。而且容易验证：两个光子沿同一根轴自旋分量为零的态（如 $H_z\otimes H_z$）给出非零结果，一个光子与另一个光子自旋相反的态也如此，而两个光子具有相同非零自旋的态则给出零。

描述反方向的缠绕算符也不难多少：它取一个自旋为 0 的粒子，给出两光子系统希尔伯特空间中的一个态，该态像自旋-0 粒子一样在转动下不变：
\begin{equation*}
T_5(z)=\frac{z}{2}\bigl(H_x\otimes H_x+H_y\otimes H_y+H_z\otimes H_z+H_t\otimes H_t\bigr).
\end{equation*}
我们选取归一化，使得当 $z$ 是单位复数时，$T_5(z)$ 是单位矢量。看出 $T_5(z)$ 在转动下不变的最简办法，是把它想成平坦欧氏空间中的度规张量：一个对角线上全是 $1$ 的矩阵。正如度规在转动下不变，这个态也不变。（严格地说，若把四元数看作矢量，这个张量态就像矢量线性函数之对偶空间上的度规；但那也像矢量空间上的度规一样在转动下不变。）

我们此前的缠绕算符 $T_4$，可以看成给出归一化不变矢量 $T_5(1)$ 与其自变量之间的内积：
\begin{equation*}
T_4(q_{V1}\otimes q_{V2})=\bigl\langle T_5(1),\ q_{V1}\otimes q_{V2}\bigr\rangle_{V\otimes V}.
\end{equation*}
我们把每个张量中各项之间的内积相乘，就得到张量积空间上的一个内积：
\begin{equation*}
\langle a\otimes b,\ c\otimes d\rangle_{V\otimes V}=\langle a,c\rangle_V\,\langle b,d\rangle_V .
\end{equation*}
把这两个缠绕算符复合，得到
\begin{align*}
T_4\bigl(T_5(z)\bigr)
&= \tfrac14\,z\,(H_x\cdot H_x+H_y\cdot H_y+H_z\cdot H_z+H_t\cdot H_t)\\
&= z .
\end{align*}

两个自旋量的不变态又如何？我们定义
\begin{align*}
T_6(z) &= \frac{z}{\sqrt2}\,(\text{UP}\otimes\text{DOWN}-\text{DOWN}\otimes\text{UP})\\
&= \frac{z}{\sqrt2}\,(H_y\otimes H_t-H_t\otimes H_y).
\end{align*}
这个态既可解释为两个左手旋量，也可解释为两个右手旋量；无论哪种情形，它都在 $\mathrm{SO}(4)$ 下不变。转动的作用是把这里的每个四元数都在左侧乘以同一个单位四元数；我们把不变性的一般证明留作练习，只完整走一个具体例子，其中来自转动的四元数为 $H_x$。
\begin{align*}
& H_y\otimes H_t-H_t\otimes H_y\\
&\to (H_x H_y)\otimes(H_x H_t)-(H_x H_t)\otimes(H_x H_y)\\
&= (H_z)\otimes(H_x)-(H_x)\otimes(H_z)\\
&= -(H_z H_z)\otimes(H_x H_z)+(H_x H_z)\otimes(H_z H_z)
\quad[\text{把第一个张量因子乘以 }-\ii\text{、第二个乘以 }\ii]\\
&= -H_t\otimes H_y+H_y\otimes H_t .
\end{align*}

最后，我们需要一个缠绕算符，它取一对相同手性的自旋量（都是左手或都是右手），给出一个在转动自旋量时不变的复数。一个显然的候选者是该双自旋量态与我们那个转动不变的态 $T_6(1)$ 的内积：
\begin{align*}
T_7(q_{L1}\otimes q_{L2})
&= \bigl\langle T_6(1),\ q_{L1}\otimes q_{L2}\bigr\rangle_{S\otimes S}\\
&= \frac{1}{\sqrt2}\Bigl[\langle H_y,q_{L1}\rangle_S\,\langle H_t,q_{L2}\rangle_S
-\langle H_t,q_{L1}\rangle_S\,\langle H_y,q_{L2}\rangle_S\Bigr].
\end{align*}
虽然逐步变换地化简成下面这个形式颇为费力，但它容易验证。两个表达式对 $q_{L1}$ 与 $q_{L2}$ 都是复线性的，且在四个基矢量 $H_t\otimes H_t$、$H_y\otimes H_y$、$H_y\otimes H_t$、$H_t\otimes H_y$ 上求值都给出相同结果。
\begin{align*}
T_7(q_{L1}\otimes q_{L2})
&= \frac{1}{\sqrt2}\bigl\langle q_{L2} H_y,\ q_{L1}\bigr\rangle_S\\
&= \frac{1}{\sqrt2}\Bigl[\operatorname{Re}_H(q_{L1}^*\,q_{L2} H_y)+\ii\,\operatorname{Re}_H(q_{L1}^*\,q_{L2} H_x)\Bigr].
\end{align*}
若把自旋 0 与自旋 $\tfrac12\otimes$ 自旋 $\tfrac12$ 之间的这两个缠绕算符复合，得到
\begin{align*}
T_7\bigl(T_6(z)\bigr)
&= \tfrac12\,z\,\Bigl[\operatorname{Re}_H(H_y^* H_t H_y)+\ii\,\operatorname{Re}_H(H_y^* H_t H_x)
-\operatorname{Re}_H(H_t^* H_y H_y)-\ii\,\operatorname{Re}_H(H_t^* H_y H_x)\Bigr]\\
&= \tfrac12\,z\,\bigl[1+0-(-1)-0\bigr]\\
&= z .
\end{align*}

\subsection{电子发射光子（Electron Emits Photon）}

\orthofig{images/07__014.png}{自旋量发射光子}

若要描述一个自旋量发射光子的过程，我们可以利用如下技巧：我们假装这个自旋量伴随一个自旋-0 粒子，后者衰变为两个光子。自旋量吸收其中一个光子，另一个则径自离去。由于我们只关心自旋表示、而不关心能量与动量，一个自旋量与一个自旋-0 粒子的初始态正好像单个自旋量那样变换，而我们最终得到一个光子和一个相反手性的自旋量，就仿佛是这个自旋量发射了光子。我们不必把这一切\emph{按字面}当真；关键在于它给出了我们所需空间之间一个具有恰好正确变换性质（transformation properties）的映射。

于是我们定义
\begin{equation*}
T_8(q_R)=T'_1\bigl(T_5(1)\otimes q_R\bigr).
\end{equation*}
$T_5(1)$ 给出了一对光子，我们把它与右手旋量作张量积，然后用一个新的缠绕算符 $T'_1$ 描述其中任一光子的吸收，另一个则不受触动。我们通过扩展 $T_1$ 的定义以容纳一个额外的光子，来得到 $T'_1$：
\begin{equation*}
T'_1(q_{V1}\otimes q_{V2}\otimes q_R)=\tfrac12\bigl(q_{V1}\otimes(q_{V2} q_R)+q_{V2}\otimes(q_{V1} q_R)\bigr).
\end{equation*}
把它写开，得到
\begin{align*}
T_8(q_R)
&= \tfrac12\,T'_1\bigl((H_x\otimes H_x+H_y\otimes H_y+H_z\otimes H_z+H_t\otimes H_t)\otimes q_R\bigr)\\
&= \tfrac12\bigl[H_x\otimes(H_x q_R)+H_y\otimes(H_y q_R)+H_z\otimes(H_z q_R)+H_t\otimes(H_t q_R)\bigr]\\
&= \tfrac12\sum_\mu\bigl[H_\mu\otimes(H_\mu q_R)\bigr].
\end{align*}
这里最终张量积中的第一个空间是被发射的光子，第二个空间是此时已变为左手性的自旋量。

例如，假定我们原来的自旋量自旋为 UP，$q_R=H_t$，则结果为
\begin{align*}
T_8(\text{UP})
&= \tfrac12\bigl[H_x\otimes H_x+H_y\otimes H_y+H_z\otimes H_z+H_t\otimes H_t\bigr]\\
&= \tfrac12\bigl[-\ii\,H_x\otimes\text{DOWN}-H_y\otimes\text{DOWN}+\ii\,H_z\otimes\text{UP}+H_t\otimes\text{UP}\bigr]\\
&= \tfrac12\bigl[-\ii\,(H_x-\ii H_y)\otimes\text{DOWN}+\ii\,(H_z-\ii H_t)\otimes\text{UP}\bigr].
\end{align*}
这里第一项 $(H_x-\ii H_y)\otimes\text{DOWN}$ 描述一个 $m_z=1$ 的光子和一个自旋 DOWN（即 $m_z=-\tfrac12$）的自旋量，所以它们总自旋 $z$ 分量为 $\tfrac12$，与原自旋量相同。第二项 $(H_z-\ii H_t)\otimes\text{UP}$ 中，光子部分是 $J_z$ 的一个本征态、本征值为零，因此我们再次得到自旋 $z$ 分量的守恒，尽管这种光子偏振只对运动中的粒子才可能。

要描述一个 Dirac 自旋量发射光子，我们把使用 $T'_1$ 与一个类似定义的 $T'_2$ 的两种构造组合起来：
\begin{align*}
T_9(q_L,q_R)
&= \tfrac12\Bigl(T'_1\bigl((H_x\otimes H_x+H_y\otimes H_y+H_z\otimes H_z+H_t\otimes H_t)\otimes q_R\bigr),\\
&\qquad\quad T'_2\bigl((H_x\otimes H_x+H_y\otimes H_y+H_z\otimes H_z+H_t\otimes H_t)\otimes q_L\bigr)\Bigr)\\
&= \tfrac12\sum_\mu\bigl[H_\mu\otimes(H_\mu q_R,\ -H_\mu^* q_L)\bigr].
\end{align*}
如果我们让一个 Dirac 自旋量用 $T_9$ 发射一个光子，再把结果喂给 $T_3$，使该自旋量重新吸收它自己发射的光子，会怎样？
\begin{align*}
T_3\bigl(T_9(q_L,q_R)\bigr)
&= \tfrac14\sum_\mu\bigl(-H_\mu H_\mu^*\,q_L,\ -H_\mu^* H_\mu\,q_R\bigr)\\
&= -(q_L,q_R).
\end{align*}
于是我们重新得到原来的 Dirac 自旋量乘以 $-1$。

\subsection{电子与正电子湮灭（Electron and Positron Annihilate）}

\orthofig{images/07__015.png}{两个自旋量湮灭产生光子}

当一个电子与一个正电子湮灭时，守恒能量和动量的唯一方式是产生\emph{两个}光子。事实上，同样的判据也适用于自由电子吸收或发射光子：电子必须散射一个入射光子，它不能吸收或发射单个光子。

但在本节中我们只关心自旋，而不关心能量和动量，所以我们就着手构造一个描述电子与正电子相合产生单个光子的缠绕算符。这其实是一件有用的事，因为这种只涉及三个基本粒子的相互作用，可以看作一个能量与动量\emph{确实}守恒的更大过程（一个更大费曼图中的一个顶点）中的一步。

要把两个自旋量组合成一个矢量，必须一个自旋量是左手、另一个是右手。我们可以用此前构造的两个缠绕算符来满足需要：$T_8$ 把右手旋量变成左手旋量和光子，然后 $T_7$ 把这个新的左手旋量与原来的左手旋量组合成一个自旋-0 粒子——一个我们可以直接乘进结果的复数。

回顾：
\begin{align*}
T_8(q_R) &= \tfrac12\sum_\mu\bigl[H_\mu\otimes(H_\mu q_R)\bigr],\\
T_7(q_{L1}\otimes q_{L2})
&= \Bigl\langle (H_y\otimes H_t-H_t\otimes H_y)/\sqrt2,\ q_{L1}\otimes q_{L2}\Bigr\rangle_{S\otimes S}.
\end{align*}
不过，与其简单地组合这两个缠绕算符，我们将构造一个平等对待左手和右手旋量的版本，让最初为左手的旋量也在发射光子中起作用。
\begin{align*}
T_{10}(q_L\otimes q_R)
&= \frac{1}{2\sqrt2}\sum_\mu
\Bigl\langle H_y\otimes H_t-H_t\otimes H_y,\ q_L\otimes(H_\mu q_R)-q_R\otimes(H_\mu^* q_L)\Bigr\rangle_{S\otimes S}H_\mu\\
&= \frac{1}{2\sqrt2}\sum_\mu\Bigl[
\langle H_y,q_L\rangle_S\,\langle H_t,H_\mu q_R\rangle_S
-\langle H_t,q_L\rangle_S\,\langle H_y,H_\mu q_R\rangle_S\\
&\qquad\qquad\qquad
-\langle H_y,q_R\rangle_S\,\langle H_t,H_\mu^* q_L\rangle_S
+\langle H_t,q_R\rangle_S\,\langle H_y,H_\mu^* q_L\rangle_S\Bigr]H_\mu .
\end{align*}
化简上式是件费力的事，但我们其实可以根据一些基本原理走得很远，来预判答案。唯一能像复化矢量四元数那样变换的左手、右手旋量组合，将是
\begin{equation*}
q_L\,(a+\ii\,b)\,q_R^* .
\end{equation*}
为使该表达式复线性，若把 $q_L$ 或 $q_R$ 右乘 $H_z$，结果必须与整体乘以 $\ii$ 相同，于是有
\begin{equation*}
H_z\,b=-b\,H_z=a,\qquad H_z\,a=-a\,H_z=-b .
\end{equation*}
在第一式中令 $b=H_x$，得 $a=H_y$，而这两个值也满足第二式。取同一选择的任意实倍数都可行，而且不难验证：若令
\begin{equation*}
T_{10}(q_L\otimes q_R)=-\frac{\ii}{\sqrt2}\,q_L\,(H_x-\ii H_y)\,q_R^*,
\end{equation*}
就能精确地与我们原来的定义吻合。例如：
\begin{equation*}
T_{10}(\text{UP}\otimes\text{UP})=T_{10}(H_t\otimes H_t)=-\frac{\ii}{\sqrt2}(H_x-\ii H_y).
\end{equation*}
结果是 $-\ii$ 乘以我们那个 $m_z=1$ 的归一化光子态。我们还有
\begin{align*}
T_{10}(\text{DOWN}\otimes\text{DOWN})
&= T_{10}\bigl((-H_y)\otimes(-H_y)\bigr)\\
&= \frac{\ii}{\sqrt2}\,H_y\,(H_x-\ii H_y)\,H_y\\
&= \frac{\ii}{\sqrt2}\,(H_x+\ii H_y),
\end{align*}
所以结果恰为 $\ii$ 乘以我们那个 $m_z=-1$ 的归一化光子态。对于其它旋量组合，我们得到
\begin{align*}
T_{10}(\text{UP}\otimes\text{DOWN})
&= T_{10}\bigl(H_t\otimes(-H_y)\bigr)\\
&= -\frac{\ii}{\sqrt2}\,(H_x-\ii H_y)\,H_y
= -\frac{\ii}{\sqrt2}\,(H_z+\ii H_t),\\
T_{10}(\text{DOWN}\otimes\text{UP})
&= T_{10}\bigl((-H_y)\otimes H_t\bigr)\\
&= \frac{\ii}{\sqrt2}\,H_y\,(H_x-\ii H_y)
= -\frac{\ii}{\sqrt2}\,(H_z-\ii H_t).
\end{align*}
两种情形下结果都是 $m_z=0$ 的归一化本征态。乍看上去，下面这个结果似乎是个错误或麻烦：
\begin{equation*}
T_{10}\bigl((\text{UP}\otimes\text{DOWN}-\text{DOWN}\otimes\text{UP})/\sqrt2\bigr)=H_t,
\end{equation*}
因为那个初始态看起来应当是一个总自旋为 0 而非 $1$ 的本征态！但一个自旋量是左手、另一个是右手这一事实在此起了关键区别；手性表示的张量积是 $(\tfrac12,0)\otimes(0,\tfrac12)=(\tfrac12,\tfrac12)$，即整个空间的总自旋为 $1$。

$T_{10}$ 是两个单个自旋量的缠绕算符。对于两个 Dirac 粒子，我们有
\begin{equation*}
T_{11}\bigl((q_{L1},q_{R1})\otimes(q_{L2},q_{R2})\bigr)
=-\frac{\ii}{\sqrt2}\Bigl(q_{L1}(H_x-\ii H_y)q_{R2}^*+q_{L2}(H_x-\ii H_y)q_{R1}^*\Bigr).
\end{equation*}

\subsection{对产生（Pair Production）}

\orthofig{images/07__016.png}{对产生}

上一节过程的逆过程涉及一个光子变成两个自旋量，一个左手、一个右手。与我们其它例子一样，这不是自由光子能在守恒能量和动量的同时做到的事，但我们仍然有兴趣理解这些自旋之间的关系。

我们用来得到这个缠绕算符的技巧，是想象一个自旋-0 粒子衰变成一对右手旋量，其中一个径自离去，另一个吸收我们的光子并变成左手旋量。我们再把它与一个版本组合：那里得到一对左手旋量，其中一个径自离去，另一个吸收光子并变成右手。与我们其它构造一样，这不应按字面理解；我们只是利用了一个事实：把初始态与一个自旋-0 表示作张量积，不会影响它的变换方式。

来自我们虚构的自旋-0 粒子的那对自旋量（无论都是左手还是都是右手），由如下归一化态给出：
\begin{align*}
T_6(1)
&= \frac{1}{\sqrt2}(\text{UP}\otimes\text{DOWN}-\text{DOWN}\otimes\text{UP})\\
&= \frac{1}{\sqrt2}(H_y\otimes H_t-H_t\otimes H_y).
\end{align*}
我们以模仿缠绕算符 $T_1$ 与 $T_2$ 的过程吸收光子，留下一个含一个左手自旋量和一个右手自旋量的态：
\begin{equation*}
T_{12}(q_V)=\frac{1}{2\sqrt2}\Bigl((q_V H_y)\otimes H_t-(q_V H_t)\otimes H_y
+H_y\otimes(q_V^* H_t)-H_t\otimes(q_V^* H_y)\Bigr).
\end{equation*}
一些例子：
\begin{align*}
T_{12}(\text{Photon with }m_z=1)
&= T_{12}\bigl((H_x-\ii H_y)/\sqrt2\bigr)\\
&= \tfrac14\Bigl[((H_x-\ii H_y)H_y)\otimes H_t-((H_x-\ii H_y)H_t)\otimes H_y\\
&\qquad\qquad -H_y\otimes((H_x-\ii H_y)H_t)+H_t\otimes((H_x-\ii H_y)H_y)\Bigr]\\
&= \tfrac14\Bigl[(H_z+H_z)\otimes H_t-(H_x-H_x)\otimes H_y\\
&\qquad\qquad -H_y\otimes(H_x-H_x)+H_t\otimes(H_z+H_z)\Bigr]\\
&= \tfrac12\bigl[H_z\otimes H_t+H_t\otimes H_z\bigr]\\
&= \ii\,\text{UP}\otimes\text{UP}.
\end{align*}
\begin{align*}
T_{12}(\text{Photon with }m_z=-1)
&= T_{12}\bigl((H_x+\ii H_y)/\sqrt2\bigr)\\
&= \tfrac14\Bigl[((H_x+\ii H_y)H_y)\otimes H_t-((H_x+\ii H_y)H_t)\otimes H_y\\
&\qquad\qquad -H_y\otimes((H_x+\ii H_y)H_t)+H_t\otimes((H_x+\ii H_y)H_y)\Bigr]\\
&= \tfrac14\Bigl[(H_z-H_z)\otimes H_t-(H_x+H_x)\otimes H_y\\
&\qquad\qquad -H_y\otimes(H_x+H_x)+H_t\otimes(H_z-H_z)\Bigr]\\
&= -\tfrac12\bigl[H_x\otimes H_y+H_y\otimes H_x\bigr]\\
&= -\ii\,\text{DOWN}\otimes\text{DOWN}.
\end{align*}
\begin{align*}
T_{12}(\text{Photon with }m_z=0)
&= T_{12}(H_z)\\
&= \frac{1}{2\sqrt2}\Bigl[(H_z H_y)\otimes H_t-(H_z H_t)\otimes H_y
-H_y\otimes(H_z H_t)+H_t\otimes(H_z H_y)\Bigr]\\
&= -\frac{1}{2\sqrt2}\Bigl[H_x\otimes H_t+H_z\otimes H_y+H_y\otimes H_z+H_t\otimes H_x\Bigr]\\
&= -\frac{1}{2\sqrt2}\Bigl[(-\ii\,\text{DOWN})\otimes\text{UP}+(\ii\,\text{UP})\otimes(-\text{DOWN})\\
&\qquad\qquad\qquad +(-\text{DOWN})\otimes(\ii\,\text{UP})+\text{UP}\otimes(-\ii\,\text{DOWN})\Bigr]\\
&= \frac{\ii}{\sqrt2}\bigl(\text{DOWN}\otimes\text{UP}+\text{UP}\otimes\text{DOWN}\bigr).
\end{align*}
\begin{align*}
T_{12}(H_t)
&= \frac{1}{\sqrt2}(H_y\otimes H_t-H_t\otimes H_y)\\
&= \frac{1}{\sqrt2}(\text{UP}\otimes\text{DOWN}-\text{DOWN}\otimes\text{UP}).
\end{align*}

若把缠绕算符 $T_{10}$ 作用于以上任何具体结果，我们都会重新得到原来的光子态。由于这些结果构成一组基矢量，这一点对任何态都必定成立：
\begin{equation*}
T_{10}\bigl(T_{12}(q_V)\bigr)=q_V .
\end{equation*}

\subsection{缠绕算符的不变性（Invariance of the Intertwiners）}

\orthofig{images/07__017.png}{三价自旋量-自旋量-光子顶点}

我们已经构造出可以用来描述以下过程的缠绕算符：
\begin{itemize}
  \item 一个电子吸收光子；
  \item 一个电子发射光子；
  \item 一个电子与一个正电子湮灭产生一个光子；
  \item 一个光子衰变成一个电子与一个正电子。
\end{itemize}
我们应当再次强调，我们只看这些粒子的\emph{自旋}；守恒能量和动量是另一回事。

现在，从一个拒绝偏爱任何特定观测者的完全客观的观点来看，所有这些过程都是同一件事的例子：两个自旋量和一个光子之间的相互作用。通过转动我们的参照系，我们可以把这些过程中的任何一个变成其它任何一个。

在两个观测者对相互作用前后存在哪些粒子拥有一致描述的情形中，结果的转动不变性很容易看出。例如，假定一个观测者 Carla 说一个电子吸收了一个光子，并把态 $q_V$ 归给光子、把 Dirac 旋量 $(q_{L1},q_{R1})$ 归给电子。如果她关心电子最终成为 Dirac 旋量 $(q_{L2},q_{R2})$ 的振幅，她将用缠绕算符 $T_3$ 把它算作
\begin{align*}
A &= \bigl\langle (q_{L2},q_{R2}),\ T_3\bigl(q_V\otimes(q_{L1},q_{R1})\bigr)\bigr\rangle_D\\
&= \tfrac12\bigl\langle (q_{L2},q_{R2}),\ (q_V q_{R1},\ q_V^* q_{L1})\bigr\rangle_D\\
&= \tfrac12\Bigl[\langle q_{L2},q_V q_{R1}\rangle_S+\langle q_{R2},q_V^* q_{L1}\rangle_S\Bigr].
\end{align*}
这里我们在 Dirac 旋量上定义了一个内积 $\langle\cdot,\cdot\rangle_D$，即左手与右手旋量之间旋量内积 $\langle\cdot,\cdot\rangle_S$ 之和。

若另一个观测者 Patrizia 同意是一个光子吸收了一个电子，但使用一组相对于 Carla 的坐标转动了单位四元数对 $(g,h)$ 的坐标，她将把不同的自旋态归给这些粒子：
\begin{equation*}
q'_V=g\,q_V\,h^{-1},\qquad
(q'_{L1},q'_{R1})=(g\,q_{L1},\ h\,q_{R1}),\qquad
(q'_{L2},q'_{R2})=(g\,q_{L2},\ h\,q_{R2}).
\end{equation*}
使用同一个缠绕算符，但应用于她自己版本的自旋态，Patrizia 将预言一个与 Carla 完全相同的振幅：
\begin{align*}
A' &= \tfrac12\Bigl[\langle q'_{L2},q'_V q'_{R1}\rangle_S+\langle q'_{R2},q'^*_V q'_{L1}\rangle_S\Bigr]\\
&= \tfrac12\Bigl[\langle g\,q_{L2},g\,q_V h^{-1}h\,q_{R1}\rangle_S+\langle h\,q_{R2},h\,q_V^* g^{-1}g\,q_{L1}\rangle_S\Bigr]\\
&= \tfrac12\Bigl[\langle g\,q_{L2},g\,q_V q_{R1}\rangle_S+\langle h\,q_{R2},h\,q_V^* q_{L1}\rangle_S\Bigr]\\
&= \tfrac12\Bigl[\langle q_{L2},q_V q_{R1}\rangle_S+\langle q_{R2},q_V^* q_{L1}\rangle_S\Bigr]\\
&= A .
\end{align*}
从倒数第三行到倒数第二行，我们用了一个此前已确立的事实：当两个自变量都在左侧乘以同一个单位四元数时，旋量内积不变。

\orthofig{images/07__018.png}{离开相互作用时的自旋}

不过，当不同观测者对相互作用\emph{之前}存在哪些粒子、\emph{之后}存在哪些粒子出现分歧时，情形就变得棘手了。要比较这类观测者的测量与预言，我们不仅需要从我们四个缠绕算符中的一个切换到另一个，还需要做比用通常表示去变换粒子自旋态更多的事情。

为什么？经典地看，\emph{自旋}取决于当你看着物体\emph{随时间向前运动}时它旋转的方向。正如不存在绝对的能动量矢量，也不存在绝对的自旋矢量（或张量）；一个从上方看绕 $z$ 轴逆时针旋转的物体，在具有相反时间箭头的观测者看来将是顺时针旋转的。

然而，我们\emph{确实}可以客观地谈论一个物体在接近某个特定相互作用时其旋转的方向。这会与把该物体视为存在于相互作用\emph{之前}的观测者的自旋方向一致，但会与把它视为存在于相互作用\emph{之后}的观测者的自旋方向相反。如果我们同意用这种措辞来表达\emph{所有}自旋，那么（颇像处理能动量守恒时那种观测者无关的做法）我们可以看到角动量在每一个相互作用处都守恒。

量子力学上，我们可以用此前描述过的 $\operatorname{rev}_S$ 函数来反转我们四元数自旋量的自旋：
\begin{equation*}
\operatorname{rev}_S(q)=q\,H_y .
\end{equation*}
对于用作光子态的复化矢量四元数，要反转自旋，我们只需在复结构意义上共轭它们：
\begin{equation*}
\operatorname{rev}_V(p+q\,\ii)=\operatorname{conj}(p+q\,\ii)=p-q\,\ii .
\end{equation*}
一旦我们把原来看作\emph{离开}相互作用而非接近相互作用的所有态都作了自旋反转，我们就应发现三个态的整个张量积具有总自旋 $0$。我们可以用一个从该张量积到自旋 0 的缠绕算符来量化这一点，其构造办法是想象这两个自旋量湮灭，然后取所得光子与原来光子的点积：
\begin{align*}
T_{\mathrm{vert}}(q_V\otimes q_L\otimes q_R)
&= T_4\bigl(q_V\otimes T_{10}(q_L\otimes q_R)\bigr)\\
&= -\frac{\ii}{2\sqrt2}\,q_V\cdot\bigl(q_L\,(H_x-\ii H_y)\,q_R^*\bigr).
\end{align*}
我们这样构造的缠绕算符 $T_{\mathrm{vert}}$ 是这样归一化的：对张量积空间的任何正交归一基求和，结果的模平方之和为 $1$。这些结果可以被看作张量积空间中位于此处的单个归一化矢量（一个总自旋本征值为 $j=0$ 的本征态）与我们喂给缠绕算符的任何三粒子态之间的内积。不过，当我们把 $T_{\mathrm{vert}}$ 与其它缠绕算符比较时，会得到一些额外的因子：
\begin{align*}
\langle q_L,\ T_1(q_V\otimes q_R)\rangle_S
&= (\sqrt2)\,T_{\mathrm{vert}}\bigl(q_V\otimes\operatorname{rev}_S(q_L)\otimes q_R\bigr)
&&\text{[右手旋量吸收光子变成左手旋量]}\\
\langle q_R,\ T_2(q_V\otimes q_L)\rangle_S
&= (\sqrt2)\,T_{\mathrm{vert}}\bigl(q_V\otimes q_L\otimes\operatorname{rev}_S(q_R)\bigr)
&&\text{[左手旋量吸收光子变成右手旋量]}\\
\langle q_V\otimes q_L,\ T_8(q_R)\rangle_{V\otimes S}
&= (\sqrt2)\,T_{\mathrm{vert}}\bigl(\operatorname{rev}_V(q_V)\otimes\operatorname{rev}_S(q_L)\otimes q_R\bigr)
&&\text{[右手旋量发射光子变成左手旋量]}\\
\langle q_V\otimes q_R,\ T_{8L}(q_L)\rangle_{V\otimes S}
&= (\sqrt2)\,T_{\mathrm{vert}}\bigl(\operatorname{rev}_V(q_V)\otimes q_L\otimes\operatorname{rev}_S(q_R)\bigr)
&&\text{[左手旋量发射光子变成右手旋量]}\\
\langle q_V,\ T_{10}(q_L\otimes q_R)\rangle_V
&= 2\,T_{\mathrm{vert}}\bigl(\operatorname{rev}_V(q_V)\otimes q_L\otimes q_R\bigr)
&&\text{[左手与右手旋量湮灭产生光子]}\\
\langle q_L\otimes q_R,\ T_{12}(q_V)\rangle_{S\otimes S}
&= 2\,T_{\mathrm{vert}}\bigl(q_V\otimes\operatorname{rev}_S(q_L)\otimes\operatorname{rev}_S(q_R)\bigr)
&&\text{[光子变成一对左手与右手旋量]}
\end{align*}
[我们在此没有费心去显式构造上面称为 $T_{8L}$ 的东西，即 $T_8$ 的一个以左手旋量开始的版本；但其精确定义可见于完整 Dirac 旋量的 $T_9$ 的右手分量，就在「电子发射光子」一节的末尾。]

$\sqrt2$ 与 $2$ 这些因子出现，是因为我们根据\emph{初始}存在哪些粒子来归一化这些不同的缠绕算符。例如，我们归一化 $T_1$ 与 $T_2$ 使其给出初始态总自旋为 $\tfrac12$、末态为某个特定自旋量的振幅。由于初始态空间中总自旋为 $\tfrac12$ 的子空间维数为 $2$，若对\emph{所有}初始态与末态的正交归一基求和，我们得到所有模平方之和为 $2$。但 $T_{\mathrm{vert}}$ 给出的是总自旋为 $0$ 的振幅，而初始态空间中只有一维子空间满足该条件，于是同类的求和只给出 $1$。因子 $\sqrt2$ 就是这两者之比的平方根。

为了从中提取出某种完全\emph{观测者无关}的东西，假定一个观测者看到一个确实涉及一个左手旋量、一个右手旋量以及一个光子的相互作用。他们测量了所有三个粒子的态，然后对他们视为出射的态施加自旋反转函数。给定三次测量所用的基，以及其中任意两个粒子的实际结果，我们应当能够为第三个粒子的所有基元素计算相对概率，无论原观测者把它看作入射粒子还是出射粒子。

我们计算相对概率为 $|T_{\mathrm{vert}}(q_V\otimes q_L\otimes q_R)|^{2}$，固定三个态中的两个为已知结果，同时把第三个态依次取为测量基的每个元素。然后我们除以这些值之和，以得到真正的概率。由于归一化中任何整体因子都会抵消，用上面列出的六个缠绕算符中的任何一个执行同样的程序，都会给出与 $T_{\mathrm{vert}}$ 完全相同的概率。缠绕算符之间除整体归一化外的所有差别，都由自旋反转来处理掉了。

\subsection{四元数拉格朗日量（Quaternionic Lagrangian）}

既然我们已经看到如何把矢量和旋量都当作四元数处理，把黎曼量子电动力学中其它定义和方程转成四元数形式就不难了。作为一个简单的例子，我们来看黎曼 QED 拉格朗日量中的相互作用项：
\begin{equation*}
L_{\mathrm{inter}}=-e\,A_\mu\,\psi^\dagger\,\gamma^\mu\,\psi .
\end{equation*}
用四元数并且工作在 Weyl 基中，我们有
\begin{equation*}
\psi=(q_L\bv_0,\ q_R\bv_0),\qquad
\psi^\dagger=(\bv_0^\dagger q_L^*,\ \bv_0^\dagger q_R^*),\qquad
q_V=A_\mu H_\mu,\qquad
\gamma^\mu=\begin{pmatrix}0&H_\mu\\ H_\mu^*&0\end{pmatrix}.
\end{equation*}
这里 $\bv_0$ 是 $\mathbb{C}^{2}$ 中一个单位矢量，我们用它把 $\mathbb{C}^{2}$ 映射到四元数；我们通常的选择是 $\bv_0=(0,1)$，以与「右乘 $H_z$ 作为四元数上的复结构」相容。我们用 $\dagger$ 表示复矢量与复矩阵的共轭转置，而我们所用于四元数共轭的 $*$，若把四元数看作 $2\times2$ 复矩阵，实际上与之相同。

于是我们有
\begin{align*}
L_{\mathrm{inter}}
&= -e\,A_\mu\Bigl[\bv_0^\dagger q_L^* H_\mu q_R \bv_0+\bv_0^\dagger q_R^* H_\mu^* q_L \bv_0\Bigr]\\
&= -e\,\bv_0^\dagger\Bigl[q_L^* q_V q_R+q_R^* q_V^* q_L\Bigr]\bv_0 .
\end{align*}
上一行方括号里的量作为一个四元数是纯实数，因为它等于自身的四元数共轭，所以作为一个 $2\times2$ 矩阵，它只是单位矩阵的一个倍数。于是我们有
\begin{align*}
L_{\mathrm{inter}}
&= -e\,(\bv_0^\dagger \bv_0)\Bigl[q_L^* q_V q_R+q_R^* q_V^* q_L\Bigr]\\
&= -e\Bigl[q_L^* q_V q_R+q_R^* q_V^* q_L\Bigr].
\end{align*}
不难看出，给定 $\mathrm{SU}(2)\times\mathrm{SU}(2)$ 中的单位四元数对 $(g,h)$，它在通常的变换下不变：
\begin{equation*}
q_V\to g\,q_V\,h^{-1},\qquad q_L\to g\,q_L,\qquad q_R\to h\,q_R .
\end{equation*}
你也许会疑惑，为什么我们不能把这个函数 $L_{\mathrm{inter}}$ 用作从「光子与两个自旋量」态到自旋 0 的缠绕算符。但拉格朗日量是一个实数，而在这一语境中 $q_V$ 只是一个普通四元数，而非复化四元数。我们此前构造的缠绕算符 $T_{\mathrm{vert}}$ 对三个自旋态都是复线性的。它是唯一的（至多相差一个整体常数因子），我们不能用拉格朗日量来代替它。

作为对结果的检验，假定左右手旋量四元数确实满足能动量矢量为 $Q_V=k^\mu H_\mu$ 的平面波四元数 Dirac 方程，即
\begin{equation*}
Q_V q_R-m q_L=0 .
\end{equation*}
我们可以把相互作用能写成拉格朗日量的相反数，然后作代换 $q_L=Q_V q_R/m$：
\begin{align*}
H_{\mathrm{inter}}
&= -L_{\mathrm{inter}}
= e\Bigl[q_L^* q_V q_R+q_R^* q_V^* q_L\Bigr]\\
&= \frac{e}{m}\Bigl[q_R^* Q_V^* q_V q_R+q_R^* q_V^* Q_V q_R\Bigr]\\
&= \frac{e}{m}\,q_R^*\Bigl[Q_V^* q_V+q_V^* Q_V\Bigr]q_R\\
&= (q_R^* q_R)\,\frac{e}{m}\Bigl[Q_V^* q_V+q_V^* Q_V\Bigr]\\
&= \tfrac12\,\frac{e}{m}\Bigl[k^\mu H_\mu^* A_\nu H_\nu+A_\nu H_\nu^* k^\mu H_\mu\Bigr]\\
&= \tfrac12\,\frac{e}{m}\,k^\mu A_\nu\Bigl[H_\mu^* H_\nu+H_\nu^* H_\mu\Bigr]\\
&= \tfrac12\,\frac{e}{m}\,k^\mu A_\nu\,(2\delta^{\nu}_{\ \mu})\\
&= \frac{e}{m}\,k^\mu A_\mu,
\end{align*}
其中从第四行到第五行，我们用了方括号中的表达式是实数这一事实（它是 $Q_V^* q_V$ 与其共轭之和），因而与 $q_R$ 相交换，随后我们又用了 $q_R^* q_R=|q_R|^{2}=\tfrac12$，这是当我们想要 $\psi^\dagger\psi=1$ 时所应取的归一化。

这个结果显然是不变量，而在电子静止、$\bk=m\be_t$ 的参照系中它变为
\begin{equation*}
H_{\mathrm{inter}}=e\,A_t=-e\,V,
\end{equation*}
其中 $V$ 是势能。我们预期会出现一个负号，因为在黎曼宇宙中，真实的相对论能量是势能的相反数。

% 本节局部宏（preamble 未定义，用 \providecommand 以防重复定义）
\providecommand{\bchi}{\boldsymbol{\chi}}
\providecommand{\bdelta}{\boldsymbol{\delta}}
\providecommand{\by}{\mathbf{y}}
\providecommand{\Chi}{\mathcal{X}}

\section{进阶：黎曼量子场论（Riemannian Quantum Field Theory）}

本节讨论黎曼宇宙中量子场论（quantum field theory, QFT）的构造。核心思路是：
把谐振子（harmonic oscillator）的算符代数推及自由标量场（scalar field）与自由 Dirac 场，
并借助四维环面 $T^{4}$ 上有限的传播模式集合来完成正则量子化（canonical quantisation）。
我们将看到，与洛伦兹情形最深刻的差别在于：黎曼宇宙中模式的数目是\emph{有限}的，
因而真空能是有限值；同时“粒子”与“反粒子”的划分依赖于观测者的时间轴，
使得建立在三维球面上的 $SO(4)$ 不变描述比通常的洛伦兹描述更为自然。

\subsection{量子力学谐振子（The Quantum Mechanical Harmonic Oscillator）}

我们从最简单的非相对论量子力学谐振子出发，先回顾其经典与量子版本的关键结果，
为后续场论量子化提供算符代数的模板。

\subsubsection*{经典力学准备}

设一个质量为 $M$ 的质点在一维坐标 $x$ 上运动，受到回复力 $F=-Kx$。
该力与质点到原点的距离成正比，对应势能 $V=\tfrac12 Kx^{2}$，即一个抛物线形的“能量谷”。
质点将在谷中以角频率
\begin{equation*}
\omega=\sqrt{\frac{K}{M}}
\end{equation*}
来回振动。以坐标 $x$ 与相应动量 $p$ 表示的哈密顿量为
\begin{equation}
H_{\text{SHO}}=\frac{p^{2}}{2M}+\tfrac12 M\omega^{2}x^{2},
\tag{Hamiltonian}
\end{equation}
其中我们把“弹簧常数” $K$ 写成 $M\omega^{2}$，因为二者完全等同，而 $\omega$ 比 $K$ 更值得我们关心。

以量上加点表示其对时间的导数，本系统的哈密顿方程（Hamilton's equations）为
\begin{equation*}
p^{\bullet}=-\partial_{x}H_{\text{SHO}}=-M\omega^{2}x,
\qquad
x^{\bullet}=\partial_{p}H_{\text{SHO}}=\frac{p}{M}.
\end{equation*}
合起来便得到
\begin{equation*}
\partial_{t}^{2}x=-\omega^{2}x,
\end{equation*}
其通解为
\begin{equation*}
x(t)=A\cos(\omega t)+B\sin(\omega t).
\end{equation*}

若改用拉格朗日量（Lagrangian）——动能减势能，以坐标 $x$ 及其时间变化率 $x^{\bullet}$ 表示：
\begin{equation}
L_{\text{SHO}}=\tfrac12 M\bigl(x^{\bullet}{}^{2}-\omega^{2}x^{2}\bigr),
\tag{Lagrangian}
\end{equation}
则运动方程由欧拉-拉格朗日方程（Euler-Lagrange equation）给出：
\begin{equation*}
\partial_{t}\bigl(\partial_{x^{\bullet}}L_{\text{SHO}}\bigr)=\partial_{x}L_{\text{SHO}}
\;\Longrightarrow\;
\partial_{t}(M x^{\bullet})=-M\omega^{2}x
\;\Longrightarrow\;
\partial_{t}^{2}x=-\omega^{2}x,
\end{equation*}
与由哈密顿量所得方程自然一致。

\subsubsection*{量子化与阶梯算符}

在量子力学中，描述质点位置与动量的数 $x$、$p$ 被替换为“可观测量”：
希尔伯特空间上的自伴算符。一种看法是把该希尔伯特空间视为复值波函数 $\psi(x,t)$ 的空间。
数 $x$ 被替换为“乘以 $x$”的算符，而数 $p$ 被替换为微分算符 $-\ii\partial_{x}$。
（这里做的是\emph{非相对论}量子力学，故不像黎曼相对论量子力学那样改变动量算符的符号。）
要求系统由某给定哈密顿量描述，即要求波函数满足薛定谔方程（Schrödinger equation）：
\begin{equation*}
\ii\partial_{t}\psi(x,t)=H\psi(x,t)
=\Bigl(\frac{p^{2}}{2M}+\tfrac12 M\omega^{2}x^{2}\Bigr)\psi(x,t)
=-\frac{1}{2M}\partial_{x}^{2}\psi(x,t)+\tfrac12 M\omega^{2}x^{2}\psi(x,t).
\end{equation*}
不过，与其直接求该微分方程的解，对我们更有用的是把量子化后谐振子的希尔伯特空间
看作由具有确定能量的矢量（即哈密顿量的\emph{本征矢量}）张成的，而不去操心对应波函数的细节。

首先注意到动量与位置算符的对易子为
\begin{equation}
[x,p]=xp-px=\ii .
\tag{commutator}
\end{equation}
这在波函数表示中极易检验：对任意 $x$ 的函数 $f$，
\begin{equation*}
[x,p]f=[x,-\ii\partial_{x}]f
=x(-\ii\partial_{x}f)-(-\ii\partial_{x}(xf))
=-\ii x\partial_{x}f+\ii(f+x\partial_{x}f)=\ii f .
\end{equation*}

由 $x$ 与 $p$ 可构造两个新算符，称为\emph{阶梯算符}（ladder operators）：
\begin{equation}
a=\frac{\ii\,(p-\ii M\omega x)}{\sqrt{2M\omega}},
\qquad
a^{\dagger}=\frac{-\ii\,(p+\ii M\omega x)}{\sqrt{2M\omega}}.
\tag{ladder}
\end{equation}
这里 $a^{\dagger}$ 是 $a$ 的厄米共轭；由于 $p$、$x$ 自伴，求 $a^{\dagger}$ 只需把 $a$ 的定义中
每个 $\ii$ 反号。由 $x$ 与 $p$ 的对易子不难证明
\begin{equation*}
[a,a^{\dagger}]=aa^{\dagger}-a^{\dagger}a=1 .
\end{equation*}
由此可轻松导出若干我们将反复用到的结果：
\begin{equation*}
aa^{\dagger}=a^{\dagger}a+1,\qquad
a^{\dagger}aa^{\dagger}=a^{\dagger}(a^{\dagger}a+1),\qquad
a^{\dagger}a=aa^{\dagger}-1,\qquad
a^{\dagger}aa=(aa^{\dagger}-1)a=a(a^{\dagger}a-1).
\end{equation*}
哈密顿量可用阶梯算符极为简洁地表示：
\begin{equation*}
H_{\text{SHO}}=\frac{p^{2}}{2M}+\tfrac12 M\omega^{2}x^{2}=\omega\bigl(a^{\dagger}a+\tfrac12\bigr).
\end{equation*}

设 $\psi_{E}$ 为哈密顿量的任一本征矢量，满足 $H_{\text{SHO}}\psi_{E}=E\psi_{E}$。于是
\begin{align*}
H_{\text{SHO}}(a^{\dagger}\psi_{E})
&=\omega(a^{\dagger}a+\tfrac12)(a^{\dagger}\psi_{E})
=a^{\dagger}\omega(a^{\dagger}a+\tfrac12+1)\psi_{E}\\
&=a^{\dagger}H_{\text{SHO}}\psi_{E}+\omega(a^{\dagger}\psi_{E})
=(E+\omega)(a^{\dagger}\psi_{E}),
\end{align*}
即 $a^{\dagger}$ 把哈密顿量的任一本征矢量变为能量高出 $\omega$ 的另一本征矢量。类似地
\begin{align*}
H_{\text{SHO}}(a\psi_{E})
&=\omega(a^{\dagger}a+\tfrac12)(a\psi_{E})
=a\,\omega(a^{\dagger}a+\tfrac12-1)\psi_{E}\\
&=aH_{\text{SHO}}\psi_{E}-\omega(a\psi_{E})
=(E-\omega)(a\psi_{E}),
\end{align*}
即 $a$ 把任一本征矢量变为能量低 $\omega$ 的另一本征矢量。但这一过程不能无限进行下去，
因为能量本征值必须非负。故必存在某态 $\psi_{0}$ 使
\begin{equation*}
a\psi_{0}=0 .
\end{equation*}
这是谐振子的“基态”（ground state）。尽管用了下标 $0$，该基态的能量并不为零，而是
\begin{equation*}
H_{\text{SHO}}\psi_{0}=\omega\bigl(a^{\dagger}a+\tfrac12\bigr)\psi_{0}=\tfrac{\omega}{2}\,\psi_{0}.
\end{equation*}
故 $\psi_{0}$ 的能量为 $\omega/2$，其上还有无穷多个高能本征态，能量为 $\omega$ 的半整数倍：
\begin{equation*}
\psi_{n}=(a^{\dagger})^{n}\psi_{0},\qquad
H_{\text{SHO}}\psi_{n}=\bigl(n+\tfrac12\bigr)\omega\,\psi_{n}.
\end{equation*}
当振子处于能量本征态（即具有确定能量）时，该能量可看作基态能量 $\omega/2$
加上某整数 $n$ 个额外“量子”（quanta），每个能量为 $\omega$。
（若你习惯把量子看作具有能量 $\hbar\omega$，请记住我们用的是 $\hbar=1$ 的单位制。）
算符 $a$ 可视为“湮灭”掉一个量子，因为它把能量本征态变为少一个量子的态；
类似地 $a^{\dagger}$ 可视为“产生”一个量子。因此我们称 $a$ 为\emph{湮灭算符}（annihilation operator），
$a^{\dagger}$ 为\emph{产生算符}（creation operator）。这些名称虽形象生动，却不可按字面理解：
它们\emph{不}描述随时间发生的物理过程，只提供一种把确定量子数的态映射到少一个或多一个量子的态的简便方法。

上述全部分析虽然针对一个作简谐运动的非相对论点粒子，但只要系统的自由度由同一类哈密顿量
（坐标平方与相应动量平方的线性组合）支配，同样的分析便可套用。由于我们将大量使用这些结果，
现将其汇总如下。

\begin{keybox}[简谐振子汇总]
$M$ 为粒子质量，$\omega$ 为振动频率，取 $\hbar=1$ 的单位制。
\begin{align*}
H_{\text{SHO}}&=\frac{p^{2}}{2M}+\tfrac12 M\omega^{2}x^{2}
&&\text{(Hamiltonian)}\\
L_{\text{SHO}}&=\tfrac12 M\bigl(x^{\bullet}{}^{2}-\omega^{2}x^{2}\bigr)
&&\text{(Lagrangian)}\\
\partial_{t}^{2}x&=-\omega^{2}x
&&\text{(Equation of Motion)}\\
a&=\frac{\ii\,(p-\ii M\omega x)}{\sqrt{2M\omega}}
&&\text{(Annihilation Operator)}\\
a^{\dagger}&=\frac{-\ii\,(p+\ii M\omega x)}{\sqrt{2M\omega}}
&&\text{(Creation Operator)}
\end{align*}
用产生、湮灭算符表示的可观测量：
\begin{equation*}
H_{\text{SHO}}=\omega\bigl(a^{\dagger}a+\tfrac12\bigr),\qquad
x=\frac{a+a^{\dagger}}{\sqrt{2M\omega}},\qquad
p=-\ii\sqrt{\frac{M\omega}{2}}\,(a-a^{\dagger});
\end{equation*}
对易子：
\begin{equation*}
[x,p]=\ii,\qquad [a,a^{\dagger}]=1 .
\end{equation*}
\end{keybox}

\subsection{自由黎曼标量场的正则量子化（Canonical Quantisation of a Free Riemannian Scalar Field）}

我们要“量子化”的第一个经典场，由黎曼标量波方程的复值解组成。
事实上，“黎曼标量波方程”（Riemannian Scalar Wave equation, RSW）就是通常所称
Klein-Gordon 方程的黎曼版本；几乎所有量子场论教材都讨论过其洛伦兹版本的量子化，
或针对实值解、或针对复值解。

如前面所论，驯服黎曼波方程、避免出现指数增长解的唯一办法，是假定黎曼宇宙在四个维度上都是有限的：
例如四维环面 $T^{4}$ 或四维球面 $S^{4}$。为简单起见，我们假定一个平坦的四维环面，
但结论不会过分依赖具体的形状或拓扑。我们在有关电磁学的章节中已描述过这类宇宙的若干方面。

自由（即无源）RSW 方程若要有解，四维环面的尺寸 $L_{x},L_{y},L_{z},L_{t}$
与最大频率 $\nu_{\max}$ 之间须有特殊关系，使得存在整数 $n_{x},n_{y},n_{z},n_{t}$ 满足
\begin{equation*}
\Bigl(\frac{n_{x}}{L_{x}}\Bigr)^{2}+\Bigl(\frac{n_{y}}{L_{y}}\Bigr)^{2}
+\Bigl(\frac{n_{z}}{L_{z}}\Bigr)^{2}+\Bigl(\frac{n_{t}}{L_{t}}\Bigr)^{2}=\nu_{\max}^{2}.
\end{equation*}
这允许波沿各维呈现整数个周期，且各频率平方之和等于所需值。对一般的环面尺寸与 $\nu_{\max}$，
此类解不存在；但若尺寸恰好选得使解\emph{存在}，且宇宙相对于最小波长远为巨大，
则会有大量而有限的整数解。

我们不深入计数这类解所需的数论；只假定解数极多，以致对多数目的而言其离散性无法被实验察觉，
并假定相应频率的平面波方向在一切方向上均匀分布。换言之，自由波的可传播四矢量 $\bk$
是 $\mathbb{R}^{4}$ 中一个三维球面上均匀散布的有限离散集合，球半径为 $\omega_{m}$。
我们将利用该集合的离散性来简化某些公式与计算，也会在有用时放心地把对这些矢量的求和近似为积分。

我们称自由波全部可能传播四矢量的集合为 $K_{4}$，并记 $N_{\text{modes}}$ 为该集合中矢量的个数。
在环面上选取某一特定坐标作为我们的时间坐标，并把 $K_{4}$ 中矢量投影到除 $t$ 外的三个维度所得的集合
命名为 $K_{3}$。对 $K_{3}$ 中每个 $\bk$，$K_{4}$ 中有两个矢量投影到它，一个时间分量为正、
一个为负。我们把 $K_{4}$ 中时间分量为正的子集记为 $K_{4}^{+}$、为负的记为 $K_{4}^{-}$，
且为简单起见假定没有时间分量为零的模式。我们把时间分量的绝对值记为 $E_{\bk}$：
\begin{equation}
E_{\bk}=|k^{t}|,\quad \bk\in K_{4};
\qquad
E_{\bk}=\sqrt{\omega_{m}^{2}-|\bk|^{2}}=\sqrt{m^{2}-|\bk|^{2}},\quad \bk\in K_{3}.
\tag{E\_k}
\end{equation}
这里 $m$ 是与该标量场相联系的粒子质量，在我们的单位制下
$m=\hbar\omega_{m}=\omega_{m}=2\pi\nu_{\max}$。

首先，我们把黎曼标量波方程及相应的拉格朗日量与哈密顿量写成两种形式：
第一种以空间、时间坐标的函数 $\phi(\bx,t)$ 表示，第二种以在某些空间周期函数上作模式展开的
（含时）系数 $\phi_{\bk}(t)$ 表示。

\subsubsection*{四维坐标形式}

复值波函数 $\phi(\bx,t)$ 的经典黎曼标量波方程及其拉格朗日量、动量与哈密顿量为
\begin{align}
0&=\partial_{\mu}\partial^{\mu}\phi(\bx,t)+m^{2}\phi(\bx,t)
&&\text{(Equation of Motion)}\\
L_{\text{RSW}}&=\int\bigl(\partial_{\mu}\phi(\bx,t)\bigr)\bigl(\partial^{\mu}\phi^{*}(\bx,t)\bigr)
-m^{2}\phi(\bx,t)\phi^{*}(\bx,t)\,\dd^{3}\bx
&&\text{(Lagrangian)}\\
\Pi(\bx,t)&=\partial_{t}\phi^{*}(\bx,t),\qquad
\Pi^{*}(\bx,t)=\partial_{t}\phi(\bx,t)
&&\text{(Momenta)}\\
H_{\text{RSW}}&=\int\Pi(\bx,t)\Pi^{*}(\bx,t)-\bigl(\partial_{i}\phi(\bx,t)\bigr)\bigl(\partial^{i}\phi^{*}(\bx,t)\bigr)
+m^{2}\phi(\bx,t)\phi^{*}(\bx,t)\,\dd^{3}\bx
&&\text{(Hamiltonian)}
\end{align}
第一组方程中，我们把 $\phi(\bx,t)$ 与其复共轭 $\phi^{*}(\bx,t)$ 当作两个独立场，
而拉格朗日密度 $L_{\text{RSW}}(\bx)$（对空间积分即得总拉格朗日量的量）经由复共轭场的
欧拉-拉格朗日方程直接给出运动方程：
\begin{equation*}
\partial^{\mu}\bigl(\partial_{\partial_{\mu}\phi^{*}(\bx,t)}L_{\text{RSW}}(\bx)\bigr)
=\partial_{\phi^{*}(\bx,t)}L_{\text{RSW}}(\bx)
\;\Longrightarrow\;
\partial^{\mu}\partial_{\mu}\phi(\bx,t)=-m^{2}\phi(\bx,t).
\end{equation*}
与场及其复共轭相联系的动量，定义为拉格朗日密度对相应场时间导数的导数：
\begin{equation*}
\Pi(\bx,t)=\partial_{\partial_{t}\phi(\bx,t)}L_{\text{RSW}}(\bx)=\partial_{t}\phi^{*}(\bx,t),
\qquad
\Pi^{*}(\bx,t)=\partial_{\partial_{t}\phi^{*}(\bx,t)}L_{\text{RSW}}(\bx)=\partial_{t}\phi(\bx,t).
\end{equation*}
哈密顿密度由对两个场 $\phi(\bx,t)$ 及其复共轭求“动量乘场的时间导数”之和，再减去拉格朗日密度得到：
\begin{equation*}
H_{\text{RSW}}(\bx)=\Pi(\bx,t)\partial_{t}\phi(\bx,t)+\Pi^{*}(\bx,t)\partial_{t}\phi^{*}(\bx,t)-L_{\text{RSW}}(\bx)
=2\Pi(\bx,t)\Pi^{*}(\bx,t)-L_{\text{RSW}}(\bx).
\end{equation*}
在上表哈密顿量的公式中，重复指标 $i$ 只对空间坐标求和，时间导数已由动量承担。

\subsubsection*{模式展开形式}

接下来把 $\phi(\bx,t)$ 写成空间周期函数 $f_{\bk}(\bx)$ 的有限和，
对每个 $\bk\in K_{3}$ 有 $f_{\bk}(\bx)\propto\exp(-\ii\,\bk\cdot\bx)$，并称展开系数为 $\phi_{\bk}(t)$。
我们用因子 $\sqrt{1/V}$ 归一化 $f_{\bk}$，其中 $V$ 是四维环面上空间的体积，
$V=L_{x}L_{y}L_{z}$，于是
\begin{equation*}
\int f_{\bk}(\bx)f_{\bq}^{*}(\bx)\,\dd^{3}\bx=\delta_{\bk,\bq}.
\end{equation*}
把模式展开代入 $L_{\text{RSW}}$ 的积分，再利用模式函数的正交性，即得模式系数形式的拉格朗日量：
对每一模式而言，空间导数平方之和恰为 $\bk^{2}$，它与 $m^{2}$ 项合并给出 $E_{\bk}^{2}$。
动量与哈密顿量随后可由拉格朗日量求得。汇总如下：
\begin{align}
f_{\bk}(\bx)&=\sqrt{\tfrac{1}{V}}\exp(-\ii\,\bk\cdot\bx),&
\phi(\bx,t)&=\sum_{\bk\in K_{3}}\phi_{\bk}(t)f_{\bk}(\bx)
&&\text{(Mode Expansion)}\\
\partial_{t}^{2}\phi_{\bk}(t)&=-E_{\bk}^{2}\phi_{\bk}(t)
&&\text{(Equation of Motion)}\\
L_{\text{RSW}}&=\sum_{\bk\in K_{3}}\bigl(\partial_{t}\phi_{\bk}(t)\partial_{t}\phi^{*}_{\bk}(t)
-E_{\bk}^{2}\phi_{\bk}(t)\phi^{*}_{\bk}(t)\bigr)
&&\text{(Lagrangian)}\\
\Pi_{\bk}(t)&=\partial_{t}\phi^{*}_{\bk}(t),\qquad
\Pi^{*}_{\bk}(t)=\partial_{t}\phi_{\bk}(t)
&&\text{(Momenta)}\\
H_{\text{RSW}}&=\sum_{\bk\in K_{3}}\bigl(\Pi_{\bk}(t)\Pi^{*}_{\bk}(t)
+E_{\bk}^{2}\phi_{\bk}(t)\phi^{*}_{\bk}(t)\bigr)
&&\text{(Hamiltonian)}
\end{align}

\subsubsection*{实值自由度}

模式系数版哈密顿量的每一项都\emph{几乎}与简谐振子的哈密顿量形式一致，
但这些自由度 $\phi_{\bk}(t)$、$\phi^{*}_{\bk}(t)$ 及其动量是复数，
而简谐振子的位置与动量是实数。故需把哈密顿量改写为实值自由度的函数。
这并不难：每一项都是一个复数与其共轭之积，等于该复数模的平方，
也就等于实部平方与虚部平方之和（其中我们把虚部当作实数，即除掉因子 $\ii$）。于是可用
\begin{equation*}
\operatorname{Re}_{\bk}(t)=\operatorname{Re}\phi_{\bk}(t),\qquad
\operatorname{Im}_{\bk}(t)=\bigl[\operatorname{Im}\phi_{\bk}(t)\bigr]/\ii
\end{equation*}
重述一切。先在拉格朗日量中做此改写，再求出相应的动量，最终得到\emph{恰为}谐振子形式的哈密顿量，
且每一模式有两个实自由度：
\begin{align}
\phi_{\bk}(t)&=\operatorname{Re}_{\bk}(t)+\ii\,\operatorname{Im}_{\bk}(t)
&&\text{(Real Variables)}\\
\partial_{t}^{2}\operatorname{Re}_{\bk}(t)&=-E_{\bk}^{2}\operatorname{Re}_{\bk}(t),\qquad
\partial_{t}^{2}\operatorname{Im}_{\bk}(t)=-E_{\bk}^{2}\operatorname{Im}_{\bk}(t)
&&\text{(Equations of Motion)}\\
L_{\text{RSW}}&=\sum_{\bk\in K_{3}}\bigl((\partial_{t}\operatorname{Re}_{\bk}(t))^{2}
-E_{\bk}^{2}\operatorname{Re}_{\bk}(t)^{2}
+(\partial_{t}\operatorname{Im}_{\bk}(t))^{2}-E_{\bk}^{2}\operatorname{Im}_{\bk}(t)^{2}\bigr)
&&\text{(Lagrangian)}\\
\Pi_{\operatorname{Re}\bk}(t)&=2\partial_{t}\operatorname{Re}_{\bk}(t),\qquad
\Pi_{\operatorname{Im}\bk}(t)=2\partial_{t}\operatorname{Im}_{\bk}(t)
&&\text{(Momenta)}\\
H_{\text{RSW}}&=\sum_{\bk\in K_{3}}\bigl(\tfrac14\Pi_{\operatorname{Re}\bk}(t)^{2}
+E_{\bk}^{2}\operatorname{Re}_{\bk}(t)^{2}
+\tfrac14\Pi_{\operatorname{Im}\bk}(t)^{2}+E_{\bk}^{2}\operatorname{Im}_{\bk}(t)^{2}\bigr)
&&\text{(Hamiltonian)}
\end{align}
把上表末行哈密顿量的每一项与简谐振子的哈密顿量比对，可见只要把谐振子的粒子质量 $M$ 取为 $2$
（这个 $M$ 与我们用标量场描述的粒子质量无关），把谐振子频率 $\omega$ 对每一模式取为 $E_{\bk}$，
并令谐振子的坐标 $x$ 对应 $\operatorname{Re}_{\bk}$ 或 $\operatorname{Im}_{\bk}$、
动量 $p$ 对应 $\Pi_{\operatorname{Re}\bk}$ 或 $\Pi_{\operatorname{Im}\bk}$，二者便精确吻合。
于是我们对每个 $\bk$ 的实部与虚部自由度都定义产生、湮灭算符。

\begin{keybox}[量子化后的黎曼标量场]
实值自由度 $\operatorname{Re}_{\bk}$ 与 $\operatorname{Im}_{\bk}$，$\bk\in K_{3}$：
\begin{align*}
a_{\operatorname{Re}\bk}&=\frac{\ii}{2\sqrt{E_{\bk}}}
\bigl(\Pi_{\operatorname{Re}\bk}-2\ii E_{\bk}\operatorname{Re}_{\bk}\bigr),&
a_{\operatorname{Im}\bk}&=\frac{\ii}{2\sqrt{E_{\bk}}}
\bigl(\Pi_{\operatorname{Im}\bk}-2\ii E_{\bk}\operatorname{Im}_{\bk}\bigr)\\
a^{\dagger}_{\operatorname{Re}\bk}&=\frac{-\ii}{2\sqrt{E_{\bk}}}
\bigl(\Pi_{\operatorname{Re}\bk}+2\ii E_{\bk}\operatorname{Re}_{\bk}\bigr),&
a^{\dagger}_{\operatorname{Im}\bk}&=\frac{-\ii}{2\sqrt{E_{\bk}}}
\bigl(\Pi_{\operatorname{Im}\bk}+2\ii E_{\bk}\operatorname{Im}_{\bk}\bigr)
\end{align*}
用产生、湮灭算符表示的可观测量与坐标：
\begin{align*}
H_{\text{RSW}}&=\sum_{\bk\in K_{3}}E_{\bk}
\bigl(a^{\dagger}_{\operatorname{Re}\bk}a_{\operatorname{Re}\bk}
+a^{\dagger}_{\operatorname{Im}\bk}a_{\operatorname{Im}\bk}+1\bigr),\\
\operatorname{Re}_{\bk}&=\frac{a_{\operatorname{Re}\bk}+a^{\dagger}_{\operatorname{Re}\bk}}{2\sqrt{E_{\bk}}},&
\Pi_{\operatorname{Re}\bk}&=-\ii\sqrt{E_{\bk}}\bigl(a_{\operatorname{Re}\bk}-a^{\dagger}_{\operatorname{Re}\bk}\bigr),\\
\operatorname{Im}_{\bk}&=\frac{a_{\operatorname{Im}\bk}+a^{\dagger}_{\operatorname{Im}\bk}}{2\sqrt{E_{\bk}}},&
\Pi_{\operatorname{Im}\bk}&=-\ii\sqrt{E_{\bk}}\bigl(a_{\operatorname{Im}\bk}-a^{\dagger}_{\operatorname{Im}\bk}\bigr);
\end{align*}
对易子：
\begin{equation*}
[\operatorname{Re}_{\bk},\Pi_{\operatorname{Re}\bq}]=\ii\delta_{\bk,\bq},\quad
[a_{\operatorname{Re}\bk},a^{\dagger}_{\operatorname{Re}\bq}]=\delta_{\bk,\bq},\quad
[\operatorname{Im}_{\bk},\Pi_{\operatorname{Im}\bq}]=\ii\delta_{\bk,\bq},\quad
[a_{\operatorname{Im}\bk},a^{\dagger}_{\operatorname{Im}\bq}]=\delta_{\bk,\bq}.
\end{equation*}
\end{keybox}

这里的产生、湮灭算符与简谐振子的作用方式完全相同：产生算符
$a^{\dagger}_{\operatorname{Re}\bk}$、$a^{\dagger}_{\operatorname{Im}\bk}$
把能量为 $E$ 的哈密顿量本征态变为能量为 $E+E_{\bk}$ 的新本征态，
而湮灭算符 $a_{\operatorname{Re}\bk}$、$a_{\operatorname{Im}\bk}$
把能量为 $E$ 的本征态变为能量为 $E-E_{\bk}$ 的新本征态。
我们定义系统的基态 $\psi_{0}$ 为满足
\begin{equation*}
a_{\operatorname{Re}\bk}\psi_{0}=0,\qquad
a_{\operatorname{Im}\bk}\psi_{0}=0
\qquad(\text{对一切 }\bk\in K_{3})
\end{equation*}
的矢量。把这些方程与哈密顿量结合，可见 $\psi_{0}$ 的能量为
\begin{equation}
E_{0}=\sum_{\bk\in K_{3}}E_{\bk}
\approx\Bigl[\frac{\tfrac12 N_{\text{modes}}}{\pi^{2}}\Bigr]
\int_{\gamma=0}^{\pi/2}m\cos\gamma\;4\pi\sin^{2}\gamma\,\dd\gamma
=\frac{2N_{\text{modes}}m}{3\pi},
\tag{E\_0}
\end{equation}
其中近似是把 $K_{3}$ 中的 $\tfrac12 N_{\text{modes}}$ 个矢量当作半径 $m$ 的三维球面 $t>0$ 半球上
均匀分布矢量连续统的投影。洛伦兹情形下 $E_{0}$ 为无限大，需要把它减去、重新定义能量零点；
这听起来可疑，但由于量子力学中多数量子化问题里真正重要的只是能量\emph{差}，所以效果尚可。
在我们的情形中并不真的需要这么做；无论描述量子态能量时是否计入 $E_{0}$，
它取的具体数值都会以自身的形式出现在某些公式中。

场中各个模式——对应 $K_{3}$ 中的各个矢量 $\bk$ 及实部、虚部——的作用恰似完全独立的谐振子。
量子力学中，独立系统由各系统希尔伯特空间的张量积所描述的希尔伯特空间来刻画，
故可把本场的希尔伯特空间视为 $N_{\text{modes}}$ 份谐振子希尔伯特空间的张量积。
所有带 $\operatorname{Re}\bk$ 或 $\operatorname{Im}\bk$ 指标的算符都只作用在其特定的因子上，
而在其余因子上充当恒等算符。任何模式都不影响其它模式的演化。

这当然极不现实，因为我们最终要处理的电磁场与 Dirac 型场彼此耦合。
但处理一个无源场，是攻克更难问题前有用的热身练习。

在找到可按谐振子类比量子化的自由度之后，我们现把一切译回更易作物理解释的原始复值自由度。
实值与复值自由度之间的关系此前已定；这虽不完全告诉我们如何为复值模式定义产生、湮灭算符，
但可以检验以下选择确实给出正确的对易子并把哈密顿量化为正确形式。
（虽然直接取 $a_{\bk}=(1/\sqrt2)(a_{\operatorname{Re}\bk}+\ii\,a_{\operatorname{Im}\bk})$
与 $b_{\bk}=(1/\sqrt2)(a_{\operatorname{Re}\bk}-\ii\,a_{\operatorname{Im}\bk})$
也许更简洁，但下面诸式的右端额外乘了一个 $-\ii$ 的相位因子，
这是为了与 Weinberg 的场变量约定一致，便于对照。）

\begin{keybox}[量子化后的黎曼标量场（复值自由度）]
复值自由度 $\phi_{\bk}$ 与 $\phi^{\dagger}_{\bk}$，$\bk\in K_{3}$：
\begin{align*}
a_{\bk}&=\frac{1}{\sqrt2}\bigl(a_{\operatorname{Im}\bk}-\ii\,a_{\operatorname{Re}\bk}\bigr)
=\frac{1}{\sqrt{2E_{\bk}}}\bigl(\Pi^{\dagger}_{\bk}-\ii E_{\bk}\phi_{\bk}\bigr),\\
b_{\bk}&=-\frac{1}{\sqrt2}\bigl(a_{\operatorname{Im}\bk}+\ii\,a_{\operatorname{Re}\bk}\bigr)
=\frac{1}{\sqrt{2E_{\bk}}}\bigl(\Pi_{\bk}-\ii E_{\bk}\phi^{\dagger}_{\bk}\bigr),\\
a^{\dagger}_{\bk}&=\frac{1}{\sqrt2}\bigl(a^{\dagger}_{\operatorname{Im}\bk}+\ii\,a^{\dagger}_{\operatorname{Re}\bk}\bigr)
=\frac{1}{\sqrt{2E_{\bk}}}\bigl(\Pi_{\bk}+\ii E_{\bk}\phi^{\dagger}_{\bk}\bigr),\\
b^{\dagger}_{\bk}&=-\frac{1}{\sqrt2}\bigl(a^{\dagger}_{\operatorname{Im}\bk}-\ii\,a^{\dagger}_{\operatorname{Re}\bk}\bigr)
=\frac{1}{\sqrt{2E_{\bk}}}\bigl(\Pi^{\dagger}_{\bk}+\ii E_{\bk}\phi_{\bk}\bigr)
\end{align*}
可观测量（哈密顿量与坐标）：
\begin{align*}
H_{\text{RSW}}&=\sum_{\bk\in K_{3}}E_{\bk}
\bigl(a^{\dagger}_{\bk}a_{\bk}+b^{\dagger}_{\bk}b_{\bk}+1\bigr),\\
\phi_{\bk}&=\frac{\ii\,(a_{\bk}-b^{\dagger}_{\bk})}{\sqrt{2E_{\bk}}},&
\phi^{\dagger}_{\bk}&=\frac{-\ii\,(a^{\dagger}_{\bk}-b_{\bk})}{\sqrt{2E_{\bk}}},\\
\Pi_{\bk}&=\sqrt{\frac{E_{\bk}}{2}}\bigl(a^{\dagger}_{\bk}+b_{\bk}\bigr),&
\Pi^{\dagger}_{\bk}&=\sqrt{\frac{E_{\bk}}{2}}\bigl(a_{\bk}+b^{\dagger}_{\bk}\bigr);
\end{align*}
对易子：
\begin{equation*}
[\phi_{\bk},\Pi_{\bq}]=\ii\delta_{\bk,\bq},\quad
[a_{\bk},a^{\dagger}_{\bq}]=\delta_{\bk,\bq},\quad
[\phi^{\dagger}_{\bk},\Pi^{\dagger}_{\bq}]=\ii\delta_{\bk,\bq},\quad
[b_{\bk},b^{\dagger}_{\bq}]=\delta_{\bk,\bq}.
\end{equation*}
\end{keybox}

此时 $\phi_{\bk}$ 与 $\Pi_{\bk}$ 是我们希尔伯特空间上的\emph{非厄米算符}，
而不再是一般的含时复值函数；与这些函数的复共轭相联系的算符即厄米共轭 $\phi^{\dagger}_{\bk}$
与 $\Pi^{\dagger}_{\bk}$。我们通常认为可观测量是厄米算符，但这只在它们测量实值量时成立；
对这些为经典场引入的复值模式系数而言，非厄米算符完全可以是合适的可观测量。

我们现以名为 $a_{\bk}$、$b_{\bk}$ 的湮灭算符及其厄米共轭（产生算符）取代原先两种实值自由度的算符。
哈密顿量用这些新算符表示时，与用旧算符时形式完全相同，故同样会有能量为 $E_{0}$ 的基态 $\psi_{0}$。
与先前一样，用任一产生算符 $a^{\dagger}_{\bk}$ 或 $b^{\dagger}_{\bk}$ 作用于能量本征态，
都给出能量增加 $E_{\bk}$ 的新本征态。

\subsubsection{守恒荷、粒子与反粒子（Conserved Charge, Particles and Antiparticles）}

用 $a^{\dagger}_{\bk}$ 或 $b^{\dagger}_{\bk}$ 往基态上添加一个能量量子，所得到的态有何物理含义？
为回答此问题，需回看以复模式系数表示的经典场拉格朗日量：
\begin{equation*}
L_{\text{RSW}}=\sum_{\bk\in K_{3}}\bigl(\partial_{t}\phi_{\bk}(t)\partial_{t}\phi^{*}_{\bk}(t)
-E_{\bk}^{2}\phi_{\bk}(t)\phi^{*}_{\bk}(t)\bigr).
\end{equation*}
若把每个 $\phi_{\bk}$ 乘以相位 $\exp(-\ii\alpha)$，则每个 $\phi^{*}_{\bk}$ 乘以 $\exp(\ii\alpha)$，
拉格朗日量本身不变。因此，对整个场作均匀相旋转是\emph{拉格朗日量的对称性}。
任何此类对称性都给出一个守恒量；适用于离散变量的诺特定理（Noether's theorem）告诉我们该量为
\begin{align*}
C&=\sum_{\bk\in K_{3}}\Bigl[
\bigl(\partial_{\partial_{t}\phi_{\bk}}L_{\text{RSW}}\bigr)
\bigl(\partial_{\alpha}\exp(-\ii\alpha)\phi_{\bk}(t)\bigr)\big|_{\alpha=0}
+\bigl(\partial_{\partial_{t}\phi^{*}_{\bk}}L_{\text{RSW}}\bigr)
\bigl(\partial_{\alpha}\exp(\ii\alpha)\phi^{*}_{\bk}(t)\bigr)\big|_{\alpha=0}\Bigr]\\
&=-\ii\sum_{\bk\in K_{3}}\bigl(\partial_{t}\phi^{*}_{\bk}(t)\,\phi_{\bk}(t)
-\partial_{t}\phi_{\bk}(t)\,\phi^{*}_{\bk}(t)\bigr)
=-\ii\sum_{\bk\in K_{3}}\bigl(\Pi_{\bk}(t)\phi_{\bk}(t)-\Pi^{*}_{\bk}(t)\phi^{*}_{\bk}(t)\bigr).
\end{align*}
把 $C$ 译成用产生、湮灭算符表示的算符，得
\begin{equation*}
C=\sum_{\bk\in K_{3}}\bigl(a^{\dagger}_{\bk}a_{\bk}-b^{\dagger}_{\bk}b_{\bk}\bigr).
\end{equation*}
用与求能量本征态十分类似的论证容易证明，$a^{\dagger}_{\bk}a_{\bk}$ 计数以 $a^{\dagger}_{\bk}$
添加到基态上的量子数，即
\begin{equation*}
(a^{\dagger}_{\bk}a_{\bk})\psi_{0}=0,\qquad
(a^{\dagger}_{\bk}a_{\bk})(a^{\dagger}_{\bk})^{n}\psi_{0}=n\,(a^{\dagger}_{\bk})^{n}\psi_{0}.
\end{equation*}
类似地，$b^{\dagger}_{\bk}b_{\bk}$ 计数以 $b^{\dagger}_{\bk}$ 添加到基态上的量子数：
\begin{equation*}
(b^{\dagger}_{\bk}b_{\bk})(b^{\dagger}_{\bk})^{n}\psi_{0}=n\,(b^{\dagger}_{\bk})^{n}\psi_{0}.
\end{equation*}
二者之差对所有 $\bk$ 求和，即是以 $a^{\dagger}$ 算符产生的量子总数减去以 $b^{\dagger}$ 算符产生的量子总数。
该数守恒；要检验它在量子版本中成立，只需验证 $C$ 与哈密顿量对易。
（事实上两个数算符各自都守恒，但那是因为我们处理的是几乎不发生什么事的自由场论。
数算符之差源于 $\phi$ 的均匀相位变化，即使在相互作用场论中也守恒——只要相互作用项同时含
$\phi$ 与 $\phi^{*}$ 的因子。）

这告诉我们，$a^{\dagger}_{\bk}$ 算符产生粒子，而 $b^{\dagger}_{\bk}$ 算符产生相应的\emph{反粒子}
（antiparticles），其含义是：这两类量子携带的“荷”符号相反，且其净值随时间保持不变。
若 $\phi$ 与电磁场耦合，此荷便是\emph{电}荷；但此刻它只是一个尚无更广含义的守恒量子数。
对最一般的能量本征态 $\psi$，有
\begin{align*}
\psi&=(a^{\dagger}_{\bk_{1}})^{q_{1}}(a^{\dagger}_{\bk_{2}})^{q_{2}}(a^{\dagger}_{\bk_{3}})^{q_{3}}\cdots
(b^{\dagger}_{\bk_{1}})^{n_{1}}(b^{\dagger}_{\bk_{2}})^{n_{2}}(b^{\dagger}_{\bk_{3}})^{n_{3}}\cdots\psi_{0},\\
H\psi&=\bigl(E_{0}+(q_{1}+n_{1})E_{\bk_{1}}+(q_{2}+n_{2})E_{\bk_{2}}
+(q_{3}+n_{3})E_{\bk_{3}}+\cdots\bigr)\psi,\\
C\psi&=\bigl((q_{1}+q_{2}+q_{3}+\cdots)-(n_{1}+n_{2}+n_{3}+\cdots)\bigr)\psi .
\end{align*}
回过头看，现在可以明白算符 $a^{\dagger}_{\operatorname{Re}\bk}$ 与 $a^{\dagger}_{\operatorname{Im}\bk}$
所产生的是何种量子：它们是粒子态与反粒子态的等量组合，净荷为零！

在洛伦兹宇宙中，相对运动的观测者对粒子与反粒子的分类会取得一致，但黎曼宇宙并非如此。
设 $\bk$ 是一个三矢量，使 $E_{\bk}$ 相当小；这意味着在测量这些量的参考系中，
四矢量 $\bk=(E_{\bk},\bk)$ 指向几乎纯空间的方向，故稍有一点相对运动，
另一观测者便会使其时间分量变为负。

我们可以通过定义以 $K_{4}$ 中四矢量为指标、较少依赖参考系的湮灭算符来表述：
\begin{equation*}
c_{(E,\bk)}=
\begin{cases}
a_{\bk}, & E>0,\\
b_{-\bk}, & E<0,
\end{cases}
\qquad
c^{\dagger}_{(E,\bk)}=
\begin{cases}
a^{\dagger}_{\bk}, & E>0,\\
b^{\dagger}_{-\bk}, & E<0,
\end{cases}
\qquad
\langle\psi_{0},\psi_{0}\rangle=1 .
\end{equation*}

那么，态 $\psi_{\bq}=c^{\dagger}_{\bq}\psi_{0}$（对某模式 $\bq\in K_{4}$ 往场中放入一个量子）
的归一性如何？计算
\begin{align*}
\langle\psi_{\bq},\psi_{\bq}\rangle
&=\langle c^{\dagger}_{\bq}\psi_{0},c^{\dagger}_{\bq}\psi_{0}\rangle
=\langle\psi_{0},c_{\bq}c^{\dagger}_{\bq}\psi_{0}\rangle\\
&=\langle\psi_{0},(c^{\dagger}_{\bq}c_{\bq}+1)\psi_{0}\rangle
=\langle\psi_{0},\psi_{0}\rangle=1,
\end{align*}
其中用到了伴随与内积的一般性质 $\langle A^{\dagger}u,v\rangle=\langle u,Av\rangle$、
湮灭与产生算符的对易子，以及任何湮灭算符作用于基态都为零这一事实。

此刻我们先略过如下细节：这些态（薛定谔绘景）或算符（海森堡绘景）应被视为随时间变化；
在我们所考察的情形中，时间变化只相当于一个振荡相位，对期望值没有影响。

态 $\psi_{\bq}$ 下 $\phi_{\bk}$ 的期望值为
\begin{align*}
\langle\psi_{\bq},\phi_{\bk}\psi_{\bq}\rangle
&=\langle c^{\dagger}_{\bq}\psi_{0},\phi_{\bk}c^{\dagger}_{\bq}\psi_{0}\rangle
=\langle\psi_{0},c_{\bq}\phi_{\bk}c^{\dagger}_{\bq}\psi_{0}\rangle\\
&=\frac{\ii}{\sqrt{2E_{\bk}}}\langle\psi_{0},c_{\bq}(c_{\bk}-c^{\dagger}_{-\bk})c^{\dagger}_{\bq}\psi_{0}\rangle=0 .
\end{align*}
这必然为零，因为任何\emph{奇数}个产生、湮灭算符都无法把基态带回基态。
此外，这个期望值为零也不足为奇，因为它有点像问简谐振子中粒子的平均位置，自然应是 $x=0$。
但位置的\emph{平方}应有非零平均值；对本复值场振幅而言，
即 $\phi^{\dagger}_{\bk}\phi_{\bk}$（也可写作 $\phi_{-\bk}\phi_{\bk}$）的期望值应非零。

先对基态计算。注意 $\phi_{\bk}$ 与 $\phi^{\dagger}_{\bk}$ 对易，故书写次序无关：
\begin{align*}
\langle\psi_{0},\phi_{-\bk}\phi_{\bk}\psi_{0}\rangle
&=\Bigl\langle\psi_{0},
\Bigl(\tfrac{\ii(c_{-\bk}-c^{\dagger}_{\bk})}{\sqrt{2E_{\bk}}}\Bigr)
\Bigl(\tfrac{\ii(c_{\bk}-c^{\dagger}_{-\bk})}{\sqrt{2E_{\bk}}}\Bigr)\psi_{0}\Bigr\rangle\\
&=-\frac{1}{2E_{\bk}}\langle\psi_{0},(c_{-\bk}-c^{\dagger}_{\bk})(c_{\bk}-c^{\dagger}_{-\bk})\psi_{0}\rangle
=\frac{1}{2E_{\bk}} .
\end{align*}
再对态 $\psi_{\bq}$ 计算：
\begin{align*}
\langle\psi_{\bq},\phi_{-\bk}\phi_{\bk}\psi_{\bq}\rangle
&=\langle\psi_{0},c_{\bq}\phi_{-\bk}\phi_{\bk}c^{\dagger}_{\bq}\psi_{0}\rangle\\
&=-\frac{1}{2E_{\bk}}\langle\psi_{0},
c_{\bq}(c_{-\bk}-c^{\dagger}_{\bk})(c_{\bk}-c^{\dagger}_{-\bk})c^{\dagger}_{\bq}\psi_{0}\rangle
=\frac{1+\delta_{\bq,\bk}+\delta_{\bq,-\bk}}{2E_{\bk}} .
\end{align*}
与模式相联系的动量平方，其期望值与上述只差因子 $E_{\bk}^{2}$：
\begin{equation*}
\langle\psi_{0},\Pi_{-\bk}\Pi_{\bk}\psi_{0}\rangle=\frac{E_{\bk}}{2},\qquad
\langle\psi_{\bq},\Pi_{-\bk}\Pi_{\bk}\psi_{\bq}\rangle
=\frac{E_{\bk}}{2}\bigl(1+\delta_{\bq,\bk}+\delta_{\bq,-\bk}\bigr).
\end{equation*}
这里我们看到，这些平方振幅与平方动量可观测量对基态能量敏感：
即使场在该可观测量所调谐的模式中\emph{没有}量子，它们的期望值也非零。
而诸如 $\psi_{\bq}$ 这样的能量本征态高度非经典：能量被锐确定，代价是该模式的相位完全不确定。

虽不展开细节，但值得简述半经典态（semi-classical state）如何构造。
对简谐振子，可取基态波函数——一个高斯波包——并把它中心任意偏离能量谷中心。
平移后的波包含无穷多个能量本征函数的成分，但随其随时间演化，波包只是在谷中来回运动而不改变形状，
极像经典粒子。我们对复标量场的某一模式也可做类似的事；若改用实部与虚部，便易于模仿谐振子的构造。
通过在实部与虚部设置相位相差 $90^{\circ}$ 振荡的高斯波包，结果是一个在复平面中旋转的复模式振幅：
一个以恒定距离绕原点旋转的二维高斯。振幅平方足够大时，该态将酷似相位良好的经典波。

\subsubsection{定域态与可观测量（Localised States and Observables）}

在洛伦兹量子场论中，本学科的常规发展包括把场 $\phi$、其厄米共轭及其共轭动量写成空间坐标的算符值函数，
使用原始的空间模式与量子化后的模式系数算符。黎曼情形同样可行，但结果须谨慎对待——
因为这样定义的 $\phi(\bx)$ 会涉及对 $K_{3}$ 的求和，从而依赖于观测者对 $t$ 坐标的选取。
洛伦兹对应物是正质量双曲面（满足 $\bk\cdot\bk=-m^{2}$ 且 $k^{t}>0$ 的类时矢量集合）上一切传播矢量的积分，
它只与时间坐标整体\emph{符号}的选取有关，且在相对运动的观测者之间变换合理。
而黎曼场 $\phi(\bx)$ 却被锁定于特定的时间\emph{轴}，因为即使时间轴只有微小差异的观测者，
也会对何为粒子、何为反粒子意见不一。

黎曼情形还有一个麻烦，甚至在经典版本中也可见。理想情况下，我们希望空间各点的场与模式系数之间的变量变换
是正则变换（canonical transformation）。哈密顿力学中，泊松括号（Poisson bracket）定义为作用于函数对的微分算符：
\begin{equation*}
\{f,g\}_{\text{PB}}=(\partial_{q_{j}}f)(\partial_{p_{j}}g)-(\partial_{p_{j}}f)(\partial_{q_{j}}g),
\end{equation*}
其中对重复指标 $j$ 求和。坐标 $q_{i}$ 与动量 $p_{i}$ 满足泊松括号
\begin{equation*}
\{q_{i},q_{k}\}_{\text{PB}}=0,\qquad
\{p_{i},p_{k}\}_{\text{PB}}=0,\qquad
\{q_{i},p_{k}\}_{\text{PB}}=\delta_{ik}.
\end{equation*}
对场值这类连续自由度，离散指标被空间坐标取代、求和被积分取代、克罗内克 $\delta$ 被狄拉克 $\delta$ 函数取代。
在变量间的正则变换中，\emph{新}坐标与动量在泊松括号用\emph{原}变量定义时也应满足这些关系。
具体地，若新坐标为 $Q_{i}$、新动量为 $P_{i}$，则应有
\begin{equation*}
\{Q_{i},Q_{k}\}_{\text{PB}}=0,\qquad
\{P_{i},P_{k}\}_{\text{PB}}=0,\qquad
\{Q_{i},P_{k}\}_{\text{PB}}=\delta_{ik}.
\end{equation*}
当一组变量连续、另一组离散时情况稍棘手，但有时仍完美奏效。
例如，若我们按仅对空间满足周期边界条件的可数无穷傅里叶模式展开，
配以适当归一化，从空间坐标变为模式系数时上述方程组中出现克罗内克 $\delta$
（对指标 $j$ 的求和换为对空间坐标的积分），反向变换时则出现狄拉克 $\delta$ 函数。

但在本情形中我们会得到什么？新坐标、动量用旧变量表示为
\begin{equation*}
\phi(\bx,t)=\sum_{\bk\in K_{3}}\phi_{\bk}(t)f_{\bk}(\bx),\qquad
\Pi(\bx,t)=\partial_{t}\phi^{*}(\bx,t)
=\sum_{\bk\in K_{3}}\Pi_{\bk}(t)f^{*}_{\bk}(\bx),
\end{equation*}
于是新旧变量间的偏导数为（为简洁起见此后略去显式的含时性）
\begin{equation*}
\partial_{\phi_{\bk}}\phi(\bx)=f_{\bk}(\bx),\qquad
\partial_{\Pi_{\bk}}\Pi(\bx)=f^{*}_{\bk}(\bx),
\end{equation*}
坐标与动量间的泊松括号则为
\begin{align*}
\{\phi(\bx),\Pi(\by)\}_{\text{PB}}
&=\sum_{\bk\in K_{3}}\bigl[(\partial_{\phi_{\bk}}\phi(\bx))(\partial_{\Pi_{\bk}}\Pi(\by))
-(\partial_{\Pi_{\bk}}\phi(\bx))(\partial_{\phi_{\bk}}\Pi(\by))\bigr]\\
&=\sum_{\bk\in K_{3}}f_{\bk}(\bx)f^{*}_{\bk}(\by)
=\frac{1}{V}\sum_{\bk\in K_{3}}\exp\bigl(\ii\,\bk\cdot(\by-\bx)\bigr)\\
&=\frac{1}{V}\sum_{\bk\in K_{4}^{+}}\exp\bigl(\ii\,\bk\cdot(\by-\bx)\bigr)
\approx\frac{N_{\text{modes}}}{2\pi^{2}V}\int_{\bu\in S^{3+}}\exp\bigl(\ii m\,\bu\cdot(\by-\bx)\bigr),
\end{align*}
其中倒数第二行用了 $\bx$、$\by$ 纯属空间向量，故与 $K_{3}$ 中某 $\bk$ 或投影到它的 $K_{4}^{+}$ 中 $\bk$
作点积并无差别；最后一行把对 $K_{4}^{+}$ 中矢量的求和近似为对 $\mathbb{R}^{4}$ 中单位三维球面正半部分的积分。
我们可以利用该积分对三空间旋转的不变性，选取坐标使 $\by-\bx=|\bx-\by|\be_{z}$。
在 $S^{3+}$ 上的极坐标 $(\gamma,\theta,\phi)$ 中（$0\le\gamma\le\pi/2$ 等），有
\begin{equation*}
\bu=\bigl(\cos\gamma,\ \sin\gamma\sin\theta\cos\phi,\ \sin\gamma\sin\theta\sin\phi,\ \sin\gamma\cos\theta\bigr),
\end{equation*}
故所需点积为
\begin{equation*}
\bu\cdot(|\bx-\by|\be_{z})=|\bx-\by|\sin\gamma\cos\theta .
\end{equation*}
由于被积函数与方位角坐标 $\phi$（注意勿与场振幅混淆）无关，可立即把它积掉，余下
\begin{align*}
\{\phi(\bx),\Pi(\by)\}_{\text{PB}}
&\approx\frac{N_{\text{modes}}}{\pi V}
\int_{\theta=0}^{\pi}\int_{\gamma=0}^{\pi/2}
\exp\bigl(\ii m|\bx-\by|\sin\gamma\cos\theta\bigr)\sin^{2}\gamma\sin\theta\,\dd\gamma\,\dd\theta\\
&=\frac{2N_{\text{modes}}}{\pi V}\frac{1}{m|\bx-\by|}
\int_{\gamma=0}^{\pi/2}\sin\bigl(m|\bx-\by|\sin\gamma\bigr)\sin\gamma\,\dd\gamma\\
&=\frac{N_{\text{modes}}}{V}\frac{J_{1}(m|\bx-\by|)}{m|\bx-\by|},
\end{align*}
这里 $J_{1}$ 是第一类贝塞尔函数（Bessel function）。$J_{1}$ 大致呈周期振荡，在原点为零；
但除以自变量后得到的函数在原点取峰、随后在振荡中衰减。其第一个零点出现在 $m|\bx-\by|\approx 4$；
回忆 $m=\omega_{m}$，可见峰的宽度与 RSW 方程相关的最小波长为同一量级。
然而尽管该函数有峰，它显然不是狄拉克 $\delta$ 函数，因此我们\emph{不}拥有变量的正则变换。

问题在于：我们可用的模式无法组合成完全定域的场。这与量子力学无关，也不是因为模式集合有限；
即便假装拥有一个传播矢量铺满三维球面半边的连续模式集合，这些模式仍有最小波长，
故不能指望由它们构造出任意的窄峰函数。

因此，与其试图处理空间中每一点的场值，不如为模式函数 $f_{\bk}(\bx)$ 张成的空间寻找一个新基，
由在该空间内尽可能定域的新函数组成。定义一族函数 $g_{\bchi}(\bx)$：
\begin{equation*}
g_{\bchi}(\bx)=\sum_{\bk\in K_{3}}\sqrt{\frac{E_{\bk}}{E_{0}}}\,f_{\bk}(\bx-\bchi)
=\sum_{\bk\in K_{3}}\sqrt{\frac{E_{\bk}}{E_{0}}}\,\exp(\ii\,\bk\cdot\bchi)f_{\bk}(\bx),
\end{equation*}
其中 $\bchi$ 暂时是任意三矢量。任何 $g_{\bchi}$ 都是模式函数 $f_{\bk}$ 的线性组合，
故与狄拉克 $\delta$ 函数不同，它落在我们真能构造的函数子空间内。可用积分近似 $g_{\bzero}$：
\begin{align*}
g_{\bzero}(\bx)&=\sum_{\bk\in K_{3}}\sqrt{\frac{E_{\bk}}{E_{0}}}f_{\bk}(\bx)
\approx\frac{N_{\text{modes}}}{2\pi^{2}\sqrt{VE_{0}}}\int_{\bu\in S^{3+}}\exp(-\ii m\,\bu\cdot\bx)\sqrt{m\,u^{t}}\\
&=\sqrt{\frac{3N_{\text{modes}}}{2\pi V}}\int_{\theta=0}^{\pi}\int_{\gamma=0}^{\pi/2}
\exp\bigl(-\ii m|\bx|\sin\gamma\cos\theta\bigr)\sin^{2}\gamma\sqrt{\cos\gamma}\sin\theta\,\dd\gamma\,\dd\theta\\
&=\sqrt{\frac{6N_{\text{modes}}}{\pi V}}\frac{1}{m|\bx|}
\int_{\gamma=0}^{\pi/2}\sin\bigl(m|\bx|\sin\gamma\bigr)\sin\gamma\sqrt{\cos\gamma}\,\dd\gamma\\
&=\sqrt{\frac{3N_{\text{modes}}}{V}}\Bigl[\frac{\Gamma(3/4)}{2^{1/4}}\Bigr]
\frac{J_{5/4}(m|\bx|)}{(m|\bx|)^{5/4}} .
\end{align*}

\orthofig{images/07__010.png}{空间定域的波（Spatially localised wave）}

该贝塞尔函数除以其自变量的幂次，形状与我们为泊松括号所求的函数非常相似，
峰的宽度与最小波长同量级。因此它是一个相当定域的函数，
我们不应指望用任何模式线性组合做得更好。

至于 $g_{\bchi}(\bx)=g_{\bzero}(\bx-\bchi)$ 只是把这个峰从原点平移到另一位置，
故所有此类函数在空间中同样定域。现假定可选取 $\tfrac12 N_{\text{modes}}$ 个不同的点 $\bchi$，
使所有这些 $g_{\bchi}$ 互相正交；我们把这一组点的选取记为 $\Chi$。
若把这些函数归一化，它们便构成 $f_{\bk}$ 张成函数空间的另一个正交基。

若该假定不成立又如何？我们确实知道，若选取\emph{任何} $\tfrac12 N_{\text{modes}}$ 个互异的点 $\bchi$，
$g_{\bchi}$ 仍给出函数空间的一个基。随后可用 Gram-Schmidt 正交化过程把它化为正交基。
所得函数不再全部取 $g_{\bchi}(\bx)$ 的形式，但若把点均匀散布于空间，结果仍应是一组与 $g_{\bchi}$
定性上非常相似的定域函数。

由于 $f_{\bk}$ 确实正交，$g_{\bchi}$ 的正交性假定等价于
\begin{equation*}
\langle g_{\bchi_{1}},g_{\bchi_{2}}\rangle
=\sum_{\bk\in K_{3}}\frac{E_{\bk}}{E_{0}}\exp\bigl(\ii\,\bk\cdot(\bchi_{2}-\bchi_{1})\bigr)
=\delta_{\bchi_{1},\bchi_{2}}\Bigl[\sum_{\bk\in K_{3}}\frac{E_{\bk}}{E_{0}}\Bigr]
=\delta_{\bchi_{1},\bchi_{2}} .
\end{equation*}
我们以 $\phi_{\bchi}$ 表示场中 $g_{\bchi}(\bx)$ 的系数，替代任意点处的场 $\phi(\bx)$：
\begin{equation*}
\sum_{\bk\in K_{3}}\phi_{\bk}f_{\bk}(\bx)=\sum_{\bchi\in\Chi}\phi_{\bchi}g_{\bchi}(\bx).
\end{equation*}
利用 $g_{\bchi}$ 的正交归一性，用相应 $g_{\bchi}$ 取内积便可提取一个系数：
\begin{equation*}
\phi_{\bchi}=\sum_{\bk\in K_{3}}\sqrt{\frac{E_{\bk}}{E_{0}}}\exp(-\ii\,\bk\cdot\bchi)\phi_{\bk},
\end{equation*}
逆公式为
\begin{equation*}
\phi_{\bk}=\sum_{\bchi\in\Chi}\sqrt{\frac{E_{\bk}}{E_{0}}}\exp(\ii\,\bk\cdot\bchi)\phi_{\bchi}.
\end{equation*}
由这些公式与正交性假定不难证明，经典变量从模式系数 $\phi_{\bk}$ 到空间系数 $\phi_{\bchi}$ 的变换是正则变换：
\begin{equation*}
\{\phi_{\bchi_{1}},\Pi_{\bchi_{2}}\}_{\text{PB}}=\delta_{\bchi_{1},\bchi_{2}} .
\end{equation*}
现在，只需使用 $\phi_{\bk}$ 的算符形式，便可写出这些新变量对应的算符。

\begin{keybox}[量子化后的黎曼标量场（薛定谔绘景）]
$\bchi$ 是取自集合 $\Chi$ 的空间三矢量。算符不随时间变化，作用于含时态矢量。
\begin{align*}
\phi_{\bchi}&=\sum_{\bk\in K_{3}}\sqrt{\frac{E_{\bk}}{E_{0}}}\exp(-\ii\,\bk\cdot\bchi)\phi_{\bk}
=\frac{\ii}{\sqrt{2E_{0}}}\sum_{\bk\in K_{4}^{+}}\bigl(c_{\bk}-c^{\dagger}_{-\bk}\bigr)
\exp(-\ii\,\bk\cdot\bchi),\\
\phi^{\dagger}_{\bchi}&=\sum_{\bk\in K_{3}}\sqrt{\frac{E_{\bk}}{E_{0}}}\exp(\ii\,\bk\cdot\bchi)\phi^{\dagger}_{\bk}
=\frac{\ii}{\sqrt{2E_{0}}}\sum_{\bk\in K_{4}^{-}}\bigl(c_{\bk}-c^{\dagger}_{-\bk}\bigr)
\exp(-\ii\,\bk\cdot\bchi);
\end{align*}
平面波模式系数的逆公式与正交性假定：
\begin{equation*}
\phi_{\bk}=\sum_{\bchi\in\Chi}\sqrt{\frac{E_{\bk}}{E_{0}}}\exp(\ii\,\bk\cdot\bchi)\phi_{\bchi},
\qquad
\sum_{\bk\in K_{3}}E_{\bk}\exp\bigl(\ii\,\bk\cdot(\bchi_{2}-\bchi_{1})\bigr)
=\delta_{\bchi_{1},\bchi_{2}}E_{0}.
\end{equation*}
\end{keybox}

在写出以 $K_{4}^{+}$ 与 $K_{4}^{-}$ 中四矢量为指标的求和时，我们利用了纯空间矢量 $\bchi$
与这些四矢量的点积，等于与其在 $K_{3}$ 中投影的点积这一事实；
而在写出厄米共轭 $\phi^{\dagger}_{\bchi}$ 的第二种形式时，我们把指标矢量从 $K_{4}^{+}$ 中的 $\bk$
换成其在 $K_{4}^{-}$ 中的相反矢量。

你会注意到，模式系数算符中出现的因子 $1/\sqrt{E_{\bk}}$ 在这些空间定域算符中被抵消了。
把因子 $\sqrt{E_{\bk}}$ 插入函数 $g_{\bchi}$ 的定义，正是为了实现这一抵消！
洛伦兹理论中的相应公式是对一切空间三矢量的积分，它们需要因子 $1/\sqrt{E_{\bk}}$
以确保某些乘积为洛伦兹不变量——因子 $1/E_{\bk}$ 恰好能把对空间矢量的积分变成对质量双曲面的洛伦兹不变积分。
但我们的求和是对\emph{已经}均匀分布于三维球面半边的四矢量进行的，
故不希望留下任何依赖参考系的能量函数，以免使表述 $SO(4)$ 不变量的量变得困难。

这组算符场定义于薛定谔绘景，其中某些算符被视为不随时间变化（如我们原初谐振子中的 $x$ 与 $p$），
而态矢量随时间变化。但在量子场论中更常用海森堡绘景，其中态矢量被视为固定不变，一切算符都成为时间的函数。

为切换到海森堡绘景，若态 $\psi$ 须固定为 $\psi(0)$，则算符须随时间按
\begin{equation*}
A(t)=\exp(\ii Ht)A(0)\exp(-\ii Ht)
\end{equation*}
变化。现考虑湮灭算符 $a_{\bk}$ 与我们的哈密顿量 $H_{\text{RSW}}$ 的具体情形。
除 $E_{\bk}a^{\dagger}_{\bk}a_{\bk}$ 一项外，算符 $a_{\bk}$ 与 $H_{\text{RSW}}$ 中其它一切对易，故
\begin{equation*}
[a_{\bk},H_{\text{RSW}}]=[a_{\bk},E_{\bk}a^{\dagger}_{\bk}a_{\bk}]
=E_{\bk}[a_{\bk},a^{\dagger}_{\bk}]a_{\bk}
\;\Longrightarrow\;
a_{\bk}H_{\text{RSW}}-H_{\text{RSW}}a_{\bk}=E_{\bk}a_{\bk},
\end{equation*}
从而
\begin{equation*}
H_{\text{RSW}}a_{\bk}=a_{\bk}\bigl(H_{\text{RSW}}-E_{\bk}\bigr),
\qquad
H_{\text{RSW}}^{n}a_{\bk}=a_{\bk}\bigl(H_{\text{RSW}}-E_{\bk}\bigr)^{n}.
\end{equation*}
把哈密顿量的指数视为幂级数，即知
\begin{equation*}
\exp(\ii H_{\text{RSW}}t)a_{\bk}=a_{\bk}\exp\bigl(\ii(H_{\text{RSW}}-E_{\bk})t\bigr),
\qquad
\exp(\ii H_{\text{RSW}}t)a_{\bk}\exp(-\ii H_{\text{RSW}}t)=a_{\bk}\exp(-\ii E_{\bk}t).
\end{equation*}
对上式两端取厄米共轭，得到相应产生算符的变换性质：
\begin{equation*}
\exp(\ii H_{\text{RSW}}t)a^{\dagger}_{\bk}\exp(-\ii H_{\text{RSW}}t)=a^{\dagger}_{\bk}\exp(\ii E_{\bk}t).
\end{equation*}
显然 $b_{\bk}$ 与 $b^{\dagger}_{\bk}$ 以同样方式变换。我们全部 $c$ 算符不过是 $a$ 或 $b$ 的更名；
只要明确 $E_{\bk}$ 是与 $K_{4}$ 中四矢量 $\bk$ 相联系时间分量的\emph{绝对值}，
同样的规则也适用于 $c_{\bk}$ 与 $c^{\dagger}_{\bk}$。利用这些变换规则，
便可把场算符及相应的动量写成海森堡绘景中空间与时间两者的函数。

\begin{keybox}[量子化后的黎曼标量场（海森堡绘景）]
$\bchi$ 是取自集合 $\Chi$ 的空间三矢量。算符含时，作用于不变的态矢量。
\begin{align*}
\phi_{\bchi}(t)&=\frac{\ii}{\sqrt{2E_{0}}}\sum_{\bk\in K_{4}^{+}}
\bigl(c_{\bk}\exp(-\ii E_{\bk}t)-c^{\dagger}_{-\bk}\exp(\ii E_{\bk}t)\bigr)\exp(-\ii\,\bk\cdot\bchi),\\
\phi^{\dagger}_{\bchi}(t)&=\frac{\ii}{\sqrt{2E_{0}}}\sum_{\bk\in K_{4}^{-}}
\bigl(c_{\bk}\exp(-\ii E_{\bk}t)-c^{\dagger}_{-\bk}\exp(\ii E_{\bk}t)\bigr)\exp(-\ii\,\bk\cdot\bchi),\\
\Pi_{\bchi}(t)&=\partial_{t}\phi^{\dagger}_{\bchi}(t)
=\frac{1}{\sqrt{2E_{0}}}\sum_{\bk\in K_{4}^{-}}E_{\bk}
\bigl(c_{\bk}\exp(-\ii E_{\bk}t)+c^{\dagger}_{-\bk}\exp(\ii E_{\bk}t)\bigr)\exp(-\ii\,\bk\cdot\bchi)\\
&=\frac{1}{\sqrt{2E_{0}}}\sum_{\bk\in K_{4}^{+}}E_{\bk}
\bigl(c_{-\bk}\exp(-\ii E_{\bk}t)+c^{\dagger}_{\bk}\exp(\ii E_{\bk}t)\bigr)\exp(\ii\,\bk\cdot\bchi),\\
\Pi^{\dagger}_{\bchi}(t)&=\partial_{t}\phi_{\bchi}(t)
=\frac{1}{\sqrt{2E_{0}}}\sum_{\bk\in K_{4}^{+}}E_{\bk}
\bigl(c_{\bk}\exp(-\ii E_{\bk}t)+c^{\dagger}_{-\bk}\exp(\ii E_{\bk}t)\bigr)\exp(-\ii\,\bk\cdot\bchi)\\
&=\frac{1}{\sqrt{2E_{0}}}\sum_{\bk\in K_{4}^{-}}E_{\bk}
\bigl(c_{-\bk}\exp(-\ii E_{\bk}t)+c^{\dagger}_{\bk}\exp(\ii E_{\bk}t)\bigr)\exp(\ii\,\bk\cdot\bchi).
\end{align*}
\end{keybox}

我们来看 $\phi_{\bchi}(t)$ 这类算符的一些期望值。由于它们是模式振幅的线性组合，
它们会有与模式振幅相同的某些性质：每个算符自身在 $\psi_{0}$、$\psi_{\bq}$ 这类态上的期望值为零。
故再次考察平方振幅，它由 $\phi^{\dagger}_{\bchi}(t)\phi_{\bchi}(t)$ 给出：
\begin{equation*}
\langle\psi_{0},\phi^{\dagger}_{\bchi}(t)\phi_{\bchi}(t)\psi_{0}\rangle
=\frac{N_{\text{modes}}}{4E_{0}}\approx\frac{3\pi}{8m} .
\end{equation*}
故与模式振幅一样，这些空间系数算符对基态能量敏感。对在模式 $\bq$ 中含一个量子的态，有
\begin{equation*}
\langle\psi_{\bq},\phi^{\dagger}_{\bchi}(t)\phi_{\bchi}(t)\psi_{\bq}\rangle
=\frac{N_{\text{modes}}+2}{4E_{0}}\approx\frac{3\pi(1+2/N_{\text{modes}})}{8m}.
\end{equation*}
可见空间系数完全分辨不出存在的是\emph{哪一个}单量子平面波模式，几乎无法把这些态与基态区分开。

这并不太令人意外；我们期望这些空间算符能分辨的，是空间定域态之间的差别。
我们可以模仿由平面波模式 $f_{\bk}$ 构造空间定域模式 $g_{\bchi}$ 的方式，
构造一组正交归一的空间定域态，定义 $\psi_{\bchi}$ 为
\begin{equation*}
\psi_{\bchi}\overset{?}{=}\sum_{\bk\in K_{3}}\sqrt{\frac{E_{\bk}}{E_{0}}}\exp(\ii\,\bk\cdot\bchi)\psi_{\bk}.
\end{equation*}
然而这些态给不出平方空间系数算符的简单期望值，也不易由此构造 $SO(4)$ 不变量。
但事实证明，把空间系数算符本身作用于基态，可以得到一些有用的定域态：
\begin{align*}
\psi^{(+)}_{\bchi,t}&=2\sqrt{\frac{E_{0}}{N_{\text{modes}}}}\,\phi^{\dagger}_{\bchi}(t)\psi_{0}
=-\ii\sqrt{\frac{2}{N_{\text{modes}}}}\sum_{\bk\in K_{4}^{+}}
\exp\bigl(\ii(E_{\bk}t+\bk\cdot\bchi)\bigr)\psi_{\bk},\\
\psi^{(-)}_{\bchi,t}&=2\sqrt{\frac{E_{0}}{N_{\text{modes}}}}\,\phi_{\bchi}(t)\psi_{0}
=-\ii\sqrt{\frac{2}{N_{\text{modes}}}}\sum_{\bk\in K_{4}^{-}}
\exp\bigl(\ii(E_{\bk}t+\bk\cdot\bchi)\bigr)\psi_{\bk}.
\end{align*}
其意图是该态在特定时刻 $t$ 定域于三矢量 $\bchi$ 附近，在其他时刻预期会空间弥散开来；
上标 $(\pm)$ 标记粒子或反粒子。

我们来计算粒子态之间的内积：
\begin{align*}
\langle\psi^{(+)}_{\bchi_{1},t_{1}},\psi^{(+)}_{\bchi_{2},t_{2}}\rangle
&=\frac{4E_{0}}{N_{\text{modes}}}\langle\psi_{0},
\phi_{\bchi_{1}}(t_{1})\phi^{\dagger}_{\bchi_{2}}(t_{2})\psi_{0}\rangle\\
&=\frac{2}{N_{\text{modes}}}\sum_{\bk\in K_{4}^{+}}
\exp\bigl(\ii\bigl(E_{\bk}(t_{2}-t_{1})+\bk\cdot(\bchi_{2}-\bchi_{1})\bigr)\bigr)
=\frac{2}{N_{\text{modes}}}\sum_{\bk\in K_{4}^{+}}\exp(\ii\,\bk\cdot\bdelta),
\end{align*}
其中定义了四矢量
\begin{equation*}
\bdelta=(t_{2}-t_{1},\ \bchi_{2}-\bchi_{1}).
\end{equation*}
不同 $\bchi$ 或 $t$ 的态并不正交，但内积在位置与时间都相同时取峰，
并随空间间隔增大而呈 $J_{1}(m|\bchi_{2}-\bchi_{1}|)/(m|\bchi_{2}-\bchi_{1}|)$ 的衰减。

反粒子态之间则给出
\begin{align*}
\langle\psi^{(-)}_{\bchi_{1},t_{1}},\psi^{(-)}_{\bchi_{2},t_{2}}\rangle
&=\frac{4E_{0}}{N_{\text{modes}}}\langle\psi_{0},
\phi^{\dagger}_{\bchi_{1}}(t_{1})\phi_{\bchi_{2}}(t_{2})\psi_{0}\rangle\\
&=\frac{2}{N_{\text{modes}}}\sum_{\bk\in K_{4}^{-}}
\exp\bigl(\ii\bigl(E_{\bk}(t_{2}-t_{1})+\bk\cdot(\bchi_{2}-\bchi_{1})\bigr)\bigr)\\
&=\frac{2}{N_{\text{modes}}}\sum_{\bk\in K_{4}^{-}}
\exp\bigl(\ii\bigl((-k^{t})(t_{2}-t_{1})-\bk_{(3)}\cdot(\bchi_{2}-\bchi_{1})\bigr)\bigr)\\
&=\frac{2}{N_{\text{modes}}}\sum_{\bk\in K_{4}^{-}}\exp(-\ii\,\bk\cdot\bdelta),
\end{align*}
其中从倒数第三行到倒数第二行，用到了对 $K_{4}^{-}$ 中每个四矢量，
取相反空间部分、相同时间分量的矢量仍在 $K_{4}^{-}$ 中这一事实。该结果与粒子态间内积完全相同。

全部粒子态与反粒子态彼此正交：
\begin{equation*}
\langle\psi^{(+)}_{\bchi_{2},t_{2}},\psi^{(-)}_{\bchi_{1},t_{1}}\rangle=0 .
\end{equation*}

\subsubsection{标量场的传播子（Propagator for a Scalar Field）}

量子场论中最重要的事情之一是场的传播子（propagator），
它给出“已知粒子在 $t_{1}$ 时刻确定位于 $\bx_{1}$，则在 $t_{2}$ 时刻位于 $\bx_{2}$”的振幅。
洛伦兹 QFT 中，这以清楚区分粒子与反粒子的方式完成，并可表述为遵从通常时间之箭的形式。
例如，事件 $(t_{1},\bx_{1})$ 与 $(t_{2},\bx_{2})$ 的\emph{费曼传播子}（Feynman propagator）
定义如下：若 $t_{2}>t_{1}$，它是粒子从 $(t_{1},\bx_{1})$ 到 $(t_{2},\bx_{2})$ 的振幅；
若 $t_{2}<t_{1}$，则是反粒子沿相反方向传播的振幅。黎曼情形中并无时间之箭的类似概念，
故我们简单地把传播子取为两种可能性的对称组合：
\begin{align*}
&\tfrac12\Bigl(\langle\psi^{(+)}_{\bchi_{1},t_{1}},\psi^{(+)}_{\bchi_{2},t_{2}}\rangle
+\langle\psi^{(-)}_{\bchi_{2},t_{2}},\psi^{(-)}_{\bchi_{1},t_{1}}\rangle\Bigr)\\
&=\frac{2E_{0}}{N_{\text{modes}}}\Bigl\langle\psi_{0},
\bigl(\phi_{\bchi_{1}}(t_{1})\phi^{\dagger}_{\bchi_{2}}(t_{2})
+\phi^{\dagger}_{\bchi_{2}}(t_{2})\phi_{\bchi_{1}}(t_{1})\bigr)\psi_{0}\Bigr\rangle\\
&=\frac{1}{N_{\text{modes}}}\sum_{\bk\in K_{4}}\exp(\ii\,\bk\cdot\bdelta)
\approx\frac{2J_{1}(m|\bdelta|)}{m|\bdelta|},\qquad
\bdelta=(t_{2}-t_{1},\ \bchi_{2}-\bchi_{1}).
\end{align*}

\orthofig{images/07__011.png}{$SO(4)$-不变的传播子（SO(4)-invariant propagator）}

显然它是 $SO(4)$ 不变的：它只是四矢量 $\bdelta$ 长度的函数，而这一量所有观测者都会认同。
若把 $\bchi_{1}$ 与 $t_{1}$ 视为固定、把 $\bchi_{2}$ 与 $t_{2}$ 当作坐标，
则它是绕 $(t_{1},\bchi_{1})$ 具有四维旋转对称性的无源 RSW 方程的解。

当然，贝塞尔函数近似仅是近似。它有 $SO(4)$ 不变性与易计算的优势，但不满足 $T^{4}$ 边界条件，
故若沿不同方向绕环面走一圈会得到不同的值。实践中，取任意两个事件间最短的四维距离会给出恰当的答案。

若我们在传播子中交换“粒子”与“反粒子”的定义，这相当于交换两个内积中两项的次序，
从而取整个表达式的复共轭。但由于最终值是实数，这样做没有差别。基于同样的理由，
交换两个事件上的标签也没有差别。

\orthofig{images/07__012.png}{传播子的物理解释（Propagator interpreted）}

左图说明了传播子所测量的是什么。两个事件由传播矢量指向一切可能方向的平面波相连，
把这些波分为粒子态与反粒子态并无不变的意义。

然而，\emph{对给定观测者}而言，我们可以把传播矢量指向时间背向的波描述为沿时间正向行进的反粒子，
把传播矢量指向时间正向的波描述为同样沿时间正向行进的粒子。

关键在于我们始终把两种可能性都计入。若照搬洛伦兹对费曼传播子的处方，
因某观测者把事件 1 判为早于事件 2 而不计入反粒子态，
那么即便另一个认同该顺序的观测者，也会选出不同的半三维球面上的粒子态（图中红箭头所示），
从而得到一个不同的振幅。

\subsection{自由黎曼狄拉克场的正则量子化（Canonical Quantisation of a Free Riemannian Dirac Field）}

假设我们把为量子化复标量场所作的初步构造重做一遍，只是用黎曼狄拉克方程取代黎曼标量波方程。
在此前讨论狄拉克方程时，我们把它当作描述量子波函数的方程；但眼下我们先假装在研究一个“经典”狄拉克场，
其中场由四个复分量组成并满足狄拉克方程，却与量子力学无关。
满足该方程的场构成欧几里得群的一个表示；我们实际上正是按群表示的观点\emph{推导}出狄拉克方程的。
故起初我们只是为这个新表示模仿先前对标量表示所做的一切。

为记号简洁，我们再次把场的最大角频率（很快会成为粒子质量）记作 $m$，
并把与该质量及 $T^{4}$ 宇宙尺寸相兼容的全部可能传播四矢量集合记作 $K_{4}$，
把 $K_{4}$ 中矢量投影到三维所得集合记作 $K_{3}$，
把与 $K_{3}$ 中任一 $\bk$ 相联系的唯一正能量记作 $E_{\bk}$：
\begin{equation*}
E_{\bk}=\sqrt{m^{2}-|\bk|^{2}} .
\end{equation*}
若我们要描述不止一种场的情形，就需区分各场对应的所有这些量的不同版本，
因为它们在各种情形下当然没有理由相同。

为把狄拉克场 $\psi$ 作傅里叶模式展开，我们将使用下列狄拉克方程的解，
它们是此前讨论过的平面波解的推广：

\begin{keybox}[黎曼狄拉克方程在 Weyl 基中的平面波解]
$\bk\in K_{3}$，$\{\xi_{s}\}$ 是 $\mathbb{C}^{2}$ 的一组正交归一基，
$H_{\mu}$ 是作为矩阵的单位四元数，$\sum k^{j}H_{j}$ 对空间指标 $x,y,z$ 求和。

正频解：
\begin{align*}
\psi_{+,\bk,s}(\bx,t)&=u(\bk,s)\exp\bigl(-\ii(\bk\cdot\bx+E_{\bk}t)\bigr)
&&\text{(Plane wave)}\\
u(\bk,s)&=\frac{\bigl([(m+E_{\bk})H_{t}+k^{j}H_{j}]\xi_{s},\
[(m+E_{\bk})H_{t}-k^{j}H_{j}]\xi_{s}\bigr)}{2\sqrt{m(m+E_{\bk})}},\\
u(\bk,s_{1})^{\dagger}u(\bk,s_{2})&=\delta_{s_{1},s_{2}},\qquad
u(\bk,s_{1})^{\dagger}\gamma^{t}u(\bk,s_{2})=\delta_{s_{1},s_{2}}\frac{E_{\bk}}{m}.
\end{align*}
负频解：
\begin{align*}
\psi_{-,\bk,s}(\bx,t)&=v(\bk,s)\exp\bigl(+\ii(\bk\cdot\bx+E_{\bk}t)\bigr)
&&\text{(Plane wave)}\\
v(\bk,s)&=\frac{\bigl([(m+E_{\bk})H_{t}+k^{j}H_{j}]\xi_{s},\
-[(m+E_{\bk})H_{t}-k^{j}H_{j}]\xi_{s}\bigr)}{2\sqrt{m(m+E_{\bk})}},\\
v(\bk,s_{1})^{\dagger}v(\bk,s_{2})&=\delta_{s_{1},s_{2}},\qquad
v(\bk,s_{1})^{\dagger}\gamma^{t}v(\bk,s_{2})=-\delta_{s_{1},s_{2}}\frac{E_{\bk}}{m}.
\end{align*}
$u$ 与 $v$ 旋量的乘积：
\begin{align*}
u(\bk,s_{1})^{\dagger}v(\bk,s_{2})=0,&\qquad
v(\bk,s_{1})^{\dagger}u(\bk,s_{2})=0,\\
u(\bk,s_{1})^{\dagger}\gamma^{t}v(-\bk,s_{2})=0,&\qquad
v(-\bk,s_{1})^{\dagger}\gamma^{t}u(\bk,s_{2})=0 .
\end{align*}
\end{keybox}

若你学过这些内容通常的洛伦兹版本，这些正交关系会显得有些奇怪！
洛伦兹 QFT 中，你需要在任意一对狄拉克旋量之间插入 $\gamma^{t}$ 才能得到洛伦兹不变的内积，此时有
\begin{equation*}
u(\bk,s_{1})^{\dagger}\gamma^{t}u(\bk,s_{2})=\delta_{s_{1},s_{2}}\quad\text{[Lorentzian]}.
\end{equation*}
但在这里，由于旋转群 $SO(4)$ 在狄拉克旋量上有一个酉表示，
我们可以简单地在 $u$ 旋量之间或 $v$ 旋量之间取普通内积，它便是 $SO(4)$ 不变的。

在下面给出的场的模式展开中，我们把这两个平面波解 $\psi_{\pm,\bk,s}(\bx,t)$ 的纯空间部分
提取为两组函数 $f_{\bk,s}(\bx)$ 与 $g_{\bk,s}(\bx)$，
然后把一般解 $\psi(\bx,t)$ 展开为这些函数乘以含时系数之和。
为使 $\psi(\bx,t)$ 是狄拉克方程的解，这些含时系数应满足与 $\psi_{\pm,\bk,s}(\bx,t)$
的含时部分相同的方程，即 $\exp(-\ii(\pm E_{\bk}t))$。

\begin{keybox}[“经典”黎曼狄拉克方程]
$\sum\gamma^{j}\partial_{j}$ 对空间指标 $x,y,z$ 求和。
\begin{align*}
0&=\ii\gamma^{\mu}\partial_{\mu}\psi(\bx,t)-m\psi(\bx,t)
&&\text{(Equation of Motion)}\\
L_{\text{RD}}&=\int\psi^{\dagger}(\bx,t)\bigl(\ii\gamma^{\mu}\partial_{\mu}\psi(\bx,t)-m\psi(\bx,t)\bigr)\dd^{3}\bx
&&\text{(Lagrangian)}\\
\Pi(\bx,t)&=\ii\psi^{\dagger}(\bx,t)\gamma^{t},\qquad
\Pi^{\dagger}(\bx,t)=0
&&\text{(Momenta)}\\
H_{\text{RD}}&=\int\psi^{\dagger}(\bx,t)\bigl(-\ii\gamma^{j}\partial_{j}\psi(\bx,t)+m\psi(\bx,t)\bigr)\dd^{3}\bx
&&\text{(Hamiltonian)}
\end{align*}
模式展开形式（$\bk\in K_{3}$，$s$ 为自旋 $\pm\tfrac12$）：
\begin{align*}
f_{\bk,s}(\bx)&=\sqrt{\tfrac{1}{V}}\,u(\bk,s)\exp(-\ii\,\bk\cdot\bx),\\
&\qquad g_{\bk,s}(\bx)=\sqrt{\tfrac{1}{V}}\,v(\bk,s)\exp(+\ii\,\bk\cdot\bx)
&&\text{(Spatial modes)}\\
\psi(\bx,t)&=\sum_{\bk\in K_{3},\ s=\pm\tfrac12}
\bigl(A_{\bk,s}(t)f_{\bk,s}(\bx)+B_{\bk,s}(t)g_{\bk,s}(\bx)\bigr)
&&\text{(Mode Expansion)}\\
\partial_{t}A_{\bk,s}(t)&=-\ii E_{\bk}A_{\bk,s}(t),\\
&\qquad\partial_{t}B_{\bk,s}(t)=+\ii E_{\bk}B_{\bk,s}(t)
&&\text{(Equations of Motion)}\\
L_{\text{RD}}&=\sum_{\bk\in K_{3},\ s=\pm\tfrac12}\frac{E_{\bk}}{m}
\Bigl[\ii\bigl(A^{*}_{\bk,s}(t)\partial_{t}A_{\bk,s}(t)-B^{*}_{\bk,s}(t)\partial_{t}B_{\bk,s}(t)\bigr)\\
&\qquad -E_{\bk}\bigl(A^{*}_{\bk,s}(t)A_{\bk,s}(t)+B^{*}_{\bk,s}(t)B_{\bk,s}(t)\bigr)\Bigr]
&&\text{(Lagrangian)}\\
\Pi_{A,\bk,s}(t)&=\ii\frac{E_{\bk}}{m}A^{*}_{\bk,s}(t),\\
&\qquad\Pi_{B,\bk,s}(t)=-\ii\frac{E_{\bk}}{m}B^{*}_{\bk,s}(t),\\
\Pi_{A^{*},\bk,s}(t)&=0,\qquad
\Pi_{B^{*},\bk,s}(t)=0
&&\text{(Momenta)}\\
H_{\text{RD}}&=\sum_{\bk\in K_{3},\ s=\pm\tfrac12}\frac{E_{\bk}^{2}}{m}
\bigl(A^{*}_{\bk,s}(t)A_{\bk,s}(t)+B^{*}_{\bk,s}(t)B_{\bk,s}(t)\bigr)\\
&=-\ii\sum_{\bk\in K_{3},\ s=\pm\tfrac12}E_{\bk}
\bigl[\Pi_{A,\bk,s}(t)A_{\bk,s}(t)\\
&\qquad-\Pi_{B,\bk,s}(t)B_{\bk,s}(t)\bigr]
&&\text{(Hamiltonian)}
\end{align*}
\end{keybox}

用模式系数 $A_{\bk,s}(t)$、$B_{\bk,s}(t)$ 导出拉格朗日量与哈密顿量需要相当多的工作，
但验证结果并不难：只需检验它们分别经欧拉-拉格朗日方程与哈密顿方程给出所需的运动方程。

要把经典哈密顿量化为场量子态希尔伯特空间上的算符，我们需要修改量子化谐振子与标量场模式时所用的程序。
在那里我们发现产生、湮灭算符之间的对易子为
\begin{equation*}
[a_{\alpha},a^{\dagger}_{\beta}]=\delta_{\alpha,\beta},\qquad
[a_{\alpha},a_{\beta}]=0,\qquad
[a^{\dagger}_{\alpha},a^{\dagger}_{\beta}]=0,
\end{equation*}
其中 $\alpha$、$\beta$ 标记模式。但由于狄拉克场描述费米子（fermions），
我们希望把产生算符作用于真空态 $\psi_{0}$ 所得的态对模式指标反对称：
\begin{equation*}
a^{\dagger}_{\alpha}a^{\dagger}_{\beta}\psi_{0}=-a^{\dagger}_{\beta}a^{\dagger}_{\alpha}\psi_{0}.
\end{equation*}
只要把其中一个对易子换成反对易子（anticommutator）便可实现：
\begin{equation*}
\{a^{\dagger}_{\alpha},a^{\dagger}_{\beta}\}
=a^{\dagger}_{\alpha}a^{\dagger}_{\beta}+a^{\dagger}_{\beta}a^{\dagger}_{\alpha}=0 .
\end{equation*}
对上式取厄米共轭，得
\begin{equation*}
\{a_{\alpha},a_{\beta}\}=0 .
\end{equation*}
沿用 Zee 的做法，我们还可为第三个反对易子给出理由：
\begin{equation*}
\{a_{\alpha},a^{\dagger}_{\beta}\}=\delta_{\alpha,\beta},
\end{equation*}
其办法是要求数算符 $N=\sum_{\alpha}a^{\dagger}_{\alpha}a_{\alpha}$ 满足
\begin{equation*}
N a^{\dagger}_{\beta}=a^{\dagger}_{\beta}N+a^{\dagger}_{\beta}.
\end{equation*}
这是一个合理的要求：把它作用于含 $n$ 个量子（其中尚无一个处于模式 $\beta$，
故 $a^{\dagger}_{\beta}\psi_{n}=\psi_{n+1}$）的态 $\psi_{n}$，便化为
\begin{equation*}
N a^{\dagger}_{\beta}\psi_{n}=a^{\dagger}_{\beta}N\psi_{n}+a^{\dagger}_{\beta}\psi_{n}
\;\Longrightarrow\;
N\psi_{n+1}=(n+1)\psi_{n+1}.
\end{equation*}
给定所提议的反对易子，连同前面“一切产生算符都反交换”的论断，便可导出所需恒等式：
\begin{align*}
N a^{\dagger}_{\beta}
&=\sum_{\alpha}a^{\dagger}_{\alpha}a_{\alpha}a^{\dagger}_{\beta}
=\sum_{\alpha}a^{\dagger}_{\alpha}\bigl(\delta_{\alpha,\beta}-a^{\dagger}_{\beta}a_{\alpha}\bigr)
=a^{\dagger}_{\beta}-\sum_{\alpha}a^{\dagger}_{\alpha}a^{\dagger}_{\beta}a_{\alpha}\\
&=a^{\dagger}_{\beta}+a^{\dagger}_{\beta}\sum_{\alpha}a^{\dagger}_{\alpha}a_{\alpha}
=a^{\dagger}_{\beta}N+a^{\dagger}_{\beta}.
\end{align*}

我们已求得的经典哈密顿量
\begin{equation*}
H_{\text{RD}}=-\ii\sum_{\bk\in K_{3},\ s=\pm\tfrac12}E_{\bk}
\bigl[\Pi_{A,\bk,s}(t)A_{\bk,s}(t)-\Pi_{B,\bk,s}(t)B_{\bk,s}(t)\bigr]
\end{equation*}
这一形式，若把这些坐标与动量或多或少直接等同于各模式的湮灭、产生算符，便说得通：
\begin{equation*}
A_{\bk,s}(t)\to\sqrt{\frac{m}{E_{\bk}}}\,a_{\bk,s},\qquad
B_{\bk,s}(t)\to\sqrt{\frac{m}{E_{\bk}}}\,b^{\dagger}_{\bk,s},
\end{equation*}
\begin{equation*}
\Pi_{A,\bk,s}=\ii\frac{E_{\bk}}{m}A^{*}_{\bk,s}(t)\to\ii\sqrt{\frac{E_{\bk}}{m}}\,a^{\dagger}_{\bk,s},
\qquad
\Pi_{B,\bk,s}=-\ii\frac{E_{\bk}}{m}B^{*}_{\bk,s}(t)\to-\ii\sqrt{\frac{E_{\bk}}{m}}\,b_{\bk,s}.
\end{equation*}
不过，由此产生一个问题：在把坐标与动量的乘积化为算符时，我们是否有权随意选取乘积的次序？
更少一点\emph{特设}（ad hoc）的做法是取某种对称形式，把两种非对易算符次序都包含进来——
或者，由于场是费米性的，取某种\emph{反}对称形式。故我们采用规则
\begin{align*}
\Pi_{A,\bk,s}(t)A_{\bk}&\to\frac{\ii}{2}\bigl(a^{\dagger}_{\bk,s}a_{\bk,s}-a_{\bk,s}a^{\dagger}_{\bk,s}\bigr)
=\ii\bigl(a^{\dagger}_{\bk,s}a_{\bk,s}-\tfrac12\bigr),\\
\Pi_{B,\bk,s}(t)B_{\bk}&\to\frac{-\ii}{2}\bigl(b^{\dagger}_{\bk,s}b_{\bk,s}-b_{\bk,s}b^{\dagger}_{\bk,s}\bigr)
=-\ii\bigl(b^{\dagger}_{\bk,s}b_{\bk,s}-\tfrac12\bigr).
\end{align*}
把所有部分合起来，最终得到

\begin{keybox}[量子化后的黎曼狄拉克场]
\begin{equation*}
H_{\text{RD}}=\sum_{\bk\in K_{3},\ s=\pm\tfrac12}E_{\bk}
\bigl[a^{\dagger}_{\bk,s}a_{\bk,s}+b^{\dagger}_{\bk,s}b_{\bk,s}-1\bigr]
\qquad\text{(Hamiltonian)}.
\end{equation*}
\end{keybox}

坦率地说，我们这里所做的一切远非严格！但无论如何，我们得到的哈密顿量算符确实为添加到场中的每个量子
赋予正确的能量 $E_{\bk}$。

若认真对待该结果，则（仅就该场而言）真空态的能量将是\emph{负}的：
\begin{equation*}
H_{\text{RD}}\psi_{0}=E_{0}\psi_{0},\qquad
E_{0}=-\sum_{\bk\in K_{3},\ s=\pm\tfrac12}E_{\bk}
\approx-\frac{4N_{\text{modes}}m}{3\pi},
\end{equation*}
这里 $N_{\text{modes}}$ 是与粒子质量及宇宙几何相兼容的\emph{标量}模式数，
我们在其它因子中计入一倍，以计及两种自旋。

费米子的负真空能并非黎曼物理独有；洛伦兹 QFT 中也出现同样情形。
与标量场一样，黎曼情形的差别在于我们只有有限个模式，故真空能是有限的。

\section*{参考文献（References）}
\addcontentsline{toc}{section}{参考文献（References）}
\begin{itemize}[leftmargin=2em]
  \item[{[1]}] \emph{An Introduction to Quantum Field Theory}, Michael E. Peskin and Daniel V. Schroeder,
        Westview Press, 1995. 第 3.2 节。
  \item[{[2]}] 例如 \emph{Quantum Mechanics}, Leonard I. Schiff, McGraw-Hill, 1968.
        第 13 节推导谐振子波函数，第 25 节使用阶梯算符方法。
  \item[{[3]}] Peskin and Schroeder, \emph{op. cit.}, 第 2 章。
  \item[{[4]}] \emph{The Quantum Theory of Fields, Volume I: Foundations}, Steven Weinberg,
        Cambridge University Press, 1995. 第 1.2、5.2 节。
  \item[{[5]}] \emph{Mechanics}, L. D. Landau and E. M. Lifshitz, Butterworth-Heinemann, 1976.
        第 45 节。\textbf{注意}：Landau 与 Lifshitz 在其 (42.5) 式中定义的泊松括号
        与我们所用约定相反。
  \item[{[6]}] \emph{Quantum Mechanics and Path Integrals}, Richard P. Feynman, Albert R. Hibbs
        and Daniel F. Styer, Dover, 2010.
  \item[{[7]}] \emph{Spinors and Trialities}, John Baez.
  \item[{[8]}] \emph{Quantum Field Theory in a Nutshell}, A. Zee, Princeton University Press, 2010.
        第 II.2 章。
\end{itemize}

\chapter{黎曼广义相对论}

本章探讨引力的本质在黎曼宇宙中会发生怎样的变化。我们首先从直观论证出发，说明为何黎曼宇宙既不会有黑洞，也不存在（我们所理解的）引力波与大爆炸；随后在进阶部分系统推导黎曼宇宙的 Schwarzschild--Kottler 解、测试粒子的运动常数与有效势、圆轨道与逃逸速度、引力蓝移与时间膨胀、引力散射，并借助 Einstein 方程把时空曲率与物质的熵态联系起来，讨论各类均匀各向同性宇宙模型及其拓扑。全章沿用黎曼度规 $(+,+,+,+)$ 与光速可变的设定，引力常数经单位选取归一化。

\section{我们确实没有黑洞（Yes, We Have No Black Holes）}

黎曼宇宙中的引力行为与我们自身宇宙有多大不同？

一种可能的变化是：既然既没有固定的光速、对物质客体也没有任何速度上限，那么看起来就不可能有黑洞——若你能以任意高的速度运动，那么无论一个天体的逃逸速度有多大，原则上你都能达到该速度。

但另一方面，物体的动能并不会随速度无限增大。如我们此前所讨论的，当速度趋于无穷时，动能趋近于物体的静止质量。那么，若你身处一个引力势阱中、深到你的势能小于 $-m$（即静止质量的反号），又会怎样？至少在牛顿引力中，质量为 $m$ 的物体在距质量为 $M$ 的恒星或行星 $r$ 处的势能为
\[
U=-\frac{GMm}{r},
\]
其大小在 $r<GM$ 时就超过静止质量 $m$。把引力势能与黎曼运动学结合，可以证明如下量守恒：
\begin{equation*}
E^{2}=\frac{1+2GM/r}{1+v^{2}},
\end{equation*}
这正是引力场中能量守恒的黎曼广义相对论（Riemannian general relativity）版本。无论物体下落并加速，还是远离大质量天体上升而被减速，这个量都保持不变。因此，为使物体在距大质量天体无穷远处速度为 $0$——也就是它恰好以逃逸速度运动的条件——需要 $E^{2}=1$。换言之：
\begin{align*}
\frac{1+2GM/r}{1+v_{esc}(r)^{2}} &= 1,\\
v_{esc}(r)^{2} &= \frac{2GM}{r},\\
\frac{1}{2}mv_{esc}(r)^{2}-\frac{GMm}{r} &= 0.
\end{align*}

最后一行表明，牛顿形式的动能与牛顿形式的势能之和为零。这同一最终结果在我们自身宇宙中其实也成立——但在我们的宇宙里 $v$ 必须小于 $1$，于是在 $r=2GM$ 处逃逸已不可能，你便越过了黑洞的事件视界（event horizon）。而在黎曼宇宙中，正如我们最初的直觉所提示的：$v$ 没有上限，意味着引力永远不会强到无法逃脱。

\orthofig{images/06__005.png}{黎曼与牛顿引力场中靠近大质量天体时的轨迹}

尽管黎曼宇宙没有黑洞，一个足够致密的大质量天体周围的弯曲四维空间仍会产生一些奇特效应。首先是时间膨胀（time dilation），其方向与我们宇宙中的引力时间膨胀相反：在黎曼宇宙中，若你在大质量天体附近度过一段时间，然后再与一位一直保持距离的同伴会合，则\emph{你}经历的时间比同伴\emph{更多}。

另一种效应是曲率对离心力的修正。在我们宇宙中，若你离黑洞太近，离心力实际上会反向，帮你被推向洞中。而在黎曼宇宙中，离心力始终是斥力，甚至可以比牛顿版本的离心力更强。其结果是，若你从某个致密大质量天体旁经过得太近，你非但不会沿牛顿双曲轨道绕行，反而会被推离开它。

由于离心力更强，行星绕行轨道的速度存在上限——无论它与所绕恒星多近、无论恒星多重。任何圆轨道都不可能以大于 $1/\sqrt2$ 的速度存在；若行星速度更快，离心力就会压倒引力。

\section{没有引力波（No Gravitational Waves）}

黎曼宇宙中支配引力波的方程，可能面临与我们此前在光波中见过的同类问题：在某一维上频率过高的波会在另一维上指数增长。但对引力波而言情况更糟，因为这会发生在\emph{任何}频率的波上——并不需要超过某个最高频率。

其结果是，那个阻止光波指数增长的机制，会把引力波彻底消除。在我们宇宙中，两个彼此靠得很近相互绕转的大质量物体会以引力波的形式辐射能量，导致它们盘旋向内、最终碰撞。在黎曼宇宙中，任何用于产生引力波的能量都会产生相反的效果，使绕转天体相互远离……但事实上这不会发生，因为根本不可能有引力辐射。

\section{……多半也没有大爆炸（\ldots and Probably No Big Bang）}

\orthofig{images/06__006.png}{含奇点的黎曼宇宙玩具模型}
\orthofig{images/06__007.png}{不含奇点的黎曼宇宙玩具模型}

在洛伦兹物理中，若你把大多数简单宇宙模型的历史回溯到过去，最终会到达一个万物密度无穷、广义相对论的经典定律必然失效的时刻：大爆炸（Big Bang）。

在黎曼物理中，这类宇宙模型是可能的，但它非常人为、造作。它要求你假设宇宙充满了世界线大致都指向同一方向的物质。在洛伦兹宇宙中这或许是个自然的假设，因为那里的世界线被限制在类时方向的锥内；但在黎曼宇宙中，它们可以指向四维空间中的\emph{任意}方向。

在黎曼情形下，若作更自然的假设——在宇宙尺度上物质的世界线对任何特定方向都无偏好——那么当我们选定某个方向并称之为"时间"时，物质在回溯过去时的平均密度根本无须改变。宇宙仍然可以收缩到零体积，但发生收缩的那个特定事件只是我们选择时间方向所造成的人为产物——正如纬线圈在南、北极收缩到零，只是我们选择坐标所造成的人为产物。

为简单起见，右图中的两个模型都画在球面上，但只有当存在一个足够大的正宇宙学常数（cosmological constant）时，四维空间才真正在所有方向上正弯曲——而在黎曼情形下，这等价于一个大的\emph{负}真空能。没有宇宙学常数时，第一个模型在"时间"方向上正弯曲、在其余三个方向上负弯曲。这意味着：若一团自由落体粒子的世界线与宇宙中其它物质大致对齐，它们会被挤入更小体积；但若它们相对该背景物质以无穷大速度运动，则反而会变得更加散开。

第二个模型在没有宇宙学常数时，在所有方向上都是负弯曲。

关于本主题的额外材料，提供给不介意稍高数学门槛的读者。

\section{进阶：黎曼 Schwarzschild 解（The Riemannian Schwarzschild Solution）}

把时空曲率与其中物质相联系的 Einstein 方程为
\begin{equation*}
G+\Lambda g=8\pi T,
\end{equation*}
其中：

\begin{itemize}
  \item $g$ 是时空度规（metric）；
  \item $G$ 是 Einstein 张量（Einstein tensor），由度规导出，量化时空曲率的某些方面；
  \item $\Lambda$ 是宇宙学常数；
  \item $T$ 是任何存在物质（包括电磁场等场）的应力--能量张量（stress-energy tensor）。
\end{itemize}

我们沿用广义相对论中的通行做法，采用使时间、空间乃至质量都以相同单位度量的单位制，并选取单位使引力常数（常被赋予与 Einstein 张量相同的符号 $G$，颇易混淆）等于 $1$。

这个方程没有理由需要为适用于黎曼宇宙而改变。引力相对于其它力的强弱在黎曼宇宙中或许不同、或许相同，但既然我们已把引力常数吸收进单位选取之中（无论其取值为何），就再无别的可做。Einstein 方程最重要的特征是 Einstein 张量的散度必为零，而度规乘以任一常数亦然。这使得这些量可以被合理地等同于应力--能量张量的某个常数倍，后者的散度必须为零，才能保证能量与动量的局域守恒。所有这一切在黎曼宇宙中同样成立、同样恰当。

于是唯一的真正区别在于：我们要求 Einstein 方程的\emph{解}中度规 $g$ 的号差（signature）为 $(++++)$，即所有维度都是类空的解。

当然，还有一种可能的改动不会破坏方程左端的无散度性质：我们可以要求引力常数为负，这样即便其大小经单位选取被设为 $1$，方程右端也会带一个负号。但正如我们将看到的，若想让黎曼宇宙中的引力在物体缓慢运动时本质上表现得像牛顿引力，保持引力常数为正才是正确的选择。

在我们自身宇宙的许多情境中，宇宙学常数 $\Lambda$ 可取为零，但我们将尽量保持一般性，在计算中保留它，只在求某些具体结果时才令其为零。因此我们从"Schwarzschild--Kottler"度规\textsuperscript{[1]}开始探索引力，它是 Schwarzschild 度规推广到包含宇宙学常数的版本。

\subsection*{Schwarzschild--Kottler 度规}

$r$、$\theta$、$\phi$ 为以质量 $M$ 的天体为中心的极坐标。度规为
\begin{align*}
g &= \left(1+\frac{2M}{r}-\frac{\Lambda r^{2}}{3}\right)\dd t^{2}
     +\frac{1}{1+\dfrac{2M}{r}-\dfrac{\Lambda r^{2}}{3}}\,\dd r^{2}
     +r^{2}\left[\dd\theta^{2}+(\sin\theta)^{2}\,\dd\phi^{2}\right]
     \qquad(\text{黎曼}),\\
g &= -\left(1-\frac{2M}{r}-\frac{\Lambda r^{2}}{3}\right)\dd t^{2}
     +\frac{1}{1-\dfrac{2M}{r}-\dfrac{\Lambda r^{2}}{3}}\,\dd r^{2}
     +r^{2}\left[\dd\theta^{2}+(\sin\theta)^{2}\,\dd\phi^{2}\right]
     \qquad(\text{洛伦兹}).
\end{align*}

该度规描述一个球对称、质量为 $M$ 的天体外部真空中的几何。眼下我们对黎曼宇宙中 $\Lambda$ 的取值乃至符号都不作承诺，只假设 $|\Lambda|$ 小到在所关心的区域内 $1+2M/r-\Lambda r^{2}/3$ 始终为正。

\orthofig{images/06__008.png}{黎曼 Schwarzschild 度规中到中心的距离}
\orthofig{images/06__009.png}{黎曼 Schwarzschild 度规中球面的表面积}

洛伦兹度规与黎曼度规的一个显著差别是：在黎曼情形下，沿 $t$、$\theta$、$\phi$ 为常数的路径测得的径向距离都是有限的。对 $\Lambda=0$，我们甚至能得到从坐标中心出发的距离的显式公式（假设该解一直适用到 $r=0$）：
\begin{equation*}
\operatorname{dist}(r)=\sqrt{r(r+2M)}-2M\,\operatorname{arcsinh}\!\sqrt{\frac{r}{2M}}.
\end{equation*}

该度规中 $r$ 为常数的曲面面积为 $4\pi r^{2}$，但 $\operatorname{dist}(r)$ 的值总小于 $r$ 本身，所以这些球面相对于其半径而言表面积比欧氏几何中的更大。这是负曲率空间的一个特征。

与此相反，洛伦兹 Schwarzschild 度规中的径向距离在 $r>2M$ 之前甚至不是类空的，而且此后它总\emph{大于} $r$ 坐标之差。

\subsection{运动常数与有效势（Constants of Motion and the Effective Potential）}

我们希望理解一个小的"测试粒子"（test particle）在此度规中的运动——该粒子小到对几何本身的影响可以忽略。我们采用的一般方法与 Misner、Thorne 和 Wheeler\textsuperscript{[2]}处理普通 Schwarzschild 几何中运动物体的方式相近。

测试粒子的世界线是一条测地线，我们用符号 $\bu$ 表示其四维速度（4-velocity），即与其世界线相切的单位矢量。矢量 $\bu$ 在坐标中的分量就是沿测地线的固有时 $\tau$ 下各坐标的变化率：
\begin{align*}
u^{t} &= \partial_{\tau} t, &
u^{r} &= \partial_{\tau} r, &
u^{\phi} &= \partial_{\tau} \phi, &
u^{\theta} &= \partial_{\tau} \theta.
\end{align*}

由于几何球对称，我们总可以选取坐标使测试粒子的世界线位于 $\theta=\pi/2$ 的"赤道面"内，从而 $u^{\theta}=\partial_{\tau}\theta=0$。

由于度规与坐标 $t$ 和 $\phi$ 无关，坐标矢量场 $\partial_{t}$ 与 $\partial_{\phi}$ 是 Killing 矢量场（Killing vector field）：即描述"把图形在几何中滑动而保持其形状与大小不变"的方式，好比我们在不扭曲一座岛屿地貌的前提下，给岛上每一点经度加 $5$ 度。$\partial_{t}$ 与 $\partial_{\phi}$ 为 Killing 矢量场的一个推论是：任一测地线的切矢量与这些场的点积沿测地线保持不变。于是对测试粒子的整条世界线都有
\begin{align*}
\bu\cdot\partial_{t} &= u_{t}=E, \\
\bu\cdot\partial_{\phi} &= u_{\phi}=L,
\end{align*}
其中我们引入 $E$ 与 $L$ 作为两个常数的名称。低标号分量 $u_{t}$ 与 $u_{\phi}$ 是对偶矢量的分量。$\bu$ 是单位矢量、其对偶矢量长度也为 $1$，这一事实用对偶分量表出即为
\begin{equation*}
g^{ij}u_{i}u_{j}=1
\end{equation*}
（此处使用 Einstein 求和约定）。由于度规对角，上标度规分量 $g^{ij}$ 就是 Schwarzschild--Kottler 度规定义中给出的下标配元的倒数，于是
\begin{align*}
g^{tt}u_{t}^{2}+g^{rr}u_{r}^{2}+g^{\phi\phi}u_{\phi}^{2} &= 1,\\
\frac{E^{2}}{1+2M/r-\Lambda r^{2}/3}+\left(1+\frac{2M}{r}-\frac{\Lambda r^{2}}{3}\right)u_{r}^{2}+\frac{L^{2}}{r^{2}} &= 1.
\end{align*}

若把 $u_{r}$ 用 $u^{r}=\partial_{\tau}r$ 表出，则有
\begin{equation*}
u_{r}=g_{rr}u^{r}=\frac{\partial_{\tau} r}{1+2M/r-\Lambda r^{2}/3},
\end{equation*}
于是
\begin{equation*}
\frac{E^{2}+(\partial_{\tau} r)^{2}}{1+2M/r-\Lambda r^{2}/3}+\frac{L^{2}}{r^{2}}=1,
\end{equation*}
或
\begin{equation*}
(\partial_{\tau} r)^{2}=\left(1-\frac{L^{2}}{r^{2}}\right)\left(1+\frac{2M}{r}-\frac{\Lambda r^{2}}{3}\right)-E^{2}.
\tag{1}
\end{equation*}

于是我们得到一个联系 $r$ 的变化率（相对于测试粒子测得的时间 $\tau$）与粒子 $r$ 坐标的方程。对它关于 $\tau$ 求导（再把两边除以 $2\partial_{\tau}r$），得到二阶导数
\begin{equation*}
\partial_{\tau}^{2} r=\frac{L^{2}(r+3M)}{r^{4}}-\frac{M}{r^{2}}-\frac{\Lambda r}{3}.
\tag{2}
\end{equation*}

第一项代表离心加速度，对给定角动量 $L$ 而言是一个近似反三次方的斥力。在洛伦兹版本中此处分子为 $r-3M$，故该项在 $r=3M$ 处变号，这正是黑洞周围不存在比此更近的轨道的原因。但在这里，离心力处处都是离心力，而且总比洛伦兹与牛顿版本更强。

第二项正是反平方的引力吸引，由此可以看出，若要引力具有通常的牛顿极限，我们把 $M$ 的符号取对是正确的。

第三项是源于宇宙学常数的吸引或排斥，它在大尺度上依 $\Lambda$ 的符号增强或削弱引力。

\orthofig{images/06__001.png}{黎曼 Schwarzschild 度规的有效势}

右图给出测试粒子径向运动在 $\Lambda=0$、不同 $L$ 取值下的有效势。函数
\begin{equation*}
V(r)=\left(\frac{L^{2}}{r^{2}}-1\right)\left(1+\frac{2M}{r}\right)
\end{equation*}
通过方程 (1) 决定径向运动，后者变为
\begin{equation*}
(\partial_{\tau} r)^{2}+V(r)=-E^{2}.
\end{equation*}

对每个非零的 $L$，有效势 $V$ 都有一个极小值；$E^{2}>1$ 且 $L\neq0$ 的粒子将处于稳定的束缚轨道，在 $\partial_{\tau}r=0$（即 $V(r)=-E^{2}$）所对应的两个 $r$ 值之间来回运动。图中那条红色虚线给出 $E^{2}=1.3$、$L=M$ 的一个例子。

若 $E^{2}<1$（例如粒子从无穷远处入射的情形），有效势在 $r\to\infty$ 处低于 $-E^{2}$，粒子将不被束缚，从无穷远驰入、在最近心点折返后逃逸回无穷远——这正是下一小节引力散射所讨论的轨道。

\subsection{圆轨道（Circular Orbits）}

假设测试粒子绕中心质量作圆轨道运动，则 $r=R$ 为常数，$r$ 的一阶与二阶导数都为零，由此得到关于 $L^{2}$ 与 $E^{2}$ 的两个方程（以轨道半径 $R$ 表出）。解出后代入四维速度 $\bu$，便可确定测试粒子相对于同一距离处静止观测者的普通速度。结果为
\begin{equation*}
v^{2}=\frac{3M+\Lambda R^{3}}{3R+6M-\Lambda R^{3}}.
\end{equation*}
对 $\Lambda=0$ 的情形，这变为
\begin{equation*}
v^{2}=\frac{M}{R+2M}.
\end{equation*}

当 $R$ 很大且 $\Lambda=0$，上式给出牛顿极限，其动能是引力势能大小的二分之一。

在牛顿物理中，绕行天体越靠近中心质量，其速度就越无界增大。而黎曼轨道则不同：当 $R$ 变小时，所有轨道的 $v^{2}$ 都趋于 $1/2$。由于离心力更强（在 $r$ 小时呈反四次方特征），$1/\sqrt2$ 就是任何物体为束缚于圆轨道所能具有的最快速度（至少就近距离而言）。

若存在宇宙学常数，足够大的轨道会有修正。若 $\Lambda$ 为正，引力在大尺度上变强，$v^{2}$ 在降到极小之后又开始无界上升。若 $\Lambda$ 为负，宇宙学项附加的排斥最终会使得无论物体运动多慢，都无法再被引力束缚于轨道。

\orthofig{images/06__002.png}{黎曼圆轨道的速度平方}

\subsection{径向运动与逃逸速度（Radial Motion and Escape Velocity）}

假设测试粒子径直朝向或背离中心质量运动，则 $L=0$，我们可以将另一个运动常数 $E^{2}$ 用它在某个特定 $r$ 处的速度（由同一距离处静止观测者判定）表出：
\begin{equation*}
E^{2}=\frac{1+2M/r-\Lambda r^{2}/3}{1+v^{2}},
\end{equation*}
对 $\Lambda=0$ 变为
\begin{equation*}
E^{2}=\frac{1+2M/r}{1+v^{2}}.
\end{equation*}

在 $\Lambda=0$ 情形中，一个恰好能逃逸到无穷远的粒子的 $E^{2}$，就是当 $r\to\infty$ 时 $v\to0$ 所取的值，即 $1$。因此粒子的逃逸速度为
\begin{equation*}
v_{esc}^{2}=\frac{2M}{r}.
\end{equation*}
这与逃逸速度的牛顿形式一致，也清楚地表明不存在无法逃脱的距离。

在更一般的 $\Lambda\neq0$ 情形中，"逃逸速度"其实不是个有意义的概念，因为宇宙学项在大距离上要么排斥、要么吸引粒子。但无论如何，中心质量本身永远不是无法逃脱的——黎曼宇宙没有黑洞。

\subsection{引力蓝移与时间膨胀（Gravitational Blueshift and Time Dilation）}

若一束离中心质量很远的、以速度 $v_{\infty}$ 接近它的光（假设 $\Lambda=0$，使 $r\to\infty$ 处的行为有良好定义），当它到达距离 $r$ 处时必定运动得更快，且有
\begin{align*}
\frac{1}{1+v_{\infty}^{2}} &= \frac{1+2M/r}{1+v(r)^{2}},\\
v(r)^{2} &= \left(1+\frac{2M}{r}\right)v_{\infty}^{2}+\frac{2M}{r}.
\end{align*}

更快的光波长更短，用《正交》（\emph{Orthogonal}）的语言说就是蓝移（blue-shifted）。但与洛伦兹情形不同，这一更短的波长会伴随一个\emph{更长}的周期 $\tau$，光的每个周期在大质量天体附近测得耗时更长。具体地说，
\begin{equation*}
\tau(r)=\tau_{\infty}\sqrt{1+\frac{2M}{r}}.
\end{equation*}

由此可知，一个在空间中绕行、接近了大质量天体的人，总体上会比绕开这段路的人经历更长的时间。与平直四维空间中运动引起的其它时间膨胀效应一样，黎曼宇宙中的引力时间膨胀意味着旅行者经历\emph{更多}时间，而非更少。

\subsection{引力散射（Gravitational Scattering）}

\orthofig{images/06__003.png}{黎曼引力散射}
\orthofig{images/06__004.png}{黎曼引力散射}

在牛顿引力物理中，一个从远处向大质量天体运动、且偏移得足以避免碰撞的测试粒子，将沿一条朝中心质量弯曲的双曲轨道运动。当粒子再次逃逸到无穷远时，其运动方向已经改变；我们说粒子被\textbf{散射}（scattered）了某个角度，记作 $\delta$。

在洛伦兹引力中类似的轨道是常态，但在黑洞的极端情形下，粒子可以在逃逸或被俘获之前绕黑洞数圈。

在黎曼引力中不存在黑洞之类的东西，但当相遇涉及大质量致密天体或高速粒子时，偏转是\emph{背离}天体而非朝向天体。

右图给出测试粒子在各种初速度 $v$、以及两个碰撞参数（impact parameter）$b$ 取值下的轨迹；$b$ 定义为轨迹相对于正碰的侧向位移（在粒子位于无穷远时测量）。上图中红色与橙色曲线——对应运动相当缓慢、且不会太靠近大质量天体的粒子——酷似牛顿双曲线，同一图中黄色曲线也仅略微偏离其原方向，但其它所有轨迹都被推离中心质量。

把这些粒子甩开的是一个非牛顿的离心力。正如我们在方程 (2) 中看到的，对小的 $r$，给定角动量 $L$ 下的离心加速度大致正比于 $r$ 的反四次方，而牛顿物理中是反三次方。

\section{进阶：宇宙学（Cosmology）}

\subsection{曲率与熵（Curvature and Entropy）}

从 Einstein 方程出发，我们可以导出用应力--能量张量表示 Ricci 张量的有用表达式。Einstein 张量 $G$ 用 Ricci 张量 $R$ 与度规 $g$ 定义为
\begin{equation*}
G=R-\tfrac{1}{2}R_{c}^{c}g,
\end{equation*}
此处用 Einstein 求和约定取 Ricci 张量的迹。把 Einstein 方程写成分量形式，有
\begin{align*}
G_{ab}+\Lambda g_{ab} &= 8\pi T_{ab},\\
R_{ab}-\tfrac{1}{2}R_{c}^{c}g_{ab}+\Lambda g_{ab} &= 8\pi T_{ab}.
\end{align*}

应力--能量张量的迹由此可求得为
\begin{align*}
8\pi T_{e}^{e} &= R_{e}^{e}-\tfrac{1}{2}R_{c}^{c}g_{e}^{e}+\Lambda g_{e}^{e}.
\end{align*}
黎曼度规的迹 $g_{e}^{e}$ 恰为空间的维数，即 $4$。于是
\begin{align*}
8\pi T_{e}^{e} &= R_{e}^{e}-2R_{c}^{c}+4\Lambda = -R_{e}^{e}+4\Lambda.
\end{align*}

[当然，对四维\emph{洛伦兹}空间，采用 $(-+++)$ 号差约定时 $g_{e}^{e}=2$。故上式及其后诸式\emph{不}适用于洛伦兹宇宙。]

利用这一点，可以写出
\begin{align*}
8\pi\bigl(T_{ab}-\tfrac{1}{2}T_{e}^{e}g_{ab}\bigr) &= 8\pi T_{ab}+\left[\tfrac{1}{2}R_{e}^{e}-2\Lambda\right]g_{ab}.
\end{align*}
再把 Einstein 方程给出的 $8\pi T_{ab}$ 代入，便得
\begin{align*}
8\pi\bigl(T_{ab}-\tfrac{1}{2}T_{e}^{e}g_{ab}\bigr)
&= R_{ab}-\tfrac{1}{2}R_{c}^{c}g_{ab}+\Lambda g_{ab}
 +\left[\tfrac{1}{2}R_{e}^{e}-2\Lambda\right]g_{ab} \nonumber\\
&= R_{ab}-\Lambda g_{ab}.
\end{align*}

这就给出用应力--能量张量与度规表示的 Ricci 张量：
\begin{equation*}
R_{ab}=8\pi\bigl(T_{ab}-\tfrac{1}{2}T_{e}^{e}g_{ab}\bigr)+\Lambda g_{ab}.
\tag{3}
\end{equation*}

整理我们此前得到的、联系 Ricci 张量与应力--能量张量之迹的结果：
\begin{equation*}
R_{c}^{c}=-8\pi T_{c}^{c}+4\Lambda.
\tag{4}
\end{equation*}

除 Ricci 张量的迹之外，我们还关心其各个对角元。若在局域正交归一坐标中工作，使每个 $a$ 都有 $g_{aa}=1$，则有（重复低标对不作求和）
\begin{equation*}
R_{aa}=8\pi\bigl(T_{aa}-\tfrac{1}{2}T_{c}^{c}\bigr)+\Lambda.
\tag{5}
\end{equation*}

$R_{aa}$ 告诉我们的是：一束初始平行的、沿 $a$ 方向排列的测地线将如何聚拢或散开——即它们所围体积（用垂直于 $a$ 方向的超曲面切片）的二阶变化率为正、为负还是为零。更精确地说：
\begin{equation*}
\frac{\partial_{a}\partial_{a}V}{V}=-R_{aa}.
\end{equation*}

大致说来，我们可以把黎曼宇宙中的物质分成两大范畴。在\textbf{低熵态}（low-entropy state）中，大多数物质沿大致相同的方向运动（一块固体、一团低负温度气体、或一个引力束缚系统），若把时间坐标 $t$ 与这些物质的世界线对齐，应力--能量张量将有较大的 $T_{tt}$，其余分量可忽略。一个例子是密度为 $\rho$、压强可忽略的尘埃云：
\begin{equation*}
T=\rho\,\bu\otimes\bu,
\end{equation*}
在尘埃云静止的正交归一参考系中 $T_{tt}=\rho$，其余分量为零。于是
\begin{align*}
T_{c}^{c} &= \rho, \\
R_{tt} &= 4\pi\rho+\Lambda, \\
R_{xx}=R_{yy}=R_{zz} &= -4\pi\rho+\Lambda, \\
R_{c}^{c} &= -8\pi\rho+4\Lambda.
\end{align*}

若 $\Lambda=0$，则 $R_{tt}>0$，于是沿 $t$ 方向排列的一束测地线所围体积会收缩。换言之，对初始相对尘埃云静止的测试粒子团，引力是吸引的，会把它们聚拢。而 $R_{xx}$ 等则为负，所以以无穷大速度穿过尘埃云的测试粒子团会散开。Ricci 张量的迹、即标量曲率，将为负。

一个足够大的正 $\Lambda$ 会使测试粒子团无论朝何方向都收缩，并给出正的标量曲率；而足够负的 $\Lambda$ 则可能反转所有这些结果。

那么，若物质处于\textbf{高熵态}（high-entropy state）、其世界线在四维空间中指向任何方向的概率都相等，又会怎样？在这种态下，
\begin{equation*}
T_{tt}=T_{xx}=T_{yy}=T_{zz}=\rho.
\end{equation*}
黎曼宇宙中的高熵物质表现得就像一个宇宙学常数：其应力--能量张量就是度规的一个倍数。把 $\rho$ 与 $\Lambda$ 的效应合并，得
\begin{align*}
T_{c}^{c} &= 4\rho, \\
R_{tt}=R_{xx}=R_{yy}=R_{zz} &= -8\pi\rho+\Lambda, \\
R_{c}^{c} &= -32\pi\rho+4\Lambda.
\end{align*}

若 $\Lambda$ 为零，则标量曲率为负，各个 Ricci 曲率 $R_{aa}$ 也都为负，于是沿\emph{任意}方向排列的一束测地线都会散开。若 $\rho$ 在四维空间中为常数，那会使黎曼宇宙成为一个四维双曲流形。

一个足够大的正 $\Lambda$ 可以克服 $\rho$ 的效应而给出正曲率，而负的 $\Lambda$ 则相当于增大 $\rho$、加剧负曲率。

\subsection{简单模型（Simple Models）}

\subsubsection{大爆炸模型（Big Bang Models）}

黎曼宇宙学中看来\emph{有可能}存在大爆炸奇点。但与洛伦兹物理（那里奇点似乎不可避免）不同，它只在给宇宙物质强加某种对称性时才会发生。

若假设一个完全均匀且各向同性的空间宇宙——形式为 3-球面、欧氏空间或双曲空间，其长度尺度 $a(t)$ 随时间变化（如 Misner、Thorne 和 Wheeler\textsuperscript{[3]}对洛伦兹情形所发展的那样），则 Einstein 方程成为关于 $a(t)$ 的微分方程：

\begin{keybox}[均匀各向同性宇宙的动力学]
（Dynamics of Homogeneous, Isotropic Universe）

\begin{gather*}
k=1\ \text{对应 3-球面};\qquad
k=0\ \text{对应欧氏空间};\qquad
k=-1\ \text{对应双曲空间}. \notag\\[4pt]
3\,\frac{a'(t)^{2}}{a(t)^{2}}-3\,\frac{k}{a(t)^{2}}+\Lambda=8\pi\rho\bigl(a(t)\bigr)
\qquad(\text{黎曼}), \\
3\,\frac{a'(t)^{2}}{a(t)^{2}}+3\,\frac{k}{a(t)^{2}}-\Lambda=8\pi\rho\bigl(a(t)\bigr)
\qquad(\text{洛伦兹}).
\end{gather*}
\end{keybox}

这里 $\rho\bigl(a(t)\bigr)$ 是质量--能量密度，作为长度尺度 $a(t)$ 的函数。

注意与洛伦兹版本相比，黎曼版本中 $\Lambda$ 与 $k$ 的符号都换了。在洛伦兹宇宙学中，暂时抛开 $\Lambda$ 的效应、假设它为零，最著名的结果之一是：当空间宇宙为 3-球面时，它会膨胀到最大尺寸后重新收缩；而若它是平坦或双曲的，则会永远膨胀下去。由于黎曼版本中 $k$ 的符号相反，其它条件相同时结论正好相反：3-球面空间宇宙会永远膨胀下去，而平坦或双曲的则会膨胀到最大尺寸后重新收缩。

在黎曼情形中能量密度 $\rho$ 的形式会有什么不同吗？在洛伦兹情形中，一种常见的粗略近似是谈论"物质"与"辐射"各自对 $\rho$ 贡献一项。物质被认为压强相对其静止质量可忽略，其随长度的标度为
\begin{equation*}
\rho_{m}a^{3}=\text{常数}.
\end{equation*}
换言之，空间宇宙任一给定区域中（由 3-球面上的角坐标或其它情形中的等价物描绘的边界所划出）有一个常数质量，该区域缩小时密度增大、增大时密度减小。各向同性辐射总是施加等于其能量密度三分之一的压强，由此可知当给定体积 $V$ 的空间区域膨胀时，辐射做功并损失能量：
\begin{equation*}
\frac{\dd(\rho_{r}V)}{\dd\tau}=-p_{r}\frac{\dd V}{\dd\tau}=-\frac{1}{3}\rho_{r}\frac{\dd V}{\dd\tau}.
\end{equation*}
取 $V=a^{3}$，此式由下式满足：
\begin{equation*}
\rho_{r}a^{4}=\text{常数}.
\end{equation*}
所以当 $a$ 小时，辐射的能量密度占主导。而当 $a\to0$ 时，能量密度的两个分量当然都趋于无穷。

在黎曼物理中，既然不存在零矢量（null vector），物质与辐射之间就没有同样深刻的区分。不过，像我们在上一节中描述的那样，区分低熵态与高熵态仍然有用。低熵物质的压强相对其静止质量可忽略，因此其贡献正如洛伦兹分析中的 $\rho_{m}$。高熵物质的压强与其静止质量相当，更像洛伦兹的辐射，但有两个关键区别。一是不存在 $1/3$ 这一因子；在真正 SO(4) 不变的态中（例如我们在无限温度稀薄气体中见过的），三个空间方向上的压强都\emph{等于}能量密度。二是压强下膨胀所做的功在黎曼情形中与洛伦兹情形相反（例如，一个黎曼系统对抗活塞膨胀时会给活塞动能，这意味着系统本身获得总能量）。于是
\begin{equation*}
\frac{\dd(\rho_{h}V)}{\dd\tau}=p_{h}\frac{\dd V}{\dd\tau}=\rho_{h}\frac{\dd V}{\dd\tau},
\end{equation*}
此式由下式满足：
\begin{equation*}
\rho_{h}=\text{常数}.
\end{equation*}
所以高熵物质的能量密度保持常数！细想之下这其实并不很令人惊讶：这种形式的物质在四维空间中根本没有任何偏好方向，因此它对我们选作 $t$ 坐标的那个方向完全无动于衷。

现在，若黎曼宇宙中均匀分布着某种\emph{低熵}物质，其世界线都正交于 3-球面或其它超曲面（由于均匀性，我们称后者为"某一瞬间的空间"），那么当我们沿这些世界线回溯到万物更紧密的时刻时，该物质的密度会趋于无穷，正如理想化自行车轮辐的密度在理想化轮毂处会趋于无穷一样确定。但在洛伦兹情形中已被证明，在某些相当合理的假设下\textsuperscript{[4]}这种奇点不可避免；而在黎曼情形中没有光锥来约束世界线，所以让它们全都汇聚到一个事件是一种\emph{选择}而非\emph{必需}——而且是一种相当人为的选择。诚然，很容易举出完全不含奇点的均匀黎曼宇宙。

\subsubsection{无奇点模型（Models Without Singularities）}

假设我们有一个充满物质的均匀宇宙，这些物质在宇宙学尺度上处于高熵、SO(4) 不变的状态。这并不妨碍任何局域的熵梯度与时间箭头，但宇宙中任何平行世界线束的取向都将指向四维空间中的随机方向。

我们可以把上一节的框架应用于此模型。我们应当能够把四维空间切片为均匀超曲面，它们或是 3-球面、或具有欧氏或双曲几何。能量密度将是某个常数 $\rho$。与我们所选超曲面正交的时间坐标 $t$ 的尺度因子 $a(t)$ 必须满足方程
\begin{equation*}
a'(t)^{2}-k=\pm K^{2}a(t)^{2},
\tag{6}
\end{equation*}
其中我们定义了一个正实常数 $K$，使得
\begin{equation*}
\pm K^{2}=\frac{8\pi\rho-\Lambda}{3}.
\end{equation*}
它给出物质与宇宙学常数的净效应，而参数 $k$ 告诉我们所拥有的是何种均匀超曲面：
\begin{equation*}
k=1\ \text{对应 3-球面};\qquad
k=0\ \text{对应欧氏空间};\qquad
k=-1\ \text{对应双曲空间}.
\end{equation*}

方程 (6) 有哪些可能的解？

\textbf{(A)} $8\pi\rho-\Lambda>0$。双曲四维空间，Ricci 标量 $R_{c}^{c}=-12K^{2}$。
\begin{itemize}
  \item $k=1$：$a(t)=\dfrac{\sinh(Kt)}{K}$，切片为 3-球面，其半径从 $0$ 开始增长，先是线性、随后指数增长。
  \item $k=0$：$a(t)=\exp(Kt)$，切片为平坦超曲面，随 $t$ 指数增长。
  \item $k=-1$：$a(t)=\dfrac{\cosh(Kt)}{K}=\dfrac{1}{2K}\bigl[\exp(Kt)+\exp(-Kt)\bigr]$，切片为双曲超曲面，随 $t$ 与 $-t$ 指数增长。
\end{itemize}

\textbf{(B)} $8\pi\rho-\Lambda=0$。欧氏四维空间，Ricci 标量 $R_{c}^{c}=0$。
\begin{itemize}
  \item $k=1$：$a(t)=t$，切片为 3-球面，随 $t$ 从半径 $0$ 线性增长（即 $t$ 是径向极坐标）。
  \item $k=0$：$a(t)=1$，切片为恒定尺度的平坦超曲面（即 $t$ 是笛卡儿坐标）。
  \item $k=-1$：无解。
\end{itemize}

\textbf{(C)} $8\pi\rho-\Lambda<0$。四维球面 $S^{4}$，Ricci 标量 $R_{c}^{c}=12K^{2}$。
\begin{itemize}
  \item $k=1$：$a(t)=\dfrac{\sin(Kt)}{K}$，切片为 3-球面，其半径从 $0$ 增长、达到最大后又重新收缩。
  \item $k=0$ 或 $-1$：无解。
\end{itemize}

于是四维空间的几何有三种可能：双曲、平坦或球面。有各种方式把这些四维空间切片为均匀的三维超曲面，但四维空间本身不受这一选择的影响。在那些宇宙从尺度因子为零开始和/或结束的情形中，这\emph{并不}意味着奇点！能量密度 $\rho$ 保持常数，尺度因子趋于零的位置只是我们把四维空间切成超曲面这一方式的产物。

如我们此前所讨论的，我们有理由偏好一个\emph{有限}的四维空间，而这里没有任何东西排除这一点。在正曲率情形中，整个四维空间必然是有限的；这可由 Myers 定理推出。然而，有些拓扑与正曲率不相容；例如，任何环面（torus）都不可能具有正标量曲率的度规\textsuperscript{[5]}。

对零曲率与负曲率，则既有有限空间也有无限空间；例如 $\mathbb{R}^{4}$（无限四维空间）与 $T^{4}$（四维环面）都可以赋予均匀平坦的度规。不过对拓扑仍有某些限制。Cartan--Hadamard 定理意味着，一个负曲率的均匀各向同性空间不可能具有 4-球面的拓扑。

\section{进阶：引力波（Gravitational Waves）}

把引力波\textsuperscript{[6]}作为平直度规中的小扰动来处理的标准方法，可以像在洛伦兹版本中那样应用于黎曼广义相对论。所得方程（假设 $\Lambda=0$）为：

\begin{keybox}[引力波方程]
（Gravitational Wave Equations，$\Lambda=0$）

\begin{gather*}
g_{ab}=\delta_{ab}+h_{ab},\qquad
w_{ab}=h_{ab}-\tfrac{1}{2}h_{c}^{c}\,\delta_{ab}, \notag\\
\partial_{x}^{2}\mathbf{w}+\partial_{y}^{2}\mathbf{w}+\partial_{z}^{2}\mathbf{w}+\partial_{t}^{2}\mathbf{w}=-16\pi\mathbf{T}
\qquad(\text{黎曼}), \\[6pt]
g_{ab}=\eta_{ab}+h_{ab},\qquad
w_{ab}=h_{ab}-\tfrac{1}{2}h_{c}^{c}\,\eta_{ab}, \notag\\
\partial_{x}^{2}\mathbf{w}+\partial_{y}^{2}\mathbf{w}+\partial_{z}^{2}\mathbf{w}-\partial_{t}^{2}\mathbf{w}=-16\pi\mathbf{T}
\qquad(\text{洛伦兹}), \\[6pt]
\partial_{x}w^{ax}+\partial_{y}w^{ay}+\partial_{z}w^{az}+\partial_{t}w^{at}=0
\qquad(\text{横波条件，Transverse}).
\end{gather*}
\end{keybox}

我们先陈述：在我们所选坐标中度规分量 $g_{ab}$ 等于 Kronecker delta（$\delta_{ab}$ 在 $a=b$ 时为 $1$、$a\neq b$ 时为 $0$，即欧氏空间笛卡儿坐标中平坦度规的分量）加上某个小扰动 $h_{ab}$，后者像度规本身一样是对称张量。

在洛伦兹版本中，我们使用 Minkowski 度规 $\eta_{ab}$，它有 $\eta_{tt}=-1$，其余分量与 $\delta_{ab}$ 相同。

我们定义一个新的对称张量 $w_{ab}$，它等于 $h_{ab}$ 减去 $\mathbf{h}$ 的迹之半乘以相应的平直空间度规：黎曼情形用 $\delta$，洛伦兹情形用 $\eta$。[通常的记号是在原扰动 $\mathbf{h}$ 上加一横，但在网页上很难做到。] 当我们在张量 $h_{ab}$ 上抬升指标以取迹时，我们用的是 $\delta$ 或 $\eta$ 而非完整度规，因为我们只关心 $h$ 的一阶量。同理，在操作 $w_{ab}$ 的指标时我们也用平直空间度规。

这些定义的目的在于：在每种情形中，对扰动及其导数取一阶的 Einstein 方程都变成关于扰动 $\mathbf{w}$ 的简单波动方程（如上所示），只要 $\mathbf{w}$ 还满足一个横波条件——类似于我们此前对黎曼矢量波施加的那个条件。

在黎曼情形中，这个"波动方程"就是 Laplace 方程，任何 $\mathbf{T}=0$ 的非平凡解都将无界。$\mathbf{w}$ 的二阶导数之和必须为零，所以若解在某个方向上呈周期性（这给出等于原函数负倍数的二阶导数），它就必须在另一个方向上指数增长（给出等于原函数\emph{正}倍数的二阶导数）。若施加周期性边界条件，例如要求宇宙是 $T^{4}$，指数解将无法满足边界条件，唯一的解将是 $\mathbf{w}=\text{常数}$。

因此，在没有宇宙学常数、且假设我们为避免电磁波中指数解所需的那类边界条件下，黎曼宇宙中不可能有无源引力波。不仅如此，与一个振荡源相关的解会随离源的距离指数衰减；其理由与如下情形基本相同：在 $T^{4}$ 边界条件下，若交变电流的频率超过最高频率 $\nu_{max}$，其周围的场会随距离指数衰减——但在引力波这里，最高频率是零！

若有一个非零的宇宙学常数会怎样？看起来这似乎会把引力波方程从 Laplace 方程的张量版本变成更像 Helmholtz 方程的东西——引入波函数的一个常数倍，就像黎曼矢量波方程（Riemannian Vector Wave Equation）中的 $\wms^{2}$ 项那样，从而若 $\Lambda$ 符号合适，就能允许某些最高频率以内的周期解。

但结果并非如此。完整分析很复杂，但若把度规扰动取为相对于给定 $\Lambda$ 的真空解，则曲率相关（因而依赖 $\Lambda$）的项会以多种方式进入扰动方程，此外 Einstein 方程本身也明显含有 $\Lambda$。最终它们全部抵消，留下一个不含波函数常数倍的方程。

所以看起来，那个为驯服电磁波所需的边界条件，会把引力波完全排除。

\section{参考文献}

\begin{itemize}
  \item[{[1]}] \emph{Exact Solutions of Einstein's Field Equations}, Hans Stephani, Dietrich Kramer, Malcolm MacCallum, Cornelius Hoenselaers, Eduard Herlt, Cambridge University Press, Cambridge, 2003, 第 15.4 节。
  \item[{[2]}] \emph{Gravitation}, Charles Misner, Kip Thorne, John Wheeler, W.H. Freeman, San Francisco, 1973, 第 25.2、25.3 节。
  \item[{[3]}] Misner, Thorne 和 Wheeler, \emph{同上}, 第 27 章。
  \item[{[4]}] \emph{The Large Scale Structure of Space-time}, S. W. Hawking, G. F. R. Ellis, Cambridge University Press, Cambridge, 1973, 第 10 章。
  \item[{[5]}] "On the structure of manifolds with positive scalar curvature", R. Schoen, S. T. Yau, \emph{Manuscripta Mathematica}, 1979, 第 28 卷, 第 1--3 期, 第 159--183 页。
  \item[{[6]}] Misner, Thorne 和 Wheeler, \emph{同上}, 第 18、35 章。
\end{itemize}

\chapter{世界观与文明设定}

\section{导语：物理如何塑造文明}

前几章所确立的正交宇宙物理——一个没有普适光速、发光反而释放能量、时空度规处处正定的世界——并不只是一堆反直觉的理论习题；它同样决定了生活于其中的文明能做什么、不能做什么，以及他们如何理解自己的宇宙。本节先综述这些物理结论如何塑造文明形态，再以一个具体的运动学问题——\emph{软性拦截}（Soft Interception）——展示变分法（Calculus of Variations）在求解最优化物理过程中的运用。

其一，正交宇宙\emph{不存在普适光速}。不同的波长以不同的速度传播，红（长波）光比紫（短波）光更慢。这一点彻底改写了星空的面貌：一颗恒星发出的各种颜色不会同时抵达，于是每颗星在天空中拖出一条"色谱拖尾"（colour trail），其红、紫两端之间的张角直接度量该恒星的横向速度。对文明而言，天象观测、测距与导航都须按一套全然不同的规则去理解——光不再是一个统一而瞬时的信使。

其二，\emph{发光过程释放能量}。在我们所处的宇宙中，产生光需要消耗能量；而在正交宇宙中，由于能量—动量矢量的"上下颠倒"，发光体系的动能反而增加，表现为放热。这一反转最深刻的文化后果体现在生命本身：这里的植物不靠吸收光来进行光合作用，而是向漆黑一片的夜空发出自己的光，借此把能量转化为化学能、制造养分。夜与昼、黑与亮对这个文明的意义，与我们恰好相反。

其三，能源与推进。既然产生光会释放（而非消耗）能量，正交宇宙的飞行器无需携带任何燃料或推进剂，便能凭借\emph{发光引擎}持续地把能量转化为动能、长时间地加速。星际航行因此在原则上变得可行，而"燃料"与"推进剂"这类概念根本不会出现在这个文明的工程语汇里。

其四，威胁与世代航行。雅尔达的故乡正遭受一群被称作"伤人者"（Hurtlers）的奇异流星的侵袭，它们以空前的速度闯入行星系统，而应对这一威胁所需的知识与技术远超该文明已有的水平。唯一看似可行的出路，是以极高的速度派出一艘飞船远航：由于正交宇宙没有速度上限，飞船可以远超我们宇宙的"光速"行进；又由于黎曼几何中"绕道使时间更长"，这趟航行对船上的旅行者而言要跨越许多代人的岁月，对故乡却只过去短短数年。旅行者便有机会发现母星急需的科学，并在灾难降临之前把它带回。

需要说明的是，本章第二节所讨论的"软性拦截"并不是关于正交宇宙特有物理的问题。在该问题所涉及的（低速）速度范围内并不会出现相对论效应，正交宇宙的另类物理恰恰只在"引擎可以无燃料、无推进剂地发光加速"这一点上起作用；除此之外，我们可以沿用与我们宇宙低速情形相同的牛顿（Newtonian）力学规则。它因此成为一座很好的桥梁：既有正交宇宙的文明底色，又是两个宇宙共享的经典运动学。

\section{软性拦截（Soft Interception）}

在《时间之箭》（\emph{The Arrows of Time}）第三章中，拉米罗（Ramiro）与塔尔奎尼亚（Tarquinia）需要拦截一艘沿固定轨迹加速飞行的叛逃飞船（rogue spacecraft）。他们的目标是登上它，因此仅仅横穿其轨迹是不够的：他们还须与它\emph{速度匹配}。这就是所谓的\emph{软性拦截}（soft interception）。

既然他们可以让自己的引擎无限期地全功率运转，那么应当选择怎样的轨迹，才能在最短时间内完成这次拦截？考虑到涉及的速度并不高、不足以引发相对论效应，正交宇宙的另类物理并不会登场（除了使航天器的引擎成为可能——这种引擎靠发光来产生动能，且不需要任何燃料或推进剂）。因此，我们可以沿用与我们宇宙低速情形相同的牛顿力学规则。

这类问题——求一个使某物理量取极大或极小的数学函数——正是由\emph{变分法}（Calculus of Variations）处理的。在我们的情形里，所关心的函数将描述一条随时间穿过空间的轨迹，而要最小化的量，则是沿这条轨迹从追击飞行器（gnat）的初始位置走到叛逃飞船轨迹上某一点所需的时间。

\orthofig{images/E3__Fig002.png}{拦截叛逃飞船的飞行方案}

\subsection{问题陈述}

设拉米罗与塔尔奎尼亚的追击飞行器开始追击时的位置为 $\bp_0$、速度为 $\bv_0$，其引擎能产生的最大加速度大小为 $A$。叛逃飞船正沿一个固定方向做恒定加速；若把该加速度矢量记作 $\mathbf{a}$，则叛逃飞船随时间的轨迹为

\begin{equation}
\mathbf{r}(t) = \mathbf{r}_0 + \bw_0\,t + \half\,\mathbf{a}\,t^{2} .
\end{equation}

我们约定从追击开始的那一刻起计时，因此在 $t=0$ 时，追击飞行器位于位置 $\bp_0$、速度为 $\bv_0$，而叛逃飞船位于位置 $\mathbf{r}_0$、速度为 $\bw_0$。

对于追击飞行器的任一条候选轨迹 $\bp(t)$，我们提出如下要求：

\begin{itemize}
  \item $\bp(0)=\bp_0$：轨迹从正确的位置出发。
  \item $\bp'(0)=\bv_0$：轨迹以正确的速度出发。
  \item $\bp(T)=\mathbf{r}(T)$，在某个时刻 $T$：追击飞行器与叛逃飞船相遇。
  \item $\bp'(T)=\mathbf{r}'(T)$，对于同一个 $T$：两艘航天器速度匹配。
  \item 对所有 $0\le t\le T$ 都有 $|\bp''(t)|\le A$：追击飞行器在任何时刻都不被要求做出超出其能力的加速。
\end{itemize}

我们的目标，是在满足这些条件的前提下，找到使拦截时刻 $T$ 尽可能小的那条轨迹 $\bp(t)$。

为使分析稍稍简化，我们不加证明地假定：追击飞行器在整个航程中都应以最大推力工作，仅仅改变引擎喷推的\emph{方向}。直觉上，以尽可能大的加速度加速总能带来某种好处（更详尽的分析会证实这一点）。这使我们能把上面最后一个不等式条件改写成等式：

\begin{itemize}
  \item 对所有 $0\le t\le T$ 都有 $|\bp''(t)|=A$：追击飞行器始终以可能的最大速率加速，不多也不少。
\end{itemize}

这个约束也可以写成

\begin{equation}
\bp''(t)\cdot\bp''(t)-A^{2}=0 .
\end{equation}

\subsection{拉格朗日乘子}

先回顾一下普通微积分里的情形：要在某个约束（例如 $G(x,y,z)=0$）下极小化函数 $F(x,y,z)$，我们会用到所谓的\emph{拉格朗日乘子}（Lagrange multiplier）技巧。设想我们设法把约束方程 $G(x,y,z)=0$ 对其中一个变量（例如 $z$）解出，便可写出

\begin{equation}
f(x,y)=F\bigl(x,y,z(x,y)\bigr),
\end{equation}

再令 $f(x,y)$ 对 $x$ 与 $y$ 的偏导数为零来求其极小值：

\begin{align}
\partial_{x} f(x,y) &= \partial_{x}F + \partial_{z}F\,\partial_{x}z = 0, \\
\partial_{y} f(x,y) &= \partial_{y}F + \partial_{z}F\,\partial_{y}z = 0 .
\end{align}

我们也可以对 $G(x,y,z(x,y))$ 求关于 $x$、$y$ 的偏导数，由于 $z=z(x,y)$ 正是 $G(x,y,z)=0$ 的解，这些偏导数当然必须为零：

\begin{align}
\partial_{x}G + \partial_{z}G\,\partial_{x}z &= 0, \\
\partial_{y}G + \partial_{z}G\,\partial_{y}z &= 0 .
\end{align}

合起来一共有四个方程。我们可以从后两个方程解出 $\partial_{x}z$ 与 $\partial_{y}z$，再代回前两个方程，得到

\begin{align}
\partial_{x}F - \frac{\partial_{z}F}{\partial_{z}G}\,\partial_{x}G &= 0, \\
\partial_{y}F - \frac{\partial_{z}F}{\partial_{z}G}\,\partial_{y}G &= 0 .
\end{align}

而这恰恰也是下面三个方程所给出的结果：

\begin{align}
\partial_{x}F + \lambda\,\partial_{x}G &= 0, \\
\partial_{y}F + \lambda\,\partial_{y}G &= 0, \\
\partial_{z}F + \lambda\,\partial_{z}G &= 0,
\end{align}

只要把其中最后一个方程解出 $\lambda$，即得

\begin{equation}
\lambda = -\frac{\partial_{z}F}{\partial_{z}G} .
\end{equation}

因此，原问题等价于求

\begin{equation}
S(x,y,z,\lambda)=F(x,y,z)+\lambda\,G(x,y,z)
\end{equation}

对 $x$、$y$、$z$ 与 $\lambda$ 的偏导数全部为零的点；其中对 $\lambda$ 的偏导数恰好把我们带回原来的约束 $G(x,y,z)=0$。这个多出来的变量 $\lambda$ 便是拉格朗日乘子。

从几何上可以更清楚地理解这一点。函数 $F$ 的梯度为 $\nabla F=(\partial_{x}F,\partial_{y}F,\partial_{z}F)$，它描述了与 $F$ 的"等值面"（即 $(x,y,z)$ 空间中使 $F$ 取常数的点集）正交的一个矢量。类似地，$G$ 的梯度为 $\nabla G=(\partial_{x}G,\partial_{y}G,\partial_{z}G)$，它与 $G$ 取常数的点集正交——当然也包括满足约束方程 $G(x,y,z)=0$ 的那个点集。前面写下的、联系 $F$ 与 $G$ 的各偏导数的三个方程，可以用梯度归纳为

\begin{equation}
\nabla F = -\lambda\,\nabla G .
\end{equation}

若我们从约束点集中的某一点出发，任何\emph{始终留在}该点集内的路径，其切矢量都正交于 $\nabla G$。但若 $\nabla F$ 是 $\nabla G$ 的倍数，则同一个切矢量也正交于 $\nabla F$，于是我们同时也沿一条 $F$ 取常数的路径运动。因此在这样的点上，满足约束就保证了 $F$ 的导数为零——这正是局部极大或极小点附近必然成立的性质。

这可以轻易推广到任意多个约束函数 $G_1,G_2,G_3,\dots$（它们都必须为零）、以及任意多变量的空间。从满足所有这些约束的某一点出发，沿一条同时正交于 $\nabla G_1,\nabla G_2,\nabla G_3,\dots$ 的路径移动，就能使全部约束继续得到满足。现在假设在约束点集中的某一点处有

\begin{equation}
\nabla F = -\lambda_{1}\nabla G_{1} - \lambda_{2}\nabla G_{2} - \lambda_{3}\nabla G_{3} - \cdots .
\end{equation}

由于在该点穿过约束点集的任一条路径，其切矢量都正交于每个 $\nabla G_i$，同一个切矢量也就正交于 $\nabla F$。这意味着 $F$ 沿任何此类路径的导数都为零，从而该点就是 $F$ 在约束 $G_1=G_2=G_3=\cdots=0$ 下取局部极大或极小的候选点。

\subsection{连续约束下的变分}

回到原问题。我们要在追击飞行器的整条轨迹上施加约束

\begin{equation}
\bp''(t)\cdot\bp''(t)-A^{2}=0 .
\end{equation}

若把 $t$ 的每个取值都看作需要一个单独的约束，就须对这些不同时刻，把"单独的拉格朗日乘子 $\times$ 相应的约束"之积全部"加起来"。但对于这样一段连续的取值，我们用积分而非求和。于是我们写下一个积分，其中既包含沿轨迹流逝的时间（即我们真正想最小化的那个函数），也包含拉格朗日乘子乘以约束的项——此时的拉格朗日乘子已是时间的函数：

\begin{equation}
I(\bp,\lambda)=\int_{0}^{T(\bp)}\Bigl[\,1+\lambda(t)\bigl(\bp''(t)\cdot\bp''(t)-A^{2}\bigr)\Bigr]\,\dd t .
\end{equation}

这里，沿轨迹对 $1$ 积分给出总历时。我们把积分上限记作 $T(\bp)$，以强调追击者拦截到叛逃飞船的时刻取决于所选的轨迹。$T(\bp)$ 的实际取值就是下式的第一个正解：

\begin{equation}
\bp(T)=\mathbf{r}(T).
\end{equation}

我们的目标是找到一条在满足全部要求的同时、使拦截时间最小的轨迹 $\bp_M(t)$；为此，我们要找到 $\bp_M(t)$ 以及相应的拉格朗日乘子 $\lambda_M(t)$，使得积分 $I(\bp,\lambda)$ 对两个函数中任一个的任意合适变分的导数都为零。所谓\emph{合适}的变分，是指它能让轨迹继续满足必要的边界条件：它必须仍具有相同的初始位置与速度，并且仍须在与叛逃飞船相遇之处匹配其位置与速度。我们将考察如下形式的轨迹：

\begin{equation}
\bp(t)=\bp_M(t)+\varepsilon\,\bq(t),
\end{equation}

并在 $t=0$ 处施加限制

\begin{equation}
\bq(0)=0,\qquad \bq'(0)=0 ,
\end{equation}

以满足边界条件。由于 $\bp_M$（依定义）已具有正确的初始位置与速度，这些限制便保证了 $\bp_M+\varepsilon\bq$ 同样如此。在 $t=T$ 处，我们要求

\begin{align}
\bp_M(T)+\varepsilon\,\bq(T) &= \mathbf{r}(T), \\
\bp_M'(T)+\varepsilon\,\bq'(T) &= \mathbf{r}'(T).
\end{align}

请记住，$T$ 会随我们改变轨迹而依赖于轨迹，因此这里的第一个方程并不是对 $\bq$ 的限制，而是\emph{确定} $T$ 取值的手段——$T$ 就是这个方程的第一个正解。

若把这两个方程对 $\varepsilon$ 求导、再令 $\varepsilon$ 等于零（并记 $T_M\equiv T(\bp_M)$ 作为简写），便得到

\begin{align}
\bq(T_M) &= \bigl[\mathbf{r}'(T_M)-\bp_M'(T_M)\bigr]\left.\frac{\dd T}{\dd\varepsilon}\right|_{\varepsilon=0} = 0, \\
\bq'(T_M) &= \bigl[\mathbf{r}''(T_M)-\bp_M''(T_M)\bigr]\left.\frac{\dd T}{\dd\varepsilon}\right|_{\varepsilon=0}
= \bigl[\mathbf{a}-\bp_M''(T_M)\bigr]\left.\frac{\dd T}{\dd\varepsilon}\right|_{\varepsilon=0}.
\end{align}

我们还要对拉格朗日乘子作变分，令

\begin{equation}
\lambda(t)=\lambda_M(t)+\varepsilon\,\mu(t).
\end{equation}

现在施加条件：我们的积分对 $\varepsilon$ 的导数为零，

\begin{equation}
\left.\frac{\dd}{\dd\varepsilon} I\bigl(\bp_M+\varepsilon\bq,\ \lambda_M+\varepsilon\mu\bigr)\right|_{\varepsilon=0}=0 ,
\end{equation}

即

\begin{equation}
\left.\frac{\dd}{\dd\varepsilon}\int_{0}^{T(\bp_M+\varepsilon\bq)}
\Bigl[\,1+\bigl[\lambda_M(t)+\varepsilon\mu(t)\bigr]\bigl([\bp_M''(t)+\varepsilon\bq''(t)]\cdot[\bp_M''(t)+\varepsilon\bq''(t)]-A^{2}\bigr)\Bigr]\,\dd t\right|_{\varepsilon=0}=0 .
\end{equation}

展开后得

\begin{equation}
\left.\frac{\dd T}{\dd\varepsilon}\right|_{\varepsilon=0}\Bigl[1+\lambda_M(T)\bigl(\bp_M''(T)\cdot\bp_M''(T)-A^{2}\bigr)\Bigr]
+\int_{0}^{T_M}\!\Bigl[2\lambda_M(t)\,\bp_M''(t)\cdot\bq''(t)+\mu(t)\bigl(\bp_M''(t)\cdot\bp_M''(t)-A^{2}\bigr)\Bigr]\,\dd t=0 .
\end{equation}

由于 $\bp_M$ 依定义必须在各处都满足"加速度大小等于 $A$"这一约束，上式化为

\begin{equation}
\left.\frac{\dd T}{\dd\varepsilon}\right|_{\varepsilon=0}
+2\int_{0}^{T_M}\lambda_M(t)\,\bp_M''(t)\cdot\bq''(t)\,\dd t=0 .
\end{equation}

我们可以用分部积分法变换这个积分，其一般法则是

\begin{equation}
\int_{a}^{b}u(t)\,v'(t)\,\dd t
= u(t)\,v(t)\Big|_{a}^{b}-\int_{a}^{b}u'(t)\,v(t)\,\dd t .
\end{equation}

把它用两次，将关于 $t$ 的导数从 $\bq''(t)$ 移到 $\lambda_M(t)\,\bp_M''(t)$ 上，得到

\begin{equation}
\left.\frac{\dd T}{\dd\varepsilon}\right|_{\varepsilon=0}
+2\Biggl[
\lambda_M(t)\,\bp_M''(t)\cdot\bq'(t)\Big|_{0}^{T_M}
-\frac{\dd}{\dd t}\bigl(\lambda_M(t)\,\bp_M''(t)\bigr)\cdot\bq(t)\Big|_{0}^{T_M}
+\int_{0}^{T_M}\frac{\dd^{2}}{\dd t^{2}}\bigl(\lambda_M(t)\,\bp_M''(t)\bigr)\cdot\bq(t)\,\dd t
\Biggr]=0 .
\end{equation}

关于 $\bq$ 的边界条件给出 $\bq(0)=0$、$\bq'(0)=0$，以及 $\bq(T_M)=0$。代入其中唯一非零的边界条件 $\bq'(T_M)=[\mathbf{a}-\bp_M''(T_M)]\left.\dfrac{\dd T}{\dd\varepsilon}\right|_{\varepsilon=0}$，得到

\begin{equation}
\left.\frac{\dd T}{\dd\varepsilon}\right|_{\varepsilon=0}
\Bigl(1+2\lambda_M(T_M)\,\bp_M''(T_M)\cdot\bigl[\mathbf{a}-\bp_M''(T_M)\bigr]\Bigr)
+2\int_{0}^{T_M}\frac{\dd^{2}}{\dd t^{2}}\bigl(\lambda_M(t)\,\bp_M''(t)\bigr)\cdot\bq(t)\,\dd t=0 .
\end{equation}

这个方程需要对所有变分 $\bq$ 都成立。不难看出：只要轨迹的加速度与拉格朗日乘子之积满足微分方程

\begin{equation}
\frac{\dd^{2}}{\dd t^{2}}\bigl(\lambda_M(t)\,\bp_M''(t)\bigr)=0,
\end{equation}

我们就能对任何 $\bq$ 使该积分为零。若满足这个方程，就还需要另外保证边界项单独为零。看出这一点的最简便办法，就是把这个方程直接解出来：

\begin{equation}
\lambda_M(t)\,\bp_M''(t)=\mathbf{C}\,t+\mathbf{D},
\end{equation}

其中 $\mathbf{C}$ 与 $\mathbf{D}$ 是两个常矢量。由于轨迹 $\bp_M$ 具有恒定的加速度大小 $A$，我们有

\begin{align}
|\lambda_M(t)|\,|\bp_M''(t)| &= |\mathbf{C}\,t+\mathbf{D}|, \\
|\lambda_M(t)|\,A &= |\mathbf{C}\,t+\mathbf{D}|, \\
\lambda_M(t) &= \pm\frac{|\mathbf{C}\,t+\mathbf{D}|}{A}, \\
\bp_M''(t) &= \pm\frac{A\,(\mathbf{C}\,t+\mathbf{D})}{|\mathbf{C}\,t+\mathbf{D}|}.
\end{align}

若我们从某一对 $\mathbf{C}\,t+\mathbf{D}$ 出发、再把这二者同乘一个常数，这一因子会在 $\bp_M''(t)$ 的解中完全消去，因此丝毫不影响我们选出满足问题各项具体要求（初态匹配追击飞行器的位置与速度、拦截点匹配叛逃飞船）的轨迹的能力。现在，把我们的解代回边界项，并略去整体因子 $\left.\dfrac{\dd T}{\dd\varepsilon}\right|_{\varepsilon=0}$，得到

\begin{align}
&1+2\lambda_M(T_M)\,\bp_M''(T_M)\cdot\bigl[\mathbf{a}-\bp_M''(T_M)\bigr] \nonumber\\
&\qquad = 1+2\,(\mathbf{C}\,T_M+\mathbf{D})\cdot\Bigl[\mathbf{a}\mp\frac{A\,(\mathbf{C}\,T_M+\mathbf{D})}{|\mathbf{C}\,T_M+\mathbf{D}|}\Bigr] \nonumber\\
&\qquad = 1+2\Bigl[(\mathbf{C}\,T_M+\mathbf{D})\cdot\mathbf{a}\mp A\,|\mathbf{C}\,T_M+\mathbf{D}|\Bigr].
\end{align}

若引入简写

\begin{equation}
\mathbf{c}=\mathbf{C}\,T_M+\mathbf{D},
\end{equation}

并作替换 $\mathbf{c}\to\alpha\,\mathbf{c}$，边界项变为

\begin{equation}
1+2\bigl[\alpha\,\mathbf{c}\cdot\mathbf{a}\mp A\,|\alpha|\,|\mathbf{c}|\bigr].
\end{equation}

若假定 $\alpha$ 为正，则可通过令

\begin{equation}
\alpha=\frac{1}{2\bigl(\pm A\,|\mathbf{c}|-\mathbf{c}\cdot\mathbf{a}\bigr)}
\end{equation}

使其为零，其中还要对符号作恰当的选择（回忆一下：这个符号决定了轨迹的加速度与方向 $\mathbf{C}\,t+\mathbf{D}$ 之间的整体关系）。我们预期追击飞行器在拦截时刻的加速方向几乎与叛逃飞船一致，因此若 $\mathbf{c}\cdot\mathbf{a}$ 为正，取正号就是恰当的。既然必须满足 $A>|\mathbf{a}|$ 才能让追击飞行器追上叛逃飞船，右端的分母便为正，从而给出我们假定的正 $\alpha$。若 $\mathbf{c}\cdot\mathbf{a}$ 为负，我们只需把最初对 $\mathbf{C}$ 与 $\mathbf{D}$ 的选择取反即可。

于是，只要轨迹的加速度取如下形式，我们就能使边界项与积分同时为零：

\begin{equation}
\bp_M''(t)=A\,\frac{\mathbf{C}\,t+\mathbf{D}}{|\mathbf{C}\,t+\mathbf{D}|}.
\end{equation}

\subsection{最优轨迹的显式构造}

为求出上式的通解，我们先从一个特解入手。（以下部分内容改编自文献 [2]，那篇文献处理的是棒球场上各垒之间跑动的最快路径问题。）

\emph{双曲正弦}函数 $\sinh$ 定义为

\begin{equation}
\sinh(x)=\half\bigl(\ee^{x}-\ee^{-x}\bigr).
\end{equation}

它的反函数 $\operatorname{arcsinh}$，可以通过对 $\ee^{x}$ 解一个二次方程再取对数求得：

\begin{align}
y &= \sinh(x), \\
2y &= \ee^{x}-\ee^{-x}, \\
(\ee^{x})^{2}-2y\,\ee^{x}-1 &= 0, \\
\operatorname{arcsinh}(y) = x &= \ln\bigl(y+\sqrt{1+y^{2}}\bigr).
\end{align}

自然对数的导数就是其反函数，所以 $\operatorname{arcsinh}$ 的导数为

\begin{equation}
\operatorname{arcsinh}'(y)=\frac{1+y/\sqrt{1+y^{2}}}{y+\sqrt{1+y^{2}}}=\frac{1}{\sqrt{1+y^{2}}} .
\end{equation}

这使我们能找到一个简单的矢量函数，其一阶导数具有恒定的模长 $A$：

\begin{align}
\mathbf{V}(t) &= \mathbf{V}_0+A\bigl(\operatorname{arcsinh}(t),\,\sqrt{1+t^{2}}\bigr), \\
\mathbf{V}'(t) &= A\left(\frac{1}{\sqrt{1+t^{2}}},\,\frac{t}{\sqrt{1+t^{2}}}\right)=A\,\frac{(1,t)}{|(1,t)|}, \\
|\mathbf{V}'(t)| &= A .
\end{align}

再多花一点功夫，即可求出这个速度的积分：

\begin{equation}
\mathbf{X}(t)=\mathbf{X}_0+\mathbf{V}_0\,t+A\Bigl(t\,\operatorname{arcsinh}(t)-\sqrt{1+t^{2}},\ \half\bigl(\operatorname{arcsinh}(t)+t\sqrt{1+t^{2}}\bigr)\Bigr).
\end{equation}

为把它推广，我们可作替换 $t\to\beta\,t+\gamma$，同时把 $A$ 换成 $A/\beta^{2}$ 以保持加速度大小不变。我们还自由地把坐标旋转任意角度 $\theta$。这样便留下八个可以调整的参数：来自矢量 $\mathbf{X}_0$ 与 $\mathbf{V}_0$ 的四个分量，外加 $\beta$、$\gamma$、$\theta$ 与拦截时刻 $T$。在 $t=0$ 处需要满足四个方程，以拟合追击飞行器的初始位置与速度；在 $t=T$ 处还需满足另外四个方程，以匹配叛逃飞船的位置与速度。

这个方程组可以用数值方法求解。所得结果即为小说中出现的那条曲线：其加速度的方向随时间改变，但大小保持不变：

\orthofig{images/E3__Fig003.png}{拦截叛逃飞船的飞行方案，已标出加速度方向}

关于变分法这一整门学科，文献 [1] 给出了远为严谨的处理，可供进一步阅读。

\section{参考文献}

\begin{itemize}
  \item [{[1]}] \emph{The Calculus of Variations}, M. Bendersky, Lecture Notes, City University of New York.
  \item [{[2]}] ``Baserunner's Optimal Path'', Davide Carozza, Stewart Johnson, Frank Morgan, \emph{The Mathematical Intelligencer}, 2009 年 10 月。
\end{itemize}

\begin{keybox}[本章要点]
\begin{itemize}
  \item \textbf{物理塑造文明}：正交宇宙没有普适光速，使得天象呈现"色谱拖尾"；发光释放（而非消耗）能量，使植物靠向夜空发光制造养分；发光引擎无需燃料与推进剂即可持续加速，使星际航行成为可能；面对"伤人者"（Hurtlers）的威胁，文明以"船上历许多代、故乡只过数年"的世代远航作为出路。
  \item \textbf{软性拦截}（soft interception）要求追击飞行器不仅与叛逃飞船相遇，还要\emph{与其速度匹配}。在所涉低速下沿用牛顿力学；正交宇宙的另类物理仅体现在引擎的运转上。
  \item 设追击飞行器初始位置 $\bp_0$、速度 $\bv_0$、最大加速度大小 $A$，叛逃飞船轨迹 $\mathbf{r}(t)=\mathbf{r}_0+\bw_0 t+\half\mathbf{a}t^{2}$。要在满足 $\bp(0)=\bp_0$、$\bp'(0)=\bv_0$、$\bp(T)=\mathbf{r}(T)$、$\bp'(T)=\mathbf{r}'(T)$、$|\bp''(t)|\le A$ 的前提下最小化拦截时刻 $T$。
  \item 假定全程以最大推力工作，把约束取为等式 $|\bp''(t)|=A$，即 $\bp''(t)\cdot\bp''(t)-A^{2}=0$。借助连续形式的拉格朗日乘子，构造泛函
        $I(\bp,\lambda)=\int_{0}^{T(\bp)}[\,1+\lambda(t)(\bp''(t)\cdot\bp''(t)-A^{2})\,]\dd t$，
        令其对一切合适变分的导数为零。
  \item 变分、两次分部积分并消去边界项后，得到 $\dfrac{\dd^{2}}{\dd t^{2}}\bigl(\lambda_M(t)\bp_M''(t)\bigr)=0$，故 $\lambda_M(t)\bp_M''(t)=\mathbf{C}t+\mathbf{D}$，从而最优轨道的加速度为
        $\bp_M''(t)=A(\mathbf{C}t+\mathbf{D})/|\mathbf{C}t+\mathbf{D}|$——方向可变、大小恒为 $A$。
  \item 边界项可由 $\mathbf{c}=\mathbf{C}T_M+\mathbf{D}$ 与 $\alpha=1/\bigl(2(\pm A|\mathbf{c}|-\mathbf{c}\cdot\mathbf{a})\bigr)$ 配平；追上叛逃飞船须 $A>|\mathbf{a}|$。
  \item 借助 $\operatorname{arcsinh}$ 可得显式特解 $\mathbf{V}(t)=\mathbf{V}_0+A(\operatorname{arcsinh}(t),\sqrt{1+t^{2}})$ 及其积分 $\mathbf{X}(t)$；再经 $t\to\beta t+\gamma$、$A\to A/\beta^{2}$ 与旋转 $\theta$ 推广，共有八个参数（$\mathbf{X}_0,\mathbf{V}_0$ 的四个分量及 $\beta,\gamma,\theta,T$），由 $t=0$ 与 $t=T$ 各四个方程数值求解。
\end{itemize}
\end{keybox}

\chapter{术语表}

本表汇总全书正文中出现的数学与物理术语，按英文原名词的字母顺序排列；每条给出中文译名、英文原词与简洁释义，便于中英对照、随时查阅。术语的中文译名与本书其余各章保持一致。

\begin{description}[style=nextline, leftmargin=0em, itemsep=0.8ex]

\item[基（Basis）]
\emph{n} 维矢量空间 $V$ 的一组基，是 $n$ 个矢量 $\{\mathbf{e}_1,\mathbf{e}_2,\dots,\mathbf{e}_n\}$ 的集合，使得 $V$ 中任意矢量 $\mathbf{v}$ 都可写成这些基矢量的倍数之和：
\[
\mathbf{v} = v^{1}\mathbf{e}_{1} + v^{2}\mathbf{e}_{2} + \cdots + v^{n}\mathbf{e}_{n} .
\]
数 $v^{1},v^{2},\dots,v^{n}$ 称为矢量 $\mathbf{v}$ 相对于该基的\emph{坐标}或\emph{分量}；同一个矢量相对于不同的基具有不同的分量。本书主要涉及四维矢量空间，基矢量常不以 $1,2,3,4$ 编号，而记作 $\mathbf{e}_x,\mathbf{e}_y,\mathbf{e}_z,\mathbf{e}_t$。我们所使用的基大多是\textbf{正交归一}（orthonormal）的："归一"指所有矢量长度为 $1$，"正交"指它们两两垂直。例如在实数四元组构成的矢量空间 $\mathbb{R}^4$ 中，配以标准点积，一组正交归一基就是\textbf{标准基}：
\[
\mathbf{e}_x=(1,0,0,0),\quad \mathbf{e}_y=(0,1,0,0),\quad \mathbf{e}_z=(0,0,1,0),\quad \mathbf{e}_t=(0,0,0,1).
\]

\orthofig{images/99__001.png}{二维空间中的基（Bases in two dimensions）}

\item[行列式（Determinant）]
$n\times n$ 矩阵 $A$ 的行列式记作 $\det(A)$，它给出与 $A$ 对应的线性函数对 $n$ 维立方体体积的改变因子。若 $A$ 把 $n$ 维立方体压扁到低于 $n$ 维，行列式为零；若 $A$ 反转 $n$ 维立方体的取向，行列式为负。例如在二维情形，点集 $\{(0,0),(1,0),(1,1),(0,1)\}$ 按\emph{逆时针}顺序给出边长 $1$ 的正方形顶点；把 $A$ 作用于这些点得到四个新点，$|\det(A)|$ 就是以这些新点为顶点的平行四边形的面积，而当新点按\emph{顺时针}顺序排列时 $\det(A)$ 为负。

\orthofig{images/99__002.png}{行列式示例（Determinant example）}

设 $A$ 为 $2\times2$ 矩阵
\[
A=\begin{pmatrix} a & b \\ c & d \end{pmatrix},
\]
则平行四边形的两条边是矢量 $\mathbf{v}=(a,c)$ 与 $\mathbf{w}=(b,d)$。由于矢量 $\mathbf{p}=(-c,a)$ 垂直于 $\mathbf{v}$，若以 $\mathbf{v}$ 为"底"，则平行四边形的高 $h=(\mathbf{w}\cdot\mathbf{p})/|\mathbf{p}|$；面积等于高乘以底长，于是
\[
\det(A)=h|\mathbf{v}|=h|\mathbf{p}|=\mathbf{w}\cdot\mathbf{p}=ad-bc .
\]
一般地，求 $n\times n$ 矩阵 $A$ 的行列式，先列出 $\{1,2,\dots,n\}$ 的所有排列；每个排列对应一个称为其符号的数，当该排列相对原序经历偶数次交换时为 $+1$，经历奇数次交换时为 $-1$。对每个排列，用其符号乘以 $n$ 个矩阵分量之积（每行取一个，列号由排列中对应的数确定），所有这些乘积之和就是行列式：
\[
\det(A)=\sum_{p\ \text{为}\ \{1,2,\dots,n\}\ \text{的排列}} \sgn(p)\, A^{1}_{p(1)} A^{2}_{p(2)} \cdots A^{n}_{p(n)} .
\]
本例中 $\{1,2\}$ 的排列只有 $\{1,2\}$（符号 $1$）与 $\{2,1\}$（符号 $-1$），故
\[
\det(A)=A^{1}_{1}A^{2}_{2}-A^{1}_{2}A^{2}_{1}=ad-bc .
\]
另一个很有用的特例是\textbf{对角矩阵} $D$——除对角元外分量全为零的矩阵——的行列式：唯一始终落在对角线上的排列是保持 $\{1,2,\dots,n\}$ 原序的那个，故
\[
\det(D)=D^{1}_{1}D^{2}_{2}\cdots D^{n}_{n} .
\]

\item[点积（Dot product）]
在四维黎曼几何中，两个矢量 $\mathbf{v}$ 与 $\mathbf{w}$ 的点积记作 $\mathbf{v}\cdot\mathbf{w}$，定义为
\[
\mathbf{v}\cdot\mathbf{w}=v^{x}w^{x}+v^{y}w^{y}+v^{z}w^{z}+v^{t}w^{t},
\]
其中 $v^{x}$ 等与 $w^{x}$ 等是两矢量在某正交归一基下的分量。它以显然的方式推广到 $n$ 维情形；用爱因斯坦求和约定即写作 $\mathbf{v}\cdot\mathbf{w}=v^{i}w^{i}$。点积还与两矢量的长度及其夹角 $\theta$ 有关：$\mathbf{v}\cdot\mathbf{w}=|\mathbf{v}|\,|\mathbf{w}|\cos\theta$。

这里从\textbf{正交归一}基出发给出的定义适用于物理应用——我们理所当然地认为可以测量矢量的长度，并能从物理上判断它是否垂直于另一矢量。而在纯数学中，则需\emph{先}给出点积才能定义正交归一基：正是 $\mathbf{e}_x\cdot\mathbf{e}_y=0$ 告诉我们 $\mathbf{e}_x$ 与 $\mathbf{e}_y$ 正交，而 $\mathbf{e}_x\cdot\mathbf{e}_x=1$ 告诉我们 $\mathbf{e}_x$ 的长度为 $1$。例如在处理实数四元组空间 $\mathbb{R}^4$ 时，一般会使用\textbf{标准点积}，其定义为
\[
(v^{x},v^{y},v^{z},v^{t})\cdot(w^{x},w^{y},w^{z},w^{t})=v^{x}w^{x}+v^{y}w^{y}+v^{z}w^{z}+v^{t}w^{t},
\]
由此 $\mathbb{R}^4$ 的标准基恰为正交归一基。（详见"点积"专题注记。）

\item[对偶（Dual）]
若 $V$ 是实矢量空间，则 $V$ 上的\textbf{对偶矢量}是从 $V$ 到实数集 $\mathbb{R}$ 的线性函数 $f$。$V$ 的\textbf{对偶空间}记作 $V^{*}$，是 $V$ 上全体对偶矢量构成的集合，其本身也是一个矢量空间。若 $\{\mathbf{e}_1,\mathbf{e}_2,\dots,\mathbf{e}_n\}$ 是 $V$ 的一组基，则 $V^{*}$ 的\textbf{对偶基} $\{\mathbf{e}^{1},\mathbf{e}^{2},\dots,\mathbf{e}^{n}\}$ 是 $V$ 上一组满足下式的函数：
\[
\mathbf{e}^{i}(\mathbf{e}_{j})=\delta^{i}_{j},
\]
其中 $\delta^{i}_{j}$ 称为克罗内克（Kronecker）$\delta$ 符号，当 $i=j$ 时为 $1$，当 $i\neq j$ 时为 $0$。（详见"对偶矢量"专题注记。）

\item[爱因斯坦求和约定（Einstein summation convention）]
爱因斯坦求和约定是书写含矢量与矩阵分量的表达式中各项之和的一种简便记法。按此约定，凡是同一指标变量在一个乘积中出现两次，就意味着该乘积实际上是对该指标的所有适用取值求和。例如
\[
M^{i}_{j}v^{j},
\]
其中 $v^{j}$ 表示 $n$ 维矢量 $\mathbf{v}$ 的一个分量，指标 $j$ 出现两次，求和约定告诉我们把它展开为
\[
M^{i}_{j}v^{j}=M^{i}_{1}v^{1}+M^{i}_{2}v^{2}+\cdots+M^{i}_{n}v^{n} .
\]
注意，写成 $v^{x}w^{x}$ 这类项时，重复的"$x$"只是标记这两个矢量各自的\emph{特定}分量，故求和约定并不适用；每当 $x,y,z,t$ 作为上标或下标出现时，都只表示对应于传统坐标轴的标记，而不是需要求和的指标变量。

\item[欧几里得群（Euclidean group）]
欧几里得群 $E(n)$ 是 $n$ 维欧几里得空间所有对称变换构成的群，包括平移、转动与反射。本书主要涉及 $E(4)$，即四维欧几里得空间的对称变换。只含平移与转动、不含反射的欧几里得群子群，称为\textbf{特殊欧几里得群} $SE(n)$。（详见"对称性"专题注记。）

\item[欧几里得宇宙（Euclidean universe）]
本书用"欧几里得宇宙"一词，指一种理想化的四维黎曼宇宙：它是\emph{无限}且\emph{平直}的，因而服从欧几里得几何的公设。在我们自己的宇宙中，远离任何强引力场的\emph{空间}的微小区域，其几何非常接近欧几里得几何，但这类区域的\emph{时空}几何却并非如此，因为时间间隔不服从勾股定理；在我们宇宙中，近似的平直时空由闵可夫斯基（Minkowski）时空描述。

\textbf{注意}：在物理学文献中，"欧几里得"一词常被用于经威克转动（Wick rotation）得到的物理定律版本。这些定律\emph{并非}\emph{正交}宇宙所适用的定律。

\item[群（Group）]
在数学中，群是集合 $G$ 连同一种"运算"——把 $G$ 中两个元素结合为第三个元素的方式。这种群运算有时称为"乘法"，把元素 $g$ 与 $h$ 结合所得的元素通常记作 $gh$。它必须服从与正数普通乘法类似的规则，但不必是\textbf{可交换}的，即可能有 $gh\neq hg$。具体地：
\begin{itemize}
  \item 群"乘法"必须\textbf{满足结合律}：$(gh)k=g(hk)$。
  \item 群必须含有\textbf{单位元}，通常记作 $e$（矩阵群中有时记作 $1$ 或 $I$），满足 $ge=eg=g$。
  \item 群中每个元素 $g$ 都必须有逆元 $g^{-1}$，满足 $gg^{-1}=g^{-1}g=e$。
\end{itemize}
可交换群的例子：正实数集 $\mathbb{R}^{+}$（以普通乘法为群运算）；实数 $n$ 元组构成的矢量空间 $\mathbb{R}^{n}$（以矢量加法为群运算、零矢量为单位元）。不可交换群的例子：全体行列式非零的 $n\times n$ 实矩阵集合 $GL(n)$（以矩阵乘法为群运算）；全体转置等于逆的 $n\times n$ 实矩阵集合 $O(n)$（以矩阵乘法为群运算）。

\item[单位矩阵（Identity matrix）]
$n\times n$ 单位矩阵记作 $I_n$，是除对角元为 $1$ 外其余分量全为零的矩阵。若 $n$ 由上下文显然可知，就简写为 $I$ 而不写 $I_n$。例如 $I_4$ 为
\[
I_{4}=\begin{pmatrix}1&0&0&0\\0&1&0&0\\0&0&1&0\\0&0&0&1\end{pmatrix}.
\]
用单位矩阵作乘法不会改变任何其他矩阵：若 $A$ 是任意 $n\times n$ 矩阵，则
\[
AI_{n}=I_{n}A=A .
\]

\item[逆矩阵（Inverse matrix）]
$n\times n$ 矩阵 $A$ 的逆是满足下式的 $n\times n$ 矩阵，记作 $A^{-1}$：
\[
AA^{-1}=A^{-1}A=I_{n},
\]
其中 $I_n$ 是 $n\times n$ 单位矩阵。$A$ 的逆存在，当且仅当 $A$ 的行列式不为零。矩阵乘积的逆等于各个矩阵之逆的\emph{逆序}乘积：
\[
(AB)^{-1}=B^{-1}A^{-1}.
\]

\item[线性函数（Linear function）]
线性函数是从一个矢量空间到另一个矢量空间（或许就是同一个）的函数，其特点是：先作函数还是先对矢量作加法或数乘，结果没有区别。也就是说，函数 $f$ 是线性的，当且仅当
\[
f(\mathbf{v}+\mathbf{w})=f(\mathbf{v})+f(\mathbf{w}),\qquad f(s\mathbf{v})=s\,f(\mathbf{v}).
\]
若 $f$ 是从 $m$ 维矢量空间 $V$ 到 $n$ 维矢量空间 $W$ 的线性函数，并为 $V$ 取基 $\{\mathbf{e}_1,\dots,\mathbf{e}_m\}$、为 $W$ 取基 $\{\mathbf{e}'_1,\dots,\mathbf{e}'_n\}$，就可用一个 $n\times m$ 矩阵 $M$ 来描述 $f$：矩阵分量 $M^{i}_{j}$ 是矢量 $f(\mathbf{e}_j)$ 相对于 $W$ 的基的第 $i$ 个分量。于是（用爱因斯坦求和约定）
\[
f(\mathbf{v})=f(v^{j}\mathbf{e}_{j})=v^{j}f(\mathbf{e}_{j})=M^{i}_{j}v^{j}\mathbf{e}'_{i}.
\]

\item[洛伦兹型（Lorentzian）]
本书用形容词"洛伦兹型"来区分\textbf{适用于我们自身宇宙的物理}与适用于\emph{正交}宇宙的物理。（详见"洛伦兹几何与黎曼几何"专题注记。）

\emph{[为何称"洛伦兹型"？亨德里克·安东·洛伦兹（Hendrik Antoon Lorentz）是一位荷兰物理学家，他建立了爱因斯坦在狭义相对论中所使用的数学框架。"洛伦兹型"一词用来指这样一种几何：其中一个维度作为\emph{时间}维度而与其余维度相区别。相比之下，在黎曼空间中，所有维度在根本上是彼此平等的。]}

\item[矩阵（Matrix）]
$n\times m$ 矩阵 $M$ 是数的网格或表格（本书中为实数或复数），有 $n$ 行、每行 $m$ 个数。各个数称为矩阵的\emph{分量}；$M^{i}_{j}$ 表示矩阵 $M$ 第 $i$ 行第 $j$ 个数。若 $A$ 是 $n\times m$ 矩阵、$B$ 是 $m\times p$ 矩阵，则可把 $A$ 与 $B$ 相乘得到 $n\times p$ 矩阵 $C=AB$，其中
\[
C^{i}_{j}=A^{i}_{1}B^{1}_{j}+A^{i}_{2}B^{2}_{j}+\cdots+A^{i}_{m}B^{m}_{j},
\]
或直接用爱因斯坦求和约定写作 $C^{i}_{j}=A^{i}_{k}B^{k}_{j}$。注意，要使矩阵乘积 $AB$ 存在，$A$ 的\emph{行}数必须等于 $B$ 的\emph{列}数。若 $A$ 与 $B$ 是同阶方阵，则 $AB$ 与 $BA$ 都存在，但一般并不相等。

可以把 $n\times m$ 矩阵 $M$ 视为从 $m$ 维矢量空间 $V$ 到 $n$ 维矢量空间 $W$ 的线性函数：把 $V$ 中任意矢量 $\mathbf{v}$ 看作一个 $m\times 1$ 矩阵，其单列的数就是 $\mathbf{v}$ 在某选定基下的 $m$ 个分量；则矩阵乘积 $M\mathbf{v}$ 给出一个 $n\times 1$ 矩阵，可视为 $W$ 中某矢量在某选定基下的 $n$ 个分量。于是 $M$ 连同 $V$、$W$ 中各选定的基，给出从 $V$ 到 $W$ 的线性函数 $f$。若 $V$ 的基为 $\{\mathbf{e}_1,\dots,\mathbf{e}_m\}$、$W$ 的基为 $\{\mathbf{e}'_1,\dots,\mathbf{e}'_n\}$，用爱因斯坦求和约定可写
\[
f(\mathbf{v})=f(v^{j}\mathbf{e}_{j})=v^{j}f(\mathbf{e}_{j})=M^{i}_{j}v^{j}\mathbf{e}'_{i}.
\]

\item[零矢量（Null Vector）]
沿普通时空中光所描绘的世界线指向的矢量。（详见"洛伦兹几何与黎曼几何"专题注记。）

\item[O(4)]
全体转置等于逆的 $4\times4$ 实矩阵集合称为 $O(4)$。它以矩阵乘法为群运算，构成一个群，另一名称为四维\emph{正交群}。若把这些矩阵视为四维空间上的线性函数，它们描述该空间中所有的转动与反射。（详见"对称性"专题注记。）

\item[黎曼型（Riemannian）]
本书用形容词"黎曼型"来区分\textbf{适用于\emph{正交}宇宙的物理}与适用于我们自身宇宙的物理。（详见"洛伦兹几何与黎曼几何"专题注记。）

\emph{[为何称"黎曼型"？格奥尔格·弗里德里希·伯恩哈德·黎曼（Georg Friedrich Bernhard Riemann）是微分几何领域的先驱之一，这门学科把欧几里得几何推广到弯曲曲面与空间。在该领域中，数学家用"黎曼型"来描述我们通常视为各类\emph{空间}（无论平直或弯曲）、其中所有维度被视为根本平等的几何。相比之下，在我们自身的洛伦兹型\emph{时空}中，其中一个维度——\emph{时间}——受到特殊对待。]}

\item[黎曼标量波（RSW）方程（Riemannian Scalar Wave (RSW) Equation）]
这是标量波 $A$ 在黎曼宇宙真空中传播所满足的方程：
\begin{equation}
\partial_{x}^{2}A+\partial_{y}^{2}A+\partial_{z}^{2}A+\partial_{t}^{2}A+\wms^{2}A=0 . \tag{RSW}
\end{equation}
（详见"标量波"专题注记。）

\item[黎曼矢量波（RVW）方程（Riemannian Vector Wave (RVW) Equations）]
这些是矢量波 $\mathbf{A}$ 在黎曼宇宙真空中传播所满足的方程：
\begin{equation}
\partial_{x}^{2}\mathbf{A}+\partial_{y}^{2}\mathbf{A}+\partial_{z}^{2}\mathbf{A}+\partial_{t}^{2}\mathbf{A}+\wms^{2}\mathbf{A}=\mathbf{0}, \tag{RVW}
\end{equation}
\begin{equation}
\partial_{x}A^{x}+\partial_{y}A^{y}+\partial_{z}A^{z}+\partial_{t}A^{t}=0 . \tag{Transverse}
\end{equation}
（详见"矢量波"专题注记。）可以把第一式推广，以包含源项：
\begin{equation}
\partial_{x}^{2}\mathbf{A}+\partial_{y}^{2}\mathbf{A}+\partial_{z}^{2}\mathbf{A}+\partial_{t}^{2}\mathbf{A}+\wms^{2}\mathbf{A}+\mathbf{j}=\mathbf{0} . \tag{RVWS}
\end{equation}
（详见"黎曼电磁学"专题注记。）

\item[标量（Scalar）]
所有观察者都会一致同意的数，因为它不依赖于观察者对坐标系的选择。例如气压是标量，而"电场在 $x$ 方向的分量"则不是。

\item[SO(4)]
全体转置等于逆且行列式为 $1$ 的 $4\times4$ 实矩阵集合称为 $SO(4)$。它以矩阵乘法为群运算，构成一个群，另一名称为四维\emph{特殊正交群}。若把这些矩阵视为四维空间上的线性函数，它们描述该空间中所有的转动。（详见"对称性"专题注记。）

\item[子空间（Subspace）]
矢量空间 $V$ 的子空间是 $V$ 的一个子集，它本身也是一个矢量空间。例如在二维及以上的矢量空间中，任何过原点的直线都是原矢量空间的一维子空间。但并非矢量空间的所有子集都是\emph{子空间}：所有长度为 $1$ 的矢量构成的球面就\emph{不}构成子空间，因为把两个长度为 $1$ 的矢量相加会得到另一个长度不同的矢量，落在所选子集之外。

\item[转置（Transpose）]
$n\times m$ 矩阵 $A$ 的转置是 $m\times n$ 矩阵，记作 $A^{\mathrm{T}}$，其各列就是原矩阵的各行。矩阵乘积的转置等于各个矩阵转置之积的\emph{逆序}乘积：
\[
(AB)^{\mathrm{T}}=B^{\mathrm{T}}A^{\mathrm{T}}.
\]

\item[矢量空间（Vector space）]
矢量空间 $V$ 是一组对象构成的集合，其中对象可以彼此相加，也可以乘以某种数（本书中这些数为实数或复数）。它须满足以下各条：
\begin{itemize}
  \item 每个矢量空间都含有一个\emph{零矢量} $\mathbf{0}$，把它加到任何矢量上都不改变该矢量：$\mathbf{v}+\mathbf{0}=\mathbf{v}$。
  \item 每个矢量 $\mathbf{v}$ 都有一个相反量（或称\emph{加法逆元}）$-\mathbf{v}$，满足 $\mathbf{v}+(-\mathbf{v})=\mathbf{0}$。
  \item 乘法对加法\emph{满足分配律}：$s(\mathbf{v}+\mathbf{w})=s\mathbf{v}+s\mathbf{w}$ 且 $(r+s)\mathbf{v}=r\mathbf{v}+s\mathbf{v}$。
  \item 还可定义矢量之间的减法：$\mathbf{v}-\mathbf{w}=\mathbf{v}+(-\mathbf{w})$。
  \item 矢量空间对这些运算必须\emph{封闭}：若选定了集合 $V$ 并在其上定义这些运算，则运算结果决不允许落到 $V$ \emph{之外}。
\end{itemize}
例如取 $V$ 为全体实数四元组构成的集合 $\mathbb{R}^4$，按显然的方式定义加法、数乘、零矢量与相反量：
\[
(x,y,z,t)+(a,b,c,d)=(x+a,y+b,z+c,t+d),\qquad s(x,y,z,t)=(sx,sy,sz,st),
\]
\[
\mathbf{0}=(0,0,0,0),\qquad -(x,y,z,t)=(-x,-y,-z,-t).
\]

\item[威克转动（Wick rotation）]
威克转动是现代物理学中用来分析\emph{我们自身宇宙}中某些问题的一种手段。在威克转动中，把时间 $t$ 替换为虚数 $\ii\tau$，而\emph{其余一切都不改变}。与之相反，小说\emph{正交}所涉及的物理，是改变各方程中某些参数的正负号，由此得到的结果与威克转动大不相同。

例如，在自然单位下，普通的洛伦兹型克莱因—戈尔登（Klein-Gordon）方程为
\begin{equation}
\partial_{x}^{2}\psi+\partial_{y}^{2}\psi+\partial_{z}^{2}\psi-\partial_{t}^{2}\psi-m^{2}\psi=0 . \tag{1}
\end{equation}
（常可看到它写作 $\partial^{\mu}\partial_{\mu}\psi+m^{2}\psi=0$，其中对 $\mu$ 用爱因斯坦求和约定求和；$m^{2}\psi$ 项取正号，是因为作者采用了闵可夫斯基度规的 $(+,-,-,-)$ 号差，即 $\partial^{\mu}\partial_{\mu}\psi=\partial_{t}^{2}\psi-\partial_{x}^{2}\psi-\partial_{y}^{2}\psi-\partial_{z}^{2}\psi$。）

经威克转动（$t\to\ii\tau$）后的版本为
\begin{equation}
\partial_{x}^{2}\psi+\partial_{y}^{2}\psi+\partial_{z}^{2}\psi+\partial_{\tau}^{2}\psi-m^{2}\psi=0 . \tag{2}
\end{equation}
适用于\emph{正交}宇宙的黎曼型版本与之相似，但 $m^{2}$ 项的符号相反：
\begin{equation}
\partial_{x}^{2}\psi+\partial_{y}^{2}\psi+\partial_{z}^{2}\psi+\partial_{t}^{2}\psi+m^{2}\psi=0 . \tag{3}
\end{equation}
普通的克莱因—戈尔登方程 \textbf{(1)} 具有有界的平面波解：
\begin{equation}
\psi(\mathbf{x})=\sin(k^{x}x+k^{y}y+k^{z}z-k^{t}t), \tag{4}
\end{equation}
其中 $\mathbf{k}$ 是任一满足 $(k^{x})^{2}+(k^{y})^{2}+(k^{z})^{2}-(k^{t})^{2}=-m^{2}$ 的类时矢量。适用于\emph{正交}宇宙的 \textbf{(3)} 式同样具有有界的平面波解：
\begin{equation}
\psi(\mathbf{x})=\sin(\mathbf{k}\cdot\mathbf{x}), \tag{5}
\end{equation}
其中 $|\mathbf{k}|^{2}=m^{2}$。但威克转动版本 \textbf{(2)} 的\emph{全部}解都会沿至少一个方向指数增长：若把解 \textbf{(5)} 代入方程 \textbf{(2)}，将得到 $|\mathbf{k}|^{2}=-m^{2}$，这只有在 $\mathbf{k}$ 至少有一个分量为虚数时才可能成立——而虚数的正弦会指数增长。

\end{description}


\appendix
\chapter{数学预备}

本附录汇列全书反复用到的数学工具与记号约定，供读者随时查阅。这些内容都是标准的线性代数、
群表示论、傅里叶分析与变分法知识，本身不涉及正交宇宙的任何物理假设；把它们集中放在这里，
是为了让正文可以专注于物理本身，而不必在每一章重复铺垫。

\section{度规、指标与对偶矢量}

全书采用自然单位制，取 $\hbar=1$、$c=1$，时空坐标记作 $(x,y,z,t)$。四维矢量的分量以
上标书写：$v^x,v^y,v^z,v^t$，其中 $x,y,z,t$ 是标号而非幂次。两个四维矢量之间的内积由度规
（metric）$g_{ab}$ 给出，
\begin{equation}
A\cdot B = g_{ab}\,A^a B^b ,
\end{equation}
这里对重复指标求和（Einstein 求和约定），指标取遍 $x,y,z,t$。

本书涉及两种度规。其一是我们自身宇宙的\emph{洛伦兹度规}（Lorentzian metric），本书统一取
\begin{equation}
g^{\Lor}_{ab}=\diag(-1,+1,+1,+1),
\tag{A.1}
\end{equation}
于是
\begin{equation}
A\cdot B \big|_{\Lor} = -A^t B^t + A^x B^x + A^y B^y + A^z B^z .
\end{equation}
其二是正交宇宙的\emph{黎曼度规}（Riemannian metric），取全正号
\begin{equation}
g^{\Riem}_{ab}=\diag(+1,+1,+1,+1),
\tag{A.2}
\end{equation}
于是
\begin{equation}
A\cdot B \big|_{\Riem} = A^x B^x + A^y B^y + A^z B^z + A^t B^t .
\end{equation}
两种度规的唯一差别，正是时间项前面那个符号；而正是这一符号之差，造成了本书所讲述的全部
物理差异。

由度规可以升降指标。若 $A^a$ 是逆变（contravariant）分量，则协变（covariant）分量为
\begin{equation}
A_a = g_{ab} A^b ,\qquad A^a = g^{ab} A_b ,
\end{equation}
其中 $g^{ab}$ 是 $g_{ab}$ 的逆矩阵。在黎曼度规下，$g_{ab}$ 与 $g^{ab}$ 都是单位矩阵，
升降指标不改变分量；而在洛伦兹度规下，时间分量的符号会发生翻转。

对四维矢量 $v$，其“长度平方”为 $v\cdot v$。在洛伦兹度规下，$v\cdot v$ 可正可负可为零，
据此把矢量分为类空（spacelike，$v\cdot v>0$）、类时（timelike，$v\cdot v<0$）与类光
（lightlike / null，$v\cdot v=0$）三类，后者的集合构成光锥（light cone）。在黎曼度规下，
度规正定，$v\cdot v\ge 0$ 且当且仅当 $v=0$ 时取零：\textbf{不存在非零的类光矢量，也不存在
光锥}。这正是正交宇宙“没有普适光速”的几何根源，我们在正文第 1 章与《几何与运动》一章中
已反复用到。

\section{对偶空间与对偶矢量}

给定矢量空间 $V$（这里即四维时空），其\emph{对偶空间}（dual space）$V^*$ 由所有从 $V$ 到
实数的线性映射（线性泛函）组成。每个矢量 $v\in V$ 都自然地对应一个对偶矢量
$\tilde v\in V^*$，其作用规则为
\begin{equation}
\tilde v(w) = v\cdot w .
\end{equation}
在黎曼度规下，$v\mapsto \tilde v$ 是一个同构（isomorphism）：不同的矢量给出不同的对偶矢量，
且每个对偶矢量都恰好来自一个矢量。于是可以无歧义地“把矢量与对偶矢量视为同一事物的两种
记法”。

在洛伦兹度规下，这个映射虽仍是同构（$g_{ab}$ 非退化），但矢量的分类使得几何图景显著复杂化：
类光矢量对应的对偶矢量为零模，因而“同时”“垂直”等概念都变得与观察者相关。正文《双勾股
定理》一章正是借助对偶空间的这套语言，来说明正交宇宙中许多几何关系的优雅来源。

\section{正交群、四元数与旋量}

黎曼度规 $(+,+,+,+)$ 在坐标变换下的保持群是\emph{正交群}（orthogonal group）$O(4)$；
其中行列式为 $+1$ 的子群 $\mathrm{SO}(4)$ 称为特殊正交群（special orthogonal group），
它给出四维空间中的“转动”——在正交宇宙里，四个维度地位平等，转动可以同时混合空间与时间
方向。此外，交换两个坐标轴的反射（reflection）也属于 $O(4)$。

$\mathrm{SO}(4)$ 有一个著名的二对一同态（two-to-one homomorphism）
\begin{equation}
\mathrm{SU}(2)\times \mathrm{SU}(2)\ \longrightarrow\ \mathrm{SO}(4),
\tag{A.3}
\end{equation}
即 $\mathrm{SU}(2)\times \mathrm{SU}(2)$ 是 $\mathrm{SO}(4)$ 的\emph{双覆盖}
（double cover）。这一事实可以用四元数（quaternion）非常简洁地表述。

\subsection{四元数}

四元数是复数的推广：除实数单位 $1$ 外，还有三个平方为 $-1$ 的“虚单位” $i,j,k$，它们满足
\begin{equation}
i^{2}=j^{2}=k^{2}=-1,\qquad ij=k,\quad jk=i,\quad ki=j,\qquad ijk=-1 .
\end{equation}
一般四元数写作 $a+bi+cj+dk$（$a,b,c,d\in\mathbb{R}$），其加法与乘法除乘法不满足交换律
（non-commutative）外，与我们熟悉的实数、复数运算一致。

四元数可表示为 $2\times 2$ 复矩阵，且矩阵乘法恰对应四元数乘法。把四维实空间记为 $H$，
其中四个基矩阵 $H_x,H_y,H_z,H_t$ 满足与 $i,j,k,1$ 完全相同的乘法表；于是任一时空矢量
$(x,y,z,t)$ 都可以写成一个四元数 $xH_x+yH_y+zH_z+tH_t$。$\mathrm{SO}(4)$ 的转动作用，
便对应到在左侧（或右侧）乘上单位四元数。这一套表述是附录 B 中推导黎曼 Dirac 方程的
基础，详见正文《黎曼量子力学》各章。

\subsection{旋量}

$\mathrm{SU}(2)\times \mathrm{SU}(2)\to \mathrm{SO}(4)$ 是二对一的：两个不同的
$\mathrm{SU}(2)\times\mathrm{SU}(2)$ 元素对应同一个 $\mathrm{SO}(4)$ 转动。因此
$\mathrm{SO}(4)$ 的表示并不能完全描述量子态——量子态所在的表示是这一双覆盖的表示，
其中的元素称为\emph{旋量}（spinor）。旋量在四维转动一周（$2\pi$）后并不回到自身，而要
在转动两周（$4\pi$）后才复原。这一“自旋”（spin）结构在正交宇宙中依然存在，并直接进入
电子等粒子的量子描述。

\section{傅里叶方法与色散关系}

本书大量使用平面波（plane wave）解。对场 $A(x)$，设其具有形式
\begin{equation}
A(x)=A_0\, \ee^{\,\ii\, k\cdot x},\qquad k=(k_x,k_y,k_z,\omega),
\end{equation}
其中 $k$ 为四维波矢，$\omega$ 为角频率（angular frequency），$k_x,k_y,k_z$ 为空间波数
（wave number）分量。在黎曼度规下 $k\cdot k=k_x^{2}+k_y^{2}+k_z^{2}+\omega^{2}$；
在洛伦兹度规下则为 $k\cdot k=-\omega^{2}+k_x^{2}+k_y^{2}+k_z^{2}$。

把平面波代入波动方程，即得\emph{色散关系}（dispersion relation），它把频率与波数联系
起来。以质量为 $m$ 的自由场为例：

在洛伦兹宇宙中，Klein–Gordon 方程给出
\begin{equation}
\omega^{2}-k_x^{2}-k_y^{2}-k_z^{2}=m^{2},
\tag{A.4}
\end{equation}
相应地，相位速度 $v_{\mathrm p}=\omega/|\mathbf{k}|$ 随波长而变，但在 $|\mathbf{k}|\to 0$
极限下趋于统一的常数 $c=1$——这正对应“光速是普适常量”。

在正交宇宙中，黎曼标量波方程（Riemannian Scalar Wave equation, RSW）为
\begin{equation}
\partial_x^{2}A+\partial_y^{2}A+\partial_z^{2}A+\partial_t^{2}A+\wms^{2}A=0 ,
\tag{A.5}
\end{equation}
其中 $\wms$ 是质量频率（mass frequency），$\wms=mc^{2}/\hbar$。代入平面波得
\begin{equation}
\omega^{2}+k_x^{2}+k_y^{2}+k_z^{2}=\wms^{2},
\tag{A.6}
\end{equation}
即 $\omega^{2}+|\mathbf{k}|^{2}=\wms^{2}$。与 (A.4) 相较，空间项与频率项的符号关系发生了
整体翻转：正交宇宙中的相位速度
\begin{equation}
v_{\mathrm p}=\frac{\omega}{|\mathbf{k}|}=\frac{\sqrt{\wms^{2}-|\mathbf{k}|^{2}}}{|\mathbf{k}|}
\end{equation}
\textbf{依赖于波长}，不存在一个所有波长共有的极限速度。长波（$|\mathbf{k}|\to 0$）
趋于无穷快，短波（$|\mathbf{k}|\to \wms$）趋于静止。这正是正文《几何与波》《黎曼电磁学》
诸章中“没有普适光速”“发光释能”等现象的数学根源。

\section{特殊函数与常用积分}

黎曼黑体谱（正文《黎曼热力学》一章）的推导涉及若干特殊函数与定积分，这里一并列出。

\subsection{Gamma 函数与 Riemann zeta 函数}

Gamma 函数 $\Gamma(z)$ 由积分定义
\begin{equation}
\Gamma(z)=\int_{0}^{\infty} t^{\,z-1}\ee^{-t}\,\dd t ,\qquad \Re z>0 ,
\end{equation}
满足 $\Gamma(n)=(n-1)!$（$n$ 为正整数）与 $\Gamma(1/2)=\sqrt{\pi}$ 以及递推关系
$\Gamma(z+1)=z\,\Gamma(z)$。Riemann zeta 函数由级数定义
\begin{equation}
\zeta(s)=\sum_{n=1}^{\infty}\frac{1}{n^{s}},\qquad \Re s>1 ,
\end{equation}
并可由解析延拓推广到更大的区域。

\subsection{多对数与玻色型积分}

\emph{多对数}（polylogarithm）定义为
\begin{equation}
\mathrm{Li}_s(z)=\sum_{n=1}^{\infty}\frac{z^{n}}{n^{s}},
\end{equation}
它是 $\zeta(s)$ 的推广：$\mathrm{Li}_s(1)=\zeta(s)$。处理黑体谱时反复出现形如
\begin{equation}
\int_{0}^{\infty}\frac{x^{\,n}}{\ee^{x}-1}\,\dd x=\Gamma(n+1)\,\zeta(n+1)=n!\,\zeta(n+1)
\end{equation}
的玻色型积分；例如 $n=3$ 时结果为 $\pi^{4}/15$。在正交宇宙中，色散关系 (A.6) 把积分
上限截断在 $\wms$ 处，从而给出与洛伦兹情形不同的黑体谱，具体形式见正文《黎曼热力学》。

\subsection{贝塞尔函数}

含 $\sqrt{\cdot}$ 型色散关系（如 (A.6)）的方程，在分离变量后常化为贝塞尔方程
\begin{equation}
x^{2}y''+xy'+(x^{2}-n^{2})y=0 ,
\end{equation}
其解为第一类与第二类贝塞尔函数 $J_n(x)$、$Y_n(x)$；在虚宗量（imaginary argument）下则
出现修正贝塞尔函数 $I_n(x)$、$K_n(x)$。正文讨论弯曲黎曼空间中的波动与黑体谱时会用到它们。

\section{变分法基础}

正文《世界观与文明设定》一章中求解“软性拦截”问题，以及《几何与运动》一章导出自由粒子的
拉格朗日量与哈密顿量，都用到了\emph{变分法}（calculus of variations）。这里给出最基本
的形式。

给定泛函
\begin{equation}
S[\mathbf{x}]=\int_{t_0}^{t_1} L\big(\mathbf{x}(t),\dot{\mathbf{x}}(t),t\big)\,\dd t ,
\end{equation}
其中 $L$ 为拉格朗日量（Lagrangian），$\mathbf{x}(t)$ 是一条候选轨迹。$S$ 取极值
（极小或极大）的必要条件是 $\mathbf{x}(t)$ 满足 Euler–Lagrange 方程
\begin{equation}
\frac{\dd}{\dd t}\left(\frac{\partial L}{\partial \dot{\mathbf{x}}}\right)
-\frac{\partial L}{\partial \mathbf{x}}=0 .
\tag{A.7}
\end{equation}
若被积函数不含 $t$ 显含项，则另有守恒量（Beltrami 恒等式）
\begin{equation}
\dot{\mathbf{x}}\cdot\frac{\partial L}{\partial\dot{\mathbf{x}}}-L=\text{const}.
\end{equation}

对“最大化某量”“最小化某量”的约束极值问题，可以引入 Lagrange 乘子（Lagrange multiplier），
把所有约束并入一个增广泛函，再用 (A.7) 求解。正文《世界观与文明设定》一章中的软性拦截
问题，正是把“与目标飞船速度匹配”作为约束、以“追击时间最短”为目标函数的一次典型应用。

\chapter{关键推导补全}

正文各章按原作者的叙述节奏展开，若干结论给出了结果而略去了中间步骤。本附录把这些跳步
集中补齐，便于读者（尤其是偏好自己复算的读者）验证。每节开头都注明其对应正文的哪一章、
以及补充的是哪一处推导；已在前文完整给出的步骤此处不再重复。

\section{质量频率 \texorpdfstring{$\wms$}{omega\_m} 的引入}
\emph{补充第 1 章与《几何与运动》中“质量频率”定义的来源。}

在洛伦兹宇宙中，自由粒子的能量--动量关系为（取 $c=1$）
\begin{equation}
E^{2}-p^{2}=m^{2},
\end{equation}
即四维动量 $P=(p_x,p_y,p_z,E)$ 在洛伦兹度规下的自内积等于 $m^{2}$。在正交宇宙中，度规
变为全正号，同一物理内容的关系式相应地变为
\begin{equation}
E^{2}+p^{2}=m^{2},
\tag{B.1}
\end{equation}
即 $P\cdot P=m^{2}$（黎曼内积）。空间中静止的粒子对应 $p=0$，此时 $E=m$。

在量子力学中，能量与角频率由 de Broglie 关系 $E=\hbar\omega$ 联系。把静止粒子的能量
$E=m$ 代入，得 $\hbar\omega=m$，即
\begin{equation}
\wms=\frac{m}{\hbar}.
\end{equation}
还原量纲（把 $c$ 放回）即
\begin{equation}
\wms=\frac{mc^{2}}{\hbar},
\end{equation}
这就是\emph{质量频率}（mass frequency）：它是以频率为量纲的粒子静止质量。把 (B.1) 除以
$\hbar^{2}$ 并写作 $\omega^{2}+|\mathbf{k}|^{2}=\wms^{2}$（其中 $p=\hbar|\mathbf{k}|$），
正是正文各章反复使用的色散关系。

\section{黎曼 Proca 方程与电荷守恒}
\emph{补充《黎曼电磁学》中由 (RVWS) 与横波条件导出 Proca 方程、以及导出 (1) 的步骤。}

从带源项的黎曼矢量波方程出发：
\begin{align}
\partial_{x}^{2}\bA+\partial_{y}^{2}\bA+\partial_{z}^{2}\bA+\partial_{t}^{2}\bA+\wms^{2}\bA+\bj&=\mathbf{0}, \tag{B.2}\\
\partial_{x}A^{x}+\partial_{y}A^{y}+\partial_{z}A^{z}+\partial_{t}A^{t}&=0. \tag{B.3}
\end{align}
在黎曼空间的直角坐标中，度规为单位矩阵，故 $\partial_{b}\partial^{b}=\partial_{x}^{2}+\partial_{y}^{2}
+\partial_{z}^{2}+\partial_{t}^{2}$。用 $A^{a}$ 的某个分量记 (B.2) 左端，可写为
\begin{equation}
\partial_{b}\partial^{b}A^{a}=-\wms^{2}A^{a}-j^{a}.
\tag{B.4}
\end{equation}

下一步把电磁场张量 $F_{ab}=\partial_{a}A_{b}-\partial_{b}A_{a}$ 的散度与 (B.4) 联系
起来。直接计算：
\begin{align}
\partial_{b}F^{ab}
&=\partial_{b}\bigl(\partial^{a}A^{b}-\partial^{b}A^{a}\bigr)\nonumber\\
&=\partial^{a}\bigl(\partial_{b}A^{b}\bigr)-\partial_{b}\partial^{b}A^{a}\nonumber\\
&=-\partial_{b}\partial^{b}A^{a},
\tag{B.5}
\end{align}
其中最后一步用到了横波条件 (B.3)：$\partial_{b}A^{b}=0$。把 (B.4) 代入 (B.5)，
\begin{equation}
\partial_{b}F^{ab}=\wms^{2}A^{a}+j^{a},
\end{equation}
即黎曼 Proca 方程
\begin{equation}
\partial_{b}F^{ab}-\wms^{2}A^{a}-j^{a}=0 . \tag{B.6}
\end{equation}
与洛伦兹麦克斯韦方程 $\partial_{b}F^{ab}-j^{a}=0$ 相比，只多出质量项 $-\wms^{2}A^{a}$；
正是它把库仑势由指数衰减改为空间振荡（见正文《黎曼电磁学》）。

再对 (B.6) 取散度 $\partial_{a}$。由于 $F^{ab}$ 反对称，$\partial_{a}\partial_{b}F^{ab}=0$
恒成立（对 $a,b$ 交换后求和相消），于是
\begin{equation}
0=\wms^{2}\,\partial_{a}A^{a}+\partial_{a}j^{a}=\partial_{a}j^{a},
\end{equation}
其中又用了 (B.3)。这就得到正文的方程 (1)：
\begin{equation}
\partial_{x}j^{x}+\partial_{y}j^{y}+\partial_{z}j^{z}+\partial_{t}j^{t}=0,
\end{equation}
即\textbf{电荷守恒}（conservation of charge）：电荷密度随时间的变化率等于电流密度的负
散度。

\section{黎曼 Dirac 方程的四元数形式}
\emph{补充《黎曼量子力学》中将平面波 Dirac 方程约化为四元数方程的一步。}

在 Weyl 基下，黎曼 Dirac 方程的平面波解写作一对左手、右手的二分旋量 $(\xi_{L},\xi_{R})$，
它们与能量--动量矢量 $k=(k_x,k_y,k_z,\omega)$ 通过
\begin{equation}
k^{\mu}H_{\mu}\,\xi_{R}-m\,\xi_{L}=0
\tag{B.7}
\end{equation}
相联系，其中 $H_{\mu}$（$\mu=x,y,z,t$）是张开 $2\times2$ 复矩阵空间 $\mathbb{H}$ 的四个
基矩阵，其乘法表与四元数基 $i,j,k,1$ 一致。

若把 $\mathbb{H}$ 与 $\mathbb{C}^{2}$ 都等同于四元数，则 (B.7) 中出现的三个对象分别按
$\mathrm{SU}(2)\times\mathrm{SU}(2)$ 的不同表示变换。记
\begin{equation}
Q_{V}=k^{\mu}H_{\mu},\qquad q_{L}\leftrightarrow \xi_{L},\qquad q_{R}\leftrightarrow\xi_{R},
\end{equation}
则在元素 $(g,h)\in \mathrm{SU}(2)\times\mathrm{SU}(2)$ 下
\begin{align}
\rho_{(\frac12,\frac12)}(g,h)\,Q_{V}&=g\,Q_{V}\,h^{-1},\\
\rho_{(\frac12,0)}(g,h)\,q_{L}&=g\,q_{L},\\
\rho_{(0,\frac12)}(g,h)\,q_{R}&=h\,q_{R}.
\end{align}
于是乘积 $Q_{V}q_{R}$ 的变换为
\begin{equation}
Q_{V}q_{R}\ \longmapsto\ (g\,Q_{V}\,h^{-1})(h\,q_{R})=g\,(Q_{V}\,q_{R}),
\end{equation}
恰好与 $q_{L}$ 的变换方式 $q_{L}\mapsto g\,q_{L}$ 相同。这就解释了为什么 (B.7) 是一个
\emph{不变}方程：方程两端属于同一个表示 $\rho_{(\frac12,0)}$。据此，把旋量认同为四元数后，
(B.7) 便是极其简洁的三四元数线性关系
\begin{equation}
Q_{V}\,q_{R}-m\,q_{L}=0 . \tag{B.8}
\end{equation}
这正是正文给出的“四元数形式的黎曼 Dirac 方程”。

\section{黎曼 Schwarzschild 解与径向距离}
\emph{补充《黎曼广义相对论》中 Schwarzschild--Kottler 度规的求解思路与 $\operatorname{dist}(r)$ 的积分。}

按球对称、静态的要求设度规
\begin{equation}
\dd s^{2}=A(r)\,\dd t^{2}+B(r)\,\dd r^{2}+r^{2}\bigl(\dd\theta^{2}+\sin^{2}\theta\,\dd\phi^{2}\bigr).
\end{equation}
在黎曼情形下，$t$ 方向与空间方向同号，故 $A,B>0$。把此度规代入真空 Einstein 方程
$G+\Lambda g=0$，得到关于 $A,B$ 的两个方程：其一是
\begin{equation}
\frac{1}{r^{2}}\left(1-\frac{1}{B}\right)-\frac{B'}{rB^{2}}+\Lambda=0,
\end{equation}
其解为
\begin{equation}
\frac{1}{B(r)}=1+\frac{2M}{r}-\frac{\Lambda r^{2}}{3}.
\end{equation}
其二给出 $A(r)=1/B(r)$，即
\begin{equation}
A(r)=1+\frac{2M}{r}-\frac{\Lambda r^{2}}{3}.
\end{equation}
这正是正文给出的黎曼 Schwarzschild--Kottler 度规（洛伦兹版仅 $2M/r$ 项前多一个负号）。

特别地，当 $\Lambda=0$ 时，$t,\theta,\phi$ 固定、只有 $r$ 变化的路径的固有长度为
\begin{equation}
\dd(\text{dist})=\sqrt{1+\frac{2M}{r}}\,\dd r ,
\end{equation}
于是从坐标中心到半径 $r$ 的距离为
\begin{equation}
\operatorname{dist}(r)=\int_{0}^{r}\frac{\dd r'}{\sqrt{1+2M/r'}} .
\tag{B.9}
\end{equation}
做代换 $r'=2M\sinh^{2}u$，则 $\dd r'=4M\sinh u\cosh u\,\dd u$，
$\sqrt{1+2M/r'}=\cosh u/\sinh u$，被积函数化为
\begin{equation}
\frac{\dd r'}{\sqrt{1+2M/r'}}=4M\sinh^{2}u\,\dd u=2M(\cosh 2u-1)\,\dd u .
\end{equation}
积分得
\begin{equation}
\operatorname{dist}(r)=M\sinh 2u-2M u=2M\sinh u\cosh u-2M u .
\end{equation}
由 $\sinh u=\sqrt{r/2M}$、$\cosh u=\sqrt{1+r/2M}$、$u=\operatorname{arcsinh}\sqrt{r/2M}$，
回代即得正文结果
\begin{equation}
\operatorname{dist}(r)=\sqrt{r(r+2M)}-2M\,\operatorname{arcsinh}\!\sqrt{\frac{r}{2M}} .
\end{equation}
由于 $\operatorname{dist}(r)<r$，半径 $r$ 的球面“看起来”比欧氏几何中更大，这正是负曲率
空间的标志，也是黎曼 Schwarzschild 解中不存在视界（故没有黑洞）的几何根源。

\section{黎曼黑体谱：态密度与积分}
\emph{补充《黎曼热力学》中态密度 $\rho(\omega)$、经典总能量密度与黑体的计算。}

在边长为 $L$、体积 $V=L^{3}$ 的立方腔中，周期性边界条件给出平面波波矢
\begin{equation}
\bk_{i,j,k}=\frac{2\pi}{L}(i\,\be_{x}+j\,\be_{y}+k\,\be_{z})+\omega_{i,j,k}\,\be_{t},
\qquad
\omega_{i,j,k}=\sqrt{\wms^{2}-\frac{4\pi^{2}}{L^{2}}(i^{2}+j^{2}+k^{2})} .
\end{equation}
色散关系给出 $\omega^{2}=\wms^{2}-\dfrac{4\pi^{2}}{L^{2}}r^{2}$，其中 $r^{2}=i^{2}+j^{2}+k^{2}$。
频率落在 $[\omega,\omega+\dd\omega]$ 内的模式数，等于（三种偏振乘以）球壳体积
$4\pi r^{2}|\dd r|$。由 $r^{2}=\dfrac{L^{2}}{4\pi^{2}}(\wms^{2}-\omega^{2})$ 对 $\omega$
微分，得
\begin{equation}
2r\,\dd r=-\frac{L^{2}}{2\pi^{2}}\,\omega\,\dd\omega
\quad\Longrightarrow\quad
r\,|\dd r|=\frac{L^{2}}{4\pi^{2}}\,\omega\,\dd\omega .
\end{equation}
于是
\begin{equation}
4\pi r^{2}|\dd r|=\frac{4\pi r^{2}}{r}\bigl(r\,|\dd r|\bigr)
=4\pi r\cdot\frac{L^{2}\omega}{4\pi^{2}}\,\dd\omega
=\frac{L^{2}r\,\omega}{\pi}\,\dd\omega .
\end{equation}
代入 $r=\dfrac{L}{2\pi}\sqrt{\wms^{2}-\omega^{2}}$、$V=L^{3}$，得到
\begin{equation}
4\pi r^{2}|\dd r| = \frac{L^{3}}{2\pi^{2}}\,\omega\sqrt{\wms^{2}-\omega^{2}}\,\dd\omega
= V\,\frac{\omega\sqrt{\wms^{2}-\omega^{2}}}{2\pi^{2}}\,\dd\omega .
\end{equation}
再乘偏振因子 $3$，即得正文的态密度
\begin{equation}
\rho(\omega)\,\dd\omega=3V\,\frac{\omega\sqrt{\wms^{2}-\omega^{2}}}{2\pi^{2}}\,\dd\omega .
\tag{B.10}
\end{equation}

\textbf{经典总能量密度。} 能均分定理给出每个模式平均能量 $kT$，故
\begin{equation}
u_{classical}(\omega)\,\dd\omega=3kT\,\frac{\omega\sqrt{\wms^{2}-\omega^{2}}}{2\pi^{2}}\,\dd\omega .
\end{equation}
对全频段积分。做代换 $\omega=\wms\sin\theta$：
\begin{equation}
\int_{0}^{\wms}\omega\sqrt{\wms^{2}-\omega^{2}}\,\dd\omega
=\wms^{3}\int_{0}^{\pi/2}\sin\theta\cos^{2}\theta\,\dd\theta
=\frac{\wms^{3}}{3},
\end{equation}
因此
\begin{equation}
u_{classical,\,total}=\frac{3kT}{2\pi^{2}}\cdot\frac{\wms^{3}}{3}=\frac{\wms^{3}kT}{2\pi^{2}},
\end{equation}
与正文一致：黎曼经典总能量密度按温度\emph{一次}标度，而非洛伦兹情形的 $T^{4}$。

\textbf{发射功率谱。} 腔内辐射并非各向同性地逃逸：单位时间、单位面积、给定取向的孔所
接收到的辐射，要乘速度 $v(\omega)=\sqrt{\wms^{2}-\omega^{2}}/\omega$，并对方向积分得因子
$1/4$。量子化后每模平均能量为 $\hbar\omega/(\ee^{\hbar\omega/kT}-1)$，与 (B.10) 相乘并
乘 $v(\omega)/4$，即得正文的发射谱
\begin{equation}
I_{quantum}(\omega)\,\dd\omega=\frac{3\hbar\omega\,(\wms^{2}-\omega^{2})}{8\pi^{2}\bigl(\ee^{\hbar\omega/kT}-1\bigr)}\,\dd\omega .
\end{equation}

\section{黎曼矢量波方程的三偏振}
\emph{补充《几何与波》中“横波条件等价于剔除标量波”的论证。}

考虑矢量平面波 $\bA(x)=\bA_{0}\sin(\bk\cdot\bx)$，其中 $\bk\cdot\bx=k_x x+k_y y+k_z z+\omega t$。
只要 $\bk\cdot\bk=\wms^{2}$，$\bA$ 的每个分量都满足黎曼标量波方程（RSW）。然而，可以由此
构造一个所有观测者都认同的\emph{标量}波
\begin{equation}
D(x)=\bk\cdot\bA(x)=(\bk\cdot\bA_{0})\sin(\bk\cdot\bx).
\end{equation}
为使 $\bA$ 表示纯矢量波（对应自旋 $1$ 的量子，如黎曼光子）而不含标量成分（自旋 $0$），
须令
\begin{equation}
\bk\cdot\bA_{0}=0 .
\tag{B.11}
\end{equation}
这一条件把 $\bA_{0}$ 限制在垂直于传播矢量 $\bk$ 的三维子空间内，故黎曼光有\emph{三}个
独立偏振（对比我们宇宙的光只有两个）。

把 (B.11) 写成对任意场 $\bA(x)$ 都适用的局域形式。对上述平面波，
\begin{align}
\partial_{x}A^{x}+\partial_{y}A^{y}+\partial_{z}A^{z}+\partial_{t}A^{t}
&=(A_{0}^{x}k^{x}+A_{0}^{y}k^{y}+A_{0}^{z}k^{z}+A_{0}^{t}k^{t})\cos(\bk\cdot\bx)\nonumber\\
&=(\bk\cdot\bA_{0})\cos(\bk\cdot\bx)=0,
\end{align}
即\emph{横波条件}
\begin{equation}
\partial_{x}A^{x}+\partial_{y}A^{y}+\partial_{z}A^{z}+\partial_{t}A^{t}=0 .
\tag{B.12}
\end{equation}
(B.12) 与 RSW 一起构成正文的黎曼矢量波方程（RVW）。三个偏振在量子化后即黎曼光子的三个
自旋态，这一点在《黎曼量子力学》与《黎曼热力学》的态计数中反复使用。

\chapter{两个宇宙对照}

本附录以表格形式，汇总正交宇宙（黎曼型时空）与我们自身宇宙（洛伦兹型时空）在物理各层面的
系统性差异。表中每一行都可在正文相应章节找到详细讨论。为便于对照，我们把“洛伦兹宇宙”列
置于左，“黎曼宇宙”列置于右。

\begin{longtable}{p{2.9cm} p{5.7cm} p{5.9cm}}
\toprule
\textbf{项目} & \textbf{洛伦兹宇宙（我们的）} & \textbf{黎曼宇宙（正交）} \\
\midrule
\endfirsthead
\toprule
\textbf{项目} & \textbf{洛伦兹宇宙（我们的）} & \textbf{黎曼宇宙（正交）} \\
\midrule
\endhead
\bottomrule
\endlastfoot

时空度规 &
$\dd s^{2}=-\dd t^{2}+\dd x^{2}+\dd y^{2}+\dd z^{2}$，号差 $(-,+,+,+)$ &
$\dd s^{2}=\dd t^{2}+\dd x^{2}+\dd y^{2}+\dd z^{2}$，号差 $(+,+,+,+)$ \\

坐标地位 &
时间方向特殊；光锥把类时与类空分开 &
四个维度完全平等；无光锥、无类光矢量 \\

光速 &
存在普适常量 $c$，一切观察者一致 &
\textbf{无普适光速}；相速依赖波长，$\omega^{2}+|\mathbf{k}|^{2}=\wms^{2}$ \\

速度上限 &
任何速度不超过 $c$ &
\textbf{无速度上限} \\

能量--动量关系 &
$E^{2}-p^{2}=m^{2}$ &
$E^{2}+p^{2}=m^{2}$；(B.1) 节 \\

光子静止质量 &
零 &
\textbf{非零}；由 Proca 型方程描述 \\

静电场（库仑势） &
$1/r$，无质量光子 &
随空间\emph{振荡}（质量项使其不指数衰减） \\

光的偏振态 &
两个（自旋 $\pm1$） &
\textbf{三个}（垂直于波矢的三维子空间） \\

发光过程 &
产生光需\emph{消耗}能量 &
\textbf{产生光释放能量}（放热） \\

高速远行的时间 &
旅行者时间膨胀，比家乡人更\emph{年轻} &
旅行者经历更\emph{长}的岁月（绕行使时间更长） \\

引力：Schwarzschild 解 &
有视界，$r=2M$ 处出现黑洞 &
径向距离有限，\textbf{无视界、无黑洞} \\

引力波 &
存在（四维时空中的横波） &
\textbf{不存在}（黎曼度规下无相应自由度） \\

宇宙学 &
大爆炸、膨胀宇宙为常见图景 &
\textbf{很可能没有大爆炸}；均匀各向同性解另有形态 \\

热力学：温度 &
温度无上界（上限趋于“无穷热”） &
存在\emph{负温度}，可“比无穷热更热” \\

黑体谱 &
普朗克谱，总功率按 $T^{4}$ 标度 &
谱被截止于 $\wms$；经典总能量按 $T$ 线性标度，经典/量子差异很小 \\

生命与光合作用 &
植物\emph{吸收}光能制造养分 &
植物向夜空\emph{发出}自己的光来制造养分 \\

星际航行 &
需携带燃料/推进剂 &
\textbf{发光引擎}可无燃料、无推进剂地持续加速 \\

量子力学 &
Dirac 方程、Klein--Gordon 方程（洛伦兹号差） &
黎曼 Dirac 方程；旋量由 $\mathrm{SU}(2)\times\mathrm{SU}(2)$ 双覆盖描述 \\

\end{longtable}

\section*{结语}

正交宇宙与我们的宇宙之别，归根到底只源于时空度规中时间项的一个符号。然而这一个符号的翻转，
却连锁地改写了光学、电磁学、热力学、引力与量子力学的面貌：光失去了普适速度、获得了静止质量；
发光由吸能变为释能；黑洞、引力波与大爆炸都不复存在；而“向黑夜发光”的植物与“无燃料的引擎”
则重塑了整个文明的生存方式。Egan 的《正交》三部曲，正是一次以如此单一的几何改动为起点、
把物理推论一路推到底的系统性思想实验。


\end{document}
